\documentclass[a4paper,10pt]{article}
\usepackage[top=25mm,bottom=24mm,left=27mm,right=27mm,headheight=15pt,headsep=8mm,footskip=12mm]{geometry}
\usepackage{fontspec,amsmath,amsthm,unicode-math}
\usepackage[spanish,es-nodecimaldot,es-noquoting]{babel}
\usepackage{microtype,xcolor,fancyhdr,titlesec,booktabs,longtable,array,tabularx}
\definecolor{accent}{HTML}{184C63}
\definecolor{muted}{HTML}{52616B}
\definecolor{ink}{HTML}{172A36}
\usepackage{graphicx,tikz,enumitem,float,needspace,fvextra}
\usetikzlibrary{arrows.meta,calc,positioning,decorations.pathreplacing,matrix}
\usepackage[unicode,hidelinks]{hyperref}
\usepackage[nameinlink,noabbrev]{cleveref}
\setmainfont{STIXTwoText-Regular.otf}[BoldFont=STIXTwoText-Bold.otf,ItalicFont=STIXTwoText-Italic.otf,BoldItalicFont=STIXTwoText-BoldItalic.otf]
\setmathfont{STIXTwoMath-Regular.otf}
\hypersetup{pdftitle={Relaciones estructurales entre las constantes físicas},pdfauthor={Oumar Haidara Fall; Rubén Ramos Balsa},pdfsubject={Acción, reciprocidad geométrica y memoria de las transformaciones}}
\newcommand{\RR}{\mathbb R}\newcommand{\ZZ}{\mathbb Z}\newcommand{\QQ}{\mathbb Q}\newcommand{\CC}{\mathbb C}\newcommand{\FF}{\mathbb F}
\newcommand{\Id}{\operatorname{id}}\newcommand{\tr}{\operatorname{tr}}\newcommand{\Tr}{\operatorname{Tr}}
\newcommand{\Spec}{\operatorname{Spec}}
\newcommand{\Span}{\operatorname{span}}\newcommand{\Cyl}{\operatorname{Cyl}}\newcommand{\dd}{\mathrm d}
\newcommand{\square}{\mdlgwhtsquare}
\newcommand{\TPK}{\mathrm{TPK}}\newcommand{\APP}{\mathrm{APP}}\newcommand{\TRIT}{\{-1,0,+1\}}\newcommand{\HMT}{\mathrm{HMT}}
\newtheorem{theorem}{Teorema}[section]
\newtheorem{lemma}[theorem]{Lema}\newtheorem{proposition}[theorem]{Proposición}\newtheorem{corollary}[theorem]{Corolario}
\theoremstyle{definition}\newtheorem{definition}[theorem]{Definición}\newtheorem{axiom}[theorem]{Axioma}
\theoremstyle{remark}\newtheorem{remark}[theorem]{Observación}
\numberwithin{equation}{section}
\setlength{\parindent}{1em}\setlength{\parskip}{3pt plus 1pt minus .5pt}
\setlength{\emergencystretch}{3em}
\clubpenalty=10000\widowpenalty=10000\displaywidowpenalty=10000
\renewcommand{\normalsize}{\fontsize{10.5}{13.4}\selectfont}\normalsize
\setlist{nosep,topsep=5pt}
\titleformat{\section}{\fontsize{13}{16}\selectfont\bfseries}{\thesection}{.65em}{}
\titleformat{\subsection}{\fontsize{11.2}{14}\selectfont\bfseries}{\thesubsection}{.65em}{}
\titleformat{\subsubsection}{\normalsize\bfseries}{\thesubsubsection}{.65em}{}
\AddToHook{cmd/section/before}{\Needspace{10\baselineskip}}
\AddToHook{cmd/subsection/before}{\Needspace{7\baselineskip}}
\AddToHook{cmd/subsubsection/before}{\Needspace{6\baselineskip}}
\setcounter{tocdepth}{2}
\makeatletter
\renewcommand*\l@subsection{\@dottedtocline{2}{1.5em}{3.9em}}
\makeatother
\titlespacing*{\section}{0pt}{17pt}{9pt}\titlespacing*{\subsection}{0pt}{13pt}{6pt}
\pagestyle{fancy}\fancyhf{}
\fancyhead[L]{\fontsize{8}{10}\selectfont Relaciones estructurales entre las constantes físicas}
\fancyfoot[C]{\thepage}\renewcommand{\headrulewidth}{.2pt}\raggedbottom
\makeatletter
\@ifundefined{example}{\theoremstyle{definition}\newtheorem{example}[theorem]{Ejemplo}}{}
\@ifundefined{notatabla}{\newcommand{\notatabla}[1]{\par\smallskip{\footnotesize #1\par}}}{}
\makeatother
\providecommand{\FF}{\mathbb F}

% Apertura independiente del índice: corrección quirúrgica 20260928.
\let\HMTIndiceOriginal\tableofcontents
\renewcommand{\tableofcontents}{\clearpage\begingroup\setlength{\parskip}{0pt}\fontsize{9.2}{10.5}\selectfont\HMTIndiceOriginal\endgroup\clearpage}
\begin{document}
\thispagestyle{plain}
\begin{center}
{\fontsize{10}{12}\selectfont Manuscrito de recopilación para transferencia académica\par}
\vspace{1.5mm}
{\fontsize{17}{20.5}\selectfont\bfseries Relaciones estructurales\par entre las constantes físicas\par}
\vspace{1.5mm}
{\fontsize{11}{13.5}\selectfont Acción, reciprocidad geométrica\par y memoria de las transformaciones\par}
\vspace{1.5mm}
{\fontsize{11.5}{14}\selectfont Oumar Haidara Fall \textperiodcentered\ Rubén Ramos Balsa\par}
\vspace{1.5mm}
{\fontsize{8}{10}\selectfont Coautoría sin prelación de contribución\par}
\vspace{1.5mm}
{\fontsize{8}{10}\selectfont Integración y recopilación documental generadas por ChatGPT - Astra.\par}
\end{center}
\vspace{1.5mm}
\begingroup\fontsize{9.5}{10.7}\selectfont\setlength{\parskip}{1.5pt}\renewenvironment{abstract}{\begin{center}\bfseries\abstractname\end{center}\noindent}{\par}\begin{abstract}
Se establecen relaciones de compatibilidad, reciprocidad y reconstrucción
entre las estructuras de acción, masa y respuesta constitutiva obtenidas
mediante Holografía Modular Triádica y Mecánica Dimensional. El desarrollo
parte de la composición de la estructura aritmética de caminos
\emph{All Possible Paths} (APP), la firma ternaria orientada TRIT y el
\emph{Topological Phase Kernel} (TPK), que actualiza y transporta el estado
enriquecido. Esta composición conserva orientación, acarreo, memoria y
retorno. Las constantes intervienen como evaluaciones de las
estructuras previamente construidas; su composición permite determinar qué
propiedades comparten y qué información recupera una observación conjunta.
En la colección constitutiva normalizada se demuestra que la velocidad de
propagación y el funcional de Barbero--Immirzi determinan los autovalores
etiquetados de los dos canales. Los logaritmos de velocidad e impedancia
constituyen el gradiente de un potencial espectral estrictamente cóncavo.
Su Hessiano caracteriza la sensibilidad de la inversión y proporciona cotas
explícitas de estabilidad. El potencial comparte un momento espectral de
profundidad dos con la construcción orientada cuyo valor extremo es Catalán;
la respuesta constitutiva restituye esa colección de manera única con su
condición de frontera en el origen.

La composición de cuartos de giro asociados a formas elípticas conjugadas
produce un transporte relativo unimodular cuyas dilataciones recíprocas
permiten reconstruir el coeficiente adimensional de acción, conservadas la
coordenada de clausura y la orientación. El levantamiento dodecafásico y la
iteración del transporte caracterizan la memoria del orden de las
operaciones. El análisis cronológico distingue el crecimiento de la
información de la tasa entrópica por posición, y la reentrada de memoria
determina su contribución a la respuesta operatoria. El retorno inscrito y
su composición polar determinan una longitud $L_\alpha$. Conservar
conjuntamente acción y área fija $G=c^3L_\alpha^2/\hbar$ en la realización
radial, también bajo variaciones no conmutativas y reducción compatible
de memoria.

La memoria electrónica se restituye mediante invariantes lineales y
espectrales que conservan el orden del registro; el transporte entre
multisecciones admite un criterio explícito de conservación de masa.
La composición térmica reúne capacidades, multiplicidades y memoria
espectral en una traza exacta, y construye su respuesta termométrica
para una lectura marcada con referencia fija. Esta realización enlaza
energía, información, acción y área, conservando las condiciones de
normalización y de observación de cada magnitud.
La composición transportada de las respuestas y la reconstrucción
Higgs--impedancia enlazan los canales constitutivos con el sector
electrodébil. La sustitución electrónica conserva su escala en Fermi,
la cancela en los cocientes angulares y transporta la longitud inversa
mediante una congruencia operatoria, también sin conmutación.

Finalmente, bajo las
condiciones circular y térmica especificadas, una misma coordenada de masa
determina la frecuencia propia, dos longitudes recíprocas de producto
invariante y la ponderación térmica del ciclo. Los resultados reúnen así
producción, compatibilidad y restitución en una descripción matemática
común, con demostraciones explícitas y con la distinción entre cambio de
representación, transporte relativo y realización física.


El lector dodecafásico uniforme determina además una sucesión de primeras raíces de retorno crítico. Se demuestra que su cociente de separaciones converge a la constante de Feigenbaum, conservando la selección original, la orientación y la memoria de la construcción.

% Adición al resumen de VII; no sustituye su contenido anterior.
La misma composición se prolonga a una acción total de geometría,
calibre y materia: incorpora conjuntamente la gravitación y las
interacciones fuerte, débil y electromagnética, conservando las
representaciones quirales y los acoplamientos previamente generados.
Se obtienen su corriente de espín, la fuente métrica total y los
términos torsionales cruzados entre especies. La transformación de
Legendre y las identidades de simetría producen restricciones
conjuntas de primera clase y su propagación. Sobre su superficie
regular se calcula la curvatura de la conexión canónica total y se
demuestra su anulación en las direcciones de deformación
características. El resultado conserva las condiciones de borde y
la posible holonomía global del estado enriquecido.



\par\smallskip\noindent\textbf{Palabras clave:} acción; reciprocidad geométrica; memoria operatoria; constantes físicas; Holografía Modular Triádica; Mecánica Dimensional.
\end{abstract}


\endgroup

\clearpage
\begingroup
\setlength{\parskip}{0pt}
\fontsize{10}{11.5}\selectfont
\tableofcontents
\endgroup
\clearpage
\clearpage
\begingroup\setlength{\parskip}{2pt}\fontsize{10}{11.8}\selectfont
\section{Introducción}

\label{md:introduccion}
La lectura de una constante como valor escalar representa el término final
de una operación. Cuando se conserva la estructura que produce esa lectura,
se hacen accesibles otras preguntas: qué relaciones comparte con las demás
lecturas, qué transformación conserva su orientación y qué parte de la
historia puede restituirse a partir de observaciones conjuntas. El presente
desarrollo sitúa esas preguntas en la continuidad entre la estructura discreta
HMT y sus realizaciones en Mecánica Dimensional.

El origen común conserva las operaciones de APP, el régimen orientado del
TRIT y la composición efectiva del TPK. El retorno de fase se distingue de
la evolución del estado enriquecido. Las coordenadas regionales y su
clausura dodecafásica se transmiten después a los lectores de acción,
incidencia y masa. Cada transmisión conserva los objetos y las condiciones
que utiliza; una salida puede intervenir en la construcción siguiente sin
perder su procedencia. El núcleo inicial expone esa arquitectura con sus
operaciones, prolongaciones y pruebas. Los antecedentes sectoriales que
siguen materializan las dependencias empleadas por las composiciones del
presente trabajo.

La cuestión central puede formularse en tres términos complementarios:
\emph{producción}, \emph{compatibilidad} y \emph{restitución}. La
producción fija el objeto y su genealogía. La compatibilidad determina qué
relaciones deben satisfacer sus distintas lecturas. La restitución establece
qué datos del objeto pueden recuperarse desde ellas. El recorrido inverso
es una propiedad de las lecturas ya construidas y conserva sus bases,
orientaciones y dominios; no reemplaza su orden generador.

El primer eje del manuscrito reúne la acción con el par angular. El
coeficiente adimensional de la sección de acción participa en el
contraángulo; este determina, con el ángulo, los canales que alimentan la
respuesta constitutiva y el funcional de Barbero--Immirzi. La misma sección
interviene en el cociente de retorno de la lectura electrónica. De esta
forma, sustituir las estructuras en sus ecuaciones revela dependencias
compartidas que una enumeración de constantes separadas deja implícitas.
La distinción entre coeficiente y magnitud dimensional es esencial:
restituir una razón de forma y restituir una acción dimensional requieren
datos de referencia diferentes.

El segundo eje estudia la información conservada en la respuesta. La
velocidad y la impedancia normalizadas retienen aspectos distintos del
reparto entre canales. La composición permite demostrar tanto una inversión
global en la colección declarada como una relación diferencial entre ambas
respuestas. Un potencial espectral organiza esa relación, y su Hessiano
controla la sensibilidad de la inversión. Así se distinguen dos propiedades:
que los datos exactos determinen unívocamente el objeto y que una precisión
finita permita reconstruirlo con un error acotado. La colección orientada cuyo
valor extremo es Catalán se incorpora mediante su operación espectral,
conservando sus argumentos y su profundidad.

La composición transportada de las respuestas proporciona su coordenada
aditiva; las firmas angulares de calibre y de Higgs permiten reconstruir
la misma pareja de canales. Los momentos del defecto cíclico completan
esa lectura con cotas de convergencia y una representación de Mellin.
Más adelante, la sección electrónica se sustituye en Fermi y en la escala
de Higgs: la prueba conserva qué factores permanecen, cuáles se cancelan
y cómo se transforma el inverso del operador de masa sin suponer conmutación.

El tercer eje traslada la discusión al orden de las operaciones. El plano
de cuarto de giro del ciclo dodecafásico y las formas conjugadas de la
elipse permiten construir un conmutador explícito. El balance de los
incrementos de fase se cierra mientras la acción sobre el estado conserva
una dilatación recíproca. Dos programas con las mismas multiplicidades de
operaciones pueden producir respuestas diferentes por su orden. Esa
diferencia admite detectores y recuperadores definidos: la traza registra
ciertos invariantes, mientras las respuestas sobre ejes orientados
distinguen sentidos que la traza identifica.

La lectura física se presenta junto con las condiciones de cada
realización. En la sección gravitatoria circular, la longitud reducida
de Compton y el radio gravitatorio reducido se dilatan recíprocamente y
mantienen un área común. En la realización térmica del ciclo, la misma
coordenada estructural interviene como exponente de un peso de equilibrio.
En el transporte elíptico, la conservación de área simpléctica y la
conservación de una energía cuadrática común se examinan mediante sus
operadores respectivos. La eventual realización experimental del programa
relativo exige especificar conjuntamente sistema, referencia, operaciones
disponibles y controlador. Estas condiciones pertenecen al significado
del resultado y permanecen contiguas a sus pruebas.

La normalización gravitatoria se desarrolla antes de utilizar sus
consecuencias escalares. El retorno de Hadamard, la inscripción de incidencias
y la norma local de $\alpha$ producen una sección de longitud. Su
transporte polar y el mismo carácter de ruta producen dos lectores
recíprocos. La energía conserva a su vez la longitud de acción. Al reunir
$R\Lambda=L_\alpha^2I$ con $H\Lambda=\hbar cI$, la constante gravitatoria
queda fijada como coeficiente común de todos los modos positivos de la
realización. El reloj circular se expresa después desde longitud y
velocidad. Esta precedencia permite explicar la normalización sin utilizar
el valor experimental de $G$ ni una longitud de Planck tabulada.

La composición electrónica y la térmica precisan cómo interviene la
memoria en esas relaciones. En la primera, se reconstruye el registro
dodecafásico a partir de su suma, sus diferencias y su orientación;
el espectro distingue órdenes que comparten multiplicidades. En la
segunda, el espectro geométrico ya construido por el transporte de
memoria aporta la razón $1/10$. Las capacidades cronológicas y los
sectores de ocupación determinan una traza marcada que se compone
con un estado de Gibbs. Se especifican por separado la energía
condicionada, la información marcada y la temperatura resultante;
esta elección de observaciones permite demostrar el transductor
termométrico y seguir su relación con la escala de acción y el área.
Las secciones correspondientes conservan los datos de preparación,
las graduaciones y las hipótesis de variación utilizadas.

La exposición desarrolla primero los antecedentes, conserva después la
narrativa explicativa del recorrido y presenta las composiciones completas.
El estudio cronológico de las vacancias, el episodio regional de reiteración
visible y la reentrada de memoria permiten seguir el origen de las
preguntas sobre compatibilidad. Los desarrollos energéticos y gaussianos
se conservan en los apéndices como realizaciones auxiliares, con sus
hipótesis y pruebas completas.
Las conclusiones reúnen las caracterizaciones estructurales obtenidas.
El suplemento digital preserva las fuentes de trabajo literales, los
programas y sus resultados; el manuscrito contiene las demostraciones
empleadas, y distingue su lectura matemática de los controles finitos que
las acompañan.

Esta organización permite formular con precisión la aportación transversal
del trabajo. Una constante conserva su valor numérico, pero su definición
estructural especifica además la operación que lo produce, el dominio en
que actúa y las transformaciones que conserva. Dos magnitudes que comparten
una dependencia pueden, por ello, restringirse mutuamente: conocer una
respuesta y su orientación permite reconstruir un reparto que una lectura
escalar aislada no distingue. El problema inverso se plantea sobre la imagen
ya generada y conserva sus datos de referencia. Así, generación y
reconstrucción se articulan sin invertir la precedencia de la construcción.

La misma distinción ordena la relación entre memoria e información. El
estado presente, el balance acumulado y el operador que responderá a una
continuación son objetos diferentes. El retorno de una fase puede coexistir
con una transformación persistente del estado; dos sucesiones con idénticos
recuentos pueden diferenciarse por el orden de composición. Las pruebas
desarrollan esa diferencia mediante operadores, invariantes y observaciones
recuperadoras. Su interés físico reside en precisar qué conserva cada
descripción y qué operación permite acceder a la información retenida.

El paso a las magnitudes dimensionales se realiza manteniendo las bases y
normalizaciones correspondientes. La reconstrucción de un coeficiente de
acción, la conservación de un área y la caracterización de una temperatura
pertenecen a resultados relacionados, con tipos distintos. La exposición
reúne esos resultados en sus condiciones de validez y muestra cómo se
componen. Las conclusiones explicitan tanto la caracterización obtenida
para cada objeto como la relación que establece con las demás estructuras.

La constante de Feigenbaum adquiere una definición estructural precisa: es el invariante asintótico de separación de los parámetros seleccionados por los retornos críticos del lector dodecafásico uniforme. La construcción conserva la procedencia APP--TRIT--TPK, el continuo conjunto y el orden de primeras raíces. La identificación de los centros seleccionados con los centros analíticos permite aplicar el escalado a esa misma sucesión; las pruebas de \ref{hmtfg:escalado:15}--\ref{hmtfg:escalado:17} y \eqref{hmtfg:articulacion:limite} establecen el resultado paramétrico. La representación de Jacobi restituye el perfil desde su medida de doce nodos, mientras la inversión constitutiva recupera los canales orientados y el momento cuyo extremo es Catalán. Estas operaciones muestran relaciones entre lectores del mismo antecedente, conservando el dominio y la información propios de cada una.

La composición se prolonga también a una acción común de geometría, calibre y materia. La gravitación y las interacciones fuerte, débil y electromagnética intervienen mediante las realizaciones del mismo estado enriquecido, con los acoplamientos recibidos de la construcción anterior. La variación conjunta determina las fuentes y la respuesta torsional; la transformación de Legendre produce un álgebra común de restricciones de primera clase. Su propagación y el cierre de la conexión canónica se demuestran sobre la superficie regular de restricciones, manteniendo las condiciones de borde y la holonomía global. Así, la unificación variacional y canónica aporta una composición dinámica explícita a las relaciones de compatibilidad y restitución que articulan el manuscrito.

\par\endgroup
\clearpage
% Núcleo común de I conservado literalmente; registro ampliado por inclusión.
\begingroup
\makeatletter
\def\input@path{{base_articulo_I/}}
\makeatother
\section{Núcleo formal de la Holografía Modular Triádica}
\label{sec:nucleo}

La Holografía Modular Triádica, abreviada HMT, comienza con una estructura finita de posiciones y operaciones. El orden de construcción importa: las posiciones reciben evaluaciones aritméticas, las transiciones adquieren una firma ternaria y la dinámica conserva la información necesaria para prolongarlas. Sólo después se forman las realizaciones numéricas. Una constante reconocida al final de este proceso no sustituye a la regla que produjo su región ni a la historia que determina su evaluación.

La denominación \emph{All Possible Paths}, APP, expresa el papel de los caminos sobre el soporte. Se adopta como denominación de la construcción, con una referencia conceptual a la formulación por caminos de la mecánica cuántica; no supone que su ley aritmética sea la integral de caminos de Feynman. El \emph{Topological Phase Kernel}, TPK,\footnote{Los autores dedican asimismo las iniciales TPK a Tan Pengkiang. Esta dedicación es independiente de la definición matemática.} designa la composición de selección, transporte, actualización de memoria y prolongación. Las siglas se mantienen únicamente después de haber definido los objetos que representan.

\subsection{APP: soporte posicional y evaluaciones aritméticas}

Consideremos las nueve marcas interiores de una recta graduada entre cero y diez,
\[
 D_9=\{1,2,3,4,5,6,7,8,9\}.
\]
Su disposición horizontal y vertical produce las ochenta y una posiciones de \(D_9^2\). Las marcas interiores no deben confundirse con los diez intervalos unitarios de la recta de referencia. Esta disposición fija un alfabeto y dos direcciones; todavía no utiliza una expansión real para decidir una trayectoria.

La identificación cíclica de las posiciones conduce al soporte
\[
 X_9=(\ZZ/9\ZZ)^2.
\]
La coordenatización de representantes positivos conserva el símbolo nueve para la clase nula. Para todo entero positivo se definen
\begin{equation}
 \rho_9(n)=1+((n-1)\bmod9),\qquad
 q_9(n)=\frac{n-\rho_9(n)}9,
 \qquad n=\rho_9(n)+9q_9(n).
 \label{eq:residuo-cociente}
\end{equation}
La raíz digital positiva es, por tanto, una elección de representante de la clase residual. No conserva por sí sola el entero anterior a la reducción; la coordenada \(q_9\) proporciona precisamente la información que falta.

En cada posición \((i,j)\), la acumulación de \(i\) repetida \(j\) veces produce \(ij\). Las dos evaluaciones enteras y sus proyecciones son
\[
 S(i,j)=i+j,\qquad P(i,j)=ij,\qquad
 \Sigma(i,j)=\rho_9(i+j),\quad \Pi(i,j)=\rho_9(ij).
\]
La simetría \((i,j)\mapsto(j,i)\) intercambia las direcciones de acumulación y conserva ambas evaluaciones. Es una propiedad de un único soporte, no una identificación posterior entre dos matrices independientes.

\begin{definition}[Evaluación aritmética levantada]
La evaluación de APP que conserva las dos hojas es
\begin{equation}
 \mathcal A(i,j)=
 \bigl((R^+_{ij},q^+_{ij}),(R^\times_{ij},q^\times_{ij})\bigr)
 =\bigl((\rho_9(i+j),q_9(i+j)),(\rho_9(ij),q_9(ij))\bigr).
 \label{eq:app-levantada}
\end{equation}
El término \emph{hoja} indica aquí una evaluación sobre el mismo soporte: aditiva o multiplicativa. Las dos coordenadas de cociente pertenecen al estado aritmético y no se descartan al reducir los residuos.
\end{definition}

\input{figures/app.tex}

\begin{proposition}[Reconstrucción de las evaluaciones]
Los dos pares de \eqref{eq:app-levantada} determinan exactamente \(i+j\) e \(ij\). Los residuos aislados no determinan esas evaluaciones enteras ni sus cocientes.
\end{proposition}
\begin{proof}
La primera afirmación es \eqref{eq:residuo-cociente} aplicada a cada hoja. Para la segunda, \(10=1+9\) y \(19=1+2\cdot9\) tienen el mismo residuo y distinto cociente. En el soporte bidimensional, las posiciones \((5,5)\) y \((8,2)\) dan la misma suma y el mismo residuo producto, pero productos \(25=7+18\) y \(16=7+9\). La proyección residual identifica datos que la evaluación levantada distingue.
\end{proof}

\subsubsection{Grafo de Cayley, topología toroidal y grado de enrollamiento}

Las traslaciones cardinales están dadas, en coordenadas de fila y columna, por
\[
 \nu_N=(-1,0),\quad\nu_E=(0,1),\quad
 \nu_S=(1,0),\quad\nu_O=(0,-1),\qquad T_d(x)=x+\nu_d.
\]
La primera coordenada aumenta hacia el sur y la segunda hacia el este. El grafo
\[
 \mathscr G_9=\operatorname{Cay}(X_9,\{\nu_N,\nu_E,\nu_S,\nu_O\})
\]
es la retícula cíclica bidimensional: tiene ochenta y un vértices y cuatro aristas orientadas salientes por vértice. Es un grafo de Cayley para la estructura \emph{aditiva} del soporte. La multiplicación de residuos tiene divisores de cero y no convierte todas las filas de APP en un grupo multiplicativo.

Una palabra cardinal \(w=d_1\cdots d_r\) y un origen \(x_0\) determinan
\[
 x_k=x_0+\sum_{j=1}^k\nu_{d_j}\pmod9.
\]
Existen \(81\cdot4^r\) caminos orientados de longitud \(r\), incluidos retornos y retrocesos. Cuando el camino se cierra, su desplazamiento entero previo a la reducción define
\begin{equation}
 \operatorname{wind}(w)=\frac19
 \bigl(\#S(w)-\#N(w),\#E(w)-\#O(w)\bigr)\in\ZZ^2.
\end{equation}
La inversión del camino cambia el signo de este vector. El retorno a la misma posición residual no obliga, por tanto, a que el enrollamiento sea nulo. Éste es el primer ejemplo de una diferencia que reaparecerá en el retorno de fase: una coordenada puede cerrarse mientras otra registra una evolución no trivial.

\subsubsection{Cociente de transporte e identidad de cociclo}

Sea \(s_9:\ZZ/9\ZZ\to D_9\) la sección de representantes positivos. El cociente de transporte asociado a la composición aditiva, denominado también \emph{acarreo} en aritmética posicional, es
\[
 c_9(a,b)=\frac{s_9(a)+s_9(b)-s_9(a+b)}9.
\]
El carácter composicional de esta memoria queda expresado por una identidad exacta.

\begin{lemma}[Identidad de cociclo]
Para \(a,b,c\in\ZZ/9\ZZ\),
\[
 c_9(a,b)+c_9(a+b,c)=c_9(b,c)+c_9(a,b+c).
\]
\end{lemma}
\begin{proof}
Ambas sumas son iguales a
\[
 \frac{s_9(a)+s_9(b)+s_9(c)-s_9(a+b+c)}9.
\]
La asociatividad de la suma de representantes levantados requiere precisamente este término adicional al componer en el cociente.
\end{proof}

Con representantes positivos, este cociclo no está normalizado en la identidad: \(c_9(0,a)=1\). Si \(s_0\) toma representantes en \(\{0,\ldots,8\}\), su acarreo normalizado \(c_0\) satisface
\[
 c_9(a,b)=c_0(a,b)+\mathbf1_{a=0}+\mathbf1_{b=0}-\mathbf1_{a+b=0}.
\]
La diferencia es el cambio de sección; no altera la identidad de cociclo. La descomposición posterior del calendario en fase y ventana utilizará los representantes de cero a ocho.

Los totales de ambas evaluaciones distinguen la contribución visible y la conservada en el cociente:
\[
 \sum S=810,\quad\sum P=2025,\quad
 \sum\Sigma=405,\quad\sum\Pi=459,
\]
\[
 \sum q^+=45,\qquad\sum q^\times=174.
\]
De aquí se obtiene
\begin{equation}
 2025-810=(459-405)+9(174-45)=54+9\cdot129.
 \label{eq:defecto-integrado}
\end{equation}
La diferencia suma--producto no se agota en el defecto visible \(54\). La identidad conserva también el defecto de cociente \(129\). Suprimirlo cambiaría la aritmética que se transporta, aunque la tabla residual permaneciera idéntica.

\subsubsection{Reducción ternaria y radical multiplicativo}

El anillo \(R=\ZZ/9\ZZ\) tiene grupo de unidades \(\{1,2,4,5,7,8\}\) e ideal máximo
\[
 \mathfrak m=(3)=\{0,3,6\},\qquad\mathfrak m^2=0.
\]
Toda fila aditiva es una permutación de los nueve representantes. Una fila multiplicativa lo es exactamente cuando su índice es una unidad. Las filas tres y seis contienen tres valores repetidos; la fila nueve contiene únicamente el representante de la clase nula. Por ello el plegamiento multiplicativo ya distingue, antes de cualquier geometría continua, las unidades del sector nilpotente.

La aplicación \(x\mapsto3x\) tiene núcleo e imagen iguales a \(\mathfrak m\). Induce un isomorfismo de espacios vectoriales sobre \(\FF_3\),
\[
 R/\mathfrak m\longrightarrow\mathfrak m,\qquad [x]\longmapsto3x,
\]
pero no un isomorfismo de anillos. En particular, la secuencia visible \(3,6,9\) admite dos lecturas diferentes: su reducción directa módulo tres es \(0,0,0\), mientras que la extracción del factor tres seguida de la reducción de su coordenada interna da \(1,2,0\). Esta segunda operación conserva una coordenada del ideal; no es la misma proyección que reducir directamente el número original.

\subsubsection{Realización latina de la evaluación aditiva}

La diferencia entre las dos evaluaciones puede expresarse también mediante
operadores. Sea \((e_r)_{r\in\ZZ/9\ZZ}\) la base ortonormal de \(\CC^9\).
La asignación \(v_{ab}=e_{a+b}\) produce una base ortonormal en cada fila
y columna. Los proyectores \(P_{ab}=v_{ab}v_{ab}^{*}\) satisfacen
\[
 \sum_bP_{ab}=I_9=\sum_aP_{ab},\qquad P_{ab}^2=P_{ab}.
\]
En efecto, cada traslación \(b\mapsto a+b\) permuta los nueve residuos.
Se obtiene una realización por proyectores de rango uno de un cuadrado
latino, en la terminología de los cuadrados latinos cuánticos
\cite{mustovicary,delascuevas2023}. Esta construcción parte de la hoja
aditiva ya definida y explicita el mapa hacia su representación operatoria.

La evaluación multiplicativa conserva otro contenido: para \(a=3\), los
vectores \(e_{3b}\) repiten los residuos \(0,3,6\), y
\[
 \sum_{b\in\ZZ/9\ZZ}|e_{3b}\rangle\langle e_{3b}|
 =3\bigl(|e_0\rangle\langle e_0|+|e_3\rangle\langle e_3|
                    +|e_6\rangle\langle e_6|\bigr).
\]
La multiplicidad queda registrada en el operador y caracteriza el radical
multiplicativo. Por ello el soporte toroidal, la propiedad latina y la
resolución operatoria de la identidad son propiedades que requieren
definiciones diferentes. Los cuadrados mágicos cuánticos estudiados por
De las Cuevas, Drescher y Netzer exigen entradas positivas y sumas de filas
y columnas iguales a la identidad; su teoría distingue además las
dilataciones con entradas conmutativas\cite{delascuevas2020}. Esta referencia
proporciona un lenguaje preciso para comparar realizaciones de APP.

\subsection{TRIT: orientación y regímenes cuadráticos}

\begin{definition}[Firma ternaria y levantamiento local]
La firma TRIT utiliza la identificación balanceada
\[
 \FF_3\longrightarrow\TRIT,\qquad 0\mapsto0,\quad1\mapsto+1,\quad2\mapsto-1.
\]
Para \(n\in\ZZ\), su levantamiento conserva \(n=r+3c\), donde \(r\in\TRIT\) y \(c\in\ZZ\). En una transición de amplitud local acotada, los tres estados \((0,-1),(0,0),(0,+1)\) pueden presentar el mismo residuo nulo y distintos acarreos.
\end{definition}

La firma no reemplaza a la trayectoria. Un cero residual no significa ausencia de información de orientación, fase o cociente. Tampoco introduce por sí solo una magnitud física. Su función es distinguir algebraicamente las clases de transporte que se realizarán a continuación.

\subsubsection{Tres regímenes cuadráticos y una exponencial común}

Una vez obtenida la firma ternaria, su realización cuadrática se define por
\begin{equation}
 \mathbb A_\tau=\RR[X]/(X^2+\tau),\qquad \tau\in\TRIT.
 \label{eq:algebras-trit}
\end{equation}
Si \(J_\tau\) es la clase de \(X\), entonces \(J_\tau^2=-\tau I\). La estructura distingue el tipo elíptico \(\tau=+1\), el parabólico \(\tau=0\) y el hiperbólico \(\tau=-1\). La realización real no genera la firma: expresa los tipos que la reducción orientada ya ha producido.

\begin{proposition}[Exponenciación de los tres tipos]
En cada álgebra \eqref{eq:algebras-trit},
\begin{equation}
 \exp(tJ_\tau)=C_\tau(t)I+S_\tau(t)J_\tau,
\end{equation}
donde
\[
 C_\tau(t)=\sum_{n\ge0}\frac{(-\tau)^nt^{2n}}{(2n)!},\qquad
 S_\tau(t)=\sum_{n\ge0}\frac{(-\tau)^nt^{2n+1}}{(2n+1)!}.
\]
Las tres realizaciones son
\[
 \begin{array}{c|c|c}
 \tau&J_\tau^2&\exp(tJ_\tau)\\\hline
 +1&-I&\cos t\,I+\sin t\,J_\tau\\
 0&0&I+tJ_\tau\\
 -1&I&\cosh t\,I+\sinh t\,J_\tau.
 \end{array}
\]
\end{proposition}
\begin{proof}
Separar los términos pares e impares de la serie exponencial y sustituir \(J_\tau^{2n}=(-\tau)^nI\) y \(J_\tau^{2n+1}=(-\tau)^nJ_\tau\) da las fórmulas. Para \(\tau=0\), todos los términos de grado al menos dos se anulan. En los otros dos casos se recuperan las series trigonométricas o hiperbólicas.
\end{proof}

La exponencial es común a los tres regímenes: no se asigna exclusivamente a la rama parabólica. Los invariantes
\[
 C_\tau(t)^2+\tau S_\tau(t)^2=1
\]
se deducen de \(\exp(tJ_\tau)\exp(-tJ_\tau)=I\). El mismo principio algebraico conserva la unidad a través de realizaciones que difieren por su tipo cuadrático.

En el régimen hiperbólico, \(P_\pm=(I\pm J_{-1})/2\) son proyectores complementarios y
\[
 \exp(tJ_{-1})=e^tP_++e^{-t}P_-.
\]
El crecimiento y la contracción aparecen conjuntamente en dos direcciones propias. En el elíptico, el parámetro describe una rotación; en el parabólico, una traslación nilpotente. La interpretación temporal adoptada por HMT asocia \(+1\) al retorno con memoria, \(0\) al umbral y \(-1\) a la apertura bidireccional. Esta interpretación exige un orden de parametrización; no añade un cuarto tipo algebraico.

\input{sections/trit_desarrollo.tex}

\input{figures/regimenes_trit.tex}

\newpage
\subsection{TPK: núcleo topológico de fase}

La dinámica utiliza el selector de fase
\[
 M_{\rm ph}(r)=
 \begin{cases}
 +1,&r\in\{1,4,7\},\\
 -1,&r\in\{2,5,8\},\\
 0,&r\in\{3,6,9\}.
 \end{cases}
\]
Su valor decide qué evaluación opera. El transporte cardinal desplaza el cursor correspondiente; el registro conserva el valor anterior, el nuevo cociente, la orientación y la transición. Por tanto, selección y transporte son funciones distintas: un signo ternario no es un desplazamiento del toro.

Denotaremos por \(\widehat X\) el espacio de estados aritméticos con memoria de trayectoria. Un elemento contiene las posiciones y direcciones de los dos cursores, el modo activo, la firma TRIT, la fase, los residuos y cocientes de ambas evaluaciones, los acarreos, el prefijo construido y los datos de frontera. Cuando se introduce un cilindro de refinamiento, éste forma parte del estado. La notación evita sustituir el objeto por una versión reducida que conserve únicamente el símbolo visible.

\begin{definition}[Actualización del TPK]
Sean \(\mathsf{Sel}_t\) la selección de modo mediante \(M_{\rm ph}\), \(\mathsf{Tra}_t\) el transporte cardinal levantado y \(\mathsf{Mem}_t\) la actualización de los registros. La transición es
\begin{equation}
 \mathcal U_t=\mathsf{Mem}_t\circ\mathsf{Tra}_t\circ\mathsf{Sel}_t.
 \label{eq:actualizacion}
\end{equation}
Cada aplicación conserva los campos que no modifica y actualiza los restantes mediante la división euclídea y la concatenación de la trayectoria.
\end{definition}

La notación \eqref{eq:actualizacion} expresa el orden de composición, no reemplaza las reglas efectivas. La siguiente sección especifica completamente su realización finita de dos cursores. Después, la prolongación no se obtendrá borrando su memoria y repitiendo la primera emisión, sino transportando los prefijos y sus fronteras compatibles. De este modo, APP proporciona las evaluaciones, TRIT distingue sus regímenes y TPK compone su evolución.

\input{sections/tpk_desarrollo_integrado.tex}

\input{sections/memoria_resolvente.tex}

% BEGIN NUCLEO_ADITIVO_20260921_CENSO

\section{Núcleo topológico de fase y generación del catálogo de semillas}
\label{act21:tpk:capitulo}

Las dos evaluaciones de APP adquieren eficacia generativa cuando se especifica
qué evaluación actúa, en qué posición se efectúa, cómo se desplaza esa posición
y qué información conserva cada transición. El \emph{Topological Phase Kernel},
abreviado TPK y denominado en español \emph{núcleo topológico de fase}, reúne
estas operaciones en una dinámica sobre estados enriquecidos. Su definición
comprende selección, transporte, actualización aritmética y conservación del
registro. La construcción finita desarrollada en este capítulo proporciona
las semillas sobre las que actuarán las elevaciones regionales y la prolongación
coinductiva.

El resultado principal de esta etapa es un catálogo generado y clasificado:
las condiciones iniciales de los dos cursores producen emisiones enteras;
su lectura ternaria determina palabras orientadas; la acción diedral organiza
esas palabras en órbitas. Cada paso conserva sus preimágenes y las relaciones
que permiten regresar a la descripción precedente. La cadena de cardinales
\[
 104\,976\longrightarrow468\longrightarrow243\subset729,
 \qquad 243\longrightarrow43,
\]
abrevia objetos diferentes. A continuación se construye cada objeto y se
justifica cada recuento, antes de utilizar cualquiera de ellos como dato de
la selección regional.

\subsection{Selección de fase, transporte y estado enriquecido}
\label{act21:tpk:operaciones}

\subsubsection{Selector trítico de fase operativa}

La fase nonádica toma sus valores en \(D_9=\{1,\ldots,9\}\). El selector
trítico es la función
\begin{equation}
 M_{\mathrm{ph}}:D_9\longrightarrow\{-1,0,+1\},\qquad
 M_{\mathrm{ph}}(g)=
 \begin{cases}
 +1,&g\in\{1,4,7\},\\
 -1,&g\in\{2,5,8\},\\
 0,&g\in\{3,6,9\}.
 \end{cases}
 \label{act21:tpk:selector}
\end{equation}
En la realización emisora, el valor \(+1\) activa la evaluación aditiva,
\(-1\) activa la evaluación multiplicativa y \(0\) registra el evento
neutral. La orientación cardinal del cursor es otra coordenada del estado.
Esta separación permite fijar simultáneamente qué operación se lee y en qué
dirección se transporta su punto de evaluación.

Para las transiciones numeradas por \(t\geq1\), la fase previa a la lectura es
\begin{equation}
 g_t=1+[(t-1)]_9,\qquad \epsilon_t=M_{\mathrm{ph}}(g_t),
 \label{act21:tpk:fase}
\end{equation}
donde \([n]_b\) designa el representante de \(n\) módulo \(b\) en
\(\{0,\ldots,b-1\}\). Una ventana de nueve transiciones contiene la sucesión
\((+1,-1,0,+1,-1,0,+1,-1,0)\).

\subsubsection{Cursores orientados y transporte cardinal}

Se utilizan índices \(I_9=\{0,\ldots,8\}\) para las posiciones y etiquetas
positivas \((a,b)=(i+1,j+1)\) para las evaluaciones APP. Un cursor finito es
\[
 c=(i,j;d)\in\mathcal S:=I_9^2\times\mathcal D,
 \qquad \mathcal D=\{N,E,S,O\}.
\]
Las cuatro direcciones actúan mediante
\[
 v_N=(-1,0),\quad v_E=(0,1),\quad
 v_S=(1,0),\quad v_O=(0,-1).
\]
La aplicación cardinal correspondiente a \(d\) es
\[
 T_d(i,j)=\bigl([i+(v_d)_1]_9,[j+(v_d)_2]_9\bigr).
\]
Así, \(M_{\mathrm{ph}}\) y \(T_d\) tienen dominios, codominios y acciones
diferentes. La activación trítica selecciona uno de los dos cursores;
el transporte cardinal modifica su posición en la dirección que ese cursor
ya conserva.

El levantamiento entero del cursor es
\(\widetilde c=(x,y;d)\in\mathbb Z^2\times\mathcal D\), con posición finita
\(([x]_9,[y]_9)\). Sus números de vuelta se recuperan en
\begin{equation}
 x=9\lfloor x/9\rfloor+[x]_9,\qquad
 y=9\lfloor y/9\rfloor+[y]_9.
 \label{act21:tpk:vueltas}
\end{equation}
Estas dos divisiones conciernen al transporte espacial. La división de una
evaluación aditiva o multiplicativa constituye un registro aritmético
distinto, que se conservará junto a ellas.

\subsubsection{Composición de la transición}

El estado de la realización finita incluye los dos cursores levantados,
sus direcciones, la fase, los acumuladores de la ventana activa, la lista
ordenada de bloques ya emitidos y el registro de transiciones. Cada registro
local contiene las posiciones previas, las evaluaciones enteras de ambas
hojas, sus residuos y cocientes, el régimen activo y las orientaciones.
Los campos de prefijo y frontera que aparezcan en una prolongación se
incorporarán al mismo estado con sus reglas de compatibilidad.

La transición se ordena como
\begin{equation}
 \mathcal U_t=\mathsf{Mem}_t\circ\mathsf{Tra}_t\circ\mathsf{Sel}_t.
 \label{act21:tpk:composicion}
\end{equation}
La selección determina el régimen y conserva la muestra previa a todo
desplazamiento. El transporte aplica el movimiento activo y la inversión
cardinal prescrita por el calendario. La actualización incorpora esa muestra
a los acumuladores y al registro; al terminar una ventana, emite su bloque.
La composición utiliza por tanto una muestra anterior al transporte aunque
su archivo definitivo se efectúe al final de la transición.

\subsection{Condiciones iniciales y enumeración reversible}
\label{act21:tpk:semillas}

Las hojas aditiva y multiplicativa disponen de cursores independientes. El
dominio de condiciones iniciales es el producto ordenado
\[
 \Omega_{\mathrm{seed}}=\mathcal S_+\times\mathcal S_\times,
 \qquad |\mathcal S|=9^2\cdot4=324.
\]
Por consiguiente,
\begin{equation}
 |\Omega_{\mathrm{seed}}|=324^2=104\,976.
 \label{act21:tpk:censo-inicial}
\end{equation}
El número \(81\) cuenta posiciones; \(324\) cuenta posiciones orientadas
de un cursor; \(104\,976\) cuenta parejas ordenadas de cursores. La elección
de una posición en APP conserva únicamente una parte de una condición
inicial del emisor.

Se fija la numeración
\[
 \operatorname{ind}(N)=0,\quad\operatorname{ind}(E)=1,\quad
 \operatorname{ind}(S)=2,\quad\operatorname{ind}(O)=3,
\]
y se define
\begin{equation}
 \nu(i,j;d)=4(9i+j)+\operatorname{ind}(d).
 \label{act21:tpk:indice-semilla}
\end{equation}

\begin{proposition}[Enumeración exacta de las condiciones iniciales]
\label{act21:tpk:enumeracion}
La aplicación \(\nu\) es una biyección de \(\mathcal S\) sobre
\(\{0,\ldots,323\}\). La pareja de índices se codifica de manera reversible
por \(324\nu_++\nu_\times\in\{0,\ldots,104\,975\}\).
\end{proposition}
\begin{proof}
La división por cuatro recupera
\[
 \operatorname{ind}(d)=[\nu]_4,\qquad
 j=[\lfloor\nu/4\rfloor]_9,\qquad i=\lfloor\nu/36\rfloor.
\]
Las cotas del dominio garantizan que estas tres cantidades pertenecen a
sus conjuntos de índices y reconstruyen \(\nu\) mediante
\eqref{act21:tpk:indice-semilla}. Para la pareja, la división euclídea por
\(324\) recupera cociente \(\nu_+\) y residuo \(\nu_\times\).
\end{proof}

La enumeración recorre el dominio completo antes de cualquier agrupación
regional. Los índices se utilizan para localizar las preimágenes de una
emisión, manteniendo explícita la condición que produjo cada trayectoria.

\subsection{Evaluación de las hojas y registro aritmético}
\label{act21:tpk:lecturas}

En una posición positiva \((a,b)\), las evaluaciones enteras son
\(A_+(a,b)=a+b\) y \(A_\times(a,b)=ab\). Para \(n>0\), se mantienen
\begin{equation}
 \rho_9(n)=1+[n-1]_9,\qquad
 q_9(n)=\lfloor(n-1)/9\rfloor,\qquad
 n=\rho_9(n)+9q_9(n).
 \label{act21:tpk:division-evaluacion}
\end{equation}
La activación aditiva incorpora \(\rho_9(A_+)\) al acumulador
correspondiente. La activación multiplicativa utiliza un observable del
residuo \(\rho_9(A_\times)\), definido a continuación.

\subsubsection{Exponente discreto en las unidades módulo nueve}

Las potencias de \(2\) módulo \(9\) recorren
\[
 2^0,2^1,\ldots,2^5\equiv1,2,4,8,7,5\pmod9,
 \qquad 2^6\equiv1\pmod9.
\]
El logaritmo discreto en este grupo de seis unidades se representa por
\begin{equation}
 \ell_9(1)=0,\quad\ell_9(2)=1,\quad\ell_9(4)=2,
 \quad\ell_9(8)=3,\quad\ell_9(7)=4,\quad\ell_9(5)=5.
 \label{act21:tpk:logaritmo}
\end{equation}
La ley emisora extiende el observable mediante
\(\ell_9(3)=\ell_9(6)=\ell_9(9)=0\). Para conservar su procedencia, la
muestra registrada incluye el residuo original y el indicador de pertenencia
al grupo de unidades. De ese modo, un valor observado cero puede proceder
de la unidad \(1\) o de alguno de los tres residuos no invertibles, y su
origen permanece recuperable.

\subsubsection{Orden de lectura y movimiento}

Las evaluaciones se efectúan en las posiciones anteriores a la transición.
Si \(\epsilon_t=+1\), se acumula la lectura aditiva y se desplaza su cursor.
Si \(\epsilon_t=-1\), se acumula el exponente discreto multiplicativo y se
desplaza el segundo cursor. Si \(\epsilon_t=0\), se incrementa el contador
neutral y se conservan ambas posiciones.

Al finalizar los pasos \(27\) y \(54\), las dos direcciones se sustituyen
por sus opuestas. Las ventanas se encadenan conservando posiciones y
orientaciones. Al inicio de cada ventana sólo se ponen a cero sus tres
acumuladores locales; el registro anterior permanece disponible.

\begin{proposition}[Inversión del transporte levantado]
\label{act21:tpk:inversion}
Sean \(a,b\in\{0,1\}\) los indicadores de activación e inversión de un
cursor. Si \(\operatorname{opp}\) intercambia \(N\leftrightarrow S\) y
\(E\leftrightarrow O\), la transformación
\[
 F_{a,b}(z,d)=\bigl(z+a v_d,\operatorname{opp}^{\,b}(d)\bigr)
\]
tiene inversa
\[
 F_{a,b}^{-1}(z',d')=
 \bigl(z'-a v_{\operatorname{opp}^{\,b}(d')},
       \operatorname{opp}^{\,b}(d')\bigr).
\]
La reducción espacial módulo nueve entrelaza esta transformación con el
transporte finito.
\end{proposition}
\begin{proof}
Aplicar dos veces la oposición recupera la dirección inicial. Una vez
recuperada esa dirección, la resta de \(a v_d\) deshace el desplazamiento.
Para la proyección espacial basta la igualdad
\([z+a v_d]_9=[[z]_9+a v_d]_9\), componente a componente. La inversión de
dirección conserva esta igualdad.
\end{proof}

La fase determina los indicadores utilizados por la inversa. El registro
espacial y el aritmético pueden así conservarse conjuntamente, incluso al
atravesar el borde del representante \(I_9^2\).

\subsection{Seis ventanas y emisión decimal por bloques}
\label{act21:tpk:emision}

\subsubsection{Acumulación en una ventana nonádica}

La ventana \(m\), con \(1\leq m\leq6\), comprende los pasos
\(9m-8,\ldots,9m\). Sean \(S_m\) la suma de sus residuos positivos
aditivos activos, \(E_m\) la suma de sus exponentes discretos activos y
\(Z_m\) su número de eventos neutrales.

\begin{lemma}[Multiplicidades de activación]
\label{act21:tpk:tres-eventos}
Cada ventana contiene tres activaciones aditivas, tres multiplicativas y
tres neutrales. En particular, \(Z_m=3\).
\end{lemma}
\begin{proof}
Las fases de una ventana recorren exactamente \(1,\ldots,9\). Cada uno
de los tres conjuntos de \eqref{act21:tpk:selector} tiene cardinal tres.
\end{proof}

Las seis ventanas forman \(54\) transiciones y producen seis bloques.
Esta longitud pertenece al emisor inicial aquí definido. La continuación
en profundidad actúa después sobre las palabras y sus registros compatibles.

\subsubsection{Regla de codificación decimal}

La ley emisora especifica las cifras
\begin{align}
 d_{m,1}&=[S_m]_{10},\notag\\
 d_{m,2}&=[S_m+E_m]_{10},\notag\\
 d_{m,3}&=[3S_m+5E_m+7Z_m]_{10},\notag\\
 U_m&=100d_{m,1}+10d_{m,2}+d_{m,3}.
 \label{act21:tpk:bloque-decimal}
\end{align}
Los pesos \(100,10,1\) son las posiciones centena, decena y unidad.
El subíndice \(10\) indica el residuo decimal, y cada bloque satisface
\(0\leq U_m\leq999\). La palabra emitida es la sucesión ordenada
\(\mathbf U=(U_1,\ldots,U_6)\), con ancho fijo de tres cifras por bloque.

Los coeficientes \(3,5,7\) constituyen los pesos explícitos de esta ley
emisora. Una vez fijada la regla, cada cifra se calcula a partir de los
acumuladores, y los resultados de este capítulo se deducen de esa regla.
El número \(1\) que aparece en su forma reducida tiene una procedencia
adicional concreta: \(7Z_m=21\equiv1\pmod{10}\).

Si \(s_m=[S_m]_{10}\) y \(e_m=[E_m]_{10}\), se obtiene
\begin{equation}
 G(s_m,e_m)=100s_m+10[s_m+e_m]_{10}
                   +[3s_m+5e_m+1]_{10}=U_m.
 \label{act21:tpk:acoplamiento-firmas}
\end{equation}
La letra \(e_m\) denota aquí una componente de la firma de exponentes
discretos. Su definición pertenece al emisor finito.

\begin{proposition}[Recuperación de las dos firmas]
\label{act21:tpk:recuperacion-firmas}
La aplicación \(G:\{0,\ldots,9\}^2\to\{0,\ldots,999\}\) es inyectiva.
Para un bloque \(U=100a+10b+c\) de su imagen, la inversa es
\[
 s=a,\qquad e=[b-a]_{10},\qquad
 c=[3s+5e+1]_{10}.
\]
En consecuencia, la palabra \(\mathbf U\) determina las dos firmas de
seis componentes \(\mathbf s\) y \(\mathbf e\).
\end{proposition}
\begin{proof}
Los dos últimos sumandos de \eqref{act21:tpk:acoplamiento-firmas} pertenecen a
\(\{0,\ldots,99\}\), de modo que la centena recupera \(s\). La decena
es \([s+e]_{10}\); al restar \(s\) y reducir módulo diez se recupera
\(e\). La unidad verifica la tercera congruencia. La reconstrucción se
aplica independientemente a las seis posiciones de la palabra.
\end{proof}

Los totales enteros \(S_m,E_m\) se recuperan de las lecturas registradas;
las centenas y decenas recuperan sus residuos. Por ejemplo, un total
\(S_m=15\) produce \(s_m=5\). Esta diferencia identifica exactamente
qué información conserva cada nivel de codificación.

\subsubsection{Codificación posicional de la palabra de seis bloques}

Los seis bloques admiten una codificación entera de ancho fijo,
\begin{equation}
 \mathscr N(\mathbf U)=\sum_{m=1}^{6}U_m1000^{6-m},
 \qquad0\leq\mathscr N(\mathbf U)<1000^6.
 \label{act21:tpk:entero-base1000}
\end{equation}
El primer bloque ocupa la posición de mayor peso y el sexto, la de unidades
de base mil. Un bloque \(058\) ocupa una posición completa con valor
entero \(58\); su ancho de tres cifras forma parte de la presentación
decimal de esa posición.

\begin{proposition}[Recuperación posicional de las emisiones]
\label{act21:tpk:inversa-base1000}
La codificación \eqref{act21:tpk:entero-base1000} determina la palabra completa
mediante
\[
 U_m=\left[\left\lfloor
 \frac{\mathscr N(\mathbf U)}{1000^{6-m}}\right\rfloor\right]_{1000},
 \qquad1\leq m\leq6.
\]
\end{proposition}
\begin{proof}
Al dividir por \(1000^{6-m}\), los bloques anteriores producen múltiplos
enteros de \(1000\), el bloque \(m\) produce \(U_m\), y la contribución
de los posteriores pertenece a \([0,1)\). La parte entera elimina esta
última contribución y el residuo módulo \(1000\) elimina las primeras.
\end{proof}

Esta codificación conserva las seis emisiones. Su dominio posee posiciones
ordenadas de base mil, mientras el dominio APP posee posiciones espaciales
módulo nueve. La composición mantiene ambos índices: \(m\) identifica una
ventana de emisión y \((i,j)\) identifica la celda visitada durante una
transición de esa ventana.

\subsection{Cálculo completo de una emisión que comienza por 501}
\label{act21:tpk:ejemplo501}

Se fijan las condiciones iniciales
\[
 c_{+,0}=(0,4;N),\qquad c_{\times,0}=(2,3;N),
 \qquad \nu_+=16,\quad\nu_\times=84.
\]
Las posiciones positivas iniciales son \((1,5)\) y \((3,4)\). El índice
cronológico previo es \(t=0\), la primera fase leída es \(g_1=1\) y el
registro está vacío. Estas coordenadas son las entradas
del cálculo; los bloques se obtienen ejecutando la recurrencia.

\begin{table}[htbp]
\caption{Primera ventana: lecturas previas y acumulación activa.}
\label{act21:tpk:tabla-501}
\small
\begin{tabular}{@{}rcccrrr@{}}
\toprule
\(t\)&\(\epsilon_t\)&posición \(+\)&posición \(\times\)&
\(\Delta S\)&\(\Delta E\)&\((S,E,Z)\)\\
\midrule
1&\(+1\)&\((1,5)\)&\((3,4)\)&6&0&\((6,0,0)\)\\
2&\(-1\)&\((9,5)\)&\((3,4)\)&0&0&\((6,0,0)\)\\
3&0&\((9,5)\)&\((2,4)\)&0&0&\((6,0,1)\)\\
4&\(+1\)&\((9,5)\)&\((2,4)\)&5&0&\((11,0,1)\)\\
5&\(-1\)&\((8,5)\)&\((2,4)\)&0&3&\((11,3,1)\)\\
6&0&\((8,5)\)&\((1,4)\)&0&0&\((11,3,2)\)\\
7&\(+1\)&\((8,5)\)&\((1,4)\)&4&0&\((15,3,2)\)\\
8&\(-1\)&\((7,5)\)&\((1,4)\)&0&2&\((15,5,2)\)\\
9&0&\((7,5)\)&\((9,4)\)&0&0&\((15,5,3)\)\\
\bottomrule
\end{tabular}
\notatabla{Las posiciones son etiquetas positivas de las celdas anteriores
al movimiento. Los eventos neutrales incrementan únicamente \(Z\).}
\end{table}

Las tres evaluaciones aditivas activas son
\[
 1+5=6,\qquad9+5=14=5+9,\qquad8+5=13=4+9,
\]
de modo que \(S_1=6+5+4=15\). Las tres evaluaciones multiplicativas activas
son \(3\cdot4=12\), \(2\cdot4=8\) y \(1\cdot4=4\). Sus residuos
positivos son \(3,8,4\), con observables \(0,3,2\); por tanto,
\(E_1=5\). Finalmente \(Z_1=3\). La regla decimal da
\[
 d_{1,1}=[15]_{10}=5,\quad
 d_{1,2}=[15+5]_{10}=0,\quad
 d_{1,3}=[45+25+21]_{10}=1,
\]
y se obtiene \(U_1=501\).

La segunda lectura aditiva anterior tiene cociente aritmético \(1\),
mientras su cursor levantado está en \((-1,4;N)\), con número de vuelta
vertical \(-1\). Son dos datos diferentes del mismo registro. Ambos
intervienen en la recuperación íntegra de la evaluación y de su posición.

\begin{table}[htbp]
\caption{Las seis ventanas de la condición inicial \((16,84)\).}
\label{act21:tpk:tabla-seis501}
\begin{tabular}{@{}rrrrrrrr@{}}
\toprule
\(m\)&pasos&\(S_m\)&\(E_m\)&\(Z_m\)&\(s_m\)&\(U_m\)&\([-U_m]_3\)\\
\midrule
1&1--9&15&5&3&5&501&0\\
2&10--18&6&5&3&6&614&1\\
3&19--27&24&5&3&4&498&0\\
4&28--36&21&5&3&1&169&2\\
5&37--45&12&5&3&2&272&1\\
6&46--54&12&5&3&2&272&1\\
\bottomrule
\end{tabular}
\end{table}

La palabra resultante y sus firmas son
\begin{align*}
 \mathbf U&=(501,614,498,169,272,272),\\
 \mathbf s&=(5,6,4,1,2,2),&
 \mathbf e&=(5,5,5,5,5,5).
\end{align*}
La inversión cardinal tras el paso \(27\) modifica los recorridos de las
tres ventanas finales. El bloque repetido \(272\) corresponde a dos
ventanas sucesivas con posiciones y registro cronológico propios.

Como segundo caso, la condición inicial mínima \((\nu_+,\nu_\times)=(0,0)\)
produce
\[
 \mathbf U=(252,109,252,903,850,850),\quad
 \mathbf s=(2,1,2,9,8,8),\quad
 \mathbf e=(3,9,3,1,7,7).
\]
En este caso, ambos cursores regresan a su posición y orientación iniciales
tras \(54\) transiciones. La fase nonádica completa seis vueltas, el
contador cronológico avanza \(54\) unidades y el registro contiene
\(54\) muestras. La lectura posicional del retorno y
la historia registrada describen propiedades distintas de la ejecución.

\subsection{Factorización del censo finito}
\label{act21:tpk:censo}

Las firmas \(f_+(\nu)=\mathbf s\) y \(f_\times(\nu)=\mathbf e\) se
calculan ejecutando el cursor correspondiente durante las seis ventanas.
La independencia de posiciones y direcciones entre los dos cursores da
\begin{equation}
 T_{\mathrm{em}}(\nu_+,\nu_\times)
   =G^{(6)}\bigl(f_+(\nu_+),f_\times(\nu_\times)\bigr),
 \label{act21:tpk:factorizacion}
\end{equation}
donde \(G^{(6)}\) aplica \eqref{act21:tpk:acoplamiento-firmas} en cada posición.
La factorización permite enumerar primero \(324+324\) trayectorias y
combinar después las firmas distintas, conservando la multiplicidad de
sus preimágenes.

\par\noindent\begin{minipage}{\linewidth}
\subsubsection{Imágenes de las dos firmas}

Las tablas siguientes incluyen todas las firmas. Su multiplicidad es el
número de condiciones iniciales que las producen; el índice mínimo localiza
una de esas condiciones. La preimagen completa se define por
\[
 f_\sigma^{-1}(a)=\{\nu\in\{0,\ldots,323\}:f_\sigma(\nu)=a\}.
\]
La recurrencia de las secciones anteriores determina cada elemento de este
conjunto mediante un cálculo entero finito.

\begin{table}[H]
\centering\fontsize{9.5}{11.4}\selectfont
\setlength{\abovecaptionskip}{8pt}\setlength{\belowcaptionskip}{6pt}
\renewcommand{\arraystretch}{1.12}
\caption{Las dieciocho firmas aditivas.}
\label{act21:tpk:firmas-aditivas}
\begin{tabular}{@{}lrr@{\qquad}lrr@{}}
\toprule
firma&multiplicidad&\(\nu_{\min}\)&firma&multiplicidad&\(\nu_{\min}\)\\
\midrule
\texttt{122564}&18&17&\texttt{564122}&18&16\\
\texttt{122898}&18&24&\texttt{645212}&18&4\\
\texttt{212645}&18&5&\texttt{654889}&18&33\\
\texttt{212988}&18&0&\texttt{889221}&18&13\\
\texttt{221456}&18&29&\texttt{889654}&18&32\\
\texttt{221889}&18&12&\texttt{898122}&18&25\\
\texttt{456221}&18&28&\texttt{898465}&18&20\\
\texttt{465898}&18&21&\texttt{988212}&18&1\\
\texttt{546988}&18&9&\texttt{988546}&18&8\\
\bottomrule
\end{tabular}

\par\vspace{12pt}

\caption{Las veintiséis firmas multiplicativas.}
\label{act21:tpk:firmas-multiplicativas}
\begin{tabular}{@{}lrr@{\qquad}lrr@{}}
\toprule
firma&multiplicidad&\(\nu_{\min}\)&firma&multiplicidad&\(\nu_{\min}\)\\
\midrule
\texttt{000000}&108&8&\texttt{555555}&24&84\\
\texttt{177393}&8&1&\texttt{555717}&8&14\\
\texttt{177555}&8&54&\texttt{555771}&8&18\\
\texttt{177771}&8&11&\texttt{555933}&8&24\\
\texttt{339555}&8&6&\texttt{717555}&8&12\\
\texttt{339771}&8&15&\texttt{717717}&8&21\\
\texttt{339933}&8&59&\texttt{717933}&8&19\\
\texttt{393177}&8&0&\texttt{771177}&8&9\\
\texttt{393393}&8&45&\texttt{771339}&8&13\\
\texttt{393555}&8&40&\texttt{771555}&8&16\\
\texttt{555177}&8&52&\texttt{933339}&8&57\\
\texttt{555339}&8&4&\texttt{933555}&8&26\\
\texttt{555393}&8&41&\texttt{933717}&8&17\\
\bottomrule
\end{tabular}
\end{table}
\end{minipage}\par


\begin{proposition}[Imagen y multiplicidades del emisor]
\label{act21:tpk:468}
La imagen \(\mathcal U_6=\operatorname{im}T_{\mathrm{em}}\) contiene
\(468\) palabras distintas. Sus preimágenes se distribuyen en \(432\)
fibras de cardinal \(144\), \(18\) de cardinal \(432\) y \(18\) de
cardinal \(1944\).
\end{proposition}
\begin{proof}
La evaluación exhaustiva del algoritmo de la sección~\ref{act21:tpk:algoritmo}
sobre los índices \(0,\ldots,323\) da las dos tablas anteriores.
La primera partición tiene \(18\) clases de cardinal \(18\), y la segunda
\(24\) clases de cardinal \(8\), una de cardinal \(24\) y una de
cardinal \(108\). Las sumas \(18\cdot18=324\) y
\(24\cdot8+24+108=324\) comprueban la masa total de ambas particiones.
La proposición~\ref{act21:tpk:recuperacion-firmas} hace inyectivo el acoplamiento
de sus imágenes; por tanto, hay \(18\cdot26=468\) emisiones.
Por \eqref{act21:tpk:factorizacion}, la preimagen de una pareja de firmas es
el producto de sus dos preimágenes. Se obtienen \(18\cdot24=432\)
productos de cardinal \(18\cdot8=144\), \(18\) de cardinal
\(18\cdot24=432\) y \(18\) de cardinal \(18\cdot108=1944\).
Finalmente,
\[
 432\cdot144+18\cdot432+18\cdot1944=104\,976.
\]
La parte enumerativa es una comprobación finita exacta de la recurrencia;
la factorización explica por qué ese recuento cubre todas las parejas.
\end{proof}

La emisión del ejemplo \ref{act21:tpk:ejemplo501} tiene una fibra de cardinal
\(18\cdot24=432\). El ejemplo de índices \((0,0)\) tiene una fibra de
cardinal \(18\cdot8=144\). En ambos casos la palabra emitida recupera
las firmas, mientras la identificación de una condición inicial particular
utiliza su índice o su registro de trayectoria.

\subsection{Lectura ternaria orientada y conservación de las fibras}
\label{act21:tpk:reduccion-ternaria}

La reducción orientada actúa separadamente sobre cada bloque:
\begin{equation}
 \Psi:\mathcal U_6\longrightarrow\mathbb F_3^6,\qquad
 \Psi(U_1,\ldots,U_6)=([-U_1]_3,\ldots,[-U_6]_3).
 \label{act21:tpk:psi}
\end{equation}
Su signo fija una orientación del catálogo. La identificación balanceada
\(0\mapsto0\), \(1\mapsto+1\), \(2\mapsto-1\) se aplica después,
componente a componente.

El módulo nueve organiza las evaluaciones y posiciones de APP; los residuos
módulo diez forman cada cifra; la base mil reúne tres posiciones decimales;
el módulo tres determina la firma ternaria de cada bloque. Estos cuatro usos
tienen objetos y funciones definidos. En particular, \(\Psi\) lee seis
bloques y produce seis trits; una conversión posicional del entero formado
por dieciocho cifras sería otra operación.

Para las dos condiciones calculadas se obtiene
\begin{align*}
 \Psi(501,614,498,169,272,272)&=(0,1,0,2,1,1),\\
 \Psi(252,109,252,903,850,850)&=(0,2,0,0,2,2).
\end{align*}
El emisor decimal está definido antes de esta reducción. La función de
\(\Psi\) es producir el objeto ternario sobre el que actúan la simetría
orbital y los selectores regionales. Esa función explica su posición en la
genealogía: enlaza una emisión ya calculada con una clasificación estructural.

La imagen \(\mathcal W_6:=\Psi(\mathcal U_6)\) contiene exactamente
\(243\) palabras dentro del conjunto ambiente de \(3^6=729\) palabras.
La enumeración del algoritmo final da el histograma
\begin{equation}
\begin{array}{c|rrrrrr}
 |\Psi^{-1}(w)|&1&2&3&4&5&6\\\hline
 \text{número de palabras }w&132&66&18&3&6&18.
\end{array}
\label{act21:tpk:fibras-ternarias}
\end{equation}
Las sumas de ambas lecturas son
\[
132+66+18+3+6+18=243,
\qquad132+2\cdot66+3\cdot18+4\cdot3+5\cdot6+6\cdot18=468.
\]
La primera cuenta palabras ternarias; la segunda recupera el número de
emisiones distribuidas entre sus fibras.

\begin{example}[Dos emisiones con la misma firma ternaria]
\label{act21:tpk:colision}
Las emisiones
\[
 (498,501,614,252,252,109),\qquad
 (870,810,923,252,252,109)
\]
son distintas y ambas tienen imagen \((0,0,1,0,0,2)\) por \(\Psi\).
Cada una tiene \(144\) condiciones iniciales. La palabra ternaria conserva
una clase de lectura; las emisiones y sus preimágenes distinguen los datos
que confluyen en esa clase.
\end{example}

La conservación de fibras se expresa concretamente mediante
\[
 \mathcal F_w=\{\mathbf U\in\mathcal U_6:\Psi(\mathbf U)=w\},
 \qquad
 \Omega_{\mathbf U}=T_{\mathrm{em}}^{-1}(\mathbf U).
\]
Por tanto, la preimagen completa de \(w\) en el dominio inicial es la
unión disjunta de los conjuntos \(\Omega_{\mathbf U}\), con
\(\mathbf U\in\mathcal F_w\). El registro enriquecido puede conservar la
emisión, sus firmas y la condición inicial junto a su lectura ternaria.
La inversa exacta se refiere entonces al dato registrado con su tipo;
la palabra de seis trits aislada mantiene el alcance preciso de
\eqref{act21:tpk:psi}.

\subsection{Acción diedral y las cuarenta y tres órbitas}
\label{act21:tpk:orbitas}

Sobre seis coordenadas se definen las permutaciones
\begin{align}
 r(u_1,u_2,u_3,u_4,u_5,u_6)
   &=(u_3,u_1,u_2,u_5,u_6,u_4),\\
 s(u_1,u_2,u_3,u_4,u_5,u_6)
   &=(u_4,u_5,u_6,u_1,u_2,u_3).
 \label{act21:tpk:generadores-diedrales}
\end{align}
La primera rota los dos grupos de tres coordenadas con sentidos conjugados;
la segunda intercambia ambos grupos. Su composición satisface
\(r^3=s^2=1\) y \(srs=r^{-1}\). Generan así una acción del grupo diedral
\(D_3\), de seis elementos.

\begin{proposition}[Equivariancia y estabilidad del catálogo]
\label{act21:tpk:equivarianza}
La reducción \(\Psi\) conmuta con la acción de \(D_3\). Los conjuntos
\(\mathcal U_6\) y \(\mathcal W_6\) son estables bajo sus generadores.
\end{proposition}
\begin{proof}
La igualdad \(\Psi(g\mathbf U)=g\Psi(\mathbf U)\) se verifica en cada
coordenada porque \(g\) permuta posiciones y \(\Psi\) aplica la misma
reducción a cada posición. La estabilidad del catálogo constituye otra
comprobación: al generar sus \(468\) elementos mediante
\eqref{act21:tpk:factorizacion}, se verifica que \(r\mathbf U\) y
\(s\mathbf U\) pertenecen de nuevo a ese conjunto. El algoritmo de la
sección~\ref{act21:tpk:algoritmo} efectúa estas comprobaciones de pertenencia.
La equivariancia transmite después la estabilidad a su imagen ternaria.
\end{proof}

\begin{proposition}[Descomposición orbital completa]
\label{act21:tpk:43}
El conjunto \(\mathcal W_6\) se descompone en \(43\) órbitas: \(38\)
de cardinal seis y \(5\) de cardinal tres.
\end{proposition}
\begin{proof}
Para cada \(w\in\mathcal W_6\) se forma
\[
 \mathcal O(w)=\{w,rw,r^2w,sw,srw,sr^2w\}.
\]
Se escoge como representante la palabra lexicográficamente mínima de ese
conjunto. Dos palabras tienen el mismo representante exactamente cuando
pertenecen a la misma órbita. El algoritmo genera la lista de representantes
de la tabla~\ref{act21:tpk:representantes}; sus cardinales suman
\(38\cdot6+5\cdot3=243\). Como cada palabra del catálogo participa en
esta operación y el catálogo es estable, la lista cubre toda la imagen.
\end{proof}

\begin{table}[!htbp]
\caption{Representantes mínimos y cardinales de las cuarenta y tres órbitas.}
\label{act21:tpk:representantes}
\small
\begin{tabular}{@{}lr@{\qquad}lr@{\qquad}lr@{}}
\toprule
representante&cardinal&representante&cardinal&representante&cardinal\\
\midrule
\texttt{001002}&6&\texttt{002112}&6&\texttt{012222}&6\\
\texttt{001011}&6&\texttt{002211}&6&\texttt{021022}&6\\
\texttt{001012}&6&\texttt{002220}&6&\texttt{021121}&6\\
\texttt{001021}&6&\texttt{011021}&6&\texttt{021202}&6\\
\texttt{001022}&6&\texttt{011112}&6&\texttt{022022}&3\\
\texttt{001100}&3&\texttt{011120}&6&\texttt{022112}&6\\
\texttt{001101}&6&\texttt{011201}&6&\texttt{022121}&6\\
\texttt{001110}&6&\texttt{011211}&6&\texttt{022202}&3\\
\texttt{001112}&6&\texttt{011220}&6&\texttt{022211}&6\\
\texttt{001202}&6&\texttt{011222}&6&\texttt{022222}&6\\
\texttt{001210}&6&\texttt{012112}&6&\texttt{112112}&3\\
\texttt{001211}&6&\texttt{012121}&6&\texttt{112211}&3\\
\texttt{001220}&6&\texttt{012202}&6&\texttt{112222}&6\\
\texttt{001222}&6&\texttt{012210}&6&&\\
\texttt{002012}&6&\texttt{012220}&6&&\\
\bottomrule
\end{tabular}
\end{table}

Las órbitas heredan invariantes enteros de sus fibras. Para cada palabra,
\(n(w)=|\mathcal F_w|\) cuenta emisiones, mientras
\[
 M(w)=\sum_{\mathbf U\in\mathcal F_w}|\Omega_{\mathbf U}|
\]
cuenta condiciones iniciales. El censo de símbolos
\(c(w)=(n_0,n_1,n_2)\) registra cuántas veces aparece cada clase ternaria.
Estos datos, junto con la orientación y el soporte de los recorridos,
proporcionan los predicados de la selección regional posterior. La
clasificación conserva así más información que la forma de una palabra
representante.

\subsection{Algoritmo autocontenido de enumeración y control}
\label{act21:tpk:algoritmo}

La enumeración utiliza enteros, listas finitas y las operaciones ya definidas.
El siguiente procedimiento precisa su ejecución; \(\sigma\) toma los valores
\(+\) y \(\times\).

\begin{enumerate}
\item Para cada \(\nu=0,\ldots,323\), recuperar \((i,j;d)\) mediante
la inversa de \eqref{act21:tpk:indice-semilla}. Iniciar una lista vacía de firmas.
\item Para cada ventana \(m=1,\ldots,6\), iniciar el acumulador \(A=0\).
Recorrer \(t=9m-8,\ldots,9m\), calculando \(\epsilon_t\) por
\eqref{act21:tpk:fase}.
\item Si \(\sigma=+\) y \(\epsilon_t=+1\), añadir
\(\rho_9(i+j+2)\) a \(A\) y actualizar
\((i,j)\leftarrow T_d(i,j)\). Si \(\sigma=\times\) y
\(\epsilon_t=-1\), añadir
\(\ell_9(\rho_9((i+1)(j+1)))\) y efectuar el mismo transporte.
\item Al terminar \(t=27\) o \(t=54\), sustituir \(d\) por su opuesta.
Al concluir la ventana, añadir \([A]_{10}\) a la lista de firmas.
Conservar la posición y la dirección para la siguiente ventana.
\item Agrupar los \(324\) índices por igualdad de sus listas de seis
componentes. Estas dos agrupaciones dan \(f_+\), \(f_\times\) y todas
sus preimágenes.
\item Para cada pareja de firmas distintas, aplicar \(G^{(6)}\).
Asignar a la emisión la multiplicidad producto de las dos preimágenes.
Comprobar la recuperación de firmas en cada bloque y la suma de
multiplicidades \(104\,976\).
\item Aplicar \(\Psi\) a cada emisión y agrupar por igualdad de palabras.
Contar los \(243\) valores y las multiplicidades de
\eqref{act21:tpk:fibras-ternarias}.
\item Aplicar \(r,s\) a cada emisión y verificar pertenencia al catálogo.
Para cada palabra ternaria, formar sus seis imágenes diedrales, eliminar
repeticiones y registrar su mínimo lexicográfico. Contar las \(43\)
órbitas y los dos tamaños de la tabla~\ref{act21:tpk:representantes}.
\end{enumerate}

La ejecución termina: el primer tramo evalúa \(648\) trayectorias de
\(54\) transiciones; el segundo combina \(468\) parejas de firmas; el
tercero agrupa un conjunto de \(243\) palabras. Una comprobación directa
adicional puede ejecutar las \(104\,976\) parejas y contrastar la
factorización \eqref{act21:tpk:factorizacion}; su alcance coincide con el mismo
dominio finito.

\subsection{Conservación del registro y función de las semillas}
\label{act21:tpk:conclusion}

Sea \(\mathsf{Sample}(q_t)\) la muestra previa a la transición y
\(\mathcal M_t\) el registro acumulado. La actualización de memoria es
\begin{equation}
 \mathcal M_{t+1}=\mathcal M_t\mathbin{\Vert}\mathsf{Sample}(q_t),
 \label{act21:tpk:archivo}
\end{equation}
donde \(\Vert\) denota concatenación ordenada. La inducción da
\[
 \mathcal M_n=\mathcal M_0\mathbin{\Vert}
 (\mathsf{Sample}(q_0),\ldots,\mathsf{Sample}(q_{n-1})).
\]
Los acumuladores se reconstruyen sumando las contribuciones activas de
cada segmento de nueve registros. Sus cocientes, las vueltas de los
cursores y las orientaciones permanecen asociados a las mismas transiciones.
La ejecución determina ese registro mediante la recurrencia; el archivo
preserva la procedencia concreta de las lecturas.

Al finalizar esta etapa están construidos el dominio completo de condiciones
iniciales, la emisión decimal, sus dos firmas recuperables, la imagen ternaria,
sus fibras y su clasificación orbital. Los censos \(104\,976\), \(468\),
\(243\), \(729\) y \(43\) tienen ahora dominio y mecanismo de obtención.
Las selecciones regionales recibirán estas estructuras con sus invariantes,
y las elevaciones actuarán sobre palabras orientadas con procedencia
conservada. Esa continuidad convierte el catálogo finito en antecedente
constructivo del desarrollo coinductivo.

% END NUCLEO_ADITIVO_20260921_CENSO
\section{Desarrollo del núcleo formal: condiciones iniciales, geometría y refinamiento}
\label{sec:extension}

El soporte de ochenta y una posiciones no coincide con el dominio de condiciones iniciales de la dinámica. Una condición inicial debe indicar además una orientación; y las hojas aditiva y multiplicativa disponen de cursores independientes. Esta distinción permite reconstruir exhaustivamente la emisión sin seleccionar de antemano una región numérica.

\subsection{Los estados iniciales y la recurrencia de emisión}

Para la enumeración utilizamos índices \(I_9=\{0,\ldots,8\}\), relacionados con las etiquetas positivas por \(i\mapsto i+1\). Sea
\[
 \mathcal S=I_9^2\times\{N,E,S,O\},\qquad
 \Omega_{\rm seed}=\mathcal S_+\times\mathcal S_\times.
\]
Cada condición inicial \(\omega=(s_+,s_\times)\) determina dos posiciones y dos direcciones. El dominio comprende todas las parejas ordenadas, por lo que
\begin{equation}
 |\mathcal S|=324,\qquad |\Omega_{\rm seed}|=324^2=104\,976.
 \label{eq:semillas}
\end{equation}
El anexo independiente enumera todas estas condiciones; el cuerpo presenta la ley que las transforma. Las filas de una tabla de emisiones son fibras de este dominio de condiciones iniciales.

La realización finita utiliza seis ventanas de nueve transiciones. En cada instante \(t\), la firma de fase es \(M_{\rm ph}(1+((t-1)\bmod9))\). Si vale \(+1\), se lee el residuo positivo \(\Sigma(i+1,j+1)=\rho_9(i+j+2)\) del primer cursor y después se desplaza ese cursor; si vale \(-1\), se lee el residuo positivo \(\Pi(i+1,j+1)=\rho_9((i+1)(j+1))\) del segundo cursor y después se desplaza el segundo. El valor cero incrementa el contador neutral. Al finalizar las transiciones veintisiete y cincuenta y cuatro, se invierten las orientaciones cardinales. Las evaluaciones enteras y sus cocientes se conservan en el registro, pero el emisor utiliza las lecturas residuales aquí especificadas.

En el sector de unidades, el lector multiplicativo utiliza el logaritmo discreto de base dos en \((\ZZ/9\ZZ)^\times\):
\[
 \ell_9(1)=0,\quad\ell_9(2)=1,\quad\ell_9(4)=2,\quad
 \ell_9(8)=3,\quad\ell_9(7)=4,\quad\ell_9(5)=5.
\]
Para la emisión finita se extiende con valor cero al sector no unitario. Esta extensión es un observable y no una identificación de los elementos \(3,6,9\) con la unidad multiplicativa. La hoja y su pertenencia al sector radical permanecen en el registro de trayectoria.

Sean \(S_m\) la suma de los tres residuos positivos aditivos \(\Sigma\) de la ventana \(m\), \(E_m\) la suma de los tres exponentes discretos de los residuos multiplicativos \(\Pi\) y \(Z_m=3\) el número de eventos neutrales. Así, \(S_m\) no suma las evaluaciones enteras \(S=i+j\), cuyos cocientes permanecen en el registro. Con \([a]_{10}\in\{0,\ldots,9\}\), la regla de emisión decimal es
\begin{align}
 d_{m,1}&=[S_m]_{10},\notag\\
 d_{m,2}&=[S_m+E_m]_{10},\notag\\
 d_{m,3}&=[3S_m+5E_m+7Z_m]_{10},\notag\\
 U_m&=100d_{m,1}+10d_{m,2}+d_{m,3}.
 \label{eq:emisor-seis}
\end{align}
Los coeficientes enteros y el calendario de esta recurrencia son parte explícita de la ley emisora declarada. No se reciben valores de \(\pi,\varphi,e\) ni de \(\alpha\). La distinción entre una regla de emisión y una constante reconocida a partir de su salida evita que la genealogía quede escondida en una tabla de decimales.

Si \(s_m=[S_m]_{10}\) y \(e_m=[E_m]_{10}\), la misma regla se escribe
\begin{equation}
 U_m=100s_m+10[s_m+e_m]_{10}+[3s_m+5e_m+1]_{10}.
 \label{eq:acoplamiento}
\end{equation}
El uno final es el residuo de \(7Z_m=21\) módulo diez. La reducción ternaria orientada del catálogo adopta la convención
\begin{equation}
 \Psi(U_1,\ldots,U_6)=([-U_1]_3,\ldots,[-U_6]_3).
 \label{eq:psi-catalogo}
\end{equation}
El signo de esta coordenatización es explícito: cambiarlo modifica las palabras y exige transportar las orientaciones posteriores. La identificación balanceada de cada componente de \(\FF_3\) con el TRIT se realiza después de \eqref{eq:psi-catalogo}.

\begin{proposition}[Factorización inyectiva de la emisión]
La palabra \(U=(U_1,\ldots,U_6)\) determina las dos firmas \(s=(s_1,\ldots,s_6)\) y \(e=(e_1,\ldots,e_6)\). La imagen de la emisión tiene \(18\cdot26=468\) elementos. La proyección ternaria tiene \(243\) valores distintos en el ambiente de \(729\) palabras.
\end{proposition}
\begin{proof}
Las centenas de \(U_m\) recuperan \(s_m\). Si \(b_m\) es su cifra de decenas, entonces \(e_m=[b_m-s_m]_{10}\). La cifra de unidades comprueba la tercera congruencia de \eqref{eq:acoplamiento}. Por ello el acoplamiento es inyectivo en el producto de firmas.

La enumeración de las \(324\) condiciones de un cursor mediante la recurrencia anterior produce dieciocho firmas aditivas, cada una con dieciocho preimágenes; y veintiséis multiplicativas, de las cuales veinticuatro tienen ocho preimágenes, una tiene veinticuatro y una tiene ciento ocho. El anexo conserva estas firmas y sus condiciones iniciales, permitiendo repetir la enumeración finita sin un catálogo numérico de referencia. La inyectividad del acoplamiento da \(18\cdot26\) emisiones y multiplicidades
\[
 \underbrace{432}_{18\cdot24}\text{ fibras de }144,
 \qquad18\text{ fibras de }432,
 \qquad18\text{ fibras de }1944.
\]
La suma es \(432\cdot144+18\cdot432+18\cdot1944=104\,976\). Aplicar \eqref{eq:psi-catalogo} a las \(468\) emisiones produce exactamente \(243\) palabras. Esta última cifra es el cardinal de la imagen, no el del espacio \(\FF_3^6\).
\end{proof}

Esta factorización explica la diferencia entre exhaustividad finita y profundidad ilimitada. El dominio \eqref{eq:semillas} es completo para el soporte y el emisor declarados; no contiene anticipadamente cada prefijo de cada prolongación. La profundidad pertenece al sistema de extensiones compatibles que se construye sobre esas condiciones.

\subsection{Clausura del calendario y avance de memoria}

El calendario de ciento ocho transiciones consta de doce ventanas nonádicas. Su estructura de grupo es más precisa que un producto informal de dos períodos:
\begin{equation}
 0\longrightarrow C_{12}\xrightarrow{\iota}C_{108}
 \xrightarrow{p}C_9\longrightarrow0,
 \quad\iota([b])=[9b],\quad p([t])=[t]_9.
 \label{eq:extension108}
\end{equation}
Si \(t=9b+r\) con \(0\le r<9\), la suma de dos fases produce el acarreo
\[
 (b,r)+(b',r')=
 \left(b+b'+\left\lfloor\frac{r+r'}9\right\rfloor,\ [r+r']_9\right),
\]
donde la primera coordenada se reduce módulo doce. La composición conserva el término que enlaza las dos coordenadas.

\begin{proposition}[No escisión del calendario]
La sucesión \eqref{eq:extension108} es exacta y no admite una sección que sea homomorfismo de grupos.
\end{proposition}
\begin{proof}
El núcleo de la reducción módulo nueve es el subgrupo de múltiplos de nueve, imagen de \(\iota\). Si la extensión se escindiera, por conmutatividad se tendría \(C_{108}\cong C_{12}\times C_9\). El primer grupo tiene exponente \(108\), mientras el segundo tiene exponente \(\operatorname{mcm}(12,9)=36\), contradicción.
\end{proof}

Nueve transiciones restauran la fase nonádica y avanzan una ventana; ciento ocho restauran el calendario finito. Ninguno de estos hechos obliga a borrar la trayectoria o a restaurar el prefijo de una prolongación. La holonomía nonádica \(\Gamma_9\) designa el transporte de una vuelta sobre estados con memoria, de modo que
\begin{equation}
 g(\Gamma_9x)=g(x),\qquad M(\Gamma_9x)=M(x)+1.
 \label{eq:holonomia-memoria}
\end{equation}
La igualdad de fase y el incremento de memoria son compatibles y constituyen dos observaciones del mismo transporte.

\input{sections/temporalidad_trit.tex}

La representación monodrómica de este retorno se desarrolla en la
sección~\ref{mono-ampl:conexion}. Allí se distingue la holonomía del
estado de su representación sobre fase y contador. La suspensión
simbólica de la sección~\ref{mono-ampl:cristal} define el modelo de
cristal temporal aperiódico, y la sección~\ref{mono-ampl:euler}
compone esa temporalidad con las tres realizaciones de Euler.

\subsection{Compresión observable y conservación ortogonal}

Una realización isométrica permite expresar de manera cuantitativa qué pierde una observación reducida. Sea \(U:\mathcal H\to\mathcal H\) una realización isométrica del transporte de una vuelta, \(U^*U=I\), y \(J:\mathcal H_{\rm vis}\to\mathcal H\) una inclusión isométrica. Definimos
\[
 T=J^*UJ,\qquad D=(I-JJ^*)UJ.
\]
Entonces
\begin{equation}
 UJ=JT+D,\qquad J^*D=0,\qquad I-T^*T=D^*D.
 \label{eq:conservacion-ortogonal}
\end{equation}
La última igualdad se deduce tomando el producto adjunto de la primera y usando la ortogonalidad. La contracción de la componente visible no constituye por sí sola destrucción de información del estado total.

\begin{proposition}[Conservación a horizonte finito]
Si \(I-T^*T=D^*D\), entonces, para todo \(n\ge1\),
\begin{equation}
 I=\sum_{r=0}^{n-1}T^{*r}D^*DT^r+T^{*n}T^n.
 \label{eq:telescopia}
\end{equation}
\end{proposition}
\begin{proof}
Multiplicar \(D^*D=I-T^*T\) por \(T^{*r}\) y \(T^r\) y sumar produce una suma telescópica. Se cancelan los términos interiores y queda \(I-T^{*n}T^n\).
\end{proof}

El último sumando conserva el estado terminal. Suprimirlo antes de justificar un límite alteraría la igualdad. Esta distinción es el contenido operativo de la conservación de la unidad en la compresión: la norma se distribuye entre la componente visible, la memoria y el terminal, sin ser reemplazada por un único residuo.

\subsection{Desarrollo local: conmutación, firma lorentziana y curvatura discreta}

El bloque diagonal adyacente con diagonales iguales módulo nueve es único. En efecto,
\[
 k^2\equiv(k+1)^2\pmod9\quad\Longleftrightarrow\quad2k+1\equiv0\pmod9
\]
da \(k=4\). Su matriz visible es
\[
 B_c=\begin{pmatrix}7&2\\2&7\end{pmatrix}=7I+2R,
 \quad R=\begin{pmatrix}0&1\\1&0\end{pmatrix},\quad R^2=I.
\]
Con \(P_\pm=(I\pm R)/2\), se obtiene \(B_c=9P_++5P_-\). La unidad de la normalización estocástica y la conservación de una forma indefinida son dos realizaciones de este bloque:
\[
 \frac{B_c}{9}=P_++\frac59P_-,\qquad
 \frac{B_c}{\sqrt{45}}=\exp(\eta R),\quad\eta=\frac12\log\frac95.
\]
Para \(Q=\operatorname{diag}(1,-1)\), la segunda matriz satisface \(B_c^{\mathsf T}QB_c=45Q\). De aquí procede una realización local en \(SO^+(1,1)\); no se requiere introducir de antemano un espacio-tiempo de cuatro dimensiones para demostrarla.

El promediado uniforme sobre las clases \(C_1=\{1,4,7\}\), \(C_2=\{2,5,8\}\), \(C_0=\{3,6,9\}\) da
\[
 M_\Pi=\begin{pmatrix}4&5&6\\5&4&6\\6&6&9\end{pmatrix},\quad
 M_\Sigma=\begin{pmatrix}5&6&4\\6&4&5\\4&5&6\end{pmatrix}.
\]
El polinomio característico de \(M_\Pi\) es
\[
 (\lambda+1)((\lambda-9)^2-72),
\]
de donde su forma cuadrática tiene firma \((2,1)\). El salto espectral local \(9-5=4\), el coeficiente \(2\), la diferencia agregada \(51-45=6\) y el defecto integrado \(9\cdot6=54\) describen niveles distintos de la misma comparación aritmética. No son índices intercambiables de un mismo operador.

La curvatura discreta puede describirse directamente antes de esta realización espectral. En residuos, sean \(S(x,y)=x+y\), \(P(x,y)=xy\) y \(\Delta_x,\Delta_y\) las diferencias progresivas. Las variables de enlace son diferencias de potencial, \(A_x^T=\Delta_xT\), \(A_y^T=\Delta_yT\), y se fija
\[
 F(T_x,T_y)=\Delta_x\Delta_yT_y-\Delta_y\Delta_xT_x.
\]
Como \(\Delta_xS=\Delta_yS=1\), \(\Delta_xP=y\) y \(\Delta_yP=x\),
\begin{equation}
 F(S,S)=F(P,P)=0,\qquad F(S,P)=+1,\quad F(P,S)=-1\pmod9.
\end{equation}
La suma y el producto, tomados separadamente, inducen conexiones planas; la mezcla ordenada cambia el signo de la curvatura de plaqueta. Al sumar sobre un rectángulo de lados \(a,b\), las aristas interiores se cancelan y la suma orientada de borde vale \(\pm ab\) módulo nueve. Así aparece una relación exacta entre orden de composición y orientación de frontera.

\subsection{Completación cúbica y dos nociones de frontera}

La ampliación del soporte cúbico de lado nueve al de lado diez admite una estratificación posicional exacta:
\begin{equation}
 10^3=9^3+3\cdot9^2+3\cdot9+1=729+243+27+1.
 \label{eq:cubo}
\end{equation}
El primer término corresponde a las tres coordenadas interiores; el segundo, a una coordenada nueva; el tercero, a dos; y el último, al vértice de tres coordenadas nuevas. La corona de ampliación contiene \(271\) posiciones. En cambio, la frontera interna del cubo discreto de lado nueve contiene \(9^3-7^3=386\). La diferencia entre estos recuentos es una diferencia de dominio, no un conflicto de valores.

Dividir \eqref{eq:cubo} por \(1000\) da una partición normalizada de la unidad. La incidencia del cubo distingue seis caras, doce aristas y ocho vértices. El enlace ulterior con hexadas y octadas de Steiner requiere los mapas combinatorios correspondientes; los cardinales del cubo no los sustituyen. Esta completación aporta aquí la relación entre las bases \(729\) y \(1000\), mientras la prolongación excepcional se desarrollará después de la generación numérica y del electrón.

\input{sections/revision_emrh.tex}

\input{sections/geometria_modular.tex}
\input{sections/aritmetica_historias.tex}

\section{Generación coinductiva de \texorpdfstring{\(\pi,\varphi,e\)}{pi, phi, e} a profundidad arbitraria}
\label{sec:generacion}

La emisión de seis bloques es una etapa finita de la construcción, no una expansión decimal completa. Su función consiste en producir regiones con multiplicidad, orientación y reglas de transporte. La prolongación debe conservar esos datos y determinar una colección compatible de prefijos de longitud arbitraria. La palabra \emph{generación} designa esta operación; \emph{genealogía} designa la reconstrucción de sus dependencias. Ambas nociones intervienen, pero cumplen funciones distintas.

\subsection{Órbitas, regiones y caracteres de clausura, propagación y autoescala}

Sobre una palabra \(w=(w_1,\ldots,w_6)\), definimos
\[
 r(w)=(w_3,w_1,w_2,w_5,w_6,w_4),\qquad
 s(w)=(w_4,w_5,w_6,w_1,w_2,w_3).
\]
Se cumple \(r^3=s^2=I\) y \(srs=r^{-1}\); por tanto estas permutaciones realizan el grupo diedral \(D_3\). La acción conserva el catálogo y las multiplicidades y conmuta con \(\Psi\). El cociente de las \(243\) palabras visibles tiene \(43\) órbitas: treinta y ocho de cardinal seis y cinco de cardinal tres.

La clase de clausura se reconoce por un predicado de fibra: sus seis orientaciones tienen cinco emisiones por palabra, multiplicidad total \(1008\) y distribución \(432+4\cdot144\). Existe una sola órbita con esas propiedades en el catálogo. Para precisar las otras dos selecciones, escribimos \(f(w)=\#\Psi^{-1}(w)\), \(m(w)=\sum_{U\in\Psi^{-1}(w)}\operatorname{mult}(U)\) y \(\nu(w)=(n_0,n_1,n_2)\). La clase diagonal aditiva de una emisión es \(d(U)=[i+j]_9\) en los índices de cero a ocho; su firma aditiva conserva esta clase y la orientación \(ES\) o \(NO\). Sea \(\mathcal D(\mathcal O)\) el conjunto de esas clases para todas las emisiones sobre una órbita.

Entre las órbitas de tamaño seis, el predicado conjunto
\[
 f(w)=1,\qquad m(w)=144\quad(w\in\mathcal O),\qquad
 \mathcal D(\mathcal O)=\{0,3,6\}
\]
deja exactamente seis órbitas. En este sector, el censo \(\nu=(2,3,1)\), igual al de clausura, selecciona únicamente la propagación; el censo \(\nu=(2,2,2)\) selecciona únicamente la autoescala. El soporte diagonal no puede omitirse: sin él hay cuatro órbitas unipuntuales equilibradas de multiplicidad \(144\), no una. Finalmente, la existencia de una emisión con \(d=6\) y orientación \(ES\) selecciona un representante de cada órbita y conduce a
\begin{equation}
 w^{\rm cl}=010211,\qquad w^{\rm pr}=201101,\qquad w^{\rm au}=121200.
 \label{eq:palabras-regionales}
\end{equation}
Estos símbolos designan primero regiones del catálogo. La identificación escalar se demostrará mediante sus lectores; no se eligen porque coincidan con cifras de una constante conocida.

La palabra de clausura \(010211\) tiene las cinco preimágenes siguientes:
\begin{center}\small
\begin{tabular}{@{}lll r@{}}\toprule
Emisión de seis bloques&Firma multiplicativa&Papel&Multiplicidad\\\midrule
\(501,614,498,169,272,272\)&\(555555\)&transversal&432\\
\(810,923,870,169,272,272\)&\(339555\)&orientación positiva&144\\
\(810,983,810,169,272,272\)&\(393555\)&orientación positiva&144\\
\(870,923,810,169,272,272\)&\(933555\)&orientación positiva&144\\
\(870,923,810,418,674,521\)&\(933717\)&retorno orientado&144\\\bottomrule
\end{tabular}
\end{center}
La multiplicidad \(432\) y la firma constante \(555555\) caracterizan la región transversal. No identifican la posición central \((5,5)\) de APP. Confundir una firma de seis ventanas con una posición del soporte eliminaría la distinción entre región numérica y estructura electrónica.

La microfibra posee además una morfología bipartita precisa. Cada emisión es un producto de una clase de cursor aditivo y una clase multiplicativa. Las tres regiones positivas comparten la misma clase aditiva de dieciocho condiciones; juntas contienen tres clases multiplicativas de ocho condiciones cada una. La región transversal posee un producto \(18\times24\); las tres positivas reunidas, otro \(18\times24\); y el retorno, un producto \(18\times8\). La descomposición por emisiones tiene cinco partes, mientras la descomposición por componentes conexas del grafo bipartito tiene tres. Ambas conservan las mismas \(1008\) semillas sin reducir la forma a un cardinal.

\subsection{Transporte finito y frontera de prolongación}

Los endomorfismos \(L_0,L_1\) del espacio de palabras ternarias actúan antes de la realización arquimediana. Su determinación comienza por las transiciones del estado aritmético, no por la prescripción de dos matrices. Si \(U_t\) es el paso compuesto de selección, transporte y actualización, el bloque observado satisface
\[
 b_j^\chi=\Pi_6(x_j^\chi),\qquad
 x_{j+1}^\chi=U_{9j+9}\cdots U_{9j+1}(x_j^\chi),
 \qquad \chi\in\{\mathrm{cl},\mathrm{pr},\mathrm{au}\}.
\]
Se forman seis transiciones por cada régimen: dos consecutivas en cada una de las tres regiones. Los pares de matrices de orígenes y destinos obtenidos en el registro son
\[
\begin{array}{c|c|c|c}
X_0&Y_0&X_1&Y_1\\\hline
010211&012222&010211&002111\\
012222&010211&002111&110221\\
201101&121221&102011&012222\\
121221&102011&012222&102011\\
121200&112202&121020&010210\\
112202&121020&010210&010200
\end{array}
\]
Cada cadena de seis dígitos representa una fila en \(\FF_3^6\). Puesto que \(\det X_0=1\) y \(\det X_1=-1\) en \(\FF_3\), las ecuaciones \(X_aL_a=Y_a\) determinan unívocamente \(L_a=X_a^{-1}Y_a\). En esa coordenatización de vectores fila resultan
\[
 L_0=\begin{pmatrix}
 2&2&2&1&2&1\\2&1&2&2&1&1\\1&0&0&1&0&2\\
 0&0&1&1&1&0\\2&2&2&0&2&1\\2&1&2&1&0&0
 \end{pmatrix},\quad
 L_1=\begin{pmatrix}
 2&1&2&0&2&2\\1&2&1&1&2&0\\1&2&1&1&1&2\\
 0&1&0&0&2&1\\2&0&2&1&0&1\\0&2&2&2&1&1
 \end{pmatrix}\quad(\bmod3).
\]
El calendario \((L_0,L_0,L_1,L_1)\) produce cuatro nuevos bloques de seis componentes. Sus concatenaciones forman
\[
 w_6\longrightarrow w_{12}\longrightarrow w_{18}
 \longrightarrow w_{24}\longrightarrow w_{30}.
\]
Los subíndices de \(w\) indican longitud acumulada. La frontera \(R_{36}\) conserva la estructura terminal y abre la relación de prolongación; no se renombra como un último bloque periódico que pudiera repetirse indefinidamente.

La identidad \(L=X^{-1}Y\) demuestra unicidad una vez fijadas las transiciones; no sustituye la producción de esas transiciones por \(U_t\). La procedencia de su registro y el examen de los dieciséis calendarios binarios se documentan en el suplemento técnico. El calendario indicado es el único de esos dieciséis que reproduce conjuntamente las tres sucesiones de cinco bloques.

\subsection{Cinco regiones de clausura y compensación angular}

La región de clausura contiene una componente transversal y cuatro acompañantes. El lector simétrico actúa en el espacio de sus cinco posiciones mediante
\[
 P_5=\frac15\mathbf1\mathbf1^{\mathsf T},\qquad
 \mathbf1=(1,1,1,1,1)^{\mathsf T}.
\]
La invariancia \(P_5\lambda=\lambda\) y la conservación \(\mathbf1^{\mathsf T}\lambda=1\) fuerzan \(\lambda=\mathbf1/5\). Esta normalización distribuye la unidad entre posiciones del lector; no sustituye las multiplicidades de semillas \(432,144,144,144,144\) por frecuencias uniformes. La coordenatización elíptica del lector realiza cada coordenada mediante \(1+\lambda_jJ\); su pendiente es \(\lambda_j\), por lo que la pendiente primitiva resulta \(1/5\). Se hace explícita así la regla que pasa de una coordenada normalizada a una pendiente, en lugar de identificar ambas nociones sin mapa. El lector compone las cuatro posiciones acompañantes por la ley
\begin{equation}
 a\oplus b=\frac{a+b}{1-ab}.
 \label{eq:suma-pendientes}
\end{equation}
Esta ley se obtiene al multiplicar \((1+aJ)(1+bJ)\) en el régimen \(J^2=-I\) y dividir por su componente escalar. Por ello la composición de pendientes es una realización de la estructura cuadrática de TRIT, no una suma ordinaria de cuatro fracciones.

\begin{proposition}[Compensador de la clausura]
La composición de cuatro pendientes \(1/5\) es \(120/119\). La condición de retorno a pendiente uno determina un único compensador positivo de la forma \(1/q\), con \(q>1\), y da \(q=239\).
\end{proposition}
\begin{proof}
La primera duplicación da \((1/5)\oplus(1/5)=5/12\); la segunda, \((5/12)\oplus(5/12)=120/119\). La condición
\[
 \frac{120/119-1/q}{1+(120/119)/q}=1
\]
equivale a \((120-119)q=120+119\), por lo que \(q=239\).
\end{proof}

El carácter de clausura resultante tiene realización
\begin{equation}
 \Pi_{\rm cl}=4\left(4\arctan\frac15-\arctan\frac1{239}\right).
 \label{eq:lector-pi}
\end{equation}
La identidad angular de Machin reconoce \(\Pi_{\rm cl}=\pi\). El sentido de la construcción es preciso: la pendiente y el número de composiciones están declarados por el lector de la pentafibra; la ecuación de compensación fuerza \(239\). Una vez producido este carácter, su evaluación por series alternadas no selecciona de nuevo la órbita de semillas.

Para \(0<x<1\), las sumas alternadas
\[
 A_n(x)=\sum_{j=0}^n\frac{(-1)^jx^{2j+1}}{2j+1}
\]
encierran \(\arctan x\) entre dos truncamientos consecutivos, con diferencia \(x^{2n+3}/(2n+3)\). Aplicadas a las dos pendientes de \eqref{eq:lector-pi}, proporcionan extremos racionales de error explícito. El valor numérico se calcula así como consecuencia del carácter exacto, sin introducir un prefijo objetivo.

\subsection{Lector de propagación y normalización de la unidad}

La región de propagación se realiza mediante una colección formal \(P_t\in\mathcal A[[t]]\), donde \(\mathcal A\) es un álgebra conmutativa unitaria sobre \(\QQ\). Su composición satisface \(P_{t+s}=P_tP_s\), con unidad \(P_0=1\), y el transporte elemental orientado fija el generador escalar \(P'_0=+1\). La condición fija su signo, no sólo su magnitud. Escribiendo
\[
 P_t=\sum_{n\ge0}a_nt^n,
\]
la normalización determina \(a_0=a_1=1\). La comparación del coeficiente de \(s^nt\) en \(P_{s+t}=P_sP_t\) da \((n+1)a_{n+1}=a_na_1\); equivalentemente, \(P'_t=P_t\). Así,
\begin{equation}
 (n+1)a_{n+1}=a_n,\qquad a_n=\frac1{n!}.
 \label{eq:propagacion}
\end{equation}
El carácter en una unidad de parámetro es \(\Pi_{\rm pr}=P_1\), reconocido posteriormente como \(e\).

La normalización del generador importa. La ley semigrupal sin ese dato admitiría \(P_t=e^{ct}\) para distintos \(c\); por ello no se omite la condición que fija la unidad del parámetro en el lector. Esta condición es anterior a la comparación decimal y pertenece a la exposición de la regla operacional, no a la elección de un valor experimental.

\begin{proposition}[Cotas racionales de propagación]
Para \(n\ge1\), sea \(E_n=\sum_{j=0}^n1/j!\). Entonces
\[
 E_n<\Pi_{\rm pr}<E_n+\frac{n+2}{(n+1)(n+1)!}.
\]
\end{proposition}
\begin{proof}
El primer término omitido es \(1/(n+1)!\). Cada cociente posterior es a lo sumo \(1/(n+2)\). La suma geométrica mayorante de la cola es
\[
 \frac1{(n+1)!}\frac1{1-1/(n+2)}
 =\frac{n+2}{(n+1)(n+1)!}.
\]
La desigualdad estricta se sigue de la disminución de los cocientes sucesivos.
\end{proof}

\subsection{Incidencia modal y generación de la autoescala}

El selector trítico produce el ciclo \((+1,-1,0)\). En los modos \(\Sigma,\Pi\), la regla es \(+1:m\mapsto\Sigma\), \(-1:m\mapsto\Pi\), \(0:m\mapsto m\). Con cierre cíclico, el modo observado es \((\Sigma,\Pi,\Pi)\). Sus aristas son \(\Sigma\to\Pi\), \(\Pi\to\Pi\) y \(\Pi\to\Sigma\). Por tanto su incidencia, en ese orden de modos, es
\begin{equation}
 F_{\rm av}=\begin{pmatrix}0&1\\1&1\end{pmatrix}.
 \label{eq:fav}
\end{equation}
Las cuatro entradas se obtienen del calendario, no de una ecuación que haya recibido previamente el valor de \(\varphi\). La repetición del calendario da conteos \(3F_{\rm av},9F_{\rm av},36F_{\rm av}\) en longitudes nueve, veintisiete y ciento ocho. Estos conteos no son las potencias de la matriz.

\begin{theorem}[Carácter positivo de autoescala]
La incidencia \eqref{eq:fav} posee un único rayo propio estrictamente positivo. Al normalizarlo como \((1,x)^{\mathsf T}\), su coordenada \(x\) satisface \(x^2-x-1=0\), \(x>0\), y es \(\varphi=(1+\sqrt5)/2\).
\end{theorem}
\begin{proof}
La ecuación propia da \(x=\lambda\) y \(1+x=\lambda x\). Eliminar \(\lambda\) produce el polinomio. De sus dos raíces sólo una es positiva. Como \(F_{\rm av}^2\) tiene todas sus entradas positivas, el rayo positivo es único y simple. La expresión radical reconoce la coordenada obtenida; no participa en la producción de \(F_{\rm av}\).
\end{proof}

Con \(F_0=0,F_1=1,F_{n+1}=F_n+F_{n-1}\), la incidencia induce, para \(n\geq1\),
\[
 F_{\rm av}^n=\begin{pmatrix}F_{n-1}&F_n\\F_n&F_{n+1}\end{pmatrix}.
\]
Los cocientes consecutivos \(F_{n+1}/F_n\) alternan alrededor de \(\varphi\). La identidad de determinante
\[
 F_{n+1}^2-F_nF_{n+2}=(-1)^n
\]
da una anchura exacta \(1/(F_nF_{n+1})\) entre dos cocientes consecutivos. Su crecimiento hace arbitrariamente pequeña esta horquilla racional.

La medida sobre todas las historias compatibles no coincide con la frecuencia de modos del calendario de tres eventos. La frecuencia cronológica es \((1/3,2/3)\). La medida de máxima entropía del envolvente simbólico definido por \(F_{\rm av}\) se construye con sus vectores propios y tiene razón de pesos \(\varphi^{-2}\). Esta última interviene en la lectura areal--volumétrica; no se obtiene sustituyendo sin más una frecuencia de conteo por otra.

\input{sections/retorno_areal_volumetrico.tex}

Después de generados los caracteres, las dos orientaciones del lector
\[
 R(x,y)=\frac{x}{y^2}
\]
proporcionan
\begin{equation}
 \rho_A=\frac\pi{\varphi^2},\qquad
 \rho_E=\frac\varphi{\pi^2},\qquad
 \varphi^3\rho_A^2\rho_E=1.
 \label{eq:areal-electron}
\end{equation}
La segunda expresión aparecerá en el exponente electrónico. La primera se obtiene también como límite de la lectura marcada \(\pi F_{n-1}/F_{n+1}\). El intercambio de argumentos relaciona ambas, pero no convierte una reparametrización de la ecuación electrónica en una segunda derivación independiente de todos sus coeficientes.

\input{sections/euler_trit_articulacion.tex}

\subsection{Cilindros compatibles y profundidad arbitraria}

La realización de un bloque ternario \(b=(b_1,\ldots,b_6)\) es el entero
\[
 \nu(b)=\sum_{j=1}^6b_j3^{6-j}\in\{0,\ldots,728\}.
\]
Una sucesión de bloques define prefijos enteros por \(N_{n+1}=729N_n+\nu(b_{n+1})\), comenzando en \(N_0=m\), donde \(m\) es la parte entera determinada por el lector. Las horquillas anteriores sitúan los caracteres de clausura, propagación y autoescala en \((3,4)\), \((2,3)\) y \((1,2)\), respectivamente; por tanto, sus valores iniciales son \(m=3,2,1\). Los bloques siguientes refinan la parte fraccionaria. El cilindro correspondiente es
\begin{equation}
 I_n=\left[\frac{N_n}{729^n},\frac{N_n+1}{729^n}\right],
 \qquad |I_n|=729^{-n}.
 \label{eq:cilindro}
\end{equation}
La convención semiabierta se utiliza para seleccionar un representante cuando un valor cae en una frontera. Para el paso al límite se utilizan las clausuras de estos intervalos y se identifican las dos expansiones fronterizas equivalentes. Esta precisión impide suponer que toda intersección de intervalos semiabiertos encajados sea automáticamente no vacía.

\begin{theorem}[Determinación de la realización escalar]
Una historia de bloques compatible determina un único límite de los extremos izquierdos de \eqref{eq:cilindro}. Si un carácter exacto dispone de horquillas racionales de anchura arbitrariamente pequeña y se separa de la frontera de cada cilindro decimal examinado, cada prefijo decimal de longitud finita se determina después de un cálculo finito. Los prefijos así obtenidos son compatibles entre sí.
\end{theorem}
\begin{proof}
La extensión de un bloque conserva el prefijo anterior por división euclídea. Los cilindros cerrados están encajados y su diámetro tiende a cero; sus extremos izquierdos son de Cauchy y tienen un único límite. Para una longitud decimal fijada, un valor separado de las fronteras de la partición posee distancia positiva a ellas. Una horquilla suficientemente estrecha queda contenida en una única celda y determina su prefijo. La unicidad de la celda y el encajamiento de las particiones prueban la compatibilidad por truncamiento. En los puntos con expansión terminal debe añadirse la decisión exacta de frontera y su convención de representación; no se deduce esa decisión de una anchura numérica solamente.
\end{proof}

\begin{corollary}[Obtención finita y compatibilidad de los prefijos regionales]
\label{gen:prefijos-regionales}
Para cada una de las realizaciones \(\pi,\varphi,e\) y para toda longitud
decimal finita, las horquillas construidas determinan el prefijo correspondiente
mediante un cálculo finito. Los prefijos obtenidos a distintas profundidades
son compatibles por truncamiento.
\end{corollary}
\begin{proof}
Los caracteres de clausura, propagación y autoescala poseen las horquillas
anteriores. Tras su reconocimiento, la irracionalidad de \(\pi,e,\varphi\)
garantiza su separación de las fronteras racionales de cada partición decimal
finita. Se aplica el teorema precedente a cada carácter.
\end{proof}
La profundidad arbitraria expresa el cuantificador del corolario: para todo
horizonte finito existe un cálculo que determina su prefijo. El objeto
matemático es la colección compatible completa; cada ejecución obtiene una
proyección finita de ella.

\subsection{Involución especular y sincronización de las regiones}

La conexión nonádica transporta conjuntamente los tres caracteres. En la coordenatización de la novena fase, el eje de autoescala y la suma ternaria de clausura y propagación permanecen fijados bajo el intercambio especular. La inversión helicoidal cambia el sentido del transporte; el espejo de los dos canales cambia su orientación relativa. Son involuciones diferentes, aunque ambas actúen sobre el mismo sistema de trayectorias.

Las coincidencias de censo ilustran esa coordinación:
\begin{center}
\begin{tabular}{@{}c rrr l@{}}\toprule
Fase&\(N_\pi\)&\(N_e\)&\(N_\varphi\)&Igualdad de conteos\\\midrule
4&207&207&122&\(N_\pi=N_e\)\\
5&39&151&151&\(N_e=N_\varphi\)\\
6&110&110&69&\(N_\pi=N_e\)\\
7&81&29&81&\(N_\varphi=N_\pi\)\\
8&59&59&59&sincronización triple\\
9&43&43&43&sincronización triple\\\bottomrule
\end{tabular}
\end{center}
Estos enteros cuentan coordenadas compatibles de las regiones. No afirman una igualdad entre \(\pi\), \(e\) y \(\varphi\) como números reales. Si
\[
 \partial(a,b,c)=(a-b,b-c,c-a),
\]
su imagen pertenece al retículo \(A_2=\{(u,v,w)\in\ZZ^3:u+v+w=0\}\). En las fases ocho y nueve la componente relativa se anula, pero la componente diagonal cambia de \((59,59,59)\) a \((43,43,43)\). La sincronización relativa coexiste con un desplazamiento de \(-16\) por canal. El cierre que conducirá a \(\alpha\) debe conservar ambas informaciones: coincidencia relativa y evolución de la memoria.

\subsection{Suspensión helicoidal esturmiana}

La relación entre las bases nueve y diez de la completación define el parámetro adimensional
\begin{equation}
 \delta=\frac{\log(10/9)}{\log10}.
 \label{eq:delta-sturm}
\end{equation}
No se obtiene de una constante física objetivo. Es la escala logarítmica relativa de dos bases ya presentes en la construcción. Su empleo como lectura simbólica del refinamiento tiene una consecuencia exacta.

\begin{lemma}[Irracionalidad de la pendiente]
El número \(\delta\) es irracional.
\end{lemma}
\begin{proof}
Si \(\delta=p/q\), \(q>0\), entonces \((10/9)^q=10^p\) y \(9^q=10^{q-p}\). La factorización del primer miembro contiene el primo tres y la del segundo sólo los primos dos y cinco. La igualdad es imposible.
\end{proof}

La palabra mecánica inferior, con intercepto \(\theta\), es
\[
 z_n(\theta)=\lfloor(n+1)\delta+\theta\rfloor-\lfloor n\delta+\theta\rfloor\in\{0,1\}.
\]
Se obtiene al observar la rotación \(\theta\mapsto\theta+\delta\pmod1\) mediante el intervalo \([1-\delta,1)\). La sucesión de eventos tiene densidad \(\delta\), y la suma sobre cualquier ventana de longitud \(L\) vale \(\lfloor L\delta\rfloor\) o \(\lceil L\delta\rceil\). Como \(21<1/\delta<22\), las separaciones entre eventos son veintiuno o veintidós. La palabra que indica cuáles son los intervalos cortos es otra palabra mecánica de pendiente \(\sigma=22-1/\delta\); sus eventos están separados por seis o siete intervalos primarios. La segunda escala es inducida por la primera y no un segundo parámetro libre.

\input{sections/vacancias_capacidad.tex}

\begin{theorem}[Complejidad de las lecturas de fase]
Para \(m\ge1\), la palabra mecánica tiene complejidad factorial \(p(m)=m+1\). Al conservar la fase \(n\bmod9\), la complejidad es \(p_9(m)=9(m+1)\); al conservar sólo su clase ternaria, es \(p_3(m)=3(m+1)\). Las tres entropías topológicas son nulas.
\end{theorem}
\begin{proof}
Los predecesores de los dos extremos del intervalo de codificación forman, para factores de longitud \(m\), \(m+1\) puntos distintos de la circunferencia: la coincidencia entre extremos desplazados elimina las duplicaciones y la irracionalidad impide coincidencias adicionales. Esos puntos determinan \(m+1\) intervalos de itinerario constante y distintos. Toda órbita irracional visita cada intervalo, de donde \(p(m)=m+1\).

Fijada una fase módulo nueve, las posiciones de esa fase recorren una rotación de paso \(9\delta\), también irracional. Por densidad aparecen todos los factores en cada fase, y la primera coordenada distingue sus nueve copias. El mismo argumento con paso \(3\delta\) da tres copias tras la reducción ternaria. Finalmente, \(m^{-1}\log(c(m+1))\to0\) para \(c=1,3,9\).
\end{proof}

El producto de fases \(C_9\times\mathbb T\), trasladado por \((1,\delta)\), tiene caracteres de frecuencia
\[
 \frac a9+k\delta\pmod1,\qquad a\in\ZZ/9\ZZ,\ k\in\ZZ.
\]
El módulo espectral no es \(\log10\): es el módulo aditivo generado por \(1/9\) y \(\delta\), módulo los enteros. El factor \(\log10\) pertenece a la normalización de \eqref{eq:delta-sturm}.

Una presentación operatoria hace explícita la memoria. En \(\ell^2(\ZZ)\), sea \(T|n\rangle=|n+1\rangle\), \(V_9|n\rangle=\zeta_9^n|n\rangle\), \(W_\delta|n\rangle=e^{2\pi i n\delta}|n\rangle\) y \(D_M|n\rangle=\lfloor n/9\rfloor|n\rangle\). Sobre vectores de soporte finito,
\begin{equation}
 V_9T=\zeta_9TV_9,\qquad
 W_\delta T=e^{2\pi i\delta}TW_\delta,\qquad
 [D_M,T^9]=T^9.
 \label{eq:weyl-relojes}
\end{equation}
Las dos primeras relaciones son covariancias de Weyl; la tercera expresa que una vuelta nonádica aumenta el grado de memoria. La rotación de base es abeliana, mientras el álgebra de observables y traslaciones no lo es. La inversión \(|n\rangle\mapsto|-n\rangle\) conjuga el avance con el retroceso. Esta realización proporciona una suspensión helicoidal aperiódica con retorno finito de fase, no una palabra periódica de longitud nueve.

\input{figures/holonomia.tex}

Para este recuento fijamos el intercepto \(\theta=0\) y los extremos acumulados \((N_0,\ldots,N_4)=(0,729,972,999,1000)\), correspondientes al orden cúbico \((729,243,27,1)\). Los conteos inferiores son \(\lfloor N_j\delta\rfloor-\lfloor N_{j-1}\delta\rfloor\), y dan \((33,11,1,0)\). Los superiores son \(\lceil N_j\delta\rceil-\lceil N_{j-1}\delta\rceil\), con la misma frontera inicial cero, y dan \((34,11,1,0)\). Estos valores dependen del intercepto y del orden de los estratos; no se afirman para una fase inicial arbitraria. Sus totales son cuarenta y cinco y cuarenta y seis. El incremento corresponde a una única unidad de frontera, situada en el primer estrato de esta partición ordenada. Las cifras treinta y cuatro y once aparecen así en un mismo recuento previo a toda coordenatización de unidades; identificarlas con exponentes físicos exige, además, la derivación dimensional correspondiente.

La denominación \emph{cristal aperiódico} se emplea aquí para este orden simbólico, con complejidad, espectro y retorno especificados. La entropía topológica nula no equivale por sí sola a disipación física nula: una realización termodinámica necesita una dinámica energética y una operación de borrado definidas. Del mismo modo, esta lectura esturmiana no sustituye a todos los operadores del TPK ni demuestra por sí sola la periodicidad o aperiodicidad de la salida completa de \(\alpha\). Su función es hacer explícita una estructura de orden que el simple retorno del calendario no describe.

\input{sections/revision_monodromia.tex}

\input{sections/primos_retornos.tex}

% BEGIN NUCLEO_ADITIVO_20260921_SELECCION

% Ampliación demostrativa DF003--DF007. No modifica la fuente integral.
% Propietario común de los archivos Lean citados en estos comentarios:
% output/PAQUETE_ARTICULO_I_INTEGRACION_ESTADO_CAMPOS_20260919/.
% Residencia causal de la primera sección: después de los lectores regionales
% coinductivos y antes de emplear sus registros como entrada de la selección K.
% Residencia causal de la segunda sección: después de la construcción marcada
% Paley--Witt--Golay--Leech y antes de la reconstrucción de la estructura VOA.
% Los dos cortes expositivos conservan el mismo origen APP--TRIT--TPK.

\section[Generación regional y refinamiento]{Compatibilidad integral de la generación regional, el refinamiento cilíndrico y la incidencia}
\label{act21:sec:df-lectores-cilindros}

La prolongación regional produce sucesiones ordenadas de bloques, junto con
sus profundidades de lectura. La proyección arquimediana y la incidencia
ternaria se calculan sobre esas mismas sucesiones. El vínculo entre ambas
realizaciones puede formularse enteramente mediante operaciones enteras:
selección por intervalos racionales, división euclídea, elevación de seis
trits, lectura de acarreo y transformación de Witt. Esta formulación permite
seguir cada coeficiente desde su bloque generador hasta su empleo posterior.

La construcción utiliza los tres lectores regionales ya obtenidos del estado
APP--TRIT--TPK. Los denotaremos por los índices
\(c\in\{\mathrm P,\mathrm E,\mathrm F\}\), correspondientes, respectivamente,
a las regiones de clausura, propagación y autoescala. Su realización real se
escribirá \(v_c\). Las identificaciones posteriores
\(v_{\mathrm P}=\pi_{\mathrm{HMT}}\),
\(v_{\mathrm E}=e_{\mathrm{HMT}}\) y
\(v_{\mathrm F}=\varphi_{\mathrm{HMT}}\) conservan ese orden causal.
La operación que emite un bloque inspecciona los intervalos racionales del
lector regional; el valor real interviene en la demostración de corrección
de dicha operación.

\subsection{Emisión posicional mediante separación racional de celdas}
\label{act21:subsec:df-emision-racional}

% Fuente: lean/biblioteca/RegionalPublicationComposition.lean,
% produced_brackets, eventual_publication, joint_eventual_publication,
% search_terminates, stopped_cell_correct, publications_compatible.

Para cada región se dispone de dos sucesiones racionales
\(\ell_c(j),u_c(j)\), construidas por su lector, que cumplen
\[
 \ell_c(j)\leq v_c\leq u_c(j).
\]
La propiedad de separación de celdas demostrada para esos lectores afirma
que, para cada entero \(s>0\), existe \(J\) tal que, para todo \(j\geq J\),
\begin{equation}
 \lfloor s\ell_c(j)\rfloor
 =\lceil su_c(j)\rceil-1
 =\lfloor sv_c\rfloor.
 \label{act21:eq:df-separacion-celdas}
\end{equation}
La evitación de las fronteras racionales por los valores regionales es la
hipótesis de los teoremas previos que garantiza esta separación. Aquí se
emplea su conclusión operativa, con los lectores y los intervalos que la
producen.

\begin{definition}[Índice de separación y prefijo emitido]
Para \(s>0\), sea
\[
 j_c(s)=\min\{j\in\mathbb N:
 \lfloor s\ell_c(j)\rfloor=\lceil su_c(j)\rceil-1\}.
\]
Para una base entera \(b\geq2\) y una profundidad \(n\geq0\), se define
\begin{equation}
 P_c(b,n)=\lfloor b^n\ell_c(j_c(b^n))\rfloor.
 \label{act21:eq:df-publicacion-prefijo}
\end{equation}
El prefijo contiene la parte entera y las primeras \(n\) posiciones en base
\(b\); por consiguiente, \(P_c(b,0)\) es la parte entera regional.
\end{definition}

\begin{theorem}[Terminación y corrección de la emisión]
\label{act21:thm:df-emision-correcta}
El índice \(j_c(s)\) existe para todo \(s>0\). Los prefijos emitidos cumplen
\begin{equation}
 P_c(b,n)=\lfloor b^n v_c\rfloor,
 \qquad
 P_c(b,n)\leq b^n v_c<P_c(b,n)+1.
 \label{act21:eq:df-prefijo-correcto}
\end{equation}
Para cada escala \(s\) puede elegirse además un único índice de suficiencia
\(J(s)\) válido simultáneamente para las tres regiones.
\end{theorem}

\begin{proof}
La igualdad eventual de \eqref{act21:eq:df-separacion-celdas} hace no vacío el
conjunto que define \(j_c(s)\). Su mínimo existe por buen orden de
\(\mathbb N\). En cualquier índice donde se produzca la igualdad, la cota
inferior racional proporciona
\(\lfloor s\ell_c(j)\rfloor\leq\lfloor sv_c\rfloor\).
La cota superior, junto con la evitación de una frontera de celda, da
\(\lfloor sv_c\rfloor\leq\lceil su_c(j)\rceil-1\).
Al coincidir ambos extremos enteros, el entero intermedio coincide con
ellos. Esto demuestra \eqref{act21:eq:df-prefijo-correcto}. Para la afirmación
conjunta basta tomar el máximo de los tres índices de suficiencia obtenidos
para la misma escala. El índice conjunto controla simultáneamente las tres
emisiones, mientras que cada región conserva sus propios intervalos.
\end{proof}

\begin{proposition}[Compatibilidad exacta de los prefijos]
\label{act21:prop:df-prefijos-compatibles}
Para \(n,m\geq0\),
\begin{equation}
 \left\lfloor\frac{P_c(b,n+m)}{b^m}\right\rfloor=P_c(b,n).
 \label{act21:eq:df-truncamiento}
\end{equation}
En particular, el bloque siguiente
\(d_c(b,n)=P_c(b,n+1)\bmod b\) satisface
\begin{equation}
 P_c(b,n+1)=bP_c(b,n)+d_c(b,n),\qquad 0\leq d_c(b,n)<b.
 \label{act21:eq:df-paso-posicional}
\end{equation}
\end{proposition}

\begin{proof}
Escribamos \(b^n v_c=N+\theta\), con
\(N=P_c(b,n)\) y \(0\leq\theta<1\). Entonces
\(P_c(b,n+m)=b^mN+\lfloor b^m\theta\rfloor\), y el segundo sumando
pertenece a \(\{0,\ldots,b^m-1\}\). La división euclídea por \(b^m\)
recupera \(N\). La especialización \(m=1\) produce
\eqref{act21:eq:df-paso-posicional} y su cota residual.
\end{proof}

\subsection{Bloques regionales de seis trits a profundidad arbitraria}
\label{act21:subsec:df-bloques-arbitrarios}

% Fuente: deltas/integracion/JointRegionalFrontier.lean,
% publish_base729_eq_base3, generated_block_from_longer_prefix,
% generated_block_eq_finite, generated_signature_eq_finite.

La igualdad entera \(729=3^6\) determina la agrupación de seis posiciones
ternarias en un bloque. Para \(t\geq0\), definimos
\begin{equation}
 u_c(t)=P_c(729,t+1)\bmod729,
 \qquad
 B_{c,j}(t)=
 \left\lfloor\frac{u_c(t)}{3^{5-j}}\right\rfloor\bmod3,
 \quad 0\leq j<6.
 \label{act21:eq:df-bloque-y-banda}
\end{equation}
De este modo, \(B_c(t)\) es una palabra ordenada de seis trits y
\begin{equation}
 u_c(t)=\sum_{j=0}^{5}B_{c,j}(t)3^{5-j}.
 \label{act21:eq:df-elevacion-banda}
\end{equation}
La palabra y su elevación entera son dos representaciones mutuamente
recuperables del mismo bloque emitido.

\begin{theorem}[Estabilidad de extracción a cualquier profundidad]
\label{act21:thm:df-bloque-estable}
Se cumplen, para todo \(n\),
\[
 P_c(729,n)=P_c(3,6n).
\]
Si \(t<n\), entonces
\begin{equation}
 u_c(t)=
 \left\lfloor\frac{P_c(729,n)}{729^{n-t-1}}\right\rfloor\bmod729.
 \label{act21:eq:df-bloque-prefijo-largo}
\end{equation}
En consecuencia, cualquier ampliación del prefijo preserva todos los
bloques previamente emitidos y todas las firmas que se calculan sobre ellos.
\end{theorem}

\begin{proof}
La primera igualdad resulta de que ambos prefijos se separan a la misma
escala, \(729^n=3^{6n}\), y del
teorema~\ref{act21:thm:df-emision-correcta}. Para la segunda aplicamos
\eqref{act21:eq:df-truncamiento} a las profundidades \(n\) y \(t+1\), y después
tomamos el residuo módulo \(729\). La extracción ternaria de
\eqref{act21:eq:df-bloque-y-banda} es determinista. Cualquier función explícita
de esas dieciocho coordenadas conserva, por tanto, el mismo resultado al
ampliar el prefijo.
\end{proof}

La construcción está cuantificada sobre \(t\in\mathbb N\). Una tabla que
materialice cien bloques constituye una evaluación finita de esta
construcción; la fórmula \eqref{act21:eq:df-bloque-prefijo-largo} determina
exactamente su compatibilidad con cualquier extensión posterior.

\begin{example}[Primer bloque conjunto]
La evaluación de los lectores regionales seleccionados da las palabras
\[
 B_{\mathrm P}(0)=010211,\qquad
 B_{\mathrm E}(0)=201101,\qquad
 B_{\mathrm F}(0)=121200.
\]
Su elevación por \eqref{act21:eq:df-elevacion-banda} produce
\[
 u_{\mathrm P}(0)=103,\qquad
 u_{\mathrm E}(0)=523,\qquad
 u_{\mathrm F}(0)=450.
\]
Por ejemplo,
\(201101_3=2\cdot243+27+9+1=523\).
Las seis posiciones, incluidos los ceros iniciales, se conservan en la
representación de banda. El cero inicial de \(010211\) tiene una posición
definida aunque la escritura decimal de \(103\) carezca de él.
\end{example}

\subsection{Registros de transición y reconstrucción de las elevaciones lineales}
\label{act21:subsec:df-registros-transicion}

% Fuentes: deltas/integracion/RegionalTransitionRecords.lean y
% GeneratedTransitionRecords.lean; lean/biblioteca/TPKLifts.lean.
% Los enlaces reconstruidos son R(0)->R(1) y R(2)->R(3).

Sobre \(\mathbb F_3\), reunimos dos emisiones consecutivas de cada región
en la matriz
\begin{equation}
 R(t)=
 \begin{pmatrix}
 B_{\mathrm P}(t)\\ B_{\mathrm P}(t+1)\\
 B_{\mathrm E}(t)\\ B_{\mathrm E}(t+1)\\
 B_{\mathrm F}(t)\\ B_{\mathrm F}(t+1)
 \end{pmatrix}.
 \label{act21:eq:df-registro-matricial}
\end{equation}
Esta definición fija tanto el orden regional como el orden temporal. Los
cuatro primeros registros son
\[
 X_0=R(0),\quad Y_0=R(1),\quad X_1=R(2),\quad Y_1=R(3).
\]
Sus valores, obtenidos por extracción de los prefijos regionales, son
\begin{align*}
 X_0&=\begin{pmatrix}
0&1&0&2&1&1\\0&1&2&2&2&2\\2&0&1&1&0&1\\
1&2&1&2&2&1\\1&2&1&2&0&0\\1&1&2&2&0&2
\end{pmatrix},&
 Y_0&=\begin{pmatrix}
0&1&2&2&2&2\\0&1&0&2&1&1\\1&2&1&2&2&1\\
1&0&2&0&1&1\\1&1&2&2&0&2\\1&2&1&0&2&0
\end{pmatrix},\\[1ex]
 X_1&=\begin{pmatrix}
0&1&0&2&1&1\\0&0&2&1&1&1\\1&0&2&0&1&1\\
0&1&2&2&2&2\\1&2&1&0&2&0\\0&1&0&2&1&0
\end{pmatrix},&
 Y_1&=\begin{pmatrix}
0&0&2&1&1&1\\1&1&0&2&2&1\\0&1&2&2&2&2\\
1&0&2&0&1&1\\0&1&0&2&1&0\\0&1&0&2&0&0
\end{pmatrix}.
\end{align*}

\begin{theorem}[Reconstrucción única de los dos enlaces iniciales]
\label{act21:thm:df-elevaciones-reconstruidas}
Las ecuaciones \(X_0L_0=Y_0\) y \(X_1L_1=Y_1\) tienen, cada una, una
única solución en \(M_6(\mathbb F_3)\). Estas soluciones son
\begin{align*}
 L_0&=\begin{pmatrix}
2&2&2&1&2&1\\2&1&2&2&1&1\\1&0&0&1&0&2\\
0&0&1&1&1&0\\2&2&2&0&2&1\\2&1&2&1&0&0
\end{pmatrix},&
 L_1&=\begin{pmatrix}
2&1&2&0&2&2\\1&2&1&1&2&0\\1&2&1&1&1&2\\
0&1&0&0&2&1\\2&0&2&1&0&1\\0&2&2&2&1&1
\end{pmatrix}.
\end{align*}
Ambas matrices son invertibles.
\end{theorem}

\begin{proof}
Las matrices
\begin{align*}
 A_0&=\begin{pmatrix}
1&1&2&0&1&2\\0&0&1&0&0&1\\2&1&2&1&0&1\\
0&2&0&1&0&2\\2&1&1&0&2&2\\2&1&1&1&1&2
\end{pmatrix},&
 A_1&=\begin{pmatrix}
1&0&1&2&0&0\\0&0&1&1&2&2\\0&1&2&0&1&1\\
1&1&0&2&0&0\\1&1&2&1&1&2\\1&0&0&0&0&2
\end{pmatrix}
\end{align*}
satisfacen \(A_iX_i=X_iA_i=I\), por multiplicación en \(\mathbb F_3\).
Por tanto, las soluciones están forzadas por \(L_i=A_iY_i\); la
multiplicación produce las dos matrices mostradas. Sus inversas son
\begin{align*}
 L_0^{-1}&=\begin{pmatrix}
1&2&0&2&0&2\\2&0&2&2&0&0\\2&1&2&0&2&0\\
1&0&0&0&2&0\\0&2&1&1&2&0\\2&2&2&2&2&2
\end{pmatrix},&
 L_1^{-1}&=\begin{pmatrix}
2&0&2&1&2&1\\2&1&0&0&2&0\\2&0&2&2&0&2\\
2&0&0&1&1&0\\1&1&1&1&1&0\\2&0&1&2&2&0
\end{pmatrix}.
\end{align*}
La verificación bilateral \(L_iL_i^{-1}=L_i^{-1}L_i=I\) demuestra la
recuperación exacta de cada palabra transformada. Todos los productos
utilizan residuos \(0,1,2\), con su interpretación en \(\mathbb F_3\).
\end{proof}

La lectura causal de este resultado es precisa: las emisiones regionales
producen \(R(t)\); los registros \(R(0),\ldots,R(3)\) determinan las dos
elevaciones anteriores; sus inversas recuperan las palabras transformadas.
La continuidad para todo \(t\) pertenece a los lectores coinductivos y al
teorema~\ref{act21:thm:df-bloque-estable}. La reconstrucción matricial aquí
demostrada tiene por dominio los dos enlaces especificados en su enunciado.

\subsection{Refinamiento simultáneo en las bases ternaria y decimal}
\label{act21:subsec:df-cilindros-mixtos}

% Fuente: deltas/integracion/RegionalCylinderTransition.lean.

Una emisión ternaria de seis trits y una emisión decimal de tres cifras
constituyen dos lecturas de distinta base. Para conservar su relación se
registra el estado
\[
 s=(n,k,N,D),
\]
donde \(N\) es un prefijo a profundidad \(n\) en base \(729\), y \(D\) es
un prefijo a profundidad \(k\) en base \(1000\). Sus cilindros
arquimedianos son
\begin{equation}
 I_{729}(s)=\left[\frac{N}{729^n},\frac{N+1}{729^n}\right),\qquad
 I_{1000}(s)=\left[\frac{D}{1000^k},\frac{D+1}{1000^k}\right).
 \label{act21:eq:df-cilindros}
\end{equation}
El extremo superior semiabierto asigna cada valor a una única celda de una
base y una profundidad dadas.

\begin{definition}[Defectos enteros de compatibilidad]
Definimos
\begin{align}
 A(s)&=N1000^k-D729^n,\label{act21:eq:df-defecto-inferior}\\
 B(s)&=(N+1)1000^k-D729^n=A(s)+1000^k.
 \label{act21:eq:df-defecto-superior}
\end{align}
El estado es compatible cuando
\begin{equation}
 A(s)<729^n,\qquad B(s)>0.
 \label{act21:eq:df-compatibilidad}
\end{equation}
\end{definition}

\begin{lemma}[Criterio entero de intersección]
Los cilindros de \eqref{act21:eq:df-cilindros} tienen intersección no vacía si y
sólo si se cumple \eqref{act21:eq:df-compatibilidad}.
\end{lemma}

\begin{proof}
Dos intervalos semiabiertos \([a,b)\) y \([c,d)\), con \(a<b\) y
\(c<d\), se cortan exactamente cuando \(a<d\) y \(c<b\). Aplicando esta
condición a \eqref{act21:eq:df-cilindros} y multiplicando por los denominadores
positivos, obtenemos
\(N1000^k<(D+1)729^n\) y
\(D729^n<(N+1)1000^k\). Son las dos desigualdades de
\eqref{act21:eq:df-compatibilidad}.
\end{proof}

\begin{definition}[Actualización de profundidades y prefijos]
Para \(0\leq u<729\) y \(0\leq d<1000\), definimos
\begin{align*}
 T_u(n,k,N,D)&=(n+1,k,729N+u,D),\\
 T_{u,d}(n,k,N,D)&=(n+1,k+1,729N+u,1000D+d).
\end{align*}
La primera operación refina únicamente el cilindro ternario. La segunda
refina ambas lecturas y conserva las dos profundidades en el estado.
\end{definition}

\begin{proposition}[Recurrencias exactas de frontera]
\label{act21:prop:df-rec-frontera}
Las actualizaciones anteriores satisfacen
\begin{align}
 A(T_u s)&=729A(s)+u1000^k,\label{act21:eq:df-rec-ternaria}\\
 A(T_{u,d}s)&=729000A(s)+u1000^{k+1}-d729^{n+1}.
 \label{act21:eq:df-rec-conjunta}
\end{align}
\end{proposition}

\begin{proof}
Para la primera identidad,
\[
 (729N+u)1000^k-D729^{n+1}
 =729(N1000^k-D729^n)+u1000^k.
\]
Para la segunda,
\begin{align*}
 &(729N+u)1000^{k+1}-(1000D+d)729^{n+1}\\
 &\qquad=729000(N1000^k-D729^n)
      +u1000^{k+1}-d729^{n+1}.
\end{align*}
Ambas son identidades enteras anteriores a cualquier aproximación real.
\end{proof}

\begin{theorem}[Refinamiento de los cilindros regionales generados]
\label{act21:thm:df-refinamiento-generado}
Sea
\[
 s_c(n,k)=(n,k,P_c(729,n),P_c(1000,k)).
\]
Entonces \(s_c(n,k)\) es compatible para cualesquiera \(n,k\). Además,
\begin{align*}
 T_{u_c(n)}s_c(n,k)&=s_c(n+1,k),\\
 T_{u_c(n),d_c(1000,k)}s_c(n,k)&=s_c(n+1,k+1).
\end{align*}
La división euclídea de los prefijos refinados recupera exactamente los
prefijos anteriores. Para todo \(a,b\geq0\),
\[
 \frac{P_c(729,n+a)-P_c(729,n+a)\bmod729^a}{729^a}
 =P_c(729,n),
\]
y se cumple la identidad análoga en base \(1000\).
\end{theorem}

\begin{proof}
El teorema~\ref{act21:thm:df-emision-correcta} sitúa \(v_c\) simultáneamente en
los dos intervalos \eqref{act21:eq:df-cilindros}; el criterio de intersección da
la compatibilidad. Las identidades de actualización son
\eqref{act21:eq:df-paso-posicional} en cada base. La recuperación se obtiene de
\eqref{act21:eq:df-truncamiento}. Por tanto, la recurrencia del defecto y la
recuperación del prefijo son consecuencias de una misma actualización.
\end{proof}

\begin{example}[Primer par de cilindros regionales]
Para \(n=k=1\), la evaluación de las tres regiones da
\[
\begin{array}{c|rr|rr}
 c&N&D&A&B\\\hline
 \mathrm P&2290&3141&211&1211\\
 \mathrm E&1981&2718&-422&578\\
 \mathrm F&1179&1618&-522&478
\end{array}
\]
En las tres filas \(A<729\) y \(B>0\). La primera columna de prefijos
se descompone como \(2290=3\cdot729+103\),
\(1981=2\cdot729+523\) y \(1179=729+450\): reaparecen exactamente los
bloques de seis trits ya emitidos. Los signos diferentes de \(A\) expresan
la posición relativa de los dos extremos inferiores, sin alterar la
pertenencia común del valor regional.
\end{example}

\subsection{Selección del bloque siguiente mediante el lector regional}

La compatibilidad de dos intervalos expresa una relación entre lecturas.
La determinación del siguiente bloque utiliza además el intervalo racional
producido por la región. Esta operación puede formularse sin eliminar la
historia de profundidad.

\begin{definition}[Admisibilidad de la prolongación regional]
Fijados \(c,n\), sea \(s=729^{n+1}\) y \(j=j_c(s)\). Un residuo
\(u\in\{0,\ldots,728\}\) es admisible cuando
\begin{equation}
 \lfloor s\ell_c(j)\rfloor
 \leq729P_c(729,n)+u
 \leq\lceil su_c(j)\rceil-1.
 \label{act21:eq:df-admisibilidad}
\end{equation}
\end{definition}

\begin{theorem}[Unicidad de la prolongación compatible con el lector]
\label{act21:thm:df-prolongacion-unica}
La condición \eqref{act21:eq:df-admisibilidad} equivale a \(u=u_c(n)\).
Por tanto, existe exactamente una prolongación admisible de cada región
a cada profundidad.
\end{theorem}

\begin{proof}
En el índice de separación, los dos extremos enteros de
\eqref{act21:eq:df-admisibilidad} coinciden con \(P_c(729,n+1)\).
La condición equivale así a
\(729P_c(729,n)+u=P_c(729,n+1)\).
Por \eqref{act21:eq:df-paso-posicional}, el lado derecho es
\(729P_c(729,n)+u_c(n)\). La cancelación del término común da la
igualdad de residuos. Recíprocamente, esa igualdad satisface la condición.
\end{proof}

La dependencia efectiva de esta selección queda completamente determinada:
región seleccionada, lector racional, profundidad, índice de separación,
prefijo anterior y residuo siguiente. El defecto \(A\) resume una relación
de frontera; el estado \((n,k,N,D)\) y los intervalos regionales conservan
los datos empleados por la actualización.

\subsection{Firma ternaria, acarreo, incidencia de Witt y reloj de lectura}
\label{act21:subsec:df-firma-conjunta}

% Fuentes: lean/regional/N69RegionalSignature.lean,
% deltas/integracion/JointRegionalFrontier.lean y JointCylinderSignature.lean.

Fijada una profundidad \(t\), escribamos \(B_{c,j}=B_{c,j}(t)\). Para
cada columna \(j\) se calculan los siguientes enteros y residuos:
\begin{align}
 S_j&=B_{\mathrm P,j}+B_{\mathrm E,j}+B_{\mathrm F,j},
 &q_j&=S_j\bmod3,
 &c_j&=\left\lfloor S_j/3\right\rfloor,
 \label{act21:eq:df-firma-qc}\\
 a_j&=(B_{\mathrm P,j}+B_{\mathrm E,j}-B_{\mathrm F,j})\bmod3,
 &w_j&=\sum_c\mathbf1_{B_{c,j}\ne0},
 &\bar w_j&=w_j\bmod3.
 \label{act21:eq:df-firma-axial}
\end{align}
El residuo de la expresión axial se toma en los enteros. Una implementación
con naturales lo calcula como
\((B_{\mathrm P,j}+B_{\mathrm E,j}+3-B_{\mathrm F,j})\bmod3\), que conserva
exactamente ese residuo.

La matriz de incidencia ternaria utilizada por el lector es
\begin{equation}
 W=\begin{pmatrix}
0&1&1&1&1&1\\
1&0&1&1&2&2\\
1&1&0&2&1&2\\
1&1&2&0&2&1\\
1&2&1&2&0&1\\
1&2&2&1&1&0
\end{pmatrix}.
\label{act21:eq:df-matriz-witt}
\end{equation}
Para cada región se conservan dos cargas residuales:
\begin{equation}
 v_c^{\mathrm{res}}=\sum_j B_{c,j}\bmod3,
 \qquad
 v_c^{\mathrm W}=\sum_j(B_cW)_j\bmod3.
 \label{act21:eq:df-cargas-residuales}
\end{equation}
La firma completa del lector es
\(\mathcal S(B)=(q,a,c,w,\bar w,v^{\mathrm{res}},v^{\mathrm W})\).
Su estructura conserva por separado el residuo de columna, el acarreo, el
contraste axial, el peso de soporte y las dos lecturas de carga.

\begin{theorem}[Recuperación exacta y cotas de la firma]
\label{act21:thm:df-recuperacion-firma}
Para cada columna y cada región se cumplen
\begin{align}
 S_j&=q_j+3c_j,&0\leq q_j,a_j&<3,&0\leq c_j&\leq2,
 \label{act21:eq:df-recuperacion-total}\\
 B_{\mathrm F,j}&=2(q_j+3-a_j)\bmod3,&0\leq w_j&\leq3,
 &0\leq\bar w_j&<3.
 \label{act21:eq:df-recuperacion-eje}
\end{align}
Además,
\begin{equation}
 v_c^{\mathrm W}
 =\left(\sum_j\bigl((B_cW)_j\bmod3\bigr)\right)\bmod3.
 \label{act21:eq:df-witt-reduccion}
\end{equation}
Las recuperaciones enunciadas determinan el total de cada columna y su
coordenada de autoescala; el registro regional conserva adicionalmente
las tres bandas ordenadas.
\end{theorem}

\begin{proof}
La identidad \eqref{act21:eq:df-recuperacion-total} es la división euclídea de
\(S_j\) por \(3\). Cada sumando de \(S_j\) está entre \(0\) y \(2\),
de modo que \(0\leq S_j\leq6\) y \(c_j\leq2\). En \(\mathbb F_3\),
la resta de las dos congruencias que definen \(q_j\) y \(a_j\) da
\(q_j-a_j=2B_{\mathrm F,j}\). Como el inverso de \(2\) es \(2\), se
obtiene \eqref{act21:eq:df-recuperacion-eje}; el sumando \(3\) permite
escribirla con naturales. Las cotas de \(w_j\) resultan de sumar tres
indicadores. Finalmente, reducir una suma entera módulo \(3\) equivale
a reducir cada sumando y volver a reducir la suma, lo que prueba
\eqref{act21:eq:df-witt-reduccion}.
\end{proof}

\begin{example}[Firma del primer bloque conjunto]
Las tres bandas iniciales producen
\begin{align*}
 q&=(0,0,2,2,1,2),&a&=(1,2,0,1,1,2),\\
 c&=(1,1,0,1,0,0),&w&=(2,2,2,3,1,2),\\
 v^{\mathrm{res}}&=(2,2,0),&v^{\mathrm W}&=(2,1,1).
\end{align*}
La cuarta columna tiene total \(2+1+2=5\); su firma registra
\(q_3=2\) y \(c_3=1\), por lo que \(q_3+3c_3=5\).
La recuperación axial da
\(2(2+3-1)\bmod3=2\), precisamente la coordenada de autoescala.
Las imágenes de Witt, reducidas componente a componente, son
\[
 (2,0,2,1,1,2),\qquad
 (0,0,0,2,0,2),\qquad
 (2,1,1,2,1,0).
\]
Su suma residual por fila produce las tres cargas \((2,1,1)\).
\end{example}

\begin{definition}[Reloj entero de lectura decimal]
El índice de transición \(t\) y la profundidad decimal de representación se
relacionan mediante
\begin{equation}
 \kappa(t)=\max\{k\in\mathbb N:1000^k\leq3^{6(t+1)}\}.
 \label{act21:eq:df-reloj}
\end{equation}
\end{definition}

\begin{proposition}[Caracterización y monotonía del reloj]
\label{act21:prop:df-reloj}
El entero \(\kappa(t)\) queda determinado unívocamente por
\begin{equation}
 1000^{\kappa(t)}\leq729^{t+1}<1000^{\kappa(t)+1}.
 \label{act21:eq:df-reloj-cotas}
\end{equation}
La función \(\kappa\) es monótona y satisface
\(\kappa(20)=\kappa(21)=20\).
\end{proposition}

\begin{proof}
Las potencias de \(1000\) son estrictamente crecientes y no acotadas;
existe, por tanto, un único intervalo consecutivo de potencias que
contiene \(729^{t+1}\). La monotonía de esta última expresión demuestra
la de \(\kappa\). Las dos evaluaciones finales son las comparaciones
enteras
\[
 1000^{20}\leq729^{21}<1000^{21},\qquad
 1000^{20}\leq729^{22}<1000^{21}.
\]
Estas comparaciones se efectúan con potencias enteras exactas.
\end{proof}

La igualdad de profundidad decimal en dos transiciones distintas explica
la necesidad de conservar ambos índices: el reloj de lectura puede
permanecer constante mientras avanza la banda ternaria y se actualiza su
firma. La información cronológica se encuentra en el registro ordenado
de las transiciones.

\subsection{Teorema de composición del bloque, el cilindro y la firma}

\begin{theorem}[Compatibilidad conjunta de las lecturas regionales]
\label{act21:thm:df-composicion-conjunta}
Para cada profundidad \(n\), existe una única terna de bloques
\((u_{\mathrm P},u_{\mathrm E},u_{\mathrm F})\) que satisface la
admisibilidad regional \eqref{act21:eq:df-admisibilidad} en las tres regiones.
Esta terna es \((u_{\mathrm P}(n),u_{\mathrm E}(n),u_{\mathrm F}(n))\).
Simultáneamente:
\begin{enumerate}
\item cada bloque actualiza su cilindro por las recurrencias
\eqref{act21:eq:df-rec-ternaria} y \eqref{act21:eq:df-rec-conjunta};
\item sus seis trits producen la firma
\(\mathcal S(B_{\mathrm P}(n),B_{\mathrm E}(n),B_{\mathrm F}(n))\);
\item los totales de columna se recuperan como \(q+3c\), y las cargas
de Witt se obtienen de las mismas bandas mediante
\eqref{act21:eq:df-witt-reduccion};
\item todo prefijo de mayor profundidad devuelve exactamente estos
bloques, cilindros truncados y firmas.
\end{enumerate}
\end{theorem}

\begin{proof}
Aplicamos el teorema~\ref{act21:thm:df-prolongacion-unica} a cada uno de los tres
índices regionales. La unicidad coordenada a coordenada da la unicidad
conjunta. Sustituimos esos mismos residuos en las operaciones
\(T_u,T_{u,d}\), lo que permite aplicar el
teorema~\ref{act21:thm:df-refinamiento-generado}. La aplicación determinista de
\eqref{act21:eq:df-bloque-y-banda} produce las bandas utilizadas por
\eqref{act21:eq:df-firma-qc}--\eqref{act21:eq:df-cargas-residuales}. Las identidades
de recuperación son las del teorema~\ref{act21:thm:df-recuperacion-firma}.
Finalmente, el teorema~\ref{act21:thm:df-bloque-estable} y el truncamiento de
cada base demuestran la estabilidad a mayor profundidad. En cada paso se
ha conservado la identidad del bloque emitido; ninguna de las dos
lecturas requiere seleccionar otro bloque.
\end{proof}

La conclusión es una compatibilidad operacional entre la contracción
cilíndrica y la conservación de incidencia. Ambas utilizan el mismo
registro generado, mantienen las profundidades y admiten recuperación
entera. Este resultado es la forma explícita del enlace
\[
 \text{bloque regional emitido}
 \longrightarrow
 \begin{cases}
 \text{cilindro y frontera entera},\\
 \text{banda, firma, acarreo y carga de Witt}.
 \end{cases}
\]
El desarrollo posterior de la incidencia y de los campos emplea la
segunda lectura, conservando la procedencia común con la primera.

\subsection{Censo de calendarios y compatibilidad simultánea de las biografías regionales}
Las filas anteriores se reúnen, sin aplicar todavía ningún lector decimal,
en las tres biografías de bloques producidas por las transiciones:
\begin{align}
 B^{\rm cl}
  &=010211|012222|010211|002111|110221,\notag\\
 B^{\rm pr}
  &=201101|121221|102011|012222|102011,
 \label{act21:eq:biografias-tpk-tres-caracteres}\\
 B^{\rm au}
  &=121200|112202|121020|010210|010200.\notag
\end{align}
En efecto, los pares consecutivos de los dos primeros tramos son las filas
de \(X_0,Y_0\), y los de los dos últimos son las filas de \(X_1,Y_1\).
Así, las biografías son otra escritura del registro de doce transiciones, no
una colección de prefijos convencionales.

Para comprobar la palabra operativa del calendario, sea
\(c=c_1c_2c_3c_4\in\{0,1\}^4\) y, para cada régimen
\(\chi\in\{{\rm cl},{\rm pr},{\rm au}\}\), defínase
\[
 u_0^\chi=b_0^\chi,\qquad
 u_j^\chi=u_{j-1}^\chi L_{c_j}\quad(1\le j\le4).
\]
Para condensar el censo, denotemos por
\[
 N_{\mathrm{ar}}(c):=
 \#\{(j,\chi):u_j^\chi=b_j^\chi,\ 1\le j\le4\},\qquad
 N_{\mathrm{bio}}(c):=
 \#\{\chi:(u_0^\chi,\ldots,u_4^\chi)=B^\chi\}.
\]
La multiplicación exacta en \(\FF_3\) de los dieciséis calendarios da:
\[
\begin{array}{c|c|c}
 c&N_{\mathrm{ar}}(c)&N_{\mathrm{bio}}(c)\\ \hline
0000,0001&6&0\\
0010&9&0\\
0011&12&3\\
0100,0101,0110,0111&3&0\\
1000,1001,1010,1011&0&0\\
1100,1101,1110,1111&0&0
\end{array}
\tag{D}
\label{act21:eq:censo-dieciseis-calendarios}
\]
La tabla agota \(2^4=16\) posibilidades y cada una de sus entradas se obtiene
por cuatro multiplicaciones de una fila por una de las dos matrices impresas.
Por consiguiente, \(0011\) es el único calendario que reconstruye
simultáneamente las tres biografías completas. En notación operatoria,
\begin{equation}
 (L_0,L_0,L_1,L_1)
 \label{act21:eq:calendario-0011-interno}
\end{equation}
no es una prescripción independiente: es la única lectura binaria compatible
con las doce aristas ya producidas por \(U_t\).

El calendario \eqref{act21:eq:calendario-0011-interno} produce cuatro nuevos
bloques de seis componentes. Sus concatenaciones forman
\[
 w_6\longrightarrow w_{12}\longrightarrow w_{18}
 \longrightarrow w_{24}\longrightarrow w_{30}.
\]
Los subíndices de \(w\) indican longitud acumulada. La frontera \(R_{36}\) conserva la estructura terminal y abre la relación de prolongación; no se renombra como un último bloque periódico que pudiera repetirse indefinidamente.

La identidad \(L=X^{-1}Y\) demuestra unicidad una vez fijadas las
transiciones; no sustituye su producción por \(U_t\). Recíprocamente,
\eqref{act21:eq:censo-dieciseis-calendarios} prueba dentro del artículo la unicidad
del orden de los cuatro transportes. La cadena causal queda así cerrada:
el censo APP produce las regiones iniciales, el TRIT fija régimen y
orientación, el TPK produce el registro de transiciones, y ese registro
determina \(L_0,L_1\) y \(0011\) antes de toda realización arquimediana.



% Desarrollo reunido K31-07--K31-15. Inserción anterior a registro_k y alpha.
% Mantiene explícita la interfaz del repertorio terminal de ocho bloques.
% Propietarios y localizadores: metadata/K_ALPHA_PROPIETARIOS.md.
\section{Determinación regional del descriptor dodecafásico y selección del registro de acoplamiento}
\label{act21:chap:despliegue-selector-k}

La determinación del registro de acoplamiento articula dos construcciones
que conviene desplegar en su orden efectivo. La primera produce los
antecedentes del selector a partir de las representaciones regionales del TPK:
bandas de seis trits, firmas de columna, calendario de conversión,
descriptor orientado y panel funcional. La segunda opera sobre esos
antecedentes y sobre el repertorio terminal de ocho bloques, enumera sus
ordenaciones admisibles y selecciona un único registro mediante dos
condiciones simultáneas: incidencia excepcional y heterotipia de lectura.
El valor del registro aparece al término de esta composición.

El presente desarrollo presupone las representaciones regionales obtenidas
en los capítulos precedentes desde APP--TRIT--TPK. Su propósito es hacer
explícita cada operación intermedia que conduce desde tales representaciones
al selector finito. Después se incorporan la realización del mismo
registro por incidencias firmadas, la inversión de Hadamard y la
representación de la constante de estructura fina. Las tres construcciones
conservan sus respectivos dominios: selección del registro, recuperación
desde sus canales y evaluación aritmética de sus representaciones.

\subsection{Bandas regionales, codificación y firmas enteras}

\subsubsection{Dominios y codificación de las bandas regionales}

Se fija el orden de canales
\[
 \chi\in\{\pi,e,\varphi\},
 \qquad B_\chi(t)\in\mathbb F_3^6,
 \qquad t\in\mathbb N.
\]
Los símbolos de canal identifican las regiones ya generadas. Cada
$B_\chi(t)$ conserva el lugar que ocupa en su sucesión de representaciones;
el índice $t$ cuenta bloques ternarios. La evaluación natural de una
banda $b=(b_0,\ldots,b_5)$ es
\begin{equation}
 \operatorname{enc}_6(b)=
 243b_0+81b_1+27b_2+9b_3+3b_4+b_5,
 \qquad 0\leq\operatorname{enc}_6(b)<729.
 \label{act21:eq:kdes-enc6}
\end{equation}
Su inversa, con anchura fijada, es
\[
 \operatorname{dig}_j(n)=
 \left\lfloor\frac{n}{3^{5-j}}\right\rfloor\bmod3,
 \qquad 0\leq j<6.
\]
Así se retienen los ceros iniciales: las palabras $002001$ y $2001$
podrían compartir una escritura entera abreviada, pero el tipo de la
primera contiene seis posiciones. Todas las concatenaciones que siguen
utilizan esa anchura.

En la ventana de cien bloques, la representación de seiscientos trits del
canal $\chi$ determina un entero $P_\chi^{(600)}$. La extracción se efectúa
mediante
\begin{equation}
 b_\chi(t)=
 \left\lfloor\frac{P_\chi^{(600)}}{729^{99-t}}\right\rfloor\bmod729,
 \qquad 0\leq t<100,
 \qquad B_\chi(t)=\operatorname{enc}_6^{-1}(b_\chi(t)).
 \label{act21:eq:kdes-extraccion}
\end{equation}
El entero de seiscientos trits es una representación del productor regional;
la tabla de cien filas es la evaluación de
\eqref{act21:eq:kdes-extraccion}. Esta separación permite sustituir una tabla
histórica por su productor sin modificar el selector que la consume.

\subsubsection{Residuos, cocientes, ocupación y orientación}

Para cada columna $j$ y tiempo $t$, se introducen
\begin{align}
 T_j(t)&=B_\pi(t)_j+B_e(t)_j+B_\varphi(t)_j,\\
 q_j(t)&=T_j(t)\bmod3,
 &c_j(t)&=\left\lfloor T_j(t)/3\right\rfloor,\\
 a_j(t)&=\bigl(B_\pi(t)_j+B_e(t)_j-B_\varphi(t)_j\bigr)\bmod3,\\
 w_j(t)&=\mathbf1_{B_\pi(t)_j\ne0}
       +\mathbf1_{B_e(t)_j\ne0}
       +\mathbf1_{B_\varphi(t)_j\ne0},
 &\bar w_j(t)&=w_j(t)\bmod3.
 \label{act21:eq:kdes-firmas}
\end{align}
La resta de la tercera expresión se realiza en los enteros y se reduce
después. En una implementación sobre naturales, la expresión equivalente
es $(B_\pi+B_e+3-B_\varphi)\bmod3$; la resta truncada alteraría el lector.

\begin{proposition}[Reconstrucción de la suma de columna]
Para todas las bandas admisibles,
\[
 T_j(t)=q_j(t)+3c_j(t),\qquad
 0\leq q_j(t)<3,\quad 0\leq c_j(t)\leq2,
 \quad0\leq w_j(t)\leq3.
\]
\end{proposition}
\begin{proof}
Cada columna suma tres representantes de $\{0,1,2\}$, de modo que
$0\leq T_j(t)\leq6$. La división euclídea por tres proporciona el residuo
y el cociente indicados. La ocupación cuenta tres condiciones binarias;
su cota se deduce directamente de esa definición.
\end{proof}

Esta identidad conserva el cociente de la suma local. La reconstrucción
de las tres bandas individuales requiere las bandas mismas o un lector
adicional que las distinga: una suma de columna tiene menos coordenadas
que el triple que la produce. Por ello, la fila temporal utilizada más
adelante conserva conjuntamente las tres bandas y los cuatro campos
$q,a,c,\bar w$.

\subsection{Reloj de conversión y primer avance nulo}

\subsubsection{Capacidad ternaria y profundidad decimal}

Una representación de $t+1$ bandas contiene $6(t+1)$ trits. El calendario
de conversión asociado se define por el entero
\begin{equation}
 d_t=\max\{k\in\mathbb N:1000^k\leq3^{6(t+1)}\}
     =\max\{k\in\mathbb N:1000^k\leq729^{t+1}\}.
 \label{act21:eq:kdes-reloj}
\end{equation}
La definición compara capacidades enteras. El contador $t$ y la
profundidad $d_t$ registran magnitudes distintas: una transición añade
una banda de seis trits; el reloj cuenta potencias completas de la base
decimal mil contenidas en la capacidad ternaria acumulada.

\begin{lemma}[Caracterización entera del reloj]
Para $t,k\in\mathbb N$ se tiene
\[
 d_t=k\quad\Longleftrightarrow\quad
 1000^k\leq729^{t+1}<1000^{k+1}.
\]
Además, $d_t\leq d_{t+1}$.
\end{lemma}
\begin{proof}
Las potencias de mil forman una sucesión estrictamente creciente y no
acotada. La potencia $729^{t+1}$ queda entre dos términos consecutivos;
el exponente inferior es exactamente el máximo de
\eqref{act21:eq:kdes-reloj}. La desigualdad
$729^{t+1}<729^{t+2}$ implica la monotonía del máximo.
\end{proof}

\begin{proposition}[Primer avance nulo]
El menor tiempo que satisface $d_{t+1}=d_t$ es
\[
 t_*=20.
\]
En particular,
\[
 d_0=0,\quad d_{19}=19,\quad d_{20}=20,
 \quad d_{21}=20,\quad d_{22}=21.
\]
\end{proposition}
\begin{proof}
Las comparaciones enteras exactas
\[
 1000^{20}\leq729^{21},\qquad
 729^{22}<1000^{21},\qquad
 1000^{21}\leq729^{23}<1000^{22}
\]
dan los valores en los tiempos veinte, veintiuno y veintidós. Para
$0\leq t\leq20$, la expresión
$729^{t+1}/1000^t=729(729/1000)^t$ es decreciente y su valor en veinte
es mayor o igual que uno. Por tanto,
$1000^t\leq729^{t+1}<1000^{t+1}$ y $d_t=t$. Ningún tiempo menor que
veinte puede satisfacer $d_{t+1}=d_t$; en veinte sí se cumple.
Todas las comparaciones utilizadas son de potencias enteras.
\end{proof}

El número veinte es, por tanto, una salida del criterio mínimo. La regla
que selecciona el instante es ``primer avance nulo del reloj''; el
algoritmo y su prueba pueden emplear esta regla sin recibir veinte como
parámetro externo.

\subsubsection{Retorno nonádico y antecedente temporal}

La longitud del retorno de fase determina
\begin{equation}
 r=t_*-9=11.
 \label{act21:eq:kdes-retorno}
\end{equation}
Las consultas del descriptor se efectúan en $t_*+1$, $r+1$ y $r$.
Esta distribución conserva la relación entre el instante de pausa, la
traslación nonádica y la lectura de fase inmediatamente asociada. Los
tiempos once, doce y veintiuno designan posiciones del calendario
regional; su diferencia temporal permanece registrada aun cuando dos
tiempos produzcan la misma profundidad decimal.

\subsection{Construcción orientada del descriptor de veinticuatro trits}

\subsubsection{Defectos opuestos de semibanda}

Se toman los defectos orientados de la coordenatización regional
\[
 \delta_+=(0,0,0,1,1,1),\qquad
 \delta_-=(0,0,0,2,2,2)\in\mathbb F_3^6.
\]
Su oposición se verifica componente a componente:
\[
 \delta_++\delta_-=0\quad\text{en }\mathbb F_3^6.
\]
El representante $2$ realiza $-1$ en la coordenatización ternaria. La concatenación
de las colas de tres posiciones, en el orden negativo--positivo, produce
la palabra $222111$. La suma de defectos y la concatenación de sus
colas son operaciones distintas: la primera tiene codominio
$\mathbb F_3^6$, mientras que la segunda conserva el orden de dos
segmentos de longitud tres.

\subsubsection{Carga dual y selección de fase}

En la coordenatización regional, la transformación de seis coordenadas utiliza
\[
 A_{\rm reg}=
 \begin{pmatrix}
 0&1&1&1&1&1\\
 1&0&1&1&2&2\\
 1&1&0&2&1&2\\
 1&1&2&0&2&1\\
 1&2&1&2&0&1\\
 1&2&2&1&1&0
 \end{pmatrix}.
\]
Con vectores fila, la carga dual de $b\in\mathbb F_3^6$ es la suma
de las coordenadas de $bA_{\rm reg}$. Las sumas enteras de las filas de
$A_{\rm reg}$ son $(5,7,7,7,7,7)$; de ahí resulta la expresión local
\begin{equation}
 \eta^\vee(b)=2b_0+b_1+b_2+b_3+b_4+b_5\pmod3.
 \label{act21:eq:kdes-cargadual}
\end{equation}
La fase de lectura en el instante $r$ es
$\theta_r=(r-1)\bmod3$. Para cada canal se define
\[
 \varepsilon_\chi(r)=
 \begin{cases}
 1,&\eta^\vee(B_\chi(r))=\theta_r,\\
 2,&\eta^\vee(B_\chi(r))\ne\theta_r.
 \end{cases}
\]
La representación fase--carga es
\begin{equation}
 D(r)=\varepsilon_\pi(r)B_\pi(r)
      +\varepsilon_e(r)B_e(r)
      +\varepsilon_\varphi(r)B_\varphi(r)
      \quad\text{en }\mathbb F_3^6.
 \label{act21:eq:kdes-fasecarga}
\end{equation}
El coeficiente se determina por una igualdad de cargas, antes de evaluar
la palabra que resultará de la combinación.

\subsubsection{Composición del descriptor}

Se denota por $\Vert$ la concatenación de palabras de anchura fijada.
La construcción completa es
\begin{equation}
 \begin{split}
 \Omega_{24}={}&
 (B_e(t_*+1)+\delta_-)
 \Vert (\delta_-|_{3,4,5}\Vert\delta_+|_{3,4,5})\\
 &\Vert(B_\pi(r+1)+B_e(r+1))\Vert D(r).
 \end{split}
 \label{act21:eq:kdes-omega}
\end{equation}
Cada sumando de longitud seis se evalúa en $\mathbb F_3^6$; la
concatenación final tiene longitud $6+6+6+6=24$. Esta construcción
separa el descriptor $\Omega_{24}$ de las etapas regionales
$w_{24}^{\chi}$: ambos objetos comparten longitud, pero sus productores
y sus lugares en la composición son diferentes.

\begin{proposition}[Evaluación regional del descriptor]
Las representaciones regionales y el criterio del primer avance nulo
determinan
\[
 \Omega_{24}=020022\,222111\,211201\,021101,
 \qquad [\Omega_{24}]_3=66241910521.
\]
\end{proposition}
\begin{proof}
Las bandas necesarias, conservando el orden $(\pi,e,\varphi)$, son
\[
 \begin{array}{c|ccc}
 t&B_\pi(t)&B_e(t)&B_\varphi(t)\\\hline
 11&022212&110001&100021\\
 12&012012&202222&000120\\
 21&221022&020100&211021
 \end{array}
\]
y proceden de la extracción \eqref{act21:eq:kdes-extraccion}. En el tiempo
veintiuno,
$020100+000222=020022$ en $\mathbb F_3^6$. Las colas de los defectos
aportan $222111$. En el tiempo doce,
$012012+202222=211201$.

Para el tiempo once, las cargas duales son, respectivamente, $0,1,2$;
la fase vale $(11-1)\bmod3=1$. Los coeficientes resultan $(2,1,2)$ y
\[
 2(022212)+(110001)+2(100021)=021101
 \quad\text{en }\mathbb F_3^6.
\]
La concatenación de estos cuatro resultados da la palabra anunciada.
Su evaluación entera se obtiene por Horner en base tres, o por tres
concatenaciones sucesivas en base $729$.
\end{proof}

La tabla anterior desempeña una función comprobatoria: cada una de sus
entradas se obtiene del productor regional. La regla de construcción
\eqref{act21:eq:kdes-omega} es la que permite reproducir el descriptor sin
introducir como dato la cadena final de veinticuatro trits.

\subsection{Panel funcional y repertorio terminal}

\subsubsection{Nueve posiciones regionales}

El panel funcional conserva las tres primeras bandas de cada uno de los
tres canales, en el orden regional declarado:
\begin{equation}
 \mathcal P_9=
 \bigl(b_\pi(0),b_\pi(1),b_\pi(2),
       b_e(0),b_e(1),b_e(2),
       b_\varphi(0),b_\varphi(1),b_\varphi(2)\bigr).
 \label{act21:eq:kdes-panel}
\end{equation}
Su evaluación es
\[
 \mathcal P_9=(103,161,103,523,457,301,450,398,438).
\]
Las nueve posiciones conservan procedencia regional, aunque el valor
103 aparece dos veces. El panel tiene nueve incidencias ordenadas y
ocho valores distintos. La enumeración ulterior emplea las incidencias;
la deduplicación de candidatos se efectúa en una etapa posterior.

\subsubsection{Repertorio no ordenado de ocho bloques}

El repertorio terminal que interviene en esta construcción es
\begin{equation}
 \mathcal S_{\rm term}=
 \{002001,020111,200110,121012,
   010122,001001,221110,222110\}\subset\mathbb F_3^6.
 \label{act21:eq:kdes-sterm}
\end{equation}
Sus ocho elementos son distintos. La codificación entera, ordenada por
valor únicamente para fijar una representación informática, es
\[
 \{28,55,98,175,437,498,687,714\}.
\]
La procedencia documental de este repertorio es la segmentación terminal
de la palabra de setenta y dos trits y la recuperación por posiciones
mediante incidencias de los lectores fase--carga y afín. El selector que
se desarrolla a continuación recibe el repertorio sin su ordenación
histórica y somete a prueba sus $8!$ ordenaciones.

La formulación exacta de la implementación reunida mantiene aquí una
interfaz explícita: los productores regionales sustituyen materialmente
el descriptor, el panel y las filas temporales; el repertorio
\eqref{act21:eq:kdes-sterm} conserva su productor documental anterior. Esta
distinción identifica dónde entra cada antecedente. La pertenencia de
ocho palabras a un repertorio temporal de lectores, la selección de esas
ocho palabras y la determinación de su orden terminal son tres
enunciados diferentes. La prueba de selección terminal demuestra el
tercero para el repertorio declarado.

La recuperación documental por posiciones puede expresarse de forma
exacta. Sea $\mathcal I$ el registro de incidencias fase--carga y sea
$\mathcal J$ el registro de la extensión afín. Cada incidencia conserva
la posición $\ell$ de la palabra de setenta y dos trits y la palabra
emitida $u$. Para $5\leq\ell\leq12$ se forma
\[
 \mathcal S_\ell=
 \{u:(\ell,u)\in\mathcal I\cup\mathcal J\}.
\]
El recuperador comprueba que las ocho posiciones aparecen y que cada
$\mathcal S_\ell$ es un singleton. Si $s_\ell$ es su elemento único,
la salida entregada al selector es
\[
 \mathcal S_{\rm term}=\{s_5,s_6,\ldots,s_{12}\}.
\]
La prueba de singleton por posición conserva la coherencia de los
testigos históricos; la exportación posterior conserva las ocho
palabras y suprime únicamente la ordenación histórica. El selector
reconstruye el orden mediante la enumeración completa y las condiciones
terminales, en lugar de recibirlo del registro de procedencia.

% RESULTADO_RECUPERADO: K31-10/P10 del inventario REV05.
% Fuente de las incidencias:
% 16_CIERRE_GLOBAL_HMT_MD_2026-07-22/02_LECTURA_DODECAFASICA/continuacion_fases/
% verificar_obstruccion_y_dato_minimo_w72.py, lineas 45--60 y 141--233.
% RESULTADO_OBSTRUCCION_Y_DATO_MINIMO_W72.json, phase_charge_reader.
% Fuente de las bandas: N69_signatures_t0_t99.csv, tiempos impresos abajo.
% Consumidor: PAQUETE_CONTINUIDAD_K_UNIDAD_20260918_BASE_COMPARTIDA/python/
% selector_orbital_original.py, verify_historical_inputs, lineas 159--193.
% No se transfiere la restriccion del lector comparativo al generador completo.
\subsubsection{Incidencias temporales del repertorio terminal y recuperación por posición}
\label{act21:sec:repertorio-incidencias-completas}

La representación del repertorio terminal utiliza los mismos bloques regionales,
los mismos acarreos y el mismo calendario que producen el descriptor inicial
de veinticuatro trits. Conviene explicitar por separado dos operaciones de
esta composición. La primera recupera ocho palabras a partir de incidencias
temporales registradas. La segunda presenta esas ocho palabras al selector
dodecafásico como un conjunto no ordenado. El selector vuelve a establecer
el orden mediante sus condiciones de incidencia, de cilindro y de cierre.
La distinción permite conservar los testigos de procedencia sin convertir
el orden de una segmentación histórica en una entrada del selector.

Sea $B_t=(b_{t,0},b_{t,1},b_{t,2})\in(\mathbb F_3^6)^3$ la frontera
regional ya generada en el tiempo $t$, con $0\leq t<100$. En la coordenatización
utilizada por el registro histórico, la matriz de incidencia es
\[
 A_{\mathrm{hist}}=
 \begin{pmatrix}
 0&1&1&1&1&1\\
 1&0&1&1&2&2\\
 1&1&0&2&1&2\\
 1&1&2&0&2&1\\
 1&2&1&2&0&1\\
 1&2&2&1&1&0
 \end{pmatrix}.
\]
Las palabras se consideran vectores fila. Sus dos cargas y la fase
cronológica son
\[
 r_{t,i}=\sum_{j=1}^{6}b_{t,i,j},\qquad
 d_{t,i}=\sum_{j=1}^{6}(b_{t,i}A_{\mathrm{hist}})_j,
 \qquad \vartheta_t=t-1\pmod 3.
\]
Todas las operaciones de estas tres expresiones tienen lugar en
$\mathbb F_3$. El subíndice histórico identifica una coordenatización concreta de
la incidencia de Witt; mantiene explícita la transformación de coordenadas
cuando se utiliza otra presentación de esa misma incidencia.

Para $c\in\{r_t,d_t\}$, $\delta\in\mathbb F_3$ y
$\epsilon\in\{1,-1\}$, definimos
\[
 \eta_i(c,\delta,\epsilon)=
 \begin{cases}
  \epsilon,&c_i=\vartheta_t+\delta,\\
  -\epsilon,&c_i\ne\vartheta_t+\delta,
 \end{cases}
 \qquad
 L_{t,c,\delta,\epsilon}
   =\sum_{i=0}^{2}\eta_i(c,\delta,\epsilon)b_{t,i}.
\]
La orientación $-1$ se escribe como $2$ al imprimir los coeficientes
ternarios. Junto a estos doce lectores comparativos por frontera, el
registro dispone del lector afín
\[
 L^{\mathrm{af}}_t
    =\sum_{i=0}^{2}(d_{t,i}+\vartheta_t)b_{t,i}.
\]
La adición de la fase en el lector afín y la comparación de cargas en
el lector anterior son operaciones distintas, ambas sobre datos ya
producidos por la frontera TPK.

\begin{table}[htbp]
\centering\small
\begin{tabular}{rccc cc}
\toprule
$t$&$b_{t,0}$&$b_{t,1}$&$b_{t,2}$&$r_t$&$d_t$\\
\midrule
8 &200121&212212&110110&011&202\\
11&022212&110001&100021&001&012\\
30&210112&110202&211110&100&012\\
34&121022&000011&112110&220&021\\
37&120121&011202&112221&100&201\\
49&110212&200212&100122&110&201\\
58&111121&022121&122201&122&220\\
67&010012&001220&011120&122&122\\
75&220000&020111&100002&120&021\\
81&122201&020202&201122&202&001\\
90&022112&022110&121010&202&200\\
\bottomrule
\end{tabular}
\caption{Fronteras temporales que realizan las incidencias impresas del
repertorio terminal. Cada entrada de seis cifras conserva el orden de
las seis coordenadas de la frontera regional.}
\label{act21:tab:terminal-bandas-testigo}
\end{table}

La tabla siguiente despliega todas las incidencias comparativas del
registro sobre los doce bloques de la segmentación considerada. Las
letras $r$ y $d$ designan respectivamente la carga visible y la carga
dual. La columna $j$ conserva el índice de posición del testigo histórico;
la operación de entrega al selector olvida este índice después de formar
el conjunto de palabras.

\begin{table}[htbp]
\centering\small
\begin{tabular}{rrcrrc c}
\toprule
$j$&$t$&carga&$\delta$&$\epsilon$&$(\eta_0\eta_1\eta_2)$&$L_{t,c,\delta,\epsilon}$\\
\midrule
4&11&$d$&0& 1&212&021101\\
5&67&$d$&1&$-1$&211&002001\\
5&67&$d$&2& 1&211&002001\\
5&67&$r$&1&$-1$&211&002001\\
5&67&$r$&2& 1&211&002001\\
6&81&$d$&0& 1&222&020111\\
6&81&$r$&2& 1&222&020111\\
7&34&$r$&1&$-1$&111&200110\\
7&75&$d$&1&$-1$&211&200110\\
7&75&$r$&2&$-1$&211&200110\\
8&90&$r$&0& 1&121&121012\\
8&90&$r$&1&$-1$&121&121012\\
9&49&$d$&0&$-1$&121&010122\\
11&11&$d$&2&$-1$&211&221110\\
11&37&$d$&0&$-1$&121&221110\\
12&8&$d$&0&$-1$&111&222110\\
12&8&$r$&1&$-1$&111&222110\\
12&58&$d$&1&$-1$&111&222110\\
12&58&$r$&0&$-1$&111&222110\\
\bottomrule
\end{tabular}
\caption{Incidencias fase--carga completas sobre la segmentación terminal
en las cien fronteras del registro. Las coincidencias de palabras entre
filas conservan sus tiempos, orientaciones y tipos de carga.}
\label{act21:tab:terminal-incidencias-completas}
\end{table}

El décimo bloque dispone de un testigo afín. En $t=30$ se tiene
$d_{30}=(0,1,2)$ y $\vartheta_{30}=2$, de donde
\[
 (d_{30,0}+\vartheta_{30},d_{30,1}+\vartheta_{30},
 d_{30,2}+\vartheta_{30})=(2,0,1)
\]
y, usando las bandas de la tabla~\ref{act21:tab:terminal-bandas-testigo},
\[
 L^{\mathrm{af}}_{30}
 =2(210112)+0(110202)+(211110)
 =(001001)\quad\text{en }\mathbb F_3^6.
\]
La igualdad es coordenada a coordenada: antes de reducir módulo tres,
el vector de la suma es $(6,3,1,3,3,4)$.

\begin{proposition}[Recuperación íntegra del repertorio por incidencias]
\label{act21:prop:repertorio-terminal-ocho-incidencias}
Considérese el registro de incidencias comparativas de la
tabla~\ref{act21:tab:terminal-incidencias-completas}, aumentado con el testigo
afín $(j,t,L)=(10,30,001001)$. Para cada posición $5\leq j\leq12$, el
conjunto de palabras de salida de sus incidencias es unitario. La
extracción de esas ocho palabras y el olvido posterior de su orden
producen exactamente
\[
 \mathcal S_{\mathrm{term}}=
 \{002001,020111,200110,121012,010122,001001,221110,222110\}.
\]
Su codificación como enteros de base tres es
\[
 \{28,55,98,175,437,498,687,714\}.
\]
\end{proposition}
\begin{proof}
La multiplicidad de incidencias por posición $j=5,6,7,8,9,10,11,12$
es, respectivamente,
\[
 (4,2,3,2,1,1,2,4).
\]
En cada posición, todas las filas imprimen la misma palabra. Cada fila
se verifica sustituyendo sus tres bandas y los tres coeficientes en la
fórmula de $L$; el décimo caso se ha calculado explícitamente mediante
$L^{\mathrm{af}}_{30}$. Por ejemplo, en $t=67$ y con coeficientes
$(2,1,1)$ se obtiene
\[
 2(010012)+(001220)+(011120)=(002001)\pmod3.
\]
Los ocho resultados son distintos como palabras de longitud seis. Su
codificación $\sum_{k=1}^{6}w_k3^{6-k}$ da, en el orden de las
posiciones, $(55,175,498,437,98,28,714,687)$. Al olvidar el orden se
obtiene el conjunto de enteros enunciado. Así, la extracción preserva
cada palabra completa, incluidas sus posiciones nulas, y entrega ocho
elementos distintos al selector.
\end{proof}

La prueba anterior es una recuperación exacta desde incidencias
identificadas. Su dato de entrada comprende el registro de incidencias
con sus posiciones y testigos. La selección orbital posterior recibe
únicamente $\mathcal S_{\mathrm{term}}$ y el descriptor regional inicial;
recorre las $8!$ ordenaciones y determina su compatibilidad con las
fronteras funcionales y las incidencias orientadas de la rueda
dodecafásica. Ambas operaciones quedan así contiguas, con dominios
explícitos: procedencia temporal de las ocho palabras y determinación
constructiva de su orden.

El cuadro comparativo contiene las diecinueve incidencias que su fórmula
produce sobre la segmentación de doce bloques en los cien tiempos
declarados. Esa exhaustividad finita se comprueba recorriendo
$100\cdot2\cdot3\cdot2=1200$ evaluaciones. El lector afín amplía la
cobertura de las posiciones terminales de siete a ocho. Este recuento
caracteriza el repertorio local de lectores que acaba de definirse;
las demás operaciones del TPK conservan los dominios y funciones
constructivas establecidos en sus respectivas derivaciones.

\subsection{Lectores comparativos, lector afín e incidencia temporal}

\subsubsection{Transformación de Witt del selector}

El selector utiliza la coordenatización de seis coordenadas
\[
 A_{\rm ter}=
 \begin{pmatrix}
 0&1&1&1&1&1\\
 1&0&1&2&2&1\\
 1&1&0&1&2&2\\
 1&2&1&0&1&2\\
 1&2&2&1&0&1\\
 1&1&2&2&1&0
 \end{pmatrix}.
\]
Se define $W(b)=bA_{\rm ter}$ en $\mathbb F_3^6$, así como las cargas
\[
 \eta(b)=\sum_{j=0}^5 b_j\pmod3,
 \qquad\eta_{\rm ter}^\vee(b)=\eta(W(b)).
\]
La coordenatización del descriptor y la del selector se relacionan mediante la
permutación $\sigma=(0,1,2,4,5,3)$, escrita como lista de imágenes:
\begin{equation}
 (A_{\rm ter})_{ij}=(A_{\rm reg})_{\sigma(i),\sigma(j)}.
 \label{act21:eq:kdes-cambio-carta}
\end{equation}

\begin{lemma}[Compatibilidad del funcional de carga]
Para toda banda $b$, se cumple
\[
 \eta_{\rm ter}^\vee(b)=\eta^\vee(b).
\]
\end{lemma}
\begin{proof}
La suma entera de cada fila correspondiente de ambas matrices coincide:
en la primera fila es cinco y en las otras cinco es siete. Por
distributividad,
\[
 \sum_j\sum_i b_i(A_{\rm ter})_{ij}
 =\sum_i b_i\sum_j(A_{\rm ter})_{ij}
 =\sum_i b_i\sum_j(A_{\rm reg})_{ij}.
\]
La reducción módulo tres da la igualdad de cargas. La identidad de
matrices \eqref{act21:eq:kdes-cambio-carta} se verifica por sustitución de sus
treinta y seis entradas y conserva, además, el cambio de coordenadas.
\end{proof}

La igualdad del funcional total de carga permite comparar sus
representaciones. Las transformaciones vectoriales siguen identificadas
por sus matrices y por la permutación explícita; la prueba anterior
conserva esa información de coordenadas.

\subsubsection{Doce lectores comparativos por instante}

Se define la fase temporal
\[
 \theta(t)=(t+2)\bmod3.
\]
Esta escritura realiza $(t-1)\bmod3$ también para $t=0$, donde la fase
vale dos. Para cada tipo de carga
$\nu\in\{\eta,\eta_{\rm ter}^\vee\}$, desplazamiento
$o\in\{0,1,2\}$ y orientación $s\in\{1,2\}$, se toman
\[
 \epsilon_\chi^{\nu,o,s}(t)=s
 \begin{cases}
 1,&\nu(B_\chi(t))=\theta(t)+o\pmod3,\\
 2,&\nu(B_\chi(t))\ne\theta(t)+o\pmod3.
 \end{cases}
\]
La palabra comparativa correspondiente es
\begin{equation}
 C_{\nu,o,s}(t)=
 \sum_{\chi\in\{\pi,e,\varphi\}}
 \epsilon_\chi^{\nu,o,s}(t)B_\chi(t)
 \quad\text{en }\mathbb F_3^6.
 \label{act21:eq:kdes-comparativo}
\end{equation}
Los parámetros generan doce representaciones etiquetadas por fila. Dos
etiquetas pueden producir la misma palabra; los testigos temporales y
las etiquetas permanecen disponibles en el calendario completo.

\subsubsection{Lector afín de fase y carga}

El lector afín emplea coeficientes
\[
 a_\chi^{\rm af}(t)=\eta_{\rm ter}^\vee(B_\chi(t))+\theta(t)
 \pmod3,
\]
y produce
\begin{equation}
 F(t)=\sum_{\chi\in\{\pi,e,\varphi\}}
 a_\chi^{\rm af}(t)B_\chi(t)
 \quad\text{en }\mathbb F_3^6.
 \label{act21:eq:kdes-afin}
\end{equation}
La regla comparativa distingue coincidencia y diferencia de cargas;
la afín suma carga y fase antes de combinar las bandas. Esta diferencia
operativa es la que utiliza la selección heterotípica.

\subsubsection{Incidencia palabra--residuo}

Sea $R$ la rotación cíclica de las seis posiciones de una banda. En la
fila de tiempo $t$ se forma el repertorio residual
\[
 \mathcal R(t)=\{R^jv:
      v\in\{q(t),a(t),c(t),\bar w(t)\},\ 0\leq j<6\}.
\]
Las relaciones de incidencia temporal son
\begin{align}
 \mathcal A_{\rm cmp}
 &=\{(C_{\nu,o,s}(t),u):
      0\leq t<100,\ u\in\mathcal R(t),\ \nu,o,s\text{ admisibles}\},\\
 \mathcal A_{\rm af}
 &=\{(F(t),u):0\leq t<100,\ u\in\mathcal R(t)\}.
 \label{act21:eq:kdes-atlas}
\end{align}
Cada par relaciona una palabra leída con un residuo de la misma fila.
El tiempo se cuantifica de manera común en ambas coordenadas del par.
En el calendario empleado, los tiempos $0,\ldots,99$ son distintos;
identificar una misma fila e identificar un mismo tiempo son, por ello,
operaciones equivalentes.

\begin{proposition}[Compresión existencial del repertorio temporal]
La función
\[
 M(v,u)=\mathbf1_{(v,u)\in\mathcal A_{\rm cmp}}
        +2\mathbf1_{(v,u)\in\mathcal A_{\rm af}}
        \in\{0,1,2,3\}
\]
conserva exactamente la existencia de cada tipo de lectura para el par
$(v,u)$.
\end{proposition}
\begin{proof}
El primer bit vale uno exactamente cuando existe un tiempo, un campo
residual, una rotación y una etiqueta comparativa que producen el par.
El segundo bit expresa la condición correspondiente al lector afín.
La unión por disyunción de bits conserva cada condición de existencia
al recorrer todas las filas. Un par puede admitir ambos tipos, y en ese
caso la máscara vale tres.
\end{proof}

La función $M$ es un lector de existencia. Los tiempos testigos, sus
multiplicidades y las etiquetas individuales permanecen en
\eqref{act21:eq:kdes-atlas} y en el calendario que las produce. La
implementación finita emplea $M$ para decidir heterotipia, mientras que
el archivo genealógico conserva los datos de procedencia más ricos.

\subsection{Enumeración de cilindros y candidatos decimales}

\subsubsection{Concatenación de setenta y dos y setenta y ocho trits}

Para una ordenación $\tau=(s_1,\ldots,s_8)$ de
$\mathcal S_{\rm term}$ se define
\[
 P_{72}(\tau)=[\Omega_{24}\Vert s_1\Vert\cdots\Vert s_8]_3.
\]
En forma recursiva, se parte de $p_0=[\Omega_{24}]_3$ y se calcula
$p_i=729p_{i-1}+[s_i]_3$ para $1\leq i\leq8$. Esta recurrencia
conserva los ceros iniciales de cada bloque por la multiplicación
explícita por $729$. Para cada posición $b$ del panel funcional,
\[
 P_{78}(\tau,b)=729P_{72}(\tau)+b.
\]
Las longitudes son, respectivamente, $24+8\cdot6=72$ y $72+6=78$.
El cilindro ternario correspondiente es
\begin{equation}
 I_{\tau,b}=
 \left[\frac{P_{78}(\tau,b)}{3^{78}},
       \frac{P_{78}(\tau,b)+1}{3^{78}}\right).
 \label{act21:eq:kdes-cilindro}
\end{equation}

\subsubsection{Intersección exacta con la rejilla de doce tríadas}

La lectura decimal utiliza doce tríadas, es decir, la rejilla
$N/10^{36}$. Se abrevia $D=10^{36}$ y $Q=3^{78}$.

\begin{lemma}[Extremos enteros del cilindro]
Para un prefijo entero $p$ de longitud setenta y ocho, los numeradores
de los puntos de la rejilla que pertenecen a su cilindro son
\begin{equation}
 \left\{N\in\mathbb Z:
 \left\lceil\frac{pD}{Q}\right\rceil
 \leq N<
 \left\lceil\frac{(p+1)D}{Q}\right\rceil\right\}.
 \label{act21:eq:kdes-rejilla}
\end{equation}
Cada cilindro contiene a lo sumo un punto de esta rejilla.
\end{lemma}
\begin{proof}
La condición $N/D\in[p/Q,(p+1)/Q)$ equivale a
$pD/Q\leq N<(p+1)D/Q$. Para $N$ entero, la primera desigualdad
equivale a $\lceil pD/Q\rceil\leq N$; la segunda equivale a
$N<\lceil(p+1)D/Q\rceil$. La longitud del intervalo en coordenada
$N$ es $D/Q<1$, pues $10^{36}<3^{78}$. Por ello puede contener,
como máximo, un entero.
\end{proof}

Los techos se evalúan por división entera:
$\lceil a/Q\rceil=\lfloor(a+Q-1)/Q\rfloor$ para $a\geq0$.
Toda la enumeración utiliza enteros exactos; la elección de puntos de
frontera queda fijada por el extremo derecho semiabierto.

Se conserva primero la lista indexada de incidencias
\[
 \mathcal C_{\rm ind}=
 \{(\tau,j,N):N/D\in I_{\tau,\mathcal P_9(j)}\},
 \qquad1\leq j\leq9,
\]
y se obtiene después el conjunto de numeradores
\[
 \mathcal C=\{N:\exists\tau,j\ (\tau,j,N)\in\mathcal C_{\rm ind}\}.
\]
El paso a $\mathcal C$ elimina repeticiones de un numerador, conservadas
en $\mathcal C_{\rm ind}$ con su procedencia.

\begin{lemma}[Recuperación de la ordenación desde el punto del cilindro]
Si $(\tau,j,N)\in\mathcal C_{\rm ind}$, el bloque número $i$ de la
palabra de setenta y ocho trits se recupera mediante
\[
 \left\lfloor\frac{N\,729^i}{D}\right\rfloor\bmod729,
 \qquad1\leq i\leq13.
\]
En particular, los bloques quinto a duodécimo recuperan $\tau$.
\end{lemma}
\begin{proof}
La pertenencia al cilindro significa
$p\leq QN/D<p+1$, donde $p=P_{78}(\tau,\mathcal P_9(j))$.
Por tanto, $\lfloor QN/D\rfloor=p$. La extracción de los bloques
anteriores coincide con la división sucesiva del prefijo $p$ por
potencias de $729$: truncar una expansión posicional antes de extraer
un bloque conserva dicho bloque. La fórmula escrita realiza
exactamente esa extracción. Los cuatro primeros bloques forman
$\Omega_{24}$ y los ocho siguientes son la ordenación $\tau$.
\end{proof}

\begin{proposition}[Censo exacto de candidatos]
Para $\Omega_{24}$, $\mathcal S_{\rm term}$ y $\mathcal P_9$
definidos anteriormente,
\[
 \#\operatorname{Ord}(\mathcal S_{\rm term})=40320,
 \qquad\#\mathcal C_{\rm ind}=21902,
 \qquad\#\mathcal C=19446.
\]
\end{proposition}
\begin{proof}
Los ocho bloques son distintos, por lo que sus ordenaciones son las
$8!=40320$ permutaciones. Para cada permutación y cada una de las
nueve posiciones del panel se calcula $p=P_{78}(\tau,b)$, se aplican
los extremos \eqref{act21:eq:kdes-rejilla} y se añade su único entero cuando
el intervalo entero es no vacío. La suma de los cardinales de esos
intervalos es $21902$. La unión de sus enteros contiene $19446$
elementos distintos. La evaluación exhaustiva de estas operaciones
está certificada por los teoremas
\texttt{indexed\_cylinder\_count} y \texttt{candidate\_count} del
selector finito. El algoritmo realiza exactamente la enumeración
descrita: permutaciones, nueve incidencias del panel, límites por
división entera y eliminación de duplicados.
\end{proof}

Los números $21902$ y $19446$ cuentan objetos diferentes. El primero
cuenta incidencias indexadas que contienen un punto decimal; el segundo
cuenta numeradores distintos. El número total de consultas de cilindro
es $40320\cdot9=362880$, incluidos los cilindros cuya intersección con
la rejilla es vacía.

\subsection{Incidencia excepcional y condición de heterotipia}

\subsubsection{Cara excepcional producida por el panel}

Se eleva una banda a doce coordenadas mediante
\[
 \Gamma(b)=(b,W(b))\in\mathbb F_3^{12}.
\]
El soporte de una palabra es el conjunto de posiciones de sus
coordenadas no nulas, con índices humanos $1,\ldots,12$. Las primeras
bandas regionales del panel producen
\begin{equation}
 \mathcal B=\operatorname{supp}\Gamma(B_\pi(0))
          \cap\operatorname{supp}\Gamma(B_e(0))
          =\{4,6,9,10\}.
 \label{act21:eq:kdes-cara}
\end{equation}
La operación utiliza las palabras regionales $103$ y $523$ y la
matriz de Witt del selector. La cara se calcula antes de consultar las
coordenadas de los candidatos.

Para $N\in\mathcal C$, sus doce tríadas son
\begin{equation}
 K_i(N)=\left\lfloor\frac{N}{1000^{12-i}}\right\rfloor\bmod1000,
 \qquad1\leq i\leq12.
 \label{act21:eq:kdes-triadas}
\end{equation}
El soporte de corona es
\[
 \mathcal H(N)=\{i:K_i(N)\geq729\},
 \qquad
 \mathcal C_{\mathcal B}=\{N\in\mathcal C:\mathcal H(N)=\mathcal B\}.
\]
La frontera $729=3^6$ relaciona la banda de seis trits con la tríada
decimal. La evaluación exhaustiva de la igualdad de soportes da
\begin{equation}
 \#\mathcal C_{\mathcal B}=79.
 \label{act21:eq:kdes-79}
\end{equation}
El soporte conserva posiciones, mientras que el registro completo
conserva también las doce magnitudes. La etapa siguiente utiliza esa
información posicional junto con relaciones entre magnitudes.

\subsubsection{Palabras ternarias y diferencias orientadas del candidato}

La $i$-ésima banda ternaria del punto $N/D$ se obtiene por
\begin{equation}
 T_i(N)=\left\lfloor\frac{N\,729^i}{D}\right\rfloor\bmod729,
 \qquad1\leq i\leq13.
 \label{act21:eq:kdes-tercandidato}
\end{equation}
La arista orientada $a\to b$ tiene residuo
\begin{equation}
 R_{a,b}(N)=(K_a(N)-K_b(N))\bmod729.
 \label{act21:eq:kdes-resarista}
\end{equation}
La resta se realiza en $\mathbb Z$ antes de la reducción. Para esta
arista se consulta en las relaciones \eqref{act21:eq:kdes-atlas} el par
\[
 \bigl(\operatorname{enc}_6^{-1}(T_b(N)),
  \operatorname{enc}_6^{-1}(R_{a,b}(N))\bigr).
\]
Es, por tanto, una consulta
conjunta a la palabra de destino y a la diferencia orientada de las
tríadas.

Se escribe $\mathrm{Cmp}(a,b;N)$ cuando el par pertenece a
$\mathcal A_{\rm cmp}$ y $\mathrm{Af}(a,b;N)$ cuando pertenece a
$\mathcal A_{\rm af}$.

\subsubsection{Determinación de la órbita transversal}

El descriptor de veinticuatro trits determina las tres primeras
tríadas decimales antes de las filtraciones terminales. En efecto, su
cilindro satisface la inclusión exacta
\[
 \left[\frac{66241910521}{3^{24}},
       \frac{66241910522}{3^{24}}\right)
 \subseteq
 \left[\frac{234543140}{10^9},
       \frac{234543141}{10^9}\right).
\]
Las dos desigualdades de extremos se verifican por productos cruzados
enteros. Por tanto, todo candidato comparte las tres posiciones de
referencia $1,2,3$ y sus representaciones $(234,543,140)$.

La traslación de paso cuatro sobre $\mathbb Z/12\mathbb Z$ posee las
cuatro órbitas
\[
 (1,5,9),\qquad(2,6,10),\qquad(3,7,11),\qquad(4,8,12).
\]
Las tres primeras contienen, respectivamente, una de las posiciones de
referencia; la cuarta es la única disjunta de $\{1,2,3\}$. El bosque
orientado de los pasos tres y cuatro se fija mediante las aristas
\[
 \begin{array}{lll}
 1\to4,&1\to5,&2\to6,\\
 3\to7,&4\to8,&5\to9,\\
 6\to10,&7\to11,&8\to12.
 \end{array}
\]
Su primera arista tiene paso tres y las restantes paso cuatro. Las dos
aristas internas de la órbita distinguida son precisamente
$4\to8$ y $8\to12$. Así queda fijado el lugar de la condición de
lectura por la geometría de las posiciones y por el prefijo inicial,
antes de examinar los candidatos de la cara excepcional.

\subsubsection{Condición terminal de heterotipia}

La condición terminal es
\begin{equation}
 \begin{split}
 \mathcal T(N)={}&
 [\mathrm{Cmp}(4,8;N)\wedge\mathrm{Af}(8,12;N)]\\
 &\vee[\mathrm{Af}(4,8;N)\wedge\mathrm{Cmp}(8,12;N)].
 \end{split}
 \label{act21:eq:kdes-heterotipia}
\end{equation}
Se admiten las dos asignaciones de tipos a las dos aristas. La
selección exige que exista una realización con tipos distintos; las
etiquetas individuales y sus tiempos quedan disponibles en los
testigos de incidencia.

\begin{theorem}[Selección terminal del registro de acoplamiento]
Sobre el dominio de candidatos producido por
$\Omega_{24}$, $\mathcal S_{\rm term}$, $\mathcal P_9$ y el calendario
regional de cien filas, existe un único numerador que satisface
simultáneamente la incidencia excepcional y la heterotipia:
\[
 \{N\in\mathcal C_{\mathcal B}:\mathcal T(N)\}
 =\{234543140729659824621058914794146601\}.
\]
Su registro de doce tríadas es
\begin{equation}
 K=(234,543,140,729,659,824,621,058,914,794,146,601).
 \label{act21:eq:kdes-kfinal}
\end{equation}
\end{theorem}
\begin{proof}
El censo precedente produce los $79$ elementos de
$\mathcal C_{\mathcal B}$. Para cada uno se extraen sus doce
coordenadas por \eqref{act21:eq:kdes-triadas}; se calculan las bandas de
destino $T_8,T_{12}$ y los residuos $R_{4,8},R_{8,12}$; se consultan
las dos máscaras del repertorio temporal y se aplica la disyunción
\eqref{act21:eq:kdes-heterotipia}. La lista resultante contiene exactamente
el entero indicado, según la evaluación exhaustiva certificada por
\texttt{terminal\_selection}. La igualdad de lista con un singleton
implica existencia y unicidad en el dominio declarado. La extracción
base mil de ese entero da \eqref{act21:eq:kdes-kfinal}.

El procedimiento de selección recibe los productores y las reglas
enumeradas. El entero final interviene en el enunciado de la igualdad
comprobada, después de ejecutar la selección. Esta posición lógica
permite cotejar el resultado sin utilizarlo como criterio de elección.
\end{proof}

\subsubsection{Testigo aritmético de pertenencia a un cilindro}

Una de las incidencias del numerador seleccionado utiliza la ordenación
\[
 (002001,020111,200110,121012,010122,001001,221110,222110)
\]
y la frontera $438=121020_3$. Para ella,
\[
 P_{72}=5283881584882673136514102926090957,
\]
\[
 P_{78}=3851949675379468716518781033120308091.
\]
La fórmula de extremos enteros produce
\begin{align*}
 \left\lceil\frac{P_{78}10^{36}}{3^{78}}\right\rceil
 &=234543140729659824621058914794146601,\\
 \left\lceil\frac{(P_{78}+1)10^{36}}{3^{78}}\right\rceil
 &=234543140729659824621058914794146602.
\end{align*}
El único entero del intervalo es, por tanto, $N_K$. Las trece bandas
ternarias de su punto decimal son
\[
 \begin{array}{rrrr}
 020022&222111&211201&021101\\
 002001&020111&200110&121012\\
 010122&001001&221110&222110\\
 121020&&&
 \end{array}
\]
Este cálculo exhibe un testigo concreto de pertenencia al dominio
enumerado. La unicidad del registro procede del examen exhaustivo de
todos los candidatos, además de este testigo particular.

\subsection{Composición regional y continuidad de las realizaciones}

\subsubsection{Composición de los generadores regionales}

Sea $\operatorname{Sel}(\Omega,\mathcal S,\mathcal P,\mathcal N)$ el
selector definido por la enumeración y los dos filtros anteriores.
Los productores regionales satisfacen las igualdades
\[
 \Omega_{\rm reg}=\Omega_{24},\qquad
 \mathcal P_{\rm reg}=\mathcal P_9,\qquad
 \mathcal N_{\rm reg}=\mathcal N_{100},
\]
donde $\mathcal N_{\rm reg}$ se calcula por extracción de bandas y
lectura de firmas. Por congruencia de funciones,
\begin{equation}
 \begin{split}
 &\operatorname{Sel}(\Omega_{\rm reg},\mathcal S_{\rm term},
                    \mathcal P_{\rm reg},\mathcal N_{\rm reg})\\
 &\hspace{10mm}=
 \operatorname{Sel}(\Omega_{24},\mathcal S_{\rm term},
                    \mathcal P_9,\mathcal N_{100})=\{N_K\}.
 \end{split}
 \label{act21:eq:kdes-composicion}
\end{equation}
La igualdad reemplaza tres interfaces documentales por su construcción
regional efectiva. El repertorio terminal mantiene la procedencia
especificada en \eqref{act21:eq:kdes-sterm}.

El registro se forma entonces por las doce divisiones euclídeas de
\eqref{act21:eq:kdes-triadas}. Cada coordenada pertenece a
$\{0,\ldots,999\}$ y la longitud es doce por construcción. La
implementación que toma el primer elemento de la lista seleccionada
está justificada por la igualdad con el singleton: el valor por defecto
de una operación informática sobre listas vacías nunca interviene en
este resultado.

\subsubsection{Relación con las incidencias firmadas y la inversión}

La construcción presente proporciona una selección prospectiva en un
dominio finito explícito. El desarrollo ulterior del registro
dodecafásico produce sus canales orientados a partir del estado
terminal enriquecido y demuestra una inversión integral mediante
Hadamard. Aquella inversión recupera un registro desde canales
compatibles; aquí se selecciona un registro a partir de sus
antecedentes regionales e incidenciales. La igualdad de sus resultados
permite componer las dos realizaciones manteniendo el origen de cada
operación.

La unicidad de selección y la unicidad de inversión son propiedades
distintas. La primera cuantifica los candidatos construidos en este
capítulo. La segunda cuantifica los registros que comparten una firma
compatible. Su conjunción justifica que el registro seleccionado pueda
transportarse a la realización firmada y recuperarse desde ella, sin
convertir una recuperación retrospectiva en el productor del registro.

\subsubsection{Relación entre los censos terminales}

La secuencia
\[
 40320\ \longrightarrow\ 21902\ \longrightarrow\ 19446
 \ \longrightarrow\ 79\ \longrightarrow\ 1
\]
registra, sucesivamente, ordenaciones, incidencias indexadas con la
rejilla, numeradores distintos, candidatos de la cara excepcional y
supervivientes heterotípicos. La notación de flechas indica el orden de
las operaciones; la segunda cifra cuenta una clase de objetos
distinta de la primera.

El censo de catálogos regionales
$196\cdot197\cdot196\longrightarrow331\longrightarrow2
\longrightarrow1$, desarrollado en su lugar correspondiente, parte de
otra población: bloques regionales de tres catálogos y condiciones
terminales de Gram, ocupación, distancia y orientación. Ambos censos
conservan sus dominios. Su convergencia en una misma representación se
expresa mediante la identificación de sus lectores y de sus salidas;
las cifras intermedias describen operaciones diferentes y permanecen
documentadas por separado.

\subsubsection{Composición regional de la constante de estructura fina}

La salida $K$ interviene a continuación en la composición regional
ordenada y en el acarreo de base mil. Su empleo posterior conserva las
doce posiciones y la frontera de lectura. El acoplamiento con la raíz
analítica y las representaciones a profundidad arbitraria utiliza esa
misma salida, junto con las coordenadas regionales ya generadas y las
reglas de la coordenatización analítica. Esta dependencia prepara el capítulo
dedicado a las dos realizaciones correlacionadas de la constante de
estructura fina.

El alcance de este capítulo es preciso: producción del descriptor y
de sus antecedentes regionales, determinación del dominio finito,
selección única sobre dicho dominio y composición material con los
productores. La recurrencia analítica de todos los grados y la
compatibilidad de los acarreos infinitos se desarrollan después con sus
demostraciones propias.

\subsection{Certificación ejecutable y trazabilidad de las operaciones}

El desarrollo formal conserva la correspondencia entre las
definiciones matemáticas y sus realizaciones ejecutables. Las pruebas
de firmas y de reloj se encuentran en
\texttt{N69RegionalSignature.lean}; la extracción de bandas y la
reproducción de las cien filas, en
\texttt{GeneratedN69Rows.lean}; el criterio mínimo y la construcción
de $\Omega_{24}$, en \texttt{RegionalW24.lean}. La compatibilidad de
las dos matrices de Witt está reunida en
\texttt{WittReaderCompatibility.lean}.

Los lectores temporales, las rotaciones, los residuos orientados y la
condición heterotípica se formalizan en
\texttt{TerminalReaders.lean}. La enumeración y los censos se
formalizan en \texttt{Terminal}\allowbreak\texttt{Orbital}\allowbreak\texttt{Selection.lean}. El panel
regional se enlaza mediante \texttt{Regional}\allowbreak\texttt{Panel}\allowbreak\texttt{Bridge.lean}, y la
composición \eqref{act21:eq:kdes-composicion} queda expresada en
\texttt{Selected}\allowbreak\texttt{From}\allowbreak%
\texttt{Regional}\allowbreak\texttt{Inputs.lean}.

Las comprobaciones finitas de enumeración emplean la evaluación nativa
de Lean; sus informes registran \texttt{Lean.ofReduceBool}, además
de los axiomas lógicos habituales de la biblioteca. La reconstrucción
regional del descriptor utiliza pruebas del núcleo y reducción
ordinaria. Esta distinción describe los mecanismos de comprobación
realmente empleados. Los recibos de compilación acreditan sus módulos
y dependencias declarados; las pruebas matemáticas conservan los
dominios, productores e hipótesis que se han explicitado en el cuerpo
del capítulo.

\subsection{Censo completo de cilindros y selección de la frontera terminal}
\label{act21:censo:terminal-completo}

La lectura terminal utiliza conjuntamente dos informaciones: la compatibilidad
de los cilindros regionales y un perfil de incidencia con canales y posiciones
etiquetados. En este apartado se construyen los catálogos completos, se
enumeran sus ternas y se demuestra el efecto de cada condición del lector.
El punto de partida son las salidas regionales ya producidas en la
sección~\ref{sec:generacion}, no valores introducidos para definir el
transporte APP--TRIT--TPK. La reducción a una matriz pertenece a una lectura
posterior de esas salidas y conserva el orden de los tres canales.

El resultado es un teorema censal con estructura de partida explícita.
El perfil de Gram, pesos, distancias y márgenes se declara como dato del
lector terminal. La exhaustividad de la enumeración demuestra qué selecciona
ese perfil sobre los catálogos fijados; no demuestra, por sí sola, que el
perfil sea una ley universal ni que pueda reconstruirse desde la única terna
que finalmente lo satisface. Esta separación permite incorporar la prueba
finita completa sin invertir su genealogía.

\subsubsection{De las salidas regionales a los catálogos ternarios}

Sean \(x_\pi,x_e,x_\varphi\in[0,1)\) las partes fraccionarias de las
tres coordenadas regionales ya construidas. Se utilizan sus cilindros
expresados de mil posiciones decimales,
\[
 J_\chi^{(1000)}=
 \left[\frac{d_\chi}{10^{1000}},
       \frac{d_\chi+1}{10^{1000}}\right),
 \qquad \chi\in\{\pi,e,\varphi\}.
\]
Los enteros \(d_\chi\) son lecturas finitas de salida. La longitud mil
es suficiente para los enteros que siguen, como se comprueba mediante las
dos cotas de cada intervalo; no es una profundidad supuesta infinita.
Definimos
\begin{equation}
 p=30,\quad r=1050,\quad L=p+r=1080,\quad k=171,
 \qquad D=3^L,\quad Q=1000^k,
 \label{act21:censo:parametros}
\end{equation}
y
\[
 P_\chi=\lfloor3^{30}x_\chi\rfloor,\qquad
 T_\chi=\lfloor Qx_\chi\rfloor.
\]
Aquí \(k\) cuenta tríadas decimales y no designa el registro dodecafásico
\(K\). Se tiene \(1000^{171}\leq3^{1080}<1000^{172}\).

\begin{lemma}[Extracción estable de un entero]
\label{act21:censo:estabilidad}
Para \(d,m,h\in\mathbb Z\), con \(m,h>0\), el valor de
\(\lfloor mx\rfloor\) es constante en \([d/h,(d+1)/h)\) si y sólo si
\begin{equation}
 \left\lfloor\frac{md}{h}\right\rfloor
 =\left\lfloor\frac{m(d+1)-1}{h}\right\rfloor.
 \label{act21:censo:prueba-estabilidad}
\end{equation}
Cuando se cumple, cualquiera de los dos miembros da ese valor.
\end{lemma}
\begin{proof}
El mínimo de \(mx\) se alcanza en el extremo izquierdo. Su supremo
es \(m(d+1)/h\) y no se alcanza. El mayor entero que puede tomar
\(\lfloor mx\rfloor\) es
\(\lceil m(d+1)/h\rceil-1\), igual al segundo miembro porque
el numerador y el denominador son enteros. La igualdad de los extremos
del intervalo de valores prueba la afirmación.
\end{proof}

La prueba de estabilidad se cumple en los tres canales para
\(m=3^{30},Q,3^{1080}\). La última elección se utilizará únicamente
para una comparación posterior, no para seleccionar el perfil. Las
primeras treinta posiciones recuperadas son precisamente las cinco
etapas de seis trits del transporte regional:
\begin{center}\small
\begin{tabular}{@{}c l r@{}}\toprule
Canal&Prefijo de treinta trits&\(P_\chi\)\\\midrule
\(\pi\)&\texttt{010211012222010211002111110221}&29152671743887\\
\(e\)&\texttt{201101121221102011012222102011}&147887858824447\\
\(\varphi\)&\texttt{121200112202121020010210010200}&127247717616687\\\bottomrule
\end{tabular}
\end{center}
La coincidencia conserva la prolongación ya construida; la lectura decimal
no se usa para sustituir retroactivamente sus matrices de transporte.

Fijado un canal, una cola de \(r\) trits se representa por
\(0\leq R<3^r\). Su cilindro y el cilindro de las primeras \(k\)
tríadas son
\begin{equation}
 I_{\chi,R}=\left[\frac{P_\chi3^r+R}{D},
                       \frac{P_\chi3^r+R+1}{D}\right),\qquad
 J_{\chi,T}=\left[\frac{T_\chi}{Q},\frac{T_\chi+1}{Q}\right).
 \label{act21:censo:dos-cilindros}
\end{equation}
El catálogo incluye todas las colas para las que estos dos cilindros
tienen intersección no vacía.

\begin{proposition}[Catálogo por intersección exacta]
\label{act21:censo:catalogos}
Las colas compatibles son los enteros del intervalo \([a_\chi,b_\chi]\),
con
\begin{align}
 a_\chi&=\max\left(0,
       \left\lfloor\frac{T_\chi D}{Q}\right\rfloor-P_\chi3^r\right),
       \label{act21:censo:extremo-a}\\
 b_\chi&=\min\left(3^r-1,
       \left\lfloor\frac{(T_\chi+1)D-1}{Q}\right\rfloor-P_\chi3^r\right).
       \label{act21:censo:extremo-b}
\end{align}
Para los cilindros regionales fijados, los recortes no actúan y se obtiene
\begin{equation}
\begin{array}{c|r|r|r}
 \chi&a_\chi\bmod729&b_\chi\bmod729&b_\chi-a_\chi+1\\\hline
 \pi&67&262&196\\
 e&584&51&197\\
 \varphi&117&312&196
\end{array}
\label{act21:censo:tabla-catalogos}
\end{equation}
\end{proposition}
\begin{proof}
Escribamos \(N=P_\chi3^r+R\). Dos intervalos semiabiertos de
\eqref{act21:censo:dos-cilindros} se cortan exactamente cuando
\[
 NQ<(T_\chi+1)D,\qquad (N+1)Q>T_\chi D.
\]
La segunda desigualdad equivale a
\(N\geq\lfloor T_\chi D/Q\rfloor\); la primera, a
\(N\leq\lfloor((T_\chi+1)D-1)/Q\rfloor\). Restar
\(P_\chi3^r\) y conservar \(0\leq R<3^r\) prueba las fórmulas.
La evaluación por divisiones enteras de los \(T_\chi\) extraídos por
\eqref{act21:censo:prueba-estabilidad} da la tabla. Como \(r\geq6\),
\(P_\chi3^r\equiv0\pmod{729}\), de modo que el residuo del primer
extremo también se obtiene directamente de
\(\lfloor T_\chi D/Q\rfloor\bmod729\).
\end{proof}

Los últimos seis trits de una cola forman la palabra
\(\nu_6(R\bmod729)\), donde \(\nu_6\) es la escritura ternaria
de longitud seis, incluidos sus ceros iniciales. Por la tabla, cada
catálogo tiene menos de \(729\) elementos consecutivos, de manera que
esta reducción es inyectiva sobre él. Así, los tres catálogos de frontera
quedan descritos completamente, sin elegir las ternas finales:
\begin{equation}
 \begin{split}
 \mathcal B_\pi&=\{\nu_6(67+i):0\leq i<196\},\\
 \mathcal B_e&=\{\nu_6((584+j)\bmod729):0\leq j<197\},\\
 \mathcal B_\varphi&=\{\nu_6(117+\ell):0\leq\ell<196\}.
 \end{split}
 \label{act21:censo:fronteras-explicitas}
\end{equation}
Exigir que el extremo izquierdo de \(I_{\chi,R}\) pertenezca a
\(J_{\chi,T}\) sería otra condición: perdería el primer cilindro
que corta la frontera y daría \(195,196,195\). El censo utiliza la
intersección de los cilindros completos, de acuerdo con su definición.

\subsubsection{Perfil declarado e indicadores de compatibilidad}

Para \(u,v\in\mathbb F_3^6\), sean
\(g(u,v)=\sum u_iv_i\in\mathbb F_3\),
\(n(u)=g(u,u)\), \(w(u)=\#\{i:u_i\ne0\}\) y
\(h(u,v)=\#\{i:u_i\ne v_i\}\). Las dos últimas funciones toman
valores enteros; las dos primeras, valores en \(\mathbb F_3\).
El perfil local del lector, en orden \((\pi,e,\varphi)\), es
\begin{equation}
 \begin{split}
 (g_{\pi e},g_{\pi\varphi},g_{e\varphi})&=(2,2,0),\qquad
 (n_\pi,n_e,n_\varphi)=(0,0,0),\\
 (w_\pi,w_e,w_\varphi)&=(3,3,3),\qquad
 (h_{\pi e},h_{\pi\varphi},h_{e\varphi})=(2,5,3).
 \end{split}
 \label{act21:censo:perfil-local}
\end{equation}
Las palabras se ordenan como en \eqref{act21:censo:fronteras-explicitas}.
Usaremos \([E]\in\{0,1\}\) para el indicador de una condición \(E\).
Para cada par de canales \(\chi,\psi\), la matriz rectangular
\(A_{\chi\psi}\) tiene entrada \([g(u,v)=g_{\chi\psi}]\).
La matriz \(A^{h}_{\chi\psi}\) multiplica esa entrada por
\([h(u,v)=h_{\chi\psi}]\). Sean \(D_\chi^n\) la matriz diagonal
de \([n(u)=0]\), y \(D_\chi^{nw}\) la de
\([n(u)=0][w(u)=3]\).

\begin{theorem}[Enumeración exhaustiva por sumas finitas]
\label{act21:censo:teorema-conteo}
El producto de los tres catálogos contiene \(7\,567\,952\) ternas.
Las condiciones sucesivas de \eqref{act21:censo:perfil-local} dejan
\begin{equation}
 7\,567\,952\longrightarrow275\,794\longrightarrow6235
 \longrightarrow3127\longrightarrow331.
 \label{act21:censo:cadena-local}
\end{equation}
Estos cuatro conteos son, respectivamente,
\begin{align}
 N_g&=\operatorname{tr}(A_{\pi e}A_{e\varphi}A_{\varphi\pi}),
 \nonumber\\
 N_{gn}&=\operatorname{tr}(D_\pi^n A_{\pi e}
             D_e^n A_{e\varphi}D_\varphi^n A_{\varphi\pi}),
 \nonumber\\
 N_{gnw}&=\operatorname{tr}(D_\pi^{nw} A_{\pi e}
             D_e^{nw} A_{e\varphi}D_\varphi^{nw} A_{\varphi\pi}),
 \nonumber\\
 N_{gnwh}&=\operatorname{tr}(D_\pi^{nw} A^h_{\pi e}
             D_e^{nw} A^h_{e\varphi}D_\varphi^{nw} A^h_{\varphi\pi}).
 \label{act21:censo:trazas}
\end{align}
Todos los productos y trazas de esta fórmula se calculan en los enteros,
no en el cuerpo de tres elementos.
\end{theorem}
\begin{proof}
El tamaño del producto es \(196\cdot197\cdot196\). Al expandir
una traza de \eqref{act21:censo:trazas}, sus índices recorren exactamente una
terna de palabras. El sumando vale uno si satisface todas las condiciones
representadas y cero en otro caso. Cada terna se cuenta una vez.
Las matrices diagonales añaden las condiciones individuales; los factores
rectangulares añaden las condiciones por pares. Esto demuestra que las
trazas son precisamente los cardinales anunciados.

Para explicitar su evaluación, fijados los índices \(i,j\) de los dos
primeros canales, se forman los conjuntos de índices \(\ell\) del tercero
que cumplen cada condición por pares. Se intersectan sus indicadores,
se añaden los indicadores individuales pertinentes y se suma su cardinal.
La suma final recorre \(0\leq i<196\), \(0\leq j<197\).
La sustitución de las palabras de \eqref{act21:censo:fronteras-explicitas} y
del perfil \eqref{act21:censo:perfil-local} da, en ese orden,
\(275794,6235,3127,331\). El suplemento registra las cuatro contribuciones
de cada uno de los \(196\) índices \(i\) y las \(331\) ternas completas;
su programa evalúa estas sumas con indicadores binarios exactos y contrasta
escalarmente cada superviviente. No interviene una muestra ni una tolerancia.
\end{proof}

\subsubsection{Incidencia etiquetada y carga de orientación}

Para una terna superviviente, sea \(B\in\mathbb F_3^{3\times6}\) su
matriz de filas ordenadas. El segundo perfil exige
\begin{equation}
 \operatorname{rank}_{\mathbb F_3}B=3,\qquad
 w_{\mathrm{col}}(B)=(1,1,2,2,3,0),\qquad
 q_{\partial}(B)=(1,2,2,1,2,0),
 \label{act21:censo:perfil-columnas}
\end{equation}
donde \(w_{\mathrm{col}}\) cuenta los elementos no nulos de cada
columna y \(q_{\partial}\) suma sus entradas módulo tres. Es un perfil
de columnas etiquetadas; permutarlas sin transportar el perfil define otro
lector. Denotemos por \(\mathcal F\) el conjunto de las \(331\) ternas
del teorema anterior.

\begin{proposition}[Reducción del censo al par terminal]
\label{act21:censo:dos-matrices}
Las sumas de indicadores sobre \(\mathcal F\), al exigir sucesivamente
rango, pesos de columna y sumas de columna, son
\begin{equation}
 \sum_{B\in\mathcal F}[\operatorname{rank}B=3]=331,
 \qquad N_{\mathrm{rango,pesos}}=18,
 \qquad N_{\mathrm{rango,pesos,sumas}}=2.
 \label{act21:censo:conteo-columnas}
\end{equation}
Con índices de catálogo comenzando en uno, las dos ternas finales son
\((120,197,157)\) y \((129,197,148)\), y sus matrices son
\[
 B_0=\begin{pmatrix}0&2&0&2&2&0\\0&0&1&2&2&0\\1&0&1&0&1&0\end{pmatrix},
 \qquad
 B_1=\begin{pmatrix}0&2&1&0&2&0\\0&0&1&2&2&0\\1&0&0&2&1&0\end{pmatrix}.
\]
\end{proposition}
\begin{proof}
En cada terna del conjunto ya enumerado se calcula el rango por eliminación
en \(\mathbb F_3\), se cuentan los elementos no nulos de las seis
columnas y se suman éstas módulo tres. Los productos de los indicadores
de \eqref{act21:censo:perfil-columnas} dan \eqref{act21:censo:conteo-columnas}.
La sustitución de los dos índices finales en
\eqref{act21:censo:fronteras-explicitas} da las matrices expuestas. La lista
completa de supervivientes y los filtros exactos del suplemento permiten
reproducir la eliminación sin presuponer esas dos matrices como dominio.
\end{proof}

La orientación \(\sigma=(1,1,-1)=(1,1,2)\in\mathbb F_3^3\)
selecciona la recta
\(\langle\sigma\rangle=\{(0,0,0),(1,1,2),(2,2,1)\}\).
Las cargas de filas de \(B_0,B_1\) son \((0,2,0)\) y \((2,2,1)\);
sólo la segunda pertenece a esa recta. Se recupera así el último paso
del selector de la sección~\ref{act21:censo:dos-matrices}, ahora precedido por
la construcción y enumeración de su dominio completo. La lectura estable
de \(\lfloor3^{1080}x_\chi\rfloor\), efectuada después de los filtros,
da las fronteras \texttt{021020}, \texttt{001220}, \texttt{100210},
con los mismos rangos \((129,197,148)\). Este cotejo posterior no define
el perfil ni participa en las sumas de indicadores.

\subsubsection{Información conservada y alcance del resultado}

El procedimiento conserva el prefijo de treinta trits, la cola compatible,
el orden regional, la posición de cada trit y la orientación. La reducción
módulo \(729\) conserva aquí la identidad de cada candidato porque el
intervalo de colas tiene longitud menor que \(729\); esa inyectividad
no se afirma para una historia arbitraria. El censo de \(331\) demuestra
también que el perfil local aislado no selecciona una sola matriz: los
márgenes etiquetados y la carga aportan condiciones efectivas adicionales.

Se ha recuperado, por tanto, una enumeración finita íntegra y su selector
sobre cilindros de salida, con perfil de lectura explícito. La construcción
del registro firmado utiliza además su lector de memoria y sus canales,
como se especifica a continuación; la matriz visible no se identifica con
ese registro por igualdad de número de componentes. Tampoco este conteo
demuestra la universalidad del perfil, genera de nuevo las coordenadas
regionales ni decide una afirmación sobre alfa o sobre el continuo completo.


% END NUCLEO_ADITIVO_20260921_SELECCION
\section{Registro dodecafásico y constante racional de fase}
\label{sec:k-registro}

La generación de las coordenadas regionales no agota la información del
estado APP--TRIT--TPK. Una lectura arquimediana determina un valor mediante
cilindros compatibles; otra lectura conserva la oposición de las hojas, la
orientación de la ruta y el registro de las incidencias atravesadas. El
registro dodecafásico pertenece a esta segunda lectura antes de intervenir en
la clausura de alfa. Por ello no se define restando una constante física a
una combinación de valores previamente conocidos. Se reconstruye desde una
coordenada firmada del mismo estado aritmético con memoria de trayectoria que produce las coordenadas
regionales.

La distinción es necesaria incluso cuando dos objetos tienen doce
coordenadas. Una palabra de doce bloques del catálogo, una lista de doce
eventos orientados, un vector entero firmado y un registro dodecafásico en base mil son
objetos distintos. Su comparación requiere las aplicaciones que los
relacionan. La igualdad de longitud no proporciona esas aplicaciones, y la
pérdida de signos no es una mera simplificación de notación.

Este capítulo mantiene el corte causal fijado en los capítulos precedentes:
las nueve posiciones APP de los ejes horizontal y vertical producen las
hojas aritméticas; TRIT tipa régimen y orientación; TPK selecciona,
transporta, actualiza y prolonga, conservando residuo, cociente, acarreo,
hoja, ruta, frontera y memoria. Desde la estructura discreta conjunta del
continuo se obtienen las coordenadas que se utilizan aquí. La cadena
\(w_6\to w_{12}\to w_{18}\to w_{24}\to w_{30}\to R_{36}\to G_9\)
no se sustituye por un calendario finito. Tampoco se invierte el orden
\(U_t\to\Gamma_9\to K_{\mathrm{ph}}\): la última aplicación es una
lectura reducida de memoria, no el generador del registro firmado.

\input{sections/revision_k.tex}

\subsection{Extensión cíclica y graduación de memoria}
\label{subsec:k-calendario}

El calendario observable contiene ciento ocho posiciones, organizadas en
doce ventanas nonádicas. Con índices positivos, sean
\[
 E_m=\{9(m-1)+1,\ldots,9m\},\qquad 1\leq m\leq12.
\]
El soporte ordenado \(E_1\sqcup\cdots\sqcup E_{12}\) conserva posición y
pertenencia a ventana. El grupo \(C_{108}=\mathbb Z/108\mathbb Z\)
conserva, en cambio, la ley de traslación cíclica. Para el representante
\(0\leq\widetilde\theta<108\), las coordenadas
\[
 \beta(\theta)=\left[\left\lfloor\frac{\widetilde\theta}{9}\right\rfloor\right]_{12},
 \qquad r(\theta)=1+(\widetilde\theta\bmod9)
\]
separan el sector dodecafásico y la posición nonádica interior. Ninguna de
ellas conserva por sí sola el número total de vueltas.

\begin{proposition}[El acarreo del calendario]
\label{prop:k-extension}
Las aplicaciones
\[
 \iota:C_{12}\longrightarrow C_{108},\quad[b]\longmapsto[9b],
 \qquad p:C_{108}\longrightarrow C_9,\quad[\theta]\longmapsto[\theta]_9
\]
forman una sucesión exacta no escindida
\[
 0\longrightarrow C_{12}\xrightarrow{\iota}C_{108}
 \xrightarrow{p}C_9\longrightarrow0.
\]
Para la sección conjuntista que elige los representantes
\(0,\ldots,8\), el defecto de aditividad es el acarreo
\[
 c(r,r')=\left[\left\lfloor\frac{\widetilde r+\widetilde r'}9\right\rfloor\right]_{12}.
\]
\end{proposition}

\begin{proof}
El núcleo de \(p\) está formado por los múltiplos de nueve y coincide con
la imagen de \(\iota\). La reducción es sobreyectiva y la multiplicación
por nueve es inyectiva en el dominio indicado. Si la sucesión se
escindiera, el grupo central sería isomorfo a \(C_{12}\times C_9\), cuyo
exponente es \(\operatorname{mcm}(12,9)=36\); el exponente de
\(C_{108}\) es 108. La contradicción prueba la no escisión. Finalmente,
la división euclídea de \(\widetilde r+\widetilde r'\) por nueve da el
acarreo escrito y verifica la identidad de la sección.
\end{proof}

El resultado explica una diferencia que reaparecerá en alfa: regresar a la
fase no equivale a regresar al estado. Si \(q_{\mathrm{mem}}\) cuenta las
vueltas nonádicas, el calendario conserva solamente
\(\beta=[q_{\mathrm{mem}}]_{12}\). Después de ciento ocho pasos,
\(\beta\) retorna, pero \(q_{\mathrm{mem}}\) aumenta en doce. El
registro conserva además el cilindro, el libro de incidencias y la memoria
vertical. El cierre de una coordenatización finita no convierte esa historia en un
ciclo sin memoria.

\subsection{El registro orientado y la coordenatización integral \texorpdfstring{\(1+5+6\)}{1+5+6}}
\label{subsec:k-registro-integral}

Para una historia producida por TPK, sea \(\tau_m\) la subruta sobre
\(E_m\), \(\sigma_m\) su orientación, \(\kappa_m\) el acarreo de
frontera y \(\Lambda_m\) el libro local de incidencias. El registro es
\[
 \mathcal R_{12}(x)=
 \bigl((E_m,\tau_m,\sigma_m,\kappa_m,\Lambda_m)\bigr)_{m=1}^{12}.
\]
Su dominio es el espacio de historias aritméticas con memoria de trayectoria, no el conjunto de
etiquetas de fase. La proyección al calendario olvida precisamente algunos
de los datos que el lector firmado necesita.

Una coordenatización inicial del registro puede escribirse mediante eventos
orientados. La construcción de referencia~(subsección~\ref{subsec:k-registro-integral}) fija el conjunto de códigos de evento
\(\{2,4,6,8\}\) y la orientación sectorial
\[
 \Sigma_{12}=(+,+,+,-,-,-,+,+,+,-,-,-).
\]
Si \(d_t\) es el código de dirección en el paso \(t\), se obtiene
\[
 e_i=\Sigma_{12}(i+1)
 \#\{t\in E_{i+1}:d_t\in\{2,4,6,8\}\},
 \qquad0\leq i<12.
\]
La fórmula mantiene separados el censo no negativo de eventos y la
orientación. No identifica una dirección cardinal con un signo TRIT. El
vector \(e\), cuyos componentes son censos locales firmados, tampoco es
todavía el vector \(U\) que aparecerá más abajo.

Las posiciones opuestas determinan, para \(0\leq i<6\),
\[
 p_i=e_i+e_{i+6},\qquad a_i=e_i-e_{i+6}.
\]
Con
\[
 s_C=\sum_{i=0}^{5}p_i,\qquad M_i=p_i-p_5\quad(0\leq i<5),
\]
se define la coordenatización integral
\[
 \mathcal L_{12}(e)
 =(s_C,M_0,\ldots,M_4,a_0,\ldots,a_5).
\]
La descomposición \(1+5+6\) no es un reparto de coordenadas arbitrario:
una coordenada conserva la carga, cinco comparan los pares respecto de un
par de referencia y seis conservan la oposición de las semicoronas.

\begin{lemma}[Inversión de la coordenatización integral]
\label{lem:k-carta-integral}
Una tupla entera \((s_C,M_0,\ldots,M_4,a_0,\ldots,a_5)\) pertenece a la
imagen de \(\mathcal L_{12}\) si y sólo si
\[
 s_C-\sum_{i=0}^{4}M_i\equiv0\pmod6,
 \qquad p_i\equiv a_i\pmod2\quad(0\leq i<6),
\]
donde
\[
 p_5=\frac{s_C-\sum_{i=0}^{4}M_i}{6},\qquad p_i=p_5+M_i\quad(i<5).
\]
En ese caso su preimagen es única y viene dada por
\[
 e_i=\frac{p_i+a_i}{2},\qquad e_{i+6}=\frac{p_i-a_i}{2}.
\]
\end{lemma}

\begin{proof}
La suma de los cinco márgenes es \(s_C-6p_5\), de donde se obtiene la
división por seis. Reconstruidos los \(p_i\), cada par de ecuaciones
\(p_i=e_i+e_{i+6}\), \(a_i=e_i-e_{i+6}\) tiene la solución indicada.
Esta solución es entera exactamente bajo la congruencia de paridad. Todas
las operaciones se invierten, por lo que no queda libertad residual.
\end{proof}

Esta reversibilidad muestra qué información se conserva; no autoriza a
suprimir el resto del libro. El registro completo tiene un dominio mayor
que esta coordenatización de censos. La extracción de los canales que alimentan a
\(U\) utiliza ese registro completo y no se deduce de la lista
\((e_i)\) por una identificación de nombres.

% Propuesta separada; no modifica el bloque común del artículo I.
% Procedencia: artículo II, revisión 05, registro_incidencias.tex y
% registro_imagen_integral.tex; aclaraciones de índices de revisiones 06/07.
% Inserción prevista inmediatamente antes de subsec:k-terminal.
\subsection{Evaluación de incidencias y determinación de las coordenadas transversales}
\label{iii:r25:incidencias}

El registro dodecafásico conserva una historia de visitas y trazas
orientadas. Sus coordenadas transversales se obtienen mediante una
composición de operaciones: evaluación aritmética de las visitas,
recuento por ventanas, normalización decimal y diferencias cíclicas.
La reconstrucción del registro a partir de esas diferencias constituye
otra operación, cuya existencia y unicidad se demostrarán después.
Esta separación permite precisar tanto los datos utilizados por el
extractor como la información que transmite a la lectura armónica.

Sea $\mathscr L$ un registro de incidencias producido por la selección,
el transporte y la actualización de una historia APP--TRIT--TPK.
Cada visita conserva su ruta, posición APP, hoja, instante,
multiplicidad y clase de frontera; cada traza conserva sus miembros,
orientación, multiplicidad y longitud en la unidad declarada por su
lector. Las ventanas son
\[
 E_m=\{9(m-1)+1,\ldots,9m\},\qquad 1\leq m\leq12.
\]
La evaluación de la hoja de suma o producto determina el entero
$n(v)$ de cada visita y su división positiva
\[
 n(v)=r(v)+9q(v),\qquad 1\leq r(v)\leq9.
\]
El residuo se calcula sobre esa evaluación. La clase de frontera de
una visita de residuo $9$ se conserva mediante
$\epsilon_\partial(v)=1$ en la corona y $-1$ en el interior.

La presentación incidencial del registro utiliza los tres recuentos
\begin{align}
 A_m&=\sum_{v\in E_m}w(v)
       \mathbf1_{\{1,3,5,7\}}(r(v)),\label{iii:r25:recuento-a}\\
 C_m&=\sum_{j\geq0}
       \bigl(\nu^-_{120,m}(j)-\nu^+_{120,m}(j)\bigr),
       \label{iii:r25:recuento-c}\\
 V_m&=\sum_{v\in E_m}w(v)
       \mathbf1_{\{9\}}(r(v))\epsilon_\partial(v).
       \label{iii:r25:recuento-v}
\end{align}
Aquí $w(v)$ es la multiplicidad incidencial y
$\nu^\pm_{120,m}(j)$ cuenta las trazas orientadas del canal indicado
de longitud reducida $120(1+3j)$, según la enumeración del registro.
La colección de trazas y su multiplicidad son datos de la evaluación,
con la regla de producción que les corresponde. La suma entera requiere
soporte finito o una construcción que determine su significado.
La longitud reducida de una traza y el número de instantes de la
ventana son magnitudes tipadas diferentes; la unidad del lector
especifica su relación. Los símbolos $A_m,C_m,V_m$ designan aquí
recuentos enteros, distintos del par angular y del operador constitutivo.

Sea $[x]_{10}\in\{0,\ldots,9\}$ el representante euclídeo.
La normalización conserva también el cociente en
\[
 x=[x]_{10}+10\lfloor x/10\rfloor.
\]
Se forma el bloque
\begin{equation}
 z_m=100[A_m]_{10}+10[C_m]_{10}+[V_m]_{10},
 \qquad 0\leq z_m\leq999.
 \label{iii:r25:bloque}
\end{equation}
Con índices cíclicos módulo $12$, sean
\[
 (D_a z)_m=z_m-z_{m+a},\qquad a\in\{3,4\},
 \qquad Q(z)=\sum_{m=1}^{12}z_m.
\]
La composición incidencial queda expresada por
\begin{equation}
 \mathcal E_{\rm cnt}(\mathscr L)
 =\mathcal D(z):=(D_3z,D_4z,Q(z))\in\mathbb Z^{25}.
 \label{iii:r25:composicion}
\end{equation}
Sus argumentos son las visitas y trazas con sus reglas de recuento.
La evaluación no recibe $K$, $U$, $\alpha$ ni las coordenadas
transversales esperadas. Para identificarla con
$\mathcal E_{108}^{90,120}$ se conserva además la composición con
la colección y las incidencias efectivas del estado terminal. La
invertibilidad de $\mathcal D$ determina una reconstrucción posterior;
por sí sola no selecciona esa colección de incidencias.

\begin{proposition}[Información residual determinada por el extractor]
\label{iii:r25:residuos}
Sean $N,N'\in\mathbb Z^{12\times3}$ dos registros de recuentos,
ordenados por $(A,C,V)$. Para las operaciones
\eqref{iii:r25:bloque}--\eqref{iii:r25:composicion},
\[
 \mathcal E_{\rm cnt}(N)=\mathcal E_{\rm cnt}(N')
 \quad\Longleftrightarrow\quad N-N'\in10\mathbb Z^{36}.
\]
Los $36$ residuos sectoriales determinan exactamente la salida de
$25$ coordenadas. Los cocientes y las procedencias constituyen
información adicional del registro.
\end{proposition}
\begin{proof}
La igualdad módulo $10$ produce los mismos bloques y sus mismas
coordenadas transversales. Recíprocamente, la igualdad de las
coordenadas da $D_3(z-z')=D_4(z-z')=0$. Los pasos $3$ y $4$
generan el ciclo de longitud $12$, por lo que $z-z'$ es constante.
La igualdad de $Q$ anula esa constante. La expansión de cada entero
de $0$ a $999$ en tres cifras decimales es única; se recuperan así
los tres residuos de cada ventana.
\end{proof}

La proposición es válida para todo el dominio entero indicado. No
postula que las $36$ clases sean independientes sobre el subconjunto
de registros admisibles producido por TPK.

\subsection{Imagen integral e inversión de las coordenadas transversales}
\label{iii:r25:imagen-seccion}

Para expresar la inversión se reindexan los bloques por
$i\in\mathbb Z/12\mathbb Z$, de $0$ a $11$. Sean
$(Sz)_i=z_{i+1}$, $D_3=I-S^3$ y $D_4=I-S^4$.
Dada una terna entera $(b,c,q)$, defínanse
\begin{equation}
 w_i=c_i-b_{i+1},\qquad
 P_0=0,\quad P_i=\sum_{j=0}^{i-1}w_j\ (1\leq i\leq12),
 \qquad T(w)=\sum_{i=0}^{11}P_i.
 \label{iii:r25:incrementos}
\end{equation}

\begin{theorem}[Caracterización de la imagen integral]
\label{iii:r25:imagen}
Una terna $(b,c,q)\in\mathbb Z^{25}$ pertenece a
$\operatorname{im}\mathcal D$ si y sólo si
\begin{align}
 b_i&=w_i+w_{i+1}+w_{i+2},\qquad 0\leq i<12,
       \label{iii:r25:relaciones}\\
 \sum_{i=0}^{11}w_i&=0,\label{iii:r25:cierre}\\
 q+T(w)&\equiv0\pmod {12}.\label{iii:r25:congruencia}
\end{align}
Su única preimagen entera es
\begin{equation}
 z_i=\frac{q+T(w)}{12}-P_i,\qquad0\leq i<12.
 \label{iii:r25:inversa}
\end{equation}
Las doce relaciones \eqref{iii:r25:relaciones} y la relación
\eqref{iii:r25:cierre} son linealmente independientes sobre
$\mathbb Q$.
\end{theorem}
\begin{proof}
Para una salida de $\mathcal D$,
\[
 c_i-b_{i+1}
 =(z_i-z_{i+4})-(z_{i+1}-z_{i+4})=z_i-z_{i+1}.
\]
Tres incrementos consecutivos suman $b_i$ y la suma cíclica de
todos ellos es cero. Además, $P_i=z_0-z_i$, de donde
$q=12z_0-T(w)$. Se obtienen las condiciones y la fórmula inversa.
Recíprocamente, la congruencia hace enteras las coordenadas de
\eqref{iii:r25:inversa} y el cierre permite prolongarlas
cíclicamente. Sus diferencias adyacentes son $w_i$, luego $D_3z=b$
y
\[
 (D_4z)_i=w_i+w_{i+1}+w_{i+2}+w_{i+3}
 =w_i+b_{i+1}=c_i.
\]
La suma da $Q(z)=q$. La fórmula determina también la unicidad.
Finalmente, el cambio entero invertible
$(b,c)\mapsto(b,w=c-Sb)$ transforma las doce primeras relaciones
en $b=(I+S+S^2)w$: cada una despeja una coordenada diferente de
$b$. El cierre $\sum_iw_i=0$ añade una relación independiente.
\end{proof}

La imagen racional tiene dimensión $12$. La condición integral añade
la congruencia de orden $12$ de \eqref{iii:r25:congruencia}.
Los once incrementos $w_0,\ldots,w_{10}$ y la carga $q$ son
coordenadas suficientes, con $w_{11}=-\sum_{i=0}^{10}w_i$.
Las $25$ coordenadas conservan esos $12$ grados de libertad
mediante $13$ relaciones lineales. Si la terna proviene además de
\eqref{iii:r25:bloque}, la inversa pertenece a
$\{0,\ldots,999\}^{12}$ y recupera sus tres residuos por ventana.

\subsubsection{Concordancia con la transformación armónica}

La transformación $\mathcal H_{12}$ actúa mediante $H_4$ en las
órbitas $(r,r+3,r+6,r+9)$, $r=0,1,2$, donde
\[
 H_4=\begin{pmatrix}
 1&1&1&1\\1&1&-1&-1\\1&-1&1&-1\\1&-1&-1&1
 \end{pmatrix},\qquad H_4^2=4I_4.
\]
Para una terna compatible, sean
\[
 B_r=\sum_{j=0}^3c_{r+3j},\qquad
 a_0=\frac{q+2B_0+B_1}{3},\quad
 a_1=a_0-B_0,\quad a_2=a_1-B_1.
\]
El lector armónico tiene bloques
\[
 (\Pi_H(b,c,q))^{(r)}
 =(a_r,b_{r+3}-b_{r+9},b_r+b_{r+6},b_r-b_{r+6}).
\]

\begin{proposition}[Concordancia de reconstrucciones]
\label{iii:r25:hadamard}
Sobre $\operatorname{im}\mathcal D$ se cumple
\[
 \Pi_H\mathcal D=\mathcal H_{12},\qquad
 \mathcal D^{-1}=\tfrac14\mathcal H_{12}\Pi_H.
\]
La imagen de $\Pi_H$ restringida a ese dominio es
\[
 \mathcal L_H=
 \{u\in\mathbb Z^{12}:\mathcal H_{12}u\equiv0\pmod4\}.
\]
\end{proposition}
\begin{proof}
Para $z=\mathcal D^{-1}(b,c,q)$, sean
$h_r=\sum_{j=0}^3z_{r+3j}$. El desplazamiento de cuatro posiciones
lleva cada órbita a la siguiente; por tanto $B_r=h_r-h_{r+1}$.
La identidad $q=h_0+h_1+h_2$ da $a_r=h_r$.
En una órbita escribamos $(x_0,x_1,x_2,x_3)$ para sus coordenadas
y $d_j=x_j-x_{j+1}$. Se verifican
\[
 \begin{aligned}
 d_1-d_3&=x_0+x_1-x_2-x_3,\\
 d_0+d_2&=x_0-x_1+x_2-x_3,\\
 d_0-d_2&=x_0-x_1-x_2+x_3.
 \end{aligned}
\]
Son las tres filas restantes de $H_4x$. La inversión y la
caracterización integral se deducen de $\mathcal H_{12}^2=4I$.
\end{proof}

La condición de congruencia armónica controla la salida de $\Pi_H$;
la pertenencia de los canales a $\operatorname{im}\mathcal D$
requiere además \eqref{iii:r25:relaciones}--\eqref{iii:r25:congruencia}.
Un control negativo exacto consiste en tomar $b=0$, $q=0$ y
$c=e_0-e_3$. Las tres sumas orbitales de $c$ se anulan, de modo
que $\Pi_H(b,c,q)=0$, pero se incumplen las relaciones de imagen.
La lectura armónica no sustituye el control de compatibilidad de
los canales.

\subsection{Conjugación, covariancia y memoria del registro incidencial}
\label{iii:r25:transporte}

Sea $\rho(m)=13-m$. Si la conjugación de incidencias invierte el
orden de las ventanas, conserva las clases aritmética y de frontera
de las visitas e intercambia los canales de las trazas, entonces
\[
 (A_m^\dagger,C_m^\dagger,V_m^\dagger)
 =(A_{\rho(m)},-C_{\rho(m)},V_{\rho(m)}).
\]
Con $c_m=[C_m]_{10}$, la normalización decimal da
\begin{equation}
 z_m^\dagger=z_{\rho(m)}
 +100\mathbf1_{\{c_{\rho(m)}\ne0\}}-20c_{\rho(m)}.
 \label{iii:r25:conjugacion}
\end{equation}
En efecto, $[-C]_{10}=0$ si $[C]_{10}=0$ y vale
$10-[C]_{10}$ en los demás casos. La negación del canal se
aplica antes de reconstruir el bloque decimal.

La aplicación a una involución concreta conserva también los
extremos de las ventanas. Para una ruta
$\gamma=(x_0,\ldots,x_n)$ leída antes del avance,
\begin{equation}
 \sum_{k=0}^{n-1}f(x_{n-k})
 =\sum_{k=0}^{n-1}f(x_k)+f(x_n)-f(x_0).
 \label{iii:r25:extremos}
\end{equation}
Ambas sumas contienen los mismos términos interiores; la diferencia
es el extremo final menos el inicial. En rutas cerradas esa
corrección se anula. En ventanas abiertas se incorpora antes de
reducir módulo $10$.

\begin{proposition}[Consecuencia de una reducción periódica]
\label{iii:r25:periodo}
Si el observable de una colección fija de rutas satisface
$o(t+54)=o(t)$, sus multiplicidades se conservan bajo ese retorno
y el mismo lector local actúa en ventanas de nueve instantes,
entonces
\[
 (A,C,V)_{m+6}=(A,C,V)_m,\qquad z_{m+6}=z_m.
\]
\end{proposition}
\begin{proof}
La traslación $t\mapsto t+54$ lleva cada ventana biyectivamente
a la situada seis posiciones después y conserva cada incidencia
contada y su multiplicidad. La igualdad se transmite a los recuentos
y a sus normalizaciones decimales.
\end{proof}

Una presentación con doce bloques distintos conserva, por tanto,
información que no factoriza por ese observable de período $54$:
selección de incidencias, orientación, extremos, multiplicidad o
memoria. La proposición delimita el efecto de esa reducción; no
atribuye período $54$ a la historia completa.

\subsubsection{Determinación sobre órbitas y acumulación de eventos}

Fijar solamente $Q(z)=q$ deja una clase afín del retículo
\[
 \mathcal Z_0=\{v\in\mathbb Z^{12}:\sum_i v_i=0\}
 =\bigoplus_{i=0}^{10}\mathbb Z(e_i-e_{11}).
\]
Por ejemplo, $100\mathbf1+e_0$ y $100\mathbf1+e_1$ tienen
carga $1201$ y satisfacen los controles de imagen mediante sus
respectivos canales. Son testigos de las ecuaciones de compatibilidad,
sin atribución de pertenencia a las historias terminales HMT.

\begin{proposition}[Covariancia cíclica]
\label{iii:r25:covariancia}
Sea $\rho$ la traslación de ventanas de un registro $R$, de
período mínimo $p\mid12$, y sea $S$ la traslación de sus
coordenadas. Un lector covariante sobre esa órbita queda determinado
por $z\in\operatorname{Fix}(S^p)$. La condición $Q(z)=q$
admite soluciones enteras si y sólo si $12/p$ divide a $q$;
cuando existen, forman una clase afín de rango $p-1$.
\end{proposition}
\begin{proof}
La regla $F(\rho^jR)=S^jz$ está bien definida exactamente si
$S^pz=z$. Estos vectores repiten $12/p$ veces un bloque de
longitud $p$, y su carga es $\frac{12}{p}\sum_{i=0}^{p-1}z_i$.
La fórmula prueba la divisibilidad y el rango. La covariancia de
los canales se sigue de $D_aS=SD_a$ y $QS=Q$.
\end{proof}

En la coordenatización armónica, la acción correspondiente es
$\widehat S_H=\mathcal H_{12}S\mathcal H_{12}/4$.
Si el origen está marcado, la acción transporta también esa marca.
El período de la órbita aquí indicada no se identifica con el de
la historia enriquecida. Para una lectura inicial $F_{108}$ ya
definida, la naturalidad bajo truncamiento se expresa mediante
$F_N=F_{108}\circ\tau_{N,108}$; esta identidad conserva la
evaluación inicial sin seleccionar su regla de producción.

\begin{proposition}[Composición aditiva de contribuciones incidenciales]
\label{iii:r25:eventos}
Supóngase que el estado inicial y cada evento $a_t$ producen,
mediante reglas previamente fijadas sobre el registro, contribuciones
$(\delta w_t,\delta q_t)$ tales que
\[
 \sum_i\delta w_{t,i}=0,
 \qquad \delta q_t+T(\delta w_t)\equiv0\pmod {12}.
\]
Con condiciones iniciales compatibles $(w^{(0)},q^{(0)})$,
la actualización por suma determina un único registro en cada etapa.
Su incremento es
\[
 \delta z_{t,i}=
 \frac{\delta q_t+T(\delta w_t)}{12}-P_i(\delta w_t).
\]
La construcción respeta la concatenación de segmentos y conserva
la lectura de un prefijo bajo prolongaciones posteriores.
\end{proposition}
\begin{proof}
Las condiciones de compatibilidad se conservan por suma y la
inversa \eqref{iii:r25:inversa} es aditiva en $(w,q)$.
La inducción fija cada estado. La suma sobre segmentos concatenados
es la suma de sus contribuciones; un prefijo recibe los mismos
sumandos antes y después de la prolongación.
\end{proof}

Esta proposición da un criterio suficiente para la acumulación
aditiva, no una exigencia de linealidad de todo lector TPK.
Cuando los eventos se transportan por acciones no triviales, sus
contribuciones se expresan en la coordenatización terminal mediante el cociclo
afín correspondiente. Las reglas de evento y el estado inicial
conservan en ambos casos su procedencia incidencial.

El desarrollo anterior determina el paso
\[
 \mathscr L\longmapsto(A,C,V)\longmapsto z
 \longmapsto(D_3z,D_4z,Q(z))\longmapsto\mathcal H_{12}z
\]
con las condiciones de cada aplicación. Su dominio es distinto del
registro regional del apéndice~\ref{iii:nucleo-finito}: allí las
ecuaciones $X_aL_a=Y_a$ determinan los levantamientos a partir de
transiciones ternarias. Aquí las incidencias de visitas y trazas
determinan bloques enteros y sus canales. La unicidad de una
reconstrucción no reemplaza la producción del registro que utiliza
la otra. Ambas conservan sus primitivas, operaciones y pruebas en
la cadena APP--TRIT--TPK.


\subsection{Dos lecturas del estado terminal}
\label{subsec:k-terminal}

Denotemos por \(x_{\mathrm{term}}\) el estado producido por la composición
de selección, transporte y actualización hasta el corte terminal. La
escritura operatoria es
\[
 x_{\mathrm{term}}=
 (\mathcal U_{108}\circ\cdots\circ\mathcal U_1)(x_{\mathrm{seed}}),
 \qquad \mathcal U_t=\mathsf{Upd}_t\circ\mathsf{Tra}_t\circ\mathsf{Sel}_t.
\]
Sus dos lecturas pertinentes son
\[
 x_{\mathrm{term}}
 \xrightarrow{(\pi_{\mathrm{term}},R_{\mathrm{sgn}})}
 (B_{\mathrm{term}},U).
\]
La matriz \(B_{\mathrm{term}}\) retiene la lectura ternaria visible; el
vector \(U\) conserva otra coordenada de la misma historia. No se inserta
entre ellos una flecha \(B_{\mathrm{term}}\to U\).

El selector terminal de la construcción de referencia~(subsección~\ref{subsec:k-terminal}) opera sobre catálogos
de tamaños 196, 197 y 196. Las firmas locales reducen sus
\(196\cdot197\cdot196=7\,567\,952\) ternas a 331, y el margen de
columnas \(q_{\partial}=(1,2,2,1,2,0)\) deja dos matrices:
\[
 B_0=\begin{pmatrix}0&2&0&2&2&0\\0&0&1&2&2&0\\1&0&1&0&1&0\end{pmatrix},
 \qquad
 B_1=\begin{pmatrix}0&2&1&0&2&0\\0&0&1&2&2&0\\1&0&0&2&1&0\end{pmatrix}.
\]
La orientación de filas es \(\sigma=(1,1,-1)=(1,1,2)\) en
\(\mathbb F_3^3\). Las cargas de filas de las candidatas son
\((0,2,0)\) y \((2,2,1)\). Sólo la segunda pertenece a
\(\langle\sigma\rangle=\{(0,0,0),(1,1,2),(2,2,1)\}\); por tanto
\(B_{\mathrm{term}}=B_1\). Esta comprobación prueba el último paso del
selector a partir del censo declarado. No convierte el margen de columnas
en una consecuencia de la carga de filas, ni reemplaza la enumeración
anterior por la inspección de esas dos matrices.

El lector firmado tiene la composición tipada
\begin{equation}
 R_{\mathrm{sgn}}
 =\Pi_H\circ\mathcal E_{108}^{90,120}\circ\mathcal R_{12}:
 \mathcal X_{\mathrm{TPK}}^{\mathrm{term,mem}}
 \longrightarrow\mathbb Z^{12}.
 \label{eq:k-lector-firmado}
\end{equation}
El registro de canales determina
\[
 \mathcal E_{108}^{90,120}(\mathcal R_{12}(x_{\mathrm{term}}))
 =(b^{90},b^{120},Q_{\mathrm{TPK}}),
\]
con las coordenadas
\begin{align*}
 b^{90}={}&(-495,-116,-684,108,601,-90,\\
            &\hspace{19mm}-173,-88,313,560,-397,461),\\
 b^{120}={}&(-425,-281,-481,671,-255,30,\\
             &\hspace{19mm}475,-543,680,251,6,-128),\\
 Q_{\mathrm{TPK}}={}&6263.
\end{align*}
Estas cantidades se utilizan como coordenadas del registro previo, no
como parámetros que se ajusten mediante alfa. La reconstrucción demostrada
en los apartados siguientes comienza exactamente en esta salida tipada.
El rastro material del extractor anterior y el alcance de su transcripción
autónoma se conservan en la nota técnica de procedencia; la aritmética
posterior no se presenta como sustituto de ese rastro.

\subsection{Del registro de canales al vector firmado}
\label{subsec:k-pi-h}

El lector armónico \(\Pi_H\) puede escribirse sin utilizar ni \(K\) ni
alfa. Para evitar confundir sus sumas orbitales con los bloques de alfa,
las denotaremos por \(s_1,s_2,s_3\). Sean
\[
 B_r=\sum_{j=0}^{3}b^{120}_{r+3j},\qquad
 d_{r,j}=b^{90}_{r+3j},\qquad1\leq r\leq3,
\]
y
\begin{equation}
 s_1=\frac{Q_{\mathrm{TPK}}+2B_1+B_2}{3},\qquad
 s_2=s_1-B_1,\qquad s_3=s_2-B_2.
 \label{eq:k-sumas-armonicas}
\end{equation}
La división por tres pertenece al lector de las tres órbitas. Su
integralidad se comprueba sobre las salidas compatibles del registro; no
se supone para un elemento arbitrario de \(\mathbb Z^{25}\).

Los tres bloques firmados son
\begin{equation}
 U^{(r)}=(s_r,d_{r,1}-d_{r,3},d_{r,0}+d_{r,2},d_{r,0}-d_{r,2}).
 \label{eq:k-bloques-firmados}
\end{equation}
La sustitución entera da
\[
 (B_1,B_2,B_3)=(972,-1073,101),\qquad
 (s_1,s_2,s_3)=(2378,1406,2479),
\]
y, por tanto,
\begin{align*}
 U^{(1)}&=(2378,-452,-668,-322),\\
 U^{(2)}&=(1406,998,-204,-28),\\
 U^{(3)}&=(2479,-551,-371,-997).
\end{align*}
Intercalando por columnas se obtiene
\begin{equation}
 U=(2378,1406,2479,-452,998,-551,
       -668,-204,-371,-322,-28,-997).
 \label{eq:k-u-firmado}
\end{equation}
La presencia de valores negativos y de componentes superiores a 999
muestra que \(U\) no es una palabra de dígitos en base mil. Tampoco es
\(U_6\Vert U_6\), concatenación catalogal de período seis. La diferencia
es de dominio y de información conservada, no solamente de valores.

El orden constructivo hasta aquí es
\[
 x_{\mathrm{term}}\longrightarrow\mathcal R_{12}
 \longrightarrow(b^{90},b^{120},Q_{\mathrm{TPK}})
 \longrightarrow U.
\]
La inversión que sigue produce el registro dodecafásico. Que después podamos recuperar
los canales desde el registro dodecafásico constituye un control de reversibilidad; no
invierte este orden causal.

\subsection{Hadamard por órbitas y reconstrucción entera}
\label{subsec:k-hadamard}

En el ciclo de doce posiciones, el paso tres determina tres órbitas:
\[
 (1,4,7,10),\qquad(2,5,8,11),\qquad(3,6,9,12).
\]
La separación de cuatro sectores de cada órbita es la que, una vez fijados
origen y orientación, admite la lectura angular de noventa grados. La
transformación no actúa sobre tres cuaternas contiguas del orden natural.

Sea
\[
 H_4=\begin{pmatrix}
 1&1&1&1\\1&1&-1&-1\\1&-1&1&-1\\1&-1&-1&1
 \end{pmatrix}.
\]
Si \(P_{\mathrm{orb}}\) reordena por las tres órbitas indicadas,
definimos
\[
 \mathcal H_{12}=P_{\mathrm{orb}}^{-1}
 (H_4\oplus H_4\oplus H_4)P_{\mathrm{orb}}.
\]

\begin{lemma}[Inversión de Hadamard e integralidad]
\label{lem:k-hadamard}
Se tiene
\[
 H_4^2=4I_4,\qquad H_4^{-1}=\frac14H_4,\qquad\det H_4=-16.
\]
Para \(u\in\mathbb Z^4\), la preimagen \(H_4^{-1}u\) es entera si y
sólo si \(H_4u\equiv0\pmod4\).
\end{lemma}

\begin{proof}
Las filas son ortogonales y tienen norma cuadrada cuatro; la matriz es
simétrica. Esto da \(H_4^2=4I_4\) y la inversa. Una eliminación entera,
o la expansión del determinante, da el signo negativo y el valor \(-16\).
La condición de integralidad es exactamente la divisibilidad por cuatro
de las cuatro coordenadas de \(H_4u\).
\end{proof}

\begin{definition}[Transformación integral dodecafásica]
La subred de registros firmados y su coordenatización integral son
\[
 \mathcal L_H=\{u\in\ZZ^{12}:\mathcal H_{12}u\equiv0\pmod4\},
 \qquad \mathscr K(u)=\tfrac14\mathcal H_{12}u.
\]
La aplicación \(\mathscr K:\mathcal L_H\to\ZZ^{12}\) es un isomorfismo
de grupos abelianos, con inversa \(k\mapsto\mathcal H_{12}k\).
El registro canónico \(K\) es la imagen del registro firmado generado
por las operaciones anteriores, con origen de fase y orientación fijados.
\end{definition}

\begin{proof}
La congruencia define exactamente la integralidad de \(\mathscr K(u)\).
La identidad \(\mathcal H_{12}^2=4I\) demuestra las dos composiciones
inversas. En particular, todo \(k\in\ZZ^{12}\) tiene preimagen
\(\mathcal H_{12}k\in\mathcal L_H\).
\end{proof}

\begin{proposition}[Registro integral canónico]
\label{prop:k-sello}
La inversión \(K^{(r)}=H_4U^{(r)}/4\) de los bloques de
\eqref{eq:k-bloques-firmados} produce
\begin{align*}
 K^{(1)}&=(234,729,621,794),\\
 K^{(2)}&=(543,659,58,146),\\
 K^{(3)}&=(140,824,914,601).
\end{align*}
En el orden natural, el registro dodecafásico es
\begin{equation}
 K=(234,543,140,729,659,824,621,058,914,794,146,601).
 \label{eq:k-vector}
\end{equation}
Es la única preimagen racional de \(U\), y pertenece a
\(\{0,\ldots,999\}^{12}\).
\end{proposition}

\begin{proof}
Los tres productos antes de dividir son
\begin{align*}
 H_4U^{(1)}&=(936,2916,2484,3176),\\
 H_4U^{(2)}&=(2172,2636,232,584),\\
 H_4U^{(3)}&=(560,3296,3656,2404).
\end{align*}
Todas las coordenadas son divisibles por cuatro. Dividir y deshacer la
permutación de órbitas da \eqref{eq:k-vector}. La unicidad procede de la
invertibilidad racional, no de una búsqueda entre candidatos próximos a
un valor físico.
\end{proof}

La integralidad también puede formularse en términos de redes. Los
divisores determinantes de \(H_4\) son \(1,2,4,16\): las entradas
tienen máximo común divisor uno; los menores de orden dos son pares y
alguno tiene módulo dos; los de orden tres son múltiplos de cuatro y
alguno tiene módulo cuatro. De ahí se deduce
\[
 \operatorname{SNF}(H_4)=\operatorname{diag}(1,2,2,4),
\]
y para la transformación global,
\[
 \operatorname{coker}\mathcal H_{12}
 \cong(\mathbb Z/2\mathbb Z)^6\oplus(\mathbb Z/4\mathbb Z)^3.
\]
El índice de la imagen es \(16^3=4096\). Así, no todo vector firmado
entero admite una preimagen entera; el vector obtenido satisface las
congruencias precisas que la inversión exige. El factor \(1/4\) está
forzado por la matriz, y no tiene el estatuto de un coeficiente libre.

\subsection{Diferencias cíclicas y componente uniforme}
\label{subsec:k-transversal}

Con índices cíclicos módulo doce, sean
\[
 (D_3z)_m=z_m-z_{m+3},\qquad(D_4z)_m=z_m-z_{m+4},
 \qquad Q(z)=\sum_{m=1}^{12}z_m.
\]
El primer operador compara sectores separados por tres posiciones; el
segundo, por cuatro. Los nombres de canal 90 y 120 se refieren aquí a
esas separaciones en el ciclo orientado, no a una identificación con un
operador físico de otro dominio.

\begin{proposition}[Completitud del sistema transversal]
\label{prop:k-transversal}
El operador
\[
 \mathcal D=(D_3,D_4,Q):\mathbb Z^{12}\longrightarrow\mathbb Z^{25}
\]
tiene rango doce y factores invariantes no nulos
\[
 \operatorname{diag}(1,\ldots,1,12),
\]
con once unidades. En particular, sus salidas determinan un único registro dodecafásico.
\end{proposition}

\begin{proof}
Si \(D_3z=D_4z=0\), el vector es constante a lo largo de ambos pasos.
Como \(3\) y \(4\) generan el ciclo módulo doce,
\(\ker D_3\cap\ker D_4=\langle\mathbf1\rangle\). La carga total
elimina esta última libertad porque \(Q(\mathbf1)=12\).

Para la afirmación integral, las filas de \((D_3,D_4)\) son incidencias
orientadas de un grafo conexo. Un árbol generador proporciona once
factores invariantes unitarios. Todo menor de orden doce debe contener la
fila \(Q\). Sumando las primeras once columnas a la última, esta última
se anula en las filas de incidencia y vale doce en la fila \(Q\). Todos
los menores máximos son divisibles por doce y el árbol produce uno de
módulo exactamente doce. El último factor es, por tanto, doce.
\end{proof}

La comprobación sobre \eqref{eq:k-vector} da
\[
 D_3K=b^{90},\qquad D_4K=b^{120},\qquad Q(K)=6263.
\]
Las fórmulas \eqref{eq:k-sumas-armonicas} y
\eqref{eq:k-bloques-firmados} reconstruyen a su vez \(U\) desde esas
diferencias. Se cierra así un triángulo de representaciones exactas del
mismo objeto: registro transversal, vector firmado y registro dodecafásico. La igualdad
de resultados no convierte sus codominios en uno solo; demuestra que
las aplicaciones declaradas recuperan la información pertinente.

\subsection{Dos lecturas racionales: \texorpdfstring{\(\kappa_{36}\)}{kappa36} y \texorpdfstring{\(\kappa_{\mathrm{per}}\)}{kappa periódica}}
\label{subsec:k-kappas}

Una vez construido \(K\), la coordenatización en base mil permite dos operaciones
distintas. La primera lee una sola vuelta; la segunda prolonga el registro dodecafásico
periódicamente. Sea \(B=1000\) y definamos su entero de concatenación
\begin{equation}
 N_K=\sum_{m=1}^{12}K_mB^{12-m}
 =234543140729659824621058914794146601.
 \label{eq:k-numerador}
\end{equation}
Los ceros iniciales de un bloque, como el de 058, son parte de su escritura
de tres cifras. No cambian su valor entero, pero sí permiten concatenar
sin ambigüedad los doce bloques.

\begin{definition}[Lecturas finita y periódica del registro dodecafásico]
\label{def:k-kappas}
La lectura de una vuelta es
\[
 \kappa_{36}=\sum_{m=1}^{12}K_mB^{-m}=\frac{N_K}{B^{12}}.
\]
La prolongación periódica es
\[
 K_j^{\mathrm{per}}=K_{1+((j-1)\bmod12)},\qquad
 \kappa_{\mathrm{per}}=\sum_{j\geq1}K_j^{\mathrm{per}}B^{-j}.
\]
\end{definition}

\begin{proposition}[Valor y período exactos del registro dodecafásico prolongado]
\label{prop:k-periodo}
Se cumple
\begin{equation}
 \kappa_{\mathrm{per}}=\frac{N_K}{B^{12}-1},\qquad
 \kappa_{\mathrm{per}}-\kappa_{36}
 =\frac{N_K}{B^{12}(B^{12}-1)}
 =B^{-12}\kappa_{\mathrm{per}}.
 \label{eq:k-diferencia-kappas}
\end{equation}
La palabra infinita del registro dodecafásico tiene período mínimo doce en base mil y
período mínimo treinta y seis en base diez.
\end{proposition}

\begin{proof}
La suma de la primera vuelta es \(N_KB^{-12}\); cada nueva vuelta
multiplica su contribución por \(B^{-12}\). La serie geométrica produce
\(N_K/(B^{12}-1)\), y la resta da la segunda identidad.

Los doce componentes de \(K\) son distintos. Un período propio del ciclo
identificaría dos de ellos; por tanto el período mínimo de bloques es
doce. La expansión decimal es canónica, pues el registro dodecafásico no tiene una cola
de bloques 999. Si su período mínimo decimal fuera \(d\), dividiría
36. Agrupar de tres en tres produciría un período de bloques
\(d/\gcd(d,3)\). Ningún divisor propio de 36 da doce mediante esa
operación. Luego el período decimal mínimo es 36.
\end{proof}

Ambas lecturas son racionales, pero no son el mismo racional. La primera
termina después de doce bloques; la segunda repite la vuelta. Además,
\(0<\kappa_{\mathrm{per}}-\kappa_{36}<B^{-12}\). Su cercanía no
permite sustituir una por otra en una identidad exacta. En particular,
el cierre finito de alfa utiliza \(\kappa_{36}\), mientras que la
suma de la sección prolongada utiliza \(\kappa_{\mathrm{per}}\).

\input{sections/k_reversibilidad.tex}

\subsection{Periodicidad y conúcleo del operador de transporte}
\label{subsec:k-clase-acarreo}

La periodicidad que acabamos de demostrar pertenece al registro dodecafásico. No implica
periodicidad de las coordenadas regionales, de los acarreos de la clausura
ni de alfa. Es útil precisar esta separación mediante un segundo
cociente, puramente integral.

Sea \(S\) el desplazamiento cíclico de doce coordenadas, con orientación
fijada, y sea \(b\geq2\) una base entera. El operador de acarreo
periódico es \(\partial_b=S-bI\). Si una identidad periódica se escribe
\[
 A=P+E-\Phi-K+\partial_bC,
\]
entonces las transformaciones
\[
 K\longmapsto K+\partial_bh,\qquad C\longmapsto C+h
\]
conservan su lado derecho. Esta invariancia compara representantes de una
clase de acarreo; no es una autorización para modificar el registro dodecafásico ya
seleccionado por el registro.

\begin{proposition}[Cociente del acarreo periódico]
\label{prop:k-cociente-acarreo}
Con la orientación anterior,
\[
 \operatorname{coker}(S-bI)\cong\mathbb Z/(b^{12}-1)\mathbb Z.
\]
\end{proposition}

\begin{proof}
El módulo cíclico se presenta como
\(\mathbb Z[t]/(t^{12}-1)\). Imponer \(S=bI\) equivale a imponer
\(t=b\), y deja la relación \(b^{12}-1=0\). La evaluación de los
coeficientes en \(b\), con el orden de concatenación elegido, representa
la clase resultante.
\end{proof}

El denominador \(B^{12}-1\) de \(\kappa_{\mathrm{per}}\) y este
cociente proceden de la misma longitud cíclica, pero sus objetos siguen
siendo distintos: uno es una evaluación racional y el otro un cociente
de un módulo de acarreos. La clausura finita de alfa será una cadena
orientada con frontera derecha. Su prolongación transportará una
frontera remota compatible, no un acarreo que se identifique por decreto
entre la primera y la decimotercera posición.

Queda así preparado el capítulo siguiente. El registro dodecafásico se ha fijado antes
de resolver la recurrencia de alfa; su reconstrucción por Hadamard es
entera y única; sus diferencias transversales recuperan el registro; y
sus dos evaluaciones racionales tienen dominios y usos separados. Ninguno
de estos hechos utiliza alfa como entrada, y ninguno decide todavía el
tipo aritmético de la salida de la clausura completa.

\section{Derivación de la constante de estructura fina}
\label{sec:alpha}

La constante de estructura fina recibe en este desarrollo una determinación
aritmética que reúne las tres coordenadas regionales y el registro integral
dodecafásico. La operación combina sus bloques con orientación fijada y
transporta los cocientes de la normalización posicional. El parámetro
resultante expresa la compatibilidad de esa diferencia orientada con
la prolongación de los cilindros. Su valor y los datos de incidencia
utilizados en su construcción permanecerán relacionados durante el
desarrollo excepcional.

El orden interno es importante. Las coordenadas \(\pi\),
\(e\) y \(\varphi\), ya construidas en la sección precedente, son coordenadas
previas del mismo estado aritmético con memoria de trayectoria, y pueden actuar legítimamente como
entradas de un lector posterior. No son constantes convencionales
inyectadas para seleccionar alfa. La generación correlacionada de las cuatro coordenadas comparte origen y
compatibilidad coinductiva, sin que ello imponga evaluar sus cuatro
componentes en un solo paso.

Las dos vías expresan funciones matemáticas diferentes del mismo parámetro.
La primera realiza una clausura entera por división euclídea en base
\(1000\). La segunda estudia el equilibrio analítico de vacancias mediante
el resolvente \(E_{21/22}\) y sus prolongaciones graduadas. Una vía
determina bloques y cocientes; la otra determina una raíz mediante una
ecuación y un dominio de unicidad. Su identificación se formula después
en los cilindros comunes a toda profundidad. La exposición distingue
estos tres resultados: construcción posicional, determinación analítica
y compatibilidad de ambas.

\subsection{Primera vía: determinación posicional por clausura entera}
\subsubsection{Dominio y orientación de la clausura entera}
\label{subsec:alpha-dominio}

Sean \(P_j,E_j,\Phi_j\in\{0,\ldots,999\}\) los bloques canónicos de
las tres coordenadas regionales, y prolonguemos el registro dodecafásico por
\[
 K_j=K_{1+((j-1)\bmod12)}.
\]
La orientación del canal es \((1,1,-1)\). La contribución del registro dodecafásico se
sustrae según la clausura dodecafásica ya tipada. Para \(B=1000\), la
operación local distingue el balance firmado
\[
 Z_j:=P_j+E_j-\Phi_j-K_j
\]
de la suma que incorpora el cociente entrante. La normalización es
\begin{equation}
 s_j=Z_j+c_{j+1},\qquad
 A_j=s_j-B\left\lfloor\frac{s_j}{B}\right\rfloor,
 \qquad c_j=\left\lfloor\frac{s_j}{B}\right\rfloor.
 \label{eq:alpha-recurrencia}
\end{equation}
Así, \(A_j\in\{0,\ldots,B-1\}\), incluso cuando \(s_j\) es
negativo. El acarreo que entra desde la derecha pertenece al dominio de
la operación; no puede omitirse porque la tabla se imprima de izquierda
a derecha.

\begin{definition}[Clausura de profundidad finita]
\label{def:alpha-finita}
Para una profundidad \(N\) y una frontera entera \(t\), la aplicación
\(\mathcal C_N\) recibe
\((P_{1:N},E_{1:N},\Phi_{1:N},K_{1:N},c_{N+1}=t)\) y aplica
\eqref{eq:alpha-recurrencia} en el orden \(N,N-1,\ldots,1\). Su salida
es la pareja de palabras
\[
 (A_{1:N},c_{1:N+1}).
\]
\end{definition}

La división euclídea no posee una elección de signo residual: el resto se
fija en \(\{0,\ldots,999\}\) y el cociente queda determinado. Por
ejemplo, \(-431=569-1000\) y \(-1200=800-2\cdot1000\). Los
acarreos negativos de la clausura no son anomalías; son la información
entera necesaria para mantener la igualdad orientada.

\subsubsection{Ventana dodecafásica y condiciones de acarreo}
\label{subsec:alpha-ventana}

Las primeras doce tríadas obtenidas son
\begin{align*}
 P={}&(141,592,653,589,793,238,462,643,383,279,502,884),\\
 E={}&(718,281,828,459,045,235,360,287,471,352,662,497),\\
 \Phi={}&(618,033,988,749,894,848,204,586,834,365,638,117).
\end{align*}
En el canal de clausura, las cinco regiones que determinan la coordenada
\(P\) conservan sus antecedentes distintos. Compartir esta lectura no
las identifica como regiones o como historias.

Con frontera \(c_{13}=0\), la recurrencia produce la tabla siguiente.
Se muestra cada entrada, el acarreo entrante, el total antes de normalizar,
el bloque obtenido y el acarreo saliente. La tabla es un cálculo entero
reproducible, no una lista de bloques seleccionados por su parecido con
un objetivo.

\begin{table}[htbp]
\centering
\small
\setlength{\tabcolsep}{3.5pt}
\begin{tabular}{r|rrrr|rr|rr}
\hline
\(j\)&\(P_j\)&\(E_j\)&\(\Phi_j\)&\(K_j\)&\(c_{j+1}\)&\(s_j\)&\(A_j\)&\(c_j\)\\
\hline
1&141&718&618&234&0&7&007&0\\
2&592&281&033&543&0&297&297&0\\
3&653&828&988&140&\(-1\)&352&352&0\\
4&589&459&749&729&\(-1\)&\(-431\)&569&\(-1\)\\
5&793&045&894&659&\(-2\)&\(-717\)&283&\(-1\)\\
6&238&235&848&824&\(-1\)&\(-1200\)&800&\(-2\)\\
7&462&360&204&621&0&\(-3\)&997&\(-1\)\\
8&643&287&586&058&\(-1\)&285&285&0\\
9&383&471&834&914&\(-1\)&\(-895\)&105&\(-1\)\\
10&279&352&365&794&0&\(-528\)&472&\(-1\)\\
11&502&662&638&146&0&380&380&0\\
12&884&497&117&601&0&663&663&0\\
\hline
\end{tabular}
\caption{Clausura inicial con frontera derecha nula. La dependencia se propaga de la última fila a la primera.}
\label{tab:alpha-ventana}
\end{table}

La palabra de acarreos es
\[
 (c_1,\ldots,c_{13})=(0,0,0,-1,-1,-2,-1,0,-1,-1,0,0,0),
\]
y la lectura de la ventana es el racional
\begin{equation}
 \alpha_{12}^{(0)}
 =\frac{7297352569283800997285105472380663}{10^{36}}
 =0.007297352569283800997285105472380663.
 \label{eq:alpha-ventana-cero}
\end{equation}
El superíndice recuerda la frontera elegida. Esta ventana es exacta en su
dominio finito; no se la identifica sin más con el truncamiento de la
sección infinita. La diferencia entre ambas lecturas se hará visible al
transportar el acarreo de la cola remota.

\subsubsection{Telescopado y reversibilidad en una fibra de frontera}
\label{subsec:alpha-telescopia}

La forma local de la igualdad es
\begin{equation}
 P_j+E_j-\Phi_j-K_j-A_j=Bc_j-c_{j+1}.
 \label{eq:alpha-identidad-local}
\end{equation}
El término de acarreo es un enlace entre posiciones contiguas. Si se lo
elimina coordenada a coordenada, se obtiene una igualdad falsa. Sólo una
evaluación que conserve sus extremos permite telescoparlo.

\begin{proposition}[Identidad entera de profundidad arbitraria]
\label{prop:alpha-telescopia}
Para \(\iota_N(z)=\sum_{j=1}^{N}z_jB^{N-j}\), toda ejecución de
\(\mathcal C_N\) satisface
\begin{equation}
 \iota_N(P)+\iota_N(E)-\iota_N(\Phi)-\iota_N(K)-\iota_N(A)
 =B^Nc_1-c_{N+1}.
 \label{eq:alpha-telescopia}
\end{equation}
\end{proposition}

\begin{proof}
Multiplicar \eqref{eq:alpha-identidad-local} por \(B^{N-j}\) y sumar
produce en el lado derecho
\[
 \sum_{j=1}^{N}(Bc_j-c_{j+1})B^{N-j}=B^Nc_1-c_{N+1}.
\]
Cada acarreo interior aparece dos veces con signo opuesto; sólo sobreviven
los dos extremos.
\end{proof}

Para la tabla~\ref{tab:alpha-ventana}, ambos extremos son cero. Si
\(p_{12},e_{12},f_{12}\) son las tres sumas fraccionarias truncadas,
\[
 \alpha_{12}^{(0)}=p_{12}+e_{12}-f_{12}-\kappa_{36}.
\]
La constante del registro dodecafásico que aparece es \(\kappa_{36}\), no
\(\kappa_{\mathrm{per}}\). El dominio de esta identidad contiene
truncamientos y una frontera finita; su forma no autoriza a borrar ambas
condiciones al pasar a la sección completa.

La inversión local es también exacta:
\[
 K_j=P_j+E_j-\Phi_j+c_{j+1}-A_j-Bc_j.
\]
Dados \(P,E,\Phi\), los bloques \(A\) y los acarreos completos,
recupera el registro dodecafásico. La concatenación recupera además
\[
 \iota_N(K)=\iota_N(P)+\iota_N(E)-\iota_N(\Phi)-\iota_N(A)
             -B^Nc_1+c_{N+1}.
\]
Esta reversibilidad es una propiedad de la fibra con datos y fronteras
fijados. No debe convertirse en una nueva receta causal que empiece por
alfa y fabrique después el registro. La dirección de construcción sigue
siendo registro firmado, registro dodecafásico, clausura y salida.

\subsubsection{Prolongación y normalización canónica de la sección completa}
\label{subsec:alpha-completa}

La prolongación conserva el registro dodecafásico periódico y recibe las secuencias
regionales completas. Sean
\[
 p=\sum_{j\geq1}P_jB^{-j},\qquad
 e_0=\sum_{j\geq1}E_jB^{-j},\qquad
 f=\sum_{j\geq1}\Phi_jB^{-j}.
\]
Cada serie es una realización del carácter interno correspondiente.
Las partes enteras ya obtenidas son 3, 2 y 1, de modo que
\(p=\pi-3\), \(e_0=e-2\) y
\(f=\varphi-1\).

\begin{proposition}[Existencia y unicidad de la normalización completa]
\label{prop:alpha-normalizacion}
La serie
\begin{equation}
 y=p+e_0-f-\kappa_{\mathrm{per}}
   =\sum_{j\geq1}(P_j+E_j-\Phi_j-K_j)B^{-j}
 \label{eq:alpha-serie}
\end{equation}
converge absolutamente. En el dominio de los bloques obtenidos,
\(0<y<10^{-2}\). Existe una única expansión canónica \((A_j)\) de
\(y\), sin cola finalmente constante 999, y una única sucesión entera
acotada \((c_j)\) que satisface \eqref{eq:alpha-identidad-local} con
\(c_1=0\). Estos acarreos pertenecen al conjunto invariante
\(\mathcal C=\{-2,-1,0,1,2\}\).
\end{proposition}

\begin{proof}
Los coeficientes de la serie están acotados por \(2(B-1)\) en módulo,
de modo que domina una serie geométrica. Si
\(Z_j=P_j+E_j-\Phi_j-K_j\) y
\(R_j(Z)=\sum_{r\geq0}Z_{j+r}B^{-r-1}\), cada cola es la diferencia
de dos sumas de dos colas canónicas y pertenece a \((-2,2)\).
El primer coeficiente es
\(Z_1=141+718-618-234=7\); por tanto
\(y=(7+R_2(Z))/1000\in(0.005,0.009)\).

Sea \((A_j)\) la expansión canónica de \(y\). Definamos
\[
 c_j=R_j(Z)-R_j(A).
\]
Al multiplicar la igualdad de las dos series por \(B^{j-1}\) y separar
sus primeros \(j-1\) términos se obtiene
\[
 c_j=\sum_{r=1}^{j-1}(A_r-Z_r)B^{j-1-r}\in\mathbb Z,
\]
y \(c_1=0\). Además, \(R_j(A)\in[0,1)\), luego
\(-3<c_j<2\), lo que incluye todos los acarreos en \(\mathcal C\).
Las identidades de desplazamiento de las colas dan
\(Bc_j=Z_j-A_j+c_{j+1}\), que es la recurrencia requerida.

Recíprocamente, cualquier solución acotada con salida canónica y
\(c_1=0\) telescopa en el límite a la misma serie \(y\). La expansión
canónica es única y, una vez fijada, la fórmula de las colas determina
todos los acarreos. Finalmente, si se aplica una ejecución finita con
\(c_{j+1}\in\mathcal C\), se tiene \(-2000\leq s_j\leq2000\),
por lo que su cociente euclídeo vuelve a pertenecer a \(\mathcal C\).
\end{proof}

La proposición explicita la normalización del lector posterior: no
reemplaza la generación de las tres secuencias por una serie convencional
introducida como dato. Una vez producidas esas secuencias por el estado
coinductivo, su combinación y la frontera de la normalización quedan
determinadas sin consultar alfa.

El procedimiento finito conservado en el corpus propaga las cinco
fronteras de \(\mathcal C\) y conserva solamente el prefijo común a
todas ellas. Si dos profundidades han coalescido antes de una frontera,
la recurrencia determinista de derecha a izquierda da los mismos bloques
a su izquierda. Ésa es la compatibilidad con el truncamiento. No se
escoge una frontera porque dé cifras favorables, ni se deduce la
compatibilidad a toda profundidad de haber observado una corrida larga.

La comparación inicial muestra por qué esta precaución importa. En la
prolongación estabilizada se obtiene \(c_{13}^{\infty}=-1\), mientras
que la ventana aislada imponía \(c_{13}=0\). Así,
\[
 884+497-117-601-1=662
\]
en el último bloque de la primera vuelta. Las primeras once tríadas no
cambian, pero la prolongación comienza
\[
 (007,297,352,569,283,800,997,285,105,472,380,662,999,\ldots).
\]
Una aparición del bloque 999 no viola la convención canónica: lo
excluido es una cola que sea finalmente constante 999. La ventana con
frontera cero termina en 663; el truncamiento de la sección completa
termina en 662. Ambas afirmaciones describen operaciones diferentes y
exactas. No deben fusionarse en una sola notación de proyección.

\subsubsection{Definición operativa y ecuación aditiva completa}
\label{subsec:alpha-definicion}

\begin{definition}[Coordenada alfa de la clausura]
\label{def:alpha-operativa}
La vía entera determina la coordenada
\begin{equation}
 \alpha_A=\sum_{j\geq1}A_jB^{-j}
 =p+e_0-f-\kappa_{\mathrm{per}},
 \label{eq:alpha-via-a}
\end{equation}
donde \((A_j,c_j)\) es la normalización canónica compatible de
\eqref{eq:alpha-recurrencia}. En la realización escalar de las
coordenadas HMT,
\begin{equation}
 \boxed{\alpha_A=
 \pi+e-\varphi
 -4-\kappa_{\mathrm{per}}.}
 \label{eq:alpha-identidad-completa}
\end{equation}
\end{definition}

El entero cuatro no se calibra: es
\(3+2-1\), combinación de las partes enteras de las tres coordenadas.
La base mil procede de la coordenatización cúbica y el período doce del registro dodecafásico. Los
signos son los del canal orientado y del cierre, y la frontera remota es
la que impone la sección compatible. Éstos son los orígenes efectivos de
los coeficientes de la ecuación.

En este nivel puede describirse alfa como la coordenada escalar del
defecto de clausura entre las tres realizaciones regionales y el retorno
racional dodecafásico, calculado con acarreo compatible. El término
\emph{defecto} designa una diferencia algebraica determinada, no una
pérdida arbitraria ni una imperfección del sistema. Esta descripción no
suprime la genealogía: sin las regiones, el registro firmado, la
orientación y la ley de acarreo, la igualdad escalar no sería una
exposición suficiente de la construcción.

También deben mantenerse separados tres objetos que la notación de alfa
puede acompañar en lectores posteriores: \(A\) es una palabra de
bloques; el soporte negativo de preacarreo se obtiene del signo de los
totales antes de normalizar; y un vector radial excepcional pertenece a
un codominio de incidencia. Ninguno de ellos es intercambiable con el
valor real \(\alpha_A\).

\input{sections/revision_alpha_notacion.tex}
\input{sections/revision_alpha_frontera.tex}

\input{sections/alpha_dependencias.tex}

\subsection{Segunda vía: determinación analítica mediante vacancias}
\subsubsection{Invariantes del transporte y resolvente graduado}
\label{subsec:alpha-via-b}

La vía analítica recibe invariantes del estado orientado, no la salida
\(\alpha_A\). Su truncamiento de orden seis es
\begin{equation}
 P_6(x)=2\pi x-\frac74x^2
 +\frac{x^3}{2\pi}
 +\frac{x^4}{20}-\frac{2x^5}{21}-\frac{x^6}{46},
 \label{eq:alpha-p6}
\end{equation}
y la condición de vacancia utiliza
\[
 \delta=\log_{10}\frac{10}{9}.
\]
La razón \(10/9\) compara la coordenatización decimal con las nueve posiciones
APP de la construcción nonádica. Su logaritmo pertenece a esta lectura
posterior de compactificación; no es un valor objetivo de alfa.

La genealogía de los coeficientes debe conservarse completa, pero con su
fuerza lógica exacta. Los desarrollos recientes reúnen una firma marcada
del estado. La matriz de sincronización y el censo asociado son
\[
 D_{\mathrm{syn}}=\begin{pmatrix}85&112\\41&52\end{pmatrix},
 \qquad N_9=43.
\]
De su determinante se extrae
\[
 a=-\frac{\det D_{\mathrm{syn}}}{N_9}
   =\frac{172}{43}=4.
\]
El reloj inducido sobre los huecos cortos tiene espectro secundario
\(\{6,7\}\); su extremo superior da \(b=7\). La cardinalidad
trítica \(m=3\) determina
\[
 t_*=mb=21,\qquad k_*=mb-1=20.
\]
La lectura de vacancia en la coordenatización mil da
\[
 v_-=\lfloor1000\delta\rfloor=45,\qquad
 v_+=\lceil1000\delta\rceil=46.
\]
Estos enteros son antecedentes tipados de la evaluación del truncamiento. No se
obtienen de su raíz ni de una aproximación metrológica a ella.

La rama inferior coincide además con el determinante del bloque APP
\[
 B_c=\begin{pmatrix}7&2\\2&7\end{pmatrix},\qquad\det B_c=45.
\]
El rango transversal del registro del capítulo anterior da
\(r_{\mathrm{dod}}=11\), y la semicorona de fase proporciona
\(f_{\mathrm{ph}}=8\binom62=120\). En particular,
\(45r_{\mathrm{dod}}=495\).

Con \(\omega=2\pi\), la firma
\[
 \mathcal I_9=(\omega,m,a,b,k_*,t_*,v_-,v_+,f_{\mathrm{ph}},r_{\mathrm{dod}})
\]
alimenta el lector graduado
\begin{align}
 \mathfrak E_9(\mathcal I_9)(x)
 ={}&\omega x-\frac ba x^2+\omega^{-1}x^3
 +\frac{x^4}{k_*}-\frac{2x^5}{t_*}-\frac{x^6}{v_+}\notag\\
 &-\frac{x^7}{f_{\mathrm{ph}}}+\frac{x^8}{v_-}
 +\frac{2x^9}{v_-r_{\mathrm{dod}}}.
 \label{eq:alpha-firma-coeficientes}
\end{align}
La sustitución recupera los coeficientes de \(P_6\) y de su
prolongación de orden nueve. El par \(45/46\) atraviesa la separación
entre ambos: la rama superior ocupa el grado seis y la inferior reaparece
en los grados ocho y nueve.

La factorización de \eqref{eq:alpha-firma-coeficientes} se formula
para un estado que conserva origen, orden de canales y orientación.
En ese dominio marcado, cada coeficiente queda fijado por
\((\omega,m,a,b,k_*,t_*,v_-,v_+,f_{\mathrm{ph}},r_{\mathrm{dod}})\).
Dos lectores que preserven la firma y apliquen la misma evaluación
graduada coinciden término a término hasta el grado nueve. Esta
unicidad no se extiende a lectores que olviden esos datos.

\subsubsection{Resolvente nilpotente en los grados \(7\), \(8\) y \(9\)}
\label{subsec:alpha-resolvente}

La primera prolongación admite una construcción operatoria finita
especialmente explícita. Sobre el espacio con base \(e_7,e_8,e_9\),
sea
\[
 Ne_7=-\frac83e_8,\qquad Ne_8=\frac2{11}e_9,\qquad Ne_9=0,
 \qquad v_7=-\frac1{120}e_7.
\]
El símbolo \(N\) designa aquí un operador nilpotente, no una profundidad
de truncamiento. Su dominio y graduación están fijados antes de evaluar
el polinomio en un número.

\begin{lemma}[Prolongación nilpotente exacta]
\label{lem:alpha-resolvente}
Se cumple \(N^3=0\) y
\[
 (I-N)^{-1}v_7
 =-\frac{e_7}{120}+\frac{e_8}{45}+\frac{2e_9}{495}.
\]
La evaluación \(e_n\mapsto x^n\) produce
\begin{equation}
 T_{7:9}(x)=-\frac{x^7}{120}+\frac{x^8}{45}+\frac{2x^9}{495}.
 \label{eq:alpha-t79}
\end{equation}
\end{lemma}

\begin{proof}
El operador aumenta el grado y se anula en el último vector de la base,
por lo que su cubo es cero. La inversa es el resolvente finito
\(I+N+N^2\). Además,
\[
 Nv_7=\frac1{45}e_8,\qquad N^2v_7=\frac2{495}e_9.
\]
Sumar los tres términos y evaluar da la fórmula.
\end{proof}

Aquí los denominadores 120, 45 y 495 no son ajustes decimales: proceden
de la semicorona, del determinante APP y de su producto por el rango
dodecafásico. Las transferencias \(-8/3\) y \(2/11\), junto con el
signo y el grado de la semilla, pertenecen al lector tipado. Conservar
sólo la simetría del ciclo y sustituir este operador por otro operador
simétrico no conserva necesariamente los coeficientes.

\subsubsection{Raíz de multiplicidad uno del truncamiento de orden nueve}
\label{subsec:alpha-jet}

Definamos
\[
 P_9=P_6+T_{7:9},\qquad E_{21/22}^{(9)}(x)=P_9(x)-\delta,
 \qquad I_\alpha=(0,10^{-2}).
\]
Este operador finito tiene un resultado propio de existencia y unicidad.

Como comprobación posterior de los caracteres ya construidos, el
aislamiento racional puede reproducirse sin una tabla decimal. Sean
\[
 A_n(t)=\sum_{j=0}^{n}\frac{(-1)^jt^{2j+1}}{2j+1},\qquad
 L_n(t)=2\sum_{j=0}^{n}\frac{t^{2j+1}}{2j+1}\quad(0<t<1).
\]
El criterio alternante y la suma geométrica de la cola dan
\begin{align}
 A_{2n+1}(t)&<\arctan t<A_{2n}(t),
 \label{eq:alpha-cotas-arctan}\\
 L_n(t)&<\log\frac{1+t}{1-t}
 <L_n(t)+\frac{2t^{2n+3}}{(2n+3)(1-t^2)}.
 \label{eq:alpha-cotas-log}
\end{align}
Aplicamos \eqref{eq:alpha-cotas-arctan} a la identidad de Machin
\(\pi=16\arctan(1/5)-4\arctan(1/239)\), y
\eqref{eq:alpha-cotas-log} a
\[
 \log(10/9)=\log\frac{1+1/19}{1-1/19},\qquad
 \log 10=\log(10/9)+2\log\frac{1+1/2}{1-1/2}.
\]
Escribimos
\[
 \delta:=\frac{\log(10/9)}{\log 10}.
\]
Con los órdenes de truncamiento \(80,20,40,120\), aplicados
respectivamente a las cuatro series anteriores, se obtienen intervalos
racionales para \(\pi\) y \(\delta\), ambos de anchura menor que
\(10^{-70}\). Estas cotas verifican una realización posterior; no
seleccionan el estado APP--TRIT--TPK ni reciben un valor de alfa.

\begin{proposition}[Raíz de multiplicidad uno del truncamiento]
\label{prop:alpha-raiz-jet}
El operador \(E_{21/22}^{(9)}\) tiene una única raíz
\(\xi_9\) en \(I_\alpha\), y esa raíz es simple. Además posee el aislamiento racional
\begin{equation}
 \frac{729735256928380099}{10^{20}}
 <\xi_9<\frac{729735256928380100}{10^{20}}.
 \label{eq:alpha-intervalo-jet}
\end{equation}
\end{proposition}

\begin{proof}
Las cotas de las coordenadas internas
\(3.14<\pi<3.15\) y
\(0.045<\delta<0.046\) son suficientes para la existencia. En cero,
\(E_{21/22}^{(9)}(0)=-\delta<0\). En \(10^{-2}\), la contribución
lineal menos los términos negativos es mayor que \(0.0626\), luego el
valor del operador es positivo.

Para \(0\leq x\leq10^{-2}\), al descartar los términos positivos de
la derivada se obtiene
\[
 (E_{21/22}^{(9)})'(x)
 \geq2\pi-\frac72\cdot10^{-2}
 -\frac{10}{21}\cdot10^{-8}-\frac6{46}\cdot10^{-10}
 -\frac7{120}\cdot10^{-12}>6.24.
\]
La continuidad da existencia; la cota derivativa da unicidad y
simplicidad. Para el extremo izquierdo \(l\) y el extremo derecho \(u\)
de \eqref{eq:alpha-intervalo-jet}, la sustitución racional de las cotas
\eqref{eq:alpha-cotas-arctan}--\eqref{eq:alpha-cotas-log} en el polinomio
da, con redondeo dirigido,
\[
 -4.6\,10^{-20}<E_{21/22}^{(9)}(l)<-4.5\,10^{-20}<0,
 \qquad
 0<1.6\,10^{-20}<E_{21/22}^{(9)}(u)<1.8\,10^{-20}.
\]
La monotonía ya probada sitúa la única raíz entre ambos extremos.
El intervalo se verifica después de construir el operador y no
interviene en la selección de sus coeficientes.
\end{proof}

El resultado concierne a \(\xi_9\). No autoriza a escribir
\(\xi_9=\alpha_A\) como una igualdad exacta de números completos.
La raíz de un truncamiento puede compartir un prefijo con la raíz del
operador prolongado y diferir en una posición posterior. La simplicidad
permite controlar el movimiento de la raíz bajo perturbaciones; no
anula esas perturbaciones.

\subsection{Criterio de identificación de las lecturas compatibles}
\label{subsec:alpha-limite-dos-vias}

La identificación de dos lecturas completas requiere conservar sus cilindros
y los mapas de restricción. Para una colección graduada compatible,
\[
 E_{21/22}^{(6)}\longleftarrow E_{21/22}^{(9)}
 \longleftarrow E_{21/22}^{(12)}\longleftarrow\cdots,
 \qquad
 E_{21/22}=\varprojlim_m E_{21/22}^{(6+3m)},
\]
el límite designa el lector con sus transportes, no una cola seleccionada
libremente. El criterio siguiente separa la consecuencia de la
compatibilidad de la construcción efectiva que la establece después.

\begin{proposition}[Identificación de los límites compatibles]
\label{prop:alpha-dos-limites}
Si los valores de las vías entera y analítica pertenecen a los
mismos cilindros \(C_N^\alpha\) para todo \(N\), y
\(\operatorname{diam}C_N^\alpha\to0\), entonces
\[
 \alpha_A=\alpha_B=\alpha.
\]
\end{proposition}

\begin{proof}
Para todo \(N\), la distancia entre ambos valores está acotada
por \(\operatorname{diam}C_N^\alpha\). Hacer tender \(N\) a infinito
da distancia cero. La existencia de las secciones compatibles asegura
que no se trata de la intersección vacía de un sistema de aproximaciones.
\end{proof}

La diferencia puede expresarse algebraicamente. Sea
\(F=\mathbb Q(\pi,\delta)\) y sea
\(j_9:F[[x]]\to F[x]/(x^{10})\) el truncamiento. Entonces
\[
 \ker j_9=x^{10}F[[x]].
\]
Todos los operadores
\[
 E_h(x)=E_{21/22}^{(9)}(x)+x^{10}h(x)
\]
tienen el mismo truncamiento de orden nueve. Sin embargo, tomando
\(h=0\) y \(h=1/729\),
\[
 E_{1/729}(\xi_9)=\frac{\xi_9^{10}}{729}>0,
\]
de modo que no comparten esa raíz. La observación no permite elegir una
cola ajena al transporte HMT; precisamente demuestra por qué la cola
generada y su compatibilidad deben conservarse. Tampoco basta una serie
formal arbitraria para hablar de raíces analíticas sin su dominio de
convergencia o sin la realización del sistema de cilindros.

\input{sections/alpha_publicacion_completa.tex}

\section{Prolongación excepcional: de la incidencia dodecafásica a Moonshine}
\label{exc:seccion}

La expresión numérica no agota el estado que la ha generado. En la
construcción precedente, APP conserva las dos hojas y la división con
residuo y cociente; el TRIT determina el régimen local y su orientación;
el TPK compone selección, transporte, acarreo, memoria y retorno. La
lectura arquimediana contrae cilindros compatibles. La lectura que ahora
estudiamos conserva, en cambio, la incidencia de las posiciones y su
transporte entre marcos. Ambas parten del mismo estado aritmético con
memoria de trayectoria. No se
reconstruye una geometría excepcional a partir de los decimales finales
de $\pi$, $\varphi$, $e$ o $\alpha$.

Este punto permite precisar el sentido de la prolongación
$\alpha/K\longrightarrow\text{incidencia excepcional}$. La flecha no
afirma que un número real aislado determine un código, un retículo o un
grupo. Transporta el registro dodecafásico y los datos firmados que han
intervenido en la determinación de $\alpha$, junto con las palabras y los
marcos procedentes de APP--TRIT--TPK. En particular, el registro ordenado
$K$ puede recuperarse desde su constante racional $\kappa$ con las
convenciones de base, longitud y origen ya fijadas. La proyección a un
soporte conserva, en cambio, sólo las incidencias seleccionadas.
Por su parte, $K_{\rm ph}$
designa otra lectura, de memoria reducida, posterior a la holonomía
$\Gamma_9$.

La exposición separa tres operaciones matemáticas: construir los objetos
finitos y sus mapas; prolongarlos a un retículo y a un álgebra de operadores
de vértice; reconocer después los objetos obtenidos mediante los teoremas
clásicos pertinentes. Los nombres Paley, Witt, Mathieu, Golay, Leech y
Monstruo designan esos reconocimientos y las realizaciones explícitas
que los hacen verificables. No son selectores retrospectivos de las
semillas ni de los coeficientes de la generación.

\input{sections/revision_incidencia_argumento.tex}



\subsection{Compatibilidad de las realizaciones arquimediana e incidencial}
\label{exc:doble-lectura}

Sea $X_n^{\rm cal}$ el espacio de prefijos admisibles de longitud $n$ del
estado aritmético con memoria de trayectoria, con truncaciones
$p_{n,m}:X_n^{\rm cal}\to X_m^{\rm cal}$ para $m\leq n$. El superíndice
recuerda que se conserva el marco inicial declarado y su transporte, no
que se ajuste un parámetro contra un valor experimental. Escribimos
$(w_j,f_j)$ para la palabra visible de seis trits y el marco en el paso
$j$. Los marcos satisfacen una ley causal
\[
 f_{j+1}=g_j(x)f_j,
\]
donde $g_j(x)$ depende del prefijo ya construido. En consecuencia, una
prolongación no altera los marcos anteriores.

La cadena
\[
 w_6\longrightarrow w_{12}\longrightarrow w_{18}
 \longrightarrow w_{24}\longrightarrow w_{30}
 \longrightarrow R_{36}\longrightarrow G_9
\]
se conserva en esta lectura. El retorno de la fase no reinicia ni el
marco ni la memoria del estado. Asimismo, se mantienen distintos el
censo de $468$ emisiones decimales, las $243$ palabras visibles de seis
trits y el ambiente $\mathbb F_3^6$, de $729$ palabras. El ambiente se
utilizará para construir un código lineal; no sustituye al dominio
admisible de emisiones.

La lectura arquimediana toma, para un bloque $b$ de seis trits, su
dirección entera $\nu(b)\in\{0,\ldots,728\}$. Un prefijo determina
\[
 q_n=m+\sum_{j=0}^{n-1}\frac{\nu(b_j)}{729^{j+1}},
 \qquad I_n=[q_n,q_n+729^{-n}].
\]
La compatibilidad de prefijos da cilindros encajados cuyo diámetro tiende
a cero. Su clase de Cauchy pertenece primero a la completación interna
$\mathbb R_{\rm HMT}$; su identificación posterior con la recta real
usual no interviene en la elección de los bloques.

Para la lectura incidencial introducimos, sobre el mismo prefijo, el
mapa de gráfico
\begin{equation}
 \gamma_A:\mathbb F_3^6\longrightarrow\mathbb F_3^{12},
 \qquad w\longmapsto(w,wA).
 \label{exc:gamma}
\end{equation}
Las palabras se escriben como vectores fila. Si $A^f=fAf^{-1}$, el
cambio de marco satisface la identidad concreta
\begin{equation}
 \gamma_{A^f}(wf^{-1})
   =\gamma_A(w)(f^{-1}\oplus f^{-1}).
 \label{exc:covariancia}
\end{equation}
El codominio calibrado conserva $f$ junto a la palabra codificada.
Definimos entonces
\begin{equation}
 Q_n(x)=\bigl((f_j,\gamma_{f_jA_Wf_j^{-1}}(w_j))\bigr)_{0\leq j<n}.
 \label{exc:qn}
\end{equation}

\begin{proposition}[Compatibilidad y prolongación de la lectura incidencial]
\label{exc:limite-inverso}
Si $r_{n,m}$ trunca las sucesiones codificadas, se cumple
\[
 r_{n,m}Q_n=Q_m p_{n,m}.
\]
Por tanto, los $Q_n$ determinan un único mapa $Q_\infty$ entre los
respectivos límites inversos. En cada marco, el mapa conserva exactamente
la palabra visible, pues $\operatorname{pr}_1\gamma_A=\operatorname{id}$.
\end{proposition}

\begin{proof}
Dos prefijos compatibles tienen las mismas palabras y los mismos
operadores de transporte hasta el índice $m-1$. Tienen, por inducción,
los mismos marcos en esos índices. Las primeras $m$ componentes de
\eqref{exc:qn} coinciden, lo que prueba el cuadrado. La propiedad
universal del límite inverso da $Q_\infty$. La afirmación de recuperación
es la proyección sobre las seis primeras coordenadas de cada gráfico.
\end{proof}

Esta inyectividad local tiene un dominio preciso. No demuestra que
$Q_\infty$ recupere toda la memoria profunda del estado TPK si sólo se
ha retenido la sucesión de palabras visibles y marcos. Mucho menos
demuestra inyectividad después de pasar de palabras a soportes.
Además, un cambio general de base en $\operatorname{GL}_6(\mathbb F_3)$
no conserva por sí mismo el peso de Hamming ordinario. En una coordenatización
general se transportan también la incidencia y la métrica de la coordenatización
inicial. Los cálculos de soportes que siguen se realizan en la coordenatización
Paley declarada; los relabelados monomiales conservan directamente su
interpretación coordenada.

\subsection{La realización finita y el código ternario}
\label{exc:paley-codigo}

El operador elíptico del TRIT admite la siguiente realización finita en
la coordenatización de seis posiciones. Éstas coordinatizan las seis unidades
$\{1,2,4,5,7,8\}$ módulo nueve, conservadas por el transporte modular.
Sobre $\mathbb P^1(\mathbb F_5)$, una vez
declarado el orden $\infty,0,1,2,3,4$, sea $C$ la matriz con diagonal
nula, $C_{\infty,a}=C_{a,\infty}=1$ y
$C_{a,b}=\chi(b-a)$ en las cinco posiciones finitas. Aquí
$\chi(1)=\chi(4)=1$, $\chi(2)=\chi(3)=-1$ y $\chi(0)=0$.
Es una realización en coordenadas de la estructura cuadrática ya
transportada; la identificación de las seis posiciones con la recta
proyectiva es un dato de coordenatización, no una nueva entrada numérica en el
generador.

Las identidades de la suma cuadrática dan $C^2=5I_6$. Al reducir módulo
tres obtenemos
\begin{equation}
 A_W=
 \begin{pmatrix}
 0&1&1&1&1&1\\
 1&0&1&2&2&1\\
 1&1&0&1&2&2\\
 1&2&1&0&1&2\\
 1&2&2&1&0&1\\
 1&1&2&2&1&0
 \end{pmatrix},
 \qquad A_W^2=2I_6=-I_6.
 \label{exc:aw}
\end{equation}
La igualdad cuadrática realiza el régimen elíptico; no identifica el
cuerpo $\mathbb F_3$ con una realización arquimediana de la unidad
imaginaria. Son realizaciones distintas de una relación operatoria
común, con sus correspondientes dominios.

\Needspace{8\baselineskip}
\begin{definition}[Código de la coordenatización dodecafásica]
\label{exc:codigo}
El código ternario asociado a \eqref{exc:aw} es
\[
 C_W=\{(w,wA_W):w\in\mathbb F_3^6\}
       \subset\mathbb F_3^{12}.
\]
Su conjunto de hexadas es
\[
 \mathcal H=
 \{\operatorname{supp}(c):c\in C_W,
                           \operatorname{wt}(c)=6\}.
\]
\end{definition}

\begin{proposition}[Autodualidad, distancia y diseño de Witt]
\label{exc:codigo-witt}
El código $C_W$ es autodual de parámetros $[12,6,6]_3$. Su enumerador
de pesos es
\begin{equation}
 \sum_{c\in C_W}z^{\operatorname{wt}(c)}
   =1+264z^6+440z^9+24z^{12}.
 \label{exc:enumerador}
\end{equation}
El conjunto $\mathcal H$ tiene $132$ elementos y constituye un sistema
$S(5,6,12)$: cada conjunto de cinco posiciones está contenido en una
única hexada.
\end{proposition}

\begin{proof}
La matriz generadora $[I_6\mid A_W]$ tiene rango seis y, por simetría
de $A_W$, satisface
\[
 [I_6\mid A_W][I_6\mid A_W]^{\mathsf T}
   =I_6+A_W^2=0.
\]
El código es autoortogonal de dimensión mitad de la longitud y, por
tanto, autodual. La evaluación finita de las $3^6$ palabras da
\eqref{exc:enumerador}; es una enumeración sobre la matriz explícita,
sin valores reales ni objetivos externos. En particular, la distancia
mínima es seis.

Si dos palabras de peso seis tienen el mismo soporte, entre sus seis
cocientes de coordenadas no nulas uno de los signos $1,-1$ aparece al
menos tres veces. Restando el múltiplo correspondiente se obtendría
una palabra no nula de peso a lo sumo tres, imposible. Luego un soporte
corresponde exactamente a la pareja $\{c,-c\}$ y hay $264/2=132$
soportes. Si dos soportes distintos compartieran cinco posiciones,
el mismo argumento cancelaría al menos tres coordenadas de una unión
de tamaño siete, produciendo peso a lo sumo cuatro. Por ello una
quintupla pertenece a lo sumo a una hexada. Finalmente,
\[
 132\binom65=792=\binom{12}5,
\]
de modo que pertenece a una y sólo una.
\end{proof}

El reconocimiento posterior es el código de Golay ternario extendido y
el pequeño diseño de Witt \cite{conwaysloane}. La prueba anterior explica qué objeto se
reconoce. No reemplaza la genealogía APP--TRIT--TPK por una búsqueda
entre códigos conocidos.

\subsection{Bandera dodecafásica asociada a \texorpdfstring{$K$ y $\alpha$}{K y alfa}}
\label{exc:bandera}

\input{sections/incidencia_registro.tex}

En el marco heredado del registro dodecafásico, los datos son
\begin{align}
 K={}&(234,543,140,729,659,824,621,\mathtt{058},914,794,146,601),
 \label{exc:k}\\
 w_\pi={}&010211, &w_e={}&201101,\nonumber\\
 w_\varphi={}&121200, &w_{\rm APP}={}&022020.\nonumber
\end{align}
Al aplicar $\gamma_{A_W}$ se obtienen las siguientes incidencias:
\begin{align*}
 P_\pi&=\{2,4,5,6,7,8,9,10,11\},\\
 H_e&=\{1,3,4,6,9,10\},\\
 H_\varphi&=\{1,2,3,4,7,9\},\\
 H_{\rm APP}&=\{2,3,5,10,11,12\}.
\end{align*}
La primera es el soporte de una palabra de peso nueve; las otras tres
son hexadas. El registro aporta, por su lectura de bloque alto,
\begin{equation}
 B_K=\{i:K_i\geq729\}=\{4,6,9,10\}=H_e\cap P_\pi.
 \label{exc:bk}
\end{equation}
El entero $729$ señala aquí una frontera del formato de bloques
declarado. No se identifica por el mero cardinal con el ambiente
ternario usado en \eqref{exc:codigo}. En este registro concreto, todos
los umbrales enteros entre $660$ y $729$ producen el mismo soporte:
el soporte no recupera el umbral ni los doce bloques.

La vía aritmética de $\alpha$ conserva, antes del acarreo, el vector
firmado
\begin{equation}
 Z_0=(7,297,353,-430,-715,-1199,-3,286,-894,-528,380,663).
 \label{exc:precarry}
\end{equation}
Su soporte negativo es
\begin{equation}
 H^-_\alpha=\{4,5,6,7,9,10\}.
 \label{exc:halpha}
\end{equation}
No se extraen estos signos de la expansión decimal del escalar
$\alpha$; se conservan en el estado que precede a su expresión numérica.

\input{sections/incidencia_normalizacion.tex}

\begin{proposition}[Coincidencia de los dos selectores de la bandera]
\label{exc:selector-bandera}
Dentro de $\mathcal H$, el soporte \eqref{exc:halpha} es la única hexada
$H$ que satisface
\[
 B_K\subset H\subset P_\pi,
 \qquad
 (|H\cap H_e|,|H\cap H_\varphi|,|H\cap H_{\rm APP}|)=(4,3,2).
\]
Así, la lectura firmada y la lectura incidencial seleccionan la misma
bandera $B_K\subset H^-_\alpha$.
\end{proposition}

\begin{proof}
Las cuatro hexadas que contienen $B_K$ se obtienen añadiendo,
respectivamente, los pares
\[
 \{1,3\},\qquad\{2,11\},\qquad\{5,7\},\qquad\{8,12\}.
\]
Sólo los pares segundo y tercero están contenidos en $P_\pi$. Para el
segundo, la intersección con $H_\varphi$ tiene tamaño tres y con
$H_{\rm APP}$ tamaño tres; no cumple la última condición. Para el
tercero, las tres intersecciones tienen tamaños $4,3,2$. Este último
produce precisamente el soporte negativo de \eqref{exc:precarry}.
\end{proof}

Conviene conservar tanto la bandera como su marco marcado:
\[
 \mathcal I_*=(B_K\subset H^-_\alpha),\qquad
 \widetilde{\mathcal I}_*
   =(\mathcal I_*;P_\pi,H_e,H_\varphi,H_{\rm APP}).
\]
Las operaciones de unión, intersección y diferencia son equivariantes
bajo un relabelado simultáneo del diseño. La clase marcada, y no los
números de las posiciones tomados aisladamente, es el objeto transportado.

\begin{lemma}[Origen seleccionado por la incidencia marcada]
\label{exc:origen}
Para el marco anterior,
\[
 o_*=(H_e\cap H_\varphi)
           \setminus(H^-_\alpha\cup H_{\rm APP})=\{1\}.
\]
La regla es equivariante bajo todo automorfismo aplicado simultáneamente
a los datos de $\widetilde{\mathcal I}_*$.
\end{lemma}

\begin{proof}
La primera intersección es $\{1,3,4,9\}$. La unión que se sustrae
contiene $3,4,9$ y no contiene $1$. La equivariancia se sigue de que una
biyección conmuta con las operaciones de conjuntos.
\end{proof}

Este origen seleccionará una dirección del vecino reticular cuando se
declaren la métrica $A_2$ y la cámara positiva. La incidencia por sí sola
no contiene esa métrica. El paso reticular conservará explícitamente
ambos datos.

\input{sections/incidencia_bandera.tex}

\subsection{Entrelazamiento isométrico de las representaciones de incidencia}
\label{exc:entrelazador}

La incidencia define un mapa lineal cuya acción puede comprobarse
posición por posición:
\begin{equation}
 I:\mathbb R^{12}\longrightarrow\mathbb R^{\mathcal H},
 \qquad (Iv)(H)=\sum_{p\in H}v_p.
 \label{exc:incidencia}
\end{equation}
Sea $V_{11}=\mathbf1^\perp\subset\mathbb R^{12}$. Las dos realizaciones
euclídeas son aquí lenguajes posteriores para el mismo diseño finito.

\begin{proposition}[Isometría de la pantalla dodecafásica]
\label{exc:isometria}
Se cumple
\[
 I^{\mathsf T}I=36I_{12}+30J_{12},
\]
donde $J_{12}$ es la matriz cuyas entradas son uno. En consecuencia,
\[
 U_W=\frac16 I\big|_{V_{11}}:
 V_{11}\xrightarrow{\ \sim\ }E_{\rm inc}:=I(V_{11})
\]
es una isometría equivariante para las acciones por permutación de los
automorfismos del diseño.
\end{proposition}

\begin{proof}
Cada punto pertenece a
$r=\binom{11}{4}/\binom54=66$ hexadas, y cada pareja a
$\lambda_2=\binom{10}{3}/\binom43=30$. Son, respectivamente, las
entradas diagonales y no diagonales de $I^{\mathsf T}I$. Sobre
$\mathbf1^\perp$, el término $J_{12}$ se anula y queda $36I$.
Por último, si $g$ es un automorfismo, $p\in H$ equivale a
$g(p)\in g(H)$; por ello $I\rho_{12}(g)=\rho_{\mathcal H}(g)I$.
\end{proof}

La información de incidencia también tiene una lectura espectral
concreta. Definamos
\[
 G_{H,H'}=\binom{|H\cap H'|}{4}.
\]

\Needspace{8\baselineskip}
\begin{lemma}[Acción del núcleo de intersección]
\label{exc:espectro}
Si $J_{132,12}$ es la matriz rectangular de unos,
\[
 GI=30I+15J_{132,12}.
\]
Por tanto, $G$ actúa por el escalar $30$ en $E_{\rm inc}$.
\end{lemma}

\begin{proof}
La entrada $(H,p)$ de $GI$ cuenta parejas $(T,H')$ con
$T\subset H\cap H'$, $|T|=4$ y $p\in H'$. Si $p\in H$, hay diez
cuaternas $T$ que contienen $p$, cada una contenida en cuatro hexadas,
y cinco que no lo contienen, cada una con una única extensión de la
quintupla $T\cup\{p\}$. El total es $10\cdot4+5=45$.
Si $p\notin H$, las quince cuaternas de $H$ dan una extensión cada
una: el total es $15$. De ahí la identidad. El término rectangular
se anula en $V_{11}$.
\end{proof}

Aquí sí existe un entrelazador material: tiene dominio, codominio,
entradas y normalización determinados. Por ejemplo, cualquier
proyector ortogonal $P$ ya construido sobre $V_{11}$ se transporta a
$Q=U_WPU_W^*$ sobre $E_{\rm inc}$, con
$U_WP=QU_W$. Esta afirmación no identifica por su sola traza a todos
los operadores que actúan en esos espacios.

\subsection{La estrella de Witt y la generación de \texorpdfstring{$M_{12}$}{M12}}
\label{exc:estrella}

\begin{definition}[Estrella sobre una cuaterna]
\label{exc:def-estrella}
Para $B\in\binom{[12]}4$, la estrella de Witt es la colección de las
cuatro hexadas que contienen $B$. Sus pétalos son los pares $H\setminus B$.
La permutación $\sigma_B$ fija $B$ e intercambia los dos puntos de cada
pétalo.
\end{definition}

La existencia de cuatro hexadas se deduce de
\[
 \lambda_4=\frac{\binom81}{\binom21}=4.
\]
Dos pétalos no pueden compartir un punto, pues entonces dos hexadas
contendrían la misma quintupla. Son, por tanto, cuatro pares disjuntos
que particionan $[12]\setminus B$. Esta descripción no es una estrella
de cinco puntas ni una metáfora geométrica: es una incidencia
cuatro-a-uno del diseño.

\begin{figure}[htbp]
\centering
\begin{tikzpicture}[x=1cm,y=1cm, every node/.style={font=\small}]
 \node[draw,rounded corners,minimum width=3.8cm,minimum height=.7cm]
  (b) at (0,0) {$B_K=\{4,6,9,10\}$};
 \node[draw,minimum width=2.7cm,minimum height=.65cm]
  (p1) at (-4.6,-1.7) {$\{1,3\}$};
 \node[draw,minimum width=2.7cm,minimum height=.65cm]
  (p2) at (-1.55,-1.7) {$\{2,11\}$};
 \node[draw,very thick,minimum width=2.7cm,minimum height=.65cm]
  (p3) at (1.55,-1.7) {$\{5,7\}$};
 \node[draw,minimum width=2.7cm,minimum height=.65cm]
  (p4) at (4.6,-1.7) {$\{8,12\}$};
 \draw (b) -- (p1); \draw (b) -- (p2);
 \draw[very thick] (b) -- (p3); \draw (b) -- (p4);
 \node at (-4.6,-2.35) {$H_e$};
 \node at (-1.55,-2.35) {$B_K\cup\{2,11\}$};
 \node at (1.55,-2.35) {$H^-_\alpha$};
 \node at (4.6,-2.35) {$B_K\cup\{8,12\}$};
 \node at (0,-3.05)
  {$\sigma_{B_K}=(1\ 3)(2\ 11)(5\ 7)(8\ 12)$};
\end{tikzpicture}
\caption{La estrella concreta de la bandera dodecafásica. Cada línea
indica la unión de la cuaterna central con el par de su extremo; la
línea gruesa señala la hexada seleccionada también por los signos
anteriores al acarreo de $\alpha$. Los cuatro pares agotan las ocho
posiciones exteriores.}
\label{exc:fig-estrella}
\end{figure}

\begin{proposition}[La colección de estrellas]
\label{exc:m12}
Las $\binom{12}4=495$ permutaciones $\sigma_B$ preservan $\mathcal H$.
El grupo que generan tiene orden $95040$ y es el grupo de Mathieu
$M_{12}$ en su acción sobre el pequeño diseño de Witt.
\end{proposition}

\begin{proof}
La afirmación de preservación es una comprobación finita sobre las
$132$ hexadas de \eqref{exc:codigo}: para cada cuaterna se construyen
los cuatro pétalos y se verifica
$\{\sigma_B(H):H\in\mathcal H\}=\mathcal H$.
La clausura por composición de estas permutaciones produce $95040$
elementos. En la enumeración lexicográfica basta añadir sucesivamente
seis estrellas no contenidas aún en el grupo; los órdenes intermedios
son $2,6,18,360,7920,95040$, y las $495$ estrellas pertenecen al grupo
final. El procedimiento y los seis generadores están explicitados en
la nota técnica de procedencia. Finalmente se aplica el reconocimiento
del grupo de automorfismos de $S(5,6,12)$ como $M_{12}$, cuyo orden es
$12\cdot11\cdot10\cdot9\cdot8=95040$.
\end{proof}

Una estrella individual genera un grupo cíclico de orden dos. No
genera por sí sola $M_{12}$, y menos aún el Monstruo. Tampoco la
existencia de $495$ generadores de un grupo de permutaciones prueba que
cada uno sea el transporte de un lazo específico del TPK. Esa afirmación
requiere el mapa de lazos correspondiente; el resultado aquí establecido
es la acción incidencial finita.

\begin{remark}[La igualdad de trazas no identifica operadores]
Sobre $V_{11}$, una estrella tiene espectro $+1$ con multiplicidad siete
y $-1$ con multiplicidad cuatro. Su traza normalizada es $3/11$.
Un proyector de rango tres en el mismo espacio tiene también traza
normalizada $3/11$, pero espectro $1^3,0^8$. No son conjugados: sus
polinomios mínimos y sus espectros son diferentes. El entrelazador
\eqref{exc:incidencia} transporta operadores definidos; no convierte
esta coincidencia escalar en una equivalencia de operadores.
\end{remark}

\subsection{El residuo binario y las cinco regiones de \texorpdfstring{$\pi$}{pi}}
\label{exc:cinco-regiones}

El siguiente reconocimiento utiliza el código de Golay binario
extendido $G_{24}$, de parámetros $[24,12,8]_2$. Sus pesos son
$0,8,12,16,24$, y los soportes de peso ocho forman el diseño
$S(5,8,24)$. Fijamos una dodecada $D$, soporte de una palabra de peso
doce. Este dato y el isomorfismo de incidencia que se declarará a
continuación pertenecen a la coordenatización de realización; no se usan para
generar ni seleccionar los decimales de $\pi$.

\begin{proposition}[Residuo de una dodecada y elevación de hexadas]
\label{exc:residuo-binario}
El conjunto
\[
 \mathcal H(D)=
 \{O\cap D: O\text{ es una octada},\ |O\cap D|=6\}
\]
es un sistema $S(5,6,12)$ sobre $D$. Cada una de sus hexadas tiene
una única elevación a una octada. Para cualquier cuaterna $B\subset D$,
las cinco octadas que la contienen se distribuyen en cuatro elevaciones
residuales y una única octada $O_B^\perp$ con $O_B^\perp\cap D=B$.
\end{proposition}

\begin{proof}
Una quintupla $A\subset D$ pertenece a una única octada $O$. Si
$j=|O\cap D|$, la palabra binaria suma de $O$ y $D$ tiene peso
$20-2j$. Entre los pesos permitidos, y con $0\leq j\leq8$, sólo
son posibles $j=2,4,6$. Como $A\subset O\cap D$, se tiene $j\geq5$
y por tanto $j=6$. Esto da existencia y unicidad de la hexada
residual, así como unicidad de su elevación al escoger cinco de sus
puntos. Por otra parte,
\[
 \lambda_4(S(5,8,24))=\frac{\binom{20}1}{\binom41}=5,
 \qquad
 \lambda_4(S(5,6,12))=4.
\]
Las cuatro hexadas residuales sobre $B$ dan cuatro octadas distintas.
La quinta no puede cortar $D$ en seis puntos, pues sería otra hexada
residual; como contiene $B$, su intersección tiene tamaño cuatro y
coincide con $B$.
\end{proof}

Una coordenatización $\ell:[12]\to D$ que preserve las hexadas identifica
$\mathcal H$ con $\mathcal H(D)$. Su existencia es el reconocimiento
del diseño ya construido como el pequeño diseño de Witt; una elección
concreta de $\ell$ fija el relabelado. La elevación
\[
 \iota_D:\mathcal H\longrightarrow
       \{\text{octadas de }G_{24}\},
 \qquad H\longmapsto O\text{ tal que }O\cap D=\ell(H),
\]
es inyectiva porque la intersección con $D$ recupera $\ell(H)$.
Esta inyectividad tiene por dominio las hexadas, no el registro $K$.

La relación con las cinco regiones de $\pi$ conserva datos más ricos
que el número cinco. Las emisiones y sus firmas son:
\begin{center}
\small
\begin{tabular}{@{}crll@{}}
\toprule
Región & Multiplicidad & Emisión $U_6$ & Firma \\
\midrule
$\perp$ & $432$ & $(501,614,498,169,272,272)$ & $555555$ \\
$++_1$ & $144$ & $(810,923,870,169,272,272)$ & $339555$ \\
$++_2$ & $144$ & $(810,983,810,169,272,272)$ & $393555$ \\
$++_3$ & $144$ & $(870,923,810,169,272,272)$ & $933555$ \\
$--$ & $144$ & $(870,923,810,418,674,521)$ & $933717$ \\
\bottomrule
\end{tabular}
\end{center}
En cada bloque decimal de tres cifras, la firma usa la diferencia
modular de las dos primeras cifras. Las cinco emisiones tienen la
misma palabra trítica visible $w_\pi=010211$, pero no la misma
genealogía regional. Las tres regiones positivas conservan la cola
$169|272|272$; la región negativa la cambia por $418|674|521$;
la región transversal tiene firma constante y multiplicidad mayor.

La extensión marcada de la coordenatización regional envía la región transversal
a $O_{\ell(B_K)}^\perp$, la negativa a $\iota_D(H^-_\alpha)$ y las
tres positivas a las otras tres elevaciones de la estrella. Así se
realiza la partición $4+1$ de las cinco octadas sobre $\ell(B_K)$.
La correspondencia individual entre los tres nombres positivos y
los tres pétalos restantes exige un orden de coordenatización; no la determina
la incidencia sin ese dato. Por ello la afirmación precisa es una
correspondencia de regiones marcada, compatible con sus roles y con
la bandera, no una biyección canónica deducida del cardinal cinco.

La identidad $432+4\cdot144=1008$ conserva las multiplicidades
regionales. También $1008=36\cdot28$ y una octada tiene
$\binom82=28$ parejas. Estas igualdades son controles aritméticos de
catálogo, no sustitutos del mapa $\iota_D$ ni de la coordenatización regional.

\input{sections/estrella_pentadica_recuperada.tex}

\subsection{Del código ternario al retículo de Niemeier}
\label{exc:niemeier}

El paso a dimensión veinticuatro no necesita recorrer primero el
código binario. Se realiza directamente desde $C_W$ por pegado de
discriminantes. Así quedan distinguidas dos prolongaciones compatibles:
una conserva el residuo $W_{12}\subset W_{24}$; la otra construye el
retículo sobre el que se realizará el álgebra de operadores de vértice.

Sea $A_2=\mathbb Z a_1\oplus\mathbb Z a_2$ con matriz de Gram
\[
 \begin{pmatrix}2&-1\\-1&2\end{pmatrix},
 \qquad \rho=a_1+a_2,\qquad \rho^2=2.
\]
Representamos las tres clases de $A_2^*/A_2$ mediante
\[
 \varpi_0=0,\qquad
 \varpi_1=\frac{2a_1+a_2}{3},\qquad
 \varpi_2=\frac{a_1+2a_2}{3}.
\]
La aplicación $j\mapsto\varpi_j+A_2$ identifica aditivamente
$\mathbb F_3$ con ese discriminante. Para $c\in C_W$, escribimos
$\phi(c)=(\varpi_{c_1},\ldots,\varpi_{c_{12}})$ y definimos
\begin{equation}
 N(C_W)=\bigcup_{c\in C_W}\bigl(A_2^{12}+\phi(c)\bigr).
 \label{exc:niemeier-def}
\end{equation}

\begin{theorem}[Pegado ternario]
\label{exc:niemeier-teorema}
El objeto \eqref{exc:niemeier-def} es un retículo positivo, par y
unimodular de rango $24$. Su sistema de raíces es exactamente
$A_2^{12}$; por tanto, se reconoce como el retículo de Niemeier con
ese sistema de raíces.
\end{theorem}

\begin{proof}
La linealidad del código hace que la unión de clases sea estable por
suma. La autoortogonalidad anula el emparejamiento discriminante:
éste es un múltiplo de $c\cdot d/3$ módulo $\mathbb Z$. Así, el
retículo es integral. La norma de una clase de pegado es
$2\operatorname{wt}(c)/3$ módulo $2\mathbb Z$. Los pesos de
\eqref{exc:enumerador} son múltiplos de tres, de modo que es par.

El índice de $A_2^{12}$ en $N(C_W)$ es $|C_W|=3^6$. En consecuencia,
\[
 \det N(C_W)=\frac{3^{12}}{(3^6)^2}=1.
\]
La positividad procede de la suma ortogonal de las doce métricas $A_2$.
En una clase no nula de $A_2^*/A_2$ la norma mínima es $2/3$;
toda palabra no nula tiene al menos seis posiciones no nulas.
Por tanto, un vector de una clase de pegado no trivial tiene norma
al menos $6\cdot2/3=4$. No aparecen raíces de norma dos fuera de
$A_2^{12}$. Dentro de éste, las raíces son precisamente las de sus
doce componentes.
\end{proof}

\subsection{Selección incidencial del vecino unimodular sin raíces}
\label{exc:leech}

\input{sections/incidencia_reticulo.tex}

La etiqueta $o_*=1$ de la bandera se realiza ahora como una de las
doce componentes $A_2$. Esta operación declara la métrica anterior y
la cámara radial positiva
\[
 v(a)=\sum_{i=1}^{12}a_i\rho_i,
 \qquad a_i=1+3b_i,\quad b_i\in\mathbb Z_{\geq0}.
\]
El requisito de paridad del vecino de índice tres es
$v(a)^2\equiv0\pmod{18}$. Puesto que
\[
 v(a)^2=2\sum_i(1+3b_i)^2,
\]
esa congruencia equivale a $\sum_i b_i\equiv1\pmod3$. El menor
valor de la norma en la cámara se obtiene con un único $b_i=1$ y
todos los demás nulos; es $54$. El origen marcado selecciona
\begin{equation}
 v_\alpha=4\rho_1+\rho_2+\cdots+\rho_{12},
 \qquad v_\alpha^2=54.
 \label{exc:vector}
\end{equation}
La palabra ``menor'' se refiere a esta cámara positiva. No afirma
minimalidad en toda la clase residual: el vector
$(a_1-2a_2,\rho,\ldots,\rho)$ representa la misma clase módulo tres
y tiene norma $14+11\cdot2=36$. La restricción de cámara es parte
de los datos de la construcción, no una conclusión omitible.

Sea $N=N(C_W)$. Definimos
\begin{equation}
 N_{\alpha,0}=\{x\in N:\langle x,v_\alpha\rangle\equiv0\pmod3\},
 \qquad
 N_\alpha=N_{\alpha,0}+\mathbb Z\frac{v_\alpha}{3}.
 \label{exc:vecino}
\end{equation}

\begin{theorem}[Vecino de Leech determinado por la bandera marcada]
\label{exc:teorema-leech}
Con la métrica y la cámara declaradas, $N_\alpha$ es un retículo
positivo, par y unimodular de rango $24$, sin vectores de norma dos.
Por el teorema de clasificación correspondiente, es isométrico al
retículo de Leech $\Lambda_{24}$.
\end{theorem}

\begin{proof}
El carácter $x\mapsto\langle x,v_\alpha\rangle\bmod3$ no es nulo:
una raíz simple de una componente se empareja con $a_i\rho_i$ en
$a_i\equiv1\pmod3$. Por ello $[N:N_{\alpha,0}]=3$ y
$\det N_{\alpha,0}=9$. Además,
$v_\alpha/3\in N_{\alpha,0}^*$ y $v_\alpha\in N_{\alpha,0}$.
Si $v_\alpha/3$ perteneciera a $N$, la integridad de $N$ haría que
todos los emparejamientos con $v_\alpha$ fueran divisibles por tres,
contradiciendo el carácter no nulo. La clase de $v_\alpha/3$ tiene
por tanto orden tres. Se deduce
\[
 [N_\alpha:N_{\alpha,0}]=3,
 \qquad \det N_\alpha=9/3^2=1.
\]
La extensión es integral, y para $x\in N_{\alpha,0}$,
\[
 \left(x+m\frac{v_\alpha}{3}\right)^2
   =x^2+2m\frac{\langle x,v_\alpha\rangle}{3}+6m^2
   \in2\mathbb Z.
\]
Luego es par.

Queda excluir las raíces, que es el paso decisivo. Las raíces de $N$
pertenecen a $A_2^{12}$. Sus emparejamientos con $v_\alpha$ son
congruentes con $\pm1$ o $\pm2$ módulo tres; ninguna pertenece a
$N_{\alpha,0}$. Para las otras dos clases del vecino, cada componente
pertenece a alguna de las clases locales
\[
 A_2+\varpi_j\pm\rho/3,\qquad j=0,1,2.
\]
Como $4\rho/3\equiv\rho/3\pmod{A_2}$, esta descripción vale
también en la componente distinguida. Escribiendo un vector local
como $(u a_1+v a_2)/3$, su norma es
$2(u^2-uv+v^2)/9$. Las seis clases anteriores tienen residuos no
nulos y mínimo $2/9$: los representantes de mínima forma cuadrática
son $(\pm1,0),(0,\pm1),(1,1),(-1,-1)$ en los residuos apropiados.
Un vector de una clase nueva tiene, por tanto, norma al menos
$12\cdot2/9=8/3>2$. No hay raíces en ninguna de las tres clases.
La clasificación de los retículos pares unimodulares positivos de
rango veinticuatro identifica el único caso sin raíces con
$\Lambda_{24}$.
\end{proof}

El argumento no reconoce Leech por una coincidencia de dimensión.
Construye el pegado, prueba integridad, paridad y determinante, y
elimina las raíces mediante una cota local. Sólo entonces aplica el
teorema de clasificación. La dependencia de $\alpha$ se conserva en
la bandera que selecciona $o_*$ y en el vector \eqref{exc:vector};
no se convierte el valor escalar de la constante en una longitud del
retículo.

\subsection{El álgebra de operadores de vértice y el Monstruo}
\label{exc:voa}

Sea $\Lambda=N_\alpha$, ya reconocido como Leech. El paso siguiente
utiliza su retículo integral, no una expansión decimal. La construcción
de álgebra de operadores de vértice de retículo se escribe
\[
 V_\Lambda=M(1)\otimes\mathbb C_\varepsilon[\Lambda].
\]
Aquí $M(1)$ es el espacio de osciladores de Heisenberg en rango
veinticuatro; $\mathbb C_\varepsilon[\Lambda]$ es el álgebra de grupo
torcida por el cociclo de la extensión central del retículo. El
cociclo y el levantamiento de las isometrías forman parte de la
construcción: una suma vectorial desnuda no determina los productos
de vértice.

La negación $\lambda\mapsto-\lambda$ admite el levantamiento estándar
$\theta$. Con su módulo torcido y la elección de paridad compatible,
la construcción de Frenkel--Lepowsky--Meurman produce
\begin{equation}
 V^\natural_{\rm HMT}
   =V_\Lambda^+\oplus(V_\Lambda^T)^+.
 \label{exc:flm}
\end{equation}
El signo $+$ denota la parte par de la acción levantada. La expresión
incluye el producto de vértice del orbifold; no se limita a una suma
de espacios graduados.

\begin{theorem}[Prolongación hasta el módulo Moonshine]
\label{exc:moonshine}
El álgebra \eqref{exc:flm}, construida sobre el vecino
\eqref{exc:vecino}, es una realización del álgebra Moonshine
$V^\natural$. Tiene carga central $24$, componente de peso uno nula
y grupo de automorfismos isomorfo al Monstruo $\mathbb M$. Su carácter
graduado es
\begin{equation}
 \operatorname{Tr}_{V^\natural_{\rm HMT}}q^{L_0-1}
   =J(\tau)=j(\tau)-744
   =q^{-1}+196884q+21493760q^2+\cdots,
 \qquad q=e^{2\pi i\tau}.
 \label{exc:j}
\end{equation}
\end{theorem}

\begin{proof}
El teorema \ref{exc:teorema-leech} verifica las hipótesis reticulares:
positividad, paridad, unimodularidad, rango veinticuatro y ausencia de
raíces. Se aplica entonces la construcción y el teorema de
Frenkel--Lepowsky--Meurman para el retículo de Leech. Éstos identifican
el orbifold con $V^\natural$, establecen su carácter y su grupo de
automorfismos. La composición de la bandera, el vecino y el functor
de construcción de VOA da la realización aquí indicada. No se deduce
el grupo a partir del coeficiente $196884$.
\end{proof}

El teorema utilizado es un reconocimiento clásico posterior con
hipótesis verificadas, no una nueva demostración en este artículo de
la construcción FLM \cite{flm}.
Además, la notación $q=e^{2\pi i\tau}$ es la variable convencional
de Fourier empleada para reconocer el carácter. No introduce $\pi$
ni $e$ como entradas del generador coinductivo que ya los ha determinado.

Puede verse una parte del mecanismo de traza antes del orbifold.
Sólo el vector reticular cero queda fijo por la negación; cada
oscilador contribuye con signo alterno. De ahí
\[
 \operatorname{Tr}_{V_\Lambda}\theta q^{L_0-1}
   =t_2(\tau)
   :=q^{-1}\prod_{n\geq1}(1+q^n)^{-24}.
\]
Esta traza en el álgebra de retículo no debe confundirse con el
carácter sin inserción de $V^\natural$. Tampoco identifica por sí
sola una clase de conjugación del Monstruo: el espacio, la inserción
y el levantamiento deben especificarse antes de comparar funciones.

\input{sections/incidencia_automorfismos.tex}

\subsection{Trazas graduadas de Moonshine y representación del Monstruo}
\label{exc:alcance}

Para $g\in\operatorname{Aut}(V^\natural_{\rm HMT})$, la traza
con inserción es
\[
 T_g(\tau)=
 \operatorname{Tr}_{V^\natural_{\rm HMT}}gq^{L_0-1}.
\]
El teorema de Moonshine monstruoso identifica estas series como
Hauptmoduln de los grupos de género cero correspondientes
\cite[teorema~1.1]{borcherds}.
El carácter $J$, el grupo $\mathbb M$ y la colección $g\mapsto T_g$
son realizaciones del álgebra ya construida. No forman una sucesión
causal en la que primero se adivina $J$ y después se deduce el Monstruo.

La unidad tiene aquí una localización algebraica precisa. El subespacio
de peso cero es la línea del vacío, \(V^\natural_0=\CC\mathbf1\), y
el coeficiente de \(q^{-1}\) de \(J\) es por ello \(1\). En peso uno
la dimensión es cero, de donde procede la ausencia del término constante
en \(J=j-744\). Un morfismo de álgebras de vértices conserva el vector
vacío. Esta propiedad proporciona el sentido técnico de conservación
de la unidad en este nivel; la unidad de los observables cilíndricos
se conserva por la precomposición demostrada anteriormente. La relación
entre ambos niveles debe seguir los mapas de construcción con sus unidades.

En peso dos aparece la igualdad
\[
 196884=1+196883=196560+324=54(5\cdot729+1).
\]
La primera separación corresponde a la línea conforme trivial y al
componente irreducible de dimensión $196883$ de la representación
del Monstruo. Las otras expresiones tienen lecturas reticulares o
aritméticas. Su igualdad no identifica automáticamente sus sumandos
con regiones, palabras o rutas del TPK. Para hacerlo habría que
presentar los mapas entre esos espacios y probar que respetan las
acciones relevantes.


En particular, una isometría $N_\alpha\simeq\Lambda_{24}$ induce
una realización de \eqref{exc:flm}, y una elección de isomorfismo
con $V^\natural$ transporta su grupo de automorfismos. No se ha
construido con ello una acción del Monstruo sobre los dígitos de
$\pi$, sobre los bloques de $K$ o sobre todas las historias TPK.
Tal afirmación adicional requeriría una representación
$\rho_{\rm TPK}$ y un mapa $F$ con un cuadrado explícito
\[
 F\rho_{\rm TPK}(g)=\rho_{V^\natural}(g)F,
\]
incluyendo dominio, codominio e información conservada. La ausencia
de ese paso en la afirmación presente no retira el entrelazador
incidencial $U_W$, la bandera ni el vecino ya construidos; sólo
delimita una promoción distinta, de acción sobre el estado fuente.

La cadena establecida en esta sección puede resumirse, manteniendo
visibles sus datos de transporte, como
\begin{align*}
 &\text{APP--TRIT--TPK y estado aritmético con memoria de trayectoria}
   \longrightarrow (w_j,f_j;K,Z_0)\\
 &\longrightarrow (C_W,\mathcal H;B_K\subset H^-_\alpha,
                   P_\pi,H_e,H_\varphi,H_{\rm APP})\\
 &\longrightarrow \bigl(N(C_W),o_*,\text{métrica y cámara}\bigr)
   \longrightarrow N_\alpha
   \longrightarrow V^\natural_{\rm HMT}.
\end{align*}
En el nivel de incidencia, las estrellas generan $M_{12}$ y la
coordenatización residual realiza la estructura regional $4+1$ en $W_{24}$.
En el nivel reticular y de VOA, la construcción da Leech y el módulo
Moonshine. Son prolongaciones de una misma fuente marcada, con
objetos y mapas diferentes. La continuidad causal no obliga a
identificar sus codominios ni a borrar las elecciones explícitas
de marco, cámara y levantamiento.

El resultado de la sección es, por tanto, una derivación reunida:
las pruebas finitas y reticulares hacen efectivo el tránsito desde
la incidencia dodecafásica hasta las hipótesis de los teoremas de
reconocimiento. Este resultado se refiere a las construcciones y
a las hipótesis especificadas; no añade una validación experimental
a la comparación de constantes.

\endgroup

\section{Constantes de Planck y estructura determinantal de la acción}
\label{sec:action}
\subsection{Escala determinantal de acción y valoración decimal}

La escala absoluta y el corrector relativo se originan en objetos
distintos. El primero utiliza el bloque local del registro dodecafásico
de paso tres,
\begin{equation}
 H_4=\begin{pmatrix}
 1&1&1&1\\1&1&-1&-1\\1&-1&1&-1\\1&-1&-1&1
 \end{pmatrix},\qquad
 G_H=\mathbb Z^4/H_4\mathbb Z^4.
 \label{eq:el-h4}
\end{equation}
La operación utiliza un cociente local. Los tres bloques del registro
completo no se sustituyen por una tensorización que multiplicara de
nuevo el número de hojas de este lector.

\begin{lemma}[Número de hojas del cociente local]
La forma normal de Smith de \(H_4\) es
\(\operatorname{diag}(1,2,2,4)\). En consecuencia,
\[
 G_H\simeq\mathbb Z/2\mathbb Z\oplus\mathbb Z/2\mathbb Z
 \oplus\mathbb Z/4\mathbb Z,\qquad |G_H|=16.
\]
\end{lemma}
\begin{proof}
El máximo común divisor de las entradas es uno. Los menores de orden dos
son pares y hay uno de valor \(-2\), de modo que su máximo común divisor
es dos. Como \(H_4H_4^{\mathsf T}=4I\) y \(\det H_4=-16\), los
cofactores son \(\pm4\); el tercer divisor determinantal es cuatro.
El cuarto es \(16\). Los cocientes sucesivos de los divisores
\(1,2,4,16\) son \(1,2,2,4\). Su producto da el orden del cociente.
\end{proof}

\begin{definition}[Lector determinantal de acción]
La norma multiplicativa de la sección escalar \(\alpha\), constante
sobre las hojas de \(G_H\), y una aplicación de la autoescala
\(\varphi\) producen
\begin{equation}
 \Sigma_H=\varphi\,\mathrm N_{G_H}(\alpha),\qquad
 \mathrm N_{G_H}(\alpha)=\prod_{g\in G_H}\alpha=\alpha^{16}.
 \label{eq:el-lector-determinantal}
\end{equation}
La coordenatización decimal asociada determina
\(\operatorname{ord}_{10}(\Sigma_H)=\lfloor\log_{10}\Sigma_H\rfloor\).
\end{definition}

La potencia dieciséis procede así del número de hojas del cociente, una
vez fijado este lector local. El factor \(\varphi\) actúa una vez sobre
la base de la recta de acción. El cálculo de Smith, por sí solo, no
seleccionaría cualquier funcional posible sobre el cociente: la norma
multiplicativa y la autoescala forman parte de la operación declarada.

\begin{theorem}[Década de la sección de acción]
El lector \eqref{eq:el-lector-determinantal}, aplicado a las coordenadas
HMT ya generadas, satisface
\begin{equation}
 10^{-34}<\Sigma_H<10^{-33},\qquad
 \operatorname{ord}_{10}(\Sigma_H)=-34.
 \label{eq:el-decada}
\end{equation}
\end{theorem}
\begin{proof}
\label{md:accion:intervalos}Bastan los intervalos racionales, mucho más gruesos que los cilindros
disponibles,
\[
 \frac{1618}{1000}<\varphi<\frac{1619}{1000},\qquad
 \frac{7297}{10^6}<\alpha<\frac{7298}{10^6}.
\]
La función \((x,a)\mapsto xa^{16}\) es creciente en ambos argumentos
positivos. Las dos desigualdades enteras
\[
 1618\cdot7297^{16}>10^{65},\qquad
 1619\cdot7298^{16}<10^{66}
\]
implican
\[
 10^{-34}
 <\frac{1618}{1000}\left(\frac{7297}{10^6}\right)^{16}
 <\Sigma_H
 <\frac{1619}{1000}\left(\frac{7298}{10^6}\right)^{16}
 <10^{-33}.
\]
La definición del orden decimal concluye la prueba. No interviene una
unidad SI ni una masa de referencia en estas desigualdades.
\end{proof}

La sección normalizada comienza por
\(\Sigma_H=1.0462438381\ldots\times10^{-34}\). Su mantisa no se
identifica con la coordenada de acción anterior o posterior al retorno.
El lector determinantal fija el orden decimal; las dos coordenadas de
retorno se obtienen mediante los lectores siguientes.


\subsection{Secciones de acción y cociente adimensional de retorno}

Sea \(\mathcal L_S\) una recta orientada de acción y \(\mathcal U_S\)
su base positiva. El carácter de quinto orden y el término regional de
retorno tienen las expresiones
\begin{align}
 H_5={}&\frac12\sqrt{\frac{1000\alpha}{\varphi}}-\alpha
   +\frac9{16}\alpha^2-\frac59\alpha^3
   +\frac7{48}\alpha^4-\frac1{54}\alpha^5,
   \label{eq:el-h5}\\
 D_A={}&\exp\!\left(-\frac{100\pi\alpha}{9}\right),
   \label{eq:el-da}\\
 \mathcal C_\pi(\alpha)={}&169\alpha^6+
       \frac{D_A\alpha^7}{1-D_A\alpha},
   \label{eq:el-cpi}\\
 \eta_{\mathrm{ret}}={}&H_5-\frac{90}{\pi}\mathcal C_\pi(\alpha).
   \label{eq:el-eta}
\end{align}
Estas son lecturas posteriores del cierre dodecafásico y de su
incidencia regional. El entero \(169\) se conserva como el selector
regional ya construido, y los coeficientes de \(H_5\) como los de su
carácter local; no se reemplazan por una interpolación de la constante de
Planck. Las fórmulas explicitan exactamente qué lecturas se necesitan
para la etapa electrónica. No es necesario introducir aquí otros
funcionales de la teoría de masas.

Para \(0<\alpha<1\), se tiene \(0<D_A<1\), de modo que
\(D_A\alpha<1\). El término racional de \eqref{eq:el-cpi} conserva la
suma completa de la cola geométrica
\begin{equation}
 \frac{D_A\alpha^7}{1-D_A\alpha}
 =\sum_{n=0}^{\infty}D_A^{n+1}\alpha^{n+7}.
 \label{eq:el-cola-regional}
\end{equation}
No es una truncación que pueda modificarse silenciosamente al cambiar la
precisión. Tras el término de índice \(N\), la cola restante es
\(D_A^{N+2}\alpha^{N+8}/(1-D_A\alpha)\), lo que permite controlar la
evaluación sin acudir a una magnitud objetivo.

\begin{definition}[Secciones de acción y lector de retorno]
Las secciones anterior y posterior al retorno son
\begin{equation}
 \hbar_{\mathrm{pre}}=H_5\,10^{-34}\mathcal U_S,
 \qquad
 \hbar_{\mathrm{ret}}=\eta_{\mathrm{ret}}\,10^{-34}\mathcal U_S.
 \label{eq:el-secciones-accion}
\end{equation}
Siempre que \(\eta_{\mathrm{ret}}>0\), se definen
\begin{equation}
 \delta_{\mathrm{ret}}=H_5-\eta_{\mathrm{ret}},\qquad
 R_{\mathrm{act}}=\frac{\hbar_{\mathrm{pre}}}{\hbar_{\mathrm{ret}}}
 =\frac{H_5}{\eta_{\mathrm{ret}}}.
 \label{eq:el-ratio}
\end{equation}
\end{definition}

\begin{proposition}[Positividad e invariancia de la razón]\label{md:accion:positividad}
En los intervalos usados en la prueba de \eqref{eq:el-decada},
\(0<\eta_{\mathrm{ret}}<H_5\), y por tanto
\(R_{\mathrm{act}}>1\). La razón no cambia al sustituir la base
\(\mathcal U_S\) por cualquier otra base positiva común.
\end{proposition}
\begin{proof}
Los términos de \(\mathcal C_\pi\) son positivos. Para una cota
elemental basta usar \(\alpha<0.008\), \(\varphi<1.619\),
\(\alpha>0.007\), \(D_A<1\) y \(\pi>3\). La raíz en
\eqref{eq:el-h5} es mayor que uno, mientras la suma de los términos
negativos es menor que \(0.008001\); en particular, \(H_5>0.99\).
Asimismo,
\[
 0<\frac{90}{\pi}\mathcal C_\pi
 <30\left(169(0.008)^6+
           \frac{(0.008)^7}{1-0.008}\right)<10^{-6}.
\]
Así \(\eta_{\mathrm{ret}}>0.989999>0\) y
\(\eta_{\mathrm{ret}}<H_5\). Finalmente, el cambio de base divide
ambas coordenadas de acción por el mismo factor positivo, que cancela en
su cociente.
\end{proof}

La evaluación posterior proporciona
\begin{align*}
 H_5&=1.0545718176461544696\ldots,\\
 \eta_{\mathrm{ret}}&=1.0545718169150390611\ldots,\\
 R_{\mathrm{act}}&=1.0000000006932817630\ldots.
\end{align*}
La diferencia y el cociente son dos lecturas de las mismas secciones,
una aditiva y otra multiplicativa. No representan dos constantes de
Planck independientes. El paso de acción angular a acción por vuelta
completa se expresa, en cada sección, mediante
\(h=2\pi\hbar\). La posterior coordenatización
\(\mathcal U_S\mapsto\mathrm{J\,s}\) asigna una unidad, pero no
retroactúa sobre el orden decimal ya demostrado.

El retorno regional puede escribirse también en una coordenatización angular.
Esta reparametrización no es una dependencia adicional indispensable:
\eqref{eq:el-h5}--\eqref{eq:el-eta} suministran directamente
\(H_5\), \(\eta_{\mathrm{ret}}\) y \(R_{\mathrm{act}}\). Si se
introducen ángulo y contraángulo, ambos deben transportarse en la misma
coordenatización; no se mezclan grados y radianes dentro de una diferencia.

\subsubsection{Carácter de acción y procedencia de sus coeficientes}
\label{sec:revision-planck}

La constante de Planck reducida, \(\hbar\), es el objeto físico que
motiva la realización de la sección de acción angular. Su posición en
la dependencia electrónica debe especificarse en dos niveles. La norma
determinantal \(\Sigma_H\) determina la década de la escala. El carácter
\(H_5\) y el retorno regional determinan después las coordenadas de las
secciones de acción. Esta separación permite exponer el resultado
necesario para el electrón sin anticipar los desarrollos completos de
otros sectores de la mecánica dimensional.

La procedencia del carácter de quinto orden puede explicitarse mediante
invariantes ya fijados en el soporte. En la coordenatización central se considera
\[
 B_c=\begin{pmatrix}7&2\\2&7\end{pmatrix},
 \qquad \lambda_+=9,\quad\lambda_-=5.
\]
El calendario tiene longitud \(12\), la órbita del paso \(3\) tiene
longitud \(4\), y el cociente censal pertinente es
\(104976/1944=54\). Los coeficientes del carácter satisfacen
\begin{equation}
 \frac9{16}=\frac{\lambda_+}{4^2},\quad
 \frac59=\frac{\lambda_-}{\lambda_+},\quad
 \frac7{48}=\frac{\operatorname{tr}(B_c)/2}{4\cdot12},\quad
 \frac1{54}=\frac{1944}{104976}.
 \label{eq:revision-h5-coeficientes}
\end{equation}
La asignación de estos invariantes a los grados \(2,3,4,5\), con la
alternancia de signos de \eqref{eq:el-h5}, pertenece a la definición del
carácter analítico utilizado. Explicitarla es esencial: el conjunto de
cardinales, considerado sin el carácter y sin el orden de sus grados,
admitiría otras expresiones. El contenido constructivo reside en la
composición del soporte, la orientación y la regla analítica declarada.

\subsubsection{Continuación regional y ecuación de renovación}

El selector \(169\) determina el impulso de grado seis. La prolongación
posterior se organiza por la transferencia \(D_A\) y una normalización
unitaria de la línea determinantal. Estas reglas admiten una formulación
en series formales antes de sustituir \(z=\alpha\), con \(D_A\) ya
construido. Escribamos
\[
 \mathscr C(z)=169z^6+R(z),\qquad R(z)\in z^7\mathbb R[[z]].
\]
La concatenación de una prolongación aumenta el grado en una unidad y
multiplica su peso por \(D_A\). La primera prolongación estricta tiene
peso \(D_A\). Por tanto, su ecuación de renovación es
\begin{equation}
 R(z)=D_Az^7+D_AzR(z).
 \label{eq:revision-renovacion}
\end{equation}

\begin{proposition}[Unicidad formal de la continuación normalizada]
La ecuación \eqref{eq:revision-renovacion} determina una única serie y
\[
 \mathscr C(z)=169z^6+\frac{D_Az^7}{1-D_Az}
             =169z^6+\sum_{n=7}^{\infty}D_A^{n-6}z^n.
\]
La realización es analítica en \(|D_Az|<1\).
\end{proposition}
\begin{proof}
El elemento \(1-D_Az\) tiene término independiente uno y es invertible
en el anillo de series formales. Resolver la ecuación produce la serie
exhibida. Equivalentemente, el coeficiente de grado siete es \(D_A\)
y cada coeficiente posterior es \(D_A\) veces el anterior. La serie
geométrica converge absolutamente en el disco indicado.
\end{proof}

La normalización de la primera prolongación es un dato operativo
indispensable. Si se conservaran únicamente el impulso de grado seis y
la razón de concatenación, las series
\[
 169z^6+\lambda\frac{D_Az^7}{1-D_Az}
\]
seguirían compartiendo ambos datos para cualquier \(\lambda\).
La regla unitaria fija \(\lambda=1\) antes de efectuar la evaluación.
Así queda localizada la contribución de cada regla a la unicidad del
lector, y la expresión racional conserva la continuación completa.

\subsubsection{Constante de Planck reducida y retorno de la sección de acción}

En la coordenatización dimensional de \eqref{eq:el-secciones-accion}, la sección
anterior al retorno es la construcción del modelo destinada al contraste
con la constante de Planck reducida:
\[
 \hbar_{\mathrm{pre}}=H_5\,10^{-34}\mathcal U_S.
\]
La sección posterior conserva, además, la compensación regional
\(90\mathscr C(\alpha)/\pi\). Ambas pertenecen a la misma recta de
acción y su cociente \(R_{\mathrm{act}}\) mide el efecto multiplicativo
del retorno. El paso a \(h\) mediante \(2\pi\) utiliza la relación
entre la parametrización angular y una vuelta completa.

Esta construcción interviene efectivamente en
\(R_{\mathrm{act}}^{80-54\alpha+6\Delta_4}\). El electrón conserva
su coordenada central y su trayectoria; la escala cronológica no
reemplaza esa estructura por un cuanto indiferenciado. Tampoco se
multiplica de nuevo por \(10^{-34}\) una masa ya expresada en MeV:
la década pertenece a la sección de acción, cancela en su cociente
adimensional y reaparece únicamente donde actúa la realización
dimensional correspondiente.

El contraste numérico de \(H_5\), de la sección retornada y de
\(h/(2\pi)\) se presenta en la sección~\ref{sec:revision-planck-contraste}.
Así se distingue la denominación del objeto físico, la ecuación que el
modelo le asigna y la precisión con la que dicha ecuación lo reproduce.

\subsection{Contraste de la constante de Planck reducida}
\label{sec:revision-planck-contraste}

El Sistema Internacional vigente desde el 20 de mayo de 2019 fija
exactamente \(h=6.62607015\times10^{-34}\,\mathrm{J\,s}\)
\cite{bipm2018}. Por tanto, su constante reducida es
\[
 \hbar_{\mathrm{SI}}=\frac{6.62607015}{2\pi}\,10^{-34}\,
 \mathrm{J\,s}
 =1.0545718176461563912\ldots\times10^{-34}\,\mathrm{J\,s}.
\]
Los puntos suspensivos corresponden a la expansión de una expresión
exacta, no a una incertidumbre experimental de \(h\) en el SI actual.
Este valor de referencia se utiliza aquí después de la evaluación de
los lectores de HMT; su función es comparativa.

Las dos coordenadas del retorno del artículo son
\begin{align*}
 H_5&=1.0545718176461544696\ldots,\\
 \eta_{\mathrm{ret}}&=1.0545718169150390611\ldots.
\end{align*}
En la identificación de \(\mathcal U_S\) con \(\mathrm{J\,s}\), la
diferencia relativa de la sección anterior respecto de
\(\hbar_{\mathrm{SI}}\) es aproximadamente \(-1.82\times10^{-15}\).
La de la sección posterior es aproximadamente \(-6.93\times10^{-10}\).
La segunda diferencia contiene el efecto del retorno que utiliza el
corrector electrónico; no equivale al error de evaluación de la
primera sección.

La precisión de estas comparaciones permite reconocer la proximidad
cuantitativa de la construcción de \(H_5\) y, simultáneamente, distinguir
proximidad de identidad exacta. En el presente manuscrito las expresiones
exhibidas no proporcionan una igualdad exacta con \(h/(2\pi)\).
El resultado demostrado sobre \(\Sigma_H\) es su década \(-34\);
el resultado analítico sobre \(H_5\) es la evaluación del carácter
declarado. La identificación física se examina en esta coordenatización y con las
diferencias explícitas que acabamos de dar. El valor SI no interviene
en la definición de la norma determinantal ni en la recurrencia regional.

\subsection{Coordenadas angulares, normalización de la vuelta y unidades de acción}
\label{sec:ii-angulos}

La magnitud de acción y la coordenada angular pertenecen a espacios
distintos. La sección \(\hbar_{\rm ret}\) ocupa la recta de acción
\(\mathcal L_S\); \(\eta_{\rm ret}\) es su coeficiente en la base
declarada. La representación angular actúa sobre la pareja adimensional
\((\alpha,\eta_{\rm ret})\), después de su construcción:
\begin{equation}
 \mathcal P(\alpha,\eta_{\rm ret})
 =\bigl([1000\alpha]_{360},[2(\eta_{\rm ret}+\alpha)]_{360}\bigr).
 \label{eq:ii-par-angular}
\end{equation}
La completación nonádica determina el factor 1000; la clausura
bidireccional determina el factor dos. La coordenatización circular de 360
unidades hace visibles las particiones
\(12\cdot30=4\cdot90=3\cdot120\). Sobre la rama positiva del
artículo, los representantes son
\[
 A_{\deg}=1000\alpha,\qquad
 C^*_{\deg}=2(\eta_{\rm ret}+\alpha).
\]

\begin{proposition}[Desplazamiento de la representación en base mil]
\label{prop:ii-desplazamiento-angular}
Sea \(\alpha=\sum_{n\geq1}b_n1000^{-n}\), con
\(b_n\in\{0,\ldots,999\}\), la expansión compatible ya generada.
La coordenada levantada \(A_{\deg}=1000\alpha\) satisface
\[
 A_{\deg}=b_1+\sum_{n\geq1}b_{n+1}1000^{-n}.
\]
La representación conserva la sucesión de bloques a partir del segundo;
el bloque inicial pasa a la parte entera.
\end{proposition}
\begin{proof}
La serie converge absolutamente. Multiplicar cada término por 1000
y separar el término inicial proporciona la identidad. La operación
actúa sobre la expansión construida, con sus bloques y su memoria;
la representación circular posterior toma su clase módulo 360.
\end{proof}

\begin{definition}[Coordenatizaciones de la coordenada angular]
La coordenatización en grados expresa una vuelta completa mediante 360 unidades.
La coordenatización en radianes la expresa mediante \(2\pi\). El cambio de
coordenadas y sus inversas son
\[
 x=\frac{\pi}{180}A_{\deg},\qquad
 y=\frac{\pi}{180}C^*_{\deg},\qquad
 A_{\deg}=\frac{180}{\pi}x,\qquad
 C^*_{\deg}=\frac{180}{\pi}y.
\]
Sobre un recorrido con número de vueltas conservado se usan las
coordenadas levantadas; el cociente circular expresa su clase módulo
la vuelta correspondiente.
\end{definition}

La definición caracteriza grados y radianes por su relación exacta
con la vuelta. El registro retiene además orientación, vueltas y
memoria de transporte. La representación de su clase angular conserva
la fase, mientras el levantamiento distingue recorridos con la misma
fase final. La lectura física de acción necesita precisar cuál de
estas coordenadas interviene en la integral.

\begin{proposition}[Covariancia del lector de acción]
Para \(h_a=2\pi\hbar_a\), la acción de una coordenada levantada
satisface
\begin{equation}
 \mathcal S_a(\theta)
 =\hbar_a\theta_{\rm rad}
 =h_a\frac{\theta_{\deg}}{360}.
 \label{eq:ii-accion-cartas}
\end{equation}
Una vuelta positiva expresa \(h_a\).
\end{proposition}
\begin{proof}
Sustituir \(\theta_{\rm rad}=\pi\theta_{\deg}/180\) da
\(\hbar_a\pi\theta_{\deg}/180
=(2\pi\hbar_a)\theta_{\deg}/360\). En una vuelta
\(\theta_{\deg}=360\), de donde resulta \(h_a\).
\end{proof}

La periodicidad de una representación se declara aparte. Si la
representación transportada satisface \(U(2\pi)=-I\), su restitución
operatoria ocurre en \(4\pi\). Esta propiedad del levantamiento
coexiste con la normalización por vuelta de
\eqref{eq:ii-accion-cartas}. Cada una conserva su objeto: magnitud
integrada, fase circular y estado representado.

\subsection{Covariancia de los canales y del funcional de incidencia}

Los canales se escriben de manera equivalente como
\begin{equation}
 q_\pm=\exp[-(x\pm y)]
 =\exp\!\left[-\frac{\pi}{180}(A_{\deg}\pm C^*_{\deg})\right].
\label{eq:ii-canales}
\end{equation}
\begin{proposition}[Recuperación de las coordenadas par e impar]
\label{prop:ii-inversion-canales-angulares}
En la coordenatización real levantada, los canales positivos determinan
unívocamente las dos coordenadas:
\[
 A_{\deg}=-\frac{90}{\pi}\log(q_+q_-),\qquad
 C^*_{\deg}=\frac{90}{\pi}\log\frac{q_-}{q_+}.
\]
El intercambio \((q_+,q_-)\mapsto(q_-,q_+)\) conserva
\(A_{\deg}\) y revierte \(C^*_{\deg}\).
\end{proposition}
\begin{proof}
La suma y la diferencia de los logaritmos de
\eqref{eq:ii-canales} son, respectivamente,
\(-\pi A_{\deg}/90\) y \(-\pi C^*_{\deg}/90\).
El logaritmo real es unívoco sobre los canales positivos.
El producto es simétrico y el cociente se invierte al intercambiarlos.
\end{proof}
En el registro de doce posiciones, la involución antipodal
\(r\mapsto r+6\pmod {12}\) forma seis parejas. La distribución
uniforme de la coordenada impar sobre estas parejas es
\(C^*_{\deg}/6\): sumar sus seis contribuciones recupera
\(C^*_{\deg}\). Esta coordenada sectorial se distingue de
la constante de Planck \(h\), que pertenece a la recta de acción.

La conversión se aplica una sola vez a cada argumento. La magnitud
\(\alpha\) sigue siendo adimensional y la magnitud \(\hbar\)
sigue perteneciendo a la recta de acción; sus coordenadas previas
intervienen en el mapa \(\mathcal P\).

El funcional de Barbero--Immirzi contiene, además de los canales, una
componente bilineal en los representantes angulares. Su transporte es
\begin{equation}
 \frac{\pi}{180}A_{\deg}C^*_{\deg}
 =\frac{180}{\pi}xy.
 \label{eq:ii-barbero-cartas}
\end{equation}
La identidad se obtiene sustituyendo ambos representantes en grados
por sus expresiones en radianes. Por tanto, la fórmula completa en
coordenadas analíticas es
\[
 \gamma_{\HMT}
 =\frac{180}{\pi}xy
 +12\bigl[S_{90}(e^{-x+y})-S_{90}(e^{-x-y})\bigr]
 -\bigl[S_{120}(e^{-x+y})-S_{120}(e^{-x-y})\bigr].
\]
La equivalencia de las dos coordenatizaciones transporta la componente bilineal
y las funciones exponenciales simultáneamente. Esta transformación
explicita el papel preciso del grado en la construcción y permite
reproducir la evaluación sin ambigüedad de unidades angulares.

\section{Elipse de acción y reciprocidad modular}
\label{sec:ellipse}
Fijadas las salidas $A,C^*,\hbar>0$ en la coordenatización
$|C^*/A|<1$, definimos
\[
\eta_{\rm el}=\operatorname{artanh}(C^*/A),\quad
a_S=\sqrt{2\hbar}\,e^{\eta_{\rm el}/2},\quad
b_S=\sqrt{2\hbar}\,e^{-\eta_{\rm el}/2}.
\]
Las coordenadas canónicas equilibradas tienen unidad de raíz de acción;
no se suman aquí posición y momento sin una coordenatización dimensional.
Se cumplen
\[
a_Sb_S=2\hbar,\qquad
\tau_{\rm el}=i\,\frac{b_S}{a_S},\qquad
R_{\rm el}=\sqrt{\frac{b_S}{a_S}}
=\left(\frac{A-C^*}{A+C^*}\right)^{1/4}.
\]
La raíz positiva determina la composición $\tau_{\rm el}=iR_{\rm el}^2$,
y permite recuperar la razón angular:
\[
\frac{C^*}{A}=\frac{1-R_{\rm el}^4}{1+R_{\rm el}^4}.
\]
La inversión de forma envía $R_{\rm el}$ a $R_{\rm el}^{-1}$ y
$\tau_{\rm el}$ a $-1/\tau_{\rm el}$.

\begin{proposition}[Energía coincidente, acción distinguible]
Sea
\[
D(\eta)=\begin{pmatrix}e^\eta&0\\0&e^{-\eta}\end{pmatrix},
\quad
S_2=\begin{pmatrix}0&1\\1&0\end{pmatrix},
\quad
J_2=\begin{pmatrix}0&-1\\1&0\end{pmatrix}.
\]
Ambos mapas satisfacen $T^{\mathsf T}D(-\eta)T=D(\eta)$.
Sin embargo, para $\omega=\dd Q\wedge\dd P$,
\[
S_2^*\omega=-\omega,\qquad J_2^*\omega=\omega.
\]
Si $\mathcal I(\gamma)=\oint_\gamma P\,\dd Q$ sobre una curva cerrada,
\[
\mathcal I(S_2\gamma)=-\mathcal I(\gamma),\qquad
\mathcal I(J_2\gamma)=\mathcal I(\gamma).
\]
\end{proposition}
\begin{proof}
La conjugación intercambia las entradas diagonales de $D$.
Los determinantes de $S_2,J_2$ son $-1,1$, respectivamente, lo que
prueba la transformación de $\omega$. Para $\lambda=P\,\dd Q$,
\[
S_2^*\lambda=\dd(PQ)-\lambda,\qquad
J_2^*\lambda=\lambda-\dd(PQ).
\]
Integrar la diferencial exacta sobre la curva cerrada completa la prueba.
\end{proof}

La energía cuadrática y el módulo no bastan para distinguir ambos
transportes. La orientación de la acción sí los distingue, incluso
en la forma circular. El marcado geométrico es así parte de la lectura
recuperable, no un adorno añadido al valor.

% Adición separada: no sustituye 07_elipse.tex.
% Procedencia: integral 2249, 52.4.1--2, 52.6.2--3, 53.28 y 90.4.6.
% Recuperación de resultados y explicitación geométrica; certificado local
% verificar_elipse_II.py. No usa un radio físico como entrada generadora.

\subsection{Semiejes, focos y radios interiores}
\label{sec:ii-elipse-geometria}

La figura se obtiene de las salidas angulares y de acción de las secciones
precedentes. APP conserva las hojas y sus residuos y cocientes; TRIT fija
la orientación; TPK transporta frontera, acarreo y memoria hasta las
representaciones del mismo estado enriquecido y de la estructura discreta
conjunta del continuo. Aquí se recibe el par \((A,C^*)\) y la sección
positiva \(\hbar\) después de esa construcción. La geometría explicita
qué información conservan sus lectores, sin usar una longitud observada
para seleccionar el estado fuente.

Sea \(\xi=C^*/A\in(-1,1)\), con \(A\ne0\) y ambos ángulos en la
misma unidad. Escribimos \(\eta=\operatorname{artanh}\xi\) y
\[
 a_S=\sqrt{2\hbar}\,e^{\eta/2},\qquad
 b_S=\sqrt{2\hbar}\,e^{-\eta/2},\qquad
 \gamma(\theta)=(a_S\cos\theta,b_S\sin\theta).
\]
Los semiejes están etiquetados por sus coordenadas: si \(\eta<0\),
el eje vertical es el mayor. Ambos pertenecen a la recta de raíz de
acción de la coordenatización equilibrada, no a la recta de longitudes espaciales.

\begin{proposition}[Geometría focal e interior radial]
\label{prop:ii-elipse-focos}
Sean \(a=\max(a_S,b_S)\), \(b=\min(a_S,b_S)\). La semidistancia
focal y la excentricidad son
\begin{equation}
 f_S=\sqrt{a^2-b^2}
 =\sqrt{4\hbar\sinh|\eta|},\qquad
 e_f=\frac{f_S}{a}=\sqrt{1-e^{-2|\eta|}}.
 \label{eq:ii-elipse-focal}
\end{equation}
Para \(\eta>0\), los focos son \((\pm f_S,0)\); para
\(\eta<0\), son \((0,\pm f_S)\). En \(\eta=0\) coinciden con
el centro. El radio de frontera en dirección polar \(\phi\) es
\begin{equation}
 \rho_{\partial}(\phi)
 =\left(\frac{\cos^2\phi}{a_S^2}
        +\frac{\sin^2\phi}{b_S^2}\right)^{-1/2}.
 \label{eq:ii-elipse-radio-polar}
\end{equation}
El segmento interior en esa dirección consta de
\(r(\cos\phi,\sin\phi)\), con \(0\le r<\rho_{\partial}(\phi)\).
En particular, \(b\le\rho_{\partial}(\phi)\le a\).
\end{proposition}

\begin{proof}
Los cuadrados de los semiejes mayor y menor son
\(2\hbar e^{|\eta|}\) y \(2\hbar e^{-|\eta|}\). Su diferencia
da \eqref{eq:ii-elipse-focal}. Para verificar el foco, supongamos
\(\eta>0\). En \(X=\gamma(\theta)\), las distancias a
\(F_+=(f_S,0)\) y \(F_-=(-f_S,0)\) satisfacen
\[
 |X-F_+|^2=(a_S-f_S\cos\theta)^2,\qquad
 |X-F_-|^2=(a_S+f_S\cos\theta)^2.
\]
Como \(a_S>f_S\), sus raíces positivas suman \(2a_S\).
Para \(\eta<0\), intercambiar las coordenadas da la misma prueba.
En el caso circular, \(f_S=0\). Finalmente, sustituir
\((Q,P)=r(\cos\phi,\sin\phi)\) en
\(Q^2/a_S^2+P^2/b_S^2=1\) da
\eqref{eq:ii-elipse-radio-polar}. El coeficiente de \(r^2\) está
entre \(a^{-2}\) y \(b^{-2}\), lo que prueba las cotas y la
caracterización del interior.
\end{proof}

\subsection{Fase paramétrica, ángulo polar y área de acción}
\label{sec:ii-elipse-fase-area}

\begin{proposition}[Ley areal y dos coordenadas angulares]
\label{prop:ii-elipse-area}
La fase \(\theta\) de \(\gamma\) y su ángulo polar continuo
\(\phi\), levantados desde cero, satisfacen
\begin{equation}
 \phi(\theta)=\arg(a_S\cos\theta+i b_S\sin\theta),\qquad
 \frac{d\phi}{d\theta}
 =\frac{a_Sb_S}{a_S^2\cos^2\theta+b_S^2\sin^2\theta}>0.
 \label{eq:ii-elipse-fases}
\end{equation}
En una coordenatización donde ambas tangentes están definidas,
\(\tan\phi=(b_S/a_S)\tan\theta\). El área orientada barrida es
\begin{equation}
 \mathcal A(\theta)=\frac12\int_0^\theta(QP'-PQ')\,dt
 =\hbar\theta,\qquad
 \mathcal A(1)=\hbar,\qquad \mathcal A(2\pi)=h.
 \label{eq:ii-elipse-area-fase}
\end{equation}
\end{proposition}

\begin{proof}
La derivada del argumento de un vector no nulo es
\((QP'-PQ')/(Q^2+P^2)\). Aquí el numerador es
\(a_Sb_S=2\hbar\), lo que prueba la derivada y la monotonía.
El cociente \(P/Q\) da la relación de tangentes; el levantamiento
continuo fija los cuadrantes y el número de vueltas. La integral areal
vale \(a_Sb_S\theta/2=\hbar\theta\), y
\(h=2\pi\hbar\) da la última identidad.
\end{proof}

La fase paramétrica no coincide en general con el ángulo polar. Un radián
de fase barre \(\hbar\), pero no se identifica por ello con un arco
circular ni con un sector de ángulo polar unitario. En coordenadas polares,
la misma ley se escribe
\(d\mathcal A=\rho_{\partial}(\phi)^2d\phi/2\).
Con la orientación de \(\gamma\) utilizada aquí,
\(\oint P\,dQ=-h\): el área positiva es \(h\), mientras la
integral de esa uno-forma conserva su signo. Esto concuerda con la
distinción entre los transportes simpléctico y antisimpléctico de la
sección anterior.

\begin{figure}[htbp]
 \centering
 % Geometría calculada, no coordenadas focales elegidas a ojo.
% Caso algebraico de ilustración xi=1/3; no evaluación de las constantes HMT.
\begingroup
\pgfmathsetmacro{\elXi}{1/3}
\pgfmathsetmacro{\elEta}{0.5*ln((1+\elXi)/(1-\elXi))}
\pgfmathsetmacro{\elA}{exp(\elEta/2)}
\pgfmathsetmacro{\elB}{exp(-\elEta/2)}
\pgfmathsetmacro{\elF}{sqrt(\elA*\elA-\elB*\elB)}
\pgfmathsetmacro{\elTheta}{30}
\pgfmathsetmacro{\elX}{\elA*cos(\elTheta)}
\pgfmathsetmacro{\elY}{\elB*sin(\elTheta)}
\pgfmathsetmacro{\elPolar}{atan(\elY/\elX)}
\begin{tikzpicture}[x=2.18cm,y=2.18cm,font=\small,
 every node/.style={inner sep=2pt},>=Stealth]
 \begin{scope}
  \node[font=\bfseries] at (0,1.38) {Forma y orientación};
  \fill[accent!12] (0,0)--(\elA,0)
    plot[domain=0:30,samples=50] ({\elA*cos(\x)},{\elB*sin(\x)})--cycle;
  \draw[ink,thick] (0,0) ellipse[x radius=\elA,y radius=\elB];
  \draw[->,black!55] (-1.42,0)--(1.48,0) node[right]{\(u\)};
  \draw[->,black!55] (0,-1.05)--(0,1.08) node[above]{\(v\)};
  \draw[accent,thick,-{Stealth}] (\elA,0)
    plot[domain=0:30,samples=50] ({\elA*cos(\x)},{\elB*sin(\x)});
  \draw[accent,thick] (0,0)--(\elX,\elY);
  \node[below left,font=\scriptsize] at (0,0) {\(O\)};
  \draw[accent] (.35,0) arc[start angle=0,end angle=\elPolar,radius=.35];
  \node[accent] at (.47,.065) {\(\phi\)};
  \fill[accent] (\elX,\elY) circle[radius=1.4pt];
  \node[above right] at (\elX,\elY) {\(X(\theta)\)};
  \node[accent,align=center,font=\scriptsize] at (1.26,.18) {\(\theta=\pi/6\)};
  \draw[dashed,black!45] (-\elF,0)--(\elX,\elY);
  \draw[dashed,black!45] (\elF,0)--(\elX,\elY);
  \fill[ink] (-\elF,0) circle[radius=1.5pt] node[below]{\(F_-\)};
  \fill[ink] (\elF,0) circle[radius=1.5pt] node[below]{\(F_+\)};
  \draw[<->,black!60] (-\elA,-1.04)--(0,-1.04)
    node[midway,below,font=\scriptsize]{\(a_S/\sqrt{2\hbar}\)};
  \draw[<->,black!60] (-1.32,0)--(-1.32,\elB)
    node[midway,left,font=\scriptsize]{\(\frac{b_S}{\sqrt{2\hbar}}\)};
  \node[align=center,font=\scriptsize] at (0,-1.55)
    {\(\rho_{\partial}(\phi)/\sqrt{2\hbar}=|OX|\)\\
     \(f_S^2=a_S^2-b_S^2\), \quad área total \(h\)};
 \end{scope}
 \begin{scope}[shift={(3.65,0)}]
  \node[font=\bfseries] at (0,1.38) {Forma recíproca};
  \draw[densely dashed,black!25] (0,0) ellipse[x radius=\elA,y radius=\elB];
  \draw[ink,thick] (0,0) ellipse[x radius=\elB,y radius=\elA];
  \draw[->,black!55] (-1.28,0)--(1.38,0) node[right]{\(u\)};
  \draw[->,black!55] (0,-1.30)--(0,1.26) node[right]{\(v\)};
  \fill[ink] (0,\elF) circle[radius=1.5pt] node[right]{\(F'_+\)};
  \fill[ink] (0,-\elF) circle[radius=1.5pt] node[right]{\(F'_-\)};
  \node[align=center,font=\scriptsize] at (0,-1.55)
    {\(\eta\mapsto-\eta\), \quad \(a_S\leftrightarrow b_S\)\\
     \(R_{\rm el}\mapsto R_{\rm el}^{-1}\), \quad \(a_Sb_S=2\hbar\)};
 \end{scope}
 \draw[->,accent,thick] (1.60,.80)--(2.17,.80)
   node[midway,above,font=\scriptsize]{intercambio};
\end{tikzpicture}
\endgroup

 \caption{Semiejes, focos y reciprocidad de la elipse de acción en las
 coordenadas normalizadas \(u=Q/\sqrt{2\hbar}\),
 \(v=P/\sqrt{2\hbar}\). El sector marcado tiene fase paramétrica
 \(\theta=\pi/6\), distinta de su ángulo polar \(\phi\); su área
 de acción es \(\hbar\pi/6\). La figura utiliza el caso algebraico
 \(\xi=1/3\) para hacer legibles las relaciones: no atribuye ese valor
 a la salida angular HMT. Los focos se calculan desde los semiejes.
 A la derecha, el intercambio de ejes conserva el área total \(h\)
 e invierte el radio modular adimensional.}
 \label{fig:ii-elipse-calculada}
\end{figure}

Para el caso de control de la figura, \(\xi=1/3\), las coordenadas
normalizadas son exactamente
\[
 \eta=\tfrac12\log2,\quad
 \frac{a_S}{\sqrt{2\hbar}}=2^{1/4}=1.1892071150\ldots,\quad
 \frac{b_S}{\sqrt{2\hbar}}=
 \frac{f_S}{\sqrt{2\hbar}}=R_{\rm el}=2^{-1/4}=0.8408964153\ldots.
\]
Son números de regresión del dibujo, no valores atribuidos al generador
angular. El certificado local comprueba además la forma circular y las
orientaciones \(\xi=\pm1/3,\pm3/5\).

\subsection{Reciprocidad de forma y distancia interior}

La representación rectangular recupera
\begin{equation}
 R_{\rm el}=\sqrt{b_S/a_S}=e^{-\eta/2},\qquad
 R_{\rm el}^4=\frac{1-\xi}{1+\xi},\qquad
 \xi=\frac{1-R_{\rm el}^4}{1+R_{\rm el}^4}.
 \label{eq:ii-elipse-reciproca}
\end{equation}
La prueba consiste en sustituir
\(2\eta=\log((1+\xi)/(1-\xi))\). El intercambio de ejes
envía \(\eta\) a \(-\eta\), \(R_{\rm el}\) a
\(R_{\rm el}^{-1}\) y \(\tau=iR_{\rm el}^2\) a \(-1/\tau\).
Se conserva \(a_Sb_S\). El radio \(R_{\rm el}\) es adimensional;
el radio polar \(\rho_{\partial}\) tiene unidad de raíz de acción.

El interior admite además la distancia proyectiva de Hilbert. Si
\(z_-,x,y,z_+\) están ordenados en una cuerda de la elipse, se define
\[
 d_H(x,y)=\frac12\log
 \frac{|y-z_-|\,|x-z_+|}{|x-z_-|\,|y-z_+|}.
\]
Es una distancia adimensional: las unidades cancelan en la razón doble.
La normalización \((Q,P)\mapsto(Q/a_S,P/b_S)\) lleva la elipse
al disco unidad y preserva esa razón, pues las longitudes sobre una
misma recta se multiplican por un factor común. Transporta así la métrica
al modelo de Beltrami--Klein. Ni esta distancia interior ni la distancia
entre focos se identifica con una longitud atómica.

\subsection{Radio de Bohr: consecuencia de la escala de acción}
\label{sec:ii-elipse-bohr}

En la realización electrodinámica, la escala de Bohr se expresa como
\(a_{\rm B}=\hbar/(m_ec\alpha)\) \cite{codata2022}. En la cadena
del presente artículo, esa lectura sucede a las construcciones de
acción, vacío, estructura fina y masa electrónica. Se reciben aquí
\(\hbar,m_e,c,\alpha>0\) en una misma realización dimensional.
\begin{proposition}[Expresión areal de la escala de Bohr]
\label{prop:ii-elipse-bohr}
En esa realización,
\begin{equation}
 a_{\rm B}=\frac{a_Sb_S}{2m_ec\alpha}
 =\frac{\mathcal A(2\pi)}{2\pi m_ec\alpha}
 =\frac{\bar\lambda_C(m_e)}{\alpha},\qquad
 \bar\lambda_C(m_e)=\frac{\hbar}{m_ec}.
 \label{eq:ii-elipse-bohr}
\end{equation}
\end{proposition}
\begin{proof}
Las identidades \(a_Sb_S=2\hbar\) y
\(\mathcal A(2\pi)=2\pi\hbar\), sustituidas en la expresión
de Bohr, dan los dos primeros miembros. La definición de
\(\bar\lambda_C(m_e)\) da el tercero.
\end{proof}

El numerador tiene dimensión de acción y \(m_ec\) de momento; el
cociente es una longitud. Esta composición utiliza el producto de
semiejes y las otras tres magnitudes: no afirma que Bohr sea un semieje,
un foco, una distancia de Hilbert o el radio modular. Tampoco equipara
la masa electrónica con el cuanto circular de orden 108. La notación
\(a_{\rm B}\) evita confundir esta longitud con la normalización
\(a_0\) del operador de área.

\subsection{Realización inductiva--capacitiva y continuidad constitutiva}
\label{sec:ii-elipse-lc}

El puente electrodinámico de la elipse conserva una coordenatización adicional:
un reloj \(\omega>0\) y una impedancia de referencia \(Z_*>0\).
La construcción recuperada del corpus define
\begin{equation}
 L_\eta=\frac{Z_*}{\omega}e^{-\eta},\qquad
 C_\eta=\frac{1}{Z_*\omega}e^\eta,
 \qquad Z_\eta=\sqrt{L_\eta/C_\eta}.
 \label{eq:ii-elipse-lc}
\end{equation}
Aquí \(C_\eta\) es capacitancia, distinta del contraángulo \(C^*\).
\begin{proposition}[Producto dinámico y cociente constitutivo]
\label{prop:ii-elipse-lc}
En la coordenatización anterior,
\begin{equation}
 L_\eta C_\eta=\omega^{-2},\qquad
 \frac{Z_\eta}{Z_*}=e^{-\eta}
 =\frac{b_S}{a_S}=R_{\rm el}^{2}.
 \label{eq:ii-elipse-impedancia}
\end{equation}
Con \(Q=\sqrt{Z_*}\,q\), \(P=\Phi/\sqrt{Z_*}\), la energía
\(H_{LC}=q^2/(2C_\eta)+\Phi^2/(2L_\eta)\) toma la forma
\[
 H_{LC}=\frac\omega2(e^{-\eta}Q^2+e^\eta P^2).
\]
Sobre \(\gamma(\theta)\), sus componentes eléctrica y magnética son
\(U_E=\hbar\omega\cos^2\theta\) y
\(U_B=\hbar\omega\sin^2\theta\); su suma es \(\hbar\omega\).
\end{proposition}
\begin{proof}
El producto en \eqref{eq:ii-elipse-lc} cancela los exponentes y
\(Z_*\); su cociente es \(Z_*^2e^{-2\eta}\). Tomar la raíz
positiva y usar \eqref{eq:ii-elipse-reciproca} da
\eqref{eq:ii-elipse-impedancia}. Sustituir \(q=Q/\sqrt{Z_*}\)
y \(\Phi=\sqrt{Z_*}P\) en \(H_{LC}\) da su forma equilibrada.
Finalmente, los semiejes de \(\gamma\) cancelan los factores
\(e^{\mp\eta}\), y \(\cos^2\theta+\sin^2\theta=1\).
\end{proof}

Este puente distingue la forma, el reloj y el transductor. La elección
del sentido dinámico debe concordar con la convención simpléctica: para
las ecuaciones \(\dot Q=\partial_PH\),
\(\dot P=-\partial_QH\), la parametrización anterior satisface
\(\dot\theta=-\omega\). Las identidades energéticas no dependen
de invertir ese recorrido.

El lector resolvente del vacío actúa sobre los canales ya producidos
\(q_\pm=e^{-(A_{\rm rad}\pm C^*_{\rm rad})}\) y la incidencia
\(90/120\). Sus respuestas
\(r_\pm=(1-q_\pm^{90})/(1-q_\pm^{120})\) expresan
\(\widehat Z=r_-/r_+\) y \(\widehat c=(r_+r_-)^{-1}\).
Son respuestas espectrales, no radios geométricos. La identidad
\eqref{eq:ii-elipse-impedancia} fija aquí la realización LC;
identificarla con la impedancia constitutiva del vacío conserva el mapa
entre ambos lectores y sus coordenatizaciones. No se igualan sus valores por compartir
el antecedente angular.

\section[Retorno de hoja, incidencia areal--volumétrica y orientación constitutiva]{Retorno de hoja, incidencia areal--volumétrica\\y orientación constitutiva}
\label{iii:sec:hojas}

La estructura de hojas interviene antes de la evaluación de las magnitudes
electromagnéticas. Su función consiste en distinguir el cierre de una proyección,
la restitución de una orientación y la evolución del registro transportado.
Estas operaciones actúan sobre objetos relacionados, pero de distinto tipo: una
trayectoria proyectada puede haberse cerrado mientras su levantamiento permanece
en la hoja opuesta; la orientación puede haberse restituido mientras la memoria
conserva las dos vueltas efectuadas. La construcción constitutiva debe recibir
esta información con la estructura que la determina.

\subsection{Retorno proyectado y restitución de la cubierta orientada}

Sea $\ell$ un lazo de retorno de una extensión superviviente y sea
$\varepsilon:\langle\ell\rangle\to\{+1,-1\}$ el carácter de su
levantamiento, con $\varepsilon(\ell)=-1$. El fibrado de intervalos asociado
tiene la presentación
\begin{equation}
 \mathcal M=([0,1]\times[-1,1])/{(0,u)\sim(1,-u)}.
 \label{iii:eq:mobius}
\end{equation}
El dato decisivo es el signo del transporte sobre el lazo especificado. La
identificación de los extremos realiza geométricamente ese carácter.

\begin{proposition}[Dos órdenes de clausura]
Una vuelta del lazo cierra la base e invierte la orientación de la fibra. Dos
vueltas restituyen la orientación. En la parametrización angular del lector de
borde, estos retornos corresponden, respectivamente, al cierre areal $2\pi$ y
al cierre volumétrico $4\pi$ de la cubierta orientada.
\end{proposition}
\begin{proof}
La relación de equivalencia de \eqref{iii:eq:mobius} identifica la fibra terminal
con la inicial mediante $u\mapsto-u$. La composición de dos transportes es
$u\mapsto-(-u)=u$. Parametrizar una vuelta de la base por $2\pi$ asigna $4\pi$
a su cuadrado. La afirmación se refiere a este levantamiento y a esta lectura.
\end{proof}

La memoria del TPK conserva, además, la trayectoria recorrida y las actualizaciones
del registro. El cuadrado del carácter de signo es la identidad, pero no impone
que se anulen esas actualizaciones. La representación espinorial, cuando se
utilice, deberá transportar este dato mediante su aplicación propia.

\subsection{Descenso de la incidencia desde el calendario trítico}

En el bloque primitivo $(+1,-1,0)$, el TPK actualiza el modo aritmético mediante
$+1:m\mapsto\Sigma$, $-1:m\mapsto\Pi$ y $0:m\mapsto m$. Su órbita cíclica
es $(\Sigma,\Pi,\Pi)$. El lector de hojas asigna $\Sigma$ a la hoja areal y
$\Pi$ a la volumétrica. En el orden de bases $(e_{\rm ar},e_{\rm vol})$, la
incidencia de aristas queda representada por
\begin{equation}
 F_{\rm av}=\begin{pmatrix}0&1\\1&1\end{pmatrix}.
 \label{iii:eq:fav}
\end{equation}
Los tres pares consecutivos son $\Sigma\to\Pi$, $\Pi\to\Pi$ y
$\Pi\to\Sigma$, cada uno con multiplicidad uno. Los calendarios de longitudes
$9$, $27$ y $108$ contienen $3$, $9$ y $36$ copias del bloque primitivo. Sus
matrices de incidencia son, por tanto, esos múltiplos de $F_{\rm av}$.

Esta construcción conserva la diferencia entre frecuencia cronológica y medida
del espacio de continuaciones. La primera asigna pesos $(1/3,2/3)$ al ciclo
de modos. La segunda se define sobre la envolvente simbólica mínima de memoria
uno que admite exactamente los pares anteriores. Su matriz de adyacencia es
$F_{\rm av}$; los recuentos de trayectorias pertenecen a esa envolvente.

\Needspace{10\baselineskip}
\begin{proposition}[Pesos de retorno de las dos hojas]
Sean $F_0=0$, $F_1=1$ y $F_{n+1}=F_n+F_{n-1}$. Para $n\geq1$,
\begin{equation}
 F_{\rm av}^{n}=
 \begin{pmatrix}F_{n-1}&F_n\\F_n&F_{n+1}\end{pmatrix}.
 \label{iii:eq:fav-powers}
\end{equation}
La medida de Parry de esta envolvente asigna a las hojas pesos proporcionales
a $(1,\varphi^2)$, y la razón de recuentos diagonales converge a $\varphi^{-2}$.
\end{proposition}
\begin{proof}
La fórmula de potencias se prueba por inducción utilizando la recurrencia.
El cuadrado de $F_{\rm av}$ es estrictamente positivo. Su vector positivo de
Perron es $(1,\varphi)$, con autovalor $\varphi$, pues
$\varphi^2=\varphi+1$. Al ser simétrica la matriz, los pesos de Parry son
proporcionales a los cuadrados de las componentes de ese vector. La otra raíz
del polinomio característico es $-\varphi^{-1}$; su módulo es estrictamente
menor que $\varphi$, lo que da el límite de los cocientes diagonales.
\end{proof}

La identificación espectral de esta incidencia es posterior a su construcción
desde el calendario. La coordenada regional $\varphi$ conserva su genealogía
coinductiva anterior: el reconocimiento del autovalor no selecciona el calendario.

\subsection{Marcado de la frontera y razón areal--volumétrica}

Sea $\Pi_{\rm HMT}(X_\infty)$ el valor del lector regional de clausura ya
construido. Con los proyectores diagonales de hoja, se define
\begin{equation}
 D_\pi=\Pi_{\rm HMT}(X_\infty)P_{\rm ar}+P_{\rm vol},\qquad
 \Theta_n=
 \frac{\langle e_{\rm ar},D_\pi F_{\rm av}^{n}e_{\rm ar}\rangle}
 {\langle e_{\rm vol},F_{\rm av}^{n}e_{\rm vol}\rangle}.
 \label{iii:eq:marked-return}
\end{equation}
La normalización volumétrica es unitaria; el cierre regional marca una sola vez
el retorno areal. Aplicando \eqref{iii:eq:fav-powers},
\begin{equation}
 \Theta_n=\Pi_{\rm HMT}(X_\infty)\frac{F_{n-1}}{F_{n+1}}
 \longrightarrow\frac{\pi}{\varphi^2}.
 \label{iii:eq:av-limit}
\end{equation}
La razón de retorno queda determinada por la incidencia, la medida y el marcado.
El cociente $2\pi/4\pi$ describe, en cambio, los dos órdenes de clausura de la
cubierta. Las dos relaciones participan de la lectura de hojas sin sustituirse.

\subsection{Transporte de la orientación a los canales constitutivos}

En una representación de dos canales, sea $\mathcal R$ una involución
autoadjunta y $\mathcal S$ un operador unitario de intercambio que satisface
$\mathcal S\mathcal R\mathcal S^{-1}=-\mathcal R$. Las coordenadas angulares
ya construidas se escriben $x=A_{\rm rad}$ e $y=C^*_{\rm rad}$, con $x>|y|$.
Entonces
\begin{equation}
 B=xI+y\mathcal R,\qquad T=e^{-B},\qquad
 \mathcal ST(y)\mathcal S^{-1}=T(-y).
 \label{iii:eq:channel-conjugation}
\end{equation}
En efecto, la conjugación transforma $B(y)$ en $B(-y)$ y conmuta con su serie
exponencial convergente. Los autovalores $q_\pm=e^{-x\mp y}$ se intercambian.
Para fijar el orden de las respuestas se utiliza la cámara $x>y>0$, donde
$0<q_+<q_-<1$. La aplicación $(x,y)\mapsto(x,-y)$ la transporta a la
cámara opuesta $x>-y>0$, donde $0<q_-<q_+<1$. Ambas pertenecen al dominio
contractivo $x>|y|$. En la frontera $y=0$ los dos canales coinciden.
La conjugación intercambia las cámaras, no conserva el orden estricto de
sus etiquetas; sobre los proyectores actúa mediante
$\mathcal SP_\pm\mathcal S^{-1}=P_\mp$, con
$P_\pm=(I\pm\mathcal R)/2$.
Esta aplicación concreta transmite la orientación al operador constitutivo que
se desarrolla a continuación.

\section{Incidencia hexada–octada y funcional orientado}
\subsection{Incidencia dodecafásica y funcional de retorno}

La contribución de Barbero--Immirzi utiliza el transporte angular y la
incidencia hexada--octada. Sean $\mathsf H\subset\mathsf O$ una hexada
y una octada con el par marcado de la hexada conservado. Los conjuntos de
banderas pertinentes son
\begin{equation}
 \mathcal F_6=\mathsf H\times\binom{\mathsf H}{2}
 \lhook\joinrel\longrightarrow
 \mathcal F_8=\mathsf O\times\binom{\mathsf H}{2},
 \qquad |\mathcal F_6|=90,\quad |\mathcal F_8|=120.
 \label{iii:eq:barbero-flags}
\end{equation}
Los cardinales son $6\binom62$ y $8\binom62$; la misma incidencia
conserva el defecto normalizado
$(90-120)/(24\binom62)=-1/12$. Esta cuenta identifica los dos sectores
empleados por el retorno, mientras el calendario determina sus
multiplicidades relativas.

Para $m\in\{90,120\}$, el retorno trítico sobre un canal contractivo es
\begin{equation}
 S_m(q)=\sum_{j\geq0}q^{m+3mj}
       =\frac{q^m}{1-q^{3m}},\qquad 0<q<1.
 \label{iii:eq:barbero-return-series}
\end{equation}
Cada término corresponde a la altura inicial $m$ y a $j$ retornos de
longitud $3m$. La serie geométrica conserva, por tanto, tanto la altura
inicial como el período del retorno. El calendario dodecafásico
$\mathbb Z/12\mathbb Z$ aporta doce sectores y una única incidencia
orientada de frontera, denotada $\partial_{\rm ret}$. En el módulo
libre de sucesos, su cadena fundamental es
\begin{equation}
 c_{12}^{+}
 =\sum_{j\in\mathbb Z/12\mathbb Z}[j,\mathcal F_6]
  -[\partial_{\rm ret},\mathcal F_8].
 \label{iii:eq:barbero-cycle}
\end{equation}
El lector lineal $\mathscr L_q$ asigna $S_{90}(q)$ a cada generador
$[j,\mathcal F_6]$ y $S_{120}(q)$ al generador de frontera. Se obtiene
\begin{equation}
 \mathscr L_q(c_{12}^{+})=12S_{90}(q)-S_{120}(q).
 \label{iii:eq:barbero-weights}
\end{equation}
Los pesos $12$ y $-1$ son las multiplicidades orientadas de la cadena.
La orientación conjugada satisface $c_{12}^{-}=-c_{12}^{+}$, con los
mismos tipos de incidencia.

\begin{proposition}[Coeficiente singular del retorno fundamental]
\label{iii:prop:barbero-pole}
El coeficiente de $(1-q)^{-1}$ de
$\mathscr L_q(c_{12}^{+})$ es $1/24$.
\end{proposition}
\begin{proof}
Para cada entero $m>0$, la factorización
$1-q^{3m}=(1-q)\sum_{j=0}^{3m-1}q^j$ da
\[
 \lim_{q\uparrow1}(1-q)S_m(q)=\frac1{3m}.
\]
Aplicar esta identidad al funcional ya determinado en
\eqref{iii:eq:barbero-weights} produce
\begin{equation}
 \lim_{q\uparrow1}(1-q)\mathscr L_q(c_{12}^{+})
 =\frac{12}{270}-\frac1{360}=\frac1{24}.
 \label{iii:eq:barbero-pole}
\end{equation}
En la coordenada compleja $q-1$, el residuo es $-1/24$.
\end{proof}

La evaluación singular verifica una consecuencia de la incidencia y del
retorno trítico. El orden demostrativo queda fijado por
\eqref{iii:eq:barbero-flags}--\eqref{iii:eq:barbero-weights}; el
coeficiente singular se calcula después de construir la cadena.


\section{Operador constitutivo y reconstrucción de los canales del vacío}
\label{iii:sec:constitutiva}

Las magnitudes constitutivas se obtienen mediante una operación sobre el
transporte angular. La información de partida comprende los canales, sus
proyectores y la orientación que los distingue. Esta organización permite
seguir por separado la construcción del operador, la evaluación de sus
funciones espectrales y su realización dimensional posterior.

\subsection{Respuesta positiva sobre dos hojas}

En la cámara orientada $x>y>0$ de la sección~\ref{iii:sec:hojas}, sean
$P_\pm=(I\pm\mathcal R)/2$ y
$T=q_+P_++q_-P_-$, con $0<q_+<q_-<1$. Las incidencias construidas por el
calendario fijan los sectores $90$ y $120$. La respuesta constitutiva es
\begin{equation}
 \mathcal V=(I-T^{90})(I-T^{120})^{-1}.
 \label{iii:eq:vacuum-op}
\end{equation}
Esta definición recibe el transporte procedente del TPK y conserva las
asignaciones de hoja. La regla de respuesta y su dominio son parte explícita
de la construcción.

El cociente procede de comparar las dos incidencias del calendario sobre
un mismo transporte, no de ajustar por separado cuatro magnitudes. Si
$S=T^{30}$ y $\Sigma_j(S)=I+S+\cdots+S^{j-1}$, la factorización de
$I-S^j$ da
\begin{equation}
 \mathcal V\,\Sigma_4(S)=\Sigma_3(S),\qquad
 \mathcal V=\Sigma_3(S)\Sigma_4(S)^{-1}.
 \label{iii:eq:completion-response}
\end{equation}
Los factores $\Sigma_3$ y $\Sigma_4$ expresan respectivamente los
sectores $90=3\cdot30$ y $120=4\cdot30$. Como el espectro de $S$ está
contenido en $(0,1)$, $\Sigma_4(S)$ es invertible: la ecuación de
completación determina una única respuesta para el transporte fijado.
Esta unicidad concierne a la regla constitutiva declarada y a sus
incidencias. La identificación electromagnética añade la asignación
orientada de las hojas y las rectas dimensionales que se precisan más
abajo.

\begin{proposition}[Descomposición y positividad]
La respuesta está bien definida y
\begin{equation}
 \mathcal V=r_+P_++r_-P_-,\qquad
 r_\pm=\frac{1-q_\pm^{90}}{1-q_\pm^{120}}.
 \label{iii:eq:rchannels}
\end{equation}
Sus autovalores satisfacen $3/4<r_-<r_+<1$ en esta cámara.
\end{proposition}
\begin{proof}
Cada coeficiente $1-q_\pm^{120}$ es positivo. La inversa de $I-T^{120}$
es la suma de los proyectores divididos por esos coeficientes, lo que prueba
\eqref{iii:eq:rchannels}. Con $s=q^{30}$ se obtiene
\begin{equation}
 f(s)=\frac{1+s+s^2}{1+s+s^2+s^3},\qquad
 f'(s)=-\frac{s^2(3+2s+s^2)}{(1+s+s^2+s^3)^2}<0.
 \label{iii:eq:monotone}
\end{equation}
Los límites en $0$ y $1$ son $1$ y $3/4$, respectivamente. La monotonía
y el orden de los canales concluyen la prueba.
\end{proof}

\subsection{Cuatro coordenadas de una respuesta común}

Se definen las evaluaciones normalizadas
\begin{equation}
 \widehat\varepsilon=r_+^2,\qquad
 \widehat\mu=r_-^2,\qquad
 \widehat Z=\frac{r_-}{r_+},\qquad
 \widehat c=\frac1{r_+r_-}.
 \label{iii:eq:four}
\end{equation}
La primera pareja conserva las respuestas cuadráticas asignadas a cada hoja.
El producto recoge el modo común; el cociente distingue su orientación relativa.
Se verifican exactamente
\begin{equation}
 \widehat\mu\widehat\varepsilon=\widehat c^{-2},\qquad
 \frac{\widehat\mu}{\widehat\varepsilon}=\widehat Z^2,
 \qquad
 r_+=\frac1{\sqrt{\widehat Z\widehat c}},\quad
 r_-=\sqrt{\frac{\widehat Z}{\widehat c}}.
 \label{iii:eq:four-identities}
\end{equation}
Las raíces positivas fijan la inversión. Bajo la conjugación de hojas,
leída respecto de las marcas $P_\pm$ fijas, se
intercambian $\widehat\varepsilon$ y $\widehat\mu$, se invierte
$\widehat Z$ y se conserva $\widehat c$. La respuesta se transporta entonces
a la cámara $x>-y>0$: las fórmulas siguen siendo válidas, con los órdenes
$q_-<q_+$ y $r_+<r_-$ invertidos. Por consiguiente, el modo común
y la orientación son datos recuperables mediante observables complementarios.

\begin{proposition}[Determinación de las evaluaciones cuadráticas]
Sea $W=r_+P_++r_-P_-$ una respuesta positiva sobre las dos hojas.
Fijadas las lecturas de producto y cociente
\[
 \widehat c^{-1}=\det W=r_+r_-,
 \qquad \widehat Z=r_-/r_+,
\]
existe una única pareja positiva $(\widehat\varepsilon,\widehat\mu)$
que satisface
$\widehat\mu\widehat\varepsilon=\widehat c^{-2}$ y
$\widehat\mu/\widehat\varepsilon=\widehat Z^2$.
Es precisamente $(r_+^2,r_-^2)$.
\end{proposition}
\begin{proof}
La positividad permite tomar las raíces sin ambigüedad:
\[
 \widehat\mu=\frac{\widehat Z}{\widehat c}
 =\frac{r_-}{r_+}r_+r_-=r_-^2,\qquad
 \widehat\varepsilon=\frac1{\widehat Z\widehat c}
 =\frac{r_+}{r_-}r_+r_-=r_+^2.
\]
Las expresiones verifican las dos relaciones y son las únicas raíces
positivas de ellas.
\end{proof}

Así, los cuadrados no constituyen dos elecciones independientes añadidas
a los canales: quedan determinados por las lecturas multiplicativa y
orientada, una vez fijadas las relaciones constitutivas. Esta proposición
establece la unicidad dentro de ese sistema de reglas. La interpretación
de las reglas como respuesta electromagnética conserva su contenido
adicional de representación física.

La lectura determinantal del transporte y la respuesta constitutiva
también deben conservar su orden. Para $T=e^{-xI-y\mathcal R}$ se tiene
$\det T=e^{-2x}$, mientras que
\[
 \det\mathcal V=f(q_+^{30})f(q_-^{30})=\widehat c^{-1}.
\]
Ambas proceden del mismo operador de dos hojas, pero evalúan funciones
diferentes de su espectro. En general, aplicar primero el determinante
y después una función racional no equivale a tomar el determinante de
esa función del operador. Mantener los dos autovalores hasta completar
la respuesta conserva la información de orientación que un determinante
aislado ya no distingue.

\subsection{Inversión cúbica y recuperación angular}

\begin{theorem}[Inversión constitutiva en el dominio contractivo]
Para cada $r\in(3/4,1)$ existe una única $s\in(0,1)$ tal que $f(s)=r$.
Esta $s$ es la raíz única en dicho intervalo de
\begin{equation}
 r s^3+(r-1)(1+s+s^2)=0.
 \label{iii:eq:cubic}
\end{equation}
La respuesta etiquetada reconstruye $s_\pm$, $q_\pm=s_\pm^{1/30}$ y
\begin{equation}
 x=-\frac12\log(q_+q_-),\qquad
 y=\frac12\log\frac{q_-}{q_+}.
 \label{iii:eq:recover-angular}
\end{equation}
\end{theorem}
\begin{proof}
La derivada y los límites de \eqref{iii:eq:monotone} prueban que $f$ es una
biyección decreciente entre los intervalos indicados. Multiplicar $f(s)=r$
por su denominador produce \eqref{iii:eq:cubic}. La raíz positiva de orden
$30$ recupera el canal. Finalmente, suma y diferencia de
$\log q_\pm=-x\mp y$ dan \eqref{iii:eq:recover-angular}.
\end{proof}

La reconstrucción conserva las etiquetas de hoja. Con el operador completo
$\mathcal V$ se conservan también los proyectores y se aplica la función
inversa por cálculo espectral. Recuperar dos coordenadas angulares es una
operación definida sobre esta representación; la traza histórica completa
del TPK conserva datos adicionales.

El intervalo $r\in(3/4,1)$ y la función inversa de este teorema pertenecen
a la normalización interna de \eqref{iii:eq:four}. Si se dispone de
coordenadas expresadas en otra coordenatización de unidades, se recupera primero
esta normalización mediante la transformación de la
sección~\ref{iii:sec:evaluacion}; después se aplica la inversión cúbica.

% Recuperación expositiva de INTERRELACIONES_ESTRUCTURALES_VACIO.md, §§6--8.
% Insertar después de la inversión cúbica y antes de las coordenadas logarítmicas.
\subsection{Composición transportada de las respuestas constitutivas}
\label{amp-comp-sec}

El transporte angular y las incidencias $90$ y $120$ proceden de la
misma estructura APP--TRIT--TPK: las hojas conservan las evaluaciones y
sus acarreos; el régimen trítico conserva la orientación; el calendario
TPK transporta fase y memoria hasta los canales angulares. La respuesta
constitutiva aplica a esos canales la función $f$ de
\eqref{iii:eq:monotone}. Su inversión permite expresar, en las coordenadas
de respuesta, la composición de transportes del mismo marco. Se conserva
así el orden constructivo: transporte, composición y evaluación espectral.
La operación obtenida actúa sobre esta representación del estado enriquecido;
su inversa recupera los canales, mientras la historia completa de memoria
permanece en la estructura de procedencia.

Sean $s=q^{30}$ y $\psi=f^{-1}$. El exponente $30$ es el máximo común
divisor de las incidencias $90$ y $120$. La inversión cúbica determina
$\psi:(3/4,1)\longrightarrow(0,1)$. Adjuntando la respuesta de frontera
$f(1)=3/4$, extendemos la biyección a
\[
 \psi:D\longrightarrow(0,1],\qquad D=[3/4,1).
\]
En esa frontera se utiliza la función racional $f$; el cociente
$(1-q^{90})/(1-q^{120})$ tiene allí una singularidad removible.

\begin{proposition}[Monoide de respuestas y coordenada aditiva]
\label{amp-comp-monoid}
La operación
\begin{equation}
 r\odot t=f\bigl(\psi(r)\psi(t)\bigr),\qquad r,t\in D,
 \label{amp-comp-law}
\end{equation}
hace de $D$ un monoide conmutativo y cancelativo con elemento neutro $3/4$.
La aplicación
\begin{equation}
 \ell(r)=-\frac1{30}\log\psi(r)
 \label{amp-comp-character}
\end{equation}
es un isomorfismo de ese monoide con $(\mathbb R_{\geq0},+)$.
El intervalo $(3/4,1)$ es cerrado bajo $\odot$; el neutro pertenece a
la extensión de frontera. El único elemento invertible dentro de $D$
es el neutro.
\end{proposition}
\begin{proof}
El producto de dos elementos de $(0,1]$ pertenece al mismo intervalo y
\[
 \psi(r\odot t)=\psi(r)\psi(t).
\]
Para tres respuestas, ambos órdenes de asociación tienen imagen
$\psi(r)\psi(t)\psi(u)$; la inyectividad de $\psi$ prueba la
asociatividad. La conmutatividad procede del producto escalar. Si
$r\odot t=r\odot u$, la positividad de $\psi(r)$ permite cancelar y
obtener $\psi(t)=\psi(u)$, luego $t=u$. El parámetro $1$ corresponde a
$3/4$ y proporciona la identidad. El logaritmo transforma el producto en
suma, y su rango en $(0,1]$ prueba que $\ell$ es una biyección sobre
$\mathbb R_{\geq0}$. En particular, para la respuesta de un canal $q$,
$\ell(r)=-\log q$. Si ambos parámetros son estrictamente menores que
$1$, su producto también lo es. Finalmente, dos números de $(0,1]$
tienen producto $1$ únicamente si ambos son $1$.
\end{proof}

\begin{proposition}[Cancelación admisible y prolongación inversa]
\label{amp-comp-inverse}
Para $r,u\in D$, la ecuación $r\odot t=u$ tiene solución $t\in D$
exactamente cuando $u\geq r$. Esa solución es única y viene dada por
\begin{equation}
 t=f\!\left(\frac{\psi(u)}{\psi(r)}\right).
 \label{amp-comp-cancellation}
\end{equation}
La misma función $f$ es una biyección de $(0,\infty)$ sobre $(0,1)$.
Con esta extensión de $\psi$, la ley~\eqref{amp-comp-law} transporta el
grupo multiplicativo de los parámetros positivos al intervalo $(0,1)$.
La inversión de grupo satisface
\begin{equation}
 r^{\ominus}=\psi(r)r,\qquad
 r\odot r^{\ominus}=\frac34.
 \label{amp-comp-group-inverse}
\end{equation}
\end{proposition}
\begin{proof}
La ecuación equivale a
$\psi(t)=\psi(u)/\psi(r)$. Este cociente pertenece a $(0,1]$
exactamente cuando $\psi(u)\leq\psi(r)$, que por monotonía decreciente
equivale a $u\geq r$. La biyección prueba existencia y unicidad.
La derivada de \eqref{iii:eq:monotone} es negativa para todo $s>0$;
los límites de $f$ en $0$ y en infinito son $1$ y $0$, respectivamente.
Esto establece la biyección extendida. Multiplicando numerador y
denominador por $s^3$ se obtiene
\[
 f(s^{-1})=\frac{s^3+s^2+s}{s^3+s^2+s+1}=s f(s).
\]
La inversa multiplicativa de $s=\psi(r)$ tiene, por tanto, respuesta
$\psi(r)r$. Su composición con $r$ corresponde al parámetro $1$.
\end{proof}

La prolongación a $s>1$, o equivalentemente a $q>1$, es la extensión
algebraica del lector. El sector físico contractivo del capítulo conserva
su dominio declarado. La coordenada $\ell$ se prolonga a un isomorfismo
con $(\mathbb R,+)$ y cambia de signo bajo la inversión de grupo.

\subsubsection{Realización operatoria en un marco de hojas común}

\begin{proposition}[Composición con proyectores marcados]
\label{amp-comp-operator}
Sean $\mathcal R=\mathcal R^*$, $\mathcal R^2=I$ y
$P_\pm=(I\pm\mathcal R)/2$. Para $j=a,b$, sean
\[
 T_j=q_{j,+}P_++q_{j,-}P_-,\qquad
 \mathcal V_j=f(T_j^{30}),\qquad 0<q_{j,\pm}<1.
\]
La composición $T_{a\circ b}=T_aT_b$ satisface
\begin{equation}
 \mathcal V_{a\circ b}
 =f\bigl(\psi(\mathcal V_a)\psi(\mathcal V_b)\bigr).
 \label{amp-comp-operator-law}
\end{equation}
Sus dos respuestas son $r_{a,\pm}\odot r_{b,\pm}$. Si
$T_j=\exp[-(x_jI+y_j\mathcal R)]$ con $x_j>|y_j|$, entonces
\begin{equation}
 T_aT_b=
 \exp[-((x_a+x_b)I+(y_a+y_b)\mathcal R)],
 \label{amp-comp-angular-addition}
\end{equation}
y $x_a+x_b>|y_a+y_b|$.
\end{proposition}
\begin{proof}
La ortogonalidad $P_+P_-=0$ da
\[
 T_aT_b=q_{a,+}q_{b,+}P_++q_{a,-}q_{b,-}P_-.
\]
En particular, $T_a$ y $T_b$ conmutan. El cálculo funcional en estos
proyectores proporciona
$\psi(\mathcal V_j)=T_j^{30}$ y
$(T_aT_b)^{30}=T_a^{30}T_b^{30}$; aplicar $f$ prueba
\eqref{amp-comp-operator-law}. Para toda función $g$ definida en los
dos autovalores se tiene, más explícitamente,
\[
 P_\pm g(T_j)=g(q_{j,\pm})P_\pm.
\]
Como $q_{j,\pm}=\exp[-(x_j\pm y_j)]$, sus productos prueban
\eqref{amp-comp-angular-addition}. La desigualdad triangular concluye
$x_a+x_b>|y_a|+|y_b|\geq|y_a+y_b|$.
\end{proof}

Las coordenadas recuperadas pueden escribirse directamente como
\[
 x=\frac{\ell(r_+)+\ell(r_-)}2,\qquad
 y=\frac{\ell(r_+)-\ell(r_-)}2.
\]
La composición suma estas parejas. El intercambio simultáneo de hojas
permuta los dos canales de cada transporte, actúa como
$(x,y)\mapsto(x,-y)$ y es un automorfismo de la ley componente a
componente. La inversión de ambos transportes en la extensión algebraica
actúa como $(x,y)\mapsto(-x,-y)$. Las dos involuciones conmutan y
conservan funciones distintas: una intercambia las marcas de orientación;
la otra invierte el transporte.

Para representar el intercambio mediante conjugación se utiliza un
operador $L$ que satisfaga $L\mathcal R L^{-1}=-\mathcal R$. En la
representación
\[
 \mathcal R=\begin{pmatrix}0&1\\1&0\end{pmatrix},\qquad
 L=\begin{pmatrix}1&0\\0&-1\end{pmatrix},
\]
la multiplicación directa da $LP_\pm L^{-1}=P_\mp$ y, por tanto,
$LT_jL^{-1}=q_{j,-}P_++q_{j,+}P_-$. La conjugación por
$\mathcal R$ conserva sus propios proyectores. Una traza que reúne los
canales tampoco reemplaza a los proyectores en esta identidad funcional:
la información de sus etiquetas interviene en la composición.

La hipótesis de marco común es sustantiva. Por ejemplo, los operadores
positivos definidos
\[
 A_0=\begin{pmatrix}1/2&0\\0&1/3\end{pmatrix},\qquad
 B_0=\begin{pmatrix}1/2&1/6\\1/6&1/2\end{pmatrix}
\]
tienen autovalores positivos, pero
\[
 A_0B_0=\begin{pmatrix}1/4&1/12\\1/18&1/6\end{pmatrix}
\]
carece de autoadjunción. La proposición especifica la composición de
transportes en los proyectores comunes; su dominio no comprende una mezcla
arbitraria de medios materiales. La realización de una concatenación de
rutas HMT conserva además sus datos de memoria: el cierre algebraico del
cono de parámetros describe la representación y no selecciona por sí solo
nuevas rutas primarias.

\subsubsection{Producto de canales y composición de transportes}

La velocidad normalizada de un estado sigue determinada por
\[
 \widehat c=(r_+r_-)^{-1}.
\]
Aquí el producto ordinario reúne las dos respuestas de ese mismo estado.
La ley $\odot$ evalúa otra operación: la composición de dos transportes
en cada hoja. El siguiente cálculo exacto distingue sus argumentos:
\begin{equation}
 f(1/2)=\frac{14}{15},\qquad
 \frac{14}{15}\odot\frac{14}{15}=f(1/4)=\frac{84}{85}
 \ne\frac{196}{225}=\left(\frac{14}{15}\right)^2.
 \label{amp-comp-rational-example}
\end{equation}
Las fracciones de este ejemplo son parámetros algebraicos de comprobación.
La expresión de $\widehat c$ y las identidades constitutivas anteriores
permanecen inalteradas. El contenido añadido es una ley explícita para
transportar la composición y recuperar su coordenada angular aditiva.


\subsection{Coordenadas logarítmicas y refinamiento}

Sean $\mathcal E=\log(r_+r_-)$ y $\mathcal B=\log(r_-/r_+)$. Entonces
\begin{equation}
 \log\mathcal V=\tfrac12\mathcal E I-\tfrac12\mathcal B\mathcal R.
 \label{iii:eq:logvac}
\end{equation}
Esta descomposición exhibe la componente invariante y la componente orientada.
En la amplificación $\mathcal V_m=\mathcal V\otimes I_{9^m}$, las
expectativas parciales normalizadas satisfacen
$E_m(\mathcal V_{m+1})=\mathcal V_m$. Para todo polinomio $p$,
$p(\mathcal V_m)=p(\mathcal V)\otimes I_{9^m}$; de ahí
$E_m(p(\mathcal V_{m+1}))=p(\mathcal V_m)$. La continuidad uniforme
extiende la identidad a las funciones continuas del espectro. Se conservan
así los proyectores, las potencias y el logaritmo en la torre indicada.

\subsection{Tipos dimensionales de las cuatro evaluaciones}

Con las rectas de acción, carga y velocidad, provistas de secciones
$u_S,u_Q,u_C$, la realización dimensional se escribe
\begin{align}
 Z_{\rm vac}&=\widehat Z u_Su_Q^{-2},&
 c_{\rm vac}&=\widehat c u_C,\\
 \mu_{\rm vac}&=\widehat\mu u_Su_C^{-1}u_Q^{-2},&
 \varepsilon_{\rm vac}&=\widehat\varepsilon u_S^{-1}u_C^{-1}u_Q^2.
 \label{iii:eq:dimensions}
\end{align}
Multiplicando y dividiendo estas expresiones se preservan
$\mu_{\rm vac}\varepsilon_{\rm vac}=c_{\rm vac}^{-2}$ y
$\mu_{\rm vac}/\varepsilon_{\rm vac}=Z_{\rm vac}^2$.
La coordenatización SI evalúa después las secciones. La comparación física requiere
exponer su construcción y normalización junto a la evaluación; las cuatro
coordenadas normalizadas de \eqref{iii:eq:four} conservan su tipo adimensional.

% INSERCIÓN FOCAL AUTORIZADA: no introduce un capítulo ni el concepto de radión.
% Incluir dentro de la sección constitutiva, después de la inversión cúbica
% y de la declaración de las cartas dimensionales. Etiquetas fuera de grupos
% de prefijado automático. Conserva íntegramente el bloque geométrico anterior.
\subsection{Reconstrucción de la forma de acción desde la respuesta constitutiva}
\label{iii:subsec:puente-forma-LC}

La respuesta constitutiva y la realización inductiva--capacitiva reciben
el mismo par angular construido desde APP--TRIT--TPK. La primera evalúa
funciones racionales de los dos canales; la segunda utiliza la razón de
sus exponentes para determinar la forma de la elipse de acción. La
inversión cúbica permite componer ambas lecturas sin identificar sus
cocientes numéricos.

Sean $x=A_{\rm rad}$ e $y=C^*_{\rm rad}$, con $x>|y|$, las coordenadas
levantadas ya construidas, y consérvense sus etiquetas de hoja. Escribimos
\begin{equation}
 q_\pm=e^{-(x\pm y)},\qquad s_\pm=q_\pm^{30},\qquad
 r_\pm=f(s_\pm),\qquad
 f(s)=\frac{1+s+s^2}{1+s+s^2+s^3}.
 \label{iii:eq:forma-canales}
\end{equation}
El dominio contractivo da $s_\pm\in(0,1)$ y
$r_\pm\in(3/4,1)$. En la cámara $x>y>0$ se cumplen
$s_+<s_-$ y $r_+>r_-$; el intercambio de hojas transporta esta cámara
a la orientación opuesta.

La imagen de las dos coordenadas constitutivas normalizadas es el dominio
\begin{equation}
 \mathcal U=\left\{(z,v)\in(0,\infty)^2:
   (zv)^{-1/2}\in(3/4,1),\quad
   (z/v)^{1/2}\in(3/4,1)\right\}.
 \label{iii:eq:forma-dominio}
\end{equation}
Aquí $(z,v)=(\widehat Z,\widehat c)$. Las restricciones expresan la
pertenencia de los canales reconstruidos al dominio del teorema de
inversión; no constituyen una elección ulterior de sus valores.

\Needspace{23\baselineskip}
\begin{proposition}[Composición constitutiva de la forma y de la impedancia LC]
\label{iii:prop:forma-LC}
Para $(\widehat Z,\widehat c)\in\mathcal U$, sean $s_\pm\in(0,1)$ las
raíces únicas de
\begin{equation}
 r_\pm s_\pm^3+(r_\pm-1)(1+s_\pm+s_\pm^2)=0,
 \qquad
 r_+=\frac1{\sqrt{\widehat Z\widehat c}},\quad
 r_-=\sqrt{\frac{\widehat Z}{\widehat c}}.
 \label{iii:eq:forma-inversion}
\end{equation}
La razón angular, la anisotropía de la elipse de acción y la impedancia
LC relativa se reconstruyen exactamente mediante
\begin{align}
 \frac{y}{x}
 &=\frac{\log s_+-\log s_-}{\log s_++\log s_-},
 \label{iii:eq:forma-razon}\\
 \eta_{\rm el}
 &=\frac12\log\frac{\log s_+}{\log s_-},
 \label{iii:eq:forma-eta}\\
 \frac{Z_{\rm LC}}{Z_*}
 &=e^{-\eta_{\rm el}}
  =\sqrt{\frac{\log s_-}{\log s_+}}.
 \label{iii:eq:forma-LC}
\end{align}
En la última identidad $Z_{\rm LC}=Z_{\eta_{\rm el}}$ es la impedancia
de la coordenatización LC de \eqref{eq:ii-elipse-lc}.
\end{proposition}
\begin{proof}
La inversión cúbica de \eqref{iii:eq:cubic} determina $s_\pm$ de manera
unívoca, porque $f$ es una biyección decreciente de $(0,1)$ sobre
$(3/4,1)$. Sus raíces positivas de orden $30$ recuperan $q_\pm$.
Las identidades
\[
 \log s_+=-30(x+y),\qquad \log s_-=-30(x-y)
\]
prueban \eqref{iii:eq:forma-razon} por diferencia y suma. Ambos logaritmos
son negativos; su cociente es positivo y su suma no es nula. Por ello,
\[
 \frac12\log\frac{\log s_+}{\log s_-}
   =\frac12\log\frac{x+y}{x-y}
   =\operatorname{artanh}(y/x)=\eta_{\rm el}.
\]
La relación $Z_{\eta_{\rm el}}/Z_*=e^{-\eta_{\rm el}}$, demostrada en
\eqref{eq:ii-elipse-impedancia}, da \eqref{iii:eq:forma-LC}.
\end{proof}

La composición conserva la información necesaria para cada reconstrucción.
Con $\hbar$ se recuperan
$a_S=\sqrt{2\hbar}\,e^{\eta_{\rm el}/2}$ y
$b_S=\sqrt{2\hbar}\,e^{-\eta_{\rm el}/2}$; con la coordenatización
$(\omega,Z_*)$ se recuperan la inductancia y la capacitancia ya definidas.
La acción, el reloj y la referencia dimensional acompañan la forma como
datos tipados. La inversión de dos coordenadas adimensionales no
reconstruye sus unidades ni toda la memoria del estado.

La normalización interviene antes de resolver las ecuaciones cúbicas.
Si se reciben $Z_{\rm vac}$ y $c_{\rm vac}$ en una coordenatización dimensional,
las relaciones \eqref{iii:eq:dimensions} determinan primero
\[
 \widehat Z=\frac{Z_{\rm vac}}{u_Su_Q^{-2}},\qquad
 \widehat c=\frac{c_{\rm vac}}{u_C}.
\]
Esas son las coordenadas a las que se aplica
\eqref{iii:eq:forma-inversion}. La referencia $Z_*$ pertenece a la
realización LC y se conserva expresamente en \eqref{iii:eq:forma-LC}.

Los dos cocientes están relacionados mediante los canales, pero evalúan
funciones diferentes:
\begin{equation}
 \widehat Z=\frac{f(s_-)}{f(s_+)},\qquad
 \frac{Z_{\rm LC}}{Z_*}
       =\sqrt{\frac{\log s_-}{\log s_+}}.
 \label{iii:eq:forma-dos-cocientes}
\end{equation}
La reconstrucción anterior proporciona el mapa entre ambas lecturas.
En particular, intercambiar las hojas invierte ambos cocientes, cambia
$\eta_{\rm el}$ por $-\eta_{\rm el}$ y conserva $\widehat c$.
En el límite simétrico $y=0$, los canales coinciden y ambos cocientes
valen $1$. Fuera de ese caso, la igualdad de origen no impone igualdad
de las dos funciones. La forma de acción es recuperable desde la
respuesta constitutiva conservando orientación y normalización.

% Recuperación autocontenida: sector angular de árbol y reconstrucción Higgs--vacío.
% Propietarios y alcance: higgs_procedencia.md. Incorporar después del operador constitutivo.
\section{Reconstrucción conjunta del calibre, Higgs y la respuesta constitutiva}
\label{amp-higgs-seccion}

La respuesta del vacío y el sector electrodébil admiten lecturas complementarias
del transporte angular ya construido. La primera conserva los dos canales
marcados; el determinante conserva su contracción común; la componente angular
de Higgs conserva además la orientación que desaparece en ese determinante.
La composición de estas lecturas permite reconstruir su antecedente espectral
y expresar la impedancia mediante el autoacoplamiento escalar y los acoplamientos
de calibre, dentro de una misma realización de árbol.

\subsection{Procedencia angular y firmas de las rutas}

Se conserva la cadena APP--TRIT--TPK: la prolongación superviviente determina
la ruta, ésta produce su registro orientado y el lector del registro expresa
la firma angular. Las coordenadas $\pi=\pi_{\rm HMT}$ y
$\alpha=\alpha_{\rm HMT}$ son salidas previas de la misma generación;
$A$ y $C^*$ son sus representaciones angulares posteriores. En este bloque se
usan radianes,
\begin{equation}
 x=\frac{\pi A_{\deg}}{180},\qquad
 y=\frac{\pi C^*_{\deg}}{180},\qquad
 A_{\deg}=1000\alpha,\qquad 0<y<x.
 \label{amp-higgs-angulos}
\end{equation}
Por tanto, $x=50\pi\alpha/9$. Los canales y las marcas de hoja son los
de la sección~\ref{iii:sec:hojas}:
\begin{equation}
 \begin{aligned}
 q_+&=e^{-x-y},& q_-&=e^{-x+y},& D&=q_+q_-=e^{-2x},\\
 P_\pm&=\frac{I\pm\mathcal R}{2},&
 T&=q_+P_++q_-P_-.
 \end{aligned}
 \label{amp-higgs-transporte}
\end{equation}
Aquí $\mathcal R$ es la involución autoadjunta marcada y cada proyector
tiene rango uno en la realización bicapa básica.

El lector angular de duración--perduración es
\begin{equation}
 \chi_{\rm dur}(n,s)=\exp\!\left(nx+\frac{s}{6}y\right).
 \label{amp-higgs-caracter}
\end{equation}
Para conservar la procedencia finita de los coeficientes que se utilizarán,
reunimos las huellas nominales de las tres rutas del calendario de longitud
$108$. Las columnas de desplazamientos cardinales se leen como $(N,E,S,O)$;
la pareja siguiente conserva la hoja y el residuo expresados por esa ruta.
La tabla conserva, para cada ruta identificada, la huella y la firma
angular expresadas por la cadena anterior.
\begin{center}
\small
\begin{tabular}{c c c c c}
\hline
Sector & Ruta y eje & $(N,E,S,O)$ & Hoja/residuo & $(n,s)$\\
\hline
$W$ & $15/37$, E--O & $(23,44,11,30)$ & $(41,8)$ & $(94,-1)$\\
$Z$ & $20/23$, E--O & $(33,45,10,20)$ & $(57,7)$ & $(95,-1)$\\
$H$ & $24/29$, N--S & $(40,34,22,12)$ & $(47,10)$ & $(97,9)$\\
\hline
\end{tabular}
\end{center}
Las tres huellas tipadas conservan el identificador de ruta, la orientación
y el registro de procedencia de las firmas empleadas. La evaluación angular
comienza después de la extracción desde el registro enriquecido; la tabla
reúne sus salidas para componer los caracteres en esta realización.
Las marcas $W_\pm$ que siguen designan coordenadas de enrollamiento,
distintas del nombre de la especie $W$:
\begin{equation}
 \begin{aligned}
 W_\pm&=6n\pm s, &
 n&=\frac{W_++W_-}{12},& s&=\frac{W_+-W_-}{2},\\
 (n,s)_W&=(94,-1), & (W_+,W_-)_W&=(563,565),\\
 (n,s)_Z&=(95,-1), & (W_+,W_-)_Z&=(569,571),\\
 (n,s)_H&=(97,9), & (W_+,W_-)_H&=(591,573).
 \end{aligned}
 \label{amp-higgs-firmas}
\end{equation}
La matriz de la primera aplicación es
$\bigl(\begin{smallmatrix}6&1\\6&-1\end{smallmatrix}\bigr)$, con determinante
$-12$ e imagen $\{(u,v)\in\mathbb Z^2:u+v\in12\mathbb Z\}$.
En efecto, las fórmulas inversas anteriores son enteras precisamente en
esa imagen. El carácter conserva las dos orientaciones mediante
\begin{equation}
 \chi_{\rm dur}(n,s)
 =\left(q_+^{-W_+}q_-^{-W_-}\right)^{1/12}.
 \label{amp-higgs-caracter-canales}
\end{equation}
La prueba se obtiene tomando logaritmos: el numerador es
$W_+(x+y)+W_-(x-y)=12nx+2sy$.
Así quedan identificados la raíz duodécima y el divisor $6$ del lector.

\begin{proposition}[Cocientes angulares determinados por las firmas]
En una sección de escala común, con los factores de refinamiento compartidos
que permiten su cancelación en los cocientes, se cumple
\begin{equation}
 \frac{m_W^{\angle}}{m_Z^{\angle}}=e^{-x},\qquad
 \frac{m_H^{\angle}}{m_W^{\angle}}
 =e^{3x+5y/3},\qquad
 \left(\frac{m_H^{\angle}}{m_W^{\angle}}\right)^2
 =e^{6x+10y/3}.
 \label{amp-higgs-cocientes}
\end{equation}
\end{proposition}
\begin{proof}
Dividir los caracteres de $W$ y $Z$ resta sus firmas: $(94,-1)-(95,-1)=(-1,0)$.
Para $H/W$, la diferencia es $(97,9)-(94,-1)=(3,10)$.
La ecuación~\eqref{amp-higgs-caracter} da $3x+(10/6)y$; el cuadrado
multiplica ese exponente por dos. La escala común y los refinamientos
compartidos se cancelan antes de estas operaciones.
\end{proof}
Estos cocientes conservan el identificador de la coordenatización angular. La aplicación
a masas con refinamientos sectoriales distintos conserva sus factores
adicionales; la condición de cancelación forma parte del enunciado.

\subsection{Realización racionalizada de árbol}

La lectura electrodébil de esta coordenatización define
$c_W^2=D$, $s_W^2=1-D$ y $e_g^2=4\pi\alpha$. La raíz positiva fija
$g=e_g/s_W$ y $g'=e_g/c_W$. Se obtiene
\begin{equation}
 g^2=\frac{4\pi\alpha}{1-D},\qquad
 g'^2=\frac{4\pi\alpha}{D}.
 \label{amp-higgs-calibre}
\end{equation}
La cantidad $e_g$ es un acoplamiento adimensional en la normalización
racionalizada, distinto de una sección dimensional de carga.
La convención del potencial escalar es
\begin{equation}
 \begin{gathered}
 V(H)=-\mu_H^2H^\dagger H+\lambda_H(H^\dagger H)^2,\\
 m_H^2=2\lambda_Hv^2,\qquad m_W^2=\frac{g^2v^2}{4}.
 \end{gathered}
 \label{amp-higgs-potencial}
\end{equation}
La misma $v$ interviene en ambas relaciones. Al dividirlas, su escala se
cancela y $m_H^2/m_W^2=8\lambda_H/g^2$. Con
\eqref{amp-higgs-cocientes} resulta
\begin{equation}
 \boxed{\lambda_H=\frac{g^2}{8}\exp\!\left(6x+\frac{10}{3}y\right).}
 \label{amp-higgs-lambda}
\end{equation}
El factor $8$ procede de las dos normalizaciones cuadráticas de
\eqref{amp-higgs-potencial}; los coeficientes $6$ y $10/3$ proceden de
las firmas. La distinción conserva la operación que produce cada factor.
La composición recibe las salidas HMT previas y usa esta realización física
de árbol a escala común; un cambio de normalización del campo, de escala o
de esquema se transporta también en sus relaciones de masas y acoplamientos.

\subsection{Inversa conjunta y dominio exacto}

\begin{theorem}[Reconstrucción angular mediante calibre e Higgs]
\label{amp-higgs-inversa-teorema}
Para $\pi>0$ fijo, la aplicación de
\eqref{amp-higgs-transporte}, \eqref{amp-higgs-calibre} y
\eqref{amp-higgs-lambda}, sobre $\alpha>0$, $x>y>0$, es inyectiva y tiene
la inversa
\begin{equation}
 \begin{aligned}
 D&=\frac{g^2}{g^2+g'^2},&
 x&=-\frac12\log D,&
 \alpha&=\frac{g^2g'^2}{4\pi(g^2+g'^2)},\\
 y&=\frac3{10}\left[\log\!\left(\frac{8\lambda_H}{g^2}\right)-6x\right].
 \end{aligned}
 \label{amp-higgs-inversa}
\end{equation}
Su imagen exacta está formada por las ternas $g,g',\lambda_H>0$ tales que,
con $D$ y $x$ reconstruidos como arriba,
\begin{equation}
 \frac{g^2}{8}e^{6x}<\lambda_H<\frac{g^2}{8}e^{28x/3}.
 \label{amp-higgs-dominio}
\end{equation}
\end{theorem}
\begin{proof}
Las dos ecuaciones de calibre dan
$g^2/(g^2+g'^2)=D$ y
$g^2g'^2/(g^2+g'^2)=4\pi\alpha$. El logaritmo real recupera $x$ y
la ecuación escalar recupera $y$. Las condiciones $g,g'>0$ implican
$0<D<1$ y $x>0$. La desigualdad~\eqref{amp-higgs-dominio} equivale,
mediante el logaritmo estrictamente creciente, a
$0<\log(8\lambda_H/g^2)-6x<10x/3$, es decir, $0<y<x$.
Recíprocamente, las fórmulas reconstruidas producen
$4\pi\alpha/(1-D)=g^2$, $4\pi\alpha/D=g'^2$ y
$g^2e^{6x+10y/3}/8=\lambda_H$. Las raíces positivas recuperan los
acoplamientos originales. Esto prueba necesidad, suficiencia y unicidad.
\end{proof}

La inversión anterior describe exactamente esta aplicación entre lectores.
La sección generada por HMT conserva además la relación de
\eqref{amp-higgs-angulos}, que se expresa como
\begin{equation}
 -\frac12\log\!\left(\frac{g^2}{g^2+g'^2}\right)
 =\frac{25}{18}\frac{g^2g'^2}{g^2+g'^2}.
 \label{amp-higgs-compatibilidad-alpha}
\end{equation}
Se obtiene sustituyendo la inversa de $\alpha$ en $x=50\pi\alpha/9$.
El contraángulo conserva igualmente su ecuación de retorno anterior.
Por ello, la imagen del lector y la sección producida por el generador
son objetos distintos: las compatibilidades previas seleccionan la sección
dentro del dominio reconstruible. La operación inversa actúa después de
la generación y permite comprobar sus salidas conservadas.

\subsection{Reconstrucción Higgs--impedancia}

Para expresar el contenido orientado sin alterar los canales,
definimos
\begin{equation}
 \rho_{\rm or}=\frac{q_-}{q_+}=e^{2y},\qquad
 \Xi_H=\frac{8D^3\lambda_H}{g^2}.
 \label{amp-higgs-Xi}
\end{equation}
La ecuación~\eqref{amp-higgs-lambda} y $D^3=e^{-6x}$ dan
\begin{equation}
 \begin{aligned}
 \Xi_H&=e^{10y/3}=\rho_{\rm or}^{5/3},&
 \rho_{\rm or}&=\Xi_H^{3/5},\\
 q_-&=\sqrt D\,\Xi_H^{3/10},&
 q_+&=\sqrt D\,\Xi_H^{-3/10}.
 \end{aligned}
 \label{amp-higgs-canales}
\end{equation}
Las potencias son las raíces reales positivas. El dominio exacto
de la cámara orientada es
\begin{equation}
 0<D<1,\qquad 1<\Xi_H<D^{-5/3}.
 \label{amp-higgs-dominio-Xi}
\end{equation}
En efecto, $0<y<x$ equivale a $1<e^{10y/3}<e^{10x/3}$.
El extremo inferior hace $q_+=q_-$; el superior hace $q_-=1$.
Para las dos orientaciones juntas, $|y|<x$ equivale a
$D^{5/3}<\Xi_H<D^{-5/3}$.

Sea $F(q)=f(q^{30})=(1-q^{90})/(1-q^{120})$, la respuesta ya probada en
\eqref{iii:eq:rchannels}--\eqref{iii:eq:monotone}.
\begin{corollary}[Confluencia de las lecturas escalar y constitutiva]
\label{amp-higgs-impedancia-corolario}
En el dominio~\eqref{amp-higgs-dominio-Xi},
\begin{equation}
 \boxed{\widehat Z=
 \frac{F(\sqrt D\,\Xi_H^{3/10})}
 {F(\sqrt D\,\Xi_H^{-3/10})}.}
 \label{amp-higgs-impedancia}
\end{equation}
Las respuestas completas son
\begin{equation}
 \begin{aligned}
 r_-&=F(\sqrt D\,\Xi_H^{3/10}),&
 r_+&=F(\sqrt D\,\Xi_H^{-3/10}),\\
 \widehat\mu&=r_-^2,& \widehat\varepsilon&=r_+^2,&
 \widehat c&=(r_-r_+)^{-1}.
 \end{aligned}
 \label{amp-higgs-vacio}
\end{equation}
Además, las dos reconstrucciones satisfacen
\begin{equation}
 \frac{F^{-1}(\sqrt{\widehat\mu})}
 {F^{-1}(\sqrt{\widehat\varepsilon})}
 =\rho_{\rm or}=\Xi_H^{3/5}.
 \label{amp-higgs-confluencia}
\end{equation}
\end{corollary}
\begin{proof}
Sustituir los canales de~\eqref{amp-higgs-canales} en las cuatro lecturas
de~\eqref{iii:eq:four} prueba las primeras identidades. La inversa
constitutiva de~\eqref{iii:eq:cubic} recupera $q_-$ y $q_+$; su cociente
es $\rho_{\rm or}$. La igualdad con $\Xi_H^{3/5}$ ya está probada en
\eqref{amp-higgs-canales}.
\end{proof}

En la sección compatible con $A_{\deg}=1000\alpha$, el acoplamiento
$Q(D):=-9\log D/25$ satisface $Q(D)=4\pi\alpha$ y
$g^2=Q(D)/(1-D)$. Por tanto, la misma coordenada de Higgs se escribe
\begin{equation}
 \Xi_H=\frac{8(1-D)D^3\lambda_H}{Q(D)},\qquad
 \frac1{g^2}+\frac1{g'^2}=\frac1{Q(D)}.
 \label{amp-higgs-curva}
\end{equation}
Estas expresiones relacionan el acoplamiento total, su reparto racionalizado
y la orientación escalar en la misma sección.

\subsection{Información conservada y transporte de orientación}

La pareja de calibre recupera $\alpha$ y $x$; el cociente
$\lambda_H/g^2$ recupera después $y$. Con las marcas $P_\pm$ se reconstruyen
los operadores completos $T$ y $\mathcal V=r_+P_++r_-P_-$.
Los tres escalares determinan los autovalores; los proyectores marcados
determinan además su realización operatoria. La historia enriquecida de la
ruta sigue siendo una información anterior de mayor resolución.

Bajo $y\mapsto-y$ respecto de las marcas fijas, $D,g,g'$ permanecen
invariantes y $\Xi_H\mapsto\Xi_H^{-1}$. Se intercambian las constitutivas,
$\widehat Z\mapsto\widehat Z^{-1}$ y $\widehat c$ se conserva.
Si $\lambda_H^+$ y $\lambda_H^-$ denotan las evaluaciones en las dos
orientaciones conjugadas, se cumple
\begin{equation}
 \lambda_H^+\lambda_H^-
 =\frac{g^4}{64}D^{-6}.
 \label{amp-higgs-orientacion}
\end{equation}
La prueba consiste en multiplicar~\eqref{amp-higgs-lambda} para $y$ y $-y$;
las contribuciones orientadas se cancelan. Esta comparación conserva el
lector de firmas y cambia su argumento angular. Una conjugación simultánea
de rutas y referencias debe transportar también las firmas.

En la amplificación $T_m=T\otimes I_{9^m}$, con
$\tau_m=\operatorname{Tr}/(2\cdot9^m)$ y
$\mathcal R_m=\mathcal R\otimes I_{9^m}$, las coordenadas adecuadas son
\begin{equation}
 D=\exp\bigl(2\tau_m(\log T_m)\bigr),\qquad
 \rho_{\rm or}=\exp\bigl(-2\tau_m(\mathcal R_m\log T_m)\bigr).
 \label{amp-higgs-refinamiento}
\end{equation}
En efecto, $\log T_m=(\log T)\otimes I_{9^m}$, de modo que las dos trazas
son $-x$ y $-y$, respectivamente. La normalización elimina la multiplicidad
de la amplificación. Las lecturas de calibre, Higgs y vacío permanecen así
compatibles con estos refinamientos; el determinante ordinario ampliado
sería $D^{9^m}$ y tiene una función distinta.

La conclusión de la composición es precisa: el invariante par de calibre
y la coordenada orientada escalar reconstruyen conjuntamente la pareja
espectral que gobierna el vacío. Cuando $\lambda_H$ se evalúa mediante la
misma fórmula, la confluencia es una identidad de consistencia. Una lectura
independiente de $\lambda_H$, transportada a idénticos esquema y escala,
convierte esa igualdad en un contraste adicional. En ambos usos permanecen
explícitos las firmas, las marcas de hoja y la normalización de árbol.

\section{Descomposición cíclica, cuarto de giro y momento orientado de Catalán}
\label{iii:sec:catalan}

El registro dodecafásico permite relacionar la respuesta constitutiva con
otros invariantes del transporte. El vínculo se establece sobre subespacios
y aplicaciones determinados: el defecto de rango mide una diferencia de
incidencia, el determinante conserva una respuesta escalar y un coeficiente
del resolvente retiene la orientación. La colección de momentos que culmina
en Catalán interviene a este último nivel.

\subsection{Dos subespacios del ciclo dodecafásico}

En $\mathbb R^{12}$, sea $C e_j=e_{j+1\bmod12}$ y sean
\begin{equation}
 P_3^{\mathrm{cyc}}=\frac{I+C^3+C^6+C^9}{4},\qquad
 P_4^{\mathrm{cyc}}=\frac{I+C^4+C^8}{3},\qquad
 H_d=\operatorname{Ran}P_d^{\mathrm{cyc}}.
 \label{iii:eq:projectors34}
\end{equation}
Los promedios sobre los subgrupos cíclicos son proyectores ortogonales.
$H_3$ está formado por las secuencias periódicas de período divisor de $3$,
y $H_4$ por las de período divisor de $4$. Sus dimensiones son $3$ y $4$.
La restricción de $C$ a cada espacio realiza el ciclo correspondiente.
El superíndice $\mathrm{cyc}$ identifica el promedio del ciclo dodecafásico:
$P_3^{\mathrm{cyc}}$ es el proyector sobre las secuencias de período divisor
de $3$. Su dominio y su acción son distintos de los del proyector isótipo
tridimensional de la realización icosaédrica; la coincidencia de rango
no identifica ambos operadores.

\begin{proposition}[Defecto de rango y respuesta determinantal]
Se verifican
\begin{equation}
 \frac{\operatorname{Tr}(P_3^{\mathrm{cyc}}-P_4^{\mathrm{cyc}})}{12}=-\frac1{12},\qquad
 \frac{\det_{H_3}(I-sC)}{\det_{H_4}(I-sC)}
 =\frac{1-s^3}{1-s^4}=f(s),\quad 0<s<1.
 \label{iii:eq:trace-det}
\end{equation}
\end{proposition}
\begin{proof}
La traza de un proyector ortogonal es su rango, de modo que el numerador
de la primera expresión es $3-4$. El polinomio característico de un ciclo
de longitud $d$ es $t^d-1$; por ello su determinante $\det(I-sC)$ es
$1-s^d$. El cociente produce la segunda identidad.
\end{proof}

La sustitución $s=q^{30}$ transporta esta representación al cociente
$90/120$. En $s=1$, el modo constante común se cancela antes de tomar el
límite, que vale $3/4$. Esta realización de $-1/12$ como defecto de traza
tiene un dominio finito explícito. Las realizaciones regularizadas y
modulares del mismo escalar conservan sus operadores y dominios propios.

\subsection{Representación fiel del cuarto de giro}

Definamos el proyector antipodal $P_{\rm ant}=(I-C^6)/2$ y
$Q=P_4^{\mathrm{cyc}}P_{\rm ant}$. Este proyector actúa sobre el ciclo de doce fases y es
distinto del proyector constitutivo $P_-$ de la sección~\ref{iii:sec:constitutiva}.
Los proyectores $P_4^{\mathrm{cyc}}$ y $P_{\rm ant}$ conmutan y
$Q$ tiene rango $2$. En su imagen se fija la base ortonormal
\begin{equation}
 u=\frac1{\sqrt6}(1,0,-1,0,1,0,-1,0,1,0,-1,0)^{\mathsf T},
 \qquad v=Cu.
 \label{iii:eq:uv}
\end{equation}
Se comprueba componente a componente que $Cu=v$ y $Cv=-u$. En consecuencia,
la isometría $W(a,b)=au+bv$ satisface
\begin{equation}
 CW=WJ,\qquad
 J=\begin{pmatrix}0&-1\\1&0\end{pmatrix},\qquad J^2=-I.
 \label{iii:eq:quarter-intertwiner}
\end{equation}
La estructura compleja del plano queda realizada por la acción efectiva del
ciclo. El signo depende de la orientación fijada: en la misma base,
$C^3W=-WJ$. Esta precisión permite componer el plano con el lector orientado.

\subsection{Resolvente y profundidad ilimitada}

Como $C^2=-I$ sobre $\operatorname{Ran}Q$, para $0\leq s<1$ se tiene
\begin{equation}
 (I-sC)^{-1}\big|_{\operatorname{Ran}Q}=\frac{I+sC}{1+s^2},
 \qquad k_4(s):=\langle v,(I-sC)^{-1}u\rangle=\frac{s}{1+s^2}.
 \label{iii:eq:resolvent}
\end{equation}
La sucesión $\langle v,C^nu\rangle$ es $0,1,0,-1,\ldots$. El lector
de profundidad conserva este carácter mediante el momento
\begin{equation}
 \mathcal C_2(s)=\sum_{k\geq0}\frac{(-1)^ks^{2k+1}}{(2k+1)^2},
 \qquad 0\leq s\leq1.
 \label{iii:eq:c2}
\end{equation}
La serie converge absolutamente y uniformemente en ese intervalo por
comparación con $\sum(2k+1)^{-2}$. Para $0<s<1$, las series derivadas
convergen uniformemente sobre cada intervalo compacto interior. Con
$\mathcal D=s\,d/ds$, se obtiene
\begin{equation}
 \mathcal D^2\mathcal C_2(s)=k_4(s),\qquad
 \mathcal C_2(s)=\int_0^s\frac{dt}{t}\int_0^t\frac{d\xi}{\xi}\,k_4(\xi).
 \label{iii:eq:c2-integral}
\end{equation}
Las condiciones $\mathcal C_2(0)=0$ y
$\lim_{s\downarrow0}\mathcal D\mathcal C_2(s)=0$ fijan las dos constantes
de integración. En efecto, la diferencia de dos soluciones de
$\mathcal D^2F=k_4$ tiene la forma $a\log s+b$; la segunda condición
anula $a$ y la primera anula $b$.
La convergencia uniforme permite evaluar el extremo $s=1$: el resultado es
el momento de Catalán. La representación finita del carácter conserva su
orientación; el índice de profundidad permanece ilimitado.

\subsection{Composición con la respuesta constitutiva}

\begin{theorem}[Relación diferencial constitutiva]
Para $0<s<1$, con la orientación de \eqref{iii:eq:uv},
\begin{equation}
 f(s)=\frac{1+s+s^2}{s(1+s)}\,\mathcal D^2\mathcal C_2(s).
 \label{iii:eq:catalan-vacuum}
\end{equation}
\end{theorem}
\begin{proof}
Sustituir \eqref{iii:eq:resolvent} en \eqref{iii:eq:c2-integral} y utilizar
$(1+s)(1+s^2)=1+s+s^2+s^3$ da la identidad. Todos los denominadores son
positivos en el dominio. El entrelazador \eqref{iii:eq:quarter-intertwiner}
identifica el coeficiente operatorio utilizado antes de su evaluación.
\end{proof}

Al evaluar $s_\pm=q_\pm^{30}$, la misma colección produce las dos respuestas
de canal de \eqref{iii:eq:rchannels}. El límite de Catalán, el defecto de
traza y los autovalores constitutivos mantienen sus funciones diferentes
dentro de esta composición. La inversión del ciclo con las marcas $u,v$
fijas cambia el signo de $k_4$ y conserva el determinante $1+s^2$.
Así se identifica una información de orientación que el determinante
aislado no distingue.

% Delta de §10.27. Insertar al final de 05_ciclo_catalan.tex.
% La actualización afín y su balance permanecen en sus propietarios originales.
\subsection{Trazas de retorno y momentos del cociente constitutivo}
\label{amp-afin-sec}

La actualización del registro y la respuesta del ciclo comparten el
transporte dodecafásico y conservan información de distinto tipo.
El evento firmado se obtiene de la dirección y la orientación de cada
ventana de la historia APP--TRIT--TPK. Su incremento vectorial y el
balance~\eqref{mem:balance} retienen la interacción con el registro previo.
La construcción siguiente prolonga la lectura espectral del mismo ciclo,
una vez fijados sus subespacios de incidencia. Este orden permite componer
retorno, determinante y momentos conservando la historia que los precede.

Conservamos el operador $C$ y los proyectores de
\eqref{iii:eq:projectors34}, y definimos
\begin{equation}
 E_{\rm cyc}=P_3^{\mathrm{cyc}}-P_4^{\mathrm{cyc}},\qquad
 b_n=\operatorname{Tr}(C^n E_{\rm cyc}),\qquad n\geq0.
 \label{amp-afin-trazas}
\end{equation}
El símbolo $b_n$ distingue esta traza espectral del evento firmado
$a_m$ de una ventana. La traza utiliza el operador de transporte; el
evento utiliza además la trayectoria efectivamente recorrida.

\begin{proposition}[Tabla de trazas producida por la incidencia]
\label{amp-afin-prop-tabla}
Se cumplen
\begin{equation}
 \begin{gathered}
 b_n=3\mathbf1_{3\mid n}-4\mathbf1_{4\mid n},\\
 (b_1,\ldots,b_{12})=(0,0,3,-4,0,3,0,-4,3,0,0,-1).
 \end{gathered}
 \label{amp-afin-tabla}
\end{equation}
El período mínimo es doce y la suma en un período es cero. Para
$B_N=\sum_{n=1}^{N}b_n$,
\begin{equation}
 B_N=3\lfloor N/3\rfloor-4\lfloor N/4\rfloor,\qquad |B_N|\leq3.
 \label{amp-afin-parciales}
\end{equation}
\end{proposition}
\begin{proof}
En $H_d$, el operador $C$ permuta cíclicamente las $d$ clases de
coordenadas. Su potencia $n$ tiene $d$ puntos fijos si $d$ divide a $n$
y ninguno en caso contrario. Como $P_d^{\mathrm{cyc}}$ conmuta con $C$,
$\operatorname{Tr}(C^nP_d^{\mathrm{cyc}})$ es la traza de esa restricción.
Restar los casos $d=3,4$ prueba la fórmula. El mínimo período positivo
es doce: $b_n=-1$ ocurre exactamente en los múltiplos de doce.
En un período las contribuciones son $3\cdot4-4\cdot3=0$.
La suma finita cuenta los múltiplos y da~\eqref{amp-afin-parciales}.
Sus valores para $N=0,\ldots,11$ son
$(0,0,0,3,-1,-1,2,2,-2,1,1,1)$ y se repiten al sumar doce.
\end{proof}

\begin{theorem}[Reconstrucción logarítmica del lector constitutivo]
\label{amp-afin-thm-logdet}
Para $|s|<1$, el cociente $f$ de~\eqref{iii:eq:trace-det} verifica
\begin{align}
 -\log f(s)&=\sum_{n\geq1}\frac{b_n}{n}s^n,\label{amp-afin-logdet}\\
 -\frac{s f'(s)}{f(s)}
 &=\frac{3s^3}{1-s^3}-\frac{4s^4}{1-s^4}.
 \label{amp-afin-logder}
\end{align}
La rama holomorfa queda fijada por $\log f(0)=0$.
En $0<s<1$ coincide con el logaritmo real.
\end{theorem}
\begin{proof}
Los polinomios $1-s^3$ y $1-s^4$ no se anulan en el disco abierto.
La expansión $-\log(1-w)=\sum_{k\geq1}w^k/k$, obtenida integrando
la serie geométrica desde cero, converge uniformemente en discos
compactos. Con $w=s^3,s^4$, su diferencia da
$-\log f(s)=\sum_{k\geq1}s^{3k}/k-\sum_{k\geq1}s^{4k}/k$.
El coeficiente de $s^n$ es $b_n/n$.
La convergencia local uniforme permite derivar y sumar las dos series
geométricas, lo que prueba~\eqref{amp-afin-logder}.
La condición en cero fija la constante de integración y la rama.
\end{proof}

Así, $b_0/12=-1/12$ registra el defecto normalizado de dimensión,
mientras toda la sucesión $b_n$ reconstruye una función.
Son dos lecturas de distinto alcance del mismo par de subespacios.
Al sustituir $s=q_\pm^{30}$, los argumentos siguen siendo las salidas
de la construcción angular previa; la expansión no los selecciona.

\begin{theorem}[Momentos de profundidad y representación de Mellin]
\label{amp-afin-thm-momentos}
La serie
\begin{equation}
 M_b(z)=\sum_{n\geq1}\frac{b_n}{n^z}
 \label{amp-afin-momento}
\end{equation}
converge localmente de manera uniforme para $\operatorname{Re}z>0$.
Si $\sigma=\operatorname{Re}z>0$, su resto satisface
\begin{equation}
 \left|M_b(z)-\sum_{n=1}^{N}\frac{b_n}{n^z}\right|
 \leq3\left(1+\frac{|z|}{\sigma}\right)(N+1)^{-\sigma}.
 \label{amp-afin-cota}
\end{equation}
En el dominio de convergencia absoluta $\operatorname{Re}z>1$,
\begin{equation}
 M_b(z)=(3^{1-z}-4^{1-z})\zeta(z),\qquad
 \zeta(z)=\sum_{n\geq1}n^{-z}.
 \label{amp-afin-reconocimiento}
\end{equation}
Para $\operatorname{Re}z>0$ se tiene además
\begin{equation}
 \Gamma(z)M_b(z)=
 \int_0^\infty t^{z-1}
 \left(\frac3{e^{3t}-1}-\frac4{e^{4t}-1}\right)dt,
 \label{amp-afin-mellin}
\end{equation}
donde $\Gamma(z)=\int_0^\infty u^{z-1}e^{-u}\,du$.
\end{theorem}
\begin{proof}
La sumación por partes y $|B_n|\leq3$ dan, al pasar al límite superior,
\[
 \sum_{n>N}b_n n^{-z}
 =-B_N(N+1)^{-z}
   +\sum_{n>N}B_n\bigl(n^{-z}-(n+1)^{-z}\bigr).
\]
La diferencia de potencias está acotada en módulo por
$|z|\int_n^{n+1}x^{-\sigma-1}dx$. Esto prueba
\eqref{amp-afin-cota}, la convergencia y su uniformidad en compactos.
Para $\sigma>1$, la convergencia absoluta permite agrupar
los múltiplos de tres y de cuatro por separado, dando
\eqref{amp-afin-reconocimiento}.

La serie geométrica del transporte, evaluada en $s=e^{-t}$, produce
\[
 \sum_{n\geq1}b_n e^{-nt}
 =\frac3{e^{3t}-1}-\frac4{e^{4t}-1}=:K_b(t).
\]
Para $\sigma>1$, se puede integrar término a término porque
$\sum|b_n|n^{-\sigma}<\infty$. La sustitución $u=nt$ da
$\Gamma(z)n^{-z}$ y prueba~\eqref{amp-afin-mellin} en ese semiplano.
En el origen se cancelan los términos $1/t$:
$K_b(t)=1/2+O(t)$. En el infinito, $K_b(t)=O(e^{-3t})$.
La integral es, por ello, holomorfa para $\operatorname{Re}z>0$
por convergencia local uniforme. La serie también lo es por la cota
anterior; el principio de identidad prolonga la igualdad a ese dominio.
\end{proof}

En $0<\operatorname{Re}z\leq1$, el orden creciente de profundidad
fija la lectura convergente de la serie. La realización como traza
absoluta del operador diagonal corresponde a $\operatorname{Re}z>1$.

\begin{corollary}[Evaluaciones compatibles con el extremo constitutivo]
\label{amp-afin-cor-extremos}
Se cumplen
\begin{equation}
 M_b(2)=\frac{\zeta(2)}{12},\qquad
 M_b(1)=\log(4/3)=-\log f(1),\qquad f(1)=\frac34.
 \label{amp-afin-extremos}
\end{equation}
\end{corollary}
\begin{proof}
La primera igualdad resulta de $3^{-1}-4^{-1}=1/12$.
Para la segunda, los parciales exactos son
$H_{\lfloor N/3\rfloor}-H_{\lfloor N/4\rfloor}$, con
$H_j=\sum_{k=1}^j1/k$ y $H_0=0$.
La comparación integral de $1/x$ entre esos dos índices da el límite
$\log(4/3)$. Cancelar $1-s$ antes de evaluar el cociente determinantal
da $f(1)=3/4$.
\end{proof}

El registro afín retiene los eventos y su interacción con la memoria
previa; la tabla de trazas retiene la incidencia de los dos subespacios
a todas las potencias del transporte; el momento aplica una ponderación
de profundidad a esa tabla ya producida. Su periodicidad espectral
convive con una historia de eventos que continúa después del retorno
de fase. Invertir $C$ conserva $b_n$, porque $d\mid n$ es invariante
bajo $n\mapsto-n$; el coeficiente orientado del resolvente de
\eqref{iii:eq:resolvent}, con sus marcas fijas, cambia de signo.
La lectura conjunta conserva una distinción que el determinante
escalar por sí solo no registra. La forma cuadrática y la suma de
coordenadas del registro mantienen sus tipos matemáticos; esta
composición no les asigna unidades de energía ni de carga eléctrica.

\subsection{Composición residual y producto de momentos}
\label{amp-afin-sec-productos}

La composición posterior de lectores conserva también la operación
efectuada sobre la profundidad. En la realización
$\mathcal H=\ell^2(\mathbb N_{>0})$, fijamos
$N\delta_n=n\delta_n$. Sean $a,c$ dos tablas residuales periódicas
ya producidas y $Q_a\delta_n=a(n)\delta_n$,
$Q_c\delta_n=c(n)\delta_n$. Al prolongar ambas al período común,
los dos lectores actúan sobre el mismo índice. Para
$\operatorname{Re}z>1$, escribimos
$M_a(z)=\operatorname{Tr}(N^{-z}Q_a)$ y análogamente $M_c$.

\begin{proposition}[Composición diagonal y composición tensorial]
\label{amp-afin-prop-productos}
Se tienen las identidades
\begin{align}
 \operatorname{Tr}(N^{-z}Q_aQ_c)
 &=\sum_{n\geq1}\frac{a(n)c(n)}{n^z},
 \label{amp-afin-diagonal}\\
 M_a(z)M_c(z)
 &=\operatorname{Tr}_{\mathcal H\otimes\mathcal H}
 \bigl((N\otimes N)^{-z}(Q_a\otimes Q_c)\bigr)\notag\\
 &=\sum_{k\geq1}\frac{(a*c)(k)}{k^z},
 \label{amp-afin-tensor}\\
 (a*c)(k)&=\sum_{d\mid k}a(d)c(k/d),
 \qquad \operatorname{Re}z>1.
 \label{amp-afin-convolucion}
\end{align}
Todos los operadores que figuran dentro de estas trazas son de clase traza.
\end{proposition}
\begin{proof}
El producto diagonal actúa sobre $\delta_n$ mediante
$a(n)c(n)$, lo que prueba~\eqref{amp-afin-diagonal}.
En $\delta_m\otimes\delta_n$, el operador $N\otimes N$
tiene valor propio $mn$ y $Q_a\otimes Q_c$ actúa mediante
$a(m)c(n)$. Si $\sigma=\operatorname{Re}z>1$, la suma de los
módulos de estos coeficientes espectrales está acotada por
$\|a\|_\infty\|c\|_\infty(\sum n^{-\sigma})^2<\infty$.
Esta cota prueba la condición de clase traza y permite factorizar
la suma doble o agruparla por $k=mn$. Los pares con producto $k$
son exactamente $(d,k/d)$, con $d\mid k$.
La cota análoga con una sola suma prueba el caso diagonal.
\end{proof}

En los lectores de esta sección, tomar $a(n)=b_n$ y
$c(n)=\langle v,C^nu\rangle$ da, para $n=3$,
$a(3)c(3)=-3$, mientras $(a*c)(3)=3$.
Esta diferencia de coeficientes exhibe la información que conserva
cada operación. La composición diagonal mantiene una profundidad común;
la tensorial conserva dos profundidades antes de agrupar por su producto.
La realización tensorial no atribuye independencia estadística a las
historias HMT que produjeron los lectores. Las derivadas regularizadas y
los cambios de normalización de otros momentos constituyen operaciones
posteriores distintas de las dos composiciones aquí demostradas.

\subsection{Factorización a través del operador constitutivo}

Conservemos la coordenatización angular levantada de la
sección~\ref{iii:sec:hojas}. Para hacer explícita la conversión angular,
escribamos
\begin{equation}
 A_{\deg}=\frac{180x}{\pi},\qquad
 C^*_{\deg}=\frac{180y}{\pi},\qquad
 q_\pm=e^{-x\mp y},\qquad s_\pm=q_\pm^{30}.
 \label{iii:eq:barbero-angular-lift}
\end{equation}
La orientación considerada en \eqref{iii:eq:rchannels} corresponde a
$0<y<x$; el intercambio de hojas se expresa por $y\mapsto-y$ en el
dominio más amplio $x>|y|$. Las exponenciales reales actúan sobre estos
levantamientos, que conservan los valores de $x\pm y$ durante la
composición.

El funcional orientado de Barbero--Immirzi es
\begin{equation}
 \gamma=\frac{\pi A_{\deg}C^*_{\deg}}{180}
   +12\bigl[S_{90}(q_-)-S_{90}(q_+)\bigr]
   -\bigl[S_{120}(q_-)-S_{120}(q_+)\bigr].
 \label{iii:eq:barbero-functional}
\end{equation}
La primera componente es la evaluación angular; la segunda es la
diferencia del lector de la cadena fundamental en los canales conjugados.
Los argumentos se reciben del transporte angular previamente construido.
La fórmula determina un escalar adimensional sobre esa estructura de
partida.

Sea $\psi:(3/4,1)\to(0,1)$ la inversa de $f$ construida mediante
\eqref{iii:eq:cubic}. Definamos
\begin{equation}
 \mathcal H(s)=\frac{12s^3}{1-s^9}
             -\frac{s^4}{1-s^{12}}
             -\frac{(\log s)^2}{20\pi},\qquad 0<s<1.
 \label{iii:eq:barbero-h}
\end{equation}
Esta función es real analítica en su dominio. Los exponentes $3,4,9,12$
resultan de expresar las alturas $90,120$ y los retornos $270,360$ en
unidades de $30$; el coeficiente de la parte logarítmica se fija al
transportar la contribución angular.

\begin{theorem}[Expresión constitutiva orientada de Barbero--Immirzi]
\label{iii:thm:barbero-factorization}
En la coordenatización \eqref{iii:eq:barbero-angular-lift},
\begin{equation}
 \gamma=\mathcal H(s_-)-\mathcal H(s_+)
       =\mathcal H\bigl(\psi(r_-)\bigr)
        -\mathcal H\bigl(\psi(r_+)\bigr).
 \label{iii:eq:barbero-factorization}
\end{equation}
Sobre la representación bidimensional de canales,
\begin{equation}
 \gamma=-\operatorname{Tr}
    \left[\mathcal R\,\mathcal H\bigl(\psi(\mathcal V)\bigr)\right],
 \label{iii:eq:barbero-trace}
\end{equation}
donde $\operatorname{Tr}$ es la traza ordinaria, con
$\operatorname{Tr}I=2$.
\end{theorem}
\begin{proof}
La sustitución $s=q^{30}$ en \eqref{iii:eq:barbero-return-series} da
\[
 S_{90}(q)=\frac{s^3}{1-s^9},\qquad
 S_{120}(q)=\frac{s^4}{1-s^{12}}.
\]
Además, \eqref{iii:eq:barbero-angular-lift} implica
\[
 \log s_\pm=-30(x\pm y)
 =-\frac\pi6(A_{\deg}\pm C^*_{\deg}).
\]
Por diferencia de cuadrados,
\begin{align}
 \frac{(\log s_+)^2-(\log s_-)^2}{20\pi}
 &=\frac{900\bigl[(x+y)^2-(x-y)^2\bigr]}{20\pi}\\
 &=\frac{180xy}{\pi}
 =\frac{\pi A_{\deg}C^*_{\deg}}{180}.
 \label{iii:eq:barbero-log-term}
\end{align}
La suma de estas identidades prueba la primera igualdad de
\eqref{iii:eq:barbero-factorization}. La segunda utiliza la inversión
$\psi(f(s_\pm))=s_\pm$ demostrada en la
sección~\ref{iii:sec:constitutiva}.

Para la expresión operatoria, los proyectores de hoja tienen rango uno y
$\mathcal RP_\pm=\pm P_\pm$. El cálculo espectral proporciona
\[
 \mathcal H\bigl(\psi(\mathcal V)\bigr)
 =\mathcal H(s_+)P_++\mathcal H(s_-)P_-.
\]
Multiplicar por $\mathcal R$ y tomar la traza ordinaria da
$\mathcal H(s_+)-\mathcal H(s_-)$. Su opuesto es $\gamma$.
\end{proof}

La marca de orientación $\mathcal R$ forma parte del funcional. Si se
mantiene fija esa marca y se intercambian los canales de $\mathcal V$,
cambia el signo de $\gamma$. La evaluación del espectro sin sus etiquetas conserva
el par no ordenado, mientras \eqref{iii:eq:barbero-trace} conserva la
contribución orientada. Si se emplea la traza normalizada
$\tau_2=\operatorname{Tr}/2$, la misma fórmula se escribe
$\gamma=-2\tau_2[\mathcal R\mathcal H(\psi(\mathcal V))]$.

En cambio, un cambio simultáneo de representación
$(\mathcal R,\mathcal V)\mapsto
(U\mathcal RU^{-1},U\mathcal VU^{-1})$ conserva el funcional:
el cálculo espectral se transporta por conjugación y la traza es
invariante. Intercambiar la respuesta respecto de una orientación fija
y transportar conjuntamente respuesta y orientación son, por tanto,
operaciones diferentes.

Esta normalización es compatible con las amplificaciones de
\eqref{iii:eq:logvac}. Para
$\mathcal V_m=\mathcal V\otimes I_{9^m}$ y
$\mathcal R_m=\mathcal R\otimes I_{9^m}$,
\begin{equation}
 \gamma=-\frac1{9^m}\operatorname{Tr}
 \left[\mathcal R_m\mathcal H\bigl(\psi(\mathcal V_m)\bigr)\right].
 \label{iii:eq:barbero-amplification}
\end{equation}
En efecto, el cálculo funcional conserva el producto tensorial y la traza
de $I_{9^m}$ es $9^m$. La multiplicidad de la representación se separa
de la evaluación orientada.

\subsection{Reconstrucción a partir de las coordenadas del vacío}

Las identidades \eqref{iii:eq:four-identities} permiten expresar el mismo
funcional mediante la impedancia y la velocidad normalizadas:
\begin{align}
 \gamma={}&
 \mathcal H\!\left(\psi\!\left(
     \sqrt{\frac{\widehat Z}{\widehat c}}\right)\right)
 -\mathcal H\!\left(\psi\!\left(
     \frac1{\sqrt{\widehat Z\widehat c}}\right)\right).
 \label{iii:eq:barbero-zc}
\end{align}
El dominio positivo de esta reconstrucción se describe por
\begin{equation}
 1<\widehat Z\widehat c<\frac{16}{9},\qquad
 \frac9{16}<\frac{\widehat Z}{\widehat c}<1.
 \label{iii:eq:barbero-zc-domain}
\end{equation}
En efecto, esas desigualdades equivalen a
$3/4<r_\pm<1$ al sustituir las raíces positivas de
\eqref{iii:eq:four-identities}. Para la orientación $y>0$ se conserva
además $r_-<r_+$, equivalentemente $\widehat Z<1$. Los puntos del
dominio se utilizan con la compatibilidad angular ya establecida por
el TPK; la inversión es una reconstrucción posterior de sus canales.
Las coordenadas son las internas de \eqref{iii:eq:four}: un cambio de
unidades debe deshacerse mediante \eqref{iii:eq:undo-units} antes de
aplicar $\psi$.

El intercambio $y\mapsto-y$ actúa como
$(\widehat Z,\widehat c)\mapsto(\widehat Z^{-1},\widehat c)$ y
produce $\gamma\mapsto-\gamma$. En el caso balanceado $y=0$,
ambos canales coinciden y \eqref{iii:eq:barbero-factorization} da
$\gamma=0$. La velocidad normalizada conserva el producto de los
canales; su lectura conjunta con la impedancia recupera las dos
respuestas y la orientación necesaria para evaluar el funcional.
En ese caso balanceado $\mathcal V=rI$ tiene espectro degenerado:
la marca $\mathcal R$ se conserva como dato de representación, aunque
el escalar $r$ ya no distinga sus proyectores. Para espectro simple,
en cambio,
$P_+=(\mathcal V-r_-I)/(r_+-r_-)$ y $P_-=I-P_+$.

\section{Evaluación constitutiva y realización dimensional}
\label{iii:sec:evaluacion}

La evaluación reúne ahora las construcciones anteriores en su orden efectivo.
Las expansiones regionales de $\pi$, $\varphi$ y $e$ proceden del refinamiento
nonádico; el registro $K$ conserva sus doce posiciones. Su lectura periódica
determina la sección posicional de $\alpha$. Sobre ella actúan el carácter
de acción y el retorno regional, después el par angular y, finalmente, la
respuesta constitutiva. Esta evaluación es posterior a la selección de las
reglas: ninguna de sus cifras se utiliza para elegir un canal.

\subsection{Evaluación de la cadena interna}

Para reproducir la tabla se emplean las expansiones regionales de mil
decimales generadas por el certificado nonádico y el entero $N_K$ de
\eqref{eq:k-numerador}. La lectura utilizada es
\[
 \alpha=\pi+e-\varphi-4-\frac{N_K}{1000^{12}-1}.
\]
La sección finita de doce bloques y la sección prolongada tienen sus
definiciones y su diferencia exacta en el capítulo del registro. Aquí se
mantiene fija la segunda durante toda la evaluación. No se intercambian
lecturas al aumentar la precisión, ni se sustituye el operador analítico
de la segunda vía por una coincidencia decimal con la primera.

\begin{longtable}{@{}ll@{}}
\toprule
Objeto & Evaluación decimal (24 cifras significativas)\\
\midrule\endhead
$\alpha$ & $0.00729735256928380099728511\ldots$ \\
$H_5$ & $1.05457181764615446962582\ldots$ \\
$\eta_{\rm ret}$ & $1.05457181691503906113369\ldots$ \\
$A_{\deg}$ & $7.29735256928380099728511\ldots$ \\
$C^*_{\deg}$ & $2.12373833896864572426195\ldots$ \\
$q_+$ & $0.848377942576404374945737\ldots$ \\
$q_-$ & $0.913660151150557285295519\ldots$ \\
$r_+$ & $0.999999628548249363178695\ldots$ \\
$r_-$ & $0.999724137189518364131517\ldots$ \\
$\widehat\varepsilon$ & $0.999999257096636702760442\ldots$ \\
$\widehat\mu$ & $0.999448350479326935089982\ldots$ \\
$\widehat Z$ & $0.999724508538937215455554\ldots$ \\
$\widehat c$ & $1.00027631048615752610639\ldots$ \\
$\gamma$ & $0.274007672694459274747255\ldots$ \\
$G_{\rm Cat}$ & $0.915965594177219015054604\ldots$ \\
$\rho$ & $0.267949192431122706472554\ldots$ \\
\bottomrule
\end{longtable}


El cálculo se realiza con 120 cifras de trabajo; se muestran 24 cifras
significativas para facilitar la comparación entre magnitudes de distintos
órdenes. Las 120 cifras son una precisión de evaluación, no una incertidumbre
experimental ni, por sí solas, un certificado de error mediante intervalos.
El paquete conserva las entradas generadas, sus huellas, las operaciones y
los residuos de las identidades comprobadas. Las pruebas algebraicas del
cuerpo establecen esas identidades con independencia de esta tabla.

El significado de las cuatro últimas coordenadas constitutivas se conserva
al evaluar: $\widehat\varepsilon$ y $\widehat\mu$ son las respuestas
cuadráticas por hoja, $\widehat Z$ es su orientación relativa y
$\widehat c$ es el inverso de su producto. La cercanía de algunas cifras a
la unidad expresa el pequeño defecto del cociente $3$--$4$ sobre estos
canales. Su origen se ve exactamente en
\[
 1-f(s)=\frac{s^3}{1+s+s^2+s^3}.
\]
La fórmula permite estudiar la respuesta antes de atribuirle una magnitud
dimensional: el signo del defecto, su orden cúbico y su monotonía son
consecuencias del mismo operador.

\subsection{Coordenatizaciones de unidades y covariancia}

Sean $\mathcal L_S$, $\mathcal L_Q$ y $\mathcal L_C$ las rectas de
acción, carga y velocidad. Una coordenatización física positiva evalúa las secciones
$u_S$, $u_Q$ y $u_C$ en unidades de acción, carga y velocidad. La realización
\eqref{iii:eq:dimensions} produce entonces una impedancia, una velocidad,
una permeabilidad y una permitividad, con sus dimensiones respectivas.
Las coordenadas de la tabla anterior permanecen adimensionales.

\begin{proposition}[Covariancia de la realización constitutiva]
Si las secciones de referencia cambian mediante
$u_S'=a u_S$, $u_Q'=b u_Q$, $u_C'=d u_C$, con $a,b,d>0$, las
coordenadas que representan las mismas magnitudes satisfacen
\[
 \widehat Z'=a^{-1}b^2\widehat Z,\qquad
 \widehat c'=d^{-1}\widehat c,\qquad
 \widehat\mu'=a^{-1}db^2\widehat\mu,\qquad
 \widehat\varepsilon'=adb^{-2}\widehat\varepsilon.
\]
Las relaciones $\widehat\mu'\widehat\varepsilon'=(\widehat c')^{-2}$
y $\widehat\mu'/\widehat\varepsilon'=(\widehat Z')^2$ se conservan.
\end{proposition}
\begin{proof}
Se iguala cada magnitud de \eqref{iii:eq:dimensions} a su expresión
en la nueva base y se despeja su coordenada. Los factores del producto
son $d^2$, y los del cociente son $a^{-2}b^4$, justamente los requeridos
por las coordenadas transformadas de $c$ y $Z$.
\end{proof}

Esta covariancia distingue dos actos: construir una sección y expresar
esa sección en unidades elegidas. El teorema espectral fija las coordenadas
en la normalización interna declarada. Una comparación SI debe exhibir
también la evaluación de las tres secciones de referencia. Ajustarlas
mediante las magnitudes que se desea reproducir sería una calibración de
la coordenatización, no una verificación independiente de aquellas magnitudes.

Para reconstruir los canales desde coordenadas expresadas en la coordenatización
primada se deshace ese cambio antes de utilizar la raíz cúbica:
\begin{equation}
 \widehat Z=ab^{-2}\widehat Z',\quad
 \widehat c=d\widehat c',\quad
 \widehat\mu=ad^{-1}b^{-2}\widehat\mu',\quad
 \widehat\varepsilon=a^{-1}d^{-1}b^2\widehat\varepsilon'.
 \label{iii:eq:undo-units}
\end{equation}
Por ejemplo, las entradas de la inversa espectral son
$r_+=(\widehat Z\widehat c)^{-1/2}$ y
$r_-=(\widehat Z/\widehat c)^{1/2}$, calculadas después de
\eqref{iii:eq:undo-units}. Introducir directamente las coordenadas primadas
en esas expresiones cambiaría la normalización del operador y podría
sacarlas del intervalo $(3/4,1)$. La covariancia dimensional conserva
las magnitudes físicas; la inversión cúbica reconstruye el transporte
en la coordenatización interna en que fue definido.

% Incorporación focal de dependencias anteriores; no sustituye al núcleo común.
% Procedencia detallada: technical/CIERRE_DEPENDENCIAS_ANTERIORES_VI.md.
% Los valores de prueba se publican después de sus operadores; no son entradas.
\section{Coordenadas generadas, acción y normalización electrónica}
\label{vi:dep:seccion}

La construcción de las colecciones utiliza resultados de las etapas anteriores:
los caracteres regionales, el registro dodecafásico, la recta de acción y la
sección electrónica. Reunimos aquí las operaciones que los conectan y las
pruebas necesarias para componerlos. Su orden es causal: las regiones y sus
registros se producen sobre APP mediante TRIT y TPK; las evaluaciones reales
se efectúan después; la elección de unidades interviene al realizar las
magnitudes dimensionales.

Las condiciones iniciales y la selección regional se han desarrollado en
la sección anterior a la aritmética de historias. El dominio del libro
incidencial, sus recuentos y la reconstrucción del registro se definen
expresamente a continuación. La selección de una colección concreta de
visitas y trazas es una operación adicional a la evaluación de un libro:
requiere fijar sus multiplicidades, el transporte de frontera y la
cobertura de las trazas admitidas. Las fórmulas de reconstrucción que se
demuestran aquí permiten evaluar y comprobar ese registro; su
invertibilidad no selecciona por sí sola la colección incidencial.

\subsection{Lectores regionales y compatibilidad de la evaluación}
\label{vi:dep:regiones}

La clausura, la propagación y la autoescala son tres regiones de una misma
generación. La región de clausura tiene cinco emisiones y multiplicidad
\(432+4\cdot144=1008\). En el calibre orientado del catálogo, sus emisiones
son
\[
\begin{array}{c|r}
(501,614,498,169,272,272)&432\\
(810,923,870,169,272,272)&144\\
(810,983,810,169,272,272)&144\\
(870,923,810,169,272,272)&144\\
(870,923,810,418,674,521)&144.
\end{array}
\]
La palabra ternaria común es \(010211\); los representantes de propagación
y autoescala son \(201101\) y \(121200\). El censo y la orientación
preceden a las evaluaciones que siguen.

En las cinco posiciones del lector de clausura, la proyección uniforme
\(P_5=\mathbf1\mathbf1^{\mathsf T}/5\), junto con
\(P_5\lambda=\lambda\) y \(\mathbf1^{\mathsf T}\lambda=1\), determina
\(\lambda=\mathbf1/5\). Esta normalización de las posiciones del lector
es distinta de la distribución de semillas indicada en la tabla.
En el régimen elíptico \(J^2=-I\), la multiplicación da
\[
(I+aJ)(I+bJ)=(1-ab)I+(a+b)J.
\]
Siempre que \(1-ab\ne0\), la pendiente compuesta es
\(a\oplus b=(a+b)/(1-ab)\). Dos duplicaciones de \(1/5\) dan
\(5/12\) y \(120/119\). El compensador positivo \(1/q\), \(q>1\),
que devuelve la pendiente a uno satisface
\[
\frac{120/119-1/q}{1+(120/119)/q}=1,
\qquad q=239.
\]
Por tanto, el carácter de clausura de este lector es
\begin{equation}
 \Pi_{\rm cl}=4\left(4\arctan\frac15-\arctan\frac1{239}\right).
 \label{vi:dep:lector-clausura}
\end{equation}
La regla de composición, su compensación y la elección de la rama angular
fijan este valor; la identidad angular convencional lo reconoce como \(\pi\).

El carácter de propagación se realiza en una colección formal conmutativa
\(P_t\), normalizada por \(P_0=1\), \(P'_0=1\) y
\(P_{t+s}=P_tP_s\). Si \(P_t=\sum_{n\ge0}a_nt^n\), la comparación
del coeficiente de \(s^nt\) da
\[
 (n+1)a_{n+1}=a_n,\qquad a_0=1.
\]
Así, \(a_n=1/n!\); la representación a una unidad de parámetro es
\(\Pi_{\rm pr}=\sum_{n\ge0}1/n!\), reconocida como \(e\).
La unidad y el signo del generador pertenecen al lector: la sola ley
semigrupal admitiría otros generadores.

Para la autoescala, TRIT selecciona los modos
\((\Sigma,\Pi,\Pi)\). Sus tres transiciones cíclicas son
\(\Sigma\to\Pi\), \(\Pi\to\Pi\) y \(\Pi\to\Sigma\). Su matriz
de incidencia, en orden areal--volumétrico, es
\begin{equation}
 F=\begin{pmatrix}0&1\\1&1\end{pmatrix}.
 \label{vi:dep:incidencia-hojas}
\end{equation}
El vector propio positivo normalizado como \((1,x)^{\mathsf T}\)
satisface \(x=\lambda\), \(1+x=\lambda x\); por tanto,
\(x^2-x-1=0\). Su raíz positiva es la coordenada de autoescala
\(\varphi\). Esta ecuación reconoce el carácter de la incidencia ya
producida por el calendario.

\begin{proposition}[Evaluación regional mediante intervalos compatibles]
\label{vi:dep:prefijos}
Los tres lectores anteriores disponen de intervalos racionales de anchura
arbitrariamente pequeña. Sus prefijos decimales finitos se determinan por
refinamiento y son compatibles por truncamiento.
\end{proposition}
\begin{proof}
Para \(0<x<1\), dos sumas alternadas consecutivas de
\(\sum_{j\ge0}(-1)^jx^{2j+1}/(2j+1)\) encierran \(\arctan x\).
El primer término omitido acota el error. Esto proporciona intervalos
racionales para \eqref{vi:dep:lector-clausura}.
Si \(E_n=\sum_{j=0}^n1/j!\), entonces, para \(n\ge1\),
\[
 E_n<\Pi_{\rm pr}<E_n+
 \frac{n+2}{(n+1)(n+1)!}.
\]
La cota resulta de mayorar los cocientes sucesivos de la cola por
\(1/(n+2)\). Finalmente, \(F^n\) tiene entradas
\(F_{n-1},F_n,F_{n+1}\), donde \(F_0=0,F_1=1\) y
\(F_{n+1}=F_n+F_{n-1}\). La identidad
\(F_{n+1}^2-F_nF_{n+2}=(-1)^n\) muestra que dos cocientes consecutivos
encierran \(\varphi\) con anchura \(1/(F_nF_{n+1})\).
Se pueden intersectar los intervalos sucesivos para conservar el
encajamiento. Tras el reconocimiento convencional, la irracionalidad de
\(\pi,e,\varphi\) asegura que ninguna coordenada está en una frontera
decimal racional. Para cada profundidad finita, un intervalo
suficientemente estrecho queda en una única celda decimal. Su índice
determina el prefijo y el encajamiento demuestra su compatibilidad.
\end{proof}

\subsection{Coordenatización dodecafásica firmada del estado TPK pleno}
\label{vi:dep:estado-pleno}

La evolución que produce el registro se formula sobre el estado completo del
Topological Phase Kernel. Denotemos por
\(\mathsf{TPKFullState}\) el espacio de estados que conserva conjuntamente
posición APP, régimen TRIT, ruta, hoja, orientación, acarreo, frontera,
memoria, libro incidencial y coordenadas terminales. En una coordenatización canónica,
\[
 X=(X_{\rm APP},X_{\rm TRIT},X_{\rm ruta},X_{\rm mem},
       C_{\rm term},U_{12}^{\rm sgn},X_{\partial},\ldots).
\]
La dinámica de ciento ocho pasos actúa mediante la composición
\begin{equation}
 X_{\rm term}
 =\bigl(\mathcal U_{108}\circ\cdots\circ\mathcal U_1\bigr)(X_{\rm seed}),
 \qquad
 \mathcal U_t=\mathsf{Upd}_t\circ\mathsf{Tra}_t\circ\mathsf{Sel}_t,
 \label{vi:dep:estado-terminal-108}
\end{equation}
donde \(\mathsf{Sel}_t\) aplica el selector trítico de fase,
\(\mathsf{Tra}_t\) transporta hoja, ruta y orientación, y
\(\mathsf{Upd}_t\) incorpora el acarreo y la memoria producidos en el
instante \(t\). Cada factor actúa, por tanto, sobre el estado pleno y la
composición conserva la genealogía de las doce ventanas nonádicas.

El libro incidencial de la sección siguiente es la coordenada
\(\mathcal R_{12}(X_{\rm term})\). Su evaluación por los canales
\(90/120\) define
\[
 \mathcal E_{108}^{90,120}\circ\mathcal R_{12}:
 \mathsf{TPKFullState}\longrightarrow\mathbb Z^{25},
 \qquad
 X\longmapsto C(X)=(b^{90}(X),b^{120}(X),Q(X)).
\]
La proyección dodecafásica firmada es
\[
 \mathcal S_{12}:=\operatorname{pr}_{U_{12}^{\rm sgn}}:
 \mathsf{TPKFullState}\longrightarrow\mathbb Z^{12}.
\]
Sobre la subred compatible producida por la dinámica, las dos coordenatizaciones
satisfacen la identidad operatoria
\begin{equation}
 \mathcal S_{12}
 =\Pi_H\circ\mathcal E_{108}^{90,120}\circ\mathcal R_{12}
 =\Pi_H\circ\operatorname{pr}_{C}.
 \label{vi:dep:identidad-productor-e108}
\end{equation}
Así, la firma dodecafásica se obtiene por proyección y evaluación del mismo
estado que conserva el libro; su dominio incluye la orientación y la memoria
que intervienen en cada recuento. La proyección cruda
\(\rho_{108}:\mathsf{TPKFullState}\to\mathsf{CT108RawProjection}\)
designa, en cambio, la coordenatización que olvida esas coordenadas después de la
producción. Esta separación de tipos fija el orden causal entre generación,
representación firmada y lectura reducida.

\subsection{De las incidencias temporales al registro dodecafásico}
\label{vi:dep:registro}

\begin{definition}[Dominio del libro incidencial orientado]
\label{vi:dep:libro-dominio}
Un libro \(\mathscr L\) contiene una colección finita no vacía de rutas,
sus visitas identificadas y una colección finita de trazas. Cada ruta
conserva una visita por instante durante los primeros \(108\) instantes,
y las visitas posteriores que utilicen sus trazas. Una visita conserva
\[
 v=(\gamma,t,i,j,s,\tau,w,\varepsilon_\partial),
\]
donde \(\gamma\) identifica la ruta, \(t\) su instante, \(1\le i,j\le9\)
su posición APP, \(s\in\{\Sigma,\Pi,0\}\) su hoja,
\(\tau\in\{-1,0,1\}\) su orientación trítica y \(w\) su multiplicidad
entera positiva. La coordenatización incidencial evalúa
\[
 n_v=
 \begin{cases}
 i+j,&s=\Sigma,\\
 ij,&s=\Pi,\\
 9,&s=0.
 \end{cases}
 \qquad
 r_v=1+((n_v-1)\bmod9),\quad q_v=(n_v-r_v)/9.
\]
La lectura \(n_v=9\) fija la convención de umbral de esta coordenatización.
Para \(r_v=9\), la etiqueta \(\varepsilon_\partial\) distingue
corona e interior; no se obtiene únicamente del residuo.

Una traza \(\theta=(v_1,\ldots,v_N)\) recorre instantes consecutivos de
una misma ruta. Conserva una multiplicidad \(u(\theta)\in\mathbb N_{>0}\)
y una unidad de longitud reducida
\(\lambda(\theta)\in\mathbb N_{>0}\), con
\[
 N\lambda(\theta)=120(1+3j(\theta)),\qquad j(\theta)\in\mathbb N_0.
\]
Su canal es
\(\sigma(\theta)=\operatorname{sgn}(\sum_{a=1}^N\tau(v_a))\).
La colección contada en los canales \(120\) comprende las trazas de
signo no nulo. La longitud reducida y el número de instantes quedan
relacionados por \(\lambda\); no se identifican entre sí.
\end{definition}

La definición especifica el dominio de la evaluación, incluidos los
datos que deben transmitirse al lector. Sus multiplicidades y la colección
de trazas pertenecen a la selección dinámica del libro. Una enumeración
del dominio o una comprobación de formato no sustituye esa selección.

La generación conserva un libro \(\mathscr L\) de visitas y trazas.
La coordenatización temporal lo divide en ventanas
\(E_m=\{9m-8,\ldots,9m\}\), \(1\le m\le12\),
equivalentes a las posiciones \(m=\beta+1\) del registro
\(\mathcal R_{12}\). Cada ventana retiene su subruta,
orientación, incremento de frontera e incidencias locales.
Cada visita registra posición APP, hoja aritmética, instante, multiplicidad
y orientación de frontera. En cada evaluación \(n_v\) se conserva la
división \(n_v=r_v+9q_v\), con \(1\le r_v\le9\). En la ventana
\(m\), el lector incidencial expresa
\begin{align}
 A_m&=\sum_{v\in\mathscr L_m}w(v)
      \mathbf1_{\{1,3,5,7\}}(r_v),\notag\\
 C_m&=\sum_{j\ge0}
       \bigl(\nu^-_{120,m}(j)-\nu^+_{120,m}(j)\bigr),\notag\\
 V_m&=\sum_{v\in\mathscr L_m}w(v)
      \mathbf1_{\{9\}}(r_v)\,\varepsilon_{\partial}(v).
 \label{vi:dep:incidencias}
\end{align}
Aquí las trazas de cada ventana tienen soporte finito,
\(\varepsilon_{\partial}=+1\) en la corona y \(-1\) en el interior,
y las etiquetas \(\pm,120\) pertenecen al transporte conservado.
La multiplicidad \(w\) y las trazas \(\nu\) son las obtenidas de ese
transporte; no se reconstruyen escogiendo un valor de \(\alpha\).

En términos de las trazas identificadas, la segunda suma tiene la
expresión finita
\[
 C_m=-\sum_{\theta:\,t(v_1)\in E_m}u(\theta)\sigma(\theta).
\]
En efecto, agrupar por \(j(\theta)\) y por signo produce
\(\nu^-_{120,m}(j)-\nu^+_{120,m}(j)\). Esta igualdad especifica
la operación efectivamente aplicada a cada registro; conserva
separados el canal, la multiplicidad y la longitud.

\begin{proposition}[Actualización y concatenación del libro]
\label{vi:dep:libro-actualizacion}
Fijada la colección de visitas y trazas admitidas por el transporte, sus
recuentos se actualizan mediante incrementos enteros. Al añadir una
visita \(v\) de la ventana \(m\), el incremento de \((A_m,C_m,V_m)\) es
\[
 \left(w(v)\mathbf1_{\{1,3,5,7\}}(r_v),\,0,\,
 w(v)\mathbf1_{\{9\}}(r_v)\varepsilon_\partial(v)\right).
\]
Al completarse una traza \(\theta\) cuyo origen pertenece a \(E_m\),
el incremento es \((0,-u(\theta)\sigma(\theta),0)\).
Para dos segmentos consecutivos, el recuento de su concatenación es
la suma de los recuentos de las visitas y trazas contenidas en cada
segmento, más el de las trazas que cruzan su unión.
\end{proposition}
\begin{proof}
Cada suma de \eqref{vi:dep:incidencias} se indexa por identificadores
de visitas o de trazas, por lo que añadir un identificador añade
exactamente el sumando escrito. Para la concatenación, las visitas
se reparten entre los dos segmentos. Las trazas se particionan en
las contenidas en el primero, las contenidas en el segundo y las
que atraviesan la unión. La aditividad de las sumas enteras prueba
la fórmula. Conservar los identificadores y las prolongaciones
pendientes evita contar dos veces una traza o perder una que aún
no ha terminado. La reducción decimal se aplica después del
recuento; los cocientes de esa reducción permanecen en la memoria.
\end{proof}

La reducción decimal de estos tres recuentos produce
\begin{equation}
 z_m=100[A_m]_{10}+10[C_m]_{10}+[V_m]_{10},
 \qquad 0\le z_m<1000.
 \label{vi:dep:bloque-incidencial}
\end{equation}
Sean \((D_3z)_i=z_i-z_{i+3}\),
\((D_4z)_i=z_i-z_{i+4}\), con índices en \(\ZZ/12\ZZ\), y
\(Q(z)=\sum_i z_i\). El lector conserva
\(\mathcal D(z)=(D_3z,D_4z,Q(z))\). La reducción decimal olvida
múltiplos de diez de cada recuento, mientras el libro original conserva
las visitas y los transportes que los generaron.
La lectura de los eventos de código par sobre esa misma historia,
su descomposición \(12=1+5+6\) y su inversión integral se
desarrollan en la proposición \ref{vi:prop:eventos}. Ese vector
de eventos no sustituye el presente registro de incidencias:
los dominios y las coordenadas conservadas se mantienen distintos.

\begin{proposition}[Conjugación incidencial y reducción decimal]
\label{vi:dep:conjugacion-incidencial}
Sea \(\varrho(m)=13-m\). Consideremos una conjugación del libro
que lleva las visitas y las trazas de la ventana \(m\) a las de
\(\varrho(m)\) mediante biyecciones que conservan sus multiplicidades.
Supongamos que conserva la clase aritmética y la etiqueta de frontera
de las visitas, e intercambia los canales positivo y negativo de las
trazas. Entonces los recuentos de \eqref{vi:dep:incidencias} satisfacen
\begin{equation}
 (A_m^\dagger,C_m^\dagger,V_m^\dagger)
 =(A_{\varrho(m)},-C_{\varrho(m)},V_{\varrho(m)}).
 \label{vi:dep:conjugacion-recuentos}
\end{equation}
Escribiendo \(u_m=[C_m]_{10}\), sus bloques decimales cumplen
\begin{equation}
 z_m^\dagger=z_{\varrho(m)}
       +100\mathbf1_{\{u_{\varrho(m)}\ne0\}}-20u_{\varrho(m)}.
 \label{vi:dep:conjugacion-bloque}
\end{equation}
\end{proposition}
\begin{proof}
Las biyecciones conservan cada sumando de \(A_m\) y \(V_m\).
El intercambio de los canales permuta \(\nu^-\) y \(\nu^+\),
por lo que cambia el signo de \(C_m\). Se obtiene
\eqref{vi:dep:conjugacion-recuentos}. Si \(C=10q+u\),
\(0\le u<10\), la división euclídea de \(-C\) da
\[
 [-C]_{10}=10\mathbf1_{\{u\ne0\}}-u,
 \qquad
 \left\lfloor\frac{-C}{10}\right\rfloor
 =-q-\mathbf1_{\{u\ne0\}}.
\]
En \eqref{vi:dep:bloque-incidencial}, sólo cambia el dígito
intermedio. Su diferencia es
\(10([-C]_{10}-[C]_{10})=100\mathbf1_{\{u\ne0\}}-20u\),
lo que demuestra \eqref{vi:dep:conjugacion-bloque}. La segunda
identidad conserva simultáneamente el cociente entero de la reducción.
\end{proof}

La proposición precisa la acción sobre el libro antes de transportar
su representación decimal. Su aplicación a la reversión de una ruta
conserva también el convenio con que se leen los extremos. Sea
\(\gamma=(x_0,\ldots,x_n)\) una ruta enriquecida y sea \(f\) un
observable entero evaluado antes de cada avance. La ruta inversa
\(\gamma^{-1}=(x_n,\ldots,x_0)\) verifica
\begin{equation}
 \sum_{k=0}^{n-1}f(x_{n-k})
 =\sum_{k=0}^{n-1}f(x_k)+f(x_n)-f(x_0).
 \label{vi:dep:conjugacion-frontera}
\end{equation}
En efecto, el miembro izquierdo suma los índices \(1,\ldots,n\);
el primer sumando del derecho suma \(0,\ldots,n-1\). Su diferencia
es exactamente el término de frontera indicado. La identidad se aplica
a cada observable incidencial con sus datos conservados. Los extremos
se contabilizan en los recuentos enteros antes de reducir módulo diez;
en una ruta cerrada el término se anula, mientras que en una ventana
abierta puede cambiar su bloque. Las trazas que atraviesan una unión
se conservan mediante la partición de la
proposición~\ref{vi:dep:libro-actualizacion}.

Así se conectan los recuentos con la conjugación de rutas de la
sección~\ref{vi:sec:conjugacion}, manteniendo declarada la acción
concreta sobre cada incidencia. La negación de un canal y la reversión
de una ruta son operaciones sobre el libro; su reducción decimal es la
representación que acaba de calcularse.

\begin{proposition}[Imagen integral y reconstrucción del registro]
\label{vi:dep:imagen-integral}
Sean \(b,c\in\ZZ^{12}\), \(q\in\ZZ\), con índices de cero a once.
Definamos
\[
 w_i=c_i-b_{i+1},\qquad P_0=0,\qquad
 P_i=\sum_{j=0}^{i-1}w_j,\qquad T_w=\sum_{i=0}^{11}P_i.
\]
Existe un único \(k\in\ZZ^{12}\) tal que
\(D_3k=b\), \(D_4k=c\), \(Q(k)=q\), si y sólo si
\begin{equation}
 b_i=w_i+w_{i+1}+w_{i+2},\qquad
 \sum_iw_i=0,\qquad q+T_w\equiv0\pmod{12}.
 \label{vi:dep:condiciones-imagen}
\end{equation}
En ese caso,
\begin{equation}
 k_i=\frac{q+T_w}{12}-P_i.
 \label{vi:dep:inversa-incidencial}
\end{equation}
\end{proposition}
\begin{proof}
Si existe \(k\), entonces
\(c_i-b_{i+1}=k_i-k_{i+1}\). La suma cíclica de estas diferencias
es cero, su suma de longitud tres es \(b_i\), y
\(k_i=k_0-P_i\). Al sumar las doce coordenadas,
\(q=12k_0-T_w\), lo que prueba las condiciones necesarias y la fórmula.
Recíprocamente, las condiciones hacen integral la fórmula y la vuelven
cíclica. Sus diferencias de longitud tres son \(b_i\), y las de
longitud cuatro son \(w_i+b_{i+1}=c_i\). Su suma es \(q\).
La diferencia entre dos soluciones tiene todas sus diferencias
consecutivas nulas y suma cero; por ello es nula.
\end{proof}

La reconstrucción es una operación sobre la imagen tipada, no una
selección retrospectiva de sus argumentos. En particular,
\(\mathcal D(z)=\mathcal D(z')\) implica \(z=z'\).
Como la expresión de un entero de \([0,999]\) en tres cifras es única,
dos ternas de recuentos expresan el mismo registro si y sólo si coinciden
componente a componente módulo diez.

El cambio de representación firmado utiliza el bloque
\begin{equation}
 H_4=\begin{pmatrix}
 1&1&1&1\\1&1&-1&-1\\
 1&-1&1&-1\\1&-1&-1&1
 \end{pmatrix},\qquad H_4^2=4I_4.
 \label{vi:dep:hadamard}
\end{equation}
El operador \(H_{12}\) aplica \(H_4\) a las tres órbitas ordenadas
\((1,4,7,10)\), \((2,5,8,11)\), \((3,6,9,12)\).
Sobre la subred
\(\mathcal L_H=\{U\in\ZZ^{12}:H_{12}U\equiv0\pmod4\}\),
las aplicaciones
\begin{equation}
 U=H_{12}K,\qquad K=\tfrac14H_{12}U
 \label{vi:dep:K-reversible}
\end{equation}
son inversas. La congruencia expresa exactamente la integridad de la
segunda aplicación, y \(H_{12}^2=4I\) prueba ambas composiciones.

El registro firmado también se calcula directamente de las coordenadas
transversales, antes de reconstruir \(K\). Con índices de cero a once,
sean, para \(r=0,1,2\),
\[
 S_r=\sum_{j=0}^3c_{r+3j},\quad
 a_0=\frac{q+2S_0+S_1}{3},\quad
 a_1=a_0-S_0,\quad a_2=a_1-S_1,
 \qquad d_{r,j}=b_{r+3j}.
\]
La aplicación armónica expresa por órbitas
\begin{equation}
 \Pi_H^{(r)}(b,c,q)=
 (a_r,d_{r,1}-d_{r,3},d_{r,0}+d_{r,2},d_{r,0}-d_{r,2}).
 \label{vi:dep:PiH}
\end{equation}
Para probar la fórmula sobre la imagen integral, sea \(k\) su
preimagen. Si \(s_r=\sum_jk_{r+3j}\), entonces
\(S_r=s_r-s_{r+1}\) y \(q=s_0+s_1+s_2\), con índices de órbita
módulo tres. Resolver da \(s_r=a_r\). Para las cuatro coordenadas
\((x_0,x_1,x_2,x_3)\) de una órbita,
\(d_j=x_j-x_{j+1}\). Sus tres combinaciones en
\eqref{vi:dep:PiH} son respectivamente
\(x_0+x_1-x_2-x_3\), \(x_0-x_1+x_2-x_3\) y
\(x_0-x_1-x_2+x_3\). Junto con \(a_r=\sum_jx_j\), son
exactamente \(H_4(x_0,x_1,x_2,x_3)^{\mathsf T}\).
La composición efectiva es así
\(\mathscr L\to(b,c,q)\to\Pi_H(b,c,q)=U\to H_{12}U/4=K\).
La preimagen utilizada en esta demostración verifica la identidad;
no es un argumento de la aplicación \(\Pi_H\).

Para el estado terminal de \eqref{vi:dep:estado-terminal-108}, la coordenada
producida antes de la representación firmada es
\begin{align}
 b^{90}={}&(-495,-116,-684,108,601,-90,
              -173,-88,313,560,-397,461),\notag\\
 b^{120}={}&(-425,-281,-481,671,-255,30,
               475,-543,680,251,6,-128),\notag\\
 Q_{\rm TPK}={}&6263.
 \label{vi:dep:canales-terminales}
\end{align}
Al sustituir estos veinticinco campos en \eqref{vi:dep:PiH}, las tres
órbitas firmadas son
\begin{align*}
 U_{\rm term}^{(1)}&=(2378,-452,-668,-322),\\
 U_{\rm term}^{(2)}&=(1406,998,-204,-28),\\
 U_{\rm term}^{(3)}&=(2479,-551,-371,-997).
\end{align*}
La intercalación por columnas produce
\begin{equation}
 U_{\rm term}=\mathcal S_{12}(X_{\rm term})
 =(2378,1406,2479,-452,998,-551,-668,-204,-371,-322,-28,-997).
 \label{vi:dep:U-terminal}
\end{equation}
La transformada orbital \(H_{12}/4\) está definida sobre esta subred por
la divisibilidad exacta de cada bloque de cuatro. En consecuencia, la cadena
\begin{equation}
 X_{\rm term}\xrightarrow{\mathcal R_{12}}\mathscr L(X_{\rm term})
 \xrightarrow{\mathcal E_{108}^{90,120}}C_{\rm term}
 \xrightarrow{\Pi_H}U_{\rm term}
 \xrightarrow{H_{12}/4}K
 \label{vi:dep:cadena-terminal-completa}
\end{equation}
es una composición tipada desde el estado pleno hasta su sello periódico.
Las proposiciones~\ref{vi:dep:libro-actualizacion} y
\ref{vi:dep:imagen-integral} proporcionan, respectivamente, la actualización
finita del libro y la inversa integral de la última representación.

El registro generado proporciona el testigo ordenado
\begin{equation}
 K=(234,543,140,729,659,824,621,058,914,794,146,601).
 \label{vi:dep:K-testigo}
\end{equation}
El cero inicial de \(058\) forma parte de la lectura por bloques.
Para comprobarlo mediante la proposición~\ref{vi:dep:imagen-integral}, sus diferencias
y suma son
\[
\begin{split}
b={}&(-495,-116,-684,108,601,-90,-173,-88,313,560,-397,461),\\
c={}&(-425,-281,-481,671,-255,30,475,-543,680,251,6,-128),\\
q={}&6263.
\end{split}
\]
En las tres órbitas anteriores, el registro firmado es
\[
 \begin{aligned}
 U^{(1)}&=(2378,-452,-668,-322),\\
 U^{(2)}&=(1406,998,-204,-28),\\
 U^{(3)}&=(2479,-551,-371,-997).
 \end{aligned}
\]
Aplicar \(H_4/4\) a estos tres vectores reproduce
\eqref{vi:dep:K-testigo}. Estas identidades verifican la representación de
las incidencias; su génesis sigue siendo la del libro
\eqref{vi:dep:incidencias}.

\subsection{\texorpdfstring{Representación periódica de \(K\) y normalización de \(\alpha\)}{Representación periódica de K y normalización de alfa}}
\label{vi:dep:alpha}

Fijemos \(B=1000\), el origen de fase y las doce posiciones. Escribamos
\[
 N_K=\sum_{j=1}^{12}K_jB^{12-j},\qquad
 \kappa=\frac{N_K}{B^{12}-1}.
\]
Ésta es la evaluación periódica exacta del registro. La multiplicación
por \(B^{12}-1\) y la división euclídea sucesiva recuperan sus doce
bloques, incluido cualquier cero inicial. Para
\(\rho K=(K_2,\ldots,K_{12},K_1)\), la misma concatenación da
\(\kappa(\rho K)=B\kappa(K)-K_1\). Por tanto, la reversibilidad
del registro conserva un origen de fase especificado.

\begin{proposition}[Período mínimo de la representación racional]
\label{vi:dep:periodo-minimo-k}
La prolongación periódica del registro \eqref{vi:dep:K-testigo}
tiene período mínimo de \(12\) bloques en base \(1000\) y de
\(36\) cifras en base \(10\).
\end{proposition}
\begin{proof}
Los doce componentes de \(K\) son distintos. Cualquier traslación
cíclica no nula menor que doce identificaría dos componentes si fuera
un período. Por tanto, el período mínimo de bloques es doce.

La concatenación de sus bloques de tres cifras, incluido \(058\),
produce un período decimal de longitud \(36\). Esta expansión es
canónica: no es eventualmente constante en la cifra \(9\).
Sea \(d\) su período decimal mínimo. Entonces \(d\mid36\),
porque los desplazamientos que preservan una palabra periódica forman
un subgrupo cíclico. Al agrupar las cifras de tres en tres, el período
mínimo de bloques es el menor \(k>0\) tal que \(d\mid3k\), es
decir, \(d/\gcd(d,3)\). Ningún divisor propio de \(36\)
satisface \(d/\gcd(d,3)=12\). Se sigue que \(d=36\).
\end{proof}

Este período pertenece a la coordenada racional del registro, con el
origen de fase fijado. La composición siguiente conserva además las
coordenadas regionales y los cocientes de normalización de su salida.

Sean \(P_j,E_j,\Phi_j\) los bloques fraccionarios de los tres
caracteres regionales ya construidos, y extendamos \(K_j\) con período
doce. La composición de registros es
\begin{equation}
 Z_j=P_j+E_j-\Phi_j-K_j,
 \qquad
 A_j=Z_j+c_{j+1}-Bc_j.
 \label{vi:dep:normalizacion-alpha}
\end{equation}
El cociente entero \(c_j\) transporta la continuación hacia la posición
observada. En una ventana finita se fija primero \(c_{N+1}\), y luego
\[
 c_j=\left\lfloor\frac{Z_j+c_{j+1}}B\right\rfloor,
 \qquad A_j=Z_j+c_{j+1}-Bc_j.
\]
La división euclídea determina unívocamente \(0\le A_j<B\).
Como \(-1998\le Z_j\le1998\), el conjunto
\(\{-2,-1,0,1,2\}\) es estable para estos cocientes.

\begin{theorem}[Composición exacta a profundidad ilimitada]
\label{vi:dep:alpha-completa}
La evaluación completa de la normalización, en la convención decimal
canónica, es
\begin{equation}
 \alpha_A=\pi+e-\varphi-4-\kappa.
 \label{vi:dep:alpha-A}
\end{equation}
Existe una única expansión canónica \((A_j)\), con bloques entre cero
y \(B-1\), y una sucesión acotada de cocientes enteros que satisfacen
\eqref{vi:dep:normalizacion-alpha} y \(c_1=0\).
\end{theorem}
\begin{proof}
Las partes enteras regionales son \(3,2,1\). Por convergencia absoluta,
\[
 y:=\sum_{j\ge1}Z_jB^{-j}
   =(\pi-3)+(e-2)-(\varphi-1)-\kappa.
\]
Los primeros bloques son \(141,718,618,234\), de modo que
\(Z_1=7\). La suma de la cola tiene valor absoluto a lo sumo
\(1998\sum_{j\ge2}B^{-j}=0.002\); en particular \(0<y<1\).
Tomemos la expansión canónica \(y=\sum A_jB^{-j}\), eligiendo la
expansión terminal y no la eventualmente constante \(B-1\) si se
presenta una frontera. Pongamos
\[
 c_j=\sum_{r\ge0}(Z_{j+r}-A_{j+r})B^{-r-1}.
\]
La igualdad de los valores completos muestra que cada \(c_j\) es el
negativo de una suma finita de enteros ponderados por potencias de
\(B\); es por tanto entero, y \(c_1=0\). Separar el primer término
da \(Bc_j=Z_j-A_j+c_{j+1}\). Las cotas de las colas hacen acotada
la sucesión. Más precisamente, la cola de \(Z\) está en \([-2,2]\)
y la cola canónica de \(A\) en \([0,1)\), por lo que
\(-3<c_j\le2\). Al ser entero, \(c_j\) pertenece al conjunto
\(\{-2,-1,0,1,2\}\) indicado.
Recíprocamente, la suma de la recurrencia hasta \(N\) da
\[
 \sum_{j=1}^N A_jB^{-j}
 =\sum_{j=1}^N Z_jB^{-j}-c_1+c_{N+1}B^{-N}.
\]
La acotación anula el último término en el límite. La unicidad de la
expansión canónica de \(y\) determina \(A\), y la fórmula de colas
determina entonces todos los \(c_j\).
\end{proof}

La evaluación finita \(\kappa_{36}=N_K/B^{12}\) difiere de
\(\kappa\) en \(B^{-12}\kappa\). Además, imponer
\(c_{13}=0\) en una ventana de doce bloques no sustituye el cociente
de la continuación completa. En el registro expresado, la continuación
da \(c_{13}=-1\): el duodécimo bloque es \(662\), mientras la
ventana cortada con cociente terminal cero da \(663\). La evaluación
completa comienza por
\[
 \alpha_A=0.007297352569283800997285105472380662999\ldots.
\]
El período de \(K\) no impone período a \((A_j)\), pues en
\eqref{vi:dep:normalizacion-alpha} también intervienen los tres
registros regionales y los cocientes enteros.

En las fórmulas siguientes, \(\alpha\) designa esta salida posicional
completa \(\alpha_A\). La vía analítica de vacancias tiene su operador
propio; un polinomio truncado de orden nueve y la raíz de ese polinomio
no se identifican aquí con el operador completo. Las composiciones de
masas que siguen no requieren promover aquella aproximación a una
igualdad entre las dos vías infinitas.

\subsection{Construcción tridimensional del prefactor electrónico}
\label{vi:dep:prefactor}

Sea \(\mathcal A_3=\{1,\ldots,9\}^3\), con medida uniforme, y
\(\rho_9(n)=1+((n-1)\bmod9)\). En este dominio actúan
\[
 \Sigma(x)=\rho_9(x_1+x_2+x_3),\qquad
 \Pi(x)=\rho_9(x_1x_2x_3).
\]
En \(\ell^2(\mathcal A_3)\), sean \(P\) y \(Q\) las medias
condicionales sobre las fibras de \(\Sigma\) y \(\Pi\),
respectivamente. Cada una es idempotente y autoadjunta: dentro de cada
fibra sustituye un vector por su componente constante, y conserva
ortogonalmente esa componente. Son, pues, proyectores ortogonales.

Las fibras aditivas tienen cardinal \(81\). Las multiplicativas tienen
cardinales
\[
m=(36,36,108,36,36,108,36,36,297).
\]
La incidencia conjunta
\(N_{ab}=\#\{x:\Sigma(x)=a,\Pi(x)=b\}\) es \(9N_0\), donde
\begin{equation}
 N_0=\begin{pmatrix}
0&1&1&0&1&1&0&1&4\\1&0&1&1&0&1&1&0&4\\
1&0&2&0&1&2&0&0&3\\0&1&1&0&1&1&0&1&4\\
1&0&1&1&0&1&1&0&4\\0&0&2&1&0&2&0&1&3\\
0&1&1&0&1&1&0&1&4\\1&0&1&1&0&1&1&0&4\\
0&1&2&0&0&2&1&0&3
\end{pmatrix}.
 \label{vi:dep:incidencia-cubica}
\end{equation}
Esta matriz se obtiene contando las \(729\) ternas. También se
reproduce, sin redondeo, con
\(N_{ab}=\sum_{i,j,k=1}^9
\mathbf1_{\{a\}}(\rho_9(i+j+k))
\mathbf1_{\{b\}}(\rho_9(ijk))\).
En las bases normalizadas de funciones constantes por fibra, el
entrelazamiento tiene matriz \(C_{ab}=N_{ab}/\sqrt{81m_b}\).

\begin{proposition}[Norma de la no conmutatividad de las dos hojas]
\label{vi:dep:norma-PQ}
Los proyectores anteriores satisfacen
\begin{equation}
 \|[P,Q]\|=\frac{\sqrt3}{4}.
 \label{vi:dep:prefactor-valor}
\end{equation}
\end{proposition}
\begin{proof}
En el espacio de las nueve fibras aditivas, \(PQP\) tiene matriz
\(M=CC^{\mathsf T}\). La multiplicación de las matrices de
\eqref{vi:dep:incidencia-cubica} da autovalor uno sobre el vector
constante, \(1/4\) sobre \((1,-1,0,1,-1,0,1,-1,0)\), y
\(7/132\) sobre \((1,1,-2,1,1,-2,1,1,-2)\).
Sobre las coordenadas \(\{3,6,9\}\) con suma cero actúa por
\(1/18\), con multiplicidad dos. Sobre cada uno de los conjuntos
\(\{1,4,7\}\) y \(\{2,5,8\}\), los vectores de suma cero
pertenecen al núcleo; dan cuatro dimensiones. Estos subespacios
ortogonales suman las nueve dimensiones.

Para un autovector unitario \(u\in\operatorname{ran}P\) de
\(PQP\) con autovalor \(0<\lambda<1\), el vector
\(v=(Qu-\lambda u)/\sqrt{\lambda(1-\lambda)}\) es unitario y
ortogonal a \(\operatorname{ran}P\). En la base \((u,v)\),
\[
 P=\begin{pmatrix}1&0\\0&0\end{pmatrix},\quad
 Q=\begin{pmatrix}\lambda&\sqrt{\lambda(1-\lambda)}\\
 \sqrt{\lambda(1-\lambda)}&1-\lambda\end{pmatrix}.
\]
El conmutador tiene norma \(\sqrt{\lambda(1-\lambda)}\).
Los bloques correspondientes a \(\lambda=0,1\) conmutan. De los
autovalores calculados, el máximo se alcanza en \(\lambda=1/4\),
pues \(\lambda(1-\lambda)\) es creciente en \([0,1/2]\).
Su raíz cuadrada es \(\sqrt3/4\).
\end{proof}

En el bloque extremal, \(S=2P-I\) y
\(J=4[P,Q]/\sqrt3\) satisfacen
\[
 S^2=I,\qquad J^2=-I,\qquad SJ=-JS.
\]
Las cuatro matrices \(I,S,J,SJ\) son linealmente independientes,
de modo que realizan \(\mathrm{Cl}_{1,1}\simeq M_2(\RR)\).
Los elementos \(n_\pm=(S\pm J)/2\) cumplen
\(n_\pm^2=0\) y \(n_+n_-+n_-n_+=I\), por expansión directa.
El prefactor de la ecuación electrónica procede así de dos evaluaciones
de APP en dimensión tres y de su conmutador, antes de cualquier masa
de referencia.

\subsection{Lectura de frontera y composición electrónica basal}
\label{vi:dep:basal}

La incidencia \eqref{vi:dep:incidencia-hojas} tiene vector propio
positivo \(r=(1,\varphi)^{\mathsf T}\). La normalización
\(p_{ij}=F_{ij}r_j/(\varphi r_i)\) produce
\[
 P_{\rm av}=\begin{pmatrix}0&1\\\varphi^{-2}&\varphi^{-1}\end{pmatrix},
 \qquad \mu=\frac{(1,\varphi^2)}{1+\varphi^2}.
\]
La identidad \(\varphi^{-2}+\varphi^{-1}=1\) prueba que las filas
suman uno. Como \(F\) es simétrica, los pesos proporcionales a
\(r_i^2\) satisfacen balance detallado; por ello \(\mu P_{\rm av}=\mu\).
La incidencia es irreducible y tiene una autorretención, de modo que
esta distribución estacionaria es única. La frecuencia cronológica
\((1/3,2/3)\) cuenta los tres instantes del calendario; \(\mu\)
pondera sus trayectorias admisibles de memoria uno.

El operador que conserva las hojas, fija la normalización volumétrica
en uno y marca la areal por el carácter de clausura es
\(D_\pi=\operatorname{diag}(\pi,1)\). Por inducción,
\[
 F^n=\begin{pmatrix}F_{n-1}&F_n\\F_n&F_{n+1}\end{pmatrix}.
\]
El cociente de retornos diagonales marcados es
\[
 \Theta_n=\pi\frac{F_{n-1}}{F_{n+1}}
 \longrightarrow\frac{\pi}{\varphi^2}.
\]
La convergencia se deduce del autovalor dominante \(\varphi\), o de
los cocientes consecutivos ya acotados. El lector
\(\mathcal R(x,y)=x/y^2\), bajo el intercambio de argumentos,
expresa también
\begin{equation}
 \rho_A=\frac{\pi}{\varphi^2},\qquad
 \rho_E=\frac{\varphi}{\pi^2},\qquad
 \varphi^3\rho_A^2\rho_E=1.
 \label{vi:dep:dos-orientaciones}
\end{equation}
La última igualdad se comprueba multiplicando las fracciones.
Relaciona dos orientaciones del mismo lector y no selecciona los
restantes coeficientes de la ecuación electrónica.

El retorno central tipado tiene longitud \(27\), con soporte
\(\ZZ/9\ZZ\times\FF_3\). Las cinco posiciones de la microfibra
se insertan, en el calibre regional fijado, como
\((4,0),(5,0),(6,0),(6,1),(6,2)\).
El complemento tiene \(27-5=22\) posiciones. Su incidencia con los
tres estados de fase trítica tiene \(5\cdot3=15\) pares.
Así se obtienen los enteros \(-22\) y \(15\) utilizados por el
lector electrónico; la orientación fija el signo del primer término.
La composición cúbica utiliza tres copias de la coordenada de cierre,
por lo que su evaluación es \(\alpha^3\). El soporte, la inserción
regional y esta regla de composición son datos operativos distintos;
un recuento aislado no sustituye su declaración.

El defecto global y su retorno nonádico se expresan mediante
\begin{equation}
 d_4=\frac{\pi-e\log\pi}{270},\qquad
 \Delta_4=d_4\left(1+\frac{\pi}{729}\right).
 \label{vi:dep:delta4}
\end{equation}
El denominador \(270\) pertenece a la corona sin su unidad de
clausura; \(729=9^3\) al soporte cúbico. La función
\(f(x)=x-e\log x\) alcanza su mínimo cero en \(x=e\), pues
\(f'(x)=1-e/x\) cambia allí de signo. Como \(\pi\ne e\),
se obtiene \(d_4>0\) y \(\Delta_4>0\).

\begin{definition}[Coordenada electrónica basal]
\label{vi:dep:electron-basal}
La composición de los lectores anteriores es
\begin{equation}
 \varepsilon_e^{(0)}=
 \frac{\sqrt3}{4}
 \left[\exp\left(\frac{\varphi}{\pi^2}\right)-22\alpha^3\right]
 (1+15\Delta_4).
 \label{vi:dep:formula-basal}
\end{equation}
\end{definition}
La exponencial actúa sobre la orientación \(\rho_E\) de
\eqref{vi:dep:dos-orientaciones}; la lectura de vacancia corrige esa
componente y el último factor conserva el retorno del defecto global.
Para \(0<\alpha<0.01\), todos los factores son positivos:
\(\exp(\varphi/\pi^2)>1\), \(22\alpha^3<22\cdot10^{-6}<1\)
y \(\Delta_4>0\). Ésta es una coordenada adimensional. Su presencia
en una coordenatización de masa se define después de la construcción de la escala.

\subsection{Escala determinantal y constante reducida de Planck}
\label{vi:dep:planck}

El bloque local \eqref{vi:dep:hadamard} tiene determinante \(-16\).
Su forma normal de Smith es \(\operatorname{diag}(1,2,2,4)\).
En efecto, el máximo divisor de las entradas es uno; los menores de
orden dos son pares y uno vale dos en valor absoluto. Los menores de
orden tres son, por la identidad de la adjunta, de valor absoluto cuatro,
y el determinante tiene valor absoluto dieciséis. Los divisores
determinantales son \(1,2,4,16\), cuyas razones sucesivas dan los
cuatro factores de Smith. Por consiguiente,
\begin{equation}
 G_H=\operatorname{coker}H_4
 \simeq\ZZ/2\ZZ\oplus\ZZ/2\ZZ\oplus\ZZ/4\ZZ,
 \qquad |G_H|=16.
 \label{vi:dep:conucleo}
\end{equation}

El lector normado local multiplica la misma coordenada \(\alpha\)
una vez por cada clase de ese conúcleo, y aplica una vez la coordenada
de autoescala. Su definición produce
\begin{equation}
 \Sigma_H=\varphi\prod_{g\in G_H}\alpha
         =\varphi\alpha^{16}.
 \label{vi:dep:sigma-H}
\end{equation}
El orden del conúcleo determina el exponente de este lector. La forma
normal de Smith y la elección del lector normado cumplen funciones
distintas: la primera no convierte cualquier funcional del bloque en
la expresión \eqref{vi:dep:sigma-H}.

\begin{proposition}[Década de acción]
\label{vi:dep:decada}
La salida \eqref{vi:dep:sigma-H} satisface
\(10^{-34}<\Sigma_H<10^{-33}\); su valoración decimal es \(-34\).
\end{proposition}
\begin{proof}
Las evaluaciones regionales y posicionales dan los intervalos
\(1.618<\varphi<1.619\) y \(0.007297<\alpha<0.007298\).
Por positividad,
\[
 \frac{1618\,7297^{16}}{10^{99}}
 <\Sigma_H<
 \frac{1619\,7298^{16}}{10^{99}}.
\]
Las desigualdades enteras
\(1618\,7297^{16}>10^{65}\) y
\(1619\,7298^{16}<10^{66}\) prueban el resultado sin utilizar una
constante de acción externa.
\end{proof}

La década y la mantisa son representaciones de etapas diferentes.
Antes de construir la segunda, el lector regional selecciona el retorno
persistente de la microfibra de clausura. La proyección
\(p_{\rm ret}(u_1,\ldots,u_6)=(u_4,u_5,u_6)\) aplicada a las
cinco emisiones de \ref{vi:dep:regiones} produce el bloque
\(b_\pi=(169,272,272)\) cuatro veces y
\(b_\pi^-=(418,674,521)\) una vez. Estas multiplicidades cuentan
filas regionales, no las multiplicidades de semillas \(432,144\).

\begin{proposition}[Selector equivariante del retorno regional]
\label{vi:dep:selector169}
Entre las reglas que conservan las permutaciones de las cinco regiones,
seleccionan el bloque terminal de multiplicidad máxima, son equivariantes
bajo permutaciones de sus tres posiciones y retienen una posición fija
por el estabilizador del bloque, la coordenada seleccionada es única
y vale \(\rho_\pi=169\).
\end{proposition}
\begin{proof}
La proyección de las cinco filas muestra que el bloque de multiplicidad
máxima es único y es \(b_\pi\). Su estabilizador en
\(\mathfrak S_3\) intercambia las dos posiciones de valor \(272\)
y fija únicamente la posición de valor \(169\). Retener esa posición
es compatible con cualquier permutación simultánea de coordenadas.
Equivalentemente, el valor se expresa sin elegir una posición inicial
como
\(\operatorname{sing}(b_\pi)=\sum_j(b_\pi)_j
-2\operatorname{moda}(b_\pi)=169\).
\end{proof}

La graduación por longitud expresa una palabra de longitud \(n\)
mediante el carácter \(\chi_z([w])=z^{|w|}\). La concatenación
suma longitudes, por lo que
\(\chi_z([uv])=\chi_z([u])\chi_z([v])\). Como las emisiones
regionales pertenecen a \(\mathcal U_6\), el impulso del retorno
aparece como \(\rho_\pi z^6\). El grado seis corresponde a la
longitud de la emisión, no a la longitud de un paseo cerrado que
recorra varios anclajes del grafo.

El peso de prolongación se obtiene en la línea determinante del
transporte bicapa. Para una involución real \(R\) en dimensión dos,
con \(\operatorname{tr}R=0\), sean
\[
 x_A=\frac{50\pi}{9}\alpha,\qquad
 \mathcal T(x_A,y)=\exp(-x_AI_2-yR).
\]
La coordenada orientada \(y\) todavía es formal. En una base propia
de \(R\), las dos entradas son \(e^{-x_A-y}\) y
\(e^{-x_A+y}\). Por tanto,
\[
 \bigwedge^2\mathcal T=\det\mathcal T
 =e^{-2x_A}=D_A.
\]
La dependencia en \(y\) se cancela. La multiplicatividad por
concatenación determina el carácter de prolongación \(k\mapsto D_A^k\):
cualquier carácter con valor elemental \(D_A\) debe tener ese valor
en \(k=1+\cdots+1\). Esta construcción precede a la evaluación
de la coordenada angular conjugada.

El lector regional de orden cinco emplea los autovalores \(9,5\)
del bloque central \(\bigl(\begin{smallmatrix}7&2\\2&7\end{smallmatrix}\bigr)\),
la órbita de cuatro posiciones y el calendario de doce ventanas.
Su regla graduada fija
\begin{align}
 H_5={}&\frac12\sqrt{\frac{1000\alpha}{\varphi}}-\alpha
       +\frac9{16}\alpha^2-\frac59\alpha^3
       +\frac7{48}\alpha^4-\frac1{54}\alpha^5,\notag\\
 D_A={}&\exp\left(-\frac{100\pi\alpha}{9}\right),\notag\\
 \mathcal C_\pi={}&169\alpha^6+
                  \frac{D_A\alpha^7}{1-D_A\alpha},\notag\\
 \eta_{\rm ret}={}&H_5-\frac{90}{\pi}\mathcal C_\pi.
 \label{vi:dep:mantisas}
\end{align}
En esa regla, los cocientes son
\(9/16=9/4^2\), \(5/9\), \(7/48=(\operatorname{tr}B_c/2)/(4\cdot12)\)
y \(1/54=1944/104976\). Su asignación a los grados y los signos
forma parte del lector graduado. El entero \(169\) es la representación
del selector demostrado en la proposición \ref{vi:dep:selector169}.
Su dato regional es anterior a la evaluación de \(\alpha\); el
lector analítico lo compone después con la coordenada de cierre.

La cola de \(\mathcal C_\pi\) tiene una construcción por renovación.
El impulso \(\rho_\pi z^6\) y la continuación estricta pertenecen
a sumandos distintos. Después de registrar una sola vez el impulso,
la continuación se normaliza por la unidad de la línea determinante;
cada nuevo paso la multiplica por \(D_A\). La orientación del
calibre regional coloca este retorno en el sector compensador que
se sustrae del carácter de quinto grado.
Con parámetro formal \(z\), la regla
\begin{equation}
 R(z)=D_Az^7+D_AzR(z)
 \label{vi:dep:renovacion}
\end{equation}
determina de modo único
\(R(z)=D_Az^7/(1-D_Az)\). Se puede probar por comparación de
coeficientes o por inversión de la serie formal \(1-D_Az\).
La normalización del primer retorno, además de la razón recurrente,
es esencial para esta unicidad. Como \(0<D_A<1\) y \(\alpha<1\),
la evaluación en \(z=\alpha\) converge absolutamente.

\begin{proposition}[Positividad y razón de retorno de acción]
\label{vi:dep:positividad-accion}
En la rama generada, \(0<\eta_{\rm ret}<H_5\). La razón
\begin{equation}
 R_{\rm act}=\frac{H_5}{\eta_{\rm ret}}
 \label{vi:dep:Ract}
\end{equation}
es positiva, mayor que uno y adimensional.
\end{proposition}
\begin{proof}
Los límites \(0.007<\alpha<0.008\) y \(\varphi<1.619\) dan
\(\tfrac12\sqrt{1000\alpha/\varphi}>1\). La suma del valor absoluto
de los términos negativos restantes de \(H_5\) es menor que
\(0.008001\), por lo que \(H_5>0.99\).
Como \(D_A<1\), \(\pi>3\) y \(\alpha<0.008\),
\[
 0<\frac{90}{\pi}\mathcal C_\pi
 <30\left(169(0.008)^6+
                  \frac{(0.008)^7}{1-0.008}\right)<10^{-6}.
\]
Se deduce \(0<\eta_{\rm ret}<H_5\). La razón de dos coordenadas
de acción en la misma base es adimensional y no cambia al reescalar
esa base.
\end{proof}

Sea \(\mathcal U_S\) una base de la recta de acción. Las dos
secciones, anterior y posterior al retorno, son
\begin{equation}
 \hbar_{\rm pre}=H_5\,10^{-34}\mathcal U_S,
 \qquad
 \hbar_{\rm ret}=\eta_{\rm ret}\,10^{-34}\mathcal U_S.
 \label{vi:dep:hbar}
\end{equation}
Su lectura física es la construcción propuesta de la constante reducida
de Planck, con la sección de retorno explicitada. La unidad
\(\mathcal U_S\mapsto\mathrm{J\,s}\) es una realización metrológica
posterior. El coeficiente de \(\Sigma_H\) dentro de su década no es
\(H_5\): el determinante expresa la escala y el lector regional
expresa la coordenada que ocupa esa escala. Para cada sección,
\(h=2\pi\hbar\) expresa la acción de una vuelta completa.

\subsection{Coordenadas angulares y reloj de la realización absoluta}
\label{vi:dep:reloj}

El par angular utiliza las coordenadas adimensionales ya generadas:
\begin{equation}
 A_{\deg}=1000\alpha,\qquad
 C^*_{\deg}=2(\eta_{\rm ret}+\alpha),\qquad
 x=\frac{\pi}{180}A_{\deg},\quad
 y=\frac{\pi}{180}C^*_{\deg}.
 \label{vi:dep:angulos}
\end{equation}
La primera operación desplaza un bloque la expansión en base mil.
La segunda es la representación bidireccional de la coordenada regional:
equivalentemente,
\(y=\pi(H_5+\alpha)/90-\mathcal C_\pi\).
Esta igualdad se deduce sustituyendo \eqref{vi:dep:mantisas}; no
introduce una regla adicional de selección del contraángulo.
La coordenatización de una vuelta tiene \(360\) unidades angulares o \(2\pi\)
radianes. Así, para una coordenada levantada,
\[
 \hbar\theta_{\rm rad}
 =h\frac{\theta_{\deg}}{360}.
\]
El factor \(10^{-34}\) permanece en la recta de acción y no se
inserta en las coordenadas angulares.

Los canales positivos son
\[
 q_+=e^{-x-y},\qquad q_-=e^{-x+y}.
\]
Se recupera la pareja mediante
\(x=-\tfrac12\log(q_+q_-)\),
\(y=\tfrac12\log(q_-/q_+)\). El intercambio de canales conserva
\(x\) y cambia el signo de \(y\). El operador con proyectores
orientados \(P_\pm\),
\(T=q_+P_++q_-P_-\), conserva esas etiquetas al expresar sus dos
autovalores.

Para una unidad temporal interna \(t_0>0\), el cierre de \(108\)
pasos determina
\begin{equation}
 \omega_0=\frac{2\pi}{108t_0},\qquad
 E_{\rm clk}=\hbar_{\rm ret}\omega_0,
 \qquad m_{\rm clk}=\frac{E_{\rm clk}}{c_{\rm int}^{2}}.
 \label{vi:dep:reloj-absoluto}
\end{equation}
Estas fórmulas producen dimensiones de frecuencia, energía y masa.
La velocidad y la longitud se obtienen mediante la respuesta de los
mismos canales, como precisan las identidades siguientes. Expresarlas
en segundos, julios o unidades de masa requiere las coordenatizaciones
correspondientes; esas coordenatizaciones no alteran los canales ni la sección
electrónica que se normaliza respecto al reloj.

\subsubsection{Identidad entre velocidad constitutiva y velocidad del reloj}

En el dominio \(x>|y|\), los dos canales satisfacen \(0<q_\pm<1\).
Sobre la realización de dos hojas, con proyectores de rango uno
\(P_+,P_-\), se define la respuesta temporal
\[
 \mathsf V=(I-T^{90})(I-T^{120})^{-1}
          =r_+P_++r_-P_-,
 \qquad r_\pm=\frac{1-q_\pm^{90}}{1-q_\pm^{120}}>0.
\]
El inverso existe porque ambos autovalores de \(T^{120}\) son menores
que uno. La lectura simétrica de la respuesta es
\[
 \mathcal E=
 \log[(1-q_+^{90})(1-q_-^{90})]
 -\log[(1-q_+^{120})(1-q_-^{120})]
 =\log(r_+r_-).
\]
Sean \(u_S,u_C,u_T\) las bases positivas de acción, velocidad y
tiempo, y \(u_L=u_Cu_T\) la base de longitud. Las lecturas
cinemática y constitutiva son
\[
 c_{\rm int}=e^{-\mathcal E}u_C,\qquad
 c_{\rm vac}=(r_+r_-)^{-1}u_C.
\]

\begin{proposition}[Realización común de propagación y masa absoluta]
\label{vi:dep:velocidad-comun}
En una misma coordenatización dimensional,
\[
 c_{\rm int}=c_{\rm vac}=\widehat c\,u_C,
 \qquad \widehat c=(\det\mathsf V)^{-1}=(r_+r_-)^{-1}.
\]
Para \(t_0=\tau_0u_T\), con \(\tau_0>0\), la longitud elemental es
\[
 \ell_0=c_{\rm int}t_0=\tau_0\widehat c\,u_L.
\]
La masa cronológica de \eqref{vi:dep:reloj-absoluto} conserva, por
tanto, la forma
\[
 m_{\rm clk}=
 \frac{2\pi\,\eta_{\rm ret}\,10^{-34}}
      {108\,\tau_0\,\widehat c^{\,2}}\,
 \frac{\mathcal U_S}{u_Tu_C^2}.
\]
\end{proposition}
\begin{proof}
La identidad \(\mathcal E=\log(r_+r_-)\) da
\(e^{-\mathcal E}=(r_+r_-)^{-1}\); el rango uno de cada hoja da
\(\det\mathsf V=r_+r_-\). El intercambio de hojas permuta los
factores y conserva su producto. La longitud se obtiene multiplicando
la misma velocidad por \(t_0\). Sustituir estas expresiones y
\eqref{vi:dep:hbar} en \eqref{vi:dep:reloj-absoluto} produce la
fórmula másica. El factor \(\mathcal U_S/(u_Tu_C^2)\) tiene
dimensión de masa; no interviene en la selección de los canales.
\end{proof}

Para una ruta homogénea de \(n>0\) transiciones, las lecturas aditivas
\(T(\gamma)=nt_0\) y \(L(\gamma)=n\ell_0\) dan
\(L(\gamma)/T(\gamma)=c_{\rm vac}\). Si el canal varía por
arista, se suman sus longitudes; el cociente es entonces la velocidad
media de esas aristas, no necesariamente una velocidad local común.

La igualdad de secciones permanece bajo cambios de coordenatización. Si
\(u_C'=du_C\), \(u_T'=\eta u_T\), con \(d,\eta>0\), entonces
\[
 \widehat c'=\widehat c/d,\qquad
 \tau_0'=\tau_0/\eta,\qquad u_L'=d\eta u_L,
 \qquad
 \tau_0'\widehat c'u_L'=\ell_0.
\]
En particular, elegir la base adaptada \(u_C'=c_{\rm vac}\)
hace \(\widehat c'=1\), pero la coordenada del lector en la
coordenatización interna sigue siendo \((r_+r_-)^{-1}\).

\subsection{Dos lectores de la trayectoria electrónica y evaluación beta}
\label{vi:dep:beta}

La trayectoria central empieza en \((r,c)=(5,5)\), con
orientación \((h,v)=(1,1)\). TRIT repite el ciclo \((+1,-1,0)\).
En el primer régimen se lee \(r+c\) y después se avanza
\(c\mapsto1+((c-1+h)\bmod9)\); en el segundo se lee \(rc\) y
después se avanza \(r\mapsto1+((r-1+v)\bmod9)\); en el régimen
cero se lee \(9\), sin desplazamiento. Se expresa la clase módulo
tres de la raíz digital de la lectura. Tras los pasos \(27,54,81\),
se invierten simultáneamente \(h,v\). La representación de \(108\)
pasos conserva la trayectoria completa y da
\begin{equation}
 a=100000220,\quad b=120200000,\quad
 \gamma_e=(a^3b^3)^2,
 \quad\tau_e=(+,+,+,-,-,-,+,+,+,-,-,-).
 \label{vi:dep:trayectoria-electron}
\end{equation}
La identidad se comprueba recorriendo los nueve instantes de cada
ventana; al tercer bloque la inversión cambia \(a\) por \(b\),
y el segundo recorrido de \(54\) pasos repite la representación.
El estado conserva asimismo sus coordenadas de memoria.

El primer lector mide desacuerdos cíclicos de la representación ternaria:
\[
 D_k=\#\{t\in\ZZ/108\ZZ:\gamma_e(t)\ne\gamma_e(t+k)\}.
\]
Como la palabra expresada tiene período \(54\), cada desacuerdo
aparece con su trasladado por \(54\); todos los \(D_k\) son pares.
En la mitad \(a^3b^3\), el recuento de desacuerdos para
\(k=1,3,4,27,36\) da respectivamente \(27,28,34,24,16\).
Al duplicarlo,
\[
 (D_1,D_3,D_4,D_{27},D_{36})=(54,56,68,48,32).
\]
Por tanto,
\begin{equation}
 K_D=\left(D_{27}+D_{36},D_1,\frac{D_4-D_3}{2}\right)
     =(80,54,6).
 \label{vi:dep:KD}
\end{equation}

El segundo lector actúa sobre la orientación de las doce ventanas.
Definamos, con índices cíclicos,
\[
 C_k=\sum_{r=0}^{11}\tau_r\tau_{r+k},\quad
 A_\ell=\sum_{r=0}^{11}\left(\sum_{j=0}^{\ell-1}\tau_{r+j}\right)^2,
 \quad P_6=\sum_{r=0}^{5}\tau_r\tau_{r+6}.
\]
La combinación entera primitiva orientada que cancela la parte diagonal
de \(A_3,A_4\) es \(\Omega=4A_3-3A_4\): sus coeficientes
satisfacen \(3u+4v=0\), \(\gcd(u,v)=1\), y se fija \(u>0\).
Expandir el cuadrado prueba
\[
 A_\ell=\ell C_0+
     2\sum_{k=1}^{\ell-1}(\ell-k)C_k,
 \qquad\Omega=-2C_1-4C_2-6C_3.
\]
Para \(\tau_e\), se obtiene
\(C_0=12,C_1=4,C_2=-4,C_3=-12\),
\(A_3=44,A_4=32,P_6=6\). De aquí,
\begin{equation}
 K_\Omega=(\Omega,54,P_6)=(80,54,6)=K_D.
 \label{vi:dep:KOmega}
\end{equation}
El segundo componente conserva el semirretorno de \(108\) pasos;
no se introduce por la cancelación diagonal. La igualdad de ambos
lectores se ha probado sobre la trayectoria central especificada;
no se extiende por definición a todas las colecciones del atlas.

\begin{definition}[Coordenada electrónica posterior al retorno]
La composición electrónica completa es
\begin{equation}
 \varepsilon_e^\beta
  =\varepsilon_e^{(0)}
      R_{\rm act}^{80-54\alpha+6\Delta_4}.
 \label{vi:dep:electron-beta}
\end{equation}
\end{definition}
La positividad de \eqref{vi:dep:formula-basal} y
\eqref{vi:dep:Ract} permite definir esta potencia mediante el
logaritmo real. El triple entero procede indistintamente de los dos
lectores centrales acabados de calcular. La evaluación relativa
preserva la constante reducida de Planck mediante la razón de sus
secciones, y la realización absoluta utiliza además
\eqref{vi:dep:reloj-absoluto} y la coordenatización dimensional declarada.

En las tablas heredadas, \(\mathfrak e_0\) denomina la
coordenada basal \(\varepsilon_e^{(0)}\) en la coordenatización allí
especificada. La sección \(\hbar_{\rm int}\) usada para una
evaluación de supervivencia debe declararse en la misma coordenatización:
para la realización posterior al retorno del presente desarrollo
es \(\hbar_{\rm ret}\). Las filas históricas conservan la
convención declarada en su tabla de origen y no se recalculan mediante un
cambio silencioso de esta sección.

Si \(\mathcal U_E\) es la base de energía de la evaluación
electrónica, \(E_e=\varepsilon_e^\beta\mathcal U_E\). Al emplear
el cuanto cronológico como base, el coeficiente es
\[
 \widehat\varepsilon_e
   =\varepsilon_e^\beta\frac{\mathcal U_E}{E_{\rm clk}},
 \qquad E_e=\widehat\varepsilon_e E_{\rm clk}.
\]
La igualdad conserva el objeto y transforma su coordenada. Por ello
el número que se presenta en una coordenatización de MeV no puede multiplicarse
sin conversión por cualquier otro cuanto de energía. El electrón
conserva su construcción estructural al expresarse respecto a una
escala cronológica: no se sustituye por la masa de ese cuanto.

La comparación metrológica se aplica al resultado de esta composición,
con la versión de \(\alpha\), la sección de acción, el lector
\(\Delta_4\) y la coordenatización de unidades declaradas. En particular,
\(d_4\) y \(\Delta_4\) no son intercambiables: reemplazarlos
simultáneamente en el factor basal y en el exponente cambia el
resultado por la razón positiva
\[
 \frac{1+15\Delta_4}{1+15d_4}
 R_{\rm act}^{6(\Delta_4-d_4)}>1.
\]
Las fórmulas conservadas permiten localizar esa diferencia sin
atribuirla a una masa externa ni seleccionar de nuevo la trayectoria.

\section{La partícula como clase de secciones persistentes}
\label{vi:sec:particula}

La construcción material comienza en el estado que producen APP, TRIT y
TPK. APP proporciona las posiciones y las dos operaciones aritméticas;
TRIT distingue el régimen y la orientación de sus transiciones; TPK
transporta esa información junto con los cocientes, las incidencias de
frontera y la memoria. Al prolongar el estado, una evaluación finita
permanece vinculada a las evaluaciones que la preceden. La noción de
partícula se apoya en esta compatibilidad y en su representación material.
Su masa será un observable de la sección así construida.

Esta secuencia determina el significado de los tres objetos que se
utilizarán: el estado genealógico conserva la extensión completa; la ruta
registra su historia observable; la representación material hace actuar
esa historia sobre un espacio de estados y permite evaluar observables.
La distinción resulta esencial incluso cuando los tres objetos comparten
una misma etiqueta electrónica o una misma coordenada del atlas.

\subsection{Extensión, historia observable y sección compatible}
\label{vi:subsec:extension-persistente}

Sean $\mathcal S_m$ los estados admisibles supervivientes hasta profundidad
$m$, obtenidos por la prolongación del TPK, y sean
$\rho_{n,m}:\mathcal S_n\to\mathcal S_m$, para $n\geq m$, sus
aplicaciones de restricción. La compatibilidad del refinamiento significa
\begin{equation}
 \rho_{m,m}=I,\qquad
 \rho_{n,m}\rho_{p,n}=\rho_{p,m}\quad(p\geq n\geq m).
 \label{vi:eq:sistema-estados}
\end{equation}
Una extensión superviviente es un elemento
\begin{equation}
 x=(x_m)_{m\geq0}\in\Omega^{\mathrm{surv}}_{\mathrm{HMT}}
 :=\varprojlim_m\mathcal S_m,
 \qquad \rho_{n,m}(x_n)=x_m.
 \label{vi:eq:extension-superviviente}
\end{equation}
La existencia y la admisibilidad de estos estados proceden de la
construcción del núcleo formal. La expresión de límite inverso describe
las compatibilidades que conserva un estado ya generado.

Sea $\mathcal R_m$ el espacio de historias observables a profundidad $m$.
El lector de ruta $\operatorname{sh}_m:\mathcal S_m\to\mathcal R_m$
y la restricción $\operatorname{tr}_{n,m}$ satisfacen
\begin{equation}
 \operatorname{tr}_{n,m}\operatorname{sh}_n
 =\operatorname{sh}_m\rho_{n,m}.
 \label{vi:eq:naturalidad-ruta}
\end{equation}
Así, el lector transmite las restricciones del estado a sus historias.
La ruta de $x$ es la colección
$\gamma_x=(\gamma_x^{[m]})_m$, donde
$\gamma_x^{[m]}=\operatorname{sh}_m(x_m)$.

\begin{proposition}[Compatibilidad de la historia expresada]
\label{vi:prop:historia-compatible}
Para toda extensión \eqref{vi:eq:extension-superviviente}, sus historias
finitas cumplen
$\operatorname{tr}_{n,m}(\gamma_x^{[n]})=\gamma_x^{[m]}$.
\end{proposition}
\begin{proof}
Aplicando \eqref{vi:eq:naturalidad-ruta} al estado $x_n$ se obtiene
\[
 \operatorname{tr}_{n,m}(\gamma_x^{[n]})
 =\operatorname{sh}_m(\rho_{n,m}x_n)
 =\operatorname{sh}_m(x_m)=\gamma_x^{[m]}.
\]
La segunda igualdad utiliza precisamente la condición de límite inverso.
\end{proof}

La proposición establece el descenso de la extensión a la historia.
Su contenido no incluye la inyectividad del lector: dos extensiones
pueden expresar una misma historia finita y conservar memorias distintas.
Por ello, al reutilizar una ruta se mantiene su estado de procedencia,
en lugar de reconstruirlo a partir de una cifra o de una palabra visible.

Sobre cada historia $\gamma_x^{[m]}$ se considera el fibrado discreto
$\mathcal B_{\gamma_x^{[m]}}$ de estados materiales de la representación
declarada. Sean
$\mathscr E_m(x)=\Gamma(\mathcal B_{\gamma_x^{[m]}})$ sus espacios
de secciones, con restricciones
$p_{n,m}:\mathscr E_n(x)\to\mathscr E_m(x)$ compatibles con las
del estado. Una sección persistente es una colección
\begin{equation}
 s=(s_m)_{m\geq m_0},\qquad s_m\in\mathscr E_m(x),\qquad
 p_{n,m}s_n=s_m\quad(n\geq m\geq m_0).
 \label{vi:eq:seccion-persistente}
\end{equation}
La profundidad inicial $m_0$ permite describir un estado desde cualquier
corte a partir del cual su representación esté definida. La información
anterior que intervenga en sus observables permanece en el registro
transportado por la extensión.

\subsection{Identidad bajo refinamiento y equivalencia de representaciones}
\label{vi:subsec:identidad-material}

Fijemos los grupos de simetrías internas $G_m$ que actúan en las fibras
materiales y preservan los observables declarados. Dos secciones
persistentes $s$ y $s'$ representan la misma identidad material cuando,
a partir de algún nivel $m_1$, existe una colección $g_m\in G_m$ tal que
\begin{equation}
 s'_m=g_ms_m,\qquad
 p_{n,m}g_n=g_mp_{n,m},
 \label{vi:eq:equivalencia-persistente}
\end{equation}
y coinciden los caracteres observables del registro genealógico. Las
simetrías admisibles se fijan antes de formar el cociente; su acción
transporta la hoja, la orientación, la memoria y la representación con
las leyes correspondientes.

\begin{proposition}[Equivalencia compatible de las secciones]
\label{vi:prop:equivalencia-persistente}
La relación \eqref{vi:eq:equivalencia-persistente}, sobre colecciones que
portan los mismos tipos de representación y observables, es una relación
de equivalencia. Sus clases son independientes de retirar un número
finito de niveles iniciales, siempre que se conserve su información
genealógica transportada.
\end{proposition}
\begin{proof}
La identidad de cada $G_m$ da reflexividad. La inversa de una colección
compatible vuelve a ser compatible: de
$p_{n,m}g_n=g_mp_{n,m}$ se deduce
$p_{n,m}g_n^{-1}=g_m^{-1}p_{n,m}$. Esto da simetría.
Si $s'=gs$ y $s''=hs'$, a partir del mayor de los dos niveles iniciales
se tiene $s''_m=h_mg_ms_m$ y
$p_{n,m}h_ng_n=h_mg_mp_{n,m}$. La composición preserva también los
observables, y prueba transitividad. Todas estas relaciones se formulan
sobre una cola común; cambiar el punto inicial dentro de esa cola
conserva la clase. La última afirmación se refiere al cambio de
presentación, no a borrar memoria del estado.
\end{proof}

\begin{definition}[Partícula material persistente]
\label{vi:def:particula}
Una partícula HMT--MD es una clase $[s]_{\mathrm{pers}}$ de secciones
\eqref{vi:eq:seccion-persistente} bajo
\eqref{vi:eq:equivalencia-persistente}, acompañada de su extensión y
registro de procedencia, que admite una representación material
irreducible o un bloque mínimo invariante y una salida espectral estable
bajo el acoplamiento especificado.
\end{definition}

La definición reúne condiciones de distintos niveles. La prolongación
y el registro pertenecen al estado HMT. La representación y el
acoplamiento pertenecen a su realización material. La estabilidad
espectral se predica del operador y del dominio utilizados por ese
acoplamiento. Esta organización permite estudiar conjuntamente identidad,
carga, espín y masa sin convertir un observable particular en definición
de todos los demás.

El retorno nonádico ilustra la distinción entre identidad y estado
instantáneo. Para el lector de fase $g$ y el registro de memoria
$\mathcal B$ pueden coexistir
\begin{equation}
 g(k+9)=g(k),\qquad
 p_{k+9,k}(s_{k+9})=s_k,\qquad
 \mathcal B_{k+9}\ne\mathcal B_k.
 \label{vi:eq:persistencia-fase-memoria}
\end{equation}
La primera igualdad es periódica; la segunda expresa persistencia; la
tercera registra el avance. El retorno de fase pertenece a la historia
de una misma identidad material, mientras el estado conserva que se ha
efectuado una vuelta adicional.

\subsection{La especie como fibra: alcance del atlas finito}
\label{vi:subsec:especie-fibra}

El atlas construido por las operaciones APP--TRIT--TPK tiene soporte
$\mathcal C_{81}=\{1,\ldots,9\}^2$. Sus cuatro orientaciones por
celda forman el dominio de $324$ rutas orientadas. La aplicación de
registro $\mathcal R$ distingue $56$ configuraciones de ruta, mientras el
clasificador interno
\begin{equation}
 q:\mathcal C_{81}\longrightarrow\Sigma_{13}
 \label{vi:eq:clasificador-especies}
\end{equation}
reúne $13$ tipos de colección y nivel. Estos censos y sus reglas se
desarrollan en el cuerpo dedicado al atlas; aquí se precisa qué objeto
transmiten a la definición de partícula.

Para $y\in\Sigma_{13}$, la fibra $q^{-1}(y)$ conserva todas sus
celdas y registros. Fijada la ordenación interna de cada fibra de
$\mathcal R$, sea $k(c)$ el rango de la celda dentro de ella. Su
presentación finita es
\begin{equation}
 \mathfrak S_y^{\mathrm{atlas}}
 =\{(\mathcal R(c),k(c)):q(c)=y\}.
 \label{vi:eq:multiseccion-atlas}
\end{equation}
El rango se añade para separar celdas que poseen el mismo registro de
colección; no selecciona una masa. El par $(\mathcal R(c),k(c))$
recupera la celda dentro de la coordenatización ordenada porque el segundo dato
individualiza su posición en la fibra del primero.

Si $\pi_{81}:\Omega^{\mathrm{surv}}_{\mathrm{HMT}}\to
\mathcal C_{81}$ es el lector celular, la elevación de la especie es
\begin{equation}
 \widetilde{\mathfrak S}_y
 =\{x\in\Omega^{\mathrm{surv}}_{\mathrm{HMT}}:
                  q(\pi_{81}x)=y\}.
 \label{vi:eq:multiseccion-enriquecida}
\end{equation}
Esta preimagen contiene extensiones, mientras
\eqref{vi:eq:multiseccion-atlas} contiene presentaciones finitas. La
construcción y el dominio del lector $\pi_{81}$ forman parte del paso
entre ambas. Una partícula concreta se realiza como sección persistente
dentro de la multisección enriquecida correspondiente.

\begin{theorem}[Punto fijo central y selección equivariante]
\label{vi:thm:selector-central}
Sea $\rho_c(r,s)=(10-r,10-s)$ la inversión celular. En las fibras
$\rho_c$-estables del clasificador \eqref{vi:eq:clasificador-especies},
con acción trivial sobre la etiqueta $y$, una selección puntual
equivariante sólo puede tomar valores en la celda $(5,5)$. La fibra
electrónica de nivel cero contiene ese único punto fijo; ninguna otra
fibra admite tal selección.
\end{theorem}
\begin{proof}
Resolver $\rho_c(r,s)=(r,s)$ da $r=s=5$. Si $a(y)\in q^{-1}(y)$
es una selección equivariante y $y$ queda fijo, entonces
$a(y)=a(\rho_c y)=\rho_c a(y)$. Su valor debe ser el punto fijo
$(5,5)$. El atlas sitúa esa celda en la fibra electrónica de nivel
cero, cuyo cardinal es uno. Para cualquier otra fibra, la condición
de selección produciría un punto fijo que no pertenece a ella.
\end{proof}

La consecuencia es constructiva: las demás especies se presentan por
órbitas o multisecciones, y una realización que necesite una línea
particular debe declarar su representación y su acoplamiento. El
electrón proporciona el caso central de rango uno sin convertir todas
las colecciones en copias de ese caso. Además, la celda $(5,5)$, la firma
central cruda, la firma regional $555555$ y el origen relativo de una
recta de acción conservan sus dominios propios. La unicidad del punto
no identifica esos registros entre sí.

\subsection{Del registro material al observable de masa}
\label{vi:subsec:particula-observable}

La cadena utilizada por el observable de masa conserva el orden
\begin{equation}
 \text{extensión}\longrightarrow\text{ruta}\longrightarrow
 \text{registro}\longrightarrow\text{firma entera}\longrightarrow
 \text{carácter}\longrightarrow\text{torres}\longrightarrow
 \text{operador material}.
 \label{vi:eq:cadena-material}
\end{equation}
Cada flecha retiene los datos necesarios para la siguiente. La firma es
una lectura del registro; el carácter evalúa esa firma con los
coeficientes internos ya construidos; las torres incorporan su
refinamiento; el operador actúa finalmente en la representación de la
fibra seleccionada por el acoplamiento. Su salida puede ser escalar,
un multiplete o una distribución espectral, según esa representación.
La comparación metrológica tiene lugar después de esta construcción.

Una especie queda descrita por carga, representación de espín, color,
sabor, generación, configuración de ruta y operador material. Estos datos
explican por qué dos preparaciones con una masa coincidente pueden
conservar distinta información, y por qué una misma identidad puede
expresar observables diferentes bajo acoplamientos distintos. La
identidad reside en la clase persistente y en sus transportes admisibles;
la medida es una lectura de esa identidad en el codominio declarado.

La misma jerarquía separa una extensión persistente de una sección de
horizonte finito. Si $\operatorname{Ext}_L(s)\ne\varnothing$ y
$\operatorname{Ext}_{L+1}(s)=\varnothing$, el horizonte $L$ registra
una obstrucción de prolongación. Una participación intermedia de este
tipo en una composición no crea, por sí misma, una nueva especie
asintótica. Análogamente, una resonancia exige el operador de
supervivencia y su polo, además de la identidad de la sección que
transporta. Se preserva así la diferencia entre persistencia,
prolongación finita y evaluación espectral.

\section{Conjugación, orientación y retorno de las secciones materiales}
\label{vi:sec:conjugacion}

La identidad persistente de la sección~\ref{vi:sec:particula} admite
operaciones de inversión sobre niveles diferentes de su estructura.
La inversión celular actúa en el atlas; la conjugación dodecafásica
actúa en el registro de orientación; el levantamiento de un lazo actúa
en una fibra; una rotación espinorial actúa en un espacio de
representación. Especificar esos dominios permite componer las
operaciones y explicar qué información se conserva en cada retorno.

\subsection{Involución dodecafásica y paridad de los observables}
\label{vi:subsec:conjugacion-registro}

Sea $\mathcal T_{12}=\{-1,0,1\}^{12}$ el ambiente de registros
tríticos de doce posiciones. El dominio admisible del TPK es el
subconjunto con las compatibilidades de ruta ya construidas. La
operación sobre el ambiente es
\begin{equation}
 C_M:\mathcal T_{12}\longrightarrow\mathcal T_{12},\qquad
 (C_M\tau)_j=-\tau_{11-j},\quad 0\leq j<12.
 \label{vi:eq:cm}
\end{equation}
La restricción a registros admisibles utiliza la estabilidad de ese
dominio bajo la conjugación. El orden invertido y el cambio de signo
son componentes simultáneas de la operación.

\begin{proposition}[Paridad lineal y cuadrática del registro]
\label{vi:prop:paridad-registro}
Se cumple $C_M^2=I$. Si $a_j=a_{11-j}$ y
$L_a(\tau)=\sum_ja_j\tau_j$, entonces
$L_a(C_M\tau)=-L_a(\tau)$. Si una forma cuadrática
$Q_b(\tau)=\sum_{j,k}b_{jk}\tau_j\tau_k$ satisface
$b_{jk}=b_{11-j,11-k}$, entonces $Q_b(C_M\tau)=Q_b(\tau)$.
\end{proposition}
\begin{proof}
En cada posición,
$[C_M(C_M\tau)]_j=-(-\tau_{11-(11-j)})=\tau_j$.
Para el funcional lineal, sustituir $k=11-j$ en la suma da
$L_a(C_M\tau)=-\sum_ka_{11-k}\tau_k=-L_a(\tau)$.
En la expresión cuadrática los dos cambios de signo se cancelan y
la sustitución $(r,s)=(11-j,11-k)$ recupera $Q_b$ por la
simetría de sus coeficientes.
\end{proof}

La orientación lineal y las correlaciones de memoria tienen, por tanto,
leyes distintas. Incluso una composición de estos observables debe
conservar sus paridades: si $L_a$ es impar, su carácter exponencial
satisface
\begin{equation}
 \exp(L_a(C_M\tau))=\exp(L_a(\tau))^{-1}.
 \label{vi:eq:caracter-inverso}
\end{equation}
La igualdad expresa inversión multiplicativa. La invariancia de un
operador material completo se estudia sobre ese operador, con sus
canales y su representación, y se demuestra por la covariancia
correspondiente.

La inversión celular del teorema~\ref{vi:thm:selector-central} es
$\rho_c:\mathcal C_{81}\to\mathcal C_{81}$, mientras $C_M$ tiene
el dominio \eqref{vi:eq:cm}. Ambas son involuciones, pero para
compararlas se requiere el mapa de registro de la ruta y su ley de
transporte. Este dato evita sustituir la conjugación de una historia
por la inversión de una coordenada aislada del atlas.

\subsection{Levantamiento a antisecciones y conservación de la masa}
\label{vi:subsec:antisecciones}

La conjugación se aplica a la partícula mediante dos mapas
compatibles. El mapa de base $C_m^{\mathrm b}$ transporta el estado
$x_m$ a su estado conjugado $\bar x_m$; el levantamiento
$\widetilde C_m:\mathscr E_m(x)\to\mathscr E_m(\bar x)$
transporta sus secciones materiales. Sean $t_m$ el lector del
registro dodecafásico y $\bar p_{n,m}$ las restricciones en la
colección conjugada. Sus condiciones son conservación de la
admisibilidad, involutividad y compatibilidad:
\begin{equation}
 \begin{aligned}
 \rho_{n,m}C_n^{\mathrm b}&=C_m^{\mathrm b}\rho_{n,m},
 &t_m(C_m^{\mathrm b}x_m)&=C_Mt_m(x_m),\\
 \bar p_{n,m}\widetilde C_n&=\widetilde C_mp_{n,m},
 &\widetilde C_{m,\bar x}\widetilde C_{m,x}&=I.
 \end{aligned}
 \label{vi:eq:levantamiento-cm}
\end{equation}
En la primera fila, el lector expresa el registro de orientación
de la historia. El transporte restante del estado acompaña esa
representación. La segunda fila da las leyes del levantamiento;
los subíndices $x,\bar x$ explicitan sus dos fibras de partida.
Si $s$ es persistente, su antisección es
$\bar s=(\widetilde C_ms_m)_m$.

\begin{proposition}[Persistencia de la antisección]
\label{vi:prop:antiseccion-persistente}
Las condiciones \eqref{vi:eq:levantamiento-cm} transportan secciones
persistentes a secciones persistentes y restituyen la sección al
conjugar dos veces. Si además transportan las simetrías internas por
$g_m\mapsto\widetilde C_mg_m\widetilde C_m^{-1}$, inducen una
involución en las clases materiales de la
definición~\ref{vi:def:particula}.
\end{proposition}
\begin{proof}
La compatibilidad da
$\bar p_{n,m}\bar s_n=\widetilde C_mp_{n,m}s_n
=\widetilde C_ms_m=\bar s_m$.
La involutividad restituye $s$. Para $s'_m=g_ms_m$, sus
antisecciones se relacionan mediante
$\widetilde C_mg_m\widetilde C_m^{-1}$. La ley de restricción
en \eqref{vi:eq:levantamiento-cm} conserva la compatibilidad de esta
colección, que pertenece a los grupos internos por la condición del
enunciado. Así la construcción desciende a las clases.
\end{proof}

La representación de carga utiliza la parte impar del registro
material. Si $\mathfrak q$ es el carácter de carga declarado, su
ley es $\mathfrak q(\bar s)=-\mathfrak q(s)$. La carga elemental
que interviene al expresar esa coordenada en una recta dimensional es
la salida interna ya construida; la representación se aplica después
al estado, sin utilizar una carga objetivo para elegir su registro.

Para la masa se necesita la covariancia de su operador. Sean
$\widehat M_s$ y $\widehat M_{\bar s}$ los operadores autoadjuntos
de las dos realizaciones, y sea $\mathcal C$ su entrelazador
unitario o antiunitario. Su dominio se transporta mediante
$\mathcal C\operatorname{Dom}(\widehat M_s)
=\operatorname{Dom}(\widehat M_{\bar s})$. La condición operatoria es
\begin{equation}
 \mathcal C\widehat M_s\mathcal C^{-1}
 =\widehat M_{\bar s}.
 \label{vi:eq:covariancia-masa}
\end{equation}

\begin{proposition}[Conservación espectral bajo conjugación]
\label{vi:prop:masa-conjugada}
Bajo \eqref{vi:eq:covariancia-masa}, el transporte de una sección
propia de masa $m\in\mathbb R$ conserva ese autovalor. La medida
espectral se transporta mediante
$E_{\bar s}(B)=\mathcal C E_s(B)\mathcal C^{-1}$ para todo
boreliano real $B$. En particular, una realización escalar conserva
la masa al pasar a la antisección.
\end{proposition}
\begin{proof}
Si $\widehat M_s\psi=m\psi$, la covariancia da
$\widehat M_{\bar s}\mathcal C\psi
=\mathcal C(m\psi)=m\mathcal C\psi$; la última igualdad
vale también en el caso antiunitario porque $m$ es real. El cálculo
funcional real de los operadores autoadjuntos transporta sus
proyectores espectrales y prueba la segunda afirmación.
\end{proof}

En un registro de banda puede verse la misma condición en forma
elemental. Dada una lectura real $E_0(\tau)$, sean
$E_\pm=(E_0\pm E_0\circ C_M)/2$ y sea $\xi\in\{+1,-1\}$
la orientación de hoja. Entonces
\begin{equation}
 E_{\mathrm{band}}(\tau,\xi)=E_+(\tau)+\xi E_-(\tau),\qquad
 E_{\mathrm{band}}(C_M\tau,-\xi)=E_{\mathrm{band}}(\tau,\xi).
 \label{vi:eq:banda-paridad}
\end{equation}
La igualdad se obtiene usando que $E_+$ es par y $E_-$ impar.
Explica cómo puede conservarse la energía de banda al transportar
simultáneamente registro y hoja; su identificación con la masa de una
realización concreta utiliza \eqref{vi:eq:covariancia-masa}.

El registro electrónico
$\tau_e=(+,+,+,-,-,-,+,+,+,-,-,-)$ satisface $C_M\tau_e=\tau_e$.
Esta igualdad concierne a la palabra de doce posiciones. El par
$(\tau_e,\xi)$ conserva una orientación adicional que la
conjugación cambia. La distinción permite mantener el centro
electrónico único y, a la vez, las realizaciones conjugadas de carga;
una palabra fija no se convierte en una prueba de neutralidad.

\subsection{Holonomía de signo y banda de persistencia}
\label{vi:subsec:mobius}

Sea $\ell$ un lazo de retorno del lector material de una extensión
superviviente. Su grupo de vueltas es $\langle\ell\rangle\simeq
\mathbb Z$. Sobre este grupo se declara el carácter de signo
\begin{equation}
 \varepsilon:\langle\ell\rangle\longrightarrow\{+1,-1\},
 \qquad \varepsilon(\ell)=-1.
 \label{vi:eq:holonomia-signo}
\end{equation}
El grupo de vueltas registra cuántas veces se recorre el lazo; su
tipo es distinto de la fase finita $C_9$. El fibrado de intervalos
asociado al carácter es
\begin{equation}
 \mathcal M_\ell=([0,1]\times[-1,1])/
                  \bigl((0,u)\sim(1,-u)\bigr).
 \label{vi:eq:banda-mobius}
\end{equation}

\begin{theorem}[Retorno orientado sobre la banda]
\label{vi:thm:retorno-mobius}
El fibrado \eqref{vi:eq:banda-mobius} es una banda de Möbius.
Su transporte sobre $\ell$ cambia el signo de la fibra y sobre
$\ell^2$ lo restituye. Por tanto, un vector no nulo de esa línea
requiere dos vueltas para su retorno orientado. Esto conserva la
distinción entre retorno de orientación y retorno del estado completo.
\end{theorem}
\begin{proof}
La relación de extremos de \eqref{vi:eq:banda-mobius} es el mapa
lineal $u\mapsto-u$, con determinante negativo. Da la línea real
no orientable sobre el círculo y su correspondiente banda de
intervalos. Componer dos transportes produce $u\mapsto u$.
La memoria de la extensión no forma parte de ese cociente de signo;
su actualización continúa obedeciendo el transporte del TPK y
\eqref{vi:eq:persistencia-fase-memoria}.
\end{proof}

Una parametrización de una vuelta de la base mediante $2\pi$
representa la restitución de la cubierta mediante $4\pi$. Si el
lector temporal asigna $108$ pasos a la vuelta relevante, el
levantamiento orientado cierra a $216$ pasos. Estas asignaciones
requieren el lazo y el lector indicados; no dividen la historia en
dos partículas sucesivas de $108$ pasos.

\begin{proposition}[Descomposición de observables en la cubierta doble]
\label{vi:prop:observables-paridad}
Sea $p:\widetilde X\to X$ una cubierta principal doble, con
involución $\iota$. Para $f\in C^0(\widetilde X,\mathbb R)$,
\[
 f^+=\frac{f+f\circ\iota}{2},\qquad
 f^-=\frac{f-f\circ\iota}{2}
\]
dan $f=f^++f^-$, con $f^+\circ\iota=f^+$ y
$f^-\circ\iota=-f^-$. La primera parte desciende a $X$; la
segunda representa una sección de la línea de signo asociada. En
la línea de Möbius, toda sección global continua tiene algún cero.
\end{proposition}
\begin{proof}
Las identidades se obtienen sustituyendo $\iota^2=I$. La parte
invariante es constante en cada fibra y define una función en el
cociente. La parte antiinvariante satisface la ley de transición
de la representación de signo. En el círculo, una sección de esta
línea se representa por una función real continua $u$ en $[0,1]$
con $u(1)=-u(0)$. Si el extremo vale cero ya hay un cero; en caso
contrario, los extremos tienen signos opuestos y el teorema del
valor intermedio proporciona un cero interior.
\end{proof}

La descomposición distingue dos clases de observables sobre una
misma cubierta. Permite representar información impar de orientación
y datos pares conservados. Para aplicarla a un registro material se
utiliza el levantamiento \eqref{vi:eq:levantamiento-cm}; el carácter
de signo describe la línea transportada, no todos los datos de la
partícula.

\subsection{Retículo bidireccional y conjugación de los canales}
\label{vi:subsec:reticulo-bidireccional}

El registro angular entero expresa dos coordenadas
\begin{equation}
 W_+=6n_A+s_C,\qquad W_-=6n_A-s_C,
 \qquad (n_A,s_C)\in\mathbb Z^2.
 \label{vi:eq:reticulo-ida-retorno}
\end{equation}
Los coeficientes $n_A,s_C$ son lecturas internas de ruta. El factor
$6$ pertenece a la normalización dodecafásica de ese registro.

\begin{proposition}[Índice del retículo e intercambio de las lecturas]
\label{vi:prop:indice-bidireccional}
La imagen de \eqref{vi:eq:reticulo-ida-retorno} tiene índice $12$
en $\mathbb Z^2$ y forma normal de Smith $\operatorname{diag}(1,12)$.
Se caracteriza por $W_++W_-\equiv0\pmod{12}$. La operación
$s_C\mapsto-s_C$ intercambia $W_+$ y $W_-$.
\end{proposition}
\begin{proof}
La matriz del mapa es
$\bigl(\begin{smallmatrix}6&1\\6&-1\end{smallmatrix}\bigr)$.
El máximo divisor común de sus entradas es $1$ y su determinante
es $-12$, lo que da los factores invariantes indicados. Las
coordenadas inversas son
$n_A=(W_++W_-)/12$ y $s_C=(W_+-W_-)/2$.
Si la suma es divisible por $12$, los dos enteros tienen la misma
paridad y ambas expresiones son enteras. La condición es también
necesaria por la suma de \eqref{vi:eq:reticulo-ida-retorno}.
La afirmación sobre el intercambio se verifica sustituyendo $-s_C$.
\end{proof}

En la representación angular real, las coordenadas ya generadas
$a=A_{\mathrm{rad}}$ y $b=C^*_{\mathrm{rad}}$ producen los
canales $q_\pm=\exp[-(a\pm b)]$. En el dominio $a>|b|$ ambos
pertenecen a $(0,1)$. Sea $R^*=R$, $R^2=I$ la graduación de
hojas y sea $C$ un intercambio unitario con $CRC^{-1}=-R$.
Para $T(b)=\exp[-(aI+bR)]$ se obtiene
\begin{equation}
 CT(b)C^{-1}=T(-b),\qquad
 CP_\pm C^{-1}=P_\mp,\quad P_\pm=(I\pm R)/2.
 \label{vi:eq:conjugacion-canales}
\end{equation}
En efecto, la conjugación transforma el exponente y cada término
de su serie convergente. Esta prueba transmite el intercambio a
los canales. El retículo entero y el operador real tienen dominios
distintos y expresan la misma operación de intercambio a través de
sus mapas de lectura declarados.

\subsection{Representación espinorial del bloque electrónico}
\label{vi:subsec:spin-electron}

La construcción APP tridimensional de los proyectores de suma y
producto selecciona un plano principal electrónico por el modo
trítico $(1,-1,0)^3$. Su derivación y la evaluación de los
proyectores se incluyen entre los antecedentes del presente
artículo. En una base ortonormal de ese plano, sus restricciones
son
\begin{equation}
 P=\begin{pmatrix}1&0\\0&0\end{pmatrix},\qquad
 Q=\begin{pmatrix}1/4&\sqrt3/4\\\sqrt3/4&3/4\end{pmatrix},
 \qquad S=2P-I,\quad J=\frac4{\sqrt3}[P,Q].
 \label{vi:eq:bloque-electronico}
\end{equation}
Esta representación recibe los dos proyectores construidos y
conserva su orientación. La multiplicación da
$S^*=S$, $J^*=-J$, $S^2=I$, $J^2=-I$ y $SJ=-JS$.
En concreto, $S=\operatorname{diag}(1,-1)$ y
$J=\bigl(\begin{smallmatrix}0&1\\-1&0\end{smallmatrix}\bigr)$.

Sea $\mathscr S_e=E_e\otimes_{\mathbb R}\mathbb C$ la
complejificación del plano real. La unidad escalar $i$ pertenece
al cuerpo de esta representación; $J$ sigue siendo el
endomorfismo real procedente de los proyectores.

\begin{proposition}[Terna hermítica y representación de rotaciones]
\label{vi:prop:spin}
En $\mathscr S_e$, defínanse
\begin{equation}
 \sigma_1=SJ,\qquad \sigma_2=-iJ,\qquad \sigma_3=S.
 \label{vi:eq:terna}
\end{equation}
Estos operadores son hermíticos, de traza nula, y satisfacen
\begin{equation}
 \sigma_a\sigma_b=\delta_{ab}I+
 i\sum_c\epsilon_{abc}\sigma_c,
 \qquad L_a=\sigma_a/2,\qquad
 [L_a,L_b]=i\sum_c\epsilon_{abc}L_c,
 \quad\sum_aL_a^2=\tfrac34I.
 \label{vi:eq:spin-algebra}
\end{equation}
Constituyen la representación irreducible de espín $1/2$ sobre
el bloque electrónico.
\end{proposition}
\begin{proof}
Las relaciones de $S,J$ dan $(SJ)^*=SJ$, $(-iJ)^*=-iJ$ y
$\sigma_a^2=I$. Los productos cíclicos son
$(SJ)(-iJ)=iS$, $(-iJ)S=iSJ$ y $S(SJ)=J=i(-iJ)$;
invertir el orden cambia el signo. Esto prueba la identidad de
producto y los conmutadores. En la base de
\eqref{vi:eq:bloque-electronico} las tres matrices son
\[
 \begin{pmatrix}0&1\\1&0\end{pmatrix},\quad
 \begin{pmatrix}0&-i\\i&0\end{pmatrix},\quad
 \begin{pmatrix}1&0\\0&-1\end{pmatrix}.
\]
Su traza es cero y, junto con $I$, generan $M_2(\mathbb C)$.
Un subespacio invariante por las tres es invariante por toda esa
álgebra y sólo puede ser $0$ o $\mathscr S_e$. Finalmente,
sumar $L_a^2=I/4$ da el valor $3I/4$ del operador cuadrático.
\end{proof}

El reconocimiento de esta terna como matrices de Pauli es posterior
a \eqref{vi:eq:bloque-electronico}. La norma y el producto
$\langle X,Y\rangle=\tfrac12\operatorname{tr}(XY)$ en el
espacio real de matrices hermíticas de traza nula dan su espacio
de direcciones. Para $n\in\mathbb R^3$, $|n|=1$, sean
\begin{equation}
 H_e(n)=\sum_an_a\sigma_a,\qquad
 h_e(n)=\tfrac12H_e(n),\qquad
 \Pi_\pm(n)=\tfrac12(I\pm H_e(n)).
 \label{vi:eq:helicidad}
\end{equation}

\begin{proposition}[Descomposición direccional y doble retorno]
\label{vi:prop:helicidad}
Los operadores $\Pi_\pm(n)$ son proyectores ortogonales de
rango uno y $h_e(n)\Pi_\pm(n)=\pm\tfrac12\Pi_\pm(n)$.
El grupo de rotaciones alrededor de $n$ es
\begin{equation}
 U_n(\theta)=e^{-i\theta H_e(n)/2}
 =\cos(\theta/2)I-i\sin(\theta/2)H_e(n),\qquad
 U_n(2\pi)=-I,\quad U_n(4\pi)=I.
 \label{vi:eq:retorno-spin}
\end{equation}
\end{proposition}
\begin{proof}
La parte antisimétrica de \eqref{vi:eq:spin-algebra} se anula
al contraerse con $n_an_b$, por lo que $H_e(n)^2=I$.
De aquí se siguen la idempotencia, la ortogonalidad y la suma
$\Pi_++\Pi_-=I$. Son hermíticos y de traza uno, luego su
rango es uno. Multiplicar por $H_e(n)/2$ da los autovalores.
Separar las potencias pares e impares de la exponencial produce
\eqref{vi:eq:retorno-spin} y su ley de grupo. La normalización
angular se comprueba por
\[
 U_n(\theta)H_e(v)U_n(\theta)^*
 =H_e\bigl(v\cos\theta+(n\times v)\sin\theta
       +n(n\cdot v)(1-\cos\theta)\bigr),
\]
que resulta de multiplicar con la identidad de producto de la
terna. La acción inducida en las direcciones es una rotación de
ángulo $\theta$; el levantamiento espinorial lleva su mitad.
\end{proof}

Si $n$ se realiza como dirección de propagación, $h_e(n)$ es
la helicidad adimensional. La inversión de esa dirección cumple
$H_e(-n)=-H_e(n)$ y $\Pi_\pm(-n)=\Pi_\mp(n)$.
Esta operación cambia las etiquetas direccionales; la conjugación
de carga transporta el registro material mediante
\eqref{vi:eq:levantamiento-cm}. La quiralidad se define mediante
la graduación relativista que corresponda a la representación y
sus proyectores, distintos de \eqref{vi:eq:helicidad}.

La evaluación escalar electrónica actúa como
$\varepsilon_e^\beta I$ sobre $\mathscr S_e$ y conmuta con
la terna y con $\Pi_\pm(n)$. La información rotacional conserva
la estructura de esa evaluación sin introducir otro factor en
la masa. Por su parte, la banda de Möbius transporta una línea
real alrededor de un lazo y la holonomía nonádica transporta
fase y memoria. Los tres retornos pueden intervenir en una
misma realización; sus dominios y entrelazadores determinan la
composición precisa.

\section{Del registro de trayectoria al operador de masa}
\label{vi:sec:operador-masa}

La construcción distingue dos operaciones. Primero, la dinámica produce
un registro de la trayectoria. Después, un funcional extrae coordenadas
enteras y las evalúa mediante las constantes internas ya obtenidas.
Conservar ese orden explica cómo la masa puede ser una magnitud de la
persistencia sin agotar su estructura.

\subsection{Lectura cronológica y eventos orientados}

Sea \(d_1,\ldots,d_{108}\) la lectura digital de una ruta enriquecida.
En la coordenatización declarada, el soporte de acción es \(\{1,3,5,7\}\), y
\[
 n_A=\#\{t:d_t\in\{1,3,5,7\}\}.
\]
Se conservan además los \(12\) eventos de los bloques de \(9\) pasos:
\[
 e_m=\sigma_m\#\{t:9m+1\le t\le9m+9,\ d_t\in\{2,4,6,8\}\},
 \quad 0\le m<12,
\]
donde \(\sigma=(+,+,+,-,-,-,+,+,+,-,-,-)\). El soporte de este
extractor es una parte explícita de la coordenatización de lectura; su función no
debe confundirse con el selector ternario de fase del TPK.

La distribución de los eventos es relevante incluso cuando sus sumas
coinciden. Para verla, introducimos
\[
 p_i=e_i+e_{i+6},\quad a_i=e_i-e_{i+6},\quad
 s_C=\sum_{i=0}^{5}p_i,\quad M_i=p_i-p_5\ (0\le i<5).
\]

\begin{proposition}[Reconstrucción integral de los eventos]
\label{vi:prop:eventos}
La transformación \((e_0,\ldots,e_{11})\mapsto(s_C,M_0,\ldots,M_4,a_0,\ldots,a_5)\)
es inyectiva sobre \(\mathbb Z^{12}\). Su imagen se caracteriza por
\[
 s_C-\sum_{i=0}^4M_i\equiv0\pmod6,
 \qquad p_i\equiv a_i\pmod2\quad(0\le i<6),
\]
donde \(p_5=(s_C-\sum M_i)/6\) y \(p_i=p_5+M_i\) para \(i<5\).
\end{proposition}
\begin{proof}
Sumar las seis expresiones de \(p_i\) da
\(s_C=6p_5+\sum_{i<5}M_i\). Esto recupera los \(p_i\), y después
\[
 e_i=\frac{p_i+a_i}{2},\qquad
 e_{i+6}=\frac{p_i-a_i}{2}.
\]
Las congruencias son necesarias para que esas cantidades sean enteras y
suficientes para que las fórmulas definan una preimagen integral. Sustituir
en la transformación directa recupera todas las coordenadas.
\end{proof}

La descomposición \(12=1+5+6\) adquiere así un significado operativo:
el total, las diferencias y las orientaciones recuperan los eventos.
La división por \(6\) pertenece a esa reconstrucción. Cuando el carácter
usa \(s_C/6\), \(s_C\) sigue siendo la suma entera; una coordenada ya
dividida no se divide una segunda vez.

\subsection{Memoria torsional y firma}

Para cada aparición de \(9\), la lectura conserva su predecesor.
Definimos, con índices cíclicos,
\[
 j_2(d)=\begin{cases}0,&d=9,\\1,&d\in\{2,4,6,8\},\\-1,&d\in\{1,3,5,7\},\end{cases}
 \quad z_t=\mathbf1_{\{d_t=9\}}j_2(d_{t-1}).
\]
Entonces
\[
 T_z=\sum_{t=1}^{108}z_t,\quad
 C_z=\sum_{t\in\{27,54,81,108\}}z_t,\quad
 k_{\Delta_4}=T_z-2C_z.
\]
La primera suma conserva la memoria acumulada y la segunda distingue los
cierres de corona. Eliminar los predecesores antes de evaluar estas
expresiones altera el objeto que se lee.

Con \(\tau_m=\operatorname{sgn}(e_m)\), las restantes coordenadas son
\[
 b_{120}=\sum_{a=0}^{3}\operatorname{sgn}(e_a+e_{a+4}+e_{a+8}),
 \qquad b_{\zeta}=\sum_{m=0}^{11}\tau_m\tau_{m+3},
\]
donde el segundo índice es módulo \(12\). Se obtiene la firma
\[
 L_{\mathrm{tick}}(\gamma)
 =\sigma(\gamma)=(n_A,s_C,k_{\Delta_4},b_{120},b_{\zeta})\in\mathbb Z^5.
\]
Esta aplicación es determinista en la coordenatización declarada. Los signos y las
correlaciones explican por qué la composición debe realizarse sobre el
registro antes de extraer los \(5\) enteros: la suma de firmas no
representa, en general, la composición de dos historias.

\subsection{Carácter y especialización interna}

Sea \(X=(X_A,X_C,X_4,X_{120},X_{270})\). El carácter formal es
\[
 \chi(\sigma;X)=\exp\!\left(n_AX_A+s_CX_C+k_{\Delta_4}X_4+
 b_{120}X_{120}+b_{\zeta}X_{270}\right).
\]
Para cada \(X\), satisface
\(\chi(\sigma+\sigma';X)=\chi(\sigma;X)\chi(\sigma';X)\)
porque el exponente es lineal en la firma. Es una ley sobre firmas; no
afirma que el extractor sea aditivo sobre rutas sin conservar la interfaz.

La especialización emplea las coordenadas internas construidas previamente:
\[
 X_A=A_{\mathrm{rad}},\quad X_C=C^*_{\mathrm{rad}}/6,
 \quad X_4=\Delta_4,\quad X_{120}=s_{120},\quad
 X_{270}=\zeta_{270}:=135\Delta_4^2.
\]
El carácter es adimensional y positivo. Su efecto es una homotecia sobre
la recta de masa. La forma exponencial determina cómo actúan los
coeficientes; su procedencia queda fijada en los registros y construcciones
anteriores.

\subsection{Jerarquía bicapa y conservación de las dos genealogías}

Escribimos \(x=A_{\mathrm{rad}}\), \(y=C^*_{\mathrm{rad}}\). En la
cámara \(x>y>0\), los canales son
\[
 q_+=e^{-x-y},\qquad q_-=e^{-x+y},\qquad0<q_+<q_-<1.
\]
Para una involución \(R\), con \(P_\pm=(I\pm R)/2\), el operador
de los dos canales es
\[
 T=q_+P_++q_-P_-.
\]
Sus momentos escalar y orientado se expresan como
\[
 c_n=q_-^n+q_+^n,\qquad s_n=q_-^n-q_+^n.
\]
Multiplicar estas expresiones prueba
\[
 c_n^2-s_n^2=4e^{-2nx},\quad
 c_{m+n}=\frac{c_mc_n+s_ms_n}{2},\quad
 s_{m+n}=\frac{s_mc_n+c_ms_n}{2}.
\]
Ambas sucesiones satisfacen
\[
 X_{n+2}=c_1X_{n+1}-e^{-2x}X_n,\qquad X\in\{c,s\}.
\]
Estas relaciones muestran que las alturas comparten un mismo operador.
No se define una colección distinta de momentos para cada especie.

\begin{proposition}[Semigrupo de alturas]
\[
 \langle90,120\rangle=\{90,120\}\cup\{30r:r\ge6\}.
\]
\end{proposition}
\begin{proof}
Tras dividir por \(30\), basta estudiar \(\langle3,4\rangle\).
Los enteros \(6,7,8\) son \(3+3,3+4,4+4\); sumar \(3\) produce
todos los enteros posteriores. Por debajo de \(6\), los únicos
elementos positivos son \(3,4\).
\end{proof}

Las alturas \(360=4\cdot90=3\cdot120\) ilustran la diferencia entre
igualdad de evaluación y memoria de composición. El cociente algebraico
identifica los dos exponentes; el registro conserva si el transporte
empleó cuatro etapas de un tipo o tres del otro.

La especialización refinada tiene exponente
\[
 \log\chi_{\mathrm{ref}}=n_Ax+\frac{s_C}{6}y+k_{\Delta_4}\Delta_4
 +b_{\zeta}\zeta_{270}+\sum_{h\in\langle90,120\rangle}N_hs_h,
 \qquad N_h=\eta_h+b_{120}\mathbf1_{\{h=120\}}.
\]
El par \((\eta_{120},b_{120})\) se conserva aunque la evaluación use
su suma. También se conservan separados \(s_{270}\) y
\(\zeta_{270}\): el primero es un momento bicapa y el segundo una
expresión cuadrática del refinamiento.

\begin{proposition}[Convergencia de la evaluación refinada]
Si \(\sum_h|N_hs_h|<\infty\), el carácter refinado define un escalar
positivo finito y sus evaluaciones truncadas convergen a él.
\end{proposition}
\begin{proof}
La hipótesis da convergencia absoluta del exponente. La continuidad de la
exponencial permite pasar al límite; la exponencial de un real finito es
estrictamente positiva.
\end{proof}

La convergencia evalúa una sucesión de coeficientes ya especificada.
La selección de esos coeficientes por la ruta es una operación distinta,
y debe acompañar a toda evaluación atribuida a una especie.

La notación \(\chi_{\mathrm{full}}\) conservada en algunas fórmulas
de evaluación de los sectores designa este mismo carácter
\(\chi_{\mathrm{ref}}\), con los mismos argumentos y coeficientes.
El cambio de símbolo no añade otra especialización ni otra torre.

\subsection{Lectores de desplazamiento y autocorrelación}

La misma ruta admite dos lecturas del retorno:
\[
 K_D=\left(D_{27}+D_{36},D_1,\frac{D_4-D_3}{2}\right),\qquad
 K_\Omega=(\Omega,54,P_6).
\]
El primer mapa conserva defectos de desplazamiento; el segundo, memoria de
retorno y autocorrelación. En el censo completo de \(324\) orientaciones
hay \(14\) perfiles del primero, \(9\) del segundo y \(40\) pares
conjuntos. Las \(18\) coincidencias tienen valor \((80,54,6)\).
No se infiere de ello la igualdad de ambos operadores en todo el atlas.

Para la ruta electrónica, los desplazamientos son
\[
 (D_1,D_3,D_4,D_{27},D_{36})=(54,56,68,48,32).
\]
La sustitución produce \(K_D=(80,54,6)\), coincidente con
\(K_\Omega\). Su exponente es
\[
 \kappa_e=80-54\alpha+6\Delta_4.
\]
La trayectoria electrónica se construye en
\eqref{vi:dep:trayectoria-electron}; sus lectores de desplazamiento y
autocorrelación se evalúan, respectivamente, en
\eqref{vi:dep:KD} y \eqref{vi:dep:KOmega}. La coincidencia
\(K_D=K_\Omega=(80,54,6)\) pertenece al teorema de cierre beta de la
subsección~\ref{vi:dep:beta}. Aquí esa excepción central se conserva y la
operación se extiende a las fibras del atlas.

\subsection{Operador sobre una fibra y sección dimensional}

Sea \(\mathscr H_M\) el espacio de Hilbert finito con base ortonormal
\((e_\gamma)\) etiquetada por las rutas de una multisección. Para
\(r\in\{D,\Omega\}\), el operador de retorno se define por
\(B_M^re_\gamma=\kappa_\gamma^re_\gamma\), con exponentes reales;
por tanto es autoadjunto. Sea \(M_M^{(0)}\) el operador basal positivo,
cuya conexión con el carácter se fija a continuación.
La razón adimensional de las
secciones de acción es \(R_{\mathrm{act}}>0\). Se define
\[
 M_M^{\beta,r}=R_{\mathrm{act}}^{B_M^r/2}
 M_M^{(0)}R_{\mathrm{act}}^{B_M^r/2}.
\]

\begin{proposition}[Positividad y restricción escalar]
El operador \(M_M^{\beta,r}\) es positivo. Si ambos operadores son
diagonales en la base de rutas, cada valor basal \(m_\gamma^{(0)}\)
se transforma en
\[
 m_\gamma^\beta=m_\gamma^{(0)}R_{\mathrm{act}}^{\kappa_\gamma}.
\]
\end{proposition}
\begin{proof}
Sea \(S=R_{\mathrm{act}}^{B_M^r/2}\), autoadjunto e invertible.
Para todo \(v\),
\(\langle v,SM_M^{(0)}Sv\rangle=\langle Sv,M_M^{(0)}Sv\rangle\ge0\).
En la base diagonal, los dos factores \(S\) multiplican cada valor
por \(R_{\mathrm{act}}^{\kappa_\gamma/2}\); su producto da la fórmula.
\end{proof}

La sección absoluta de masa se obtiene de la acción y el reloj:
\[
 \omega_0=\frac{2\pi}{108t_0},\qquad
m_{\mathrm{clk}}=\hbar_{\mathrm{ret}}\omega_0c_{\mathrm{int}}^{-2}.
\]
En la coordenatización normalizada respecto del electrón, el operador basal anterior es
\[
 M_M^{(0)}e_\gamma
 =m_e^{(0)}\,
   \chi_{\mathrm{ref}}(\sigma_\gamma-\sigma_e,
                       \eta_\gamma-\eta_e)e_\gamma.
\]
Aquí \(m_e^{(0)}\) es la realización de la coordenada electrónica
basal, con su prefactor, su refinamiento y su base dimensional.
Si \(E_e^{(0)}=\varepsilon_e^{(0)}\mathcal U_E\), entonces
\[
 m_e^{(0)}=\varepsilon_e^{(0)}\mathcal U_Ec_{\mathrm{int}}^{-2}
 =b_e\,m_{\mathrm{clk}},\qquad
 b_e=\varepsilon_e^{(0)}\frac{\mathcal U_E}{E_{\mathrm{clk}}}.
\]
En la coordenatización \(\mathcal U_E=E_{\mathrm{clk}}\) se obtiene
\(b_e=\varepsilon_e^{(0)}\); en otra coordenatización se conserva el cociente
de bases indicado. Si las coordenadas de torre ya son
relativas al electrón, se omite una segunda sustracción de \(\eta_e\).
Esta igualdad enlaza el operador con la firma de cada ruta y conserva
explícitamente la normalización electrónica. Una elección distinta de
sección de referencia debe transportar también esa normalización.
La escala de la constante reducida de Planck pertenece a esta construcción
dimensional. El carácter y el factor de retorno actúan después como
escalares. El electrón conserva su estructura: no se identifica con el
cuanto cronológico \(m_{\mathrm{clk}}\).

Para dos rutas en la misma coordenatización absoluta,
\[
 \frac{m_Y^\beta}{m_X^\beta}
 =\chi_{\mathrm{ref}}(\sigma_Y-\sigma_X,\eta_Y-\eta_X)
 R_{\mathrm{act}}^{\kappa_Y-\kappa_X}.
\]
La sección dimensional común cancela. La diferencia de los exponentes de
retorno no cancela salvo que sea nula. Esta identidad conserva la
distinción entre escala absoluta y estructura relativa de una colección.

\subsection{El sector de propagación nula}

La exponencial toma valores estrictamente positivos. La coordenatización nula se
construye antes de ese carácter, en la realización partida
\[
 \mathcal A_{\mathrm{sp}}=\mathbb R[R]/(R^2-1),\quad
 z=r_+P_++r_-P_-,\quad N(z)=r_+r_-.
\]
Las direcciones \(\mathbb RP_+\) y \(\mathbb RP_-\) tienen norma
nula. Una propagación transportada en una de ellas pertenece a ese sector;
la lectura material bicapa emplea el acoplamiento de ambas. De esta manera,
el fotón y el sector gluónico conservan su residencia en la teoría aunque
no constituyan una fila del carácter positivo. La realización de sus
representaciones se expone junto a los sectores correspondientes.

\section{Acción, gravedad e información térmica}
\label{sec:gravity}
\subsection{Período dodecafásico y realización circular de la acción}
El calendario ya construido determina $T_{108}=108t_0$.
En su realización circular,
\[
\omega_{108}=\frac{2\pi}{108t_0},\qquad
R_{108}=\frac{c_{\rm int}}{\omega_{108}}
=\frac{54}{\pi}\ell_0,\qquad \ell_0=c_{\rm int}t_0.
\]
La unidad $\ell_0$ y las coordenatizaciones de tiempo y velocidad se mantienen
explícitas; la coordenatización SI no selecciona el ciclo.
La composición longitudinal de la sección~\ref{md:rigidez:longitud}
fija en esta edición $R_{108}=L_\alpha$ y
$t_0=\pi_{\rm HMT}L_\alpha/(54c_{\rm int})$.
El argumento circular siguiente se conserva como especialización
del carácter unitario del teorema~\ref{md:rigidez:teorema}.
Para cada sección orientada de acción $a\in\{\mathrm{pre},\mathrm{ret}\}$,
\[
E_{108,a}=\hbar_a\omega_{108},\qquad
m_{108,a}=\frac{E_{108,a}}{c_{\rm int}^2}.
\]
La longitud reducida y la completa son respectivamente
\[
\frac{\hbar_a}{m_{108,a}c_{\rm int}}=R_{108},
\qquad
\frac{h_a}{m_{108,a}c_{\rm int}}=2\pi R_{108}.
\]
La misma acción y la misma duración intervienen en ambas lecturas.

\begin{proposition}[Selector de coincidencia de radios]
En la colección dimensional
\[
G_a(\vartheta)=\vartheta\frac{c_{\rm int}^5t_0^2}{\hbar_a},
\qquad \vartheta>0,
\]
la condición de la realización circular
$G_a(\vartheta)m_{108,a}/c_{\rm int}^2=R_{108}$
tiene la única solución
\[
\vartheta^\circ_{108}=(54/\pi)^2.
\]
\end{proposition}
\begin{proof}
El primer radio es
$r_g(\vartheta)=\vartheta(\pi/54)\ell_0$.
Compararlo con $(54/\pi)\ell_0$, con $\ell_0>0$, reduce la
igualdad a una ecuación lineal estrictamente creciente.
\end{proof}
La condición de coincidencia es parte explícita de esta realización;
el teorema resuelve su coeficiente sin usar un valor metrológico de $G$.
El radio aquí empleado es $Gm/c^2$, por lo que no se traslada
inadvertidamente el factor dos de la convención $2Gm/c^2$.

\begin{proposition}[Acción y gravedad recíprocas, área conservada]
En las dos coordenatizaciones anteriores,
\[
\frac{G_{\rm pre}}{G_{\rm ret}}=R_{\rm act}^{-1},\qquad
\hbar_aG_a=\vartheta^\circ_{108}c_{\rm int}^5t_0^2,
\qquad
\ell_P^2=\frac{\hbar_aG_a}{c_{\rm int}^3}
=\vartheta^\circ_{108}\ell_0^2.
\]
\end{proposition}
\begin{proof}
La definición de $G_a$ contiene $\hbar_a^{-1}$ y el resto de factores
es común a las dos coordenatizaciones. Tomar el cociente y el producto da
las tres igualdades.
\end{proof}

\subsection{Definición estructural y función de la constante gravitatoria}
\label{sec:ii-definicion-gravedad}

\begin{definition}[Sección gravitatoria contragrediente de la acción]
Para una sección positiva \(\hbar_a\), una velocidad
\(c_{\rm int}>0\) y el radio cronológico \(R_{108}\) de la
sección~\ref{sec:gravity}, la sección gravitatoria circular es
\begin{equation}
 G_a=\frac{c_{\rm int}^{3}R_{108}^{2}}{\hbar_a},
 \qquad R_{108}=\frac{54}{\pi}\ell_0,
 \qquad \ell_0=c_{\rm int}t_0.
 \label{eq:ii-definicion-gravedad-seccion}
\end{equation}
Es la sección recíproca de la acción que realiza la coincidencia
entre el radio gravitatorio reducido del ciclo y su radio
cronológico, en la convención \(r_g=Gm/c_{\rm int}^2\).
\end{definition}

La proposición de coincidencia de radios de la
sección~\ref{sec:gravity} demuestra su unicidad en la colección
\(G_a(\vartheta)=\vartheta c_{\rm int}^{5}t_0^2/\hbar_a\).
La condición circular fija \(\vartheta=(54/\pi)^2\) y conduce
a \eqref{eq:ii-definicion-gravedad-seccion}. Esa condición y las
secciones dimensionales positivas forman parte de la realización.
La sección~\ref{md:rigidez} construye su longitud y reúne la
prueba de unicidad para el carácter positivo general, con
$R_{108}=L_\alpha$. Por ello el selector circular y el lector
longitudinal conservan una misma normalización.
Manteniéndolas comunes entre las dos orientaciones, se obtiene
\begin{equation}
 \frac{G_{\mathrm{pre}}}{G_{\mathrm{ret}}}
 =R_{\mathrm{act}}^{-1},
 \qquad
 \frac{G_a\hbar_a}{c_{\rm int}^{3}}
 =R_{108}^{2}=\ell_P^2.
 \label{eq:ii-definicion-gravedad-area}
\end{equation}
La prueba es la sustitución directa de
\eqref{eq:ii-definicion-gravedad-seccion}; el cambio de acción
se compensa por el cambio recíproco de \(G_a\).

\section{Información reducida, energía y coordenadas térmicas}

\subsection{Lector, memoria y recuperación del estado}

Fijemos un nivel finito del desarrollo y denotemos por \(\Omega\) el
conjunto de estados enriquecidos admisibles. El registro conserva hoja,
orientación, residuo, cociente, acarreo, ruta, fase y memoria en los
dominios donde interviene cada coordenada. Sean
\[
 q:\Omega\longrightarrow Y,\qquad
 M:\Omega\longrightarrow\mathcal M
\]
la representación visible y la memoria retenida, respectivamente. Para una
distribución \(p\) sobre \(\Omega\), escribimos
\(\mathsf H(p)=-\sum_xp_x\log p_x\), con \(0\log0=0\).
El logaritmo se aplica a probabilidades adimensionales.

\begin{definition}[Recuperación genealógica de una representación]
El par \(\widehat q=(q,M)\) es un lector recuperable sobre el soporte
\(\Omega_p\) cuando existe
\[
 d:\widehat q(\Omega_p)\longrightarrow\Omega_p,
 \qquad d\circ\widehat q=\operatorname{id}_{\Omega_p}.
\]
La información de fibra de \(q\) es
\(\mathsf H_{\mathrm{fib}}(q;p)=\mathsf H(X\mid q(X))\), donde
\(X\) tiene ley \(p\).
\end{definition}

\begin{proposition}[Identidad de recuperación y coste de la reducción]
Para un lector recuperable,
\begin{align}
 \mathsf H(X\mid q(X),M(X))&=0,\\
 \mathsf H(X)&=\mathsf H(q(X))+\mathsf H(X\mid q(X)),\\
 \mathsf H(X\mid q(X))&=\mathsf H(M(X)\mid q(X)).
 \label{eq:ii-kb-aut-fibra}
\end{align}
\end{proposition}
\begin{proof}
La aplicación \(d\) recupera \(X\) del par \((q(X),M(X))\),
por lo que la primera entropía condicional es nula. Como \(q\) es
función de \(X\), la regla de la cadena da la segunda igualdad.
Para la tercera, se desarrolla de dos maneras
\(\mathsf H(X,M(X)\mid q(X))\). Una proporciona
\(\mathsf H(X\mid q(X))\), pues \(M\) es función de \(X\);
la otra proporciona
\(\mathsf H(M(X)\mid q(X))+\mathsf H(X\mid q(X),M(X))\).
\end{proof}

La división euclídea de las hojas de APP muestra el mecanismo
aritmético elemental. Con \(i,j\in\{1,\ldots,9\}\), se conserva
\[
 i+j=9q^+_{ij}+R^+_{ij},\qquad
 ij=9q^\times_{ij}+R^\times_{ij},\qquad
 1\leq R^\sigma_{ij}\leq9\quad(\sigma\in\{+,\times\}).
\]
El par residuo--cociente recupera el valor entero de la operación.
Al sumar las diferencias sobre las ochenta y una posiciones se obtiene
\[
 2025-810=54+9\cdot129.
\]
La memoria del cociente restaura la contribución que no expresa la
reducción residual. La recuperabilidad del estado completo requiere,
además, conservar las coordenadas de posición y de ruta pertinentes;
el par residuo--cociente identifica aquí el resultado de la operación.

\subsection{Actualización reversible, decisión trítica y borrado}

Sea \(U:\Omega\to\Omega\) la actualización TPK en un dominio donde
el transporte enriquecido posee inversa. La ley transportada es
\(p'_y=p_{U^{-1}y}\). La biyección conserva el multiconjunto de
probabilidades y, por tanto,
\begin{equation}
 \mathsf H(U(X))=\mathsf H(X),\qquad
 \mathsf H(X\mid U(X))=0.
 \label{eq:ii-kb-aut-reversible}
\end{equation}
Si la salida observada es \(q(U(X))\), la información no expresada
se cuantifica por \(\mathsf H(X\mid q(U(X)))\). El estado completo,
su representación y el reinicio de un registro son así tres mapas distintos.

Para los tres estados orientados del TRIT, sea
\(p=(p_{-1},p_0,p_{+1})\) una distribución normalizada. Su información es
\begin{equation}
 I_{\mathrm{TRIT}}(p)=-\sum_{a=-1}^{1}p_a\log p_a,
 \qquad 0\leq I_{\mathrm{TRIT}}(p)\leq\log3.
 \label{eq:ii-kb-aut-trit}
\end{equation}
El extremo superior corresponde a \(p_a=1/3\). En efecto,
\[
 D(p\Vert(1/3,1/3,1/3))
 =\sum_ap_a\log(3p_a)=\log3-I_{\mathrm{TRIT}}(p)\geq0;
\]
la desigualdad se obtiene de la convexidad de \(t\log t\), y la
igualdad caracteriza la distribución uniforme. El extremo inferior
corresponde a una distribución concentrada en un estado.

La realización ideal de Landauer utiliza una memoria con estados lógicos
energéticamente degenerados, un entorno isotermo de temperatura \(\Theta\)
y un balance que conserva la salida del mapa lógico y contabiliza los
registros auxiliares. En ese modelo, la disminución de entropía lógica es
\(k_B[\mathsf H(X)-\mathsf H(F(X))]
=k_B\mathsf H(X\mid F(X))\) para un mapa determinista \(F\).
El balance entrópico exige transferir al entorno al menos esa cantidad;
su expresión energética es \(Q\geq k_B\Theta\mathsf H(X\mid F(X))\).
La distinción entre evolución lógica reversible y borrado físico
corresponde a la realización de Landauer--Bennett \cite{bennett2003}.
La actualización invertible
\eqref{eq:ii-kb-aut-reversible} tiene contribución lógica nula; el
reinicio de una terna uniforme requiere transferir \(\log3\) nats
al entorno y satisface \(Q\geq k_B\Theta\log3\).
La energía de control y la disipación de una realización física
concreta pertenecen a su balance termodinámico adicional.

\subsection{Acción, frecuencia y energía de decisión}

Sean \(\mathcal L_{\rm act}\) y \(\mathcal L_T\) las rectas
dimensionales de acción y duración; la recta energética es
\(\mathcal L_E=\mathcal L_{\rm act}\otimes\mathcal L_T^*\).
El desarrollo anterior aporta las secciones positivas
\(\hbar_a\in\mathcal L_{\rm act}\), con
\(a\in\{\mathrm{pre},\mathrm{ret}\}\), y la duración elemental
\(t_0\in\mathcal L_T\). El calendario completo proporciona
\[
 T_{108}=108t_0,\qquad
 \omega_{108}=\frac{2\pi}{T_{108}},\qquad
 h_a=2\pi\hbar_a.
\]
La contracción dimensional entre acción y frecuencia define
\begin{equation}
 \epsilon_a:=\hbar_a\omega_{108}
 =\frac{h_a}{108t_0}
 =\frac{\hbar_a}{t_P},\qquad
 t_P=\frac{54}{\pi}t_0.
 \label{eq:ii-kb-aut-energia}
\end{equation}
El último miembro emplea el mismo período circular reducido que la
sección gravitatoria. Las tres expresiones corresponden a una única
sección de energía, con una orientación \(a\) mantenida en todos
los factores. La información normalizada de una decisión uniforme
produce la sección de energía por nat
\begin{equation}
 \varepsilon_a:=\frac{\epsilon_a}{\log3}\in\mathcal L_E^+.
 \label{eq:ii-kb-aut-energia-nat}
\end{equation}
La acción previamente construida determina la escala de energía al
fijar la duración; el cardinal de alternativas del TRIT determina su
normalización informacional. El valor de \(\log3\) entra por el
conteo y la medida declarados en \eqref{eq:ii-kb-aut-trit}.

\subsection{Definición dimensional de Boltzmann y selección de coordenadas}

Sea \(\mathcal L_\Theta\) la recta orientada de temperatura.
El espacio dimensional de entropía es
\(\mathcal L_{\rm ent}=\mathcal L_E\otimes\mathcal L_\Theta^*\).
Una aplicación lineal positiva
\(k_B:\mathcal L_\Theta\to\mathcal L_E\) es, equivalentemente,
un elemento positivo de \(\mathcal L_{\rm ent}\). Por eso el mismo
objeto puede actuar sobre una temperatura o multiplicar una información
adimensional.

\begin{definition}[Normalización térmica de la decisión trítica]
Un par térmico compatible con la sección orientada \(a\) consta de
un transductor positivo \(k_B\) y una sección positiva
\(\Theta_{{\rm clk},a}\in\mathcal L_\Theta\) tales que
\begin{equation}
 k_B(\Theta_{{\rm clk},a})=\varepsilon_a,
 \qquad
 k_B\Theta_{{\rm clk},a}\log3=\epsilon_a.
 \label{eq:ii-kb-aut-par-termico}
\end{equation}
La constante de Boltzmann desempeña la función de transductor
energía--temperatura de la información trítica en esta realización.
\end{definition}

\begin{theorem}[Unicidad relativa a la sección térmica y covariancia dimensional]
Fijada una sección positiva \(\Theta_{{\rm clk},a}\), existe una
única aplicación lineal positiva que satisface
\eqref{eq:ii-kb-aut-par-termico}. Para bases positivas \(u_E,u_\Theta\),
si
\[
 \epsilon_a=\epsilon_a^{\rm num}u_E,\quad
 \Theta_{{\rm clk},a}=\theta_a u_\Theta,\quad
 k_B(u_\Theta)=b\,u_E,
\]
su coordenada es
\begin{equation}
 b=\frac{\epsilon_a^{\rm num}}{\theta_a\log3}.
 \label{eq:ii-kb-aut-coordenada}
\end{equation}
Bajo \(u'_E=s_Eu_E\), \(u'_\Theta=s_\Theta u_\Theta\), con
\(s_E,s_\Theta>0\), se cumple
\[
 b'=\frac{s_\Theta}{s_E}b,\qquad
 \theta'_a=\frac{\theta_a}{s_\Theta},\qquad
 (\epsilon_a^{\rm num})'=\frac{\epsilon_a^{\rm num}}{s_E}.
\]
\end{theorem}
\begin{proof}
Una sección no nula genera una recta. Prescribir su imagen determina
la aplicación lineal única; la positividad de ambas secciones prueba
su positividad. Sustituir las coordenadas en
\eqref{eq:ii-kb-aut-par-termico} da
\(b\theta_a\log3=\epsilon_a^{\rm num}\). Las tres leyes de
transformación se obtienen expresando las mismas secciones y la misma
aplicación en las bases nuevas.
\end{proof}

Una misma constante de Boltzmann debe actuar en las dos orientaciones.
Puesto que \(t_P\) es común, la compatibilidad simultánea equivale a
\begin{equation}
 \frac{\Theta_{{\rm clk},{\rm pre}}}{\Theta_{{\rm clk},{\rm ret}}}
 =\frac{\epsilon_{\rm pre}}{\epsilon_{\rm ret}}
 =\frac{\hbar_{\rm pre}}{\hbar_{\rm ret}}.
 \label{eq:ii-kb-aut-orientaciones}
\end{equation}
En efecto, un isomorfismo de rectas conserva cocientes de secciones.
Recíprocamente, si la igualdad anterior se cumple, el isomorfismo
determinado en una orientación satisface también la normalización de
la otra. Así, el retorno modifica la sección térmica del ciclo en la
misma proporción que la energía, mientras el transductor permanece
común a ambas lecturas.

La ecuación anterior explicita la elección requerida para expresar un
valor individual. Si la sección térmica todavía no ha sido seleccionada
por un lector independiente, para cada \(\lambda>0\) el par
\[
 (k_B,\Theta_{{\rm clk},a})
 \longmapsto(\lambda k_B,\lambda^{-1}\Theta_{{\rm clk},a})
\]
conserva \(\epsilon_a\). Por consiguiente, la presentación de una
coordenada individual internamente seleccionada exige dar el lector de
\(\Theta_{{\rm clk},a}\), las bases y su transporte dimensional.
Las ecuaciones de acción e información permanecen determinadas antes
de esa especificación. La identificación posterior de las bases con
joule y kelvin pertenece a la representación metrológica del resultado.

\subsection{Entropía dimensional e información del estado reducido}

Para una distribución de rutas admitidas \(p\), o para un estado
reducido de matriz de densidad \(\rho\), sean
\[
 s(p)=-\sum_xp_x\log p_x,\qquad
 s(\rho)=-\operatorname{Tr}(\rho\log\rho).
\]
Sus representaciones dimensionales son
\begin{equation}
 S(p)=k_Bs(p),\qquad S(\rho)=k_Bs(\rho)
 \quad\text{en }\mathcal L_{\rm ent}.
 \label{eq:ii-kb-aut-entropia}
\end{equation}
La evaluación en \(\Theta\) produce una energía
\(\Theta S=k_B\Theta\,s\). Para la terna uniforme a la temperatura
del ciclo, \eqref{eq:ii-kb-aut-par-termico} da exactamente
\[
 S_{\rm TRIT}=k_B\log3,\qquad
 \Theta_{{\rm clk},a}S_{\rm TRIT}=\epsilon_a.
\]

% Recuperación documental para el artículo V; RESULTADO_RECUPERADO.
% Propietario: /Users/ruben/Documents/New project/output/REVISION_EDITORIAL_HMT_MD_20260905/01_FUENTES/INTEGRAL/tree/New project/output/ENSAMBLADO_CAUSAL_RESTAURADO_HMT_MD_20260905/fuente/manuscrito/sections/md/delta_127/04b_completaciones_cuanticas.tex
% Líneas de origen: 1--334.
% REV02: construcción algebraica; la aplicación estadística se reúne en 40.
% Correspondencia editorial: gestion/CONSOLIDACION_PAULI_FOCK_REV02.json.
\section{Construcción algebraica de las completaciones cuánticas sobre APP}
\label{v:rec:chap:completaciones-cuanticas}
\label{v:sec:completaciones-algebraicas}

APP y TPK proporcionan estados orientados, transporte y paridad de hoja.
Su realización lineal permite construir las álgebras simétrica y exterior,
con sus operadores de creación, aniquilación y ocupación. Esta sección
establece esas construcciones algebraicas, el sector celular de campos y
el acoplamiento reversible de medición. La aplicación de los espectros de ocupación al equilibrio se desarrolla en la sección~\ref{v:rec:ocupacion:section:02}.
Cada tramo conserva las condiciones de la representación que utiliza.

\subsection{Espacio lineal de estados asociado a una realización de partícula}
\label{v:rec:completaciones:section:02}

Sea \(S\) un conjunto finito de clases de rutas TPK a la profundidad elegida
y sea
\[
 \mathcal H_1=\ell^2(S)
\]
el espacio de Hilbert con base ortonormal \((e_s)_{s\in S}\). Una elección de
pesos positivos \(E_s\) define el operador diagonal
\[
 He_s=E_se_s.
\]
La elección de \(S\), de su cociente por simetrías y de los pesos forma parte
del modelo. Una vez fijada, las completaciones multipartícula son canónicas.

\subsection{Completaciones simétrica y exterior}
\label{v:rec:completaciones:section:03}

Definimos
\[
 \mathcal F_+(\mathcal H_1)=
 \bigoplus_{n\geq0}\operatorname{Sym}^n\mathcal H_1,
 \qquad
 \mathcal F_-(\mathcal H_1)=
 \bigoplus_{n\geq0}\bigwedge^n\mathcal H_1.
\]
La primera es la completación bosónica; la segunda, la fermiónica. En ambas,
el sumando de grado cero es el espacio del vacío. Las potencias exteriores
se dotan del producto escalar para el cual las cuñas ordenadas de vectores
de una base ortonormal son ortonormales; las potencias simétricas se utilizan
con la normalización usual de los operadores de creación y aniquilación.

\subsection{Graduación de hoja y álgebra de intercambio}
\label{v:rec:completaciones:section:04}

Supongamos que la hoja orientada define un carácter aditivo
\[
 p:\operatorname{Mor}(\mathsf P_{\rm TPK})\longrightarrow\mathbb Z/2\mathbb Z,
 \qquad p(\gamma\eta)=p(\gamma)+p(\eta).
\]
Sobre el espacio lineal graduado correspondiente se introduce la simetría
\[
 \tau(v\otimes w)=(-1)^{p(v)p(w)}w\otimes v.
\]

\begin{theorem}[Completación estadística determinada por la paridad]
\label{v:rec:completaciones:statement:02}
El cociente del álgebra tensorial por las relaciones
\[
 v\otimes w-(-1)^{p(v)p(w)}w\otimes v
\]
es simétrico sobre el sector par y exterior sobre el sector impar. Dos estados
impares idénticos tienen producto nulo; los estados pares admiten ocupación
arbitraria.
\end{theorem}

\begin{proof}
Si \(p(v)=p(w)=0\), la relación es \(vw=wv\). Si
\(p(v)=p(w)=1\), es \(vw=-wv\); tomando \(v=w\) y trabajando en
característica cero se obtiene \(v^2=0\). Los sectores mixtos conmutan sin
signo adicional. Esta es la propiedad universal del álgebra simétrica
graduada.
\end{proof}

\subsection{Operadores canónicos e idempotencia de la ocupación}
\label{v:pauli:car-construccion}

El carácter de paridad determina la completación; el producto de esa
completación determina a continuación el álgebra de operadores. En el
sector exterior, la creación \(a_s^\dagger\) es el producto exterior por
\(e_s\), y \(a_s\) es su adjunto.

\begin{theorem}[Principio de exclusión de Pauli en la completación exterior]
\label{v:rec:completaciones:statement:01}
\label{v:rec:thm:principio-exclusion-pauli}
Para todo modo \(s\in S\), los operadores de creación y aniquilación de la
completación exterior satisfacen
\begin{equation}
 \{a_s,a_t^\dagger\}=\delta_{st}I,
 \qquad
 \{a_s,a_t\}=\{a_s^\dagger,a_t^\dagger\}=0.
 \label{v:pauli:eq:car}
\end{equation}
En particular, el operador de ocupación
\(N_s=a_s^\dagger a_s\) es un proyector:
\begin{equation}
 N_s^2=N_s,
 \qquad \Spec(N_s)\subseteq\{0,1\}.
 \label{v:pauli:eq:proyeccion-ocupacion}
\end{equation}
Los autovalores \(0\) y \(1\) corresponden, respectivamente, a la ausencia
y a la presencia del modo \(s\).
\end{theorem}

\begin{proof}
Sobre una cuña ordenada, \(a_s\) suprime el factor \(e_s\), cuando está
presente, con el signo de su posición; si está ausente, da cero. Componer
esta operación con el producto exterior demuestra las relaciones
\eqref{v:pauli:eq:car}: para \(s=t\), las dos composiciones cubren los casos
complementarios de ausencia y presencia; para \(s\ne t\), los signos de
intercambio se cancelan. La antisimetría implica
\(e_s\wedge e_s=0\) y
\(a_s^2=(a_s^\dagger)^2=0\). Entonces
\begin{equation}
 N_s^2=a_s^\dagger a_s a_s^\dagger a_s
 =a_s^\dagger(I-a_s^\dagger a_s)a_s
 =a_s^\dagger a_s=N_s.
 \label{v:pauli:eq:prueba-idempotencia}
\end{equation}
El operador es autoadjunto, y todo autovalor de un idempotente satisface
\(\lambda^2=\lambda\). Por tanto, su espectro está contenido en
\(\{0,1\}\). El vacío y el estado \(e_s\) realizan ambos valores
\cite{Pauli1925,Dirac1926}.
\end{proof}

En el sector simétrico, los operadores bosónicos \(b_s^\dagger,b_s\)
actúan sobre la base normalizada de ocupaciones por los factores
\(\sqrt{m_s+1}\) y \(\sqrt{m_s}\), respectivamente. De sus composiciones
se obtienen
\begin{equation}
 [b_s,b_t^\dagger]=\delta_{st}I,
 \qquad [b_s,b_t]=[b_s^\dagger,b_t^\dagger]=0.
 \label{v:pauli:eq:ccr}
\end{equation}
Estas identidades se entienden sobre el dominio algebraico de partículas
finitas, estable bajo los operadores considerados.

El teorema~\ref{v:rec:completaciones:statement:02} convierte una paridad composicional TPK en una regla exacta de
ocupación. Para obtener el teorema físico de espín--estadística hay que
construir además una teoría local lorentziana positiva y demostrar que la
paridad de hoja coincide con el espín bajo rotaciones. La igualdad no se
deduce únicamente de que ambos objetos tomen valores módulo dos.

\section{Distribuciones de Bose--Einstein y Fermi--Dirac}
\label{v:rec:ocupacion:section:02}

La formulación de equilibrio parte del Hamiltoniano y del número total de
ocupaciones,
\begin{equation}
 H=\sum_kE_k^{\rm occ} n_k,
 \qquad
 N=\sum_k n_k,
 \qquad
 K_\mu:=H-\mu N.
 \label{v:rec:eq:hamiltoniano-gran-canonico-rev3}
\end{equation}
Para \(\beta>0\), el estado gran canónico es
\begin{equation}
 \rho_{\beta,\mu}
 =\frac{e^{-\beta K_\mu}}{\mathcal Z_{\beta,\mu}},
 \qquad
 \mathcal Z_{\beta,\mu}
 =\operatorname{Tr}(e^{-\beta K_\mu}).
 \label{v:rec:eq:estado-gran-canonico-rev3}
\end{equation}

\begin{theorem}[Factorización gran canónica y leyes de ocupación]
\label{v:rec:thm:factorizacion-gran-canonica-rev3}
Para niveles de energía \(E_k^{\rm occ}\) con degeneraciones \(g_k\), la
traza fermiónica factoriza como
\begin{equation}
 \mathcal Z_{\rm FD}
 =\prod_k\bigl(1+e^{-\beta(E_k^{\rm occ}-\mu)}\bigr)^{g_k},
 \qquad
 \langle n_k\rangle_{\rm FD}
 =\frac{g_k}{e^{\beta(E_k^{\rm occ}-\mu)}+1}.
 \label{v:rec:eq:factorizacion-fd-rev3}
\end{equation}
En el sector bosónico, bajo la condición
\(\mu<\inf_kE_k^{\rm occ}\),
\begin{equation}
 \mathcal Z_{\rm BE}
 =\prod_k\bigl(1-e^{-\beta(E_k^{\rm occ}-\mu)}\bigr)^{-g_k},
 \qquad
 \langle n_k\rangle_{\rm BE}
 =\frac{g_k}{e^{\beta(E_k^{\rm occ}-\mu)}-1}.
 \label{v:rec:eq:factorizacion-be-rev3}
\end{equation}
\end{theorem}

\begin{proof}
El operador \(K_\mu\) es una suma de operadores de ocupación conmutantes, de
modo que su traza factoriza por modos. En cada modo fermiónico se suman las
ocupaciones \(0\) y \(1\), y el factor correspondiente es
\(1+u_k\), donde
\(u_k=e^{-\beta(E_k^{\rm occ}-\mu)}\). En cada modo bosónico se suman todas
las ocupaciones no negativas y se obtiene la serie geométrica
\((1-u_k)^{-1}\); la condición sobre \(\mu\) garantiza \(u_k<1\). Las
potencias \(g_k\) incorporan la degeneración. Finalmente,
\[
 \langle n_k\rangle
 =-\frac1\beta\frac{\partial}{\partial E_k^{\rm occ}}
   \log\mathcal Z
\]
produce las dos leyes de ocupación
\cite{Bose1924,Einstein1924,Fermi1926,Dirac1926,Pathria2011}.
\end{proof}

% RESULTADO_RECUPERADO: funtor dimensional integral, eq:cZ-internos,
% eq:longitud-ruta y definición de ell_0; composición con artículo III.
% La prueba reunida y su control algebraico son específicos de esta contribución.
\subsection{Velocidad constitutiva y realización cinemática del reloj}
\label{v:sec:enlace-velocidad-constitutiva}

La propagación radiativa utiliza la sección de velocidad producida por el
transporte angular del mismo estado HMT. La identificación se establece
antes de elegir unidades de representación. Sean $x=A_{\mathrm{rad}}$ e
$y=C^*_{\mathrm{rad}}$ las coordenadas angulares generadas en la
sección~\ref{sec:ii-angulos}, con los canales
\eqref{eq:ii-canales}, y sea
$\mathsf R$ la involución autoadjunta de las dos hojas, con proyectores
de rango uno $P_+$ y $P_-$. En el dominio $x>|y|$, el transporte y sus
canales son
\begin{equation}
 T=\exp(-xI-y\mathsf R)=q_+P_++q_-P_-,
 \qquad q_+=e^{-x-y},\quad q_-=e^{-x+y},\quad 0<q_\pm<1.
 \label{v:eq:velocidad-transporte}
\end{equation}
La respuesta de los bloques temporales de longitudes $90$ y $120$ es
\begin{equation}
 \mathsf V=(I-T^{90})(I-T^{120})^{-1}
       =r_+P_++r_-P_-,\qquad
 r_\pm=\frac{1-q_\pm^{90}}{1-q_\pm^{120}}.
 \label{v:eq:velocidad-respuesta}
\end{equation}
La construcción parte de los canales heredados; los símbolos $x,y$ y
$q_\pm$ designan sus coordenadas y no parámetros ajustados a una velocidad
objetivo.

\begin{proposition}[Identidad de las realizaciones constitutiva y cinemática]
\label{v:prop:identidad-velocidad}
En una misma coordenatización dimensional, la sección $c_{\mathrm{int}}$ del reloj
y la sección $c_{\mathrm{vac}}$ de la respuesta constitutiva coinciden:
\begin{equation}
 c_{\mathrm{int}}=c_{\mathrm{vac}}
   =\widehat c\,u_C,
 \qquad
 \widehat c=(r_+r_-)^{-1}
            =\exp(-\mathcal E),
 \label{v:eq:identidad-velocidad}
\end{equation}
donde
\begin{equation}
 \mathcal E=
 \log\!\bigl[(1-q_+^{90})(1-q_-^{90})\bigr]
 -\log\!\bigl[(1-q_+^{120})(1-q_-^{120})\bigr].
 \label{v:eq:velocidad-lector-simetrico}
\end{equation}
La igualdad conserva el intercambio de hojas.
\end{proposition}
\begin{proof}
Los factores que aparecen en los logaritmos son positivos. Por tanto,
\[
 \mathcal E=\log(r_+r_-)
            =\log\det\mathsf V,
 \qquad \exp(-\mathcal E)=(\det\mathsf V)^{-1}.
\]
El determinante se calcula aquí en la realización de dos hojas de rango
uno. La definición cinemática $c_{\mathrm{int}}=\exp(-\mathcal E)u_C$
y la definición constitutiva $c_{\mathrm{vac}}=(r_+r_-)^{-1}u_C$
son, en consecuencia, la misma sección. El intercambio $P_+\leftrightarrow
P_-$ permuta $r_+$ y $r_-$ y conserva su producto. Finalmente,
$\det T=e^{-2x}$ corresponde al transporte elemental, mientras que
$\det\mathsf V=r_+r_-$ corresponde a su respuesta temporal; ambos
determinantes tienen funciones distintas en esta composición.
\end{proof}

\subsubsection{Longitud de trayectoria y duración elemental}

El marco dimensional positivo $(u_S,u_C,u_T,u_Q)$ induce la base de
longitud $u_L=u_Cu_T$. Para una trayectoria enriquecida $\gamma$ de
$n=|\gamma|$ transiciones y canal constante, los lectores aditivos son
\begin{equation}
 T_{\mathrm{int}}(\gamma)=n u_T,
 \qquad L_{\mathrm{int}}(\gamma)=n\widehat c\,u_L.
 \label{v:eq:velocidad-longitud-reloj}
\end{equation}
Su aditividad se refiere al conteo y a la realización de cada transición:
las rutas enriquecidas conservan adicionalmente orientación y memoria.

\begin{corollary}[Normalización de la longitud elemental]
\label{v:cor:longitud-elemental-constitutiva}
Para $n>0$, $L_{\mathrm{int}}(\gamma)/T_{\mathrm{int}}(\gamma)
=c_{\mathrm{vac}}$. En la normalización $t_0=u_T$,
\begin{equation}
 \ell_0=c_{\mathrm{int}}t_0=\widehat c\,u_L.
 \label{v:eq:longitud-elemental-constitutiva}
\end{equation}
En una coordenatización temporal general $t_0=\tau_0u_T$, con $\tau_0>0$,
la expresión correspondiente es $\ell_0=\tau_0\widehat c\,u_L$.
\end{corollary}
\begin{proof}
La división de \eqref{v:eq:velocidad-longitud-reloj} cancela el entero
$n$ y utiliza $u_L/u_T=u_C$. La fórmula elemental se obtiene
multiplicando la sección de velocidad por $t_0$.
\end{proof}

La generalización a un canal dependiente de la arista conserva la fórmula
aditiva $L_{\mathrm{int}}(\gamma)=\sum_{e\in\gamma}\widehat c(e)u_L$.
Su cociente por $nu_T$ es la media de las velocidades de esas aristas.
La realización radiativa homogénea emplea el canal constante declarado
en \eqref{v:eq:velocidad-longitud-reloj}.

\subsubsection{Cambio de coordenatización y normalización adaptada}

\begin{proposition}[Covariancia de la identificación]
\label{v:prop:covariancia-identidad-velocidad}
Sean $u_C'=d u_C$ y $u_T'=\eta u_T$, con $d,\eta>0$.
La misma velocidad y la misma duración elemental tienen coordenadas
\begin{equation}
 \widehat c'=d^{-1}\widehat c,
 \qquad \tau_0'=\eta^{-1}\tau_0,
 \qquad u_L'=d\eta u_L.
 \label{v:eq:velocidad-cambio-carta}
\end{equation}
En particular, si $c_V=\widehat c_Vu_C^V$ y
$u_C^V=d u_C^{III}$, la igualdad de secciones equivale exactamente a
$d\widehat c_V=\widehat c_{III}$.
\end{proposition}
\begin{proof}
Se expresa cada sección en las dos bases. Se obtiene
$\widehat c'u_C'=\widehat c u_C$ y
$\tau_0'u_T'=\tau_0u_T$. Su producto satisface
\[
 \tau_0'\widehat c' u_L'
 =\frac{\tau_0}{\eta}\frac{\widehat c}{d}
       d\eta u_L
 =\tau_0\widehat c u_L=\ell_0.
\]
La última equivalencia resulta de sustituir la relación entre las bases.
\end{proof}

La elección adaptada $d=\widehat c$ da $u_C'=c_{\mathrm{vac}}$ y
$\widehat c'=1$. Esta elección expresa la sección ya construida en una
base adaptada a ella. El coeficiente espectral en la coordenatización interna sigue
siendo $(r_+r_-)^{-1}$; la inversión espectral utiliza esa coordenatización interna.
Por ejemplo, manteniendo las bases de acción y carga, las coordenadas
constitutivas se transforman como
\[
 \widehat\mu'=d\widehat\mu,
 \qquad \widehat\varepsilon'=d\widehat\varepsilon,
 \qquad
 \widehat\mu'\widehat\varepsilon'=(\widehat c')^{-2}.
\]
En la coordenatización adaptada, $\widehat\mu'=r_-/r_+$ y
$\widehat\varepsilon'=r_+/r_-$, pues las coordenadas internas son
$\widehat\mu=r_-^2$ y $\widehat\varepsilon=r_+^2$.

\section{Restitución radial orientada y transporte simpléctico}
\subsection{Recuperación de Planck con orientación conservada}

Con los semiejes etiquetados de la elipse, definimos
\(R_\star=\sqrt{b_S/a_S}\). Entonces
\(q=R_\star^4=e^{-2\eta_{\rm el}}=(1-\xi)/(1+\xi)\),
donde \(\xi=C^*/A\) y \(\eta_{\rm el}=\operatorname{artanh}\xi\).
Esta razón radial conserva las marcas de ambos ejes.

\begin{theorem}[Recuperador radial de la sección de retorno]
\label{rec:thm:accion}
En la coordenatización positiva, si \(q=R_\star^4\), entonces
\begin{align}
 \eta_{\rm ret}&=\alpha\frac{499-501q}{1+q},\label{rec:eq:eta-radio}\\
 \hbar_{\rm ret}&=\alpha\frac{499-501q}{1+q}\,10^{-34}\mathcal U_S,
 \qquad\hbar_{\rm pre}=H_5\,10^{-34}\mathcal U_S,\label{rec:eq:planck-radio}\\
 R_{\rm act}&=\frac{\hbar_{\rm pre}}{\hbar_{\rm ret}}
 =\frac{H_5(1+q)}{\alpha(499-501q)}.\label{rec:eq:ratio-radio}
\end{align}
La rama \(\eta_{\rm ret}>0\) equivale a \(q<499/501\), dentro de \(q>0\).
\end{theorem}
\begin{proof}
Las identidades angulares dan \(\eta_{\rm ret}=\alpha(500\xi-1)\). Se sustituye \(\xi=(1-q)/(1+q)\); el numerador es \(499-501q\). El denominador y \(\alpha\) son positivos. La realización dimensional y el cociente son los de las secciones de acción previamente construidas. La década \(-34\) y la unidad \(\mathcal U_S\) conservan su derivación determinantal anterior.
\end{proof}

Esta fórmula explicita la sección de Planck dentro del contraángulo. En términos de las magnitudes de acción, la relación angular es
\[
 C^*=2\left(\frac{\hbar_{\rm ret}}{10^{-34}\mathcal U_S}+\alpha\right).
\]
La división por la base dimensional extrae el coeficiente adimensional; no suma una magnitud de acción a una constante sin unidades.

Para la orientación angular \(\varepsilon\in\{-1,+1\}\), definimos
\[
 C^*_{\varepsilon}=2\varepsilon(\eta_{\rm ret}+\alpha),\quad
 \xi_\varepsilon=C^*_{\varepsilon}/A,\quad
 q_\varepsilon=\frac{1-\xi_\varepsilon}{1+\xi_\varepsilon}>0.
\]
El recuperador sobre ambas coordenatizaciones es
\begin{equation}
 \eta_{\rm ret}=\alpha\left(500\varepsilon
        \frac{1-q_\varepsilon}{1+q_\varepsilon}-1\right).
 \label{rec:eq:eta-orientada}
\end{equation}

\begin{corollary}[Invariancia del recuperador conjugado]
La transformación \((\varepsilon,q)\mapsto(-\varepsilon,q^{-1})\) conserva \(\eta_{\rm ret}\), \(\hbar_{\rm ret}\) y \(R_{\rm act}\).
\end{corollary}
\begin{proof}
La razón \((1-q)/(1+q)\) cambia de signo al invertir \(q\). El producto por \(\varepsilon\) permanece invariante. Se aplica \eqref{rec:eq:eta-orientada} y después las representaciones de acción.
\end{proof}

Si se descarta la orientación y se mantiene indebidamente \(\varepsilon=+1\), la misma expresión produce \(\eta_{\rm ret}(q^{-1})=-\eta_{\rm ret}(q)-2\alpha\). El registro distingue, por tanto, dos operaciones algebraicamente distintas. La orientación angular \(\varepsilon\) tampoco se identifica sin más con el signo simpléctico \(\sigma\): el intercambio de ejes es antisimpléctico, mientras el giro que los intercambia preserva la forma de área.

\subsection{Entrelazamiento de fase y anisotropía}

Sean
\[
 S_\eta=\operatorname{diag}(e^{\eta/2},e^{-\eta/2}),\quad
 J_0=\begin{pmatrix}0&-1\\1&0\end{pmatrix},\quad
 J_\eta=S_\eta J_0S_\eta^{-1}
       =\begin{pmatrix}0&-e^\eta\\e^{-\eta}&0\end{pmatrix}.
\]
Se tienen \(\det S_\eta=1\), \(J_\eta^2=-I\) y
\(C_\eta(\theta)=e^{\theta J_\eta}=S_\eta e^{\theta J_0}S_\eta^{-1}\).
La ley \(S_{\eta+\delta}=S_\delta S_\eta\) prueba el entrelazamiento exacto
\begin{equation}
 C_{\eta+\delta}(\theta)S_\delta=S_\delta C_\eta(\theta).
 \label{rec:eq:entrelazamiento}
\end{equation}
La transformación preserva \(dQ\wedge dP\) y lleva la evolución de fase a la nueva elipse. Con la normalización de \eqref{eq:ii-elipse-area-fase}, la acción barrida es \(\hbar\theta\), o \(\sigma\hbar\theta\) en la hoja orientada. La fase \(\theta\), la anisotropía \(\eta_{\rm el}\) y el parámetro \(t\) del flujo de \(K\) conservan así funciones diferentes, relacionadas por operadores explícitos.

\section{Dominio positivo de la coordenatización elíptica}
\begin{proposition}[Pertenencia de las representaciones a la coordenatización elíptica]
\label{iv:ellipse:pertenencia}
La rama positiva de \eqref{eq:ii-par-angular}, con los
intervalos \eqref{md:accion:intervalos}, satisface
\(A>0\), \(0<C^*<A\) y \(\hbar_s>0\).
La reversión de la coordenada impar conserva \(|C^*/A|<1\).
\end{proposition}
\begin{proof}
La proposición \ref{md:accion:positividad} da
\(0<\eta_{\mathrm{ret}}<H_5\) y la positividad de ambas
secciones de acción. Para una cota superior de \(H_5\),
\eqref{md:accion:intervalos} implica
\(\alpha<1/125\) y \(\varphi>8/5\). Al conservar sólo
los términos positivos de \eqref{eq:el-h5}, se obtiene
\[
 H_5<
 \frac{\sqrt5}{2}+
 \frac{9}{16\cdot125^2}+
 \frac{7}{48\cdot125^4}.
\]
Como \(5<81/16\), se tiene \(\sqrt5/2<9/8\). Además,
\[
 \frac{9}{16\cdot125^2}+
 \frac{7}{48\cdot125^4}
 =\frac{421882}{11718750000}<\frac1{1000}.
\]
Por tanto, \(H_5<9/8+1/1000<2\). Usando ahora
\(\alpha>7/1000\),
\[
 0<C^*=2(\eta_{\mathrm{ret}}+\alpha)
 <2\left(2+\frac1{125}\right)=\frac{502}{125}
 <7<1000\alpha=A.
\]
El intercambio de canales de
\ref{prop:ii-inversion-canales-angulares} cambia \(C^*\) por \(-C^*\)
y conserva \(A\); por ello conserva la desigualdad absoluta.
\end{proof}

\section{Retornos primitivos y frecuencia levantada}
\subsection{Autómata local de extensión}
Sea \(\mathscr X_\Gamma\) el conjunto finito de estados locales compatibles con la ruta. Un estado conserva

\[
 x=(t,g,B,p,\varepsilon,c,w),
\]
donde \(t\) es el tiempo discreto, \(g\) la fase nonádica, \(B\) el bloque, \(p\) el prefijo, \(\varepsilon\) la hoja, \(c\) el acarreo y \(w\) el devanado. Existe una arista \(x\to y\) cuando \(y\) es una extensión admisible de \(x\), respeta el registro genealógico y supera la poda \(G_9\). La tabla de 36.864 composiciones de acarreo proporciona la ley local de suma; las tablas de retorno fijan la memoria.

Sea \(U_\Gamma\) el operador de transferencia normalizado de este grafo y \(R_\Gamma\) la involución de orientación. Se exige

\[
 U_{C_M\Gamma}=J U_\Gamma J^{-1},\qquad
 R_{C_M\Gamma}=-J R_\Gamma J^{-1}.
\]
\Needspace{7\baselineskip}
\subsection{Conteos primitivos orientados}
Para \(n\ge1\), la traza orientada

\[
 a_n(\Gamma)=\operatorname{Tr}(R_\Gamma U_\Gamma^n)
\]
cuenta retornos cerrados con signo. La inversión de Möbius separa órbitas primitivas:

\[
 p_n(\Gamma)
 =\frac1n\sum_{d\mid n}\mu(d)a_{n/d}(\Gamma).
\]
Se toma el electrón como origen afín de las correcciones relativas y se define

\[
 \boxed{
 \eta_h(\Gamma)=p_h(\Gamma)-p_h(\Gamma_e),
 \qquad h\in\mathfrak S.}
\]
Así se obtiene el mapa

\[
 \mathcal F_{\rm tower}:
 \Gamma^{\rm enr}\longmapsto
 \eta(\Gamma)=(\eta_h(\Gamma))_{h\in\mathfrak S}.
\]
No hay parámetros por especie. El mismo autómata, la misma involución y la misma inversión de Möbius actúan sobre las 324 rutas.

\Needspace{7\baselineskip}
\subsection{Convergencia del refinamiento primitivo}
Sea \(N_\Gamma=\dim\mathscr X_\Gamma\). Si \(U_\Gamma\) es una contracción,

\[
 |a_n(\Gamma)|\le N_\Gamma.
\]
El número de divisores de \(n\) crece sublinealmente, por lo que

\[
 |p_n(\Gamma)|
 \le\frac{N_\Gamma\tau(n)}n.
\]
Como \(0<q_+<q_-<1\),

\[
 |s_h|\le q_-^h+q_+^h\le2q_-^h.
\]
Por comparación,

\[
 \sum_{h\in\mathfrak S}|\eta_h(\Gamma)s_h|<\infty
\]
uniformemente en cualquier colección finita de fibras. El carácter refinado está, por tanto, bien definido.

\Needspace{7\baselineskip}
\subsection{Identidades de control}
La construcción debe satisfacer simultáneamente:

\begin{enumerate}
\item invariancia por renumeración de estados;
\item covarianza bajo conjugación de hoja;
\item compatibilidad con suma directa de componentes;
\item igualdad al recomputar desde registro genealógico y desde el grafo certificado;
\item conservación del origen \(\eta(\Gamma_e)=0\);
\item insensibilidad absoluta a una permutación de las columnas externas de masa;
\item determinación de las colas mediante una cota anterior al contraste.
\end{enumerate}
La representación exacta de un residual dado demuestra completitud del codominio, pero no sustituye estas siete pruebas.

\Needspace{7\baselineskip}
\section{Reloj de ruta y normalización absoluta}
\Needspace{7\baselineskip}
\subsection{Los retornos enteros no bastan para distinguir especies}
En el atlas de 324 rutas, el primer retorno de celda es 27, el retorno cinemático dentro de CT--108 es 54 y el retorno cinemático extendido es 216. Al ser comunes, estos enteros fijan la arquitectura temporal global, pero no pueden explicar por sí solos diferencias de masa entre especies.

\Needspace{7\baselineskip}
\subsection{Holonomía temporal}
La información específica reside en la holonomía del operador de extensión. Sea \(\lambda_j(\Gamma)\) un autovalor no trivial de \(U_\Gamma\) sobre el subespacio persistente. Escribimos

\[
 \lambda_j(\Gamma)=r_j(\Gamma)e^{i\theta_j(\Gamma)}.
\]
La frecuencia interna asociada es

\[
 \omega_j(\Gamma)
 =\frac{\theta_j(\Gamma)+2\pi w_j(\Gamma)}{108t_0},
\]
donde el entero \(w_j\) lo fija el devanado registrado, no una elección de rama posterior. Para un modo estable \(r_j=1\); para una resonancia \(r_j<1\) y su atenuación aporta una anchura.

El cociente

\[
 \rho_{\omega,j}(\Gamma)
 =\frac{\omega_j(\Gamma)}{\omega_0}
\]
es adimensional y puede modificar razones de masa sin alterar la década común \(-34\). La energía y la masa del modo son

\[
 E_j(\Gamma)=\hbar_{\rm ret}\omega_j(\Gamma)
 \chi_{\rm full}(\Gamma)R_{\rm act}^{\kappa_j(\Gamma)},
\]
\[
 m_j(\Gamma)=E_j(\Gamma)c_{\rm int}^{-2}.
\]
Esta ecuación reúne acción, reloj, carácter y lector sin convertir ninguno de ellos en otro. Su aplicación numérica exige determinar \(U_\Gamma\), fijar el devanado y declarar la coordenatización dimensional.


% Clausura focal de antecedentes: fuentes del artículo XI, edición 21-09-2026.
% Rangos, copias íntegras y huellas: gestion/mapa_restitucion_memoria.json.
\section{Antecedentes de restitución operatoria y memoria}
\label{md:antmem:capitulo}

El corte de estos antecedentes está después de la producción conjunta
APP--TRIT--TPK del estado enriquecido y de la estructura discreta del
continuo. Se trabaja sobre el registro dodecafásico ya generado y sobre
realizaciones declaradas de sus lectores. El portador lineal posterior
y los candidatos de una reconstrucción no reemplazan el dominio de
historias generadas. Tampoco una realización reversible posterior se
identifica, por su sola notación, con la holonomía nonádica \(\Gamma_9\).

En la coordenatización y orientación del registro,
\begin{equation}
 \begin{gathered}
 K=(234,543,140,729,659,824,621,58,914,794,146,601),\\
 \mathbf1^{*}K=6263,\qquad
 (Sx)_i=x_{i+1},\quad i\pmod {12}.
 \end{gathered}
 \label{md:antmem:registro}
\end{equation}
El registro se usa después de su generación, no como entrada
retrospectiva para seleccionar una salida. El operador \(S\) es aquí
el desplazamiento posicional marcado. Sus dominios se mantendrán
separados de los planos canónicos utilizados más adelante.

\subsection{La órbita completa como referencia de su portador}
\label{md:antmem:k:orbita}

En \(V=\mathbb R^{12}\), con su producto euclídeo canónico, consideremos
\[
 O_K=(K,SK,\ldots,S^{11}K):\mathbb R^{12}\longrightarrow V.
\]
El desplazamiento \(S\) actúa sobre posiciones del registro. No se
identifica por la sola periodicidad con la evolución completa TPK o con
una simetría excepcional posterior.

\begin{theorem}[Referencia cíclica completa y Gram exacto]
\label{md:antmem:k:marco-ciclico}
La matriz \(O_K\) es invertible. Su Gram tiene los valores propios
\[
\begin{array}{c|c}
 \text{valor}&\text{multiplicidad}\\\hline
 39225169=6263^2&1\\
 1248025-345420\sqrt3&2\\
 589609&2\\
 1407529&2\\
 1053157&2\\
 1248025+345420\sqrt3&2\\
 697225&1
\end{array}
\]
y satisface
\begin{equation}
 589609\|c\|^2\leq\|O_Kc\|^2\leq39225169\|c\|^2,
 \quad
 \|O_K^{-1}\|=\frac1{\sqrt{589609}},\quad
 \operatorname{cond}(O_K)=\frac{6263}{\sqrt{589609}}.
 \label{md:antmem:k:cotas-orbita}
\end{equation}
La órbita del registro centrado tiene rango once.
\end{theorem}
\begin{proof}
Las entradas del Gram son las autocorrelaciones
\(\langle K,S^jK\rangle\). En orden \(j=0,\ldots,11\) valen
\begin{align*}
 (&4251257,2999323,3163385,3287920,3216553,3344743,\\
  &2950064,3344743,3216553,3287920,3163385,2999323).
\end{align*}
Como \(S\) es ortogonal, el Gram es circulante. Sustituir los doce
caracteres \(\exp(2\pi i m/12)\) en esa fila da la tabla real
anterior. Todos los valores son positivos. El menor de los dos valores
con \(\sqrt3\) es mayor que \(589609\), pues
\(658416>345420\sqrt3\), desigualdad que se comprueba elevando al
cuadrado sus dos miembros positivos. Los demás valores de la tabla son
mayores o iguales a \(589609\). El modo uniforme tiene valor
\((\sum K_i)^2=6263^2\); la desigualdad triangular de la transformada
de Fourier, usando \(K_i\geq0\), acota todos los demás por ese mismo
valor. El teorema espectral da \eqref{md:antmem:k:cotas-orbita}. Centrar
el registro elimina sólo el modo uniforme; los once restantes conservan
los valores no nulos de la tabla.
\end{proof}

El determinante, con el orden de columnas y la orientación de \(S\)
fijados aquí, es
\[
 \det O_K=5483119259214020183850397930008283625.
\]
La ausencia de un período propio del vector, por sí sola, no probaría la
invertibilidad: lo decisivo es que no se anula ninguno de sus modos de
Fourier. La condición numérica es aproximadamente \(8.15643\). Si se
normaliza \(K\) a norma uno, los extremos del Gram se dividen por
\(\|K\|^2=4251257\); la amplitud de sus componentes no debe
confundirse con una mejora intrínseca frente a toda referencia de igual
norma.

\begin{corollary}[Separación polar de forma y orientación]
\label{md:antmem:k:polar}
Sean
\[
 G_K=O_K^*O_K,\qquad P_K=G_K^{1/2},\qquad
 F_K=O_KG_K^{-1/2}.
\]
Entonces \(P_K\) es positivo definido, \(F_K\) es ortogonal y
\(O_K=F_KP_K\). Una realización isométrica \(J:V\to\mathcal Y\)
conserva \(G_K\) y \(P_K\), y transporta la parte isométrica a
\(JF_K\).
\end{corollary}
\begin{proof}
La positividad del Gram permite definir su raíz e inversa mediante el
teorema espectral. La identidad
\(F_K^*F_K=G_K^{-1/2}G_KG_K^{-1/2}=I\) prueba la afirmación.
Además, \((JO_K)^*(JO_K)=O_K^*O_K\).
\end{proof}

Esta polar se define a partir del Gram completo: no se obtiene sólo de la
lista de sus valores propios ni de la norma de \(K\). El símbolo
\(P_K\) de este corolario no es el proyector dimensional de rango diez
de otra realización incidencial. Esta separación es una consecuencia de la
referencia cíclica, no una identificación de todos sus lectores.

\subsection{Identificación de operadores y holonomía relativa}
\label{md:antmem:k:identificacion}

\begin{theorem}[Doce respuestas generales o una respuesta cíclica]
\label{md:antmem:k:operadores}
Para cualquier operador lineal \(H:V\to V\), sus doce respuestas
\(Y=(HK,HSK,\ldots,HS^{11}K)\) determinan
\[
 H=YO_K^{-1}.
\]
Si \(HS=SH\), una única respuesta vectorial completa \(y=HK\)
determina
\begin{equation}
 H=O_yO_K^{-1},\qquad
 H=\sum_{j=0}^{11}h_jS^j,\qquad h=O_K^{-1}y.
 \label{md:antmem:k:filtro}
\end{equation}
Con errores \(E\) en la matriz de respuestas o \(\delta y\) en el
vector, respectivamente,
\begin{align}
 \|\widehat H-H\|&\leq\frac{\|E\|}{\sqrt{589609}},
 \label{md:antmem:k:error-matriz}\\
 \|\widehat h-h\|&\leq\frac{\|\delta y\|}{\sqrt{589609}},\qquad
 \|\widehat H-H\|\leq\sqrt{\frac{12}{589609}}\,\|\delta y\|.
 \label{md:antmem:k:error-filtro}
\end{align}
\end{theorem}
\begin{proof}
La primera igualdad es \(Y=HO_K\). Si \(H\) conmuta con \(S\),
entonces \(HS^jK=S^jy\), de modo que \(Y=O_y\). La conmutación con
el ciclo determina todas las columnas de \(H\) desde una de ellas;
las doce potencias \(I,S,\ldots,S^{11}\) son independientes y generan
ese espacio de operadores. Por tanto \(y=O_Kh\), lo que da la
segunda fórmula. Las dos primeras cotas usan
\(\|O_K^{-1}\|=1/\sqrt{589609}\). La última usa
\(\|O_{\delta y}\|\leq\|O_{\delta y}\|_{\rm F}
=\sqrt{12}\|\delta y\|\).
\end{proof}

Una respuesta significa aquí un vector de doce componentes, no una
medición escalar. Como ejemplo, si \(H=2I+3S-S^5\), la respuesta
\(y=HK\) devuelve exactamente los coeficientes \(2,3,-1\) en sus
posiciones mediante \eqref{md:antmem:k:filtro}; el inversor necesita \(y,K,S\),
no los coeficientes del filtro. En cambio, el operador no nulo cuya primera
fila es \((543,-234,0,\ldots,0)\) y cuyas restantes filas son nulas
anula \(K\). Una sola respuesta no identifica un operador arbitrario.

La función de referencia admite una generalización exacta. En un ciclo de
\(n\) posiciones, \(H\mapsto Hg\) identifica todos los operadores
que conmutan con el ciclo si y sólo si
\(O_g=(g,Sg,\ldots,S^{n-1}g)\) es invertible. La suficiencia es la
prueba anterior. Si \(O_g\) es singular, un vector no nulo de su núcleo
da un polinomio de grado menor que \(n\) que anula \(g\); el operador
es no nulo porque las potencias del ciclo son independientes. El impulso
unitario también tiene órbita ortonormal y sirve de referencia. La
consecuencia específica para \(K\) es que el mismo objeto ya usado en
clausura e incidencia satisface esa condición, sin haber sido elegido de
nuevo para optimizar el ensayo. No se deduce una identificación de
dinámicas no lineales ni de estados ocultos arbitrarios; una aplicación a
una linealización necesita que ésta se haya definido y justificado.

\begin{corollary}[Recuperación de holonomía relativa desde el complemento]
\label{md:antmem:k:holonomia}
Sean \(H_B\) conocido e invertible y \(H_C\) otro transporte del
mismo tipo. El complemento
\[
 D_\gamma=\frac{\sqrt8}{9}(H_C-H_B)
\]
determina
\[
 H_B^{-1}H_C=I+\frac9{\sqrt8}H_B^{-1}D_\gamma.
\]
Si \(D_\gamma\) conmuta con \(S\), en particular si lo hacen ambos
transportes, basta \(y=D_\gamma K\) para recuperar
\begin{equation}
 H_B^{-1}H_C
 =I+\frac9{\sqrt8}H_B^{-1}O_yO_K^{-1}.
 \label{md:antmem:k:holonomia-una-respuesta}
\end{equation}
El error en esta reconstrucción está acotado por
\[
 \frac9{\sqrt8}\|H_B^{-1}\|
 \sqrt{\frac{12}{589609}}\,\|\delta y\|.
\]
\end{corollary}
\begin{proof}
Se despeja \(H_C\) en la definición del complemento y se aplica
\cref{md:antmem:k:operadores}. La cota resulta de la submultiplicatividad.
\end{proof}

Para \(H_B=S^3\) y \(H_C=S^4\), el dato racional
\(y/\sqrt8=(S^4-S^3)K/9\) permite reconstruir la holonomía relativa
\(S\) sin proporcionar \(H_C\) al inversor. Si \(H_B\) es
unitario, el factor de error es aproximadamente \(0.01435510\). Sin
conmutación se mantienen las doce respuestas generales, no una hipótesis
ficticia de simetría. La selección de un lazo físico y su representación
siguen perteneciendo a su dinámica generadora; este protocolo no las
sustituye.

\subsubsection{Consecuencia de la referencia marcada para dos orientaciones}
\label{md:antmem:orientaciones}
La consecuencia utilizada en las notas del cuaderno conserva la tabla
anterior como antecedente. Para un operador real \(H\) que conmuta con
\(S\), sean \(y_+=HK\) e \(y_-=H^{\mathsf T}K\).
Sus doce intensidades de Fourier coinciden, pero \(y_+=y_-\) si y sólo
si \(H=H^{\mathsf T}\). Cada respuesta vectorial completa, junto con
\(K\) y la orientación de \(S\), recupera su operador.
\begin{proof}
El operador \(H\) es circulante; sus multiplicadores de Fourier son
\(\lambda_m\), y los de \(H^{\mathsf T}\) son sus conjugados.
Por ello ambas respuestas tienen intensidades
\(|\lambda_m|^2|\widehat K_m|^2\).
Si sus vectores coinciden, \((H-H^{\mathsf T})K=0\).
La conmutación implica que \(H-H^{\mathsf T}\) anula también \(S^jK\)
para todo \(j\). Esos vectores forman una base, así que
\(H-H^{\mathsf T}=0\). La recuperación es
\(H=O_{y_+}O_K^{-1}\), y análogamente para \(H^{\mathsf T}\).
\end{proof}
En particular, los desplazamientos \(S\) y \(S^{-1}\) dan respuestas
de igual espectro de potencia, pero distintas. Las autocorrelaciones
exactas de la tabla anterior dan la identidad ya utilizada en el cuaderno:
\begin{equation}
 \|SK-S^{-1}K\|^2
 =2\bigl(\|K\|^2-\langle K,S^2K\rangle\bigr)
 =2\,175\,744>0.
 \label{md:antmem:distancia-orientada}
\end{equation}
Este número es una consecuencia de las autocorrelaciones, no un dato
de partida adicional. El control racional del cuaderno recupera de ambas
respuestas los coeficientes \(e_1\) y \(e_{11}\), respectivamente, por
inversión de \(O_K\). Una intensidad espectral o un balance escalar no
sustituye la respuesta vectorial.
Este corolario concierne al portador cíclico de doce componentes:
no identifica \(S\) con la dinámica enriquecida completa ni convierte
por sí solo \(S\leftrightarrow S^{-1}\) en conjugación física de carga.
Para \(H\) general, la transposición tampoco significa inversión temporal;
en el caso \(S\), sí coincide con \(S^{-1}\).

\subsection{Acoplamiento reversible y transporte no estacionario}
\label{md:antmem:acoplamiento}
\subsubsection{Extensión reversible del balance entre dos formas}

Sean \(B,C:V\to V'\) isomorfismos que conservan una forma no degenerada \(b\) entre dos realizaciones:

\[
B^*b'B=C^*b'C=b.
\]

La forma puede ser una métrica positiva o una forma alternante. En espacios reales, \(*\) denota transposición respecto de las coordenadas elegidas; la igualdad expresa el transporte de la forma entre las realizaciones especificadas. En espacios de Hilbert, los operadores y sus inversos son acotados.

Con \(s=\sqrt8\) y \(Q_V=\frac13\left(\begin{smallmatrix}sI&-I\\I&sI\end{smallmatrix}\right)\), definimos

\begin{equation}
\boxed{
\mathcal U(B,C)=Q_{V'}^*\operatorname{diag}(B,C)Q_V
=\frac19\begin{pmatrix}8B+C&\sqrt8(C-B)\\\sqrt8(C-B)&B+8C\end{pmatrix}.}
\label{md:antmem:acop:eq:1}
\end{equation}

Escribimos \(T=(8B+C)/9\), \(D=\sqrt8(C-B)/9\) y \(R=(B+8C)/9\). La actualización completa es

\begin{equation}
\binom{x'}{m'}=\begin{pmatrix}T&D\\D&R\end{pmatrix}\binom{x}{m}.
\label{md:antmem:acop:eq:2}
\end{equation}

\begin{theorem}
La aplicación \eqref{md:antmem:acop:eq:1} es reversible y conserva la forma \(b\oplus b\):

\begin{equation}
\mathcal U(B,C)^*(b'\oplus b')\mathcal U(B,C)=b\oplus b,
\quad\mathcal U(B,C)^{-1}=\mathcal U(B^{-1},C^{-1}).
\label{md:antmem:acop:eq:3}
\end{equation}
\end{theorem}

\begin{proof}
Se cumple \(Q^*Q=I\) por \(8+1=9\). Sus bloques escalares hacen que \(Q\) preserve \(b\oplus b\) para cualquier forma \(b\). La aplicación diagonal conserva las formas por hipótesis. La composición prueba \eqref{md:antmem:acop:eq:3}, y la inversión de los tres factores da la inversa declarada.
\end{proof}

Para memoria entrante \(m=0\), se recupera el balance \(b=T^*b'T+D^*b'D\). La segunda columna determina la respuesta a memoria entrante arbitraria. La conservación de normas positivas y la conservación de área orientada son las especializaciones respectivas de \eqref{md:antmem:acop:eq:3}.

\paragraph{Realización colectiva en las nueve copias}

En \(V^9\), sea \(W_E=\operatorname{diag}(I,\ldots,I,E)\), con \(E\) en la novena posición. Sean \(A(a,b)=(a/\sqrt8,\ldots,a/\sqrt8,b)\) y \(L=AQ\). La inclusión uniforme es \(Jx=(x,\ldots,x)/3\). Entonces \(L(x,0)=Jx\) y

\begin{equation}
W_E L=L\mathcal U(I,E).
\label{md:antmem:acop:eq:4}
\end{equation}

\begin{proof}
\(W_EA=A\operatorname{diag}(I,E)\); multiplicar por \(Q\) y usar \(AQ=L\) da \eqref{md:antmem:acop:eq:4}. La ortogonalidad de la media de las ocho primeras posiciones con sus siete diferencias muestra además que el complemento queda fijo bajo \(W_E\).
\end{proof}

Sea \(P_{\rm cyc}(v_0,\ldots,v_8)=(v_8,v_0,\ldots,v_7)\). El paso cíclico original es \(S_E=P_{\rm cyc}W_E\). Por ello \(S_EL=(P_{\rm cyc}L)U(I,E)\), con coordenatizaciones de entrada y salida distintas. La reducción colectiva conserva el patrón correspondiente: \(S_E^9=I_9\otimes E\), mientras \(W_E^9=\operatorname{diag}(I_8,E^9)\). Cuando un lector utiliza las nueve etiquetas individuales, éstas se conservan en su representación original.

\subsubsection{Composición exacta y reentrada de memoria}

Para \(B_1,C_1:V_0\to V_1\) y \(B_2,C_2:V_1\to V_2\),

\begin{equation}
\boxed{\mathcal U(B_2,C_2)\mathcal U(B_1,C_1)
=\mathcal U(B_2B_1,C_2C_1).}
\label{md:antmem:acop:eq:5}
\end{equation}

\begin{proof}
En \eqref{md:antmem:acop:eq:1} se cancelan \(Q_{V_1}Q_{V_1}^*\), y los dos bloques diagonales se multiplican en el orden indicado. La inducción extiende la fórmula a cualquier número finito de transportes, conservando su orden.
\end{proof}

El bloque de lectura contiene la identidad

\begin{equation}
\boxed{T_2T_1+D_2D_1=\frac{8B_2B_1+C_2C_1}{9}.}
\label{md:antmem:acop:eq:6}
\end{equation}

Desde \((x,0)\), la primera etapa produce \((T_1x,D_1x)\). La segunda lectura es \(T_2T_1x+D_2D_1x\). El segundo término es exactamente la contribución de la memoria reintroducida. Al almacenar \(D_1x\) fuera del siguiente acoplamiento y continuar sólo con \(T_1x\) se obtiene el procedimiento de registros separados, cuya lectura es \(T_2T_1x\). La igualdad \eqref{md:antmem:acop:eq:6} precisa la diferencia entre ambos procedimientos.

La realización \(U\) conserva el estado del acoplamiento y los productos ordenados \(B_N\cdots B_1\), \(C_N\cdots C_1\). La conservación de todos los prefijos incluye además el registro de emisiones de HMT, que mantiene cada factor junto con su orden.

\paragraph{Ecuación covariante de incremento}

Para el registro separado \(x_{n+1}=T_nx_n\), \(m_n=D_nx_n\), se tiene

\begin{equation}
x_{n+1}-B_nx_n=m_n/\sqrt8.
\label{md:antmem:acop:eq:7}
\end{equation}

Con \(G_0=I\) y \(G_{n+1}=B_nG_n\), la telescopía da

\begin{equation}
x_0=G_N^{-1}x_N-\frac1{\sqrt8}\sum_{n<N}G_{n+1}^{-1}m_n.
\label{md:antmem:acop:eq:8}
\end{equation}

La memoria es el incremento respecto del transporte entre formas \(B_n\), con la normalización declarada. La igualdad \eqref{md:antmem:acop:eq:8} es una reconstrucción finita algebraica; su paso a una serie infinita requiere convergencia. El teorema siguiente construye la inversa estable mediante el límite de operadores positivos.

\subsubsection{Memoria de una sucesión no estacionaria}

En una realización de Hilbert \(H\), sean \(C_n\) unitarios, con su orden de composición declarado. Escribimos

\[
T_n=(8I+C_n)/9,\quad D_n=\sqrt8(C_n-I)/9,
\quad P_0=I,\quad P_{n+1}=T_nP_n,\quad A_N=P_N^*P_N.
\]

El índice \(N\) cuenta compresiones con almacenamiento separado, a diferencia del transporte completo con reentrada. Para métricas variables con \(B_n\) isomorfismos isométricos, la conjugación por \(G_n\) lleva los transportes a un Hilbert fijo y produce este mismo caso, con \(\widetilde C_n=G_{n+1}^{-1}C_nG_n\). La afirmación utiliza las métricas de las fibras y las isometrías declaradas entre ellas.

\begin{theorem}[Terminal positivo y recuperación no estacionaria]
\label{md:antmem:terminal-no-estacionario}
Existe un operador positivo \(0\leq A_\infty\leq I\), límite fuerte de \(A_N\), y

\begin{equation}
\boxed{I=A_\infty+\sum_{n\ge0}P_n^*D_n^*D_nP_n.}
\label{md:antmem:acop:eq:9}
\end{equation}

La aplicación

\begin{equation}
\boxed{\mathcal Vv=(A_\infty^{1/2}v,(D_nP_nv)_{n\ge0})}
\label{md:antmem:acop:eq:10}
\end{equation}

es isométrica y tiene inversa izquierda estable

\begin{equation}
\mathcal V^*(a,(m_n))=A_\infty^{1/2}a+
\sum_{n\ge0}P_n^*D_n^*m_n,\qquad\|\mathcal V^*\|=1.
\label{md:antmem:acop:eq:11}
\end{equation}

\end{theorem}
\begin{proof}
La identidad local \(T_n^*T_n+D_n^*D_n=I\) da
\(A_N-A_{N+1}=P_N^*D_N^*D_NP_N\geq0\). Las formas cuadráticas de
la sucesión positiva decreciente convergen a la de un operador positivo
\(A_\infty\). Para \(B_N=A_N-A_\infty\), la desigualdad
\(0\leq B_N^2\leq B_N\) convierte esa convergencia en fuerte.
La telescopía finita y el límite prueban \eqref{md:antmem:acop:eq:9}. Evaluar sobre \(v\)
prueba \eqref{md:antmem:acop:eq:10}. Su adjunto es \eqref{md:antmem:acop:eq:11}; la serie converge en norma porque
las truncaciones de un registro cuadrado-sumable convergen y el adjunto es acotado.
\end{proof}

La reconstrucción con el terminal límite y \(N\) registros tiene error
exacto \((A_N-A_\infty)v\). El vector \(A_\infty^{1/2}v\) es la
realización positiva canónica del terminal de norma. Su existencia procede
del límite de Gram, con independencia de la convergencia vectorial de \(P_Nv\).

\subsubsection{Tres propiedades diferentes: agotamiento, distinción y estabilidad}

Sea \(Mv=(D_nP_nv)_n\). De \eqref{md:antmem:acop:eq:9}, \(M^*M=I-A_\infty\). Por tanto:

\begin{enumerate}
\item Toda la norma de \(v\) pasa al registro exactamente cuando \(A_\infty v=0\).
\item El registro distingue cada estado exactamente cuando \(\ker(I-A_\infty)=\{0\}\).
\item La recuperación desde memoria sola tiene inversa acotada exactamente cuando
\(I-A_\infty\geq\varepsilon I\) para algún \(\varepsilon>0\).
En ese caso una inversa izquierda es \((I-A_\infty)^{-1}M^*\).
\end{enumerate}

Además,

\begin{equation}
\boxed{\ker M=\bigcap_{n\ge0}\operatorname{Fix}(C_n).}
\label{md:antmem:acop:eq:12}
\end{equation}

\begin{proof}
Si \(Mv=0\), entonces \(D_0v=0\), de donde \(C_0v=v\)
y \(P_1v=v\). Inductivamente, \(D_nP_nv=0\) da \(C_nv=v\) y
\(P_{n+1}v=v\). El recíproco es inmediato. Las tres equivalencias se
deducen del Gram \(M^*M\) y de la caracterización de los operadores
acotados inferiormente.
\end{proof}

\paragraph{Ejemplo exacto que separa las tres propiedades}

En el ejemplo algebraico \(C_0=-I\) y \(C_n=I\) para \(n\geq1\),

\begin{equation}
A_\infty=\frac{49}{81}I,\qquad M^*M=\frac{32}{81}I,
\qquad m_0=-\frac{2\sqrt8}{9}v,
\qquad v=-\frac9{2\sqrt8}m_0.
\label{md:antmem:acop:eq:13}
\end{equation}

La memoria contiene \(32/81\) de la norma al cuadrado y permite recuperar el vector completo. El terminal retiene \(49/81\) y es un operador positivo distinto de un proyector. La recuperabilidad depende del núcleo y del condicionamiento del lector. Lectura y memoria constituyen componentes correlacionadas de una isometría.

Este ejemplo separa las tres consecuencias lógicas del teorema. La
aplicación a una trayectoria física conserva además la selección y la
procedencia efectiva de sus \(C_n\) mediante el TPK.

La letra \(\eta\) de la realización gaussiana siguiente designa la
deformación de la elipse, es decir, el parámetro \(\zeta\) de los
desarrollos de transporte; no es el coeficiente de acción
\(\eta_{\rm ret}\). Se mantiene una sección positiva \(\hbar_a\) y
una misma unidad para las dos coordenadas del cociente angular.
Los semiejes conservados son
\[
 a_S=\sqrt{2\hbar_a}\,e^{\eta/2},\qquad
 b_S=\sqrt{2\hbar_a}\,e^{-\eta/2},\qquad
 \eta=\operatorname{artanh}(C^*/A).
\]

\subsection{Covarianza canónica y reciprocidad de la elipse de acción}
\label{md:antmem:covarianza-accion}

La construcción precedente determina dos invariantes complementarios:
el producto de los semiejes fija el área de acción y su cociente conserva
la anisotropía del par angular. Precisamos ahora la realización canónica
en la que esos semiejes son dos desviaciones estándar. La estructura de
partida es la elipse ya obtenida de las representaciones angulares y de
acción de APP--TRIT--TPK; la realización operatoria se establece después
de esa construcción.

En este apartado escribimos \(\eta=\eta_{\rm el}
=\operatorname{artanh}(C^*/A)\), con \(|C^*/A|<1\) y
\(\hbar_a>0\). El cociente angular utiliza una misma unidad para sus dos
componentes. En las coordenadas equilibradas \((Q,P)\), ambas de
dimensión raíz de acción, definimos
\begin{equation}
 V_\eta:=\frac14
 \begin{pmatrix}a_S^2&0\\0&b_S^2\end{pmatrix}
 =\frac{\hbar_a}{2}
 \begin{pmatrix}e^\eta&0\\0&e^{-\eta}\end{pmatrix}.
 \label{md:antmem:covarianza-geometrica}
\end{equation}
Esta matriz positiva tiene entradas de dimensión acción. Su interpretación
como covarianza cuántica queda especificada por la realización siguiente.

\begin{proposition}[Realización gaussiana de la covarianza de acción]
\label{md:antmem:covarianza-gaussiana}
Considérese la representación canónica en \(L^2(\mathbb R,dQ)\), con
\(\widehat Qf(Q)=Qf(Q)\) y
\(\widehat Pf(Q)=-i\hbar_a f'(Q)\), inicialmente sobre el espacio de
Schwartz \(\mathcal S(\mathbb R)\). El estado normalizado
\begin{equation}
 \psi_\eta(Q)=\bigl(\pi\hbar_a e^\eta\bigr)^{-1/4}
 \exp\!\left(-\frac{Q^2}{2\hbar_a e^\eta}\right)
 \label{md:antmem:gaussiana-accion}
\end{equation}
tiene medias nulas y matriz de covarianza \(V_\eta\). En particular,
\begin{equation}
 (\Delta Q)^2=\frac{\hbar_a e^\eta}{2},\qquad
 (\Delta P)^2=\frac{\hbar_a e^{-\eta}}{2},\qquad
 \operatorname{Cov}(Q,P)=0,\qquad
 \det V_\eta=\frac{\hbar_a^2}{4}.
 \label{md:antmem:varianzas-accion}
\end{equation}
Esta realización satura la desigualdad de Robertson--Schrödinger.
\end{proposition}

\begin{proof}
Los operadores anteriores son esencialmente autoadjuntos sobre
\(\mathcal S(\mathbb R)\), que constituye un núcleo común invariante;
sobre este espacio se cumple
\([\widehat Q,\widehat P]=i\hbar_a I\). Para
\(\widehat Y=(\widehat Q,\widehat P)\), la covarianza simétrica es
\[
 V_{jk}=\frac12\left\langle
 (\widehat Y_j-\langle\widehat Y_j\rangle)
 (\widehat Y_k-\langle\widehat Y_k\rangle)
 +(\widehat Y_k-\langle\widehat Y_k\rangle)
 (\widehat Y_j-\langle\widehat Y_j\rangle)
 \right\rangle.
\]
Con \(s=\hbar_a e^\eta>0\), sea
\(I_s=\int_{\mathbb R}e^{-Q^2/s}\,dQ\). El producto de dos
integrales iguales, evaluado en coordenadas polares, da
\(I_s^2=2\pi\int_0^\infty r e^{-r^2/s}\,dr=\pi s\).
Además, la integral de
\(\frac{d}{dQ}(Qe^{-Q^2/s})\) es cero, por el decaimiento en ambos
extremos; de ahí
\(\int_{\mathbb R}Q^2e^{-Q^2/s}\,dQ=sI_s/2\).
Estas identidades y la paridad prueban la normalización, la media nula y
\(\langle\widehat Q^2\rangle=s/2\). Además,
\(\widehat P\psi_\eta=i e^{-\eta}Q\psi_\eta\).
Por tanto, \(\langle\widehat P\rangle=0\),
\(\|\widehat P\psi_\eta\|^2=\hbar_a e^{-\eta}/2\) y
\(\operatorname{Re}\langle\widehat Q\psi_\eta,
\widehat P\psi_\eta\rangle=0\).

Para cualquier estado normalizado de ese dominio, Cauchy--Schwarz aplicada
a \((\widehat Q-\langle\widehat Q\rangle)\psi\) y
\((\widehat P-\langle\widehat P\rangle)\psi\) da
\[
 (\Delta Q)^2(\Delta P)^2
 \geq \operatorname{Cov}(Q,P)^2
       +\frac14\left|\langle[\widehat Q,\widehat P]\rangle\right|^2
 =\operatorname{Cov}(Q,P)^2+\frac{\hbar_a^2}{4}.
\]
Las varianzas calculadas alcanzan la igualdad. Equivalentemente, el
operador
\[
 a_\eta=\frac{e^{-\eta/2}\widehat Q
                  +i e^{\eta/2}\widehat P}{\sqrt{2\hbar_a}}
\]
satisface \([a_\eta,a_\eta^\dagger]=I\) y
\(a_\eta\psi_\eta=0\), por sustitución directa.
\end{proof}

\subsection{Covarianza, forma normal y espectro gaussiano}
\label{md:antmem:gauss:gaussianas-espectro}

La sección positiva de acción \(\hbar_a\) y la realización canónica
anterior determinan la normalización empleada aquí. Se demuestra la
correspondencia entre covarianza, espectro, pureza y entropía para un
número finito de modos; el espacio de Fock conserva todas sus ocupaciones.
La reducción simpléctica y las fórmulas gaussianas son resultados
establecidos que se reproducen para hacer autónoma su aplicación al
transporte de memoria.

\subsubsection{Coordenadas canónicas y función característica}
En \(L^2(\mathbb R^d)\), sean \(\widehat q_j=\widehat Q_j/\sqrt{\hbar_a}\),
\(\widehat p_j=\widehat P_j/\sqrt{\hbar_a}\) y
\(\widehat z=(\widehat q_1,\widehat p_1,\ldots,\widehat q_d,\widehat p_d)\).
Sobre Schwartz,
\[
 [\widehat z_j,\widehat z_k]=iJ_{jk}I,\qquad
 J=\bigoplus_{j=1}^d\begin{pmatrix}0&1\\-1&0\end{pmatrix}.
\]
Esta convención fija el signo de la forma canónica. Los operadores de
Weyl \(\mathcal W(\xi)=\exp(i\xi^{\mathsf T}\widehat z)\) se definen
por las clausuras autoadjuntas de las combinaciones lineales indicadas.
Para un estado centrado \(\rho\) de segundos momentos finitos, la
covarianza normalizada y su función característica son
\[
 W_{jk}=\operatorname{Tr}\rho(\widehat z_j\widehat z_k+\widehat z_k\widehat z_j),
 \qquad \chi_\rho(\xi)=\operatorname{Tr}(\rho\mathcal W(\xi)),\qquad
 V=\frac{\hbar_a}{2}W.
\]
El estado es gaussiano centrado cuando
\begin{equation}
 \chi_\rho(\xi)=\exp(-\xi^{\mathsf T}W\xi/4).
 \label{md:antmem:gauss:eq:caracteristica-gaussiana}
\end{equation}
La esperanza de \((c^{\mathsf T}\widehat z)^\dagger
(c^{\mathsf T}\widehat z)\) es \(c^\dagger(W+iJ)c/2\); así
\(W+iJ\geq0\). Además \(W>0\): si \(v\) es real y
\(v^{\mathsf T}Wv=0\), la positividad implica \((W+iJ)v=0\),
por tanto \(Jv=0\) y \(v=0\).

\begin{lemma}[Unicidad de la función característica]
\label{md:antmem:gauss:lem:unicidad-weyl}
Dos operadores de clase traza con la misma función característica son iguales.
\end{lemma}
\begin{proof}
Reagrupando posiciones y momentos como \(\xi=(u,v)\),
\[
 (\mathcal W(u,v)f)(x)=e^{iu\cdot(x+v/2)}f(x+v).
\]
Para un operador de rango finito con vectores de Schwartz y núcleo \(K_A\),
\[
 \operatorname{Tr}(A\mathcal W(u,v))=\int K_A(x+v/2,x-v/2)e^{iu\cdot x}\,dx.
\]
Plancherel y el cambio de variables de jacobiano absoluto uno dan
\[
 \|A\|_{\rm HS}^2=(2\pi)^{-d}\int_{\mathbb R^{2d}}
 |\operatorname{Tr}(A\mathcal W(\xi))|^2\,d\xi.
\]
Los operadores anteriores son densos en clase traza. La convergencia
en esa norma implica convergencia en Hilbert--Schmidt y convergencia
uniforme de las funciones características, porque su diferencia está
acotada por \(\|A-B\|_1\). La igualdad se extiende en \(L^2\).
Aplicarla a la diferencia de los operadores prueba la afirmación.
\end{proof}

\subsubsection{Reducción simpléctica de la forma positiva}
\begin{theorem}[Forma normal de Williamson]
\label{md:antmem:gauss:thm:williamson}
Para una matriz real simétrica \(W>0\) de orden \(2d\), existen
\(S\in\operatorname{Sp}(2d,\mathbb R)\) y \(\nu_j>0\) tales que
\[
 W=SDS^{\mathsf T},\qquad D=\bigoplus_j\nu_jI_2.
\]
Los autovalores de \(JW\) son \(\pm i\nu_j\). La condición
\(W+iJ\geq0\) equivale a \(\nu_j\geq1\) para todo \(j\).
\end{theorem}
\begin{proof}
La matriz \(B=W^{1/2}JW^{1/2}\) es antisimétrica e invertible.
Si \(-B^2u=\nu^2u\) y \(\|u\|=1\), el vector
\(v=-Bu/\nu\) es unitario y ortogonal a \(u\); cumple
\(Bu=-\nu v\), \(Bv=\nu u\). El complemento de esta pareja es
invariante. La iteración produce una matriz ortogonal \(O\) con
\(O^{\mathsf T}BO=JD\). Definiendo \(S=W^{1/2}OD^{-1/2}\),
\[
 S^{\mathsf T}JS=J,\qquad SDS^{\mathsf T}=W.
\]
La semejanza \(W^{1/2}(JW)W^{-1/2}=B\) determina los valores
\(\nu_j\), con multiplicidad. La congruencia por \(S^{-1}\)
transforma \(W+iJ\) en \(D+iJ\); cada bloque tiene autovalores
\(\nu_j+1\) y \(\nu_j-1\).
\end{proof}

\begin{lemma}[Implementación unitaria en un número finito de modos]
\label{md:antmem:gauss:lem:implementacion-simplectica}
Toda matriz simpléctica \(S\) admite un unitario \(U_S\) que satisface
\(U_S^\dagger\mathcal W(\xi)U_S=\mathcal W(S^{\mathsf T}\xi)\).
La conjugación \(\rho\mapsto U_S\rho U_S^\dagger\) transporta
la covarianza por \(W\mapsto SWS^{\mathsf T}\).
\end{lemma}
\begin{proof}
En la descomposición polar \(S=OP\), la identidad
\(J^{-1}(S^{\mathsf T}S)J=(S^{\mathsf T}S)^{-1}\), aplicada por
cálculo espectral a la raíz positiva, prueba \(PJP=J\).
Entonces \(O\) es ortogonal simpléctica. Los autoespacios de \(P\)
se emparejan mediante \(J\) con autovalores \(\lambda,\lambda^{-1}\);
el de autovalor uno es \(J\)-invariante. Por ello una matriz
ortogonal simpléctica \(K\) reduce \(P\) a
\(\bigoplus_j\operatorname{diag}(e^{r_j},e^{-r_j})\).

La dilatación se implementa por
\[
 (\mathcal D_rf)(x)=e^{-\frac12\sum_jr_j}
 f(e^{-r_1}x_1,\ldots,e^{-r_d}x_d).
\]
Un cambio de variables prueba su unitariedad y la conjugación de
\((\widehat q_j,\widehat p_j)\) a
\((e^{r_j}\widehat q_j,e^{-r_j}\widehat p_j)\).
Una transformación ortogonal simpléctica conmuta con \(J\) y
corresponde a una matriz \(u\in U(d)\) sobre
\(a_j=(\widehat q_j+i\widehat p_j)/\sqrt2\). Su segunda
cuantización \(\Gamma(u)=\bigoplus_{n\geq0}u^{\otimes_s n}\)
es unitaria en \(\mathcal F_s(\mathbb C^d)\), y la identidad sobre
productos simétricos finitos da la transformación de los \(a_j\).

La base de Hermite identifica este espacio con \(L^2(\mathbb R^d)\).
Para justificar su completitud, si \(f\) es ortogonal a todos los
polinomios multiplicados por \(e^{-|x|^2/2}\), la función entera
\(F(z)=\int f(x)e^{-|x|^2/2}e^{z\cdot x}\,dx\) tiene nulas
todas sus derivadas en cero. Así \(F=0\); restringiendo a
\(z=i\xi\), la unicidad de Fourier da \(f=0\). Las
normalizaciones y la ortogonalidad resultan de creación y aniquilación.
La composición de las implementaciones anteriores produce \(U_S\).
La igualdad de Weyl implica
\(\chi_{U_S\rho U_S^\dagger}(\xi)=\chi_\rho(S^{\mathsf T}\xi)\),
que demuestra la ley de covarianza.
\end{proof}

\subsubsection{Cálculo del espectro de la densidad isotrópica}
En un modo, los estados coherentes normalizados son
\[
 |z\rangle=e^{-|z|^2/2}\sum_{n\geq0}\frac{z^n}{\sqrt{n!}}|n\rangle.
\]
Su función de posición,
\(\pi^{-1/4}\exp[-(x-q_0)^2/2+ip_0(x-q_0/2)]\), donde
\(q_0=\sqrt2\Re z\), \(p_0=\sqrt2\Im z\), da por integración
\[
 \langle z|\mathcal W(u,v)|z\rangle
 =e^{-(u^2+v^2)/4}e^{i(uq_0+vp_0)}.
\]
Para \(s>0\), la integral
\[
 \rho^{(s)}=\frac1{\pi s}\int_{\mathbb C} e^{-|z|^2/s}
 |z\rangle\langle z|\,d^2z
\]
converge en clase traza, es positiva y tiene traza uno. La integración
angular anula sus elementos fuera de la diagonal de ocupación; la
integración radial da
\[
 \langle n|\rho^{(s)}|n\rangle=\frac{s^n}{(1+s)^{n+1}},\qquad
 \chi_{\rho^{(s)}}(u,v)=e^{-(2s+1)(u^2+v^2)/4}.
\]
Con \(\nu=2s+1\), el único estado gaussiano de covarianza
\(\nu I_2\) es, por el lema~\ref{md:antmem:gauss:lem:unicidad-weyl},
\begin{equation}
 \rho_\nu=(1-t)\sum_{n\geq0}t^n|n\rangle\langle n|,
 \qquad t=\frac{\nu-1}{\nu+1}.
 \label{md:antmem:gauss:eq:estado-geometrico}
\end{equation}
El caso \(\nu=1\) es el proyector sobre el vacío.

\begin{theorem}[Espectro, pureza y entropía]
\label{md:antmem:gauss:thm:espectro-gaussiano}
Un estado gaussiano centrado con covarianza \(W\) y valores
simplécticos \(\nu_1,\ldots,\nu_d\) es unitariamente equivalente
a \(\bigotimes_j\rho_{\nu_j}\). Sus autovalores son
\[
 \lambda_{\mathbf n}=\prod_j(1-t_j)t_j^{n_j},\qquad
 t_j=\frac{\nu_j-1}{\nu_j+1},\quad \mathbf n\in\mathbb N^d.
\]
Con \(0\log0=0\),
\begin{align}
 \operatorname{Tr}\rho^2&=\prod_j\nu_j^{-1}=(\det W)^{-1/2},
 \label{md:antmem:gauss:eq:pureza-general}\\
 -\operatorname{Tr}(\rho\log\rho)&=\sum_j
 \left[\frac{\nu_j+1}{2}\log\frac{\nu_j+1}{2}
 -\frac{\nu_j-1}{2}\log\frac{\nu_j-1}{2}\right].
 \label{md:antmem:gauss:eq:entropia-general}
\end{align}
\end{theorem}
\begin{proof}
Williamson da \(W=SDS^{\mathsf T}\). El producto de las densidades
isotrópicas tiene función característica \(e^{-\xi^{\mathsf T}D\xi/4}\).
La implementación unitaria y la unicidad de Weyl prueban la
equivalencia y la fórmula espectral. En un modo,
\[
 \sum_{n\geq0}(1-t)^2t^{2n}=\frac{1-t}{1+t}=\nu^{-1},
\]
y
\[
 -\sum_{n\geq0}(1-t)t^n\log((1-t)t^n)
 =-\log(1-t)-\frac{t}{1-t}\log t.
\]
La sustitución \(t=(\nu-1)/(\nu+1)\) prueba las fórmulas; el
límite en \(t=0\) es cero. En el producto finito, Tonelli justifica
la suma de los términos de entropía no negativos, y
\(\det W=\prod_j\nu_j^2\) da la expresión determinantal.
\end{proof}

Si el estado es la reducción de un vector puro bipartito, sus
autovalores son los cuadrados de los coeficientes de Schmidt.
La descomposición singular del operador de Hilbert--Schmidt definido
por ese vector demuestra la igualdad de espectros no nulos de ambas
reducciones. La entropía calculada es entonces la de entrelazamiento.

\subsubsection{Preparación y transporte de las dos covarianzas}
La aplicación posterior utiliza dos entradas gaussianas puras idénticas.
En general, para \(M\in\operatorname{Sp}(2d,\mathbb R)\), su covarianza es
\(V_0=(\hbar_a/2)MM^{\mathsf T}\), y la entrada conjunta tiene
covarianza \(V_0\oplus V_0\). El transporte físico se obtiene
conjugando \(U_H\) por \(M\oplus M\). En la coordenatización equilibrada,
la covarianza normalizada de entrada es \(I_{4d}\); con
\[
 T=\frac{8I+H}{9},\qquad D=\frac{\sqrt8(H-I)}9,
\]
la primera reducción tiene
\[
 W=TT^{\mathsf T}+DD^{\mathsf T}=\frac{8I+HH^{\mathsf T}}9.
\]
Su función característica se obtiene anulando los argumentos del
segundo grupo de modos. Por ello permanece gaussiana y los teoremas
anteriores calculan todo su espectro, sin truncar el registro de ocupación.

\subsection{Prolongación del transporte de memoria: realización cuántica}

La realización cuántica utiliza el acoplamiento reversible de la sección~\ref{md:antmem:acoplamiento}. Consideramos primero un plano canónico real con matriz alternante \(\Omega=\left(\begin{smallmatrix}0&1\\-1&0\end{smallmatrix}\right)\). La partición nonádica determina

\begin{equation}
Q=\frac13\begin{pmatrix}\sqrt8 I&-I\\I&\sqrt8 I\end{pmatrix},\qquad
U_H=Q^{\mathsf T}\operatorname{diag}(I,H)Q
=\frac19\begin{pmatrix}8I+H&\sqrt8(H-I)\\
\sqrt8(H-I)&I+8H\end{pmatrix}.
\label{md:antmem:cuantica:eq:5.1}
\end{equation}

El operador \(H\in\operatorname{Sp}(2,\mathbb R)\) transporta ese plano. La matriz \(Q\) conserva el producto euclídeo y la forma \(\Omega\oplus\Omega\); por ello \(U_H\) es simpléctica. La primera componente define el sistema observado y la segunda el registro de memoria; la transformación completa actúa sobre ambos.

El lazo ordenado cuyas identidades se reproducen en
\ref{md:antmem:lazo-ep} es

\begin{equation}
C=\begin{pmatrix}0&-1\\1&-1\end{pmatrix},\quad
P=\begin{pmatrix}1&1\\0&1\end{pmatrix},\quad
H_n=C^{-1}P^{-n}CP^n
=\begin{pmatrix}1+n&n^2\\n&n^2-n+1\end{pmatrix}.
\label{md:antmem:cuantica:eq:5.2}
\end{equation}

Para cada \(n\in\mathbb Z\) se cumple \(\det H_n=1\). En particular, para \(n=1\),

\begin{equation}
H_1=\begin{pmatrix}2&1\\1&1\end{pmatrix}
=F_1^2,\quad F_1=\begin{pmatrix}1&1\\1&0\end{pmatrix}
=P F_{\rm av}P^{-1}.
\label{md:antmem:cuantica:eq:5.3}
\end{equation}

Su espectro es \(\{\varphi^2,\varphi^{-2}\}\), por el reconocimiento posterior del modo áureo ya generado. El entero \(n\) especifica el número orientado de iteraciones en el lazo elegido. La colección de transportes está determinada algebraicamente; su realización como protocolo físico conserva además los datos de preparación y de evolución correspondientes.

\subsubsection{Identidades algebraicas del lazo y reconocimiento posterior}
\label{md:antmem:lazo-ep}
Para conservar la prueba de las identidades que se utilizan, se reproduce
su composición matricial. En la misma coordenatización,
\[
 N_0=\begin{pmatrix}0&1\\0&0\end{pmatrix}.
\]
Sea
\[
 P=I+N_0=\begin{pmatrix}1&1\\0&1\end{pmatrix}.
\]
Como \(N_0^2=0\), se cumple \(P^n=I+nN_0\) para todo
\(n\in\mathbb Z\).
El lazo ordenado de cuatro transportes da

\begin{equation}
H_n=C^{-1}P^{-n}CP^n
=\begin{pmatrix}1+n&n^2\\n&n^2-n+1\end{pmatrix}.
\label{md:antmem:matriz:eq:28}
\end{equation}

Para \(n\ne0\), su diferencia normalizada es

\[
F_n=\frac{H_n-I}{n}
=\begin{pmatrix}1&n\\1&n-1\end{pmatrix}.
\]

Multiplicando esas matrices se obtiene, sin introducir valores espectrales,

\begin{equation}
\boxed{F_n^2=H_n,\qquad F_n^2-nF_n-I=0.}
\label{md:antmem:matriz:eq:29}
\end{equation}

Con \(H_B=I\), la memoria es
\(D_{\gamma,n}=(\sqrt8/9)nF_n\).
En consecuencia,

\begin{equation}
\boxed{H_n=\left(\frac9{n\sqrt8}D_{\gamma,n}\right)^2.}
\label{md:antmem:matriz:eq:30}
\end{equation}

La respuesta matricial de memoria reconstruye la autoescala
hiperbólica. Para \(n>0\), el reconocimiento espectral posterior
identifica

\begin{equation}
\rho_{{\rm ep},n}=\frac{n+\sqrt{n^2+4}}2,\quad
\operatorname{Spec}F_n=\{\rho_{{\rm ep},n},-\rho_{{\rm ep},n}^{-1}\},\quad
\operatorname{Spec}H_n=\{\rho_{{\rm ep},n}^2,\rho_{{\rm ep},n}^{-2}\}.
\label{md:antmem:matriz:eq:31}
\end{equation}

El caso \(n=1\) produce el polinomio áureo.
La relación con el operador
\(F_{\rm av}=\begin{pmatrix}0&1\\1&1\end{pmatrix}\)
generado por el calendario conserva el cambio
de coordenatización:

\begin{equation}
F_1=PF_{\rm av}P^{-1},\qquad
H_1=PF_{\rm av}^{\,2}P^{-1}.
\label{md:antmem:matriz:eq:32}
\end{equation}

El transporte \(P\), de determinante uno, entrelaza ambas
matrices y conserva \(\Omega\).

\subsubsection{Estado inicial, transporte y lectura}

La proposición~\ref{md:antmem:covarianza-gaussiana} construye la realización gaussiana de la elipse de acción. En la sección positiva \(\hbar_a>0\), con \(A>0\) y \(|C^*/A|<1\), su covarianza es

\begin{equation}
V_\eta=\frac{\hbar_a}{2}M_\eta M_\eta^{\mathsf T},\qquad
M_\eta=\operatorname{diag}(e^{\eta/2},e^{-\eta/2}),\qquad
\eta=\operatorname{artanh}(C^*/A).
\label{md:antmem:cuantica:eq:5.4}
\end{equation}

En la representación de Schrödinger sobre \(L^2(\mathbb R,dQ)\), con \(\widehat Q\) multiplicación por \(Q\) y \(\widehat P=-i\hbar_a\,d/dQ\), la relación \([\widehat Q,\widehat P]=i\hbar_a I\) vale en el espacio de Schwartz. El estado puro normalizado es

\[
\psi_\eta(Q)=(\pi\hbar_a e^\eta)^{-1/4}
\exp[-Q^2/(2\hbar_a e^\eta)].
\]

La preparación es el producto \(\psi_\eta\otimes\psi_\eta\). El transporte de \eqref{md:antmem:cuantica:eq:5.1} se expresa en la misma coordenatización de covarianza por

\begin{equation}
\widetilde U=(M_\eta\oplus M_\eta)U_H
(M_\eta^{-1}\oplus M_\eta^{-1}).
\label{md:antmem:cuantica:eq:5.5}
\end{equation}

El lema~\ref{md:antmem:gauss:lem:implementacion-simplectica} proporciona una implementación unitaria de esta matriz simpléctica en dos modos. Las implementaciones que difieren por una fase global producen la misma conjugación del estado. La covarianza se transforma por \(V\mapsto\widetilde U V\widetilde U^{\mathsf T}\); el estado conjunto permanece puro y su traza parcial sobre el segundo modo determina la lectura visible, conforme a la construcción de la sección~\ref{md:antmem:gauss:gaussianas-espectro}.

\begin{theorem}[Mezcla reducida por deformación relativa]

La covarianza de la lectura visible es

\begin{equation}
\boxed{V_{\rm vis}=
\frac{\hbar_a}{2}M_\eta
\frac{8I+HH^{\mathsf T}}9
M_\eta^{\mathsf T}.}
\label{md:antmem:cuantica:eq:5.6}
\end{equation}

Su valor simpléctico normalizado es

\begin{equation}
\boxed{\nu(H):=\frac{2\sqrt{\det V_{\rm vis}}}{\hbar_a}
=\sqrt{1+\frac8{81}\left[\operatorname{tr}(HH^{\mathsf T})-2\right]}.}
\label{md:antmem:cuantica:eq:5.7}
\end{equation}
\end{theorem}

\begin{proof}
Sean \(T=(8I+H)/9\) y \(D=\sqrt8(H-I)/9\). Los dos modos iniciales son independientes y tienen la misma covarianza. En la coordenatización equilibrada del estado inicial, la covarianza visible normalizada es \(TT^{\mathsf T}+DD^{\mathsf T}\). La expansión da

\[
\begin{aligned}
TT^{\mathsf T}+DD^{\mathsf T}
&=\frac{64I+8(H+H^{\mathsf T})+HH^{\mathsf T}
+8[HH^{\mathsf T}-H-H^{\mathsf T}+I]}{81}\\
&=\frac{8I+HH^{\mathsf T}}9.
\end{aligned}
\]

Esto demuestra \eqref{md:antmem:cuantica:eq:5.6}. Puesto que \(\det M_\eta=1\) y \(\det(HH^{\mathsf T})=1\),
\(\det(8I+HH^{\mathsf T})=65+8\operatorname{tr}(HH^{\mathsf T})\).
La división por \(81\) demuestra \eqref{md:antmem:cuantica:eq:5.7}. Los dos autovalores positivos de \(HH^{\mathsf T}\) son recíprocos; su suma es al menos dos. En consecuencia, \(\nu\geq1\), con igualdad exactamente cuando \(HH^{\mathsf T}=I\).
\end{proof}

La escala \(\hbar_a\) y la deformación \(\eta\) se conservan hasta efectuar el determinante. Su cancelación es covariante porque se transportan conjuntamente estado y operador. La expresión \(HH^{\mathsf T}\) utiliza el producto euclídeo de la coordenatización equilibrada del estado inicial; un cambio de coordenatización transporta también ese producto métrico.

Un transporte \(H\) ortogonal y simpléctico puede tener \(D\neq0\) y, a la vez, conservar el producto gaussiano isotrópico, con \(\nu=1\). El complemento operatorio y el entrelazamiento del estado tienen, por tanto, criterios distintos. En esta preparación, la deformación relativa no isométrica caracteriza exactamente el caso \(\nu>1\).

\subsubsection{Evaluación exacta del lazo áureo}

Para \eqref{md:antmem:cuantica:eq:5.2},

\begin{equation}
\operatorname{tr}(H_nH_n^{\mathsf T})-2
=n^2(2n^2-2n+5),\qquad
\nu_n^2=1+\frac8{81}n^2(2n^2-2n+5).
\label{md:antmem:cuantica:eq:5.8}
\end{equation}

La identidad resulta de sumar los cuadrados de las cuatro entradas. Para \(n=1\), la suma es siete y

\begin{equation}
\boxed{\nu_1=\frac{11}{9},\quad
\operatorname{Tr}\rho_{\rm vis}^2=\frac9{11},\quad
\bar n_{\rm W}=\frac19,\quad
\frac{S_{\rm vis}}{k_B}=\frac{10}{9}\log10-\log9.}
\label{md:antmem:cuantica:eq:5.9}
\end{equation}

El teorema~\ref{md:antmem:gauss:thm:espectro-gaussiano} demuestra que el estado gaussiano centrado de un modo con valor simpléctico \(\nu\) es unitariamente equivalente a la densidad geométrica de \eqref{md:antmem:gauss:eq:estado-geometrico}, cuya ocupación media es \(\bar n_{\rm W}=(\nu-1)/2\). Para \(\nu=11/9\), sus autovalores son

\begin{equation}
\lambda_j=\frac9{10}\,10^{-j},\qquad j=0,1,2,\ldots.
\label{md:antmem:cuantica:eq:5.10}
\end{equation}

La serie suma uno. Su cuadrado suma
\((9/10)^2/(1-10^{-2})=9/11\). Finalmente,

\[
-\sum_j\lambda_j\log\lambda_j
=-\log(9/10)+\log10\sum_jj\lambda_j
=\frac{10}{9}\log10-\log9.
\]

La reducción gaussiana y su espectro están demostrados en la sección~\ref{md:antmem:gauss:gaussianas-espectro}; las dos series anteriores evalúan exactamente la pureza y la entropía de esta realización. Como el estado conjunto es puro, \(S_{\rm vis}/k_B\) es también su entropía adimensional de entrelazamiento. La transformación global es unitaria, mientras la restricción al primer modo tiene entropía positiva.

El resultado corresponde a la preparación y al transporte especificados. Su identificación con una medición del vacío requiere la realización experimental de ese protocolo. La entropía de una reducción y el calor transferido se mantienen como magnitudes distintas.

% Correspondencia editorial: gestion/mapa_transportes_auxiliares.md.
\section{Memoria operacional e invariantes de una estructura evolutiva}
\label{md:lectin:seccion}

\subsection{Valor observable, capacidad de actuación y compatibilidad}
\label{md:lectin:valor}

Conocer un objeto exige determinar qué conserva su valor observable y qué
aspectos de su capacidad de actuar requieren una descripción adicional.
En la construcción HMT, dos lecturas presentes pueden coincidir mientras
la orientación y la memoria conservadas producen respuestas distintas
ante una misma continuación. La genealogía adquiere contenido
físico-matemático cuando determina esas diferencias de respuesta.

El origen de esta investigación es la composición APP--TRIT--TPK, su
estado enriquecido y la estructura discreta conjunta. Las coordenadas
$\pi,\varphi,e,\alpha$ conservan su función de salidas correlacionadas.
Los lectores de capacidad, reciprocidad y acción intervienen después,
con sus dominios propios. Las composiciones siguientes examinan esos
lectores y mantienen el generador como antecedente.

Una caracterización conceptual de las constantes consiste en estudiarlas
como invariantes de compatibilidad de una estructura que evoluciona.
Esta caracterización exige precisar, para cada constante, las
transformaciones bajo las cuales permanece invariante, la información
que puede variar y el lector que realiza su identificación. La
reciprocidad radial proporciona un caso demostrado: la sección de
acción permanece invariante cuando se invierte el radio y se conserva
el cambio correspondiente de orientación. Una extensión de esta
invariancia a la holonomía completa requiere la composición explícita
de las operaciones que intervienen.

\subsection{Memoria conservada y memoria operativamente visible}
\label{md:lectin:memoria-visible}

En la representación conjunta de lectura y memoria, la actualización
de la lectura es $x'=Tx+Dm$. Para dos estados con una misma lectura
presente se obtiene
\begin{equation}
 \Delta x'=D\,\Delta m.
 \label{md:lectin:diferencia-memoria}
\end{equation}
La diferencia genealógica es visible en ese paso precisamente cuando
$D\Delta m\ne0$. La identidad identifica la continuación que
transforma una diferencia de memoria en una diferencia de lectura.

En sentido complementario, para un transporte estacionario $C$ de
esta representación se tiene
\begin{equation}
 D=\frac{\sqrt8}{9}(C-I),\qquad R=\frac{I+8C}{9}.
 \label{md:lectin:bloques-memoria}
\end{equation}
Si $C\Delta m=\Delta m$, entonces $D\Delta m=0$ y
$R\Delta m=\Delta m$. Esa diferencia de memoria permanece invisible
en las repeticiones del mismo protocolo. Una operación diferente puede
transferirla a la lectura.

La suficiencia de una descripción se refiere, por tanto, a las
diferencias que debe conservar para predecir las respuestas a las
operaciones admitidas. La genealogía completa puede contener más
información que la requerida por una observación particular. El
alcance del observador determina qué reducción conserva capacidad
predictiva. Esta reducción se estudia desde los transportes producidos
por HMT y mantiene la procedencia de sus estados y operadores.

En el caso estacionario, las diferencias de memoria invisibles son
exactamente $\operatorname{Fix}(C)$. Para una colección de operaciones
disponibles, las diferencias invisibles ante todas las continuaciones
forman la intersección de sus espacios fijos. La necesidad se obtiene
al tomar cada operación como primer paso. La suficiencia se demuestra
por inducción: cada operación deja fija la diferencia considerada y
su bloque de transferencia anula esa diferencia, de modo que nunca
aparece en la lectura. El resultado concierne a esta representación
lineal; el contador cronológico de la holonomía conserva su propio mapa
hacia ella.

\subsection{Regularidad de las operaciones y prolongación del estado}
\label{md:lectin:regularidad}

Las cadencias de vacancias distinguen varios niveles de cierre. Un
régimen TRIT puede retornar mientras la fase de capacidad permanece
desplazada. Una fase nonádica puede retornar mientras aumenta el
contador de vueltas. La repetición de un tipo de operación es
compatible con la producción de una historia nueva.

Esta distinción permite estudiar la regularidad de una prolongación
ilimitada: una regla finita puede ordenar la evolución sin que cada
estado nuevo reproduzca uno anterior. La finitud de la regla y la
finitud del conjunto de estados alcanzables son propiedades diferentes.
Una lectura decimal aperiódica puede proceder de una dinámica sujeta
a relaciones estrictas.

En las regiones correlacionadas, el objeto de estudio es la relación
que determina conjuntamente el siguiente bloque de cada región, las
propiedades conservadas por el espejo y la memoria que hace compatibles
las prolongaciones. Las cadencias $21/22$ y $6/7$ son observaciones
de capacidad; sus mapas concretos las relacionan con los lectores
regionales. La prueba de profundidad arbitraria y la explicación de
la estructura compartida constituyen resultados distintos y
complementarios.

La composición de las cadencias separa exactamente los niveles:
completar los tres cambios del selector deja un desplazamiento de
capacidad de tres o seis clases módulo nueve. Las duraciones $437$
y $459$ cuantifican esa distinción en la construcción correspondiente;
su función aquí es describir esas cadencias, sin promoverlas a
constantes fundamentales adicionales.

\subsection{Carácter de acción y restitución radial orientada}
\label{md:lectin:h5}

El carácter $H_5$ y la continuación regional determinan una sección
de acción. La forma radial permite restituir esa sección cuando se
conserva la orientación. A $\alpha$ fija, la coordenada
$\chi_{\mathrm{rad}}=\varepsilon\log q_\varepsilon$ clasifica la
reciprocidad y se relaciona monótonamente con la sección de retorno:
\begin{equation}
 \frac{\eta_{\mathrm{ret}}}{\alpha}
 =-500\tanh\!\left(\frac{\chi_{\mathrm{rad}}}{2}\right)-1.
 \label{md:lectin:accion-radial}
\end{equation}
La forma radial orientada y la sección de acción quedan sometidas a
una restricción conjunta. La acción incorpora una condición de
compatibilidad geométrica; sus coeficientes, su corrección regional
y su orientación deben permanecer coordinados.

Esta relación distingue dos niveles de conocimiento. La
reparametrización demuestra equivalencia de descripciones y permite
recuperar una desde la otra. Un contraste entre lectores aporta
información adicional cuando ambos se evalúan desde la estructura
anterior sin introducir como entrada el valor que se está
comprobando. Calcular el radio mediante la inversa de la fórmula de
acción y sustituirlo de nuevo verifica una identidad algebraica.
Producir la lectura geométrica mediante sus operaciones antecedentes
y compararla con el carácter de acción contrasta la composición
causal. Ambas comprobaciones tienen funciones propias.

La distinción identifica los desarrollos que deben reunirse y
ejecutarse; mantiene como raíz común la genealogía HMT. La independencia
del contraste se exige respecto del valor comprobado, no respecto
del origen compartido de sus lectores.

\subsection{Acción, información y temperatura}
\label{md:lectin:termica}

El estado de las continuaciones admisibles y sus pesos determina una
información adimensional. El transductor $k_B$ realiza su lectura en
unidades de entropía. La acción y el período intervienen en
\begin{equation}
 k_B\Theta_{\mathrm{clk}}\log3=\frac{h}{T_{108}}.
 \label{md:lectin:accion-informacion}
\end{equation}
Las tres alternativas equiprobables producen $\log3$; una distribución
no uniforme conserva sus pesos en la evaluación de la información.

La composición permite seguir cómo una misma evolución modifica la
memoria operacional, las continuaciones compatibles y la lectura
energética. La longitud de una historia y su temperatura corresponden
a lecturas diferentes. La evaluación térmica requiere identificar el
estado estadístico y la operación de intercambio que la realizan.
El problema preciso consiste en determinar qué transformaciones
generadas conservan la acción, cuáles redistribuyen la información
accesible y cuáles producen una diferencia térmica calculable.

Cuando un recorrido se cierra en una proyección mientras cambia su
memoria, la evaluación de su trabajo exige declarar el espacio en el
que se considera completo el ciclo. Éste es el problema de la
conservatividad. La no conmutatividad caracteriza, en cambio, la
sensibilidad al orden de las operaciones. Conservar la distinción
permite estudiar la relación entre ambas propiedades.

\subsection{Programa de contraste y jerarquía de las comprobaciones}
\label{md:lectin:contrastes}

\paragraph{Historias con lectura común y respuestas futuras distintas.}
El primer contraste reúne dos rutas admisibles producidas por TPK
que coinciden en un lector, conserva sus memorias y aplica la misma
continuación. La magnitud examinada es la separación exacta de sus
lecturas posteriores y el componente de memoria que la determina.
Así, la diferencia entre valor y estructura se expresa mediante una
diferencia de capacidad predictiva. Los ejemplos ya construidos
mantienen su procedencia y pueden prolongarse dentro de ese esquema.

\paragraph{Compatibilidad entre acción regional y forma geométrica.}
El segundo contraste sigue el lector de $H_5$ y su continuación, y
el lector geométrico con sus antecedentes propios. Su compatibilidad
se examina antes y después de un retorno, distinguiendo la sección
límite de sus aproximaciones finitas. El contraste identifica qué
permanece constante mientras avanza la memoria. La invariancia radial
demostrada proporciona el caso de referencia y conserva su alcance
concreto.

\paragraph{Orden de operaciones y reentrada de memoria.}
El tercer contraste emplea dos operaciones efectivas del TPK y
mantiene las coordenatizaciones intermedias. El reparto $1:8$ de la realización
auxiliar corresponde al régimen de una fibra con referencia identidad.
La aplicación debe verificar ese régimen o transportar las coordenatizaciones
correspondientes. El experimento matemático compara los dos órdenes
con y sin reentrada de memoria; la realización física debe especificar
la preparación, la interacción y el observable que representan esas
operaciones.

Los dos primeros contrastes fijan la prioridad conceptual: determinan
qué información tiene eficacia dinámica y qué invariantes sostienen
la compatibilidad de la arquitectura. El tercero proporciona una
respuesta cuantitativa directa dentro de su régimen demostrado.
Esta jerarquía conserva el papel auxiliar de la realización energética.

\subsection{Comparación científica y alcance de la interpretación}
\label{md:lectin:comparacion}

La memoria de una dinámica reducida tiene antecedentes en física
matemática. El formalismo de Mori--Zwanzig permite estudiar la
evolución y los términos de memoria de observables reducidos; un
ejemplo es el análisis de Hudson y Li.\footnote{Hudson y Li,
\emph{Coarse-graining of overdamped Langevin dynamics via the
Mori--Zwanzig formalism},
\href{https://arxiv.org/abs/1810.08175}{arXiv:1810.08175}.}
La comparación pertinente alcanza la composición concreta examinada
aquí: origen APP--TRIT--TPK, regiones correlacionadas, acción,
orientación y restitución. La mención de memoria por sí sola no
determina la prioridad de esa composición.

La aportación expositiva consiste en precisar consecuencias y
contrastes dentro de la construcción. Una evaluación histórica de
la novedad y una corroboración física conservan sus pruebas
específicas. El programa matemático establece objetos y operaciones
que pueden someterse a comparación científica sin anticipar el
resultado de aquellas evaluaciones.

La memoria almacenada y la memoria operativamente visible quedan
caracterizadas aquí mediante el núcleo de $D$ y el espacio fijo de
$C$. Esta caracterización es consecuencia de la actualización
declarada y conserva el alcance de su representación. Los tres
contrastes de la sección~\ref{md:lectin:contrastes} constituyen un
programa de trabajo; su formulación distingue la definición del
experimento matemático, la ejecución de rutas TPK y la realización
de un experimento físico.

% Fuente conservada íntegra: originales/cuaderno/NOTA.md, §§1--7.
\section{Memoria cronológica, balance y respuesta de vacancias}
\label{md:cron:seccion}

La construcción siguiente recibe el factor simbólico de capacidad y estudia
sus lectores posteriores. La razón interna $3^6/10^3=729/1000$ y su
calendario pertenecen a la construcción precedente
(secciones~\ref{mono-ampl:relojes} y~\ref{vac-capacidad-trit}).
La arquitectura APP--TRIT--TPK, las regiones, el retorno nonádico y el
registro $K$ conservan su posición causal anterior. Los logaritmos empleados
en el lector de capacidad describen dicho factor y no seleccionan de nuevo
las regiones de $\pi,\varphi,e,\alpha$.
Las proposiciones son consecuencias focales sobre observación y memoria;
su alcance se distingue del estado TPK completo y de una atribución de
prioridad histórica a resultados generales sobre sucesiones mecánicas o
información.

\subsection{Objeto, medida y dos lectores}
\label{md:cron:lectores}

Sea $\delta=\log_{10}(10/9)$ la pendiente del calendario de capacidad,
y sea $\theta$ una fase uniformemente distribuida en $[0,1)$.
La medida representa incertidumbre sobre el origen de lectura, mientras el
mecanismo permanece determinado. Definimos
\begin{equation}
 \begin{aligned}
 z_j(\theta)&=\lfloor(j+1)\delta+\theta\rfloor
             -\lfloor j\delta+\theta\rfloor,\\
 W_n(\theta)&=(z_0(\theta),\ldots,z_{n-1}(\theta)),\\
 B_n(\theta)&=\sum_{j=0}^{n-1}z_j(\theta).
 \end{aligned}
 \label{md:cron:dos-lectores}
\end{equation}
El bloque $W_n$ conserva la cronología y $B_n$ conserva el balance.
Para una fase exacta y una regla exacta conocidas, ambos objetos son
deterministas; las entropías se refieren expresamente al conjunto de fases
posibles.

La telescopía da
\begin{equation}
 B_n(\theta)=\lfloor n\delta+\theta\rfloor.
 \label{md:cron:telescopia}
\end{equation}
Por tanto, $B_n$ admite sólo los valores $\lfloor n\delta\rfloor$ y
$\lceil n\delta\rceil$. Los puntos $\{-j\delta\}$, con $j=0,\ldots,n$,
dividen el círculo en $n+1$ intervalos de fase. Cada intervalo porta una
secuencia distinta de longitud $n$. En efecto, los prefijos acumulados
$\lfloor k\delta+\theta\rfloor$, para $k=1,\ldots,n$, son monótonos en
$\theta$, y cada frontera cambia uno de ellos; la secuencia reconstruye
todos esos prefijos.

Los extremos tienen medida cero. Una convención semiabierta resuelve su
asignación y conserva las entropías.

\subsection{Información creciente y tasa entrópica nula}
\label{md:cron:crecimiento}

Sean $\ell_0,\ldots,\ell_n$ las longitudes de los intervalos de fase y
\begin{equation}
 H_n=-\sum_{j=0}^{n}\ell_j\log_2\ell_j.
 \label{md:cron:entropia-bloque}
\end{equation}

\begin{proposition}
\label{md:cron:proposicion-tasa}
Se cumplen simultáneamente
\begin{equation}
 H_n\longrightarrow+\infty,\qquad
 \frac{H_n}{n}\longrightarrow0.
 \label{md:cron:dos-limites}
\end{equation}
\end{proposition}
\begin{proof}
La cota superior es $H_n\leq\log_2(n+1)$. Si
$\ell_{\max}(n)=\max_j\ell_j$, entonces
$H_n\geq-\log_2\ell_{\max}(n)$.
La pendiente $\delta$ es irracional: si $\delta=p/q$ con $q>0$,
se tendría $(10/9)^q=10^p$, lo cual contradice la valuación del primo $3$.
Por irracionalidad, los puntos de la órbita son densos. Las particiones
sucesivas tienen malla máxima tendente a cero: para cada $\epsilon>0$
existe un tramo finito $\epsilon$-denso, y los tramos posteriores sólo
subdividen sus intervalos. Por tanto, la cota inferior diverge. La cota
superior dividida por $n$ tiende a cero.
\end{proof}

Las longitudes de los intervalos no son iguales en general. La prueba
establece las cotas anteriores, sin afirmar
$H_n=\log_2(n+1)$ ni una ley asintótica logarítmica exacta que requiriese
condiciones diofánticas adicionales.
La entropía topológica del factor también es cero porque
$\log(n+1)/n\to0$. La información del bloque y su tasa son cantidades
diferentes. El crecimiento del archivo de una trayectoria conocida no
equivale, por sí solo, a producción de información estadística independiente.

\subsection{Información cronológica no determinada por el balance}
\label{md:cron:condicional}

\begin{corollary}
\label{md:cron:corolario-condicional}
Para el mismo conjunto de fases,
\begin{equation}
 H(W_n\mid B_n)=H_n-H(B_n)
 \geq H_n-1\longrightarrow+\infty.
 \label{md:cron:entropia-condicional}
\end{equation}
\end{corollary}
\begin{proof}
$B_n$ es función de $W_n$, de modo que la regla de la cadena da la igualdad.
Su soporte tiene dos elementos, por lo que $H(B_n)\leq1$ bit.
Se aplica la proposición~\ref{md:cron:proposicion-tasa}.
\end{proof}

La entropía condicional que diverge es la media sobre los balances;
no se atribuye la misma incertidumbre a cada balance concreto.
Tampoco se atribuye esa pérdida al estado enriquecido completo: cuando la
memoria retenida recupera $W_n$, se tiene
\begin{equation}
 H(W_n\mid B_n,M_n)=0.
 \label{md:cron:restitucion-memoria}
\end{equation}
Por cardinalidad, al menos una de las dos fibras de $B_n$ contiene
$\lceil(n+1)/2\rceil$ secuencias o más. Esta cota combinatoria se distingue
de la cota entrópica.
El resultado precisa la relación entre valor y genealogía en este lector:
los dos balances posibles no agotan las historias temporales compatibles.

\subsection{Corriente centrada: balance acotado y actividad persistente}
\label{md:cron:corriente}

El lector centrado se expresa por
\begin{equation}
 J_j=z_j-\delta,\qquad
 \sum_{j=0}^{n-1}J_j=B_n-n\delta.
 \label{md:cron:corriente-definicion}
\end{equation}
En consecuencia,
$\left|\sum_{j=0}^{n-1}J_j\right|<1$ para toda ventana.
A la vez,
\begin{equation}
 \lim_{n\to\infty}\frac1n\sum_{j=0}^{n-1}J_j=0,\qquad
 \lim_{n\to\infty}\frac1n\sum_{j=0}^{n-1}J_j^2
 =\delta(1-\delta)>0.
 \label{md:cron:actividad}
\end{equation}
\begin{proof}
Se tiene $z_j^2=z_j$ y $B_n/n\to\delta$ por la cota de discrepancia.
La expansión
$(z_j-\delta)^2=z_j(1-2\delta)+\delta^2$, seguida de la sustitución del
promedio de $z_j$, da la segunda identidad.
El argumento vale para cualquier fase y no requiere un promedio aleatorio.
\end{proof}

El balance acumulado acotado, la actividad cuadrática positiva y la tasa de
entropía cero son, por tanto, compatibles. Las identidades conservan la
distinción entre el lector centrado y una representación de carga
elemental; tampoco establecen una producción energética sin balance.

Si una realización dimensional ya establecida toma
$I_j=(u_Q/t_0)J_j$, se obtiene
\begin{equation}
 I_{\mathrm{rms}}^2
 =\left(\frac{u_Q}{t_0}\right)^2\delta(1-\delta).
 \label{md:cron:eficaz}
\end{equation}
Si cada $I_j$ se mantiene durante un intervalo de duración $t_0$, este
promedio discreto coincide con el promedio temporal.
Como contraste físico posterior, una carga resistiva ideal
$R_{\mathrm{res}}>0$ tendría potencia media
\begin{equation}
 \overline P_{\mathrm{res}}
 =R_{\mathrm{res}}\left(\frac{u_Q}{t_0}\right)^2
 \delta(1-\delta)>0.
 \label{md:cron:potencia-resistiva}
\end{equation}
Es una predicción del modelo compuesto con esa carga. La resistencia y el
dispositivo conservan su especificación propia; no se derivan de la sola
sucesión de capacidad.

\subsection{Espectro de la corriente y orientación}
\label{md:cron:espectro}

La función de observación es
$f(\theta)=\mathbf1_{[1-\delta,1)}(\theta)-\delta$, con
$J_j=f(\theta+j\delta)$. Sus coeficientes de Fourier, para $m\neq0$, son
\begin{equation}
 a_m=\frac{\exp(2\pi i m\delta)-1}{2\pi i m},\qquad
 |a_m|^2=\frac{\sin^2(\pi m\delta)}{\pi^2m^2},\qquad
 a_0=0.
 \label{md:cron:fourier}
\end{equation}
Se obtienen integrando la indicatriz sobre su intervalo. La identidad de
Parseval da
\begin{equation}
 \sum_{m\neq0}|a_m|^2=\delta(1-\delta),
 \label{md:cron:parseval}
\end{equation}
en concordancia con la actividad cuadrática. Las frecuencias temporales
son $m\delta$ módulo $1$. La correlación estacionaria, con fase uniforme,
queda determinada por esas intensidades; éstas no contienen la fase inicial.
Una traslación multiplica $a_m$ por una fase unitaria, y la inversión
temporal intercambia las frecuencias opuestas.

La potencia espectral aislada no recupera la orientación ni el origen
temporal. Una lectura sensible a fase o una referencia marcada aporta
la información complementaria. La deducción conserva la distinción entre
esa referencia general y el registro $K$, así como entre ritmos simbólicos
y una realización experimental determinada.

\subsubsection{Referencia cíclica y transportes de igual intensidad espectral}
\label{md:cron:k-ciclico}

La propiedad cíclica del registro $K$ permite una consecuencia finita
independiente del espectro anterior. En $V_K=\mathbb R^{12}$, sea $S$ el
desplazamiento $(Sx)_i=x_{i+1}$ y
\begin{equation}
 O_K=(K,SK,\ldots,S^{11}K).
 \label{md:cron:matriz-ciclica}
\end{equation}
Se utiliza la invertibilidad de $O_K$ establecida por la referencia cíclica
del registro. Para cualquier operador real $H$ que conmute con $S$,
pongamos $y_+=HK$ e $y_-=H^TK$.

\begin{proposition}
\label{md:cron:proposicion-intensidades}
Las doce intensidades de Fourier de $y_+$ e $y_-$ coinciden.
Sin embargo, $y_+=y_-$ si y sólo si $H=H^T$.
Cada respuesta vectorial completa, junto con $K$ y la orientación de $S$,
recupera su operador.
\end{proposition}
\begin{proof}
$H$ es circulante. Sus multiplicadores de Fourier son $\lambda_m$ y los
de $H^T$ son sus conjugados, de modo que las dos respuestas tienen
intensidades $|\lambda_m|^2|\widehat K_m|^2$.
Si sus vectores coinciden, $(H-H^T)K=0$. La conmutación implica que
$H-H^T$ anula también $S^jK$ para todo $j$. Esos vectores forman una base;
por tanto, $H-H^T=0$. La implicación inversa es inmediata.
La recuperación se expresa por
\begin{equation}
 H=O_{y_+}O_K^{-1},\qquad
 H^T=O_{y_-}O_K^{-1}.
 \label{md:cron:restitucion-operador}
\end{equation}
\end{proof}

En particular, los desplazamientos $S$ y $S^{-1}$ producen respuestas con
el mismo espectro de potencia, pero distintas. Las autocorrelaciones exactas
del registro dan
\begin{equation}
 \|SK-S^{-1}K\|^2
 =2\bigl(\|K\|^2-\langle K,S^2K\rangle\bigr)
 =2175744>0.
 \label{md:cron:separacion-desplazamientos}
\end{equation}
La inversión racional de $O_K$ recupera de ambas respuestas los
coeficientes $e_1$ y $e_{11}$, respectivamente. $K$ se utiliza después
de su generación como referencia marcada. La intensidad espectral o el
balance escalar no sustituyen la respuesta vectorial.
El corolario describe el portador cíclico de doce componentes;
conserva la diferencia entre $S$ y la dinámica enriquecida completa,
así como los operadores adicionales de una conjugación física de carga.
Para un $H$ general, la transposición tampoco significa inversión temporal;
en el caso $S$, sí coincide con $S^{-1}$.

\subsection{Enlace operatorio electrónico}
\label{md:cron:electron}

La incidencia entre las fibras de suma y producto de APP sobre $D_9^3$
define $C$. Su vector trítico
$m_{\mathrm{ph}}=(1,-1,0)^3$ satisface
\begin{equation}
 Cm_{\mathrm{ph}}=C^Tm_{\mathrm{ph}}
 =-\frac12m_{\mathrm{ph}}.
 \label{md:cron:modo-tritico}
\end{equation}
El calendario inducido recorre los regímenes del mismo selector trítico;
la incidencia selecciona el plano principal de singularidad $1/2$.
En él,
\begin{equation}
 P=\begin{pmatrix}1&0\\0&0\end{pmatrix},\qquad
 Q=\begin{pmatrix}1/4&\sqrt3/4\\\sqrt3/4&3/4\end{pmatrix},
 \qquad J=\frac4{\sqrt3}[P,Q],\qquad J^2=-I.
 \label{md:cron:plano-electronico}
\end{equation}
La incidencia y la norma del conmutador se desarrollan en la
sección~\ref{vi:dep:prefactor}. Se recupera así una relación efectiva entre
organización temporal y bloque electrónico, conservando el selector,
sus etiquetas y la incidencia. La representación de carga y la realización
dimensional mantienen sus operadores propios: el enlace no asigna
directamente las cargas elementales opuestas a $z_j-\delta$.

\subsection{Información retenida y composición energética}
\label{md:cron:informacion-energia}

La recuperación de fibra establece
\begin{equation}
 H(X\mid q(X),M(X))=0
 \label{md:cron:inversa-enriquecida}
\end{equation}
cuando el lector enriquecido tiene inversa; una actualización biyectiva
conserva $H(X)$. Estas propiedades se reúnen en
\eqref{eq:ii-kb-aut-fibra} y~\eqref{eq:ii-kb-aut-reversible}.
La coexistencia de $H_n$ creciente y tasa cero es compatible con esas
propiedades y con las identidades de Shannon.

El trabajo de borrado depende de la información accesible retenida.
En el régimen ideal isotermo de memorias lógicas degeneradas y en el
límite asintótico apropiado, el coste por copia con información lateral es
\begin{equation}
 W_{\mathrm{borrado,ideal}}
 =k_BT\ln(2)\,H(X\mid M).
 \label{md:cron:borrado}
\end{equation}
Si $X$ se recupera desde $M$, la contribución lógica ideal puede ser cero.
El transporte físico, la lectura, el control y el ciclo completo conservan
balances propios; la igualdad $h_{\mathrm{top}}=0$ no los elimina.

La composición relevante conserva la secuencia
\[
 \begin{gathered}
 \text{cronología}\longrightarrow
 \text{registro de frontera}\longrightarrow
 \text{lector de corriente y estado}\\
 \longrightarrow
 \text{memoria retenida}\longrightarrow
 \text{balance térmico}.
 \end{gathered}
\]
El desarrollo anterior construye los dos primeros lectores y sus
consecuencias informacionales. Su alcance se distingue de la certificación
de un dispositivo físico completo y mantiene intacta la genealogía
constructiva de las constantes.

\section{Reiteración visible y diferenciación genealógica en una región}
\label{md:ejemplo:region}
Las biografías de \eqref{act21:eq:biografias-tpk-tres-caracteres}
permiten precisar la diferencia entre una reiteración de cifras y un
retorno del estado. Para el canal que se reconoce como \(e\), el
calendario \eqref{act21:eq:calendario-0011-interno} produce
\[
 201101\mid121221\mid102011\mid012222\mid102011.
\]
La elevación precede a la lectura decimal. Las mismas matrices actúan
en los otros dos canales, conservando sus emisiones iniciales y sus
estados regionales distintos. Esta estructura común permite comparar
las biografías sin identificar sus coordenadas.

Sea \(w\) una concatenación de \(L\) trits, y \(N\) su valor entero
en base tres. El cilindro correspondiente es
\[
 I_w=[N/3^L,(N+1)/3^L).
\]
Para \(d\) posiciones decimales, el prefijo es estable en todo el
cilindro exactamente cuando
\[
 \left\lfloor\frac{10^dN}{3^L}\right\rfloor
 =\left\lfloor\frac{10^d(N+1)-1}{3^L}\right\rfloor.
\]
Es la aplicación del lema~\ref{act21:censo:estabilidad} al lector
decimal. El extremo derecho semiabierto explica el término \(-1\)
de la división entera. Añadiendo la parte entera de la coordenatización regional,
las cinco lecturas sucesivas son
\begin{center}
\begin{tabular}{@{}rl@{}}
\toprule
Longitud trítica & Prefijo estable de la lectura regional\\
\midrule
6 & \(2.71\ldots\)\\
12 & \(2.71828\ldots\)\\
18 & \(2.7182818\ldots\)\\
24 & \(2.7182818284\ldots\)\\
30 & \(2.71828182845904\ldots\)\\
\bottomrule
\end{tabular}
\end{center}
En longitud veinticuatro se tiene
\(N=202864003874\) y \(3^{24}=282429536481\).
En longitud treinta se obtiene
\(N=147887858824447\) y \(3^{30}=205891132094649\).
Las divisiones enteras anteriores certifican los prefijos de la tabla.
En particular, las primeras doce cifras fraccionarias se agrupan como
\[
 718281828459
 =7\,\underbrace{1828\,1828}_{\text{reiteración posicional}}\,459.
\]
El cilindro de veinticuatro trits ya fija ambas copias y la cifra
siguiente. Esta localización distingue el episodio decimal de la
reaparición del bloque trítico \(102011\), situado en las posiciones
tercera y quinta: sus cortes y sus bases de lectura son diferentes.

La repetición del bloque visible tampoco identifica sus preestados.
Los preestados módulo veintisiete conservados en el desarrollo de
esta biografía son
\[
 (25,11,7,17,8,16),\qquad (7,8,4,23,26,7),
\]
con cocientes posteriores respectivos
\[
 (6,13,9,2,13,1),\qquad (15,0,3,19,23,25).
\]
La igualdad de la emisión coexiste así con una genealogía diferenciada.
Su función explicativa es concreta: la observación de un bloque
reiterado conserva menos información que la continuación enriquecida
que lo produce. Para atribuir una diferencia futura exclusivamente
a la memoria hay que conservar asimismo las fases y las coordenatizaciones de
comparación.

Este episodio relaciona tres niveles: las transiciones producen la
cadena trítica, los cilindros determinan las cifras estables y el
estado conserva la procedencia de una emisión reiterada. La
profundidad arbitraria pertenece a la prolongación demostrada
anteriormente; esta lectura finita precisa un episodio de esa
biografía. La regularidad que interesa estudiar reside en la
compatibilidad entre los niveles, no en prolongar indefinidamente
un fragmento decimal: una tercera copia de \(1828\) comenzaría por
\(1\), mientras la continuación ya producida comienza por \(4\).

% Fuente íntegra: originales/cuaderno/REORIENTACION_GENEALOGICA_Y_RESULTADOS.md,
% §§2--7. Gestión y procedencia conservadas aparte.
\Needspace{16\baselineskip}
\section{Regularidad genealógica, restitución operatoria y geometría de la acción}
\label{md:restop:seccion}

\subsection{Producción, compatibilidad y restitución}
\label{md:restop:jerarquia}

El fenómeno primario es una dinámica que conserva procedencia al producir
nuevas observaciones. Se reúnen tres operaciones con funciones matemáticas
diferenciadas:
\begin{enumerate}
 \item \emph{Producción.} La actualización APP--TRIT--TPK prolonga el estado
 y sus regiones correlacionadas. La fase puede retornar mientras la memoria
 avanza.
 \item \emph{Compatibilidad.} Las lecturas regionales, angulares, de acción
 y de incidencia obedecen relaciones comunes. Su variación debe respetar
 esas relaciones para corresponder a un mismo estado.
 \item \emph{Restitución.} Conservar la información complementaria permite
 recuperar valores anteriores, productos de operadores y diferencias
 entre órdenes de composición.
\end{enumerate}
La comparación entre procedimientos para obtener $\pi$ debe atender a esos
tres niveles. La tesis específica de HMT reside en la construcción conjunta
y sus relaciones, además de la producción de expansiones numéricas.
Una coincidencia decimal o la periodicidad de una fase no describen por sí
solas el alcance de la construcción. La distinción identifica el objeto
matemático pertinente para una comparación histórica; el desarrollo focal
no establece, por sí mismo, una prioridad mundial.

\subsection{Regímenes de vacancia y retorno de la fase de capacidad}
\label{md:restop:vacancias}

\subsubsection{Calendario de capacidad y selector trítico}
\label{md:restop:antecedentes-capacidad}

El lector de capacidad recibe la compactificación interna $729/1000$ y
distingue el índice de refinamiento del índice de capacidad. Con la
convención de la sección~\ref{vac-capacidad-trit},
\begin{equation}
 \delta=\log_{10}(10/9),\qquad
 p_j=\left\lfloor\frac{j}{\delta}\right\rfloor,\qquad
 v_j=p_j-j,\qquad j\geq1.
 \label{md:restop:capacidad}
\end{equation}
Los intervalos entre eventos son
\begin{equation}
 h_j=p_{j+1}-p_j\in\{21,22\},\qquad
 s_j=22-h_j.
 \label{md:restop:intervalos}
\end{equation}
La clase trítica obedece
\begin{equation}
 [v_{j+1}]_3=[v_j-s_j]_3.
 \label{md:restop:clase-tritica}
\end{equation}
Un intervalo corto cambia la clase; uno largo la conserva.
Las separaciones $6/7$ describen las mesetas de este régimen inducido.
La construcción mantiene su tipo de factor de observación del estado TPK;
ni sustituye dicho estado ni identifica esas separaciones con signos de
carga o con cifras de $\pi$ o de $e$.
El selector canónico es
$M_{\mathrm{ph}}(1+[v_j]_9)$, cuya pauta depende de la clase módulo tres.

\subsubsection{Restitución del régimen local y desplazamiento módulo nueve}
\label{md:restop:consecuencia-capacidad}

Sea
\begin{equation}
 \beta_{\mathrm{cap}}=22-\delta^{-1},\qquad
 q_k=\left\lfloor\frac{k}{\beta_{\mathrm{cap}}}\right\rfloor,\qquad k\geq1.
 \label{md:restop:pendiente-inducida}
\end{equation}
Las posiciones $q_k$ son los índices de los intervalos cortos. En efecto,
\begin{equation}
 p_j=22j-\lceil j\beta_{\mathrm{cap}}\rceil,\qquad
 s_j=\lceil(j+1)\beta_{\mathrm{cap}}\rceil
     -\lceil j\beta_{\mathrm{cap}}\rceil.
 \label{md:restop:eventos-cortos}
\end{equation}
La irracionalidad de $\delta$ se obtiene de las valuaciones de $2,3,5$
en la igualdad que impondría $(10/9)^q=10^p$; implica la de
$\beta_{\mathrm{cap}}$. Por ello, los índices positivos no presentan
ambigüedad de extremos. La discontinuidad correspondiente al entero $k$
ocurre en $j=\lfloor k/\beta_{\mathrm{cap}}\rfloor$.

\begin{proposition}
\label{md:restop:proposicion-regimen}
Entre los eventos de índices $q_k$ y $q_{k+3}$, tomando los intervalos con
índices en $[q_k,q_{k+3})$, retorna la clase TRIT, mientras cambia la fase
de capacidad módulo nueve.
\end{proposition}
\begin{proof}
Hay exactamente tres intervalos cortos en ese tramo.
Las cotas $1/7<\beta_{\mathrm{cap}}<3/20$ dan
$20<3/\beta_{\mathrm{cap}}<21$, por lo que
\begin{equation}
 m_k=q_{k+3}-q_k\in\{20,21\}.
 \label{md:restop:duracion-tritica}
\end{equation}
La suma de longitudes y la ganancia de capacidad son, respectivamente,
\begin{equation}
 \Delta p=22m_k-3\in\{437,459\},\qquad
 \Delta v=21m_k-3\in\{417,438\}.
 \label{md:restop:ganancias}
\end{equation}
En ambos casos $\Delta v\equiv0\pmod3$. En cambio,
\begin{equation}
 \boxed{\Delta v\equiv3\ \text{o}\ 6\pmod9.}
 \label{md:restop:desplazamiento-nonadico}
\end{equation}
Queda probado el enunciado.
Las cotas de $\beta_{\mathrm{cap}}$ disponen de certificados enteros:
\begin{equation}
 \begin{aligned}
 10^{146}>9^{153}&\ \Longrightarrow\
 \delta>\frac7{153},\\
 10^{417}<9^{437}&\ \Longrightarrow\
 \delta<\frac{20}{437}.
 \end{aligned}
 \label{md:restop:certificados-cotas}
\end{equation}
La monotonía de $22-1/x$ produce
$1/7<\beta_{\mathrm{cap}}<3/20$. Así, las cotas se apoyan en
desigualdades enteras y no en una sustitución decimal redondeada.
\end{proof}

El valor del selector
$\tau_j=M_{\mathrm{ph}}(1+[v_j]_9)$ retorna porque depende de la clase
módulo tres. Este retorno conserva su distinción con el estado TRIT
completo y con una vuelta cronológica de $\Gamma_9$.
Ambas duraciones aparecen infinitas veces por la densidad de la rotación
irracional de pendiente $1/\beta_{\mathrm{cap}}$.

La regularidad tiene niveles distintos: se restaura el régimen local
mientras sigue desplazada la fase de capacidad. El retorno nonádico del
estado tiene, además, su incremento propio de memoria. Cada objeto
conserva, por tanto, su noción de retorno y su duración.
Los números $437$ y $459$ son duraciones de esta observación inducida,
sin atribuirles el estatuto de constantes universales adicionales.
El resultado reúne explícitamente las fórmulas previas; su presentación
como consecuencia focal no constituye una afirmación de prioridad
histórica frente a desarrollos no comparados exhaustivamente.

\subsection{Participación de la memoria en la composición operatoria}
\label{md:restop:memoria}

\subsubsection{Representación ampliada de lectura y memoria}
\label{md:restop:estructura-partida}

La representación de lectura y memoria se obtiene separando ocho
componentes y una novena. En una misma fibra, con transporte de referencia
identidad, el transporte interno $C$ induce
\begin{equation}
 T_C=\frac{8I+C}{9},\qquad
 D_C=\frac{\sqrt8(C-I)}{9},\qquad
 R_C=\frac{I+8C}{9},
 \label{md:restop:bloques}
\end{equation}
\begin{equation}
 \mathcal U_C=
 \begin{pmatrix}
 T_C&D_C\\D_C&R_C
 \end{pmatrix}.
 \label{md:restop:transporte-ampliado}
\end{equation}
El acoplamiento procede del cambio de base de la descomposición nonádica.
Para transportes que conservan la forma declarada, conserva la forma
ampliada. La identidad de composición es
\begin{equation}
 \mathcal U_{C_2}\mathcal U_{C_1}
 =\mathcal U_{C_2C_1},\qquad
 T_2T_1+D_2D_1=\frac{8I+C_2C_1}{9},
 \label{md:restop:composicion-ampliada}
\end{equation}
con $T_i=T_{C_i}$ y $D_i=D_{C_i}$.
El segundo término corresponde a la memoria producida en la primera
operación y reintroducida en la segunda. Apartarla cambia el
procedimiento: produce $T_2T_1$ en lugar de la lectura de la composición
ampliada.

\subsubsection{Reparto de la sensibilidad al orden de composición}
\label{md:restop:conmutadores}

\begin{proposition}
\label{md:restop:proposicion-orden}
Con la convención $[A,B]=AB-BA$,
\begin{equation}
 [T_2,T_1]=\frac1{81}[C_2,C_1],\qquad
 [D_2,D_1]=\frac8{81}[C_2,C_1].
 \label{md:restop:reparto-conmutadores}
\end{equation}
En consecuencia, la diferencia de lectura entre ambos órdenes, con la
memoria reintroducida, es
\begin{equation}
 \boxed{(T_2T_1+D_2D_1)-(T_1T_2+D_1D_2)
 =\frac19[C_2,C_1].}
 \label{md:restop:orden-completo}
\end{equation}
\end{proposition}
\begin{proof}
Al expandir los productos de $T$, los términos constantes y lineales se
cancelan en la resta; queda el conmutador dividido por $81$.
Al expandir los productos de $D$ se produce la misma cancelación y el
factor es $8/81$. Su suma es $1/9$.
\end{proof}

El mecanismo de reentrada forma parte de la representación ampliada.
La comparación de los dos órdenes hace explícita una consecuencia:
en esta representación, la contribución de memoria a la respuesta
antisimétrica es ocho veces la retenida por la composición de lecturas
aisladas. La respuesta completa es nueve veces esta última. Las
proporciones corresponden a operadores de respuesta, no a porcentajes
de energía.

La restitución conserva la no conmutatividad y su intensidad de lectura.
En este caso, la compresión aislada atenúa un conmutador no nulo sin
anularlo. Una aplicación a fuerzas físicas conserva además los operadores
de interacción y el funcional de trabajo correspondiente; la identidad
algebraica no los identifica por sí sola.

\subsubsection{Contraste operatorio y normalización nonádica}
\label{md:restop:contraste}

Para el patrón algebraico general $(n-1)+1$, las mismas cuentas dan los
factores $1/n^2$, $(n-1)/n^2$ y $1/n$. La arquitectura nonádica fija
$n=9$. Se distinguen así la ley de composición y el valor concreto de
su normalización.

Cuando las operaciones efectivas de una aplicación se representan por
$C_1,C_2$, los procedimientos que apartan o reintroducen la memoria
predicen respuestas de orden distintas. El contraste recae en la
composición operacional y en la accesibilidad de la memoria, conservando
los coeficientes de la representación.

\subsubsection{Normas de lectura y memoria en la inversión central}
\label{md:restop:treinta-y-dos}

Para $C=-I$, se obtiene
\begin{equation}
 T^*T=\frac{49}{81}I,\qquad
 D^*D=\frac{32}{81}I.
 \label{md:restop:normas}
\end{equation}
El numerador $32$ procede de $8|{-1}-1|^2$.
En las nueve componentes de la memoria, la norma incluye
$8^2+8=72$, con denominador $27^2$. El cuadrado $64$ participa en
ese cálculo junto con las otras ocho contribuciones.
Una identificación con una región de $64$ casillas exige el mapa de
incidencia de la región y no se obtiene duplicando el numerador $32$.
La lectura geométrica de un soporte $8\times8$ frente a otro
$9\times9$ conserva, por tanto, su problema de conexión con la
aplicación concreta; la igualdad numérica aislada no lo sustituye.

\clearpage
\subsection{Carácter de acción y coordenada completa de reciprocidad}
\label{md:restop:h5}

\subsubsection{Coeficientes genealógicos y continuación regional}
\label{md:restop:coeficientes}

La matriz central y sus invariantes son
\begin{equation}
 B_c=\begin{pmatrix}7&2\\2&7\end{pmatrix},\qquad
 \lambda_+=9,\qquad\lambda_-=5.
 \label{md:restop:matriz-central}
\end{equation}
Intervienen además la órbita de longitud cuatro del paso tres, el
calendario de longitud doce y el cociente censal $104976/1944=54$.
Las relaciones de coeficientes son
\begin{equation}
 \begin{aligned}
 \frac9{16}&=\frac{\lambda_+}{4^2},&
 \frac59&=\frac{\lambda_-}{\lambda_+},\\
 \frac7{48}&=\frac{\operatorname{tr}B_c/2}{4\cdot12},&
 \frac1{54}&=\frac{1944}{104976}.
 \end{aligned}
 \label{md:restop:coeficientes-h5}
\end{equation}
La asignación de grados y signos pertenece a la definición del carácter;
sus cardinales aislados no la determinan. El lector completo es
\begin{equation}
 H_5=\frac12\sqrt{\frac{1000\alpha}{\varphi}}-\alpha
 +\frac9{16}\alpha^2-\frac59\alpha^3
 +\frac7{48}\alpha^4-\frac1{54}\alpha^5,
 \label{md:restop:h5-completo}
\end{equation}
\begin{equation}
 \begin{aligned}
 D_A&=\exp(-100\pi\alpha/9),\\
 \mathcal C_\pi&=169\alpha^6+
 \frac{D_A\alpha^7}{1-D_A\alpha},\\
 \eta_{\mathrm{ret}}&=H_5-\frac{90}{\pi}\mathcal C_\pi.
 \end{aligned}
 \label{md:restop:retorno-regional}
\end{equation}
La fracción conserva la continuación infinita definida por
\begin{equation}
 \mathscr R(z)=D_Az^7+D_Az\mathscr R(z),
 \label{md:restop:renovacion}
\end{equation}
con peso inicial unitario. La continuación permanece íntegra en el
lector; un polinomio truncado no sustituye ese objeto. La construcción
de los coeficientes y la renovación se reúnen en la
sección~\ref{sec:revision-planck}.

\subsubsection{Clasificación de la reciprocidad radial orientada}
\label{md:restop:geometria}

Se conservan $A_{\deg}=1000\alpha$,
$C^*_{\deg}=2(\eta_{\mathrm{ret}}+\alpha)$ y el recuperador radial del
teorema~\ref{rec:thm:accion}. Para mantener ambas orientaciones, sean
\begin{equation}
 \epsilon\in\{+1,-1\},\qquad
 q_{\mathrm{rad}}=R_\star^4>0,\qquad
 \Xi_{\mathrm{rad}}=
 \epsilon\frac{1-q_{\mathrm{rad}}}{1+q_{\mathrm{rad}}}.
 \label{md:restop:par-radial}
\end{equation}
El recuperador \eqref{rec:eq:eta-orientada} se escribe
\begin{equation}
 \frac{\eta_{\mathrm{ret}}}{\alpha}=500\Xi_{\mathrm{rad}}-1.
 \label{md:restop:recuperador}
\end{equation}
La involución
$(\epsilon,q_{\mathrm{rad}})
 \mapsto(-\epsilon,q_{\mathrm{rad}}^{-1})$
conserva esa acción. La coordenada que reúne esta estructura es
\begin{equation}
 \boxed{\chi_{\mathrm{rad}}=\epsilon\log q_{\mathrm{rad}}.}
 \label{md:restop:chi}
\end{equation}

\begin{proposition}
\label{md:restop:proposicion-orbitas}
$\chi_{\mathrm{rad}}$ clasifica exactamente las órbitas de esa involución
sobre $\{\pm1\}\times\mathbb R_{>0}$.
\end{proposition}
\begin{proof}
Cambiar simultáneamente el signo e invertir $q_{\mathrm{rad}}$ conserva
$\chi_{\mathrm{rad}}$. Si dos pares tienen el mismo valor de
$\chi_{\mathrm{rad}}$, con igual signo tienen el mismo
$q_{\mathrm{rad}}$; con signos opuestos, sus radios cuárticos son
recíprocos. Son exactamente los dos casos de pertenencia a una misma
órbita. El argumento incluye $q_{\mathrm{rad}}=1$, cuya órbita conserva
las dos orientaciones.
\end{proof}

La identidad
$(1-q)/(1+q)=-\tanh(\log q/2)$ da
\begin{equation}
 \Xi_{\mathrm{rad}}
 =-\tanh(\chi_{\mathrm{rad}}/2),\qquad
 \boxed{\eta_{\mathrm{ret}}
 =\alpha[-500\tanh(\chi_{\mathrm{rad}}/2)-1].}
 \label{md:restop:accion-chi}
\end{equation}
La compatibilidad con el carácter completo adopta la forma
\begin{equation}
 \boxed{\chi_{\mathrm{rad}}
 =-2\operatorname{artanh}
 \left(\frac{H_5+\alpha-(90/\pi)\mathcal C_\pi}{500\alpha}\right).}
 \label{md:restop:chi-caracter}
\end{equation}
El término $-\alpha$ de $H_5$ se cancela con el $+\alpha$ del
contraste angular antes de esta lectura. El denominador de la
continuación regional permanece completo. La igualdad recibe las
salidas y los lectores HMT anteriores, y reúne sus dos representaciones;
no define retroactivamente $\alpha$.

En la rama positiva de acción y radio finito,
\begin{equation}
 0<\eta_{\mathrm{ret}}<499\alpha,\qquad
 \chi_{\mathrm{rad}}
 <-2\operatorname{artanh}(1/500).
 \label{md:restop:dominio-positivo}
\end{equation}
Con $\alpha$ fija,
\begin{equation}
 \frac{\partial\eta_{\mathrm{ret}}}{\partial\chi_{\mathrm{rad}}}
 =-250\alpha\,\operatorname{sech}^2(\chi_{\mathrm{rad}}/2)<0.
 \label{md:restop:derivada-accion}
\end{equation}
Por tanto, la sección adimensional de acción determina unívocamente la
clase radial orientada en esta representación, y recíprocamente.
Para variaciones conjuntas,
\begin{equation}
 d\eta_{\mathrm{ret}}
 =\frac{\eta_{\mathrm{ret}}}{\alpha}\,d\alpha
 -250\alpha\,\operatorname{sech}^2(\chi_{\mathrm{rad}}/2)
 \,d\chi_{\mathrm{rad}}.
 \label{md:restop:diferencial-accion}
\end{equation}

El carácter $H_5$, la corrección regional y la anisotropía radial orientada
mantienen así una relación de compatibilidad y no constituyen tres
parámetros independientes. El carácter analítico y la geometría de
reciprocidad representan la misma sección de retorno.
Este nivel estructural de la ecuación de Planck permanece explícito
antes de componer la acción con una frecuencia.

La década de acción conserva su genealogía determinantal, independiente
de esta reparametrización:
\begin{equation}
 \hbar_{\mathrm{ret}}
 =10^{-34}\mathcal U_S\eta_{\mathrm{ret}},\qquad
 \hbar_{\mathrm{pre}}=10^{-34}\mathcal U_SH_5.
 \label{md:restop:secciones-accion}
\end{equation}
Se conservan además la acción por vuelta $h=2\pi\hbar$ y la lectura
térmica
\begin{equation}
 k_B\Theta_{\mathrm{clk}}\log3=\frac{h}{T_{108}}
 \label{md:restop:lectura-termica}
\end{equation}
en su ámbito. El reloj físico de una realización conserva su constructor:
la coordenada $\chi_{\mathrm{rad}}$ clasifica esta reciprocidad y no la
totalidad de historias TPK. La clasificación por
$\chi_{\mathrm{rad}}$ y la relación diferencial reúnen las representaciones
existentes de la acción; mantienen su carácter de formalización focal y no
constituyen una segunda derivación independiente de Planck.

\subsection{Circuitos de fase, memoria y conservatividad}
\label{md:restop:conservatividad}

La holonomía nonádica restaura la fase e incrementa la memoria, según
la sección~\ref{mono-ampl:vueltas}. Si un recorrido cerrado en las fases
aumenta el contador entero, ese incremento no puede escribirse como
diferencia de un potencial dependiente únicamente de la fase.

La demostración se obtiene por telescopía: la suma
\begin{equation}
 \sum_j\bigl[F(g_{j+1})-F(g_j)\bigr]=0
 \label{md:restop:potencial-fase}
\end{equation}
en un circuito de fases contrasta con el incremento de memoria, igual
a uno por vuelta. Sobre el estado ampliado, el contador $m$ registra
ese incremento. El recorrido cerrado en la proyección conservaba,
por tanto, una apertura en el estado completo.

El estudio de conservatividad comienza por declarar el espacio en el
que se cierra el recorrido y las coordenadas que participan en la
dinámica. No conmutatividad y no conservatividad son propiedades
diferentes. La sección~\ref{md:restop:memoria} controla el orden de
operaciones; esta sección distingue un circuito proyectado de un circuito
completo. Para estudiar una fuerza concreta, el funcional de acción o
trabajo debe actuar sobre el mismo espacio de estados y el mismo transporte.

El crecimiento de memoria no asigna por sí mismo una producción
termodinámica. La sección~\ref{md:cron:crecimiento} demuestra que la
información de bloques puede crecer con tasa entrópica nula.
La restitución operatoria añade una precisión: conservar información
puede significar conservar qué transformaciones se pueden componer,
además de mantener los símbolos almacenados.

\subsection{Compatibilidad dinámica entre observaciones}
\label{md:restop:sintesis}

La unidad de estudio es la compatibilidad dinámica entre observaciones,
con la memoria como parte de la operación. El valor de una constante
constituye una lectura; su estructura contiene restricciones sobre otras
lecturas y sobre las transformaciones admisibles. La prolongación del
análisis consiste en obtener relaciones que impidan modificar de manera
independiente valor, orientación, memoria, forma y acción.

Las tres elaboraciones focales recorren ese mismo eje:
\begin{itemize}
 \item Un ciclo completo de regímenes TRIT conserva todavía un
 desplazamiento nonádico de capacidad: la regularidad local no agota
 el retorno.
 \item La reentrada de memoria modifica exactamente la respuesta al orden
 de dos operaciones: el historial interviene causalmente en la composición.
 \item El carácter $H_5$ y la geometría radial orientada determinan
 mutuamente su sección de retorno: la acción admite una representación
 estructural recuperable.
\end{itemize}
Estas consecuencias continúan la genealogía de $\pi,\varphi,e,\alpha$ y
explicitan qué significa conservarla al pasar a las operaciones y la
acción. La composición posterior reúne esos niveles mediante los mapas
comunes ya construidos, manteniendo la distinción entre sus índices y
sus memorias aunque compartan una denominación.

% Focal insertion; hierarchy inherited from the host.
\begingroup
% Fragmento maestro: incluir desde el anfitrion, con \HMTFGRoot fijada a la raiz del paquete.
% Requiere amsmath, amssymb/mathrsfs o unicode-math, amsthm, booktabs, longtable, hyperref, needspace.
% El anfitrion conserva sus entornos theorem, proposition, lemma y proof.
\providecommand{\HMTFGRoot}{.}
\begingroup
\providecommand{\ZZ}{}\renewcommand{\ZZ}{\mathbb Z}
\providecommand{\RR}{}\renewcommand{\RR}{\mathbb R}
\providecommand{\FF}{}\renewcommand{\FF}{\mathbb F}
\providecommand{\APP}{}\renewcommand{\APP}{\mathrm{APP}}
\providecommand{\TPK}{}\renewcommand{\TPK}{\mathrm{TPK}}
\providecommand{\TRIT}{}\renewcommand{\TRIT}{\{-1,0,+1\}}
% unicode-math usa lBrack/rBrack; newtxmath ya proporciona llbracket/rrbracket.
\providecommand{\llbracket}{\lBrack}
\providecommand{\rrbracket}{\rBrack}
% Fragmento de inserción; se incluye desde la raíz de este paquete editorial.
% Las fuentes completas y los cortes materiales constan en ANTECEDENTES_Y_DEPENDENCIAS.md.
% Procedencia: XI REV11 ES, sections/nucleo.tex, tpk_desarrollo_integrado.tex,
% ch_tres_vueltas_ext.tex, ch_elevacion_catalogal.tex, ch_cierre_cardinal.tex,
% ch_descenso_por_fibras.tex; perfil integral public_final/feigenbaum/c37_feigenbaum.tex;
% Grassmannianos REV07 ES, sections/jacobi_spectral.tex;
% TRANSVERSALIDAD_Y_ESCALADO_LOCAL_FEIGENBAUM.md, sección 14.
% Ampliación expositiva de demostraciones existentes; no introduce resultados nuevos.
\providecommand{\ZZ}{\mathbb Z}
\providecommand{\RR}{\mathbb R}
\providecommand{\FF}{\mathbb F}
\providecommand{\APP}{\mathrm{APP}}
\providecommand{\TPK}{\mathrm{TPK}}
\providecommand{\TRIT}{\{-1,0,+1\}}

\section{Antecedentes operatorios de la familia de criticidad}
\label{hmtfg:antecedentes-operatorios}

La familia de criticidad se construye a partir de una publicación del
calendario TPK. Las dos hojas aritméticas, la orientación trítica y el
transporte con memoria determinan primero los estados y sus prolongaciones.
La lectura ordenada del continuo permite después evaluar funciones
analíticas, comparar parámetros y seleccionar raíces. En este desarrollo
se conservan el estado genealógico y sus dos publicaciones: la carta
arquimediana proporciona el espacio de evaluación, mientras la publicación
de incidencia conserva las relaciones de frontera. La estructura discreta
conjunta del continuo mantiene sus cinco construcciones consustanciales.
El presente corte expositivo utiliza su lector ordenado y conserva el
antecedente común del que procede.

\subsection{Células aritméticas, orientación y transporte}
\label{hmtfg:app-trit-transporte}

Se utilizan la célula levantada \eqref{eq:app-levantada}, las coordenadas tríticas y su ley aritmética (\ref{trit-ampl-def-coordenadas}, \ref{trit-ampl-prop-aritmetica}), y el selector~\ref{tpk-int:seleccion}. Para fijar la notación, sean $D_9=\{1,\ldots,9\}$, $X_9=(\mathbb Z/9\mathbb Z)^2$ y, para todo $n\in\mathbb Z$, $\rho_9(n)=1+((n-1)\bmod9)$ y $q_9(n)=(n-\rho_9(n))/9$. El representante balanceado es $r_3(n)\in\{-1,0,+1\}$; su acarreo es $\chi(a,b)=(a+b-r_3(a+b))/3$. Estas fórmulas conservan la extensión a enteros de las evaluaciones positivas del núcleo.

El transporte levantado y su reversibilidad están demostrados en \ref{tpk-int:prop-desplazamiento}; la configuración observable y su retorno de $108$ pasos, en \ref{tpk-int:fobs} y \ref{tpk-int:periodo}. La actualización \eqref{eq:actualizacion} conserva el estado completo. La holonomía de \ref{tpk-int:gamma} y los retornos de \ref{hist:tres-retornos} distinguen restauración de fase y avance de memoria. Sobre estas operaciones se construye a continuación la realización concreta de las tres ramas y su registro íntegro.

\input{inserciones_20260930/feigenbaum/00a_transiciones_nonadicas.tex}

\subsection{Refinamiento de cilindros y lectura ordenada}
\label{hmtfg:refinamiento-ordenado}

% XI REV11, ch_elevacion_catalogal.tex: bloques de 54 aristas, beta_blk,
% partición efectiva de 729 subcilindros; ch_cierre_cardinal.tex:375--497.
Para \(\mathbf b=(b_1,\ldots,b_6)\in\FF_3^6\), sea
\(\upsilon:\FF_3\to\{-1,0,+1\}\) el representante balanceado.
Las tres prolongaciones construidas arista por arista definen
\[
 \Gamma^{[54]}_{\mathbf b,k}
 =\gamma_{\upsilon(b_6),k+45}\circ\cdots\circ
   \gamma_{\upsilon(b_1),k}.
\]
Partiendo de \(Y_0=\{x_\star\}\), se pone
\[
 Y_{N+1}
 =\{\Gamma^{[54]}_{\mathbf b,1+54N}(x):
          x\in Y_N,\ \mathbf b\in\FF_3^6\}.
\]
El truncamiento \(\rho^{\rm blk}_{N+1,N}\) elimina las últimas
cincuenta y cuatro aristas, recuperando el estado anterior mediante
el registro de transición.

\begin{proposition}[Codificación efectiva de los bloques de prolongación]
\label{hmtfg:bloques-729}
La lectura de las seis parejas tríticas de cada bloque define biyecciones
\[
 \beta_N:Y_N\xrightarrow{\ \cong\ }(\FF_3^6)^N
\]
que conmutan con el truncamiento. Cada estado de \(Y_N\) posee
exactamente \(729\) prolongaciones de bloque en \(Y_{N+1}\).
\end{proposition}
\begin{proof}
En cada vuelta, el registro conserva la pareja trítica y su acarreo.
Las tres parejas producen las tres salidas distintas
\(\chi(\varepsilon,\varepsilon)=\varepsilon\), y el registro permite
recuperar la rama utilizada, incluida la rama de acarreo nulo.
Las seis concatenaciones producen, por tanto, todas las palabras de
\(\FF_3^6\). Dos palabras diferentes se distinguen en la primera pareja
registrada en la que difieren. La recuperación del prefijo prueba la
compatibilidad con el truncamiento. Por inducción sobre \(N\), cada
palabra de longitud \(N\) tiene un único estado antecedente; la siguiente
palabra puede ser cualquiera de las \(3^6=729\) posibilidades.
\end{proof}

Sea \(\Omega_{\Gamma}^{\rm cal}=\varprojlim_NY_N\). Las biyecciones
anteriores inducen
\[
 \Omega_{\Gamma}^{\rm cal}\cong(\FF_3^6)^{\mathbb N}.
\]
La correspondencia y su inversa son continuas: cada coordenada finita
depende de un prefijo finito. Los cilindros son abiertos y cerrados.
Cada cilindro de profundidad \(N\) se particiona en sus \(729\)
subcilindros de profundidad \(N+1\); cada uno contiene una prolongación
infinita, obtenida, por ejemplo, continuando con la rama nula en todos
los bloques posteriores.

% Antecedentes efectivos del continuo conjunto y del universo interno.
% Se imprimen en ambos destinos; la etiqueta de universo no reemplaza su regla.
\input{inserciones_20260930/feigenbaum/00b_continuo_universo.tex}

\subsection{Publicación ordenada de los cilindros de bloque}
\label{hmtfg:publicacion-cilindros-bloque}

Para el representante \([b]_3\in\{0,1,2\}\), se fija
\[
 \nu(\mathbf b)=\sum_{i=1}^6[b_i]_3\,3^{6-i}\in\{0,\ldots,728\}.
\]
En la completación ordenada HMT, un prefijo
\(s=(m;\mathbf b_0,\ldots,\mathbf b_{N-1})\), \(m\in\ZZ\), determina
\[
 a_N(s)=m+\sum_{j=0}^{N-1}\nu(\mathbf b_j)\,729^{-j-1},\qquad
 I_N(s)=[a_N(s),a_N(s)+729^{-N})_{\rm HMT}.
\]
El acarreo conserva el empalme entre los dos representantes de una
frontera. Si \(C_N(s)=\overline{I_N(s)}\), las relaciones anteriores dan
\[
 I_N(s)=\bigsqcup_{\mathbf b\in\FF_3^6}I_{N+1}(s;\mathbf b),
 \qquad \operatorname{diam}C_N(s)=729^{-N}\longrightarrow0.
\]

\begin{proposition}[Cobertura de la lectura arquimediana]
\label{hmtfg:cobertura-arquimediana}
La aplicación
\[
 P_{\rm ar}:\ZZ\times\Omega_\Gamma^{\rm cal}
       \longrightarrow\mathbb R_{\rm HMT},\qquad
 \{P_{\rm ar}(m,\xi)\}
    =\bigcap_N C_N(m;\beta_N(\rho_{\infty,N}\xi))
\]
está bien definida, es sobreyectiva y determina
\[
 (\ZZ\times\Omega_\Gamma^{\rm cal})/\!\sim_{P_{\rm ar}}
       \ \cong\ \mathbb R_{\rm HMT}.
\]
\end{proposition}
\begin{proof}
Los extremos izquierdos forman una sucesión de Cauchy en la completación
HMT, pues los intervalos son anidados y sus diámetros tienden a cero.
Su límite pertenece a todos los intervalos cerrados y es único.
Para cualquier elemento de la completación, se elige primero el intervalo
entero que lo contiene y después el único subintervalo semiabierto de
cada partición que lo contiene. La proposición anterior realiza cada
una de esas elecciones mediante una prolongación TPK efectiva.
Los prefijos compatibles definen \(\xi\), cuya lectura es el elemento
elegido. En las fronteras, el acarreo identifica las escrituras
equivalentes después de su construcción. El cociente por igualdad
de lectura resulta, por definición y sobreyectividad, biyectivo con
la completación.
\end{proof}

La construcción se realiza asimismo en el universo genealógico interno
\(\mathfrak G=\mathfrak G_{\rm HMT}(h,m_\infty)\). Aquí \(h\) codifica
las reglas operatorias y \(m_\infty\) la historia compatible; el universo
se obtiene mediante clausura transitiva inicial, definibilidad en cada
sucesor y unión en los estadios límite:
\[
 G_0=\operatorname{TC}\{\omega,u_{h,m},h,m_\infty\},\qquad
 G_{\beta+1}=\operatorname{Def}_{\Sigma_{\rm HMT}}(G_\beta),\qquad
 G_\lambda=\bigcup_{\beta<\lambda}G_\beta .
\]
Las relaciones de \(\Sigma_{\rm HMT}\) son las de las operaciones ya
construidas y sus códigos. Se mantiene, por tanto, el dominio interno
en todas las afirmaciones de cobertura. Denotemos por
\(\mathbb R_{\rm HMT}^{\mathfrak G}\) su completación y por
\(\mathbb R^{\mathfrak G}\) la realización ordenada dentro del mismo
universo.

\begin{proposition}[Carta ordenada interna]
\label{hmtfg:carta-interna}
Existe un único isomorfismo de cuerpos ordenados arquimedianos completos
\[
 \iota:\mathbb R_{\rm HMT}^{\mathfrak G}
             \xrightarrow{\ \cong\ }\mathbb R^{\mathfrak G}
\]
que fija la unidad. Conserva los extremos racionales, las raíces
cuadradas positivas y los límites de sucesiones convergentes.
\end{proposition}
\begin{proof}
La unidad fija los enteros y los racionales. A cada elemento de la
completación HMT se asocia el corte de los racionales menores que él;
la completitud de la realización ordenada proporciona el único elemento
con ese corte. La densidad racional prueba inyectividad y conservación
estricta del orden. El mismo procedimiento en sentido inverso prueba
sobreyectividad. Las aproximaciones racionales por exceso y por defecto
conservan suma y producto y, con ello, la estructura de cuerpo.
El elemento positivo cuyo cuadrado es \(a>0\) se transporta al elemento
positivo cuyo cuadrado es \(\iota(a)\). Si una sucesión converge, para
cada tolerancia racional positiva sus términos finales están en el
entorno racional del límite; la conservación del orden transporta esa
condición. Todas estas operaciones se efectúan dentro de
\(\mathfrak G\). La unicidad sigue de la conservación de los cortes
racionales.
\end{proof}

\subsection{Descenso por fibras y conservación del selector}
\label{hmtfg:descenso-selector}

% XI REV11, ch_descenso_por_fibras.tex:10--39; prueba completa pertinente.
\begin{proposition}[Criterio de descenso de una operación]
\label{hmtfg:criterio-descenso}
Sea \(q:\Omega\twoheadrightarrow J\) y
\(\star:D\subseteq\Omega^2\to\Omega\), con \(D\) saturado por las
fibras de \(q\times q\). Existe una única operación \(\bar\star\)
sobre \((q\times q)(D)\) tal que
\[
 q(x\star y)=q(x)\,\bar\star\,q(y)
\]
si y sólo si
\[
 q(x)=q(x'),\quad q(y)=q(y')
 \ \Longrightarrow\ q(x\star y)=q(x'\star y')
\]
para los pares del dominio.
\end{proposition}
\begin{proof}
Si existe \(\bar\star\), ambos resultados son la evaluación de la
misma pareja de argumentos. Recíprocamente, se eligen antecedentes
de los argumentos visibles y se define el resultado como la lectura
de su operación. La saturación fija el dominio; la condición sobre
fibras hace independiente el resultado de los antecedentes elegidos.
La sobreyectividad prueba la unicidad.
\end{proof}

En una aplicación mediante \(P_{\rm ar}\), este criterio identifica la
condición precisa para transportar una operación desde estados completos
hasta su lectura. En la carta \(\iota\), la estructura de cuerpo y
los límites ya satisfacen esa condición. A continuación se construye
el perfil analítico a partir del calendario y se utiliza esta
compatibilidad para transportar su iteración y sus raíces.

\subsection{Publicación dodecafásica y normalización cuadrática}
\label{hmtfg:perfil-dodecafasico}

% Propietario integral c37_feigenbaum.tex: comienzo hasta la familia f_lambda.
La sección orientada de criticidad sobre las doce ventanas del calendario
es
\[
 \mathcal P_{\rm crit}(m)=\varphi_m
   =(30m-10)\frac{\pi_{\rm HMT}}{180}
   =\pi_{\rm HMT}\frac{3m-1}{18},\qquad m=1,\ldots,12.
\]
Esta sección fija el origen y la orientación de sus representantes
angulares. El lector uniforme de las doce ventanas es
\[
 \eta_0(F)=\frac1{12}\sum_{m=1}^{12}F(\varphi_m).
\]
Se distingue la función \(\eta_0\), que asigna pesos a estas ventanas,
de la carta ordenada \(\iota\). La normalización uniforme y la sección
angular forman parte de los datos operatorios de este perfil.

Poniendo \(a_m=(3m-1)/18\), se obtiene
\[
 D_0=\frac1{12}\sum_{m=1}^{12}a_m^2=\frac{899}{648},
 \qquad
 s_0=\frac{36}{\pi_{\rm HMT}\sqrt{899}},\qquad
 k_m=s_0\varphi_m=\frac{2(3m-1)}{\sqrt{899}}.
\]
En efecto, \(\sum_{m=1}^{12}(3m-1)^2=5394=6\cdot899\), por lo
que \(s_0^2\pi_{\rm HMT}^2D_0=2\). El perfil es
\begin{equation}
 u_0(x)=\frac1{12}\sum_{m=1}^{12}\bigl[1-\cos(k_mx)\bigr],
 \qquad f_\lambda(x)=1-\lambda u_0(x).
 \label{hmtfg:perfil-familia}
\end{equation}
La constante \(\pi_{\rm HMT}\) entra como publicación angular del
núcleo HMT y se cancela algebraicamente en la normalización del perfil.
La escala \(\lambda\) es el parámetro de la familia, cuyo valor crítico
será seleccionado mediante la dinámica.

\begin{proposition}[Normalización y forma cerrada del perfil]
\label{hmtfg:forma-perfil}
El perfil \(u_0\) es par y analítico, satisface
\(u_0(0)=u_0'(0)=0\) y \(u_0''(0)=2\), y, para
\(y=x/\sqrt{899}\), admite la expresión
\[
 u_0(x)=1-\frac{\cos(37y)\sin(36y)}{12\sin(3y)}
       =1-\frac{\sin(73y)-\sin(y)}{24\sin(3y)}.
\]
Las singularidades aparentes se interpretan por continuidad.
Su expansión inicial es
\[
 u_0(x)=x^2-\frac{715769}{2424603}x^4+O(x^6).
\]
\end{proposition}
\begin{proof}
La suma finita de cosenos conserva paridad y analiticidad. Su derivada
segunda en cero es \(12^{-1}\sum_m k_m^2=2\).
La suma de la progresión trigonométrica
\(\sum_{m=1}^{12}\cos((6m-2)y)\) es
\(\cos(37y)\sin(36y)/\sin(3y)\); la identidad producto--suma
proporciona la segunda expresión. La serie de coseno y
\(\sum_m(3m-1)^4=4294614\) dan el coeficiente de grado cuatro
indicado. La expresión como suma finita fija la extensión en los
ceros del denominador de la forma cerrada.
\end{proof}

\subsection{Transporte de iteradas, raíces y mínimos}
\label{hmtfg:transporte-raices}

% Reunión existente: TRANSVERSALIDAD_Y_ESCALADO_LOCAL_FEIGENBAUM.md, sección 14.1.
Dentro de \(\mathbb R_{\rm HMT}^{\mathfrak G}\), la serie convergente
\[
 \cos_{\rm HMT}(t)=\sum_{j=0}^{\infty}
                   \frac{(-1)^j t^{2j}}{(2j)!}
\]
define la realización analítica correspondiente. Sean
\[
 u_{\rm HMT}(x)=\frac1{12}\sum_{m=1}^{12}
  \left[1-\cos_{\rm HMT}
   \left(\frac{2(3m-1)x}{\sqrt{899}_{\rm HMT}}\right)\right],
 \qquad F_a(x)=1-a\,u_{\rm HMT}(x),
\]
\[
 S_n^{\rm HMT}(a)=F_a^{\,2^n}(0),\qquad
 S_n(\lambda)=f_\lambda^{\,2^n}(0).
\]

\begin{proposition}[Conjugación ordenada de la familia y del selector de raíces]
\label{hmtfg:conjugacion-selector}
Para los argumentos internos de las expresiones anteriores,
\[
 \iota\circ F_a=f_{\iota(a)}\circ\iota,\qquad
 \iota(S_n^{\rm HMT}(a))=S_n(\iota(a)).
\]
La carta transporta exactamente las condiciones de retorno, periodo
exacto y orden de los parámetros. Si un conjunto de raíces admisibles
posee mínimo, la imagen de ese mínimo es el mínimo del conjunto imagen.
\end{proposition}
\begin{proof}
La carta fija los racionales, transporta la raíz cuadrada positiva y
conserva los límites de las sumas parciales de coseno. Por ello,
\(\iota(u_{\rm HMT}(x))=u_0(\iota(x))\). La compatibilidad con
suma y producto prueba la primera igualdad; la inducción sobre el
número de iteraciones prueba la segunda. Una raíz se caracteriza por
la igualdad con cero, que se conserva en ambas direcciones. El
periodo exacto \(2^n\) añade las desigualdades
\(F_a^{2^j}(0)\ne0\), \(0\le j<n\), también conservadas por la
inyectividad de \(\iota\). Si se exige que el parámetro supere un
centro precedente, se transporta esa desigualdad por conservación
del orden. Finalmente, \(a\le b\) para todo \(b\) del conjunto
equivale a \(\iota(a)\le\iota(b)\) para todo elemento de su
imagen. La sobreyectividad sobre los conjuntos así definidos
completa la afirmación acerca del mínimo.
\end{proof}

El dominio de este transporte es el continuo interno indicado.
La existencia de cada mínimo y su identificación con una rama dinámica
concreta son resultados de la construcción de criticidad que sigue;
la carta conserva esas propiedades cuando han sido establecidas.

\subsection{Representación espectral finita del lector uniforme}
\label{hmtfg:jacobi-finito}

% Grassmannianos REV07, sections/jacobi_spectral.tex:35--136,
% mecanismo de Hankel, Gram--Schmidt, recurrencia y representación espectral.
% Especialización atómica ya reunida en la sección 14.2 del desarrollo Feigenbaum.
Las frecuencias cuadradas del perfil definen la medida positiva
\[
 s_m=k_m^2=\frac{4(3m-1)^2}{899},\qquad
 \nu_0=\frac1{12}\sum_{m=1}^{12}\delta_{s_m},\qquad
 M_r=\int s^r\,d\nu_0(s).
\]
Sus doce nodos son distintos y positivos. Para \(1\le r\le12\),
la matriz de Hankel \(H_r=(M_{i+j})_{0\le i,j<r}\) es definida
positiva. En efecto, para \(p\) de grado menor que \(r\),
\[
 \mathbf p^{\mathsf T}H_r\mathbf p
   =\frac1{12}\sum_{m=1}^{12}p(s_m)^2>0
\]
salvo que \(p=0\), porque un polinomio de grado a lo sumo once
que se anula en doce nodos distintos es nulo. En grado doce aparece
el polinomio \(\prod_m(s-s_m)\), de norma nula en esta medida.
La representación espectral correspondiente tiene dimensión doce.

Sean \(p_0,\ldots,p_{11}\) los polinomios mónicos ortogonales,
\(h_j=\int p_j^2\,d\nu_0>0\) y
\(\psi_j=p_j/\sqrt{h_j}\). Como \(\nu_0\) es una probabilidad,
\(\psi_0=1\). La multiplicación por \(s\) en \(L^2(\nu_0)\)
es autoadjunta y posee una matriz tridiagonal \(J_{12}\) en esta base.

\begin{proposition}[Recurrencia y entrelazamiento espectral exactos]
\label{hmtfg:jacobi-entrelazamiento}
Con
\[
 a_j=\frac{\int s p_j(s)^2\,d\nu_0(s)}{h_j},\qquad
 \beta_j=\frac{h_j}{h_{j-1}}>0\quad(1\le j\le11),
\]
la matriz \(J_{12}\) tiene diagonal \(a_j\) y subdiagonal
\(\sqrt{\beta_j}\). Si
\[
 Q_{mj}=\frac{\psi_j(s_m)}{\sqrt{12}}
 \quad(1\le m\le12,\ 0\le j\le11),\qquad
 D=\operatorname{diag}(s_1,\ldots,s_{12}),
\]
entonces
\[
 Q^{\mathsf T}Q=I,\qquad DQ=QJ_{12},\qquad
 J_{12}=Q^{\mathsf T}DQ,
\]
y el perfil satisface
\begin{equation}
 u_0(x)=e_0^{\mathsf T}
          \bigl[I-\cos(x\sqrt{J_{12}})\bigr]e_0.
 \label{hmtfg:perfil-espectral}
\end{equation}
\end{proposition}
\begin{proof}
La positividad de los sistemas de Hankel proporciona de forma única
los polinomios mónicos ortogonales. Al desarrollar \(s p_j\) en la
base ortogonal, los coeficientes de \(p_i\), \(i\le j-2\), se
anulan porque
\(\langle sp_j,p_i\rangle=\langle p_j,sp_i\rangle=0\).
El coeficiente de \(p_{j+1}\) es uno por monicidad, el de \(p_j\)
es \(a_j\) y el de \(p_{j-1}\) es \(h_j/h_{j-1}\).
En el último grado, la misma igualdad se entiende en
\(L^2(\nu_0)\), donde el polinomio mónico de grado doce que se
anula en los nodos representa el vector cero. La normalización da
la subdiagonal positiva indicada.

La ortonormalidad equivale a \(Q^{\mathsf T}Q=I\). La igualdad
de la recurrencia en cada nodo da \(DQ=QJ_{12}\).
Como \(Q\) es cuadrada, también \(QQ^{\mathsf T}=I\), y se
obtiene la diagonalización de \(J_{12}\).
Sus autovalores positivos permiten definir su raíz cuadrada
positiva. Para toda función definida en esos doce nodos,
\[
 e_0^{\mathsf T}\Phi(J_{12})e_0
   =\frac1{12}\sum_{m=1}^{12}\Phi(s_m),
\]
pues \(Qe_0=(1,\ldots,1)^{\mathsf T}/\sqrt{12}\).
La elección \(\Phi(s)=1-\cos(x\sqrt{s})\) prueba
\eqref{hmtfg:perfil-espectral}.
\end{proof}

La representación conserva exactamente el funcional uniforme, sus
momentos y su perfil. Sus primeros coeficientes son
\[
 a_0=M_1=2,\qquad
 \beta_1=M_2-M_1^2=\frac{2493348}{808201}.
\]
La genealogía completa permanece en el estado TPK del que procede
el lector; \(J_{12}\) representa su publicación espectral concreta.
La construcción de Gram--Schmidt y la recurrencia de Jacobi se aplican
a esta medida atómica, con su terminación en dimensión doce.

La cadena operatoria utilizada en las secciones siguientes queda así
fijada: las células de APP y la orientación TRIT alimentan la
transición TPK; sus prolongaciones producen los cilindros y su
lectura ordenada; las doce ventanas orientadas y el funcional uniforme
producen \(u_0\); la familia \(f_\lambda\) determina retornos y raíces;
la carta interna conserva el orden de sus selectores. La realización
espectral describe el mismo perfil mediante una matriz autoadjunta
finita y un vector cíclico.

% Fragmento integrable. Preambulo anfitrion: amsmath, amssymb, longtable, hyperref.
\section{Perfil dodecafásico y retorno cuadrático}
\label{hmtfg:fragmento:perfil-dodecafasico-y-retorno-cuadratico}

\subsection{Genealogía y objeto de la demostración}
\label{hmtfg:desarrollo:1}

El objeto inicial es la publicación crítica de la rueda dodecafásica orientada ya construida en el corpus. APP evalúa suma y producto sobre las dos hojas de la retícula de dígitos, conservando cociente y residuo; TRIT determina régimen y orientación; el transporte TPK selecciona, transporta y actualiza el estado con sus registros de acarreo, frontera y memoria. La conexión nonádica prolonga ese estado, con retorno de fase y avance de memoria. Sus lectores actúan sobre la misma estructura discreta del continuo.

El lector focal de la fuente es
\[
\mathcal P_{\mathrm{crit}}:U_{12}^{\mathrm{sign}}\longrightarrow
\mathbb R/2\pi\mathbb Z,\qquad
m\longmapsto\phi_m=\frac{(3m-1)\pi}{18},\quad 1\le m\le12.
\]
Para evaluar el momento angular se conserva el representante orientado de esa carta: sustituirlo por otro representante módulo \(2\pi\) cambiaría el momento. La traza uniforme fija \(w_m=1/12\). La cadena focal es
\[\begin{aligned}
\text{APP–TRIT–TPK}&\longrightarrow
\text{estructura discreta del continuo}\\ &\longrightarrow U_{12}^{\mathrm{sign}}
\xrightarrow{\mathcal P_{\mathrm{crit}}}(\phi_m,w_m)
\\ &\longrightarrow u_0\longrightarrow f_\lambda
\longrightarrow\text{retornos y renormalización}.
\end{aligned}\]

La escritura abreviada conserva como antecedente el estado completo y sus prolongaciones; la familia escalar es una lectura de ese antecedente. Se toma como estructura de partida explícita el lector ya fijado, con su orden, orientación y sector uniforme. Las conclusiones se refieren a ese lector concreto; la unicidad de ese lector entre todas las publicaciones posibles del TPK queda fuera de estas pruebas.

La coordenada \(\pi_{\mathrm{HMT}}\) empleada por la carta pertenece a las salidas internas del corpus. Su factor se cancela exactamente al normalizar el segundo momento. Los valores de Feigenbaum intervienen exclusivamente en el reconocimiento posterior, nunca en la definición de la familia o en la certificación de raíces.


\subsection{Perfil exacto y criticidad cuadrática}
\label{hmtfg:desarrollo:2}

Sea \(s_m=3m-1\). La suma entera
\[
\sum_{m=1}^{12}s_m^2=5394=6\cdot899
\]
determina
\[
\kappa_0=\frac{36}{\pi\sqrt{899}},\qquad
k_m=\kappa_0\phi_m=\frac{2s_m}{\sqrt{899}},
\]
\[
u_0(x)=\frac1{12}\sum_{m=1}^{12}(1-\cos(k_mx)),\qquad
f_{\lambda,0}(x)=1-\lambda u_0(x).
\]
Por identidad del segundo momento,
\[
u_0(0)=u_0'(0)=0,\quad u_0''(0)=2,\quad
f_{\lambda,0}''(0)=-2\lambda.
\]
La representación analítica conserva coeficientes racionales:
\[
u_0(z)=\sum_{j\ge1}(-1)^{j+1}
\frac{4^j\sum_m s_m^{2j}}{12\,899^j(2j)!}z^{2j}
=z^2-\frac{715769}{2424603}z^4
+\frac{1353832978}{32695771455}z^6+\cdots.
\]
Así, el registro finito de doce sectores produce un perfil entero, con criticidad cuadrática y escala fijadas antes de iterar.


\subsection{Robustez estructural con cotas uniformes}
\label{hmtfg:desarrollo:3}

La deformación ya presente en el código de procedencia es
\[
w_m(\eta)=\frac{1+\eta p_m}{12},\qquad
p_m=12\cos(3\phi_m)-\cos(4\phi_m).
\]
Definimos
\[
M_{2,\eta}=\sum_mw_m(\eta)\phi_m^2,\quad
\kappa_\eta=\sqrt{2/M_{2,\eta}},\quad k_{m,\eta}=\kappa_\eta\phi_m,
\]
\[
u_\eta(x)=\sum_mw_m(\eta)(1-\cos(k_{m,\eta}x)),\qquad
f_{\lambda,\eta}=1-\lambda u_\eta.
\]

\textbf{Teorema de robustez estructural.} Para \(|\eta|\le1/40\), los pesos son positivos y suman uno. El perfil es par, entero, satisface \(u_\eta''(0)=2\) y es estrictamente creciente en \((0,1]\). Para \(0<\lambda\le2/u_\eta(1)\), el mapa lleva \([-1,1]\) en sí mismo, tiene un único punto crítico en ese intervalo y cumple
\[
S f_{\lambda,\eta}(x)\le-\frac{640}{47647}<0,\qquad0<|x|\le1.
\]

\textbf{Demostración.} Las sumas de los cosenos \(3\phi_m\) y \(4\phi_m\) se anulan por ciclos de longitudes cuatro y tres. Como \(|p_m|\le13\),
\[
\frac9{160}\le w_m\le\frac{53}{480},\qquad
\frac{27}{40}M_{2,0}\le M_{2,\eta}\le\frac{53}{40}M_{2,0}.
\]
Se deduce
\[
k_{\max,\eta}^2\le\frac{196000}{24273}<9<\pi^2,\qquad
k_{\min,\eta}^2\ge\frac{640}{47647}.
\]
Para \(0<x\le1\), todos los senos \(\sin(k_{m,\eta}x)\) son positivos, de donde
\[
u_\eta'(x)=\sum_mw_mk_m\sin(k_mx)>0,\qquad
u_\eta'''(x)=-\sum_mw_mk_m^3\sin(k_mx)<0.
\]
El cociente \(u_\eta'''/u_\eta'\) es el negativo de una media positiva de los \(k_m^2\). La invariancia de la schwarziana bajo cambios afines del valor da
\[
S f_{\lambda,\eta}
=\frac{u_\eta'''}{u_\eta'}-\frac32
\left(\frac{u_\eta''}{u_\eta'}\right)^2
\le-k_{\min,\eta}^2.
\]
La paridad cubre \(x<0\), y la imagen del intervalo es
\([1-\lambda u_\eta(1),1]\). \(\square\)

La anchura \(1/40\) es una elección para obtener cotas uniformes, distinta de una constante fundamental o de un umbral óptimo. El sector original sigue siendo \(\eta=0\).


\subsection{Coordenada holomorfa cuadrática global en una franja}
\label{hmtfg:desarrollo:4}

\textbf{Teorema de factorización holomorfa.} Para \(|\eta|\le1/40\), en
\[
\mathcal S=\{z\in\mathbb C:|\operatorname{Re}z|<\pi/3\}
\]
existe una función holomorfa, impar y univalente \(b_\eta\), con \(b_\eta'(0)=1\), tal que
\[
\boxed{u_\eta(z)=b_\eta(z)^2,\qquad
f_{\lambda,\eta}(z)=1-\lambda b_\eta(z)^2.}
\]

\textbf{Demostración.} Para \(z=x+iy\) en la semibanda derecha,
\[
\operatorname{Re}u_\eta'(z)
=\sum_mw_mk_m\sin(k_mx)\cosh(k_my)>0.
\]
La integral de \(u_\eta'\) sobre el segmento entre dos puntos, dividida por la diferencia de los extremos, tiene parte real positiva. La convexidad de la semibanda prueba allí la inyectividad. La paridad da el resultado correspondiente a la izquierda.

Además,
\[
\operatorname{Im}u_\eta(x+iy)
=\sum_mw_m\sin(k_mx)\sinh(k_my)
\]
tiene el signo de \(y\) para \(x>0\). Sobre el eje real positivo \(u_\eta>0\); sobre el imaginario,
\[
u_\eta(iy)=\sum_mw_m(1-\cosh(k_my))
\]
es estrictamente negativo fuera de cero y decrece estrictamente con \(|y|\). Estas separaciones, la inyectividad de cada semibanda y la paridad prueban que las fibras de \(u_\eta\) en \(\mathcal S\) son exactamente \(\{z,-z\}\), salvo la fibra del origen.

La función \(u_\eta(z)/z^2\), prolongada con valor uno en cero, es holomorfa y carece de ceros en la franja simplemente conexa. Tiene un logaritmo holomorfo \(L_\eta\) con \(L_\eta(0)=0\), que es par por unicidad de la normalización. Definimos
\[
b_\eta(z)=z\exp\!\left(\frac12L_\eta(z)\right).
\]
Entonces \(b_\eta^2=u_\eta\), \(b_\eta\) es impar y \(b_\eta'(0)=1\). La igualdad \(b_\eta(z)=b_\eta(w)\) obliga a \(w=z\) o \(w=-z\). El segundo caso sólo da igualdad en el origen. Esto prueba su univalencia. \(\square\)

Las doce fases determinan así una deformación conforme explícita del pliegue cuadrático. La conjugación dinámica completa conserva esa deformación:
\[
b_\eta\circ f_{\lambda,\eta}\circ b_\eta^{-1}(w)
=b_\eta(1-\lambda w^2),
\]
en sus dominios de composición. La factorización demostrada y una conjugación a toda la familia cuadrática son afirmaciones distintas.


\subsection{Retorno, cambio de escala y memoria}
\label{hmtfg:desarrollo:5}

Sean \(a=-f(1)>0\), \(H_a(x)=-ax\), \(J=H_a([-1,1])\). Cuando \(J\subset[-1,1]\) es admisible para el retorno, \(f^2(J)\subset J\) y las composiciones están definidas,
\[
\mathcal Rf=H_a^{-1}f^2H_a,\qquad
\mathcal Rf(x)=-a^{-1}f(f(-ax)).
\]
Si además \(f(J)\subset(0,1]\), se conserva un único máximo cuadrático central:
\[
(\mathcal Rf)(0)=1,\qquad
(\mathcal Rf)''(0)=-a f'(1)f''(0)<0.
\]
La composición de schwarzianas da, en los puntos regulares,
\[
S(\mathcal Rf)(x)=a^2
\left[(Sf)(f(-ax))f'(-ax)^2+(Sf)(-ax)\right]<0.
\]
Las inclusiones de dominios forman parte de las hipótesis y se comprueban a cada escala.

Sea \(v_\eta=(u_\eta|_{[0,1]})^{-1}\). Las ramas inversas son
\[
\beta_\sigma(y)=\sigma v_\eta((1-y)/\lambda),\quad
\sigma\in\{-1,+1\},
\]
con el registro crítico \(\beta_0(1)=0\). El paso
\[
(x,w)\longmapsto(f(x),w\mathbin{\|}\sigma(x))
\]
se invierte sobre su imagen leyendo la última hoja y aplicando \(\beta_\sigma\). La palabra conserva decisiones de rama; el valor anterior se reconstruye por la ecuación del mapa.

El retorno agrupa dos transiciones. Desde el extremo \(x'\) y sus hojas \((\sigma_0,\sigma_1)\) se recupera
\[
x=H_a^{-1}\beta_{\sigma_0}
\!\left(\beta_{\sigma_1}(H_ax')\right).
\]
También se reconstruyen los puntos intermedios. Los registros APP–TRIT–TPK de hoja, acarreo, ruta, frontera y memoria se concatenan conservando su identidad. El signo de la rama analítica es un registro adicional: no reemplaza esos datos por una etiqueta binaria.

Si \(C_N\) agrupa las historias de longitud \(2N\) en pares, las restricciones satisfacen
\[
\pi^{\mathrm{ret}}_{N,M}C_N
=C_M\pi^f_{2N,2M},\qquad M\le N.
\]
Ambos miembros suprimen los mismos pares terminales y conservan las mismas ramas inversas. Queda demostrado el transporte compatible de historias finitas.


\subsection{Escalas acumuladas y derivada del operador}
\label{hmtfg:desarrollo:6}

Sean \(c_n=f^{2^n}(0)\), \(c_0=1\) y \(H_n(x)=c_nx\). Mientras las composiciones sean admisibles y \(c_n\ne0\),
\[
\boxed{\mathcal R^nf=H_n^{-1}f^{2^n}H_n},\qquad
(\mathcal R^nf)(1)=\frac{c_{n+1}}{c_n}=-a_n.
\]
La inducción se obtiene de \(H_nH_{a_n}=H_{n+1}\). La convención \(a_n>0\) exige alternancia de signo de \(c_n\). En una raíz superestable de período \(2^N\), \(c_N=0\): esa normalización se hace singular y deben utilizarse los niveles anteriores o un límite controlado.

Para \(z=-ax\), \(w=f(z)\) y una perturbación \(h\), la derivada completa es
\[
\boxed{
D\mathcal R_f[h](x)=
-\frac{h(w)+f'(w)h(z)}a+
\frac{h(1)}a
\left[\mathcal Rf(x)-xf'(w)f'(z)\right].
}
\]
Se deriva usando \(\dot a=-h(1)\), \(\dot z=xh(1)\),
\(\dot w=h(z)+xf'(z)h(1)\) y la variación del denominador. Si \(h(0)=0\), resulta \(D\mathcal R_f[h](0)=0\). La dirección paramétrica inicial HMT es \(h=-u_\eta\).

Esta fórmula incluye la variación de escala; congelar \(a\) daría otro operador linealizado.


\subsection{Raíces e itinerarios certificados hasta período 256}
\label{hmtfg:desarrollo:7}

Sea \(F_n(\lambda)=f_{\lambda,0}^{2^n}(0)\). Se calcula
\[
x_0=p_0=0,\qquad x_{j+1}=1-\lambda u_0(x_j),\qquad
p_{j+1}=-u_0(x_j)-\lambda u_0'(x_j)p_j.
\]
Entonces \(F_n=x_{2^n}\) y \(F_n'=p_{2^n}\).

El programa adjunto usa intervalos racionales con redondeo entero hacia afuera, raíz cuadrada encerrada por raíz entera y seno/coseno con resto de Taylor explícito. La cota exacta \(|u_0''|\le2\) permite evaluación centrada con resto cuadrático.

Las aproximaciones históricas sólo proponen cajas de radio \(10^{-42}\). Se certifican signos opuestos en sus extremos, \(F_n'\) separado de cero en toda la caja y exclusión de los retornos de períodos propios divisores de \(2^n\). Valor intermedio y monotonía prueban existencia y unicidad local; la exclusión de divisores prueba período mínimo.

\begingroup\small
\begin{longtable}{p{.12\linewidth}p{.17\linewidth}p{.62\linewidth}}
\hline
Nivel & Período mínimo & Centro de la caja, redondeado \\
\hline
1 & 2 & 1.345913990471832792210949 \\
2 & 4 & 1.763379787846910127290794 \\
3 & 8 & 1.850413796950136464530567 \\
4 & 16 & 1.868979756606798125150574 \\
5 & 32 & 1.872955034435659740004205 \\
6 & 64 & 1.873806367467030119412262 \\
7 & 128 & 1.873988695221365362307169 \\
8 & 256 & 1.874027744154450672897008 \\
\hline\end{longtable}
\endgroup

La primera raíz tiene la expresión exacta \(\lambda_1=1/u_0(1)\). Las cajas completas y sus comprobaciones están en el certificado JSON.

Se han encerrado los 502 estados preclausura de estos ocho niveles, todos separados de cero, y sus palabras completas de signos. En particular,
\[
\operatorname{sign}f_{\lambda_n,0}^{2^j}(0)=(-1)^j,\qquad0\le j<n.
\]
Quedan acreditadas las orientaciones finitas de los reescalados anteriores a la clausura. Las inclusiones de intervalos restrictivos a todas las escalas constituyen una comprobación adicional.

Para \(\delta_{\mathrm F,n}=(\lambda_{n-1}-\lambda_{n-2})/(\lambda_n-\lambda_{n-1})\), un intervalo decimal exterior abreviado es
\[
\boxed{\begin{gathered}
4.66921218915020415455609621588966392804<\delta_{\mathrm F,8},\\
\delta_{\mathrm F,8}<4.66921218915020415455609621588966392863.
\end{gathered}}
\]
La anchura del intervalo completo es inferior a \(5.808\cdot10^{-37}\). Mide el error numérico de la razón finita \(\delta_{\mathrm F,8}\), distinto del error de aproximación al límite universal.


\subsection{Persistencia analítica de las raíces deformadas}
\label{hmtfg:desarrollo:8}

Para cada nivel certificado, \(F_n'(\lambda_n)\ne0\). La función implícita produce una rama analítica local \(\lambda_n(\eta)\) con
\[
f_{\lambda_n(\eta),\eta}^{2^n}(0)=0.
\]
Por continuidad, los retornos propios siguen separados de cero para \(\eta\) suficientemente pequeño; se preservan período mínimo e itinerario. La existencia de esos entornos locales queda demostrada. Su anchura común aún no se identifica con toda la franja estructural \(1/40\).

Con \(M=M_{2,\eta}\), la respuesta del perfil es
\[
v_\eta(x)=\partial_\eta u_\eta(x)
=\sum_m\frac{p_m}{12}(1-\cos(k_{m,\eta}x))
-\frac{M'}{2M}\,x u_\eta'(x).
\]
El segundo término conserva la normalización del segundo momento. Si \(q_j=\partial_\eta x_j\) a \(\lambda\) fija,
\[
q_0=0,\qquad q_{j+1}=-\lambda v_\eta(x_j)-\lambda u_\eta'(x_j)q_j,
\qquad
\boxed{\lambda_n'(\eta)=-q_{2^n}/p_{2^n}}.
\]
Así queda explícita la cadena pesos–escala–respuesta–parámetro superestable. Esta no degeneración finita es distinta de la transversalidad asintótica a la hoja estable del punto fijo universal.


\subsection{Colas analíticas y propagación del error}
\label{hmtfg:desarrollo:9}

Para \(\eta=0\), sean
\[
\mu_{2j}=\frac1{12}\sum_mk_m^{2j},\quad
p_N(z)=\sum_{j=1}^N(-1)^{j+1}\frac{\mu_{2j}}{(2j)!}z^{2j},
\quad\Omega=\frac{70}{\sqrt{899}}.
\]
Si \(q_N=\Omega^2r^2/[(2N+3)(2N+4)]<1\),
\[
\|u_0-p_N\|_r\le
\frac{\mu_{2N+2}r^{2N+2}}{(2N+2)!(1-q_N)}.
\]
La desigualdad \(\mu_{2j+2}\le\Omega^2\mu_{2j}\) domina la cola por una serie geométrica. Para \(0\le s\le2N+2\),
\[
\|(u_0-p_N)^{(s)}\|_r\le
\frac{\mu_{2N+2}r^{2N+2-s}}{(2N+2-s)!}
\left(1-\frac{\Omega^2r^2}{(2N+3-s)(2N+4-s)}\right)^{-1},
\]
cuando el denominador es positivo.

Para transportar errores por un retorno se controlan también sus dominios: si \(\|f-g\|\le\varepsilon\), \(a_f,a_g\ge a_*>0\), las composiciones quedan en un dominio convexo común y allí \(\|f'\|\le L,\ \|g\|\le M\), entonces
\[
\|\mathcal Rf-\mathcal Rg\|_r
\le\varepsilon\left[\frac{1+L+L^2r}{a_*}+\frac{M}{a_*^2}\right].
\]
La prueba suma la perturbación exterior, la interior, el cambio del argumento por \(|a_f-a_g|\le\varepsilon\) y la variación del denominador. Estas cotas evitan reutilizar la suma inicial de cosenos como si todos los perfiles renormalizados conservaran esa misma forma.

% Fragmento integrable. Preambulo anfitrion: amsmath, amssymb, longtable, hyperref.
\section{Profundidad nonádica y persistencia del itinerario}
\label{hmtfg:fragmento:profundidad-nonadica-y-persistencia-del-itinerario}

\subsection{Crecimiento exterior y resolución interior}
\label{hmtfg:prolongacion:2}

La ley de refinamiento del artículo XI, con base entera \(B\ge2\), es
\[
L_c(x,y)=(Bx-c,By+c),\qquad 0\le c<B.
\]
Para una palabra \(c_1\ldots c_m\), su registro es
\[
C_m=\sum_{j=1}^m c_jB^{m-j},\qquad
(x_m,y_m)=(B^mx_0-C_m,B^my_0+C_m).
\]
La suma exterior vale \(B^mM_0\), con \(M_0=x_0+y_0\). La lectura normalizada del registro pertenece al cilindro
\[
I_m(C_m)=\left[\frac{C_m}{B^m},\frac{C_m+1}{B^m}\right],
\qquad \operatorname{diam}I_m=B^{-m}.
\]
La pareja entera y la escala recuperan los cocientes, residuos y ceros iniciales de la palabra. Los extremos con doble expansión conservan su marca de frontera. Así se coordinan crecimiento de soporte y refinamiento de lectura sin identificar historias distintas por su sola imagen límite.

La composición es exacta:
\[
L_D^{[b]}\circ L_C^{[a]}=L_{B^bC+D}^{[a+b]}.
\]
Por asociatividad, agrupar una historia de dieciocho eventos en nueve bloques de dos o en dos bloques de nueve produce el mismo registro. A nivel de retorno \(2^N\), el horizonte común es \(9\,2^N\). La prueba es una igualdad de composición sobre la misma historia ordenada; mantiene sus acarreos y sus subbloques recuperables.

La restricción de estados y la prolongación de sus lectores se expresan mediante los límites siguientes:
\[
X_\infty=\varprojlim X_m,\qquad
r^*a=a\circ r,\qquad
r^*e_x=\sum_{r(y)=x}e_y,
\qquad \mathfrak H_{\rm cyl}=\varinjlim\mathfrak H_m.
\]
Los estados se restringen; los lectores se prolongan. La preimagen preserva productos, unidad e idempotentes. Una familia exactamente compatible induce un lector en el límite directo. Los aproximantes que sólo son compatibles hasta error requieren, además, una estimación de convergencia en su completación. La sección siguiente proporciona esa estimación para el retorno crítico.

\subsubsection{Capacidad y vacancias del registro}

El lector de capacidad de XI compara bloques de \(729\) y \(1000\) posibilidades. Con \(\rho=\log_{1000}729=\log_{10}9\), conserva
\[
\mathcal C_k=\lfloor(k+1)\rho\rfloor,\qquad
\mathcal C_k-\mathcal C_{k-1}=1-z_k,
\]
\[
\mathcal C_b-\mathcal C_a+\sum_{k=a+1}^b z_k=b-a,
\qquad T_{\rm cap}(a)=\lceil a/\rho\rceil-1\quad(a\ge1).
\]
Cada vacancia registra una transición durante la cual se mantiene la capacidad; la longitud cronológica queda conservada en el balance. La fase de capacidad y la fase cronológica son lectores distintos.

Para almacenar un prefijo de \(m\) dígitos de base nueve, una capacidad decimal suficiente es
\[
a(m)=\left\lceil\frac{m\rho}{3}\right\rceil.
\]
El cociente y el residuo permiten convertir el entero \(C_m\) entre bases, manteniendo aparte la profundidad \(m\). Se cumplen
\[
9^m\le1000^{a(m)}\le729^{T_{\rm cap}(a(m))+1}.
\]
Estas desigualdades también se deciden mediante potencias enteras, sin redondear logaritmos. El presupuesto de lectura de la sección siguiente se traduce así a capacidad del registro conservando las unidades. Esta cuenta dimensiona el almacenamiento; las reglas de supervivencia determinan, adicionalmente, qué historias nativas son admisibles.


\subsection{Presupuesto explícito de profundidad para cada retorno}
\label{hmtfg:prolongacion:3}

\textbf{Teorema.} Fijado \(|\eta|\le1/40\), normalícense
\[
t=\lambda u_\eta(1)/2\in[0,1],\qquad z=(x+1)/2,
\]
\[
F_{t,\eta}(z)=1-t\frac{u_\eta(2z-1)}{u_\eta(1)},\qquad z\in[0,1].
\]
Para todo entero \(q\ge1\) y dos parámetros \(t,s\),
\[
\left|F_{t,\eta}^{q}(1/2)-F_{s,\eta}^{q}(1/2)\right|
\le S_q|t-s|,
\qquad S_q=\sum_{j=0}^{q-1}(6\sqrt2)^j<9^q.
\]
En consecuencia, un cilindro de parámetros en base nueve, de profundidad \(m=2^N+r\), encierra la lectura del retorno \(q=2^N\) en un intervalo de diámetro inferior a \(9^{-r}\).

\textbf{Demostración.} La función \((1-\cos v)/v^2\) decrece para \(0<v<\pi\): el signo de su derivada es el de \(v\sin v-2(1-\cos v)\), negativo porque \(\tan(v/2)>v/2\). Por el segundo momento,
\[
u_\eta(1)\ge\frac{1-\cos3}{9}\sum_jw_jk_j^2
=\frac{2(1-\cos3)}9>\frac13.
\]
La última desigualdad se verifica, por ejemplo, acotando \(\cos3\) mediante su suma de Taylor alternada hasta grado ocho, que es menor que \(-1/2\). Además, Cauchy–Schwarz da \(|u_\eta'(x)|\le\sqrt2\) en la recta real. Se deduce
\[
\operatorname{Lip}_zF_{t,\eta}\le6\sqrt2<9,\qquad
\operatorname{Lip}_tF_{t,\eta}\le1.
\]
Si \(d_j\) es la diferencia entre las dos órbitas, entonces \(d_0=0\) y \(d_{j+1}\le6\sqrt2\,d_j+|t-s|\). La inducción prueba la suma geométrica. Para \(|t-s|\le9^{-m}\), el diámetro es menor que \(9^{q-m}\). Sustituir \(q=2^N\) y \(m=q+r\) termina la prueba. \(\square\)

La cota controla la incertidumbre de parámetro con \(\eta\) fijado. Una implementación añade y propaga separadamente el error de evaluación de cosenos y de redondeo; el certificado racional antecedente ofrece esas operaciones. La cota es suficiente y conservadora, con coste que crece con el horizonte.

\textbf{Lector límite.} Realícese \(t\) desde \((x_0,y_0)=(1,0)\), \(B=9\), de modo que
\[
(x_m,y_m)=(9^m-C_m,C_m),\qquad t_m=C_m/9^m.
\]
Los lectores \(E_{q,m}=F_{t_m,\eta}^{q}(1/2)\), llevados a un nivel común, cumplen
\[
\|E_{q,m+a}-r^*E_{q,m}\|_\infty\le S_q9^{-m}.
\]
Convergen uniformemente a un lector continuo \(E_q\) en la completación del álgebra cilíndrica. La versión conjuntista
\[
J_{q,m}(C)=\{F_{t,\eta}^{q}(1/2):t\in I_m(C)\}
\]
es exactamente anidada y sus diámetros tienden a cero. Este resultado da el paso al límite de la \textbf{lectura a horizonte fijo}, con control de profundidad a cualquier horizonte solicitado.


\subsection{Palabra de duplicación y lector nonádico de la cronología}
\label{hmtfg:prolongacion:4}

Los ocho itinerarios certificados en el antecedente obedecen
\[
W_1=+,\qquad W_{n+1}=W_n\,(-1)^n\,W_n.
\]
Su continuación prefijal única es
\[
\tau_j=(-1)^{v_2(j)},\qquad j\ge1;
\qquad \tau_{2j}=-\tau_j,\quad\tau_{2j+1}=+.
\]
La prueba utiliza \(v_2(2^n+r)=v_2(r)\) para \(0<r<2^n\). Es una recurrencia a toda profundidad, y la coincidencia con los ocho retornos calculados constituye su comprobación finita en la familia.

En el orden lexicográfico con paridad de un mapa con máximo crítico, el factor de orientación previo al lugar \(k\) es \(\prod_{j<k}(-\tau_j)\). La palabra \(\tau\) es estrictamente mayor que cualquiera de sus desplazamientos positivos. En efecto, el primer desacuerdo con un desplazamiento impar ocurre en \(k=1\) o \(2\); un desplazamiento par reduce la comparación a la mitad y duplica \(k\). Por tanto el primer desacuerdo está en \(k=2^r\). Allí
\[
\prod_{j<2^r}(-\tau_j)=(-1)^r=\tau_{2^r},
\]
y la diferencia orientada con el signo opuesto es positiva. Esta demostración también prueba aperiodicidad.

El registro cronológico conserva
\[
K=\rho_9(K)+9q_9(K).
\]
Su lector diádico es
\[
\ell_n(B_K)=[\rho_9(K)+9q_9(K)]_{2^n}=[K]_{2^n}.
\]
Sobre la carta de historias en que \(\Gamma_9 B_K=B_{K+9}\),
\[
\ell_n(\Gamma_9B_K)=\ell_n(B_K)+9\pmod{2^n},\qquad
\pi_{n+1,n}\ell_{n+1}=\ell_n.
\]
Como nueve es impar, el avance por nueve visita las \(2^n\) posiciones. Su conjugación con el avance unitario es \(j\mapsto9j\). Para \(0<j<2^n\),
\[
v_2(9j\bmod2^n)=v_2(j),
\]
de modo que el remuestreo nonádico conserva también la palabra orientada en cada ciclo certificado.

Este lector utiliza residuo \textbf{y} cociente. Las fases de \(K=1\) y \(K=10\) coinciden módulo nueve, mientras sus signos diádicos son opuestos. La memoria distingue las dos historias. La propiedad aritmética del remuestreo vale para todo multiplicador impar; el calendario HMT fija aquí la elección de nueve.


\subsection{Memoria bilateral compatible a profundidad infinita}
\label{hmtfg:prolongacion:5}

Sean \(\mathcal H_n=\ell^2(\mathbb Z/2^n\mathbb Z)\),
\[
C_ne_j=e_{j+1},\qquad
J_ne_j=(e_j+e_{j+2^n})/\sqrt2.
\]
Una comprobación sobre la base da \(J_n^*J_n=I\) y \(C_{n+1}J_n=J_nC_n\), incluidos los términos de borde.

El artículo XI aporta los proyectores
\[
P_0=\frac19\begin{pmatrix}8&-\sqrt8\\-\sqrt8&1\end{pmatrix},\qquad
P_1=\frac19\begin{pmatrix}1&\sqrt8\\\sqrt8&8\end{pmatrix}.
\]
Son ortogonales y complementarios. Por tanto
\[
\mathbb U_n=P_0\otimes I+P_1\otimes C_n
=\begin{pmatrix}T_n&D_n\\D_n&R_n\end{pmatrix}
\]
es unitario, con
\[
T_n=(8I+C_n)/9,\quad D_n=\sqrt8(C_n-I)/9,\quad R_n=(I+8C_n)/9.
\]
La naturalidad ya probada implica
\[
\mathbb U_{n+1}(I_2\otimes J_n)=(I_2\otimes J_n)\mathbb U_n,
\quad \mathbb U_n^a=P_0\otimes I+P_1\otimes C_n^a
\quad(a\in\mathbb Z).
\]
El límite de Hilbert se identifica con \(L^2(\mathbb Z_2,\mu_{\rm Haar})\), llevando \(e_j^{(n)}\) a \(2^{n/2}\mathbf1_{j+2^n\mathbb Z_2}\). Las identidades se prolongan por densidad y acotación. Para todo vector \(v\) y todo tiempo entero \(a\),
\[
\boxed{\|T_{\infty,a}v\|^2+\|D_{\infty,a}v\|^2=\|v\|^2,}
\]
\[
\boxed{v=T_{\infty,a}^*(T_{\infty,a}v)+D_{\infty,a}^*(D_{\infty,a}v).}
\]
Aquí \(T_{\infty,a}=(8I+C_\infty^a)/9\). La composición bilateral conserva el archivo entre pasos; componer solamente \(T_n\) describe otra operación.

La permutación \(S_ne_j=e_{9j\bmod2^n}\) satisface
\[
S_{n+1}J_n=J_nS_n,\qquad S_nC_nS_n^{-1}=C_n^9.
\]
Así se obtiene en el límite
\[
S_\infty C_\infty S_\infty^{-1}=C_\infty^9,
\quad (I_2\otimes S_\infty)\mathbb U_\infty
(I_2\otimes S_\infty^{-1})=\mathbb U_\infty^9.
\]
La representación temporal es fiel: \(C_\infty^a=C_\infty^b\) fuerza \(2^n\mid a-b\) para todo \(n\), de donde \(a=b\). Preparaciones particulares pueden conservar simetrías temporales. La identidad operatoria conserva esta distinción.

El punto proyectivo \(([K]_{2^n})_n\) y un vector de \(L^2\) tienen tipos diferentes: una historia puntual induce una Dirac, mientras las amplitudes normalizables pertenecen al espacio de Hilbert. El resultado construye un lector equivarante de la cronología; las restantes coordenadas nativas del TPK acompañan el registro completo.


\subsection{Existencia y separación de las obligaciones asintóticas}
\label{hmtfg:prolongacion:6}

\subsubsection{Realización de la palabra infinita en toda la franja}
\label{hmtfg:prolongacion:6.1}

\textbf{Teorema de existencia.} Para cada \(|\eta|\le1/40\) existe al menos un
\[
\lambda_\infty(\eta)\in
\left[\frac1{2u_\eta(1)},\frac2{u_\eta(1)}\right]
\]
tal que
\[
\boxed{\operatorname{sign}f_{\lambda_\infty(\eta),\eta}^{j}(0)
=(-1)^{v_2(j)}\quad\text{para todo }j\ge1.}
\]
La conclusión realiza el itinerario infinito de duplicación dentro de la familia fijada. El cuantificador de existencia se mantiene separado de una selección única y de la ley de distancias entre parámetros superestables.

\textbf{Demostración.} Fijamos \(\eta\). Escribimos \(K_\lambda\) para la palabra de signos de la órbita crítica; cuando ésta retorna por primera vez a cero, la palabra termina en ese cero. Usamos el orden unimodal con máximo, cuyo factor antes de comparar el símbolo \(j\) es el producto de los signos \(-K_i\), \(i<j\).

En el extremo izquierdo la imagen de todo el intervalo está en \([1/2,1]\); por tanto \(K_-=+++\cdots\). En el extremo derecho la órbita es \(0\mapsto1\mapsto-1\mapsto-1\), de donde \(K_+=+---\cdots\). Comparando los símbolos segundo y tercero, respectivamente, se obtiene
\[
K_-<\tau<K_+.
\]

Si \(K_{\lambda_0}\) es infinita y distinta de \(\tau\), la primera diferencia ocurre en un lugar finito. La continuidad de los correspondientes iterados hace localmente constante el lado de la comparación. Queda examinar un parámetro superestable, de primer período \(p\), con palabra \(W0\), donde \(|W|=p-1\).

Para \(\lambda\) cercano a \(\lambda_0\), sea \(g_\lambda=f_{\lambda,\eta}^p\) y \(a=g_\lambda(0)\). Se tiene \(g_\lambda'(0)=0\), y una cota local uniforme de la segunda derivada da
\[
g_\lambda(x)=a+O(x^2).
\]
Para \(|a|\) suficientemente pequeño y distinto de cero, \(g_\lambda\) lleva \([-2|a|,2|a|]\) en \([a-|a|/2,a+|a|/2]\). Los retornos conservan el signo de \(a\), y las \(p-1\) posiciones intermedias conservan \(W\). Las palabras vecinas son, por ello, \((W+)^\infty\) o \((W-)^\infty\); los parámetros con \(a=0\) conservan \(W0\).

Si \(\tau\) difiere de \(W\) antes del lugar \(p\), esa diferencia fija el mismo lado para las tres palabras. Si comparte \(W\), pongamos
\[
s=\tau_p,\qquad e=\prod_{j<p}(-\tau_j),\qquad A=Ws.
\]
El signo de la comparación de \(\tau\) con \(W0\) y con \((W,-s)^\infty\) es \(es\). Para compararla con \(A^\infty\), sea \(q\) el primer lugar en que \(\tau_{p+q}\ne\tau_q\). Existe por aperiodicidad. Hasta el lugar \(p+q-1\), ambas palabras coinciden, y el signo de orientación se factoriza como \(O_pO_{q-1}\), con \(O_p=-es\). La maximalidad estricta ya demostrada da
\[
O_{q-1}(\tau_q-\tau_{p+q})>0.
\]
En consecuencia, la comparación de \(\tau\) con \(A^\infty\) tiene signo \(-O_p=es\). Las tres posibilidades locales están al mismo lado de \(\tau\).

Si ningún parámetro realizara \(\tau\), los conjuntos \(\{\lambda:K_\lambda<\tau\}\) y \(\{\lambda:K_\lambda>\tau\}\) serían abiertos, disjuntos, no vacíos y cubrirían el intervalo de parámetros. La conexidad del intervalo excluye esa partición. Existe, por tanto, el parámetro anunciado. \(\square\)

Este argumento de continuidad de itinerarios se aplica después de fijar el generador HMT y su familia analítica; sus hipótesis se verifican aquí directamente. La construcción conserva la palabra infinita sin introducir como dato un parámetro de la familia cuadrática ni el valor universal.

\subsubsection{Márgenes que impiden perder el itinerario al tomar límites}
\label{hmtfg:prolongacion:6.2}

Escribamos \(c_p=f_{\lambda,\eta}^p(0)\). Una cota racional uniforme suficiente es
\[
|f'|,|f''|\le16,\qquad
B_p:=16^p\frac{16^p-1}{15}\ge\|(f^p)''\|_{[-1,1]}.
\]
En efecto, \(1-\cos v\ge v^2/2-v^4/24\ge v^2/8\) para \(|v|\le3\), de modo que \(u_\eta(1)\ge1/4\), \(\lambda\le8\) y ambas derivadas se acotan por \(16\), usando los momentos. La regla de la cadena prueba por inducción la cota \(B_p\).

La palabra \(\tau\) satisface \(\tau_{2p}=-\tau_p\) para todo \(p\). Para cualquiera de sus realizaciones, \(c_{2p}\) y \(c_p\) tienen signos opuestos. Como \((f^p)'(0)=0\), Taylor da
\[
|c_{2p}-c_p|\le\tfrac12B_p|c_p|^2.
\]
El lado izquierdo es mayor que \(|c_p|\), y por tanto
\[
\boxed{|c_p|>2/B_p>0.}
\]
Para cada horizonte fijo existe así un margen uniforme que impide que una sucesión de realizaciones pierda su signo al converger. En las coordenadas compactas \((\eta,b)\), \(b=\lambda u_\eta(1)\in[1/2,2]\), el conjunto de realizaciones de \(\tau\) es cerrado y compacto, y su proyección sobre toda la franja de \(\eta\) es sobreyectiva.

El mismo argumento sirve para un árbol de prefijos: si se realizan todos los prefijos de \(\tau\), una subsucesión compacta conserva cada signo, porque los prefijos suficientemente largos incluyen simultáneamente los lugares \(p\) y \(2p\). El teorema de existencia proporciona ahora la no vacuidad que ese argumento necesita.

\subsubsection{Intervalos invariantes de retorno a todos los niveles}
\label{hmtfg:prolongacion:6.3}

\textbf{Teorema de retornos anidados.} Para cualquiera de las realizaciones del teorema de existencia de la sección \ref{hmtfg:prolongacion:6.1}, sean \(c_j=f^j(0)\) y
\[
J_n=\operatorname{conv}\{c_{2^n},c_{2^{n+1}}\},\qquad n\ge1.
\]
Entonces \(0\in\operatorname{int}J_n\), \(J_{n+1}\subset J_n\),
\[
f^{2^n}(J_n)=J_n,
\]
y las \(2^n\) imágenes \(J_n,f(J_n),\ldots,f^{2^n-1}(J_n)\) son disjuntas. Los retornos normalizados son unimodales de máximo, con valor crítico uno y la misma palabra \(\tau\). Estos intervalos son núcleos invariantes de retorno; una normalización que exija extremos periódicos puede ampliar el núcleo al intervalo restrictivo maximal correspondiente y debe especificar ese paso.

\textbf{Demostración del primer retorno.} Sea \(F\) un mapa unimodal de máximo con \(F(0)=1\) y palabra \(\tau\), y \(d_j=F^j(0)\). Sus primeros signos son \(+,-,+,+,+,-,+\). El intervalo \(I=[d_2,1]\) es invariante: \(F(d_2)=d_3>0\), \(F(1)=d_2<0\), y el máximo es uno. Pongamos \(J=[d_2,d_4]\). Como \(F\) decrece en positivos y
\[
F(d_3)=d_4>0> d_6=F(d_5),
\]
se tiene \(d_3<d_5\), y por tanto \(F(J)=[d_3,1]\).

Además \(d_3>d_4\). Para comprobarlo, \(F^2\) es estrictamente creciente entre los dos puntos positivos \(d_3,d_4\), porque sus imágenes \(d_4,d_5\) son positivas. Sus valores son \(F^2(d_3)=d_5>0>d_6=F^2(d_4)\), lo que fuerza ese orden. Así \(J\) y \(F(J)\) quedan separados por un intervalo abierto, y
\[
F^2(J)=F([d_3,1])=[d_2,d_4]=J.
\]
El mapa \(F^2|_J\) tiene un único mínimo crítico: la imagen \(F(J)\) permanece positiva. La conjugación \(h(x)=d_2x\) invierte orientación y normaliza el retorno como máximo con valor uno. Su itinerario es
\[
\operatorname{sign}\frac{d_{2j}}{d_2}
=-\tau_{2j}=\tau_j.
\]
Se puede repetir el mismo argumento a cualquier profundidad.

\textbf{Disjunción inductiva.} Supónganse disjuntas las \(p=2^n\) imágenes del padre \(J_n\), y sea \(g=f^p\). El primer retorno aplicado al mapa normalizado produce \(B=J_{n+1}\), \(A=g(B)\), con \(A\cap B=\varnothing\), \(g(B)=A\), \(g(A)=B\). Para \(0\le j<p\), una intersección de \(f^j(A)\) y \(f^j(B)\), al aplicarle \(f^{p-j}\), produciría una intersección de \(B\) y \(A\). Por tanto esas dos imágenes son disjuntas. Las imágenes correspondientes a distintos \(j\) pertenecen a imágenes disjuntas del padre. Quedan demostradas las \(2p\) imágenes disjuntas y el retorno \(f^{2p}(B)=B\). La fórmula de los extremos resulta de componer las normalizaciones \(x\mapsto c_{2^n}x\). \(\square\)

Los conjuntos compactos
\[
\Lambda_n=\bigcup_{0\le j<2^n}f^j(J_n)
\]
son no vacíos y anidados. Su intersección \(\Lambda_\infty\) tiene un lector continuo sobre el odómetro diádico: a cada punto se le asigna la etiqueta de la única componente que ocupa a cada nivel. Todas las sucesiones compatibles de etiquetas tienen preimagen por compacidad, y el avance por \(f\) suma uno en las etiquetas. Esta aplicación es una semiconjugación sobreyectiva. Para promoverla a conjugación puntual se requiere demostrar que los diámetros de las componentes tienden a cero; esa promoción métrica se mantiene separada.

% Fragmento integrable. Preambulo anfitrion: amsmath, amssymb, longtable, hyperref.
\section{Compatibilidad de lectores y restitución del parámetro}
\label{hmtfg:fragmento:compatibilidad-de-lectores-y-restitucion-del-parametro}

\subsection{Compatibilidad de lectores, mínimos y recuperación de la historia}
\label{hmtfg:escalado:13}

Las operaciones de prolongación y los lectores ordenados conservan la información requerida por la selección de raíces. Las siguientes identidades fijan sus dominios y su composición.

\subsubsection{Regla forward recuperada}
\label{hmtfg:escalado:13.1}

La actualización determinista parte del estado de frontera R36 con carácter \(\chi\) y residuo interno. La actualización es
\begin{equation}
\mathcal U_K^\chi=
\operatorname{Upd}_K^\chi\circ
\operatorname{Tra}_K^\chi\circ
\operatorname{Sel}_K^\chi,
\qquad
N_{K+1}=729N_K+d_K^\chi.
\label{hmtfg:escalado:eq:26}
\end{equation}
El dígito y la actualización dependen del estado precedente. Dos secciones forward con el mismo estado completo R36 coinciden por inducción. La subrelación \(\operatorname{Ext}^{\chi,\rightarrow}\) es el grafo de esa actualización; la relación total \(\operatorname{Ext}\) admite multiplicidad de continuaciones según los datos y caracteres.

Las prolongaciones compatibles transportan los lectores regionales y la firma conjunta sobre los mismos prefijos.

Estas reglas transportan orientación, residuo, cociente, acarreo, frontera y memoria. Su aplicación a la cascada debe conservar además la selección ordenada del parámetro del lector crítico.

\subsubsection{Independencia de las constantes antecedentes respecto del parámetro libre}
\label{hmtfg:escalado:13.2}

Sea \(B\) el estado HMT que proporciona la rueda y sus constantes antecedentes. Escribamos \(\mathcal C(B)\) para el conjunto de esas publicaciones y \(u_B\) para el perfil crítico. En el dominio del lector crítico, la construcción posterior es
\begin{equation}
(B,\lambda)\longmapsto f_{B,\lambda}(x)=1-\lambda u_B(x),
\qquad 0<\lambda\le 2/u_B(1).
\label{hmtfg:escalado:eq:27}
\end{equation}
Para un mismo \(B\), la lectura de las constantes anteriores se factoriza por la proyección al primer componente:
\begin{equation}
\widetilde{\mathcal C}(B,\lambda)
=\mathcal C\circ\operatorname{pr}_B(B,\lambda)
=\mathcal C(B).
\label{hmtfg:escalado:eq:28}
\end{equation}
En consecuencia, dos valores distintos de \(\lambda\) conservan las mismas publicaciones antecedentes. La igualdad es una propiedad de esta construcción causal: el perfil se fija antes de variar \(\lambda\), y el factor \(\pi_{\mathrm{HMT}}\) se cancela en su normalización cuadrática.

\textbf{Lema.} Supóngase que una condición \(A\) se evalúa únicamente sobre esas publicaciones antecedentes. Para todo \(B\) fijo,
\begin{equation}
\{\lambda\in\Lambda:A(\widetilde{\mathcal C}(B,\lambda))\}
=
\begin{cases}
\Lambda,&A(\mathcal C(B))\text{ se cumple},\\
\varnothing,&A(\mathcal C(B))\text{ falla}.
\end{cases}
\label{hmtfg:escalado:eq:29}
\end{equation}
\textbf{Demostración.} Sustituir \eqref{hmtfg:escalado:eq:28} hace independiente de \(\lambda\) el predicado. Por tanto, lo satisfacen todos los parámetros del dominio o ninguno. Esto incluye una comparación de las constantes antecedentes con intervalos metrológicos.

La unicidad de la sección \(B\) y la restricción posterior \(f_{B,\lambda}^{2^n}(0)=0\) producen una fibra de soluciones en \(\lambda\). La regla de primera raíz actúa sobre esta fibra mediante el retorno y su itinerario. La selección y su prolongación se demuestran en los resultados siguientes.

Este resultado mantiene el orden HMT: la metrología contrasta salidas previamente construidas y la selección de la raíz se demuestra en el lector que la produce.

\subsubsection{Lema de transporte ordenado del selector}
\label{hmtfg:escalado:13.3}

Sea \(P_n\) un conjunto ordenado de candidatos genealógicos con mínimo \(p_n^*\). Sea \(Z_n\) el conjunto original de raíces candidatas del corpus, con el orden real, y sea \(L_n:P_n\to Z_n\) un lector monótono. Se dice que su imagen es \textbf{coinicial} cuando
\begin{equation}
\forall z\in Z_n\quad\exists p\in P_n:
L_n(p)\le z.
\label{hmtfg:escalado:eq:30}
\end{equation}
Entonces
\begin{equation}
\boxed{L_n(p_n^*)=\min Z_n
\quad\Longleftrightarrow\quad
L_n(P_n)\text{ es coinicial en }Z_n.}
\label{hmtfg:escalado:eq:31}
\end{equation}
\textbf{Demostración.} Si la imagen del mínimo es el mínimo global, se toma \(p=p_n^*\) en \eqref{hmtfg:escalado:eq:30}. Recíprocamente, dado \(z\in Z_n\), la coinicialidad proporciona \(p\) con \(L_n(p)\le z\). La monotonía da \(L_n(p_n^*)\le L_n(p)\le z\). Como \(L_n(p_n^*)\in Z_n\), es el mínimo global. La sobreyectividad de \(L_n\) constituye una condición suficiente, más fuerte que \eqref{hmtfg:escalado:eq:30}.

El isomorfismo ordenado de la carta del continuo satisface la sobreyectividad del lector. La aplicación concreta al retorno y a su conjunto completo de ceros se demuestra en la sección de transporte ordenado; la clasificación y el aislamiento determinan después su continuación dinámica.

\subsubsection{Selección mínima mediante supervivencia a profundidad arbitraria}
\label{hmtfg:escalado:13.6}

La composición con la supervivencia produce además un selector global preciso. Se mantiene la familia uniforme \(\eta=0\), construida en sección~\ref{hmtfg:escalado:1}, y se utilizan los teoremas de realización del itinerario infinito y conservación de sus signos. Sean
\[
I=\left[\frac1{2u_0(1)},\frac2{u_0(1)}\right],\qquad
c_j(\lambda)=f_\lambda^j(0),\qquad
\tau_j=(-1)^{v_2(j)}.
\]
El conjunto
\[
\Lambda_\tau=\{\lambda\in I:\operatorname{sign}c_j(\lambda)=\tau_j
\text{ para todo }j\ge1\}
\]
es compacto y no vacío. El control de derivadas de ese propietario proporciona
\begin{equation}
B_j=\frac{16^j(16^j-1)}{15},\qquad
|c_j(\lambda)|>\frac2{B_j}\quad(\lambda\in\Lambda_\tau).
\label{hmtfg:escalado:eq:32}
\end{equation}
Definamos
\begin{equation}
\varepsilon_j=\frac1{B_j},\qquad
A_N=\{\lambda\in I:\tau_jc_j(\lambda)\ge\varepsilon_j,\
1\le j\le N\}.
\label{hmtfg:escalado:eq:33}
\end{equation}
Cada \(A_N\) es compacto por continuidad de los iterados y no vacío porque contiene \(\Lambda_\tau\). Además,
\begin{equation}
A_{N+1}\subseteq A_N,\qquad
\bigcap_{N\ge1}A_N=\Lambda_\tau.
\label{hmtfg:escalado:eq:34}
\end{equation}
La inclusión de derecha a izquierda se sigue de \eqref{hmtfg:escalado:eq:32}. La inclusión inversa se sigue de los márgenes estrictamente positivos de \eqref{hmtfg:escalado:eq:33}, que fijan cada signo.

\textbf{Teorema del mínimo superviviente.} Los mínimos \(a_N=\min A_N\) satisfacen
\begin{equation}
\boxed{a_N\uparrow a_*=\min\Lambda_\tau.}
\label{hmtfg:escalado:eq:35}
\end{equation}
\textbf{Demostración.} La inclusión \eqref{hmtfg:escalado:eq:34} hace creciente la sucesión de mínimos, acotada dentro de \(I\); tiene, por tanto, un límite \(a\). Para cada \(r\), todos los términos \(a_N\), \(N\ge r\), pertenecen al cerrado \(A_r\); así \(a\in A_r\). Por \eqref{hmtfg:escalado:eq:34}, \(a\in\Lambda_\tau\). Para cualquier \(\lambda\in\Lambda_\tau\), se tiene \(a_N\le\lambda\) a todo nivel, y al pasar al límite, \(a\le\lambda\). Esto prueba \eqref{hmtfg:escalado:eq:35}.

En el lector cilíndrico HMT, el objeto \(a_*\) posee prefijos compatibles a toda profundidad, con orientación, residuo y memoria conservados. Su elección utiliza el orden de la lectura paramétrica y la supervivencia completa. La compacidad demuestra la existencia y convergencia de estos mínimos.

Los parámetros \(a_N\) resuelven restricciones de signos con margen. Sus fronteras pueden satisfacer \(\tau_jc_j=\varepsilon_j\), mientras las raíces superestables satisfacen \(c_{2^n}=0\). Por ello \eqref{hmtfg:escalado:eq:35} establece un selector global de realización infinita y preserva, como objeto distinto, el selector de primeras raíces del corpus.

\subsubsection{Recuperación del parámetro desde la historia crítica completa}
\label{hmtfg:escalado:13.7}

El segundo término de la historia crítica proporciona un lector inverso exacto:
\begin{equation}
c_1(\lambda)=1,\qquad
c_2(\lambda)=1-\lambda u_0(1),\qquad
\boxed{\lambda=\frac{1-c_2(\lambda)}{u_0(1)}.}
\label{hmtfg:escalado:eq:36}
\end{equation}
La positividad \(u_0(1)>0\), demostrada para el perfil, hace válida la división. En consecuencia, dos historias críticas completas iguales tienen el mismo parámetro. La palabra de signos \(\tau\) constituye una publicación reducida de esa historia; conserva su combinatoria y omite la magnitud de \(c_2\).

Las fórmulas \eqref{hmtfg:escalado:eq:35}–\eqref{hmtfg:escalado:eq:36} unen la selección por supervivencia y la restitución del parámetro desde la historia crítica. La igualdad de su límite con el cruce estable se demuestra por primera salida, y la identificación de los centros por el transporte analítico de sus clases híbridas.

% Fragmento integrable. Preambulo anfitrion: amsmath, amssymb, longtable, hyperref.
\section{Transporte ordenado y representación de Jacobi}
\label{hmtfg:fragmento:transporte-ordenado-y-representacion-de-jacobi}

\subsection{Transporte completo desde el continuo y realización espectral del perfil}
\label{hmtfg:escalado:14}

El transporte ordenado del continuo y la representación espectral del lector se componen antes de la selección paramétrica.

\subsubsection{Cobertura de todas las raíces y transporte de su mínimo}
\label{hmtfg:escalado:14.1}

La construcción ordenada del continuo desarrollada en los antecedentes produce la publicación arquimediana sobreyectiva desde cilindros compatibles y su cociente ordenado. Dentro de su universo declarado \(\mathfrak G\), demuestra el isomorfismo de cuerpos ordenados completos
\begin{equation}
\iota:\mathbb R_{\rm HMT}^{\mathfrak G}
   \xrightarrow{\cong}\mathbb R^{\mathfrak G}.
\label{hmtfg:escalado:eq:37}
\end{equation}
Se cambia aquí el nombre \(\eta\) del propietario por \(\iota\), para distinguirlo del parámetro de ponderación, fijado a cero. La segunda proyección conserva incidencia, frontera y memoria sobre el mismo antecedente; el cociente de evaluación y la historia completa mantienen sus dominios respectivos.

Construyamos en el cuerpo HMT, antes de seleccionar una raíz,
\[
\cos_{\rm H}(t)=\sum_{k=0}^{\infty}\frac{(-1)^kt^{2k}}{(2k)!},
\qquad
u_{\rm H}(x)=\frac1{12}\sum_{m=1}^{12}
\left[1-\cos_{\rm H}\!\left(\frac{2(3m-1)x}{\sqrt[\rm H]{899}}\right)\right],
\]
\begin{equation}
F_a(x)=1-a\,u_{\rm H}(x),\qquad S_n^{\rm H}(a)=F_a^{\,2^n}(0).
\label{hmtfg:escalado:eq:38}
\end{equation}
Los enteros y el segundo momento son los ya derivados del lector dodecafásico. La raíz cuadrada es la positiva del cuerpo ordenado. La serie convergente realiza analíticamente ese lector; conserva el perfil y el coeficiente normalizador de sección~\ref{hmtfg:escalado:1}.

\textbf{Teorema de transporte del retorno.} Para todo parámetro del dominio interno de \eqref{hmtfg:escalado:eq:38},
\begin{equation}
\iota(F_a(x))=f_{\iota(a)}(\iota(x)),\qquad
\iota(S_n^{\rm H}(a))=S_n(\iota(a)).
\label{hmtfg:escalado:eq:39}
\end{equation}
\textbf{Demostración.} El isomorfismo ordenado fija los racionales y preserva la raíz positiva. Preserva los límites convergentes: cada entorno de radio racional positivo se transporta a otro del mismo radio. Aplicado a las sumas parciales de \eqref{hmtfg:escalado:eq:38}, transporta el coseno y \(u_{\rm H}\). La primera igualdad resulta de conservar suma y producto; la segunda, por inducción sobre el número de iteraciones.

Fijado un centro anterior \(a_{n-1}\), definamos \(Z_n^{\rm H}\) por retorno cero, período crítico exacto \(2^n\) y \(a>a_{n-1}\), dentro del dominio paramétrico declarado. Sea \(Z_n\) el conjunto definido por las mismas condiciones para \(f_\lambda\), a la derecha de \(\iota(a_{n-1})\). Entonces
\begin{equation}
\boxed{\iota[Z_n^{\rm H}]=Z_n,\qquad
\iota(\min Z_n^{\rm H})=\min Z_n.}
\label{hmtfg:escalado:eq:40}
\end{equation}
La segunda igualdad se afirma cuando existe el mínimo; la primera preserva también la no existencia de candidatos.

\textbf{Demostración.} \eqref{hmtfg:escalado:eq:39} preserva el retorno cero y los retornos anteriores que determinan el período exacto. El orden conserva la posición respecto del centro previo. La sobreyectividad de \eqref{hmtfg:escalado:eq:37}, aplicada a cualquier parámetro raíz, prueba la cobertura inversa. Si \(a_*\) es el mínimo, cualquier \(z\in Z_n\) tiene la forma \(\iota(a)\), con \(a\in Z_n^{\rm H}\), de donde \(\iota(a_*)\le z\). La afirmación inversa utiliza \(\iota^{-1}\).

La sobreyectividad cubre todos los parámetros de la carta, incluidos aquellos cuyas raíces no se hayan aislado por un cálculo finito. El isomorfismo conserva el universo de interpretación declarado y el orden de selección.

\subsubsection{Representación de Jacobi del perfil dodecafásico completo}
\label{hmtfg:escalado:14.2}

Definamos la medida atómica del lector,
\begin{equation}
s_m=\frac{4(3m-1)^2}{899},\qquad
\nu_0=\frac1{12}\sum_{m=1}^{12}\delta_{s_m},\qquad
M_r=\int s^r\,d\nu_0(s).
\label{hmtfg:escalado:eq:41}
\end{equation}
Los doce nodos son distintos y positivos. Para un polinomio \(p\ne0\) de grado menor que \(r\le12\),
\begin{equation}
\sum_{i,j=0}^{r-1}p_iM_{i+j}p_j
=\frac1{12}\sum_{m=1}^{12}p(s_m)^2>0.
\label{hmtfg:escalado:eq:42}
\end{equation}
Un polinomio de grado menor que doce no se anula en los doce nodos. En grado doce aparece \(\prod_m(s-s_m)\), de norma nula.

La ortogonalización respecto de esta medida produce polinomios ortonormales \(\psi_0,\ldots,\psi_{11}\), con \(\psi_0=1\). Sean
\[
Q_{mj}=\psi_j(s_m)/\sqrt{12},\quad
D=\operatorname{diag}(s_m),\quad J_{12}=Q^{\mathsf T}DQ.
\]
Se cumplen
\begin{equation}
Q^{\mathsf T}Q=I,\quad DQ=QJ_{12},\qquad
\boxed{u_0(x)=e_0^{\mathsf T}
\bigl[I-\cos(x\sqrt{J_{12}})\bigr]e_0.}
\label{hmtfg:escalado:eq:43}
\end{equation}
\textbf{Demostración.} La ortogonalidad procede de \eqref{hmtfg:escalado:eq:42}. La multiplicación por \(s\) produce una recurrencia de tres términos: los elementos separados por más de una posición se anulan por ortogonalidad y grado. Fijar positivos los coeficientes principales hace positivas las subdiagonales. Por cálculo funcional,
\(\cos(x\sqrt{J_{12}})=Q^{\mathsf T}\cos(x\sqrt D)Q\).
El vector \(Qe_0\) es constante, de valor \(1/\sqrt{12}\); sustituirlo recupera la suma de sección~\ref{hmtfg:escalado:1}.

El artículo de Grassmannianos desarrolla este mecanismo para medidas marcadas y refinamientos. Aquí se aplica a \eqref{hmtfg:escalado:eq:41}; el rango doce pertenece a este lector. Las marcas genealógicas se conservan junto al lector, conforme a la fidelidad marcada del propietario. \(J_{12}\) representa el perfil analítico; el estado completo conserva adicionalmente rutas, orientación y memoria.

El verificador de esta ampliación comprueba con fracciones exactas los doce grados de norma positiva, la anulación del grado doce, el entrelazador \(VT=DV\) en base mónica y la autoadjunción ponderada \(HT=T^{\mathsf T}H\). Comprueba los momentos de grado cero a treinta; la identidad para todo grado se sigue del entrelazador. Sus primeros coeficientes son
\begin{equation}
a_0=2,\qquad \beta_1=\frac{2493348}{808201}.
\label{hmtfg:escalado:eq:44}
\end{equation}

\subsubsection{Sensibilidad orientada del retorno}
\label{hmtfg:escalado:14.3}

Para un retorno exacto de período \(p=2^n\), sean
\[
c_j=f_\lambda^j(0),\quad q_j=\partial_\lambda c_j,\quad
D_j=\prod_{k=1}^j f_\lambda'(c_k),\quad D_0=1.
\]
Antes del retorno crítico, los factores de \(D_{p-1}\) son distintos de cero. Diferenciación y telescopaje dan
\begin{equation}
q_{j+1}=-u_0(c_j)+f_\lambda'(c_j)q_j,\qquad
\boxed{\frac{q_p}{D_{p-1}}
=-\sum_{j=1}^{p-1}\frac{u_0(c_j)}{D_j}.}
\label{hmtfg:escalado:eq:45}
\end{equation}
En la palabra canónica, \(\operatorname{sign}c_j=(-1)^{v_2(j)}\). El signo de \(f_\lambda'(c_j)\) es el opuesto; por tanto,
\begin{equation}
\operatorname{sign}D_j=(-1)^{j+\sum_{k=1}^jv_2(k)}
=(-1)^{s_2(j)},
\label{hmtfg:escalado:eq:46}
\end{equation}
usando \(\sum_{k=1}^jv_2(k)=j-s_2(j)\), donde \(s_2(j)\) cuenta los unos binarios. La sensibilidad normalizada es
\begin{equation}
\mathcal T_n=-\sum_{j=1}^{2^n-1}(-1)^{s_2(j)}t_j,\qquad
t_j=\frac{u_0(c_j)}{|D_j|}>0.
\label{hmtfg:escalado:eq:47}
\end{equation}
La positividad nodal de \eqref{hmtfg:escalado:eq:42} determina \(u_0\); \eqref{hmtfg:escalado:eq:47} conserva además los productos orbitales y sus orientaciones.

La representación adicional por momentos exigiría una medida positiva \(\rho\) en \([0,1]\) con \(t_j=\int z^{j-1}\,d\rho\). Bajo esa condición,
\begin{equation}
\mathcal T_n=\int_0^1
\frac{1-\prod_{r=0}^{n-1}(1-z^{2^r})}{z}\,d\rho(z)>0
\label{hmtfg:escalado:eq:48}
\end{equation}
para \(\rho\ne0\), extendiendo el integrando por su valor \(1\) en cero. La identidad se obtiene expandiendo el producto binario.

El control racional en la raíz de período cuatro previamente certificada obtiene
\[
t_1\approx0.40425224267600139730,\quad
t_2\approx0.04925249167432646403,\quad
t_3\approx0.15740870098294833737
\]
y el intervalo exterior riguroso
\begin{equation}
\begin{gathered}
0.296096033367379523967764444238832345360107<\mathcal T_2,\\
\mathcal T_2<0.296096033367379523967764444238832345360112.
\end{gathered}
\label{hmtfg:escalado:eq:49}
\end{equation}
También certifica \(t_3>t_2\). Una medida positiva sobre \([0,1]\) impondría \(t_3\le t_2\); ese levantamiento concreto de los pesos orbitales a momentos positivos queda descartado. El signo positivo de \eqref{hmtfg:escalado:eq:49} permanece certificado. Este control distingue el fallo de una prueba propuesta del fallo de la propiedad que se pretendía probar.

% Fragmento integrable. Preambulo anfitrion: amsmath, amssymb, longtable, hyperref.
\section{Clasificación y continuación del selector de primeras raíces}
\label{hmtfg:fragmento:clasificacion-y-continuacion-del-selector-de-primeras-raices}

\subsection{Objeto y procedencia}
\label{hmtfg:clasificacion:1}

Este apéndice conserva la familia publicada por el lector dodecafásico HMT:

\[
u_0(x)=\frac1{12}\sum_{m=1}^{12}\left(1-\cos\frac{2(3m-1)x}{\sqrt{899}}\right),
\qquad f_\lambda(x)=1-\lambda u_0(x).
\]

La familia procede de los antecedentes operatorios y analíticos expuestos en este desarrollo. La existencia y la compacidad de sus realizaciones del itinerario infinito se han demostrado en las subsecciones de realización de la palabra y conservación de sus signos.

El argumento siguiente es propio y utiliza exclusivamente esas propiedades y las operaciones de retorno de la familia fijada. No utiliza como entrada un parámetro de otra familia ni el valor de una constante universal. Las hojas de los itinerarios aquí utilizados son las dos ramas monótonas del lector crítico; las restantes coordenadas y la memoria HMT conservan su residencia en el estado del que procede el lector.

Escribimos
\[
c_j(\lambda)=f_\lambda^j(0),\qquad
\tau_j=(-1)^{v_2(j)},\qquad
W_n=\tau_1\cdots\tau_{2^n-1}.
\]
El itinerario de una órbita crítica que retorna por primera vez al origen termina en el símbolo cero. El orden de itinerarios es el orden unimodal de máximo: antes de comparar el símbolo de índice \(j\), el factor de orientación es
\[
O_{j-1}(K)=\prod_{i=1}^{j-1}(-K_i).
\]
El orden ordinario de los símbolos es \(-1<0<1\).


\subsection{Clasificación anterior a la palabra infinita}
\label{hmtfg:clasificacion:2}

\textbf{Teorema de clasificación.} Sea \(F:I\to I\), con \(I=[\ell,1]\), \(\ell<0\), un mapa continuo, estrictamente creciente a la izquierda de cero y estrictamente decreciente a su derecha, cuyo máximo satisface \(F(0)=1\). Supóngase que cero tiene período exacto \(2^n\), con \(n\ge1\), y que su itinerario crítico \(K\) satisface \(K<\tau\). Entonces
\[
\boxed{K=W_n0.}
\]
La conclusión se refiere al itinerario, y permite que parámetros distintos de una familia tengan ese mismo itinerario.

\textbf{Demostración.} Denotemos \(d_j=F^j(0)\). Para período dos el itinerario es \(+0=W_10\). Consideremos un período exacto \(p=2^n\ge4\).

Se tiene \(d_2<0\). En efecto, \(d_2=0\) daría período dos. Si \(d_2>0\), la monotonía de la rama positiva da
\[
F([d_2,1])=[d_2,d_3]\subset[d_2,1],
\]
y la órbita crítica permanece en un intervalo estrictamente positivo, incompatible con el retorno a cero.

Los signos iniciales son entonces \(+,-\). El factor de orientación para comparar el tercer signo es \(-1\); por ello \(K<\tau\) obliga a \(d_3>0\). Un cero en ese lugar produciría período tres.

Tampoco puede ocurrir \(d_4<0\). Bajo esta hipótesis,
\[
d_2<d_4<0,
\qquad F([d_2,d_4])=[d_3,d_5]\subset(0,1],
\]
pues \(d_4=F(d_3)>F(1)=d_2\) y \(d_5=F(d_4)>F(d_2)=d_3\). La rama positiva es decreciente, de modo que
\[
F^2([d_2,d_4])=[d_6,d_4]\subset[d_2,d_4].
\]
Aquí \(d_6\ge d_2\), porque \(d_5\le1\), y \(d_6<d_4\), porque \(d_5>d_3\). Este intervalo estrictamente negativo atrapa los retornos pares de la órbita crítica, y excluye el retorno a cero. Por tanto \(d_4\ge0\). Si \(p=4\), se obtiene precisamente \(+-+0=W_20\).

Supongamos ahora \(p\ge8\). Entonces \(d_4>0\). Los factores de orientación para los signos quinto y sexto son, respectivamente, \(-1\) y \(+1\). Como \(\tau_5=+1\) y \(\tau_6=-1\), la desigualdad \(K<\tau\) obliga a
\begin{equation}
\operatorname{sign}(d_1,\ldots,d_6)=(+,-,+,+,+,-).
\label{hmtfg:clasificacion:eq:1}
\end{equation}
Los signos cero en esos lugares están excluidos por el período exacto.

Construimos el intervalo de retorno
\[
J=[d_2,d_4].
\]
La desigualdad \(F(d_3)=d_4>0>d_6=F(d_5)\), junto con \(d_3,d_5>0\), implica \(d_3<d_5\); por tanto
\[
F(J)=[d_3,1].
\]
Además \(d_4<d_3\). Para comprobarlo, la imagen por \(F\) del intervalo de extremos \(d_3,d_4\) es estrictamente positiva, pues sus imágenes extremas son \(d_4,d_5>0\). En ese intervalo \(F^2\) es estrictamente creciente. Sus valores extremos satisfacen
\[
F^2(d_3)=d_5>0>d_6=F^2(d_4),
\]
lo que fuerza \(d_3>d_4\). En consecuencia,
\begin{equation}
J\cap F(J)=\varnothing,
\qquad F^2(J)=F([d_3,1])=[d_2,d_4]=J.
\label{hmtfg:clasificacion:eq:2}
\end{equation}

El mapa \(F^2|_J\) posee un único mínimo en cero. La conjugación por \(h(x)=d_2x\), que invierte la orientación, define
\begin{equation}
G=h^{-1}\circ F^2\circ h:
\left[\frac{d_4}{d_2},1\right]\longrightarrow
\left[\frac{d_4}{d_2},1\right].
\label{hmtfg:clasificacion:eq:3}
\end{equation}
Este mapa satisface las mismas hipótesis de máximo que \(F\), con \(G(0)=1\), y su crítico tiene período exacto \(p/2\). Su itinerario es
\begin{equation}
K'_j=-K_{2j}.
\label{hmtfg:clasificacion:eq:4}
\end{equation}
Todos los símbolos impares de \(K\) son positivos, porque las posiciones pares pertenecen a \(J\) y las impares a \(F(J)\subset(0,1]\).

La transformación \eqref{hmtfg:clasificacion:eq:4} conserva el orden de comparación con \(\tau\). En efecto, \(\tau_{2j-1}=+1\) y \(\tau_{2j}=-\tau_j\). Si \(q\) es el primer lugar en que \(K'\) y \(\tau\) difieren, el primer lugar de diferencia entre \(K\) y \(\tau\) es \(2q\), y
\begin{equation}
O_{2q-1}(K)=-O_{q-1}(K'),
\qquad K_{2q}-\tau_{2q}=-(K'_q-\tau_q).
\label{hmtfg:clasificacion:eq:5}
\end{equation}
El producto que determina el orden es idéntico. Las identidades siguen siendo válidas cuando el primer símbolo distinto es el cero final. Así \(K'<\tau\).

La inducción aplicada a \(G\) da \(K'=W_{n-1}0\). Las posiciones impares positivas y \eqref{hmtfg:clasificacion:eq:4} reconstruyen exactamente \(K=W_n0\). Queda demostrado el teorema. \(\square\)

El teorema demuestra también que cada mapa clasificado posee los retornos restrictivos de duplicación anteriores a su retorno superestable, con los dominios de \eqref{hmtfg:clasificacion:eq:2}–\eqref{hmtfg:clasificacion:eq:3}. No presupone una monotonía del parámetro de una familia.

\textbf{Normalización simétrica de la familia fijada.} Cuando \(F\) es par y está definido sobre \([-1,1]\), las desigualdades anteriores dan además
\[
F(d_4)=d_5>d_3=F(d_2)=F(-d_2).
\]
Como ambos argumentos \(d_4,-d_2\) son positivos y \(F\) decrece en esa rama,
\begin{equation}
0<d_4<-d_2<1.
\label{hmtfg:clasificacion:eq:5a}
\end{equation}
Por tanto \(F^2([d_2,-d_2])=[d_2,d_4]\), y la misma conjugación \(h(x)=d_2x\) define sobre todo \([-1,1]\) un mapa par de máximo cuya imagen es \([d_4/d_2,1]\subset[-1,1]\). Coincide exactamente con el operador \(RF=-F(F(-ax))/a\), \(a=-F(1)\), utilizado en el corpus. Esta ampliación simétrica conserva el retorno crítico, su período y su itinerario.


\subsection{Primer parámetro de itinerario infinito}
\label{hmtfg:clasificacion:3}

Trabajamos sobre el intervalo compacto ya utilizado en la demostración de existencia,
\[
L=\left[\frac1{2u_0(1)},\frac2{u_0(1)}\right].
\]
Sea
\[
\Lambda_\tau=\{\lambda\in L:
\operatorname{sign}c_j(\lambda)=\tau_j\text{ para todo }j\ge1\}.
\]
Por las demostraciones de existencia y separación de signos de la prolongación, este conjunto es no vacío y compacto. Por tanto existe
\begin{equation}
a=\min\Lambda_\tau.
\label{hmtfg:clasificacion:eq:6}
\end{equation}

\textbf{Lema de comparación.} Para todo \(\lambda\in[1/(2u_0(1)),a)\), se tiene \(K_\lambda<\tau\).

\textbf{Demostración.} En el extremo izquierdo, el itinerario es \(+^\infty<\tau\). El argumento de separación de la realización de la palabra infinita demuestra que los conjuntos de comparación estricta con \(\tau\) son abiertos: en un itinerario infinito distinto de \(\tau\), la primera diferencia finita persiste por continuidad; en un parámetro superestable, las dos palabras periódicas vecinas y la palabra con cero final se sitúan en el mismo lado de \(\tau\), por su maximalidad estricta. El intervalo anterior a \(a\) carece, por definición, de parámetros con itinerario \(\tau\). Su conexidad mantiene la comparación inicial en todo ese intervalo. \(\square\)


\subsection{Existencia de cada primera raíz nueva}
\label{hmtfg:clasificacion:4}

Conservamos el selector original. La primera raíz es
\[
\lambda_1=\frac1{u_0(1)},
\]
y, para \(n\ge2\), \(\lambda_n\) es la primera raíz de \(c_{2^n}\) situada a la derecha de \(\lambda_{n-1}\) que satisface
\begin{equation}
c_{2^j}(\lambda_n)\ne0\quad(1\le j<n).
\label{hmtfg:clasificacion:eq:7}
\end{equation}
Como cualquier primer período que divida \(2^n\) es una potencia de dos, \eqref{hmtfg:clasificacion:eq:7} equivale a exigir período crítico exacto \(2^n\).

\textbf{Lema de existencia.} Todas las raíces de este selector existen y cumplen
\begin{equation}
\lambda_1<\lambda_2<\cdots<\lambda_n<a.
\label{hmtfg:clasificacion:eq:8}
\end{equation}

\textbf{Demostración.} El signo \(c_2(a)<0\) implica \(\lambda_1<a\). Supongamos construida \(\lambda_{n-1}<a\), y pongamos \(p=2^{n-1}\). La función analítica \(c_p(\lambda)\) no es idénticamente cero, puesto que \(c_p(a)\ne0\). Su conjunto de ceros en \([\lambda_{n-1},a]\) es finito, no vacío y carece de \(a\). Sea \(z<a\) su último cero.

En \((z,a]\), el signo de \(c_p\) es constante e igual a \(s=\tau_p\). Para \(g_\lambda=f_\lambda^p\), se tiene
\[
g_\lambda(0)=c_p(\lambda),\qquad g'_\lambda(0)=0.
\]
Una cota uniforme local de la segunda derivada da
\begin{equation}
c_{2p}(\lambda)
=g_\lambda(c_p(\lambda))
=c_p(\lambda)+O(c_p(\lambda)^2).
\label{hmtfg:clasificacion:eq:9}
\end{equation}
Al aproximarse \(\lambda\) a \(z\) por la derecha, el término de primer orden domina; por tanto \(c_{2p}(\lambda)\) tiene signo \(s\). En \(a\), su signo es
\[
\tau_{2p}=-\tau_p=-s.
\]
El teorema del valor intermedio proporciona un cero \(\beta\in(z,a)\) de \(c_{2p}\). Como \(c_p\ne0\) en ese intervalo, el período crítico en \(\beta\) es exactamente \(2p\).

La analiticidad y \(c_{2p}(a)\ne0\) implican que los ceros de \(c_{2p}\) en \([\lambda_{n-1},a]\) son finitos. Entre sus ceros de período exacto \(2p\), situados a la derecha de \(\lambda_{n-1}\), existe así un primero. Es el selector original \(\lambda_n\), y satisface \(\lambda_n\le\beta<a\). \(\square\)

La utilización del último cero auxiliar \(z\) pertenece únicamente a la prueba de existencia. El parámetro seleccionado continúa siendo la \textbf{primera raíz nueva} \(\lambda_n\).


\subsection{Convergencia del selector global al primer itinerario infinito}
\label{hmtfg:clasificacion:5}

\textbf{Teorema de continuación global.} Para la familia HMT fijada y el selector original, todos los itinerarios seleccionados son canónicos y
\begin{equation}
\boxed{K_{\lambda_n}=W_n0,\qquad
\lambda_n\nearrow\min\Lambda_\tau.}
\label{hmtfg:clasificacion:eq:10}
\end{equation}

\textbf{Demostración.} Los lemas de comparación y existencia sitúan cada parámetro seleccionado por debajo de \(a\), con itinerario menor que \(\tau\). El teorema de clasificación identifica ese itinerario con \(W_n0\). La sucesión es creciente y acotada por \(a\); sea \(b\le a\) su límite.

Fijemos \(p\ge1\). Para \(n\) suficientemente grande, \(2p<2^n\). Entonces los signos de \(c_p(\lambda_n)\) y \(c_{2p}(\lambda_n)\) coinciden con \(\tau_p\) y \(\tau_{2p}=-\tau_p\), respectivamente. La cota uniforme de las segundas derivadas,
\[
\|(f_\lambda^p)''\|\le B_p,
\qquad B_p=16^p\frac{16^p-1}{15},
\]
y la identidad \((f_\lambda^p)'(0)=0\) implican
\[
|c_{2p}(\lambda_n)-c_p(\lambda_n)|
\le\tfrac12B_p|c_p(\lambda_n)|^2.
\]
Como los dos términos del lado izquierdo tienen signos contrarios,
\begin{equation}
|c_p(\lambda_n)|>2/B_p.
\label{hmtfg:clasificacion:eq:11}
\end{equation}
Por continuidad, \(c_p(b)\) conserva el signo \(\tau_p\) y satisface \(|c_p(b)|\ge2/B_p>0\). Esto vale para cada \(p\), de modo que \(b\in\Lambda_\tau\). La minimalidad de \(a\) da \(b\ge a\), y por tanto \(b=a\). \(\square\)

La prueba completa conserva el selector global y produce sus itinerarios y su límite. La información de los primeros ocho niveles es compatible con el resultado, pero la demostración de \eqref{hmtfg:clasificacion:eq:10} no extrapola esos ocho niveles: utiliza inducción, intervalos invariantes y separación uniforme de cada signo.


\subsection{Controles lógicos y falsadores}
\label{hmtfg:clasificacion:7}

\par\noindent\textbf{1.}\enspace El teorema de clasificación requiere un máximo unimodal estricto y un intervalo invariante. Si alguno de los dos ramos pierde monotonía, los intervalos \eqref{hmtfg:clasificacion:eq:2} pueden dejar de ser restrictivos y la clasificación debe comprobarse de nuevo.
\par\noindent\textbf{2.}\enspace La condición \(K<\tau\) es esencial. Después de la palabra límite pueden existir centros de período \(2^n\) con otras combinatorias; el teorema no los identifica con \(W_n0\).
\par\noindent\textbf{3.}\enspace La analiticidad en el parámetro proporciona aislamiento y finitud de los ceros sobre el compacto. Para una familia solamente continua, la existencia de una primera raíz estrictamente posterior requiere un argumento adicional.
\par\noindent\textbf{4.}\enspace Un límite de raíces críticas puede, en general, perder signos al alcanzar un retorno de período menor. La desigualdad \eqref{hmtfg:clasificacion:eq:11} excluye precisamente ese fenómeno para esta sucesión y cada horizonte fijo.
\par\noindent\textbf{5.}\enspace La convergencia de las primeras raíces al primer parámetro de itinerario \(\tau\) no equivale a la unicidad de todos los parámetros de ese itinerario ni determina por sí misma una tasa métrica. El cruce transversal y la identificación de la rama aportan esas conclusiones posteriores.

% Fragmento integrable. Preambulo anfitrion: amsmath, amssymb, longtable, hyperref.
\section{Entrada transversal y escalado local}
\label{hmtfg:fragmento:entrada-transversal-y-escalado-local}

\subsection{Procedencia y resultado}
\label{hmtfg:escalado:1}

APP evalúa suma y producto sobre las dos hojas de la retícula, conservando cociente y residuo. TRIT determina régimen y orientación. El transporte TPK selecciona, transporta y actualiza el estado con sus registros de acarreo, frontera y memoria. Sus lectores actúan sobre la misma estructura discreta del continuo. La prolongación \texttt{w6→w12→w18→w24→w30→R36→G9} conserva esta procedencia y distingue retorno de fase de incremento de memoria.

El lector crítico de la rueda dodecafásica, construido en los antecedentes operatorios y transportado por el refinamiento compatible, produce
\[
\phi_j=\frac{(3j-1)\pi_{\mathrm{HMT}}}{18},\qquad
w_j=\frac1{12},\qquad
k_j=\sqrt{\frac2{\sum_{\ell=1}^{12}w_\ell\phi_\ell^2}}\,\phi_j
=\frac{2(3j-1)}{\sqrt{899}}.
\]
La coordenada \(\pi_{\mathrm{HMT}}\) conserva su origen APP–TRIT–TPK; su factor común se cancela en este segundo momento. El lector resultante es
\begin{equation}
u_0(x)=\frac1{12}\sum_{j=1}^{12}(1-\cos k_jx),\qquad
f_\lambda(x)=1-\lambda u_0(x).
\label{hmtfg:escalado:eq:1}
\end{equation}
El sector focal es el uniforme, \(\eta=0\). La ampliación antecedente conserva por separado la familia ponderada \(|\eta|\le1/40\). Las nuevas cotas numéricas corresponden a \eqref{hmtfg:escalado:eq:1}.

La familia es par, entera, con \(f_\lambda(0)=1\), crítico cuadrático y segundo momento \(\sum_jw_jk_j^2=2\). El parámetro \(\lambda\) recorre una familia posterior al generador HMT. La constante objetivo de Feigenbaum permanece fuera de las entradas de \eqref{hmtfg:escalado:eq:1}.

\textbf{Resultado principal.} Existe un único
\begin{equation}
\lambda_*\in[1.874038,1.874039]
\label{hmtfg:escalado:eq:2}
\end{equation}
para el que el cuarto retorno normalizado de \eqref{hmtfg:escalado:eq:1} pertenece a la hoja estable local del punto fijo real de duplicación. El cruce de la familia con esa hoja es transversal, con separación cuantificada. Para todo \(n\ge0\),
\begin{equation}
\boxed{\|R^{4+n}f_{\lambda_*}-g\|_{\mathcal A}
<0.001666\,(0.83996)^n.}
\label{hmtfg:escalado:eq:3}
\end{equation}
La cascada superestable local asociada a este cruce posee parámetros \(\mu_j\) y constantes \(C\ne0\), \(0<\theta<1\), tales que
\begin{equation}
\mu_j-\lambda_*=C\delta_{\mathrm F}^{-j}\bigl[1+O(\theta^j)\bigr],
\qquad
\frac{\mu_{j-1}-\mu_{j-2}}{\mu_j-\mu_{j-1}}\longrightarrow\delta_{\mathrm F}.
\label{hmtfg:escalado:eq:4}
\end{equation}
Aquí \(\delta_{\mathrm F}>1\) es el autovalor inestable del operador de duplicación. La selección de \(\mu_j\) se realiza mediante los retornos locales de sección~\ref{hmtfg:escalado:8}. La identificación con el selector global se demuestra en secciones~\ref{hmtfg:escalado:15}--\ref{hmtfg:escalado:17}: todos los términos globales conservan la combinatoria y los centros coinciden para todo nivel suficientemente alto al igualar sus períodos.


\subsection{Carta analítica y retorno efectivo}
\label{hmtfg:escalado:2}

Se emplea la carta de funciones pares
\[
s=x^2,\qquad z=\frac{s-1}{5/2},\qquad
f(x)=1-x^2V(z),\qquad V(z)=\sum_{j\ge0}v_jz^j.
\]
Las coordenadas del espacio de Banach son
\begin{equation}
u=10v_0,\qquad \nu=(v_1,v_2,\ldots),\qquad
\|(\Delta u,\Delta\nu)\|_{\mathcal A}
=|\Delta u|+\sum_{j\ge1}|\Delta v_j|.
\label{hmtfg:escalado:eq:5}
\end{equation}
En particular, el coeficiente constante de \(V\) tiene peso diez. Esta normalización coincide con la carta \(u/10\) de Hertling–Spandl, sección 2 de la fuente.

Sea \(a=-f(1)=v_0-1\). Para los retornos reales admisibles,
\begin{equation}
(Rf)(x)=-\frac{f(f(-ax))}{a}.
\label{hmtfg:escalado:eq:6}
\end{equation}
El programa evalúa una factorización exacta de \eqref{hmtfg:escalado:eq:6}. Definiendo
\[
S=\frac{a^2s-1}{5/2},\qquad b=V(S),\qquad
h=1-a^2sb,\qquad t=\frac{h^2-1}{5/2},
\]
\[
\operatorname{shift}V(z)=\sum_{j\ge0}v_{j+1}z^j,
\]
resulta
\begin{equation}
\boxed{V_R=a\,b\,(h+1)\left[V(t)+\frac25\operatorname{shift}V(t)\right].}
\label{hmtfg:escalado:eq:7}
\end{equation}
Para comprobarla, se usa \(v_0=1+a\),
\(h^2-1=-a^2sb(h+1)\) y
\(V(t)-v_0=t\operatorname{shift}V(t)\). Así,
\(a+1-h^2V(t)=a^2sb(h+1)[V(t)+(2/5)\operatorname{shift}V(t)]\),
que es precisamente \(asV_R\).

Los momentos iniciales son racionales:
\begin{equation}
\frac1{12}\sum_jk_j^{2m}
=\frac1{12}\left(\frac4{899}\right)^m
\sum_{j=1}^{12}(3j-1)^{2m}.
\label{hmtfg:escalado:eq:8}
\end{equation}
Por ello, tanto \eqref{hmtfg:escalado:eq:7} como la construcción inicial utilizan intervalos racionales; las funciones trigonométricas iniciales se representan por sus series con resto exterior.


\subsection{Cotas externas utilizadas como reconocimiento posterior}
\label{hmtfg:escalado:3}

Una vez construida \eqref{hmtfg:escalado:eq:1}, se utiliza el teorema de hiperbolicidad de Lanford como certificado analítico del operador \eqref{hmtfg:escalado:eq:6}. Sean \(g_0\) el polinomio central publicado y \(g\) el punto fijo. Las estimaciones utilizadas son
\[
\|g-g_0\|_{\mathcal A}<\varepsilon_0=4\,10^{-5},
\]
y, en la bola \(\|f-g_0\|_{\mathcal A}<0.01\), los bloques de \(DR\), escritos \(R=(U,V)\), satisfacen
\begin{equation}
U_u\ge a_0=4.521,\quad
\|U_\nu\|\le b_0=0.560,\quad
\|V_u\|\le c_0=0.756,\quad
\|V_\nu\|\le d_0=0.719.
\label{hmtfg:escalado:eq:9}
\end{equation}
El punto fijo posee una única dirección inestable, con autovalor real positivo \(\delta_{\mathrm F}>1\); el resto del espectro queda dentro del disco unidad. Las cifras de \eqref{hmtfg:escalado:eq:9} proceden de las estimaciones publicadas, no de una nueva reproducción de la prueba de Lanford en este expediente. Los coeficientes de \(g_0\) están incorporados explícitamente en el polinomio central de referencia del programa, conforme a la tabla 2 de Hertling–Spandl.

Fuentes: \href{https://repo-archives.ihes.fr/A_GARDER/20180409/Test/I_Prepublications/LANFORD/1968-2013/P_81_17/P_81_17_web.pdf}{Lanford, estimaciones de hiperbolicidad}; \href{https://arxiv.org/pdf/1410.3277}{Hertling–Spandl, carta analítica y tabla 2}.


\subsection{Certificado racional de entrada y dirección}
\label{hmtfg:escalado:4}

El certificado de entrada a la renormalización divide \eqref{hmtfg:escalado:eq:2} en dieciséis intervalos racionales iguales. Propaga cuatro retornos de \eqref{hmtfg:escalado:eq:7} y la derivada exacta respecto a \(\lambda\), usando diferenciación de la composición y redondeo exterior a cien posiciones decimales.

Cada serie consta de un polinomio de grado 32 con coeficientes intervalares y un error analítico arbitrario \(E\), \(\|E\|_1\le\epsilon\). Este error puede modificar también coeficientes bajos. Para \(\|q\|_1<1\),
\begin{equation}
\|E\circ q\|_1\le\epsilon,\qquad
\|E'\circ q\|_1\le\frac{\epsilon}{(1-\|q\|_1)^2},\qquad
\|\operatorname{shift}E\|_1\le\epsilon.
\label{hmtfg:escalado:eq:10}
\end{equation}
Los productos incluyen los términos polinomio–error y error–error. La cola inicial se acota por su primer término absoluto y una razón geométrica, utilizando \(\|1+(5/2)z\|_1=7/2\). Cada argumento compuesto permanece estrictamente dentro de la bola unidad de esta álgebra.

Para \(\Gamma(\lambda)=R^4f_\lambda=(u_\lambda,\nu_\lambda)\), el certificado da las siguientes cotas conservadoras, uniformes en \eqref{hmtfg:escalado:eq:2}:

\begingroup\small
\begin{longtable}{p{0.465\linewidth}p{0.465\linewidth}}
\hline
Cantidad & Cota certificada exterior \\
\hline
\(\|\Gamma(\lambda)-g_0\|_{\mathcal A}\) & \(<0.004993128\) \\
\(\|\nu_\lambda-\nu_{g_0}\|_1\) & \(<0.001396077\) \\
\(u'_\lambda\) & \(>3821.289115\) \\
\(\|\nu'_\lambda\|_1\) & \(<551.757000\) \\
\(\|\nu'_\lambda\|_1/u'_\lambda\) & \(<0.144386<1/4\) \\
Error analítico de la función & \(<8.482\,10^{-19}\) \\
Error analítico de la tangente & \(<4.569\,10^{-15}\) \\
\hline\end{longtable}
\endgroup

Los decimales completos, los dieciséis subintervalos y cada retorno intermedio se conservan en el certificado racional de entrada.

El cono \(\|\dot\nu\|_1\le|\dot u|/4\) es invariante bajo \eqref{hmtfg:escalado:eq:9}:
\[
|\dot u_{\rm nuevo}|\ge(4.521-0.560/4)|\dot u|
=4.381|\dot u|,
\]
\begin{equation}
\|\dot\nu_{\rm nueva}\|_1\le(0.756+0.719/4)|\dot u|
=0.93575|\dot u|<\tfrac14(4.381)|\dot u|.
\label{hmtfg:escalado:eq:11}
\end{equation}
La dirección paramétrica certificada entra, por tanto, en un cono uniformemente expansivo del operador.


\subsection{Construcción cuantitativa de la hoja estable}
\label{hmtfg:escalado:5}

Traslademos \(g\) al origen y escribamos las coordenadas como \((x,y)=(u-u_g,\nu-\nu_g)\). Fijemos
\begin{equation}
L=0.16,\qquad r=0.004,\qquad
A=a_0-Lc_0=4.40004,\qquad q=d_0+c_0L=0.83996.
\label{hmtfg:escalado:eq:12}
\end{equation}
El cilindro \(|x|\le Lr\), \(\|y\|_1\le r\) está dentro de la bola de \eqref{hmtfg:escalado:eq:9}, puesto que
\((1+L)r+\varepsilon_0=0.00468<0.01\).

Considérese el espacio completo de funciones \(\sigma:B_r\to\mathbb R\) con \(\sigma(0)=0\) y constante de Lipschitz a lo sumo \(L\), dotado de la distancia uniforme. Para cada \(y\), se define \((\mathcal G\sigma)(y)\) como la raíz en \([-L\|y\|_1,L\|y\|_1]\) de
\begin{equation}
F_\sigma(x,y)=U(x,y)-\sigma(V(x,y)).
\label{hmtfg:escalado:eq:13}
\end{equation}
En ese intervalo, \(\|V(x,y)\|_1\le q\|y\|_1<r\). Al incrementar \(x\), \eqref{hmtfg:escalado:eq:9} muestra que \(F_\sigma\) crece al menos con pendiente \(A\). En sus extremos,
\[
F_\sigma(L\|y\|_1,y)
\ge[AL-b_0-Ld_0]\|y\|_1,
\]
y la desigualdad opuesta vale en el extremo negativo. El margen es
\begin{equation}
AL-b_0-Ld_0=0.0289664>0.
\label{hmtfg:escalado:eq:14}
\end{equation}
Se obtiene una raíz única. Comparando \eqref{hmtfg:escalado:eq:13} para dos valores de \(y\),
\[
\operatorname{Lip}(\mathcal G\sigma)
\le\frac{b_0+Ld_0}{A}
=\frac{0.67504}{4.40004}<0.16.
\]
Comparando dos gráficos,
\begin{equation}
\|\mathcal G\sigma-\mathcal G\widetilde\sigma\|_\infty
\le A^{-1}\|\sigma-\widetilde\sigma\|_\infty,
\qquad A^{-1}<0.228.
\label{hmtfg:escalado:eq:15}
\end{equation}
El teorema de contracción produce un único gráfico fijo \(x=\sigma_*(y)\). Sobre él,
\begin{equation}
\|y_n\|_1\le q^n\|y_0\|_1,
\qquad
\|(x_n,y_n)\|_{\mathcal A}\le(1+L)q^n\|y_0\|_1.
\label{hmtfg:escalado:eq:16}
\end{equation}
Además, la distancia vertical \(d(x,y)=x-\sigma_*(y)\) cumple
\[
|d(R(x,y))|\ge A|d(x,y)|
\]
mientras la órbita permanece en el cilindro. Toda órbita que permanece allí indefinidamente pertenece al gráfico. Este identifica exactamente la hoja estable local y proporciona la estimación \eqref{hmtfg:escalado:eq:16}. Su regularidad local coincide con la variedad estable analítica del punto fijo hiperbólico.


\subsection{Existencia, unicidad y transversalidad del parámetro}
\label{hmtfg:escalado:6}

Defínase, para \(\lambda\) en \eqref{hmtfg:escalado:eq:2},
\[
H(\lambda)=u_\lambda-u_g-
\sigma_*(\nu_\lambda-\nu_g).
\]
El certificado asegura \(\|\nu_\lambda-\nu_g\|_1<0.001436077<r\). Para \(\lambda_2>\lambda_1\),
\begin{equation}
H(\lambda_2)-H(\lambda_1)
\ge\int_{\lambda_1}^{\lambda_2}
\bigl[u'_t-L\|\nu'_t\|_1\bigr]dt
>3733.007995\,(\lambda_2-\lambda_1).
\label{hmtfg:escalado:eq:17}
\end{equation}
Para decidir los signos extremos se utiliza \(\|g-g_0\|_{\mathcal A}<\varepsilon_0\) y \(|\sigma_*(y)|\le L\|y\|_1\). La evaluación puntual racional produce
\begin{equation}
H(1.874038)<-0.0011856110945,
\qquad H(1.874039)>0.0021866053399.
\label{hmtfg:escalado:eq:18}
\end{equation}
Por continuidad, existe una raíz; \eqref{hmtfg:escalado:eq:17} prueba su unicidad y la transversalidad del cruce. En esa raíz, \eqref{hmtfg:escalado:eq:16} y
\((1+L)(0.001396077+0.00004)<0.001666\)
demuestran \eqref{hmtfg:escalado:eq:3}.

La cota es uniforme en el número de retornos: convierte el avance de profundidad en un error operatorial geométrico explícito.


\subsection{Admisibilidad real de todos los retornos}
\label{hmtfg:escalado:7}

Los cuatro retornos iniciales se acompañan, en cada subintervalo racional, de las desigualdades
\begin{equation}
0<a=-f(1)<1,\qquad b=f(a)>a,\qquad f(b)<a.
\label{hmtfg:escalado:eq:19}
\end{equation}
El programa guarda \(a,b,f(b)\) antes de cada aplicación de \eqref{hmtfg:escalado:eq:7}. La unimodalidad, la paridad y el crítico cuadrático se conservan bajo estos retornos normalizados.

Para las iteraciones siguientes basta una verificación uniforme en la bola de \eqref{hmtfg:escalado:eq:9}. Si \(E=V-V_{g_0}\), la norma \eqref{hmtfg:escalado:eq:5} da, para \(-0.4\le z\le0\),
\[
|E(z)|\le0.004,\qquad |E'(z)|\le0.01.
\]
La segunda desigualdad usa \(\sup_{j\ge1}j(0.4)^{j-1}=1\). La evaluación racional de los coeficientes centrales proporciona
\begin{equation}
\frac65<V(z)<\frac85,\qquad |V'(z)|<\frac12,
\qquad 0.39<a<0.41.
\label{hmtfg:escalado:eq:20}
\end{equation}
En consecuencia,
\[
f'(x)=-2x\left[V(z)+\frac25x^2V'(z)\right],
\]
y el corchete supera uno para \(0\le x\le1\). El mapa es unimodal y preserva \([-1,1]\). Además,
\[
b=f(a)>1-(0.41)^2\frac85=\frac{4569}{6250}>a,
\]
\begin{equation}
f(b)<1-\left(\frac{4569}{6250}\right)^2\frac65
=\frac{35028967}{97656250}<0.39<a.
\label{hmtfg:escalado:eq:21}
\end{equation}
Las órbitas de la hoja estable satisfacen \eqref{hmtfg:escalado:eq:19} a toda profundidad. Así, las iteraciones analíticas del operador corresponden a retornos reales de duplicación, con sus dominios y orientaciones conservados.


\subsection{Escalado de la cascada superestable local}
\label{hmtfg:escalado:8}

La conclusión métrica requiere dos cruces: el de la familia con la hoja estable, demostrado en sección~\ref{hmtfg:escalado:6}, y el de la variedad inestable con el blanco superestable. El segundo está probado por Eckmann–Wittwer para
\(\Sigma_2=\{\psi:\psi(1)=0\}\), correspondiente al ciclo superestable de período dos. La normalización \eqref{hmtfg:escalado:eq:6} coincide con la del operador real par empleado allí.

El teorema general de Collet–Eckmann–Lanford, sección 6 de la fuente citada, proposición 6.1 y teoremas 6.2–6.3, aplica a un mapa de Banach \(C^2\), un punto fijo con un único autovalor inestable \(\delta_{\mathrm F}>1\), una curva transversal a su hoja estable y un blanco transversal a su variedad inestable. Las preimágenes locales del blanco intersectan la curva en puntos que satisfacen
\begin{equation}
\delta_{\mathrm F}^j(\mu_j-\lambda_*)\longrightarrow C\ne0.
\label{hmtfg:escalado:eq:22}
\end{equation}
El índice \(j\) cuenta los retornos locales. Como la curva inicial aquí es \(R^4f_\lambda\), el período físico es \(2^{j+5}\) cuando el blanco se toma en \(\Sigma_2\); cualquier número fijo de pasos empleado para traer el blanco al entorno local se incorpora en un desplazamiento fijo del índice y en \(C\).

Las condiciones reales de sección~\ref{hmtfg:escalado:7} y la rama real de duplicación del blanco fijan el período exacto y el itinerario. La sucesión se selecciona por esta continuación local orientada. Tomando diferencias en \eqref{hmtfg:escalado:eq:22} se obtiene el límite de razones de \eqref{hmtfg:escalado:eq:4}.

Fuentes de los dos resultados utilizados: \href{https://link.springer.com/article/10.1007/BF01013368}{Eckmann–Wittwer, \emph{A complete proof of the Feigenbaum conjectures}}; \href{https://people.math.harvard.edu/~knill/history/lanford/papers/ColletEckmannLanford.pdf}{Collet–Eckmann–Lanford, sección 6}.

\subsubsection{Resto asintótico de potencia}

La analiticidad de nuestra familia permite precisar \eqref{hmtfg:escalado:eq:22}. En la construcción de coordenadas normales de Collet–Eckmann–Lanford se aplanan únicamente la variedad estable, la variedad inestable y la hipersuperficie blanco. Estas tres subvariedades admiten cambios locales \(C^2\) con derivadas acotadas. El corte suave empleado se realiza en la coordenada escalar inestable; la foliación completa se obtiene después mediante el cambio construido.

En esas coordenadas, la corrección se escribe como una serie de incrementos \(w_j\) cuya norma \(C^1\) satisface
\[
\|w_j\|_{C^1}\le C_1\rho^j,\qquad 0<\rho<1.
\]
La regla de la cadena para el mapa cortado, con derivadas de órdenes uno y dos uniformemente acotadas, proporciona
\[
\|w_j\|_{C^2}\le C_2B^j
\]
para algún \(B\ge1\). La norma ponderada de la construcción controla la norma \(C^1\) ordinaria en el dominio considerado. Por interpolación,
\[
[Dw_j]_\beta\le C_3
\left(\rho^{1-\beta}B^\beta\right)^j.
\]
Se elige
\begin{equation}
0<\beta<\min\left\{1,
\frac{-\log\rho}{\log B-\log\rho}\right\}.
\label{hmtfg:escalado:eq:23}
\end{equation}
La serie converge en \(C^{1,\beta}\). Sea \(z\) la coordenada normal resultante, normalizada para que las preimágenes del blanco satisfagan \(z=\delta_{\mathrm F}^{-j}\), absorbiendo el desplazamiento fijo de índice. Como nuestra curva es analítica y transversal,
\[
z(\Gamma(\lambda))=c(\lambda-\lambda_*)
+O(|\lambda-\lambda_*|^{1+\beta}),\qquad c\ne0.
\]
La inversión local da
\begin{equation}
\mu_j-\lambda_*=C\delta_{\mathrm F}^{-j}
+O\left(\delta_{\mathrm F}^{-(1+\beta)j}\right).
\label{hmtfg:escalado:eq:24}
\end{equation}
Esto demuestra el resto de \eqref{hmtfg:escalado:eq:4}, con \(\theta=\delta_{\mathrm F}^{-\beta}\). El enunciado afirma la existencia de las constantes del resto paramétrico. La tasa operatorial explícita \(0.83996\) de \eqref{hmtfg:escalado:eq:3} controla, por su parte, la convergencia de los mapas renormalizados.

\subsubsection{Conversión del resto en error de las razones}

Si \(\mu_j-\lambda_*=C\delta_{\mathrm F}^{-j}(1+e_j)\), con \(|e_j|\le M\theta^j\), se escribe
\[
\mu_{j-1}-\mu_j=C(\delta_{\mathrm F}-1)\delta_{\mathrm F}^{-j}(1+r_j),
\quad
r_j=\frac{\delta_{\mathrm F} e_{j-1}-e_j}{\delta_{\mathrm F}-1}.
\]
Con \(K=M(\delta_{\mathrm F}+\theta)/(\delta_{\mathrm F}-1)\), se tiene \(|r_j|\le K\theta^{j-1}\). Por tanto,
\begin{equation}
\boxed{
\left|\frac{\mu_{j-2}-\mu_{j-1}}{\mu_{j-1}-\mu_j}-\delta_{\mathrm F}\right|
\le\frac{\delta_{\mathrm F} K(1+\theta)\theta^{j-2}}
{1-K\theta^{j-1}}}
\label{hmtfg:escalado:eq:25}
\end{equation}
cuando \(K\theta^{j-1}<1\). Esta fórmula separa la precisión de evaluación de cada raíz del error de truncamiento de la cascada.


\subsection{Continuación certificada de las ocho raíces iniciales}
\label{hmtfg:escalado:9}

El certificado de continuación finita conserva la familia \eqref{hmtfg:escalado:eq:1} y el certificado anterior de ocho raíces. Para \(2\le n\le8\), estudia
\[
S_n(\lambda)=f_\lambda^{2^n}(0)
\]
desde la caja certificada de \(\lambda_{n-1}\) hasta \(1.874039\). La cobertura racional completa contiene exactamente la raíz heredada \(\lambda_{n-1}\) y la raíz nueva \(\lambda_n\). En particular, \(\lambda_n\) es la única raíz nueva en \((\lambda_{n-1},1.874039]\).

Las cajas se excluyen por signo del retorno o por una derivada local de signo certificado. Las cajas ancla conservan existencia y unicidad por signos extremos y monotonía local. Sus extremos cubren exactamente cada intervalo; se verifican las adyacencias. La evaluación usa los momentos \eqref{hmtfg:escalado:eq:8}, Horner de orden 32, colas geométricas exteriores y \(|u_0''|\le2\). La cola uniforme del potencial es menor que \(1.294\,10^{-58}\) en \(|x|\le3/2\).

El certificado final contiene 650 cajas procesadas y 332 hojas de cobertura. El primer nivel tiene la expresión exacta \(\lambda_1=1/u_0(1)\). Las aproximaciones siguientes se muestran únicamente como orientación; sus cajas completas se conservan en los certificados:

\begingroup\small
\begin{longtable}{p{0.465\linewidth}p{0.465\linewidth}}
\hline
Nivel & Parámetro superestable \\
\hline
1 & 1.345913990471832792210949095… \\
2 & 1.763379787846910127290794030… \\
3 & 1.850413796950136464530566778… \\
4 & 1.868979756606798125150574309… \\
5 & 1.872955034435659740004205128… \\
6 & 1.873806367467030119412262311… \\
7 & 1.873988695221365362307168780… \\
8 & 1.874027744154450672897007573… \\
\hline\end{longtable}
\endgroup

% Fragmento integrable. Preambulo anfitrion: amsmath, amssymb, longtable, hyperref.
\section{Aislamiento del límite y escalado del selector global}
\label{hmtfg:fragmento:aislamiento-del-limite-y-escalado-del-selector-global}

\subsection{Continuación indefinida de la primera raíz y exclusión del intervalo intermedio}
\label{hmtfg:escalado:15}

La clasificación completa demostrada anteriormente conserva la familia y el selector de primera raíz nueva. El transporte ordenado conserva los ceros y su mínimo; la palabra de retornos conserva la lectura cronológica nonádica.

\subsubsection{Clasificación y existencia a toda profundidad}
\label{hmtfg:escalado:15.1}

Escribamos \(W_n=\tau_1\cdots\tau_{2^n-1}\), con \(\tau_j=(-1)^{v_2(j)}\), y \(a_*=\min\Lambda_\tau\). La existencia y compacidad de \(\Lambda_\tau\) están demostradas en el desarrollo de prolongación, \ref{hmtfg:prolongacion:6.1}–\ref{hmtfg:prolongacion:6.2}; su mínimo existe por compacidad y orden.

\textbf{Clasificación.} Todo mapa unimodal estricto de máximo, sobre un intervalo invariante, cuyo crítico tiene período exacto \(2^n\) y cuyo itinerario \(K\) satisface \(K<\tau\), tiene itinerario \(W_n0\).

La inducción se realiza mediante el retorno efectivo. Para período al menos ocho, la comparación con \(\tau\), junto con los intervalos que atraparían órbitas de signos incompatibles, fuerza el prefijo \((+,-,+,+,+,-)\). De aquí se obtienen
\[
J_1=[d_2,d_4],\quad F(J_1)=[d_3,1],\quad
d_4<d_3,\quad F^2(J_1)=J_1.
\]
El retorno normalizado por \(h(x)=d_2x\) tiene máximo uno, período mitad e itinerario \(K'_j=-K_{2j}<\tau\). El orden se conserva porque
\[
O_{2q-1}(K)=-O_{q-1}(K'),\qquad
K_{2q}-\tau_{2q}=-(K'_q-\tau_q).
\]
Los casos iniciales son \(+0\) y \(+-+0\); la inducción reconstruye \(W_n0\). Para los mapas pares, también \(0<d_4<-d_2<1\), de modo que el retorno se extiende a \([-1,1]\) y coincide con \eqref{hmtfg:escalado:eq:6}.

Antes de \(a_*\), todos los itinerarios son menores que \(\tau\): los conjuntos de comparación estricta son abiertos por la separación demostrada en \ref{hmtfg:prolongacion:6.1}; el intervalo anterior a \(a_*\) es conexo y comienza con \(+^\infty<\tau\).

\textbf{Existencia de cada primera raíz.} Supóngase construida \(\lambda_{n-1}<a_*\), y sea \(p=2^{n-1}\). Los ceros de \(c_p\) en \([\lambda_{n-1},a_*]\) forman un conjunto finito no vacío: la función es analítica y \(c_p(a_*)\ne0\). Sea \(z<a_*\) su último cero. En \((z,a_*]\), \(c_p\) tiene signo \(\tau_p\). Para \(g_\lambda=f_\lambda^p\),
\[
g_\lambda'(0)=0,\qquad
c_{2p}(\lambda)=c_p(\lambda)+O(c_p(\lambda)^2).
\]
Cerca de \(z\), por la derecha, \(c_{2p}\) tiene signo \(\tau_p\); en \(a_*\), tiene el contrario. Existe un cero de \(c_{2p}\) en \((z,a_*)\), y su período es exactamente \(2p\), pues allí \(c_p\ne0\). La finitud de los ceros proporciona la primera raíz nueva en \((\lambda_{n-1},a_*)\). El último cero auxiliar prueba existencia; la elección efectiva sigue siendo la primera raíz. La base es \(\lambda_1=1/u_0(1)<a_*\).

\subsubsection{Límite de la sucesión original}
\label{hmtfg:escalado:15.2}

\textbf{Teorema.} El selector original existe a todos los niveles y satisface
\begin{equation}
\boxed{K_{\lambda_n}=W_n0,\qquad
\lambda_n\nearrow a_*=\min\Lambda_\tau.}
\label{hmtfg:escalado:eq:50}
\end{equation}
\textbf{Demostración.} La existencia precedente da una sucesión creciente acotada por \(a_*\). La clasificación fija todos sus itinerarios. Sea \(b\le a_*\) su límite. Para cada \(p\), los términos suficientemente avanzados conservan los signos opuestos \(\tau_p,\tau_{2p}\). Taylor y la cota uniforme \(B_p=16^p(16^p-1)/15\) dan
\[
|c_{2p}(\lambda_n)-c_p(\lambda_n)|
\le\tfrac12 B_p|c_p(\lambda_n)|^2,\qquad
|c_p(\lambda_n)|>2/B_p.
\]
Al pasar al límite se conserva \(\operatorname{sign}c_p(b)=\tau_p\), con módulo al menos \(2/B_p>0\). Para todo \(p\), esto implica \(b\in\Lambda_\tau\); la minimalidad da \(b=a_*\).

El resultado resuelve existencia indefinida, conservación de la cascada simbólica y selección de su primer parámetro límite para la regla global original. La completación HMT transporta esta sucesión y su límite conservando orden y lectura.

\subsubsection{Certificado de exclusión del intervalo intermedio}
\label{hmtfg:escalado:15.3}

El programa de puente certifica
\begin{equation}
\Lambda_\tau\cap[\lambda_{8,\rm inf},1.874038]=\varnothing,
\label{hmtfg:escalado:eq:51}
\end{equation}
donde
\[
\lambda_{8,\rm inf}
=1.874027744154450672897007573595092408435133414523352839687903472138975
\]
es el extremo inferior racional del intervalo certificado de \(\lambda_8\).

La cobertura contiene 374 intervalos exactamente adyacentes y ninguna región pendiente: 319 con \(c_{512}>0\), 10 con \(c_{1024}<0\), dos con \(c_{2048}>0\), todos opuestos al signo de \(\tau\), y 43 próximos a centros de períodos 256, 512 o 1024, excluidos por su segundo retorno.

En este último caso, \(g=f_\lambda^p\), \(d=g(0)\) y una cota uniforme \(|g''|\le M\) entre cero y \(d\) satisfacen \(M|d|<2\). Entonces
\[
|g(d)-d|\le\tfrac12M|d|^2<|d|\quad(d\ne0).
\]
Ambos retornos tienen el mismo signo; para \(d=0\), ambos son cero. Las dos situaciones excluyen \(\tau_{2p}=-\tau_p\) con signos estrictos.

La evaluación utiliza momentos racionales, colas geométricas del potencial y de su segunda derivada, y \(|u_0'''|<6\). El centrado paramétrico intersecta la envolvente natural con
\[
c_j(\lambda_{\rm medio})+
\partial_\lambda c_j(I)\,(I-\lambda_{\rm medio});
\]
la derivada se propaga sobre la caja completa. La autoaplicación de \([-1,1]\) se demuestra antes de restringir a ella las envolventes.

El certificado completo conserva la cobertura y comprueba su teselado exacto, junto con cada desigualdad de signo o de Taylor.

\subsubsection{Localización conjunta}
\label{hmtfg:escalado:15.4}

El cruce local \(\lambda_*\) realiza \(\tau\) por sus retornos reales a toda profundidad, de modo que \(a_*\le\lambda_*\). Por \eqref{hmtfg:escalado:eq:50}, \(\lambda_8<a_*\); junto con \eqref{hmtfg:escalado:eq:51},
\begin{equation}
\boxed{1.874038<a_*\le\lambda_*<1.874039.}
\label{hmtfg:escalado:eq:52}
\end{equation}
La localización conserva el itinerario canónico y sitúa su primer parámetro límite en el mismo intervalo que el cruce estable.

\subsection{Aislamiento por primera salida y coincidencia de los límites}
\label{hmtfg:escalado:16}

Esta sección completa el aislamiento identificado en sección~\ref{hmtfg:escalado:15.4}. Conserva la familia HMT \eqref{hmtfg:escalado:eq:1}, la selección mínima \eqref{hmtfg:escalado:eq:50}, la carta \eqref{hmtfg:escalado:eq:5} y la hoja estable de secciones~\ref{hmtfg:escalado:5}--\ref{hmtfg:escalado:7}. La prueba utiliza el orden \(a_*\le\lambda_*\): basta excluir una separación hacia el lado negativo del cono expansivo.

\subsubsection{Cotas superior e inferior del transporte transversal}
\label{hmtfg:escalado:16.1}

Además de \eqref{hmtfg:escalado:eq:9}, el lema 2.1 de \href{https://arxiv.org/pdf/1410.3277}{Hertling–Spandl, sección 2} proporciona \(\|D\Phi\|<0.9\) en la misma bola de radio \(0.01\), donde
\[
\Phi_u(u,\nu)=u-\frac{U(u,\nu)-u}{3.669}.
\]
Por tanto,
\begin{equation}
\left|1-\frac{U_u-1}{3.669}\right|<0.9,\qquad
U_u<1+1.9(3.669)=7.9711.
\label{hmtfg:escalado:eq:53}
\end{equation}
El número \(3.669\) pertenece al precondicionador publicado de este certificado analítico; la familia generada \eqref{hmtfg:escalado:eq:1} conserva sus coeficientes previos.

Tomemos \(\kappa=0.21\). Para dos puntos de la bola unidos por una secante con
\(\Delta u<0\), \(\|\Delta\nu\|_1\le\kappa|\Delta u|\), la integración de los bloques de \(DR\) sobre el segmento da
\begin{equation}
4.4034|\Delta u|
\le|\Delta u'|\le8.0887|\Delta u|,\qquad \Delta u'<0,
\label{hmtfg:escalado:eq:54}
\end{equation}
\begin{equation}
\|\Delta\nu'\|_1
\le(0.756+0.719\kappa)|\Delta u|
=0.90699|\Delta u|
<\kappa(4.4034)|\Delta u|.
\label{hmtfg:escalado:eq:55}
\end{equation}
Aquí \(4.4034=4.521-0.560\kappa\) y \(8.0887=7.9711+0.560\kappa\). Así, el cono de secantes conserva orientación y se expande mientras el segmento permanezca en la bola.

\subsubsection{Reducción de una separación infinita a una región finita}
\label{hmtfg:escalado:16.2}

Supongamos \(a_*<\lambda_*\), y pongamos
\[
f_n=R^{4+n}f_{a_*},\qquad
s_n=R^{4+n}f_{\lambda_*},\qquad
(\Delta u_n,\Delta\nu_n)=f_n-s_n.
\]
La cota positiva de \(u'\) y \(\|\nu'\|_1/u'<0.144386<\kappa\) sobre \(\Gamma(J)\) sitúan la secante inicial en el cono negativo. La hoja estable está escrita respecto de \(g\), con
\begin{equation}
s_n=g+e_n,\quad
\|e_n\|_{\mathcal A}<0.001666(0.83996)^n,\quad
|(e_n)_u|\le0.16\|(e_n)_\nu\|_1.
\label{hmtfg:escalado:eq:56}
\end{equation}
Los controles de entrada implican
\begin{equation}
|\Delta u_0|
<0.004993128+0.00004
 +0.16(0.001396077+0.00004)
=0.00526290032< h,\qquad h=0.006.
\label{hmtfg:escalado:eq:57}
\end{equation}
Mientras \(|\Delta u_n|<h\), \eqref{hmtfg:escalado:eq:55}–\eqref{hmtfg:escalado:eq:56} dan
\begin{equation}
\|f_n-g_0\|_{\mathcal A}
<(1+\kappa)h+0.001666+0.00004
=0.008966<0.01.
\label{hmtfg:escalado:eq:58}
\end{equation}
El punto \(s_n\) y el segmento entre ambos también permanecen en esa bola convexa. La expansión inferior \eqref{hmtfg:escalado:eq:54} obliga a una primera salida finita \(N\), y la cota superior aplicada al paso precedente proporciona
\begin{equation}
0.006\le-\Delta u_N<0.0485322,\qquad
\|\Delta\nu_N\|_1\le0.21|\Delta u_N|.
\label{hmtfg:escalado:eq:59}
\end{equation}
Se utiliza la desigualdad \(8.0887h=0.0485322\); el salto completo, y no únicamente la superficie de primera salida, queda incluido.

En consecuencia, \(f_N\) pertenece a la siguiente región:
\begin{equation}
\begin{split}
\mathcal C_-=\{\,g+(d_u,d_\nu)+e:\;&
-0.0485322\le d_u\le-0.006,\\
&\|d_\nu\|_1\le0.21|d_u|,\quad
|e_u|+\|e_\nu\|_1\le0.001666,\\
&|e_u|\le0.16\|e_\nu\|_1\,\}.
\end{split}
\label{hmtfg:escalado:eq:60}
\end{equation}
Como \(a_*\) realiza \(\tau\), todos sus retornos normalizados son autoaplicaciones unimodales de \([-1,1]\) con el mismo itinerario. Esta propiedad procede de la inducción de retornos, independientemente de la distancia de \(f_N\) a \(g_0\). Por eso el certificado siguiente trata \(\mathcal C_-\) intersectada con esa clase de autoaplicaciones.

\subsubsection{Exclusión certificada de la región de primera salida}
\label{hmtfg:escalado:16.3}

El certificado racional de exclusión por primera salida, reproducido en el suplemento computacional, cubre exactamente el intervalo de \eqref{hmtfg:escalado:eq:59} mediante 64 bandas racionales adyacentes. Todas quedan excluidas; ninguna subdivisión adicional ni región pendiente intervienen en el resultado.

De las 64 bandas, 54 presentan un signo contrario y diez un retorno contractivo de período 16 o 32. La mayor cota de Lipschitz de esos retornos es menor que \(0.475146491\); el horizonte máximo de exclusión es el iterado 64.

La evaluación conserva como variables compartidas los primeros cuarenta coeficientes de la perturbación, junto con una cota explícita de su cola infinita. Para cada banda, el centro es \(g_0\) desplazado en la coordenada \(u\). Si \(a,b\ge0\) acotan las sensibilidades duales en \(u,\nu\), el soporte de la incertidumbre del punto fijo y del residuo estable es
\begin{equation}
0.00004\max(a,b)
+0.001666\max\left(b,\frac{0.16a+b}{1.16}\right).
\label{hmtfg:escalado:eq:61}
\end{equation}
La segunda expresión se obtiene maximizando \(a|e_u|+b\|e_\nu\|_1\) bajo las dos restricciones de \eqref{hmtfg:escalado:eq:56}. De este modo se utiliza la orientación de la hoja estable demostrada previamente, además de su norma.

Sea \(x_j\) la órbita del centro y \(L_j\) su sensibilidad lineal a todos los coeficientes compartidos. El modelo de Taylor tiene la forma
\[
c_j=x_j+L_j(\Delta v)+r_j,\qquad |r_j|\le E_j.
\]
Si \(\rho_j\) acota \(|L_j(\Delta v)|+E_j\), la recurrencia exterior empleada es
\begin{equation}
E_{j+1}\le |f_0'(x_j)|E_j
+\tfrac12\sup|f_0''|\rho_j^2
+\sup|\Delta f'|\rho_j.
\label{hmtfg:escalado:eq:62}
\end{equation}
Las derivadas se acotan sobre el segmento de los dos estados. La órbita central se verifica dentro de \([-1,1]\); las otras órbitas pertenecen a ese intervalo por la autoaplicación incluida en el dominio certificado. Los cálculos utilizan redondeo racional exterior, incluyendo las colas geométricas.

Cada banda se excluye por una de dos desigualdades:

\par\noindent\textbf{1.}\enspace Un valor crítico tiene signo estrictamente opuesto a \(\tau_j\).
\par\noindent\textbf{2.}\enspace Para \(p=2^k\), el retorno \(H=f^p\) tiene constante de Lipschitz \(L_p<1\) sobre el segmento entre \(0\) y \(H(0)\). Entonces
\[
|H(H(0))-H(0)|\le L_p|H(0)|.
\]
Si \(H(0)\ne0\), los dos valores tienen el mismo signo; si \(H(0)=0\), ambos se anulan. En ambos casos queda excluido el itinerario estricto \(\tau_{2p}=-\tau_p\).

La constante \(L_p\) se obtiene propagando el segmento inicial \([-r,r]\), \(r\ge|H(0)|\), junto a las cajas de la órbita crítica y multiplicando las cotas exteriores de las derivadas. Se acotan los mapas de la banda completa, manteniendo \eqref{hmtfg:escalado:eq:61}.

\subsubsection{Teorema de coincidencia del parámetro de acumulación}
\label{hmtfg:escalado:16.4}

\textbf{Teorema.} Para la familia uniforme HMT \eqref{hmtfg:escalado:eq:1}, el límite de la selección original de primeras raíces coincide con el cruce estable certificado:
\begin{equation}
\boxed{\lambda_n\nearrow a_*=\lambda_*,
\qquad 1.874038<\lambda_*<1.874039.}
\label{hmtfg:escalado:eq:63}
\end{equation}

\textbf{Demostración.} \eqref{hmtfg:escalado:eq:50}–\eqref{hmtfg:escalado:eq:52} dan \(\lambda_n\uparrow a_*\le\lambda_*\), con ambos parámetros en \(J\). Si la desigualdad fuera estricta, \eqref{hmtfg:escalado:eq:54}–\eqref{hmtfg:escalado:eq:59} producirían un retorno \(f_N\in\mathcal C_-\) que realiza \(\tau\). La cobertura completa de sección~\ref{hmtfg:escalado:16.3} excluye precisamente esa posibilidad. Luego \(a_*=\lambda_*\).

La prueba controla todas las profundidades mediante una primera salida finita y una cobertura uniforme de su región. La identidad \eqref{hmtfg:escalado:eq:63} enlaza el límite del selector global con el cruce transversal.


\subsection{Identificación de los centros originales y escalado global}
\label{hmtfg:escalado:17}

El último paso conserva la selección de \eqref{hmtfg:escalado:eq:50}. La continuación local y los parámetros originales se comparan por su período exacto y su clase híbrida; la coincidencia se obtiene mediante una unicidad en un entorno paramétrico fijo.

\subsubsection{Realización analítica de grado dos y transversalidad}
\label{hmtfg:escalado:17.1}

El punto fijo \(g\) de sección~\ref{hmtfg:escalado:3} es par, analítico, de crítico cuadrático y combinatoria estacionaria de duplicación. El teorema 2.1 de \href{https://annals.math.princeton.edu/wp-content/uploads/annals-v164-n3-p01.pdf}{de Faria–de Melo–Pinto}, aplicado a esa única combinatoria, da un conjunto límite de un solo elemento, con extensión holomorfa propia de grado dos entre discos topológicos anidados. Su atracción sobre el intervalo identifica ese elemento con \(g\), pues \(Rg=g\).

El lema 3.2 de la misma fuente extiende un número fijo de retornos a una vecindad analítica completa de \(g\), con imágenes de grado dos y módulo positivo. La norma \eqref{hmtfg:escalado:eq:5} controla la norma uniforme sobre una vecindad compleja suficientemente pequeña de \([-1,1]\), porque \(|(x^2-1)/2.5|<1\) allí. La convergencia \eqref{hmtfg:escalado:eq:3} entra en dicha vecindad; la dependencia analítica de los retornos proporciona, para cierto entero fijo \(d\), una familia
\begin{equation}
\mathcal F_\lambda=R^d f_\lambda
\label{hmtfg:escalado:eq:64}
\end{equation}
de mapas de grado dos, analítica en un entorno complejo de \(\lambda_*\). Los discos y el número de retornos se fijan antes de variar \(\lambda\) en ese entorno.

La tangente paramétrica pertenece al cono expansivo, mientras la hoja estable satisface la pendiente de sección~\ref{hmtfg:escalado:5}. Los retornos preservan esa separación. La realización de grado dos identifica localmente la hoja estable con la clase híbrida del mapa infinitamente renormalizable; por ello \eqref{hmtfg:escalado:eq:64} es transversal a esa clase.

El cambio de espacio analítico se justifica detalladamente en \ref{hmtfg:clasificacion:8.2}, dedicada al transporte de la transversalidad. Una dirección tangente a la clase híbrida produciría tangentes renormalizadas decrecientes, por la convergencia uniforme sobre un disco híbrido y la estimación de Cauchy. La dirección certificada satisface, en cambio, \(|\dot f(1)|=|\dot u|/10\) y crece bajo los retornos por el cono expansivo. Esta evaluación común de ambos representantes prueba la transversalidad sin suponer equivalencia de sus normas.

\subsubsection{Unicidad en una vecindad fija}
\label{hmtfg:escalado:17.2}

Para \(n>d\), la clasificación de sección~\ref{hmtfg:escalado:15} da
\begin{equation}
K(\mathcal F_{\lambda_n})=W_{n-d}0.
\label{hmtfg:escalado:eq:65}
\end{equation}
El enderezamiento de un mapa de grado dos conserva la órbita crítica; \eqref{hmtfg:escalado:eq:65} identifica su clase híbrida con la del centro cuadrático canónico de período \(2^{n-d}\). Las cotas a priori de la sucesión real de duplicación son las empleadas en el teorema de universalidad de Lyubich.

El \href{https://www.maths.tcd.ie/EMIS/journals/Annals/149_2/lyubich.pdf\#page=69}{teorema 7.4 de Lyubich, pp. 387–388} da una sola intersección con cada una de esas clases, a niveles suficientemente altos, sobre la transversal analítica y dentro de una vecindad fija del parámetro de acumulación.

Por \eqref{hmtfg:escalado:eq:63}, \(\lambda_n\) entra finalmente en esa vecindad. La continuación local de sección~\ref{hmtfg:escalado:8}, indexada por el mismo período, pertenece a la misma clase híbrida. La unicidad implica
\begin{equation}
\boxed{\lambda_n=\widehat\mu_n\quad(n\ge n_0),}
\label{hmtfg:escalado:eq:66}
\end{equation}
donde \(\widehat\mu_n\) designa el centro local de período físico \(2^n\). Es una reindexación de \(\mu_j\) mediante un desplazamiento fijo determinado por los retornos y el blanco de sección~\ref{hmtfg:escalado:8}. La selección original de \(\lambda_n\) permanece intacta.

El entero \(n_0\) es existencial en esta prueba. La continuación de todos los niveles y la palabra \(W_n0\) ya están demostradas en \eqref{hmtfg:escalado:eq:50}; \eqref{hmtfg:escalado:eq:66} aporta la identidad métrica necesaria para el límite.

\subsubsection{Teorema de escalado para las primeras raíces HMT}
\label{hmtfg:escalado:17.3}

\textbf{Teorema.} Para la familia uniforme \eqref{hmtfg:escalado:eq:1}, sean \(\lambda_n\) las sucesivas primeras raíces nuevas de período crítico exacto \(2^n\). Existen \(C>0\), \(0<\theta<1\) y \(M<\infty\) tales que, para todo \(n\) suficientemente grande,
\begin{equation}
\lambda_*-\lambda_n=C\delta_{\mathrm F}^{-n}(1+e_n),
\qquad |e_n|\le M\theta^n,
\label{hmtfg:escalado:eq:67}
\end{equation}
y
\begin{equation}
\boxed{\lim_{n\to\infty}
\frac{\lambda_{n-1}-\lambda_{n-2}}
{\lambda_n-\lambda_{n-1}}=\delta_{\mathrm F}.}
\label{hmtfg:escalado:eq:68}
\end{equation}

\textbf{Demostración.} \eqref{hmtfg:escalado:eq:66} transporta el escalado local \eqref{hmtfg:escalado:eq:24} a las raíces originales; el desplazamiento finito de índices se absorbe en \(C,M\). Su signo es positivo porque \(\lambda_n<\lambda_*\).

El resto de potencia admite además una justificación directa mediante la holonomía de clases híbridas. El lema 7.3 de Lyubich da regularidad \(C^{1+\beta}\), con derivada no nula, sobre la transversal en el punto de acumulación. La dinámica unidimensional inestable es analítica y tiene multiplicador \(\delta_{\mathrm F}>1\); su coordenada linealizante hace que los centros sucesivos estén a distancia proporcional a \(\delta_{\mathrm F}^{-n}\). Componer con la inversa de la holonomía da un término lineal no nulo y un resto \(O(\delta_{\mathrm F}^{-(1+\beta)n})\). Así se puede tomar \(\theta=\delta_{\mathrm F}^{-\beta}\).

Escribiendo
\[
r_n=\frac{\delta_{\mathrm F} e_{n-1}-e_n}{\delta_{\mathrm F}-1},
\qquad K=\frac{M(\delta_{\mathrm F}+\theta)}{\delta_{\mathrm F}-1},
\]
se obtiene
\[
\lambda_n-\lambda_{n-1}
=C(\delta_{\mathrm F}-1)\delta_{\mathrm F}^{-n}(1+r_n),\qquad
|r_n|\le K\theta^{n-1}.
\]
Por tanto,
\begin{equation}
\left|
\frac{\lambda_{n-1}-\lambda_{n-2}}
{\lambda_n-\lambda_{n-1}}-\delta_{\mathrm F}
\right|
\le
\frac{\delta_{\mathrm F} K(1+\theta)\theta^{n-2}}
{1-K\theta^{n-1}},
\quad K\theta^{n-1}<1.
\label{hmtfg:escalado:eq:69}
\end{equation}
El lado derecho tiende a cero, demostrando \eqref{hmtfg:escalado:eq:68}.

\subsubsection{Encadenamiento y alcance del cierre}
\label{hmtfg:escalado:17.4}

La composición efectiva es: lector dodecafásico APP–TRIT–TPK; perfil \eqref{hmtfg:escalado:eq:1}; transporte ordenado de retornos \eqref{hmtfg:escalado:eq:38}–\eqref{hmtfg:escalado:eq:40}; clasificación y existencia de las primeras raíces \eqref{hmtfg:escalado:eq:50}; localización \eqref{hmtfg:escalado:eq:52}; aislamiento por primera salida \eqref{hmtfg:escalado:eq:63}; identidad eventual de centros \eqref{hmtfg:escalado:eq:66}; escalado \eqref{hmtfg:escalado:eq:67}–\eqref{hmtfg:escalado:eq:69}. La representación de Jacobi \eqref{hmtfg:escalado:eq:43} conserva el perfil completo de este lector y sus marcas permanecen en el antecedente.

La expansión espacial del lector HMT y el refinamiento nonádico aportan la genealogía y la evaluación a profundidad arbitraria. El escalado entre parámetros se demuestra mediante la dinámica del retorno, la transversalidad y el aislamiento que acabamos de componer. Las pruebas analíticas citadas intervienen sobre la familia ya construida; sus valores de comparación no seleccionan sus fases ni sus coeficientes.

El cierre se refiere al selector global original del sector uniforme \(\eta=0\), con las dependencias teoremáticas enumeradas. La ley asintótica conserva sus cuantificadores existenciales. La anchura de una caja de raíces finitas mide precisión de evaluación; la estimación \eqref{hmtfg:escalado:eq:69} expresa la convergencia hacia el límite.


Los programas y certificados del suplemento de reproducción permiten reproducir las partes finitas de la prueba. La continuación a toda profundidad se sustenta en las inducciones, el argumento uniforme de primera salida y los teoremas de identificación, con sus dominios explícitos.

% Fragmento integrable. Preambulo anfitrion: amsmath, amssymb, longtable, hyperref.
\section{Identificación analítica de los centros originales}
\label{hmtfg:fragmento:identificacion-analitica-de-los-centros-originales}

\subsection{Realización de grado dos e identificación métrica de los centros}
\label{hmtfg:clasificacion:8}

La igualdad entre el primer parámetro del itinerario infinito y el cruce estable, demostrada por primera salida, proporciona
\begin{equation}
\boxed{a_*:=\min\Lambda_\tau=\lambda_*.}
\label{hmtfg:clasificacion:eq:H}
\end{equation}

La familia, el perfil dodecafásico, los retornos y el selector se mantienen exactamente como en secciones~\ref{hmtfg:clasificacion:1}--\ref{hmtfg:clasificacion:5}. Los resultados externos empleados a continuación reconocen la dinámica de ese objeto ya construido. Sus valores universales no intervienen en la selección de coeficientes ni de parámetros de la familia HMT.

\textbf{Teorema de transferencia métrica.} La identificación \eqref{hmtfg:clasificacion:eq:H}, junto con las cotas de entrada, admisibilidad y transversalidad anteriormente demostradas, implica que la sucesión original de primeras raíces nuevas satisface, para ciertas constantes \(C>0\) y \(0<\theta<1\),
\begin{equation}
\boxed{\lambda_*-\lambda_n=C\delta_{\mathrm F}^{-n}
\bigl(1+O(\theta^n)\bigr),}
\label{hmtfg:clasificacion:eq:14}
\end{equation}
donde \(\delta_{\mathrm F}>1\) es el autovalor inestable del operador de duplicación en el punto fijo reconocido. En particular,
\begin{equation}
\frac{\lambda_{n-1}-\lambda_{n-2}}
     {\lambda_n-\lambda_{n-1}}
=\delta_{\mathrm F}+O(\theta^n).
\label{hmtfg:clasificacion:eq:15}
\end{equation}
La sucesión coincide eventualmente con los centros de la continuación local, una vez que ambos índices se refieren al mismo período físico. La prueba mantiene el selector de primera raíz de sección~\ref{hmtfg:clasificacion:4}.

\subsubsection{Obtención de un dominio cuadrático semejante común}
\label{hmtfg:clasificacion:8.1}

Escribamos \(g\) para el punto fijo de la renormalización y
\(\Gamma(\lambda)=R^4f_\lambda\). La carta analítica utilizada en el certificado es
\[
f(x)=1-x^2V\!\left(\frac{x^2-1}{2.5}\right),\qquad
V(z)=\frac{u}{10}+\sum_{k\ge1}\nu_kz^k,
\]
con norma de diferencias \(|\Delta u|+\|\Delta\nu\|_1\). Las desigualdades de admisibilidad prueban que \(g\) es par, analítico, unimodal, cuadrático y de tipo de duplicación; además, \(g(x)=\varphi(x^2)\) con \(\varphi'<0\) en \([0,1]\).

El teorema 2.1 de de Faria–de Melo–Pinto, aplicado al conjunto de combinatorias formado únicamente por la duplicación, proporciona un conjunto límite de un solo punto con extensión cuadrática semejante. Como \(Rg=g\), la convergencia real de ese teorema identifica dicho punto con \(g\). El lema 3.2 proporciona, en un entorno analítico de \(g\), una iteración finita de renormalización con dominio cuadrático semejante y módulo uniformemente positivo. Los localizadores son secciones 2.1–2.3, pp. 736–738, y sección 3.1, lema 3.2, p. 745, de \href{https://annals.math.princeton.edu/wp-content/uploads/annals-v164-n3-p01.pdf}{de Faria–de Melo–Pinto, \emph{Global hyperbolicity of renormalization for \(C^r\) unimodal mappings}}.

Para comprobar la pertenencia al entorno de ese lema, sea \(\Omega_\epsilon\) un entorno complejo suficientemente pequeño de \([-1,1]\), con
\[
\rho_\epsilon:=\sup_{x\in\Omega_\epsilon}
\left|\frac{x^2-1}{2.5}\right|<1,
\qquad M_\epsilon:=\sup_{x\in\Omega_\epsilon}|x|^2<\infty.
\]
La restricción de la carta certificada satisface
\begin{equation}
\|\Delta f\|_{\Omega_\epsilon}
\le M_\epsilon\left(\frac{|\Delta u|}{10}
          +\rho_\epsilon\|\Delta\nu\|_1\right)
\le C_\epsilon\|\Delta f\|_{\mathcal A}.
\label{hmtfg:clasificacion:eq:16}
\end{equation}
El certificado da
\[
\|R^m\Gamma(\lambda_*)-g\|_{\mathcal A}
<0.001666(0.83996)^m.
\]
Se elige por ello un entero fijo \(m\) que sitúe esta imagen en el entorno del lema 3.2. La continuidad de la composición finita proporciona un intervalo \(I_*\) alrededor de \(\lambda_*\) donde las mismas operaciones están definidas. Si \(N\) es la iteración suministrada por dicho lema, ponemos
\begin{equation}
d=4+m+N,\qquad
\mathcal F_\lambda=R^df_\lambda,\quad \lambda\in I_*.
\label{hmtfg:clasificacion:eq:17}
\end{equation}
La familia \(\mathcal F\) posee representantes cuadráticos semejantes de grado dos, sobre un dominio común y con módulo positivo uniforme. Su dependencia analítica real admite una complejificación local del parámetro. El número \(d\) es fijo y cuenta las duplicaciones realizadas antes de aplicar el resultado paramétrico.

Esta construcción fija un dominio común y un número finito de retornos \(d\), suficientes para aplicar el teorema paramétrico.

\subsubsection{Transporte de la transversalidad certificada}
\label{hmtfg:clasificacion:8.2}

La dirección de \(\Gamma\) satisface el cono
\[
\|\dot\nu\|_1\le\tfrac14|\dot u|,\qquad \dot u\ne0.
\]
Las cotas diferenciales certificadas implican invariancia de este cono y
\begin{equation}
|\dot u_{k+1}|\ge4.381|\dot u_k|
\label{hmtfg:clasificacion:eq:18}
\end{equation}
a lo largo de la órbita de \(\Gamma(\lambda_*)\). Todos sus puntos permanecen en el dominio certificado. En particular, la dirección transportada por los pasos finitos de \eqref{hmtfg:clasificacion:eq:17} conserva una componente expansiva no nula.

La clase híbrida del punto fijo coincide con su variedad estable en el espacio de gérmenes cuadráticos semejantes: \href{https://www.maths.tcd.ie/EMIS/journals/Annals/149_2/lyubich.pdf}{Lyubich, \emph{Feigenbaum–Coullet–Tresser universality and Milnor's hairiness conjecture}, teorema 6.1, pp. 371–372}. La transversalidad sigue siendo válida al pasar a ese espacio, como puede comprobarse sin identificar normas de Banach diferentes. En efecto, si \(\partial_\lambda\mathcal F_{\lambda_*}\) fuese tangente a la clase híbrida, sería la tangente de un disco analítico contenido en ella. La convergencia exponencial uniforme de la renormalización sobre un subdisco compacto, seguida de la estimación de Cauchy en su parámetro, haría converger a cero la evaluación en \(x=1\) de sus tangentes renormalizadas. Sin embargo, en la carta anterior,
\[
\dot f(1)=-\dot u/10,
\]
y \eqref{hmtfg:clasificacion:eq:18} hace crecer su valor absoluto. Ambas derivadas representan el mismo germen y tienen la misma evaluación real. La contradicción demuestra
\begin{equation}
\partial_\lambda\mathcal F_{\lambda_*}
\notin T_{\mathcal F_{\lambda_*}}\mathcal H_{c_\infty}.
\label{hmtfg:clasificacion:eq:19}
\end{equation}
Así, \(S=\{\mathcal F_\lambda\}\) es una transversal analítica compleja a la clase híbrida del límite de duplicación.

\subsubsection{Identificación eventual de las primeras raíces por su período}
\label{hmtfg:clasificacion:8.3}

Denotemos por \(c_r\) el centro cuadrático canónico de período \(2^r\), obtenido mediante \(r\) duplicaciones sucesivas del centro de período uno. Este símbolo designa el parámetro de la familia cuadrática de reconocimiento; se distingue de las funciones orbitales \(c_j(\lambda)=f_\lambda^j(0)\) de las secciones anteriores. Se escribe \(c_r^{\rm quad}\) para conservar la distinción.

El teorema de continuación global y \eqref{hmtfg:clasificacion:eq:H} dan \(\lambda_n\to\lambda_*\), por lo que \(\lambda_n\in I_*\) para \(n\) suficientemente grande. La clasificación e intervalos restrictivos de sección~\ref{hmtfg:clasificacion:2} muestran que, cuando \(n>d\),
\begin{equation}
\mathcal F_{\lambda_n}=R^df_{\lambda_n}
\quad\text{tiene período crítico exacto }2^{n-d}
\quad\text{e itinerario }W_{n-d}0.
\label{hmtfg:clasificacion:eq:20}
\end{equation}
Su órbita crítica permanece en el intervalo real invariante, contenido en el dominio cuadrático semejante; su conjunto de Julia lleno es, por tanto, conexo. La rectificación conserva la órbita postcrítica y las ramas reales. Al recorrer sus retornos restrictivos se obtienen \(n-d\) duplicaciones canónicas y, al final, un punto crítico fijo. Este último rectifica al centro \(0\); las inversas sucesivas de la rectificación de las copias de duplicación dan el centro único \(c_{n-d}^{\rm quad}\). Esta es la aplicación concreta de la parametrización por combinatoria y del homeomorfismo de rectificación de las copias, descritos en Lyubich, secciones 5.1–5.3, pp. 362–364. Se concluye
\begin{equation}
\mathcal F_{\lambda_n}\in\mathcal H_{c_{n-d}^{\rm quad}}.
\label{hmtfg:clasificacion:eq:21}
\end{equation}

El teorema 7.4 de Lyubich, pp. 387–388, da una única intersección local de \(S\) con cada clase \(\mathcal H_{c_r^{\rm quad}}\) para \(r\) suficientemente grande. Sus hipótesis de cotas complejas a priori se cumplen para la cascada cuadrática real de duplicación, conforme al teorema 5.6, p. 367, del mismo trabajo. Escribamos \(\zeta_r\) para el parámetro de esa intersección. Las ecuaciones \eqref{hmtfg:clasificacion:eq:20}–\eqref{hmtfg:clasificacion:eq:21}, la convergencia al cruce y la unicidad local implican
\begin{equation}
\boxed{\lambda_n=\zeta_{n-d}\quad\text{para todo }n\text{ suficientemente grande}.}
\label{hmtfg:clasificacion:eq:22}
\end{equation}
La unicidad utilizada aquí corresponde a \textbf{cada clase híbrida canónica} en un entorno fijo del cruce. No se infiere de la mera existencia de una cascada local ni de la unicidad del cruce con la hoja estable. La ecuación \eqref{hmtfg:clasificacion:eq:22} identifica el selector original, ya definido globalmente, con las intersecciones locales; no introduce un selector sustitutorio.

En particular, sea \(\mu_j\) la cascada local de la sección \ref{hmtfg:escalado:8}, y sea \(s\) el entero fijo determinado por su período físico exacto:
\[
\operatorname{per}_{\rm crit}(f_{\mu_j})=2^{j+s}.
\]
En la presentación con curva \(R^4f_\lambda\) y blanco de período dos, \(s=5\); los pasos fijos adicionales usados para localizar el blanco modifican \(s\). Por \eqref{hmtfg:clasificacion:eq:21} y la misma unicidad,
\begin{equation}
\mu_j=\zeta_{j+s-d},\qquad
\boxed{\lambda_n=\mu_{n-s}\quad(n\text{ suficientemente grande}).}
\label{hmtfg:clasificacion:eq:23}
\end{equation}
El desplazamiento \(d\) de la preparación cuadrática semejante y el desplazamiento \(s\) de los períodos del blanco son objetos distintos. La identificación se efectúa por el período físico, no igualando esos dos enteros.

\subsubsection{Resto de potencia mediante holonomía transversal}
\label{hmtfg:clasificacion:8.4}

El lema 7.3 de Lyubich, p. 384 y demostración en pp. 384–387, aporta regularidad \(C^{1+\beta}\)-conforme, para algún \(\beta>0\), de la holonomía entre transversales a la clase híbrida del punto de Feigenbaum. Su definición en p. 384 se refiere a los conjuntos de conexidad sobre las transversales; contiene, en particular, los centros considerados aquí. Podemos reducir el exponente para que \(0<\beta\le1\).

Sea \(h\) la holonomía desde \(S\) hacia la variedad inestable y sea \(z\) una coordenada analítica que linealiza la restricción unidimensional de la renormalización:
\begin{equation}
z(g)=0,\qquad z(Rq)=\delta_{\mathrm F} z(q).
\label{hmtfg:clasificacion:eq:24}
\end{equation}
La existencia de esta coordenada se obtiene aplicando la linealización de un germen holomorfo repulsor, de multiplicador \(\delta_{\mathrm F}>1\), a la variedad inestable analítica. La dirección expansiva de \eqref{hmtfg:clasificacion:eq:18}, la invariancia del cono y el retorno estacionario identifican ese multiplicador con el autovalor inestable positivo empleado en el certificado.

Si \(q_r\) es el centro de la clase \(\mathcal H_{c_r^{\rm quad}}\) situado en la variedad inestable, la compatibilidad de rectificación y renormalización da \(Rq_r=q_{r-1}\). Esta identidad es la que se utiliza en la demostración del teorema 7.4, p. 388. Por \eqref{hmtfg:clasificacion:eq:24}, existe \(b\ne0\), que absorbe cualquier índice inicial fijo, tal que
\begin{equation}
z(q_r)=b\delta_{\mathrm F}^{-r}
\label{hmtfg:clasificacion:eq:25}
\end{equation}
para todo \(r\) suficientemente grande.

La regularidad de \(h\) y la transversalidad \eqref{hmtfg:clasificacion:eq:19} dan, sobre el conjunto de conexidad cercano a \(\lambda_*\),
\begin{equation}
H(\lambda):=z\bigl(h(\mathcal F_\lambda)\bigr)
=A(\lambda-\lambda_*)
+O\!\left(|\lambda-\lambda_*|^{1+\beta}\right),
\qquad A\ne0.
\label{hmtfg:clasificacion:eq:26}
\end{equation}
Como \(H(\zeta_r)=b\delta_{\mathrm F}^{-r}\), la estimación inferior
\(|H(\lambda)|\ge |A|\,|\lambda-\lambda_*|/2\) en un entorno suficientemente pequeño implica primero
\(|\zeta_r-\lambda_*|=O(\delta_{\mathrm F}^{-r})\). Sustituyendo esta cota en el resto de \eqref{hmtfg:clasificacion:eq:26}, resulta
\begin{equation}
\zeta_r-\lambda_*=\frac bA\delta_{\mathrm F}^{-r}
+O\!\left(\delta_{\mathrm F}^{-(1+\beta)r}\right).
\label{hmtfg:clasificacion:eq:27}
\end{equation}
Aplicamos \eqref{hmtfg:clasificacion:eq:22} con \(r=n-d\). El desplazamiento finito modifica el coeficiente principal y la constante del resto. Puesto que \(\lambda_n<\lambda_*\), se obtiene \eqref{hmtfg:clasificacion:eq:14}, con
\[
C=-\frac bA\delta_{\mathrm F}^d>0,
\qquad \theta=\delta_{\mathrm F}^{-\beta}\in(0,1).
\]
Finalmente,
\[
\lambda_n-\lambda_{n-1}
=C(\delta_{\mathrm F}-1)\delta_{\mathrm F}^{-n}\bigl(1+O(\theta^n)\bigr),
\]
lo que demuestra \eqref{hmtfg:clasificacion:eq:15}. \(\square\)


\subsubsection{Alcance cuantitativo del teorema}
\label{hmtfg:clasificacion:8.5}

La identificación de límites demostrada por primera salida activa la transferencia métrica para las primeras raíces globales del sector uniforme. La tasa operatorial \(0.83996\) controla los mapas renormalizados sobre la hoja estable. La tasa paramétrica \(\theta=\delta_{\mathrm F}^{-\beta}\) procede de la regularidad de la holonomía transversal. Ambas estimaciones conservan su variable y su dominio.

El teorema acredita la existencia de \(\beta,\theta,C\), de la vecindad \(I_*\), del número finito de retornos \(d\) y de un índice \(n_0\) a partir del cual coinciden los centros. Sus valores numéricos individuales no se asignan en este enunciado asintótico. Esta formulación completa el límite y su resto de potencia; la evaluación numérica de esos testigos constituye una especialización cuantitativa distinta.

\section{Retorno, memoria y universalidad de la cascada dodecafásica}
\label{sec:hmt-feig-interpretacion}

\subsection{La estructura generadora y la lectura crítica}

La construcción de Feigenbaum en Holografía Modular Triádica enlaza una
organización aritmética orientada con una ley asintótica de cambio de escala.
Su punto de partida es la composición efectiva de APP, TRIT y TPK. APP evalúa
suma y producto sobre las dos hojas de la retícula de dígitos y conserva los
cocientes y residuos de esas operaciones. TRIT determina el régimen local y
la orientación de su lectura. TPK selecciona y transporta los estados,
actualiza sus registros y prolonga sus rutas con acarreo, frontera y memoria.
La estructura discreta conjunta del continuo proporciona la residencia común
de esos estados y de los lectores que actúan sobre ellos.

La prolongación
\[
w_6\longrightarrow w_{12}\longrightarrow w_{18}
\longrightarrow w_{24}\longrightarrow w_{30}
\longrightarrow R_{36}\longrightarrow G_9
\]
organiza esta persistencia. La semilla inicial se eleva conservando sus
prefijos; las elevaciones sucesivas cambian el régimen sin perder orientación
ni memoria; la frontera coinductiva permite continuar la construcción, y las
nueve cartas reúnen sus lecturas locales. La holonomía nonádica expresa el
retorno de fase acompañado de un incremento de memoria. Así, una vuelta
visible conserva la posición que ocupa dentro de la historia completa. La
operación de retorno adquiere desde el origen dos componentes inseparables:
reiteración de una regla y conservación de su procedencia.

El lector crítico fijado sobre la rueda dodecafásica orientada convierte esta
estructura en doce fases con una traza uniforme. Sus representantes son
\[
\phi_m=\frac{(3m-1)\pi_{\mathrm{HMT}}}{18},
\qquad w_m=\frac1{12},\qquad 1\le m\le12.
\]
La elección del representante orientado importa porque el segundo momento
utiliza su magnitud. El cociente angular registra la fase, mientras el
representante conserva la información necesaria para normalizarla. La
coordenada \(\pi_{\mathrm{HMT}}\) procede de los lectores coinductivos del mismo
núcleo APP--TRIT--TPK. Dentro de su imagen correlacionada, las coordenadas
\(\pi,\varphi,e\) y el cierre dodecafásico de \(\alpha_{\mathrm{HMT}}\)
conservan su orden constructivo y sus dominios propios. En el lector crítico
aquí considerado se utiliza la coordenada angular ya generada.

La normalización se determina mediante una suma entera:
\[
\sum_{m=1}^{12}(3m-1)^2=5394=6\cdot899,
\qquad k_m=\frac{2(3m-1)}{\sqrt{899}}.
\]
El factor común \(\pi_{\mathrm{HMT}}\) se cancela al fijar el segundo momento.
Esta cancelación conserva la procedencia de las fases y produce una
representación particularmente precisa de su lectura normalizada. Se obtiene
\[
u_0(x)=\frac1{12}\sum_{m=1}^{12}\bigl(1-\cos(k_mx)\bigr),
\qquad f_\lambda(x)=1-\lambda u_0(x),
\qquad \sum_mw_mk_m^2=2.
\]
El perfil es par y entero; satisface \(u_0(0)=u_0'(0)=0\) y
\(u_0''(0)=2\). Para \(0<\lambda\le2/u_0(1)\), la familia lleva
\([-1,1]\) en sí mismo y tiene un máximo crítico cuadrático con
\(f_\lambda''(0)=-2\lambda\). El parámetro \(\lambda\) recorre esta familia
después de fijados el lector, la orientación, los pesos y la escala. La cifra
universal que se reconocerá al final queda así separada de la selección de
los coeficientes.

\subsection{Doce fases, momentos y representación espectral}

El perfil admite una segunda expresión que reúne sus fases en un operador
finito. Los nodos
\[
s_m=k_m^2=\frac{4(3m-1)^2}{899}
\]
determinan la medida positiva \(\mu_0=12^{-1}\sum_m\delta_{s_m}\). Sus
momentos \(M_r=\int s^r\,d\mu_0(s)\) definen formas de Hankel estrictamente
positivas hasta el rango doce. Esta positividad procede de la evaluación de
un polinomio sobre doce nodos distintos: un polinomio de grado menor que
doce conserva al menos una evaluación no nula. En grado doce, el producto
de los doce factores nodales anula la norma y determina exactamente el rango
del lector.

La ortogonalización de estos momentos produce una matriz de Jacobi
\(J_{12}\), real simétrica, con subdiagonales positivas. Su cálculo funcional
recupera el perfil completo:
\[
u_0(x)=e_0^{\mathsf T}
\bigl[I-\cos(x\sqrt{J_{12}})\bigr]e_0.
\]
La suma de cosenos y el operador de Jacobi son dos realizaciones del mismo
lector. Los momentos mantienen todas las frecuencias y sus pesos; las
marcas de orientación, ruta y memoria permanecen en el estado que las
genera. El rango doce identifica la dimensión espectral de esta lectura
concreta. La profundidad de sus prolongaciones pertenece a otra dimensión
de la construcción: la longitud de las historias compatibles y la resolución
con que se evalúan.

La geometría analítica también se hace explícita. Para la deformación
ponderada \(u_\eta\), en la franja \(|\eta|\le1/40\), los pesos siguen siendo
positivos y el segundo momento permanece normalizado. En
\(|\operatorname{Re}z|<\pi/3\) existe una función holomorfa impar y univalente
\(b_\eta\), normalizada por \(b_\eta'(0)=1\), tal que
\(u_\eta=b_\eta^2\). La rueda produce, por tanto, una deformación conforme
del pliegue cuadrático. Su conjugación dinámica conserva esa deformación a
través de \(b_\eta(1-\lambda w^2)\). El teorema de escalado paramétrico que
se desarrolla aquí corresponde al miembro uniforme \(\eta=0\); las
propiedades estructurales de la franja ponderada y la persistencia local de
sus raíces conservan sus cuantificadores específicos.

\subsection{Profundidad de lectura y restitución de las historias}

La profundidad nonádica refina la lectura de un parámetro, mientras la
duplicación dinámica cambia el horizonte de retorno. El refinamiento con
acarreo expresa la primera operación mediante
\[
L_c(x,y)=(Bx-c,By+c),\qquad
C_m=\sum_{j=1}^m c_jB^{m-j},
\]
\[
(x_m,y_m)=(B^mx_0-C_m,B^my_0+C_m).
\]
El soporte exterior crece como \(B^m\), y el cilindro normalizado del registro
tiene diámetro \(B^{-m}\). La pareja entera, la profundidad y la marca de
frontera conservan la reconstrucción de cocientes, residuos y prefijos.
Crecimiento exterior y resolución interior describen así operaciones
compatibles sobre una misma historia.

La composición exacta
\(L_D^{[b]}L_C^{[a]}=L_{B^bC+D}^{[a+b]}\) permite agrupar una historia en
bloques sin alterar sus datos. Un retorno de duplicación reúne dos
transiciones; el calendario nonádico reúne nueve. Ambos agrupamientos se
comparan sobre un horizonte común de \(9\,2^N\) eventos, conservando el
orden de los subbloques y sus acarreos.

La contribución de la memoria puede observarse directamente en las ramas
inversas del perfil. Si \(v=(u_0|_{[0,1]})^{-1}\), las dos ramas son
\[
\beta_\sigma(y)=\sigma v\!\left(\frac{1-y}{\lambda}\right),
\qquad \sigma\in\{-1,+1\}.
\]
El valor final, junto con la última hoja, reconstruye el valor anterior. Una
sucesión de hojas permite recuperar todos los pasos intermedios. La palabra
de ramas acompaña los registros APP--TRIT--TPK de ruta, frontera, acarreo y
memoria; cada registro conserva su función. Esta restitución concreta
explica qué información requiere la inversión de la lectura crítica.

La resolución de cada retorno dispone además de una cota explícita. Con
\(t=\lambda u_\eta(1)/2\), \(z=(x+1)/2\) y
\[
F_{t,\eta}(z)=1-t\frac{u_\eta(2z-1)}{u_\eta(1)},
\]
se obtiene
\[
\left|F_{t,\eta}^{q}(1/2)-F_{s,\eta}^{q}(1/2)\right|
\le S_q|t-s|,\qquad S_q=\sum_{j=0}^{q-1}(6\sqrt2)^j<9^q.
\]
Un cilindro de base nueve y profundidad \(m=2^N+r\) permite, por ello, leer
el retorno de horizonte \(2^N\) con diámetro inferior a \(9^{-r}\). Esta
estimación da contenido cuantitativo a la profundidad arbitraria: para cada
horizonte y precisión solicitados determina una profundidad suficiente. La
evaluación de las funciones analíticas y el redondeo poseen sus propios
restos, que se propagan junto con la incertidumbre del parámetro.

Las vacancias intervienen en el balance de capacidad del registro. La
comparación entre bloques de 729 y 1000 posibilidades distingue capacidad
acumulada y longitud cronológica, conservando los pasos en que la primera
permanece constante. Esta contabilidad determina el espacio requerido para
conservar una lectura y su profundidad. La admisibilidad de las historias
procede, a su vez, de las reglas de supervivencia. Ambas operaciones
colaboran en la construcción y mantienen objetos matemáticos distintos.

\subsection{Cronología nonádica y combinatoria de duplicación}

La palabra crítica de duplicación tiene una expresión a toda profundidad:
\[
\tau_j=(-1)^{v_2(j)},\qquad
\tau_{2j}=-\tau_j,\qquad \tau_{2j+1}=+1.
\]
Sus prefijos satisfacen \(W_{n+1}=W_n(-1)^nW_n\). La regla determina cómo
se orientan las dos copias del retorno y dónde aparece su separación. La
lectura nonádica de la cronología conserva el entero completo
\(K=\rho_9(K)+9q_9(K)\). Al reducirlo módulo \(2^n\), la holonomía que
avanza nueve eventos induce la traslación por nueve sobre ese ciclo. Puesto
que nueve es impar, esta traslación visita todas sus posiciones. La
permutación \(j\mapsto9j\pmod{2^n}\), que la conjuga con el avance unitario,
preserva la valoración diádica de cada posición no nula. El residuo y el cociente
cronológicos sostienen conjuntamente esa compatibilidad.

La memoria bilateral posee asimismo una realización operatoria compatible
con el refinamiento. Sobre los ciclos de longitud \(2^n\), el operador de
avance \(C_n\) se incorpora en
\[
\mathbb U_n=P_0\otimes I+P_1\otimes C_n,
\]
con los proyectores complementarios del lector nonádico. En el límite, sus
componentes visible y complementaria satisfacen
\[
\|T_{\infty,a}v\|^2+\|D_{\infty,a}v\|^2=\|v\|^2,
\qquad
v=T_{\infty,a}^*T_{\infty,a}v+D_{\infty,a}^*D_{\infty,a}v.
\]
Estas identidades describen conservación y restitución dentro del espacio
y la representación declarados. La contracción de los mapas renormalizados,
que aparecerá después, mide otra operación: la aproximación de sus perfiles
a una forma estable. La historia conservada por el levantamiento y la
convergencia de una lectura funcional se articulan en niveles diferentes.

\subsection{La carta ordenada y la selección global de las raíces}

La publicación arquimediana del continuo HMT transporta las operaciones
mediante un isomorfismo de cuerpos ordenados completos, dentro del universo
declarado \(\mathfrak G\):
\[
\iota:\mathbb R_{\mathrm{HMT}}^{\mathfrak G}
\longrightarrow\mathbb R^{\mathfrak G}.
\]
La construcción interna del coseno por su serie, la raíz positiva de 899,
el perfil y sus iterados precede a la selección de ceros. La conservación
de suma, producto, orden y límites convergentes da
\[
\iota\bigl(S_n^{\mathrm H}(a)\bigr)=S_n\bigl(\iota(a)\bigr),
\qquad S_n(\lambda)=f_\lambda^{2^n}(0).
\]
La sobreyectividad cubre todas las raíces del dominio de esa carta; el orden
conserva los retornos anteriores, el período exacto y la condición de
encontrarse a la derecha de la raíz precedente. De este modo se transporta
también el mínimo de cada conjunto de candidatos. La genealogía llega
hasta la selección efectiva del parámetro, con la segunda proyección y sus
marcas conservadas sobre el mismo antecedente.

El selector queda fijado por \(\lambda_1=1/u_0(1)\) y por la primera raíz
nueva de período crítico exacto \(2^n\) situada a la derecha de
\(\lambda_{n-1}\). La clasificación de itinerarios demuestra que cada raíz
seleccionada tiene palabra \(W_n0\). Los retornos restrictivos reducen su
período a la mitad conservando el orden unimodal, y la inducción restituye
la palabra completa. La analiticidad de los retornos no idénticamente nulos
aísla sus ceros en cada compacto; la
comparación de los signos antes del primer parámetro de itinerario infinito
garantiza la existencia de una raíz nueva en todos los niveles.

Sea \(\Lambda_\tau\) el conjunto compacto no vacío de parámetros que realizan
la palabra infinita, y sea \(a_*\) su mínimo. Se demuestra
\[
\lambda_n\nearrow a_*.
\]
El paso al límite conserva cada signo mediante una cota uniforme a horizonte
fijo. Si \(c_p=f_\lambda^p(0)\), los signos opuestos de \(c_p\) y
\(c_{2p}\), junto con \((f_\lambda^p)'(0)=0\), implican
\[
|c_p|>\frac2{B_p},\qquad
B_p=16^p\frac{16^p-1}{15}.
\]
Así, los prefijos de períodos crecientes preservan la palabra infinita al
acumularse. La prueba utiliza la sucesión original de primeras raíces desde
su definición global.

\subsection{La hoja estable y la identificación del límite seleccionado}

El retorno normalizado
\[
(Rf)(x)=-\frac{f(f(-ax))}{a},\qquad a=-f(1)>0,
\]
reúne dos transiciones, cambia la escala y conserva la orientación necesaria
para restituir un máximo de valor uno. Los dominios restrictivos acreditan
su interpretación como retorno real. La derivada de este operador incluye
la variación de la escala \(a\); esa contribución determina la sensibilidad
paramétrica de la transformación completa.

En la carta \(f(x)=1-x^2V((x^2-1)/(5/2))\), con coordenadas \(u=10v_0\)
y \(\nu=(v_1,v_2,\ldots)\), el cuarto retorno de la familia entra en un
entorno cuantificado del punto fijo \(g\). La curva cruza transversalmente
su hoja estable en un único parámetro local
\[
1.874038<\lambda_*<1.874039,
\]
y satisface
\[
\|R^{4+n}f_{\lambda_*}-g\|_{\mathcal A}
<0.001666(0.83996)^n.
\]
La hoja estable es una gráfica de pendiente a lo sumo \(0.16\). Las
secantes paramétricas pertenecen a un cono transversal expansivo. Ambas
propiedades distinguen la disminución de las deformaciones estables y el
crecimiento de la separación entre parámetros.

La identificación de \(a_*\) y \(\lambda_*\) reúne esta geometría local
con el mínimo global. Primero se excluye el intervalo entre la octava raíz
y el extremo inferior de la caja local. Después, el orden
\(a_*\le\lambda_*\) permite estudiar una hipotética separación mediante
el cono negativo. Mientras su coordenada transversal tiene módulo inferior
a \(0.006\), la secante permanece en el dominio analítico y crece por un
factor comprendido entre \(4.4034\) y \(8.0887\). Una separación persistente
produciría, por tanto, una primera salida de amplitud comprendida entre
\(0.006\) y \(0.0485322\).

El residuo de la órbita estable conserva información más precisa que su
norma: sus componentes satisfacen \(|e_u|\le0.16\|e_\nu\|_1\). Esta
relación permite acotar conjuntamente sus efectos sobre la órbita crítica.
La cobertura de toda la región de primera salida muestra en cada caso un
signo incompatible con \(\tau\), o un retorno contractivo que obliga a dos
valores críticos diádicos sucesivos a tener el mismo signo. La palabra
infinita exige signos opuestos. La contradicción identifica los parámetros:
\[
\lambda_n\nearrow a_*=\lambda_*.
\]
La demostración combina una inducción válida a toda profundidad con una
exclusión finita uniforme de la región a la que llegaría cualquier
separación. El alcance infinito descansa en esa composición.

\subsection{Escalado universal y significado de la precisión}

La realización holomorfa de grado dos de los retornos y la transversalidad
acreditada permiten comparar sus clases híbridas. Para cada período
suficientemente alto, la clase canónica posee una única intersección local
con la familia. Las primeras raíces globales ya convergen al cruce y
conservan la combinatoria de duplicación; por ello coinciden finalmente
con los centros locales al igualar sus períodos físicos. La prueba
transporta la selección original hasta la ley métrica universal.

Existen \(C>0\), \(M<\infty\) y \(0<\theta<1\) tales que
\[
\lambda_*-\lambda_n=C\delta_{\mathrm F}^{-n}(1+e_n),
\qquad |e_n|\le M\theta^n,
\]
y, en consecuencia,
\[
\lim_{n\to\infty}
\frac{\lambda_{n-1}-\lambda_{n-2}}
{\lambda_n-\lambda_{n-1}}=\delta_{\mathrm F}.
\]
El factor \(\delta_{\mathrm F}\) es el autovalor expansivo del operador de duplicación
reconocido al final de la construcción. La familia concreta conserva sus
doce fases, sus momentos y su historia de procedencia; el límite de sus
razones paramétricas pertenece a la misma clase universal de criticidad
cuadrática. Esta relación explica conjuntamente la especificidad del
generador y la universalidad de la salida.

La precisión posee aquí tres residencias. Los intervalos de las raíces
controlan su evaluación finita. La estimación
\(0.001666(0.83996)^n\) controla la convergencia operatorial sobre la hoja
estable. El término \(M\theta^n\) controla el traslado de esa geometría a
la escala paramétrica mediante la holonomía transversal. Cada estimación
responde a su variable y a su operación. El octavo cociente queda encerrado
alrededor de \(4.669212189150204154556096215889663928\), con anchura inferior
a \(5.808\cdot10^{-37}\); esa anchura mide la evaluación de \(\delta_{\mathrm F,8}\).
La distancia de \(\delta_{\mathrm F,8}\) al límite requiere cuantificar las constantes
del resto paramétrico, cuya existencia ya queda demostrada.

La notación mantiene separados \(\delta_{\mathrm F}\), factor de escalado paramétrico
de Feigenbaum; \(\alpha_{\mathrm F}\), constante espacial asociada a su
problema de renormalización; y \(\alpha_{\mathrm{HMT}}\), constante de
estructura fina del corpus. El resultado presente identifica \(\delta_{\mathrm F}\)
mediante las primeras raíces del sector uniforme \(\eta=0\). La franja
ponderada conserva sus teoremas estructurales y de existencia, y su extensión
métrica cuantificada constituye un enunciado específico distinto.

La aportación del encadenamiento reside en la continuidad entre generación,
lectura, memoria, selección y límite. La aritmética orientada determina el
perfil; la representación espectral conserva sus doce componentes; los
cilindros permiten resolverlo con precisión arbitraria a cada horizonte;
el transporte de historias restituye sus recorridos; la carta ordenada
conserva todas las raíces y sus mínimos; la clasificación prolonga el
selector; y la expansión transversal, junto con la contracción estable,
determina su ley asintótica. La universalidad aparece como una propiedad
demostrada de esta construcción completa y de su realización analítica.

\endgroup

% Fragmento compartido por VII y CRISTAL; incluir despues del cuerpo Feigenbaum.
% Procedencia: resultados recuperados, con demostraciones elementales reunidas.
% Perfil, Jacobi y escalado: 00_antecedentes_operatorios.tex, 03, 06 y 07
% del presente paquete; no se cambia su selector ni su dominio uniforme.
% Constitutiva: CRISTAL_ES/main_autosuficiente.tex:58738--58840,58864--58894.
% Inversion: CRISTAL_ES/sections/cr28_restitucion_constitutiva.tex:18--59.
% Catalan: el mismo propietario, 277--373; identidad integral, unicidad y serie.
% Corte material: CUMPLIMIENTO_EDITORIAL_20260928/fuentes/CRISTAL_ES.
% No contiene una nueva calibracion ni una nueva campana de calculo.
\section{Articulación de los lectores crítico, espectral y constitutivo}
\label{hmtfg:articulacion:seccion}

La composición APP--TRIT--TPK precede a los lectores de esta sección.
APP conserva en cada célula las evaluaciones aditiva y multiplicativa,
con sus residuos y cocientes; la reducción trítica determina el régimen,
la orientación y el acarreo de la transición. La composición efectiva
de selección, transporte y actualización de memoria produce el estado
enriquecido y sus prolongaciones, desarrollados en
\ref{hmtfg:app-trit-transporte} y
\ref{hmtfg:subsubsec:tres-vueltas-ext-ch}.
La holonomía nonádica compone esas transiciones: el retorno de fase
avanza la memoria y conserva la historia. Los lectores actúan sobre la
misma estructura discreta conjunta del continuo, con sus cinco
construcciones consustanciales y con las dos publicaciones del mismo
antecedente: evaluación ordenada e incidencia de frontera.

La rueda orientada y los canales angulares son salidas de esta
composición. Sus coordenadas HMT ya generadas se utilizan en operaciones
posteriores, con sus marcas y dominios; no son cifras metrológicas
introducidas para seleccionar una respuesta. La publicación ordenada
permite evaluar las fórmulas siguientes, mientras el estado de procedencia
conserva hoja, ruta, residuo, cociente, acarreo, frontera y memoria.
Las restituciones que se prueban recuperan el objeto de su dominio
declarado; la representación evaluada permanece acompañada de esos
registros.

\subsection{La separación crítica como lectura asintótica intrínseca}
\label{hmtfg:articulacion:separacion}

En el sector uniforme de la rueda dodecafásica, los representantes
orientados y sus pesos son
\[
 \phi_m=\frac{(3m-1)\pi_{\mathrm{HMT}}}{18},\qquad
 w_m=\frac1{12},\qquad 1\le m\le12.
\]
El segundo momento se normaliza mediante
\(\sum_{m=1}^{12}(3m-1)^2=5394=6\cdot899\).
De este modo, el factor angular común se cancela y quedan
\begin{equation}
 k_m=\frac{2(3m-1)}{\sqrt{899}},\qquad
 u_0(z)=\frac1{12}\sum_{m=1}^{12}\bigl(1-\cos(k_mz)\bigr),
 \qquad f_\lambda(z)=1-\lambda u_0(z).
 \label{hmtfg:articulacion:perfil}
\end{equation}
La variable \(\lambda\) recorre la familia después de fijados el
lector, sus pesos y la normalización. En
\(0<\lambda\le2/u_0(1)\), el dominio dinámico es \([-1,1]\).
La primera raíz es \(\lambda_1=1/u_0(1)\); para \(n\ge2\),
\(\lambda_n\) es la menor raíz mayor que \(\lambda_{n-1}\)
cuyo punto crítico tiene período exacto \(2^n\).
La clasificación, la existencia de cada mínimo y la identificación
métrica de estos centros se han demostrado en
\ref{hmtfg:escalado:15}--\ref{hmtfg:escalado:17}.

\begin{proposition}[Lectura estructural de la constante de escalado]
\label{hmtfg:articulacion:delta}
Para el lector uniforme \eqref{hmtfg:articulacion:perfil} y su selector
de primeras raíces, la constante de separación crítica es
\begin{equation}
 \delta_{\mathrm{HMT,crit}}
 :=\lim_{n\to\infty}
 \frac{\lambda_n-\lambda_{n-1}}{\lambda_{n+1}-\lambda_n}
 =\delta_{\mathrm F}.
 \label{hmtfg:articulacion:limite}
\end{equation}
Es el invariante asintótico de separación de parámetros de los retornos
críticos de ese lector, con su orden de selección conservado.
\end{proposition}
\begin{proof}
El teorema de escalado \eqref{hmtfg:escalado:eq:67} proporciona
\(\lambda_*-\lambda_n=C\delta_{\mathrm F}^{-n}(1+e_n)\),
con \(C>0\) y \(e_n\to0\). Restando dos términos consecutivos,
\[
 \lambda_n-\lambda_{n-1}
 =C\delta_{\mathrm F}^{-n}
   \bigl[\delta_{\mathrm F}-1+
          \delta_{\mathrm F}e_{n-1}-e_n\bigr].
\]
El cociente de esta igualdad y la del índice \(n+1\) converge a
\(\delta_{\mathrm F}\). La identidad eventual de los centros locales
con las primeras raíces, demostrada en
\eqref{hmtfg:escalado:eq:66}, permite aplicar el escalado precisamente
a la sucesión seleccionada y no a otra cascada elegida después.
\end{proof}

Esta formulación conserva más información que el numeral terminal:
identifica la rueda, su lectura, la familia y el selector cuyo límite
se evalúa. Su interpretación es intrínseca a esa composición HMT;
el reconocimiento de la constante universal utiliza después el teorema
analítico aplicado al lector ya construido.

\subsection{La medida finita y la restitución espectral del perfil}
\label{hmtfg:articulacion:jacobi}

Las frecuencias cuadradas producen la probabilidad atómica
\begin{equation}
 s_m=k_m^2=\frac{4(3m-1)^2}{899},\qquad
 \nu_0=\frac1{12}\sum_{m=1}^{12}\delta_{s_m},\qquad
 M_j=\int s^j\,d\nu_0(s).
 \label{hmtfg:articulacion:medida-finita}
\end{equation}
Aquí \(\delta_{s_m}\) es la masa de Dirac en un nodo, distinta
de la constante \(\delta_{\mathrm F}\). La medida conserva las
doce frecuencias y sus pesos. Para un polinomio no nulo de grado menor
que \(r\le12\),
\[
 \int p(s)^2\,d\nu_0(s)
 =\frac1{12}\sum_{m=1}^{12}p(s_m)^2>0,
\]
pues un polinomio de grado menor que doce no puede anularse en los
doce nodos distintos. Así, las matrices de Hankel
\((M_{i+j})_{0\le i,j<r}\) son positivas definidas.
El polinomio \(\prod_m(s-s_m)\) tiene norma nula; el espacio
\(L^2(\nu_0)\) tiene dimensión doce.

\begin{proposition}[Representación exacta del lector crítico]
\label{hmtfg:articulacion:representacion}
Sean \(\psi_0,\ldots,\psi_{11}\) los polinomios ortonormales
de coeficiente principal positivo para \(\nu_0\), con
\(\psi_0=1\). La multiplicación por \(s\) tiene matriz de
Jacobi \(J_{12}\), real simétrica y de subdiagonal positiva.
Con \(e_0=(1,0,\ldots,0)^{\mathsf T}\),
\begin{equation}
 u_0(z)=\int\bigl(1-\cos(z\sqrt{s})\bigr)\,d\nu_0(s)
 =e_0^{\mathsf T}\bigl[I-\cos(z\sqrt{J_{12}})\bigr]e_0.
 \label{hmtfg:articulacion:perfil-jacobi}
\end{equation}
\end{proposition}
\begin{proof}
La positividad anterior permite ortogonalizar hasta grado once.
Si \(i<j-1\), la autoadjunción de la multiplicación da
\(\langle s\psi_j,\psi_i\rangle
=\langle\psi_j,s\psi_i\rangle=0\).
Por simetría, sólo quedan la diagonal y sus dos vecinas; el cociente
positivo de los coeficientes principales fija la positividad de la
subdiagonal. Definamos
\[
 Q_{mj}=\frac{\psi_j(s_m)}{\sqrt{12}},\qquad
 D=\operatorname{diag}(s_1,\ldots,s_{12}).
\]
La ortonormalidad da \(Q^{\mathsf T}Q=I\), y evaluar la
recurrencia en los nodos da \(DQ=QJ_{12}\). Como \(Q\) es
cuadrada, \(J_{12}=Q^{\mathsf T}DQ\). Sus autovalores son
positivos, por lo que la raíz positiva y el coseno se calculan en
esta diagonalización. Finalmente,
\(Qe_0=(1,\ldots,1)^{\mathsf T}/\sqrt{12}\), lo que prueba
\eqref{hmtfg:articulacion:perfil-jacobi}.
\end{proof}

La realización por \((J_{12},e_0)\) recupera el mismo perfil;
por ello, al conservar la familia y el selector, recupera también
sus retornos y sus primeras raíces. La positividad espectral
justifica esta representación. La transversalidad paramétrica
utiliza además la derivada del retorno y sus productos orbitales,
desarrollados en \ref{hmtfg:escalado:14.3} y
\ref{hmtfg:escalado:6}. Son propiedades de operaciones diferentes:
la demostración del escalado conserva ambas y no sustituye el
control de la tangente por la positividad de los momentos.

\subsection{Los canales constitutivos y su inversión etiquetada}
\label{hmtfg:articulacion:constitutiva}

El lector constitutivo actúa sobre otro objeto tipado del mismo
soporte: el transporte angular de dos hojas. Sus coordenadas ya
generadas son \(x=A_{\rm rad}\) y \(y=C^*_{\rm rad}\), en la
cámara contractiva \(x>|y|\). La involución de hojas
\(R=R^*=R^{-1}\) determina los proyectores marcados
\(P_\pm=(I\pm R)/2\). La representación del transporte es
\begin{equation}
 q_+=\exp[-(x+y)],\qquad q_-=\exp[-(x-y)],\qquad
 T=q_+P_++q_-P_-,\qquad 0<q_\pm<1.
 \label{hmtfg:articulacion:canales}
\end{equation}
La evaluación exponencial realiza esos canales después de su
construcción angular. Las incidencias \(90\) y \(120\) del
transporte fijan la respuesta
\[
 \mathcal V=(I-T^{90})(I-T^{120})^{-1}.
\]
El inverso existe porque \(1-q_\pm^{120}>0\).
Escribiendo \(s=q^{30}\), con \(30=\gcd(90,120)\), se obtiene
\begin{equation}
 f(s)=\frac{1-s^3}{1-s^4}
     =\frac{1+s+s^2}{1+s+s^2+s^3},\qquad
 r_\pm=f(q_\pm^{30}),\qquad
 \mathcal V=r_+P_++r_-P_-.
 \label{hmtfg:articulacion:respuesta}
\end{equation}
En esta subsección \(f\) designa la respuesta constitutiva;
\(f_\lambda\) sigue designando el mapa dinámico de
\eqref{hmtfg:articulacion:perfil}. Sus argumentos y dominios
son distintos.

\begin{proposition}[Restitución de los canales y de la orientación]
\label{hmtfg:articulacion:inversion}
La función constitutiva \(f:(0,1)\to(3/4,1)\) es una biyección
decreciente. Su inversa \(\psi\) asigna a \(r\) la única raíz
\(s\in(0,1)\) de
\begin{equation}
 r s^3+(r-1)(1+s+s^2)=0.
 \label{hmtfg:articulacion:cubica}
\end{equation}
Las lecturas constitutivas normalizadas, con sus hojas etiquetadas,
\[
 \widehat\varepsilon=r_+^2,\qquad
 \widehat\mu=r_-^2,\qquad
 \widehat Z=\frac{r_-}{r_+},\qquad
 \widehat c=\frac1{r_+r_-},
\]
con \(\widehat\varepsilon,\widehat\mu\in(9/16,1)\), permiten recuperar
\begin{equation}
 \begin{aligned}
 r_+&=\sqrt{\widehat\varepsilon},&
 r_-&=\sqrt{\widehat\mu},\\
 s_\pm&=\psi(r_\pm),& q_\pm&=s_\pm^{1/30}.
 \end{aligned}
 \label{hmtfg:articulacion:recuperacion-canales}
\end{equation}
\begin{equation}
 x=-\frac12\log(q_+q_-),\qquad
 y=\frac12\log\frac{q_-}{q_+}.
 \label{hmtfg:articulacion:recuperacion-angular}
\end{equation}
\end{proposition}
\begin{proof}
La derivación del cociente racional da
\begin{equation}
 f'(s)=-\frac{s^2(3+2s+s^2)}{(1+s+s^2+s^3)^2}<0
 \qquad(0<s<1).
 \label{hmtfg:articulacion:monotonia}
\end{equation}
Los límites \(f(0^+)=1\) y \(f(1^-)=3/4\), junto con la
continuidad, prueban la biyección. Multiplicar \(f(s)=r\) por
el denominador positivo produce \eqref{hmtfg:articulacion:cubica}.
La positividad de \(r_\pm\) determina las raíces cuadradas
y la de \(s_\pm\) determina las raíces de orden treinta.
Por \eqref{hmtfg:articulacion:canales},
\(\log q_+=-x-y\) y \(\log q_-=-x+y\);
sumar y restar recupera \(x,y\).
Recíprocamente, las raíces recuperadas pertenecen a \((0,1)\),
de modo que \(x\pm y=-\log q_\pm>0\); la pareja obtenida
permanece en la cámara \(x>|y|\).
\end{proof}

En particular,
\(\widehat\mu\widehat\varepsilon=\widehat c^{-2}\) y
\(\widehat\mu/\widehat\varepsilon=\widehat Z^2\).
El intercambio de las etiquetas permuta \(q_+\) y \(q_-\),
conserva \(x\) y cambia el signo de \(y\). La restitución
requiere esa orientación: el producto aislado \(r_+r_-\)
conserva la propagación común, pero no determina las dos respuestas
etiquetadas. Las raíces positivas y la inversión cúbica recuperan
los canales dentro del lector fijado; el registro completo de su
generación permanece en el estado enriquecido.

\subsection{El momento de Catalán y la restitución funcional}
\label{hmtfg:articulacion:catalan}

Las mismas incidencias determinan una relación entre funciones,
antes de evaluar cualquiera de los canales. Para \(0<s<1\), sea
\(\mathcal D=s\,d/ds\), y definamos el momento orientado
\begin{equation}
 \mathcal C_2(s)
 =\int_0^s\log\frac{s}{t}\,
     \frac{1+t}{1+t+t^2}f(t)\,dt
 =\int_0^s\frac{\log(s/t)}{1+t^2}\,dt.
 \label{hmtfg:articulacion:integral}
\end{equation}
La segunda igualdad procede de la factorización
\(1+t+t^2+t^3=(1+t)(1+t^2)\). El peso continuo y positivo
\(dt/(1+t^2)\) es así producido por la composición racional
del lector constitutivo, y no por una identificación con
\(\nu_0\).

\begin{proposition}[Inversión diferencial, unicidad y lectura extrema]
\label{hmtfg:articulacion:restitucion-catalan}
El momento \(\mathcal C_2\) es la única función
\(F\in C^2((0,1))\) que satisface
\begin{equation}
 \mathcal D^2F(s)=\frac{s(1+s)}{1+s+s^2}f(s),\qquad
 \lim_{s\downarrow0}F(s)=0.
 \label{hmtfg:articulacion:problema-diferencial}
\end{equation}
La respuesta constitutiva se restituye por
\begin{equation}
 \mathcal D^2\mathcal C_2(s)=\frac{s}{1+s^2},\qquad
 f(s)=\frac{1+s+s^2}{s(1+s)}\,
          \mathcal D^2\mathcal C_2(s).
 \label{hmtfg:articulacion:inversa-diferencial}
\end{equation}
Su extensión continua al intervalo cerrado tiene la serie
\begin{equation}
 \mathcal C_2(s)=\sum_{j=0}^{\infty}
       \frac{(-1)^j s^{2j+1}}{(2j+1)^2},\qquad 0\le s\le1,
 \qquad
 \mathcal C_2(1)=\sum_{j=0}^{\infty}\frac{(-1)^j}{(2j+1)^2}
 =G_{\rm Cat}.
 \label{hmtfg:articulacion:serie-catalan}
\end{equation}
\end{proposition}
\begin{proof}
El integrando de \eqref{hmtfg:articulacion:integral} es no negativo
y
\[
 0\le\mathcal C_2(s)\le\int_0^s\log(s/t)\,dt=s.
\]
Esto prueba integrabilidad en cero y el límite requerido.
La identidad \(\log(s/t)=\int_t^s du/u\), seguida del
intercambio de integrales no negativas, proporciona
\[
 \mathcal C_2(s)
 =\int_0^s\frac1u\left(\int_0^u\frac{dt}{1+t^2}\right)du.
\]
En todo intervalo compacto del dominio abierto se aplica el
teorema fundamental del cálculo. Por tanto,
\[
 \mathcal D\mathcal C_2(s)=\int_0^s\frac{dt}{1+t^2},
 \qquad
 \mathcal D^2\mathcal C_2(s)=\frac{s}{1+s^2}.
\]
La factorización racional usada en
\eqref{hmtfg:articulacion:integral} da
\eqref{hmtfg:articulacion:problema-diferencial}; despejar \(f\)
da \eqref{hmtfg:articulacion:inversa-diferencial}, cuyos
denominadores son positivos en \((0,1)\).

Si dos soluciones tienen el límite nulo requerido, su diferencia
\(H\) cumple \(\mathcal D^2H=0\). El cambio \(z=\log s\)
transforma esta ecuación en \(d^2H(e^z)/dz^2=0\), de donde
\(H(s)=a\log s+b\). La existencia del límite nulo obliga a
\(a=b=0\), y prueba la unicidad.

Para \(s<1\), la expansión geométrica de \((1+t^2)^{-1}\)
puede integrarse término a término con el peso logarítmico, pues
\[
 \int_0^s t^{2j}\log(s/t)\,dt
 =\frac{s^{2j+1}}{(2j+1)^2},\qquad
 \sum_{j\ge0}\frac{s^{2j+1}}{(2j+1)^2}<\infty.
\]
Esto prueba la serie en el interior. La cota
\(1/(2j+1)^2\) da convergencia absoluta y uniforme de esa
serie en \([0,1]\). En la integral, el integrando extendido
por cero para \(t>s\) está dominado por \(\log(1/t)\)
en \((0,1)\). La convergencia dominada permite alcanzar
\(s=1\) y coincide con la extensión de la serie. Su valor
extremo es precisamente la serie de Catalán indicada.
\end{proof}

La inversión diferencial utiliza la familia \(\mathcal C_2(s)\),
no solamente su valor extremo \(G_{\rm Cat}\). En particular,
\(\mathcal D^2\mathcal C_2(s)\to1/2\) cuando \(s\uparrow1\).
La serie diferenciada dos veces con \(\mathcal D\) tiene allí
un límite radial de Abel; la serie del propio momento es
absolutamente convergente en la frontera. Esta distinción mantiene
las condiciones precisas de cada paso al límite.

\subsection{Composición de las restituciones y alcance conjunto}
\label{hmtfg:articulacion:conclusion}

Las interrelaciones anteriores tienen operaciones y dominios
explícitos. Los momentos de la medida atómica \(\nu_0\)
producen \(J_{12}\), y su cálculo funcional recupera el perfil
\(u_0\). El perfil y el selector de primeras raíces determinan
la sucesión cuyo escalado publica \(\delta_{\mathrm F}\).
Por otra parte, la función constitutiva \(f\) produce
\(\mathcal C_2\) mediante el momento integral, y
\eqref{hmtfg:articulacion:inversa-diferencial} recupera \(f\)
a partir de esa familia. Sus evaluaciones etiquetadas en
\(s_\pm=q_\pm^{30}\) determinan las respuestas constitutivas;
la inversión cúbica recupera ambos canales y sus coordenadas
angulares. La lectura extrema del momento da \(G_{\rm Cat}\).

La medida \(\nu_0\) es positiva, atómica y de rango doce.
El momento de Catalán integra un peso positivo continuo y posee
una expansión alternada. Los signos de esa expansión son
coeficientes de la serie geométrica; no son pesos negativos de
\(\nu_0\), ni convierten la integral positiva en la misma
representación espectral. Cada lector conserva su soporte,
orientación, normalización y operación de evaluación.

La unidad reside en la estructura APP--TRIT--TPK que genera y
transporta sus antecedentes; las relaciones funcionales se
establecen por las composiciones demostradas. La igualdad de
origen no identifica operadores de dominios diferentes.
En particular, esta articulación no postula una fórmula escalar
que determine \(\delta_{\mathrm F}\) a partir de Catalán y de
otras constantes físicas. Sitúa el escalado crítico junto a
restituciones efectivas entre lectores del mismo soporte,
conservando la información requerida por cada dirección.

\endgroup

\begingroup\allowdisplaybreaks[0]% Formalización reunida el 26 de septiembre de 2026.
% Procedencia y condiciones: COMUN/procedencia_20260926/.
\clearpage
\section{Normalización conjunta de acción, longitud y acoplamiento gravitatorio}
\label{md:rigidez}

La determinación de una magnitud se conserva al componer las lecturas del
mismo estado. Aquí se reúnen el retorno inscrito, la sección de acción, la
velocidad constitutiva y el carácter positivo de ruta. El antecedente común
es la generación APP--TRIT--TPK, con sus hojas, orientación, acarreo, memoria
y prolongaciones; las coordenadas regionales y $\alpha$ intervienen como
salidas de esa generación. La construcción siguiente utiliza esas salidas
antes de efectuar el reconocimiento dimensional de $G$.

El problema de normalización comprende tres operaciones distintas: producir
una longitud desde el retorno registrado, transportar el carácter sin perder
su complemento y comparar las lecturas de energía y longitud que resultan.
Las identidades de esta sección fijan conjuntamente esas operaciones. Sus
pruebas se refieren a fibras finitas o a operadores acotados uniformemente
positivos; el sector nulo conserva su tratamiento separado.

\subsection{Retorno inscrito y sección longitudinal}
\label{md:rigidez:longitud}

En el portador marcado del transporte se conserva el soporte
\[
 \Omega=\{(j,k,q,u,v):j\in\mathbb Z/12\mathbb Z,\ 
 k-j\in\{0,3,6,9\},\ 1\le q\le8,\
 u,v\in\{1,2,4,5,7,8\},\ u<v\}.
\]
El índice $j$ conserva la fase dodecafásica; los cuatro desplazamientos
conservan el cuarto de giro y los restantes índices la incidencia marcada.
Por tanto $N=|\Omega|=12\cdot4\cdot8\binom62=5760$ cuenta posiciones de
incidencia. El espacio de este registro se transporta junto con la memoria
genealógica; su dimensión finita corresponde a este portador particular.

Sean $B:V_U\to V_K$ el inversor de Hadamard transportado,
$B^*B=BB^*=I/4$, $u_K=N^{-1/2}(1,\ldots,1)$,
$P_K=u_Ku_K^*$ y $Q_K=I-P_K$. Definimos
\[
 C_0=Q_KB,\qquad M_0=P_KB,\qquad U=2B.
\]
La ortogonalidad de los dos proyectores demuestra
\[
 C_0^*C_0+M_0^*M_0=I/4,\qquad C_0^*M_0=0,\qquad U^*U=UU^*=I.
\]
Así, $x\mapsto(2C_0x,2M_0x)$ conserva el archivo completo de este
transporte. $U$ es su normalización unitaria; la evolución enriquecida
$U_t$ mantiene su propio dominio y su actualización.

Sea $E_i=e_ie_i^*$ la resolución del registro y
$\Delta_\Omega(A)=\sum_iE_iAE_i$ su expectativa diagonal.

\begin{lemma}[Retorno longitudinal inscrito]
\label{md:rigidez:retorno}
El retorno y su lectura inscrita satisfacen
\[
 C_0C_0^*=Q_K/4,\qquad
 \Delta_\Omega(C_0C_0^*)=r_NI,\qquad
 r_N=\frac{N-1}{4N}=\frac{5759}{23040}.
\]
Con la norma escalar local $n_H=\alpha^{16}$ del conúcleo de $H_4$,
la composición longitudinal sobre una cocadena de longitud $\beta_*$ es
\[
 \mathscr T_L(\beta_*)=
 n_H\bigl(\Delta_\Omega(C_0C_0^*)\otimes I\bigr)\beta_*
 =q_\alpha\beta_*,
 \qquad q_\alpha=\alpha^{16}r_N.
 \label{md:rigidez:qalpha}
\]
\end{lemma}
\begin{proof}
La igualdad $BB^*=I/4$ da $C_0C_0^*=Q_K/4$. Cada entrada diagonal
de $P_K$ vale $1/N$, lo que prueba la expectativa. Multiplicar una
cocadena por ese escalar conserva su tipo longitudinal y da la última
igualdad. Si $\beta_*(\bar a)=-U_a\beta_*(a)$, la imagen satisface la
misma inversión. Asimismo, de
$\beta_m(a)=\sum_jU_j^{-1}\beta_{m+1}(a_j)$ se obtiene la identidad
refinada multiplicando ambos miembros por $q_\alpha$.
\end{proof}

La expectativa posee una realización inscrita explícita:
$We_i=e_i\otimes e_i$ y
$W^*(A\otimes I)W=\Delta_\Omega(A)$. Para
$R_K=Q_K/4$ y $P_W=WW^*$, el complemento
$Z_\perp=(I-P_W)(R_K\otimes I)W$ verifica
\[
 (R_K\otimes I)W=Wr_N+Z_\perp,\qquad W^*Z_\perp=0,\qquad
 Z_\perp^*Z_\perp=\frac{N-1}{16N^2}I.
\]
En efecto, $R_K^2=R_K/4$, de modo que el cuadrado de la norma completa
es $r_N/4$, mientras la componente inscrita tiene cuadrado $r_N^2$.
La diferencia es precisamente el complemento escrito. La conservación de
ese complemento acompaña la lectura longitudinal; sumarlo a otra lectura
cambiaría el observable.

La expectativa es única entre las aplicaciones bimodulares hacia el
álgebra diagonal que fijan sus proyectores: sobre una unidad matricial
$E_{ij}$, la bimodularidad exige
$\mathcal E(E_{ij})=E_i\mathcal E(E_{ij})E_j$, que vale cero si $i\ne j$,
y $\mathcal E(E_{ii})=E_{ii}$. Esta unicidad conserva la resolución marcada.

Con la sección longitudinal de referencia $L_*$ ya especificada, queda
\begin{equation}
 L_\alpha=q_\alpha L_*,\qquad D=L_\alpha U.
 \label{md:rigidez:longitud-def}
\end{equation}
El lector usa el retorno como operador lineal sobre una longitud. Una
lectura cuadrática de la fuente produciría otro tipo de magnitud. La
sección $L_*$, la resolución inscrita y la norma local son antecedentes
explícitos de esta composición.

\subsection{Transporte polar y reciprocidad del mismo carácter}
\label{md:rigidez:polar}

Sea $X$ el carácter operatorio del sector positivo no nulo, conservando la ruta y
sus normalizaciones. En dimensión finita se exige $X>0$; en dimensión
infinita se exige $X$ acotado, autoadjunto y $X\ge\epsilon I$ para algún
$\epsilon>0$. Escribimos $Y=UXU^*$. La representación radial polar es
\[
 R=(DX^2D^*)^{1/2}.
\]
La raíz positiva es única y, al sustituir $D=L_\alpha U$, se obtiene
$R=L_\alpha Y$. Equivalentemente,
$\|Rv\|^2=\|XD^*v\|^2$ para todo $v$: la polarización recupera
$R^2=DX^2D^*$ y la positividad determina $R$. Esta condición especifica
la métrica inducida que se conserva.

La acción $\hbar$ y la velocidad constitutiva $c$ pertenecen a la misma
sección. La realización circular compatible fija
\begin{equation}
 \omega_0=\frac{c}{L_\alpha},\qquad
 t_0=\frac{\pi_{\rm HMT}L_\alpha}{54c},\qquad
 R_{108}=\frac{54}{\pi_{\rm HMT}}ct_0=L_\alpha.
 \label{md:rigidez:reloj}
\end{equation}
La longitud se ha construido antes de expresar el reloj circular. Esta
igualdad permite conservar la notación $R_{108}$ en los desarrollos
posteriores y fija su normalización. Las bases dimensionales se transportan
conjuntamente; $L_*$ y la base de arista $u_L$ tienen funciones diferentes.

\begin{theorem}[Determinación conjunta del acoplamiento radial]
\label{md:rigidez:teorema}
Los lectores sobre el mismo carácter,
\[
 H=\frac{\hbar c}{L_\alpha}Y,\qquad
 \Lambda=L_\alpha Y^{-1},\qquad R=L_\alpha Y,\qquad M=H/c^2,
\]
satisfacen
\begin{equation}
 H\Lambda=\hbar cI,\qquad R\Lambda=L_\alpha^2I,\qquad
 R=\frac{L_\alpha^2}{\hbar c}H.
 \label{md:rigidez:identidades}
\end{equation}
En la identificación dimensional de $R$ como radio gravitatorio reducido,
$R=GM/c^2$, el coeficiente común es único:
\begin{equation}
 \boxed{G=\frac{c^3L_\alpha^2}{\hbar}.}
 \label{md:rigidez:G}
\end{equation}
\end{theorem}
\begin{proof}
La multiplicación de los lectores da las dos primeras identidades.
Multiplicarlas a la derecha por $\Lambda^{-1}$ prueba la tercera.
Sustituir $M=H/c^2$ identifica el coeficiente. Si $G'$ conserva las mismas
lecturas, $(G'-G)M/c^2=0$; en la fibra no nula, la invertibilidad de $M$
implica $G'=G$.
Todos los productos están definidos en la fibra completa bajo las
hipótesis del enunciado.
\end{proof}

La prueba aplica a todos los modos positivos de esa realización, sin elegir
la masa electrónica para normalizar $G$. La identificación del radio
reducido es explícita; el radio $2GM/c^2$ corresponde a otra convención
de lectura de la misma constante. El sector nulo se mantiene fuera del
inverso. La extensión a operadores no acotados requiere un desarrollo de
dominios y límites distinto del teorema presente.

El cuadrado de $L_\alpha$ procede del producto de longitudes conjugadas.
La acción aparece inversamente al eliminar $\Lambda$ entre energía y
longitud. El paso de energía a masa y de masa a radio aporta $c^4$,
que combinado con el $c$ de $H\Lambda=\hbar cI$ produce $c^3$.
Las dos normalizaciones fijan el factor adimensional a uno. El análisis
dimensional verifica los exponentes; la composición determina el coeficiente.

\subsection{Escalas de Planck y especialización del carácter unitario}
\label{md:rigidez:escalas}

La longitud construida, la sección de acción y la velocidad determinan
\[
 t_P=\frac{L_\alpha}{c},\qquad
 m_{P,a}=\frac{\hbar_a}{cL_\alpha},\qquad
 E_{P,a}=\frac{\hbar_a c}{L_\alpha}.
\]
Sustituir \eqref{md:rigidez:G} en
$\sqrt{\hbar_aG_a/c^5}$, $\sqrt{\hbar_ac/G_a}$ y
$\sqrt{\hbar_ac^5/G_a}$ demuestra respectivamente esas identidades;
la positividad fija cada raíz. La longitud antecede al reconocimiento
convencional de estas escalas.
Para un modo de carácter $y>0$,
\[
 m(y)=m_{P,a}y,\qquad r_g(y)=L_\alpha y,\qquad
 \bar\lambda(y)=L_\alpha/y.
\]
La igualdad de las dos longitudes equivale a $y=y^{-1}$ y, por tanto,
a $y=1$. El producto areal se conserva para todos los modos positivos;
la coincidencia de radios caracteriza el modo unitario.

El paso cronológico satisface $t_0=(\pi_{\rm HMT}/54)t_P$;
el período de fase es $108t_0=2\pi_{\rm HMT}t_P$ y
$\omega_{108}=1/t_P$. Un horizonte $R_h=\zeta_hL_\alpha$ conserva
su condición geométrica propia. Para la especialización Schwarzschild
de la misma masa, $R_h=2r_g=2L_\alpha y$.
Al comparar secciones con $L_\alpha,c$ fijos y
$\rho_a=\hbar_b/\hbar_a$, se obtiene
$G_b=G_a/\rho_a$, $m_{P,b}=\rho_a m_{P,a}$ y
$E_{P,b}=\rho_a E_{P,a}$. Esta comparación de secciones conserva
$L_\alpha$ y $t_P$.

\subsection{Área y restitución de los sectores de memoria}
\label{md:rigidez:area}

La incidencia de borde conserva la agrupación de seis trits.
Si $x_j=r_{2j-1}+3r_{2j}$, $r_i\in\{0,1,2\}$, el mapa
\[
 (x_1,x_2,x_3)\longmapsto
 (r_1+3r_2+9r_3,\ r_4+3r_5+9r_6)
\]
es una biyección entre $(\mathbb Z/9\mathbb Z)^3$ y
$(\mathbb Z/27\mathbb Z)^2$. Su inversa extrae los seis dígitos en orden,
conservando el lugar donde se registra el acarreo. En el toro de borde,
el bloque $A=\{0,\ldots,8\}^2$ tiene dos colecciones de enlaces:
\begin{align*}
 \partial A_{90}&=\{((-1,j),(0,j)),((8,j),(9,j)):0\le j\le8\},\\
 \partial A_{120}&=\{((i,-1),(i,0)),((i,8),(i,9)):0\le i\le8\}.
\end{align*}
Cada colección contiene dieciocho enlaces. Sus etiquetas conservan la
procedencia de canal; cada enlace lleva un operador número con espectro
$\{0,1,2\}$. Sean $N_{90},N_{120}$ las sumas correspondientes.

La escala angular areal conserva el contraángulo ya construido:
\[
 \theta_{\rm clk}=\frac{\pi C^*_{\deg}}{1080},\qquad
 a_0=\frac{1+2\cos(\pi/18)}3\,\theta_{\rm clk}.
\]
El prefactor es la traza real normalizada de
$\operatorname{diag}(1,e^{i\pi/18},e^{-i\pi/18})$.
Así $a_0$ conserva el significado de normalización adimensional de área.
La lectura de área se toma en la rama positiva
$\theta_{\rm clk}>0$ y $\gamma>0$.
Con el funcional $\gamma$ de Barbero--Immirzi y la longitud ya determinada,
\[
 \mathcal A=\gamma a_0L_\alpha^2(N_{90}+N_{120}).
\]
Los treinta y seis operadores locales actúan en factores tensoriales
distintos. Por tanto, el espectro de $N=N_{90}+N_{120}$ es
$\{0,\ldots,72\}$ y la multiplicidad de $n$ es el coeficiente de $z^n$
en $(1+z+z^2)^{36}$. La prueba consiste en multiplicar las funciones
generatrices de los factores. La longitud nueva transporta la unidad
areal y conserva esta combinatoria.

\begin{proposition}[Lectura conjunta de área y procedencia sectorial]
La pareja
$N=N_{90}+N_{120}$ y $\mathcal D_{\partial}=90N_{90}+120N_{120}$
restituye los dos sectores:
\[
 N_{90}=4N-\frac{\mathcal D_{\partial}}{30},\qquad
 N_{120}=\frac{\mathcal D_{\partial}}{30}-3N.
\]
\end{proposition}
\begin{proof}
La matriz de la transformación es
$\left(\begin{smallmatrix}1&1\\90&120\end{smallmatrix}\right)$,
con determinante $30$. Su inversa da las fórmulas. La conmutatividad
de ambos operadores número implica la de los dos lectores.
\end{proof}
En la realización modular con
$\rho_A=Z_A^{-1}\exp(-\Delta_{\rm mod}\mathcal D_{\partial})$,
$\Delta_{\rm mod}>0$, el cálculo funcional da
$-\log\rho_A=(\log Z_A)I+\Delta_{\rm mod}\mathcal D_{\partial}$.
Con las normalizaciones conservadas, el par formado por área y
Hamiltoniano modular recupera las dos ocupaciones. La prueba de restitución
vale para cada valor construido de $\Delta_{\rm mod}$; el nuevo lector
longitudinal no redefine ese coeficiente ni el estado modular.
La suma de ocupaciones y la procedencia sectorial expresan observables
distintos. Del mismo modo, los treinta y seis enlaces de respuesta y las
ochenta y una capacidades unitarias del corte cúbico conservan sus funciones.

\subsection{Variación conjunta y reducción compatible de memoria}
\label{md:rigidez:variaciones}

Para colecciones diferenciables en norma y uniformemente positivas en un
entorno, sea $\kappa=L_\alpha^2/(\hbar c)$. Entonces
\[
 \delta R=\delta\kappa\,H+\kappa\,\delta H,\qquad
 \frac{\delta G}{G}=3\frac{\delta c}{c}
 +2\frac{\delta L_\alpha}{L_\alpha}-\frac{\delta\hbar}{\hbar}.
\]
La primera identidad resulta de diferenciar
$R=L_\alpha Y$ y $H=(\hbar c/L_\alpha)Y$; no requiere
$[Y,\delta Y]=0$. Una variación admisible de conexión que actúe en $Y$
y conserve las tres secciones mantiene $\delta G=0$.
La afirmación comprende las variaciones de esta realización radial;
la equivalencia entre funcionales de campo completos tiene sus propios
dominios y condiciones de frontera.

Sea ahora
$Y=\left(\begin{smallmatrix}A&C\\C^*&D_i\end{smallmatrix}\right)>0$
en la separación entre lectura y memoria, con
$S=A-CD_i^{-1}C^*>0$. La extensión minimizante de un vector $b$ es
$(b,-D_i^{-1}C^*b)$, como demuestra completar el cuadrado:
\[
 \left\langle\binom b m,Y\binom b m\right\rangle
 =\langle b,Sb\rangle+
 \langle m+D_i^{-1}C^*b,D_i(m+D_i^{-1}C^*b)\rangle.
\]
Por tanto la reducción coherente da
\[
 H_{\rm eff}=\frac{\hbar c}{L_\alpha}S,\qquad
 R_{\rm eff}=L_\alpha S,\qquad\Lambda_b=L_\alpha S^{-1}.
\]
La compresión de $Y^{-1}$ sobre el bloque de lectura es $S^{-1}$;
la compresión directa de $Y$ es $A$. La memoria distingue esas operaciones.

Si la eliminación dinámica produce una métrica temporal acotada
$Z\ge z_0I$, $z_0>0$, se toma $T=Z^{1/2}$, acotado e invertible.
La normalización canónica exige los transportes duales
\[
 H_c=T^{-*}H_{\rm eff}T^{-1},\qquad
 R_c=T^{-*}R_{\rm eff}T^{-1},\qquad
 \Lambda_c=T\Lambda_bT^*.
\]
Su multiplicación da $H_c\Lambda_c=\hbar cI$ y
$R_c\Lambda_c=L_\alpha^2I$, sin hipótesis de conmutatividad.
En el desarrollo resolvente de una energía con bloques $H_{bb},V,H_{ii}$,
con $H_{ii}$ positivo invertible y
$|\hbar\omega|<\|H_{ii}^{-1}\|^{-1}$,
\[
 K_b(\omega)=H_{bb}-\hbar\omega I
 -V(H_{ii}-\hbar\omega I)^{-1}V^*
 =H_{\rm stat}-\hbar\omega Z+O(\omega^2),
 \qquad Z=I+VH_{ii}^{-2}V^*.
\]
La última expresión es local alrededor de cero y conserva el orden de
aproximación; los productos duales precedentes son identidades exactas
de los lectores normalizados.

La proporcionalidad también se conserva en la eliminación dinámica
completa. Para $z$ tal que $H_{ii}-zI$ sea invertible, definimos
\[
 \mathcal K_H(z)=H_{bb}-zI-V(H_{ii}-zI)^{-1}V^*,\qquad
 \kappa=\frac{L_\alpha^2}{\hbar c}.
\]
La misma descomposición de lectura y memoria aplicada a $R=\kappa H$
produce
\begin{equation}
 \mathcal K_R(\kappa z)=\kappa\mathcal K_H(z).
 \label{md:rigidez:resolvente-exacto}
\end{equation}
En efecto, sus bloques son $R_{bb}=\kappa H_{bb}$,
$R_{bi}=\kappa V$ y $R_{ii}=\kappa H_{ii}$, y
\[
 (R_{ii}-\kappa zI)^{-1}
 =\kappa^{-1}(H_{ii}-zI)^{-1}.
\]
La sustitución en el complemento de Schur demuestra
\eqref{md:rigidez:resolvente-exacto}, conservando el orden de los factores.
Cuando ambos complementos son invertibles, sus inversas satisfacen
$\mathcal K_R(\kappa z)^{-1}=\kappa^{-1}\mathcal K_H(z)^{-1}$.
El parámetro energético es $z=\hbar\omega$ y su imagen en el lector
radial es $\kappa z$: ambos conservan sus dimensiones respectivas.
La identidad vale en todo ese dominio resolvente, por lo que transporta
conjuntamente todos los órdenes de la memoria dinámica. El desarrollo
local anterior describe sus primeros coeficientes; esta igualdad explica
la conservación exacta de la proporción radial al efectuar la reducción.

\subsection{Rigidez temporal y memoria de frontera}
\label{md:rigidez:memoria}

Conservar únicamente el área admite $R\mapsto\zeta R$,
$\Lambda\mapsto\Lambda/\zeta$. Conservar además $H,\hbar,c$ y
$H\Lambda=\hbar cI$ obliga $\zeta=1$. Si también se transforma
$H\mapsto\zeta H$, el cociente $R/H$ y el coeficiente $G$ permanecen
iguales. Las comparaciones deben conservar las mismas observaciones conjuntas.

La fase levantada aporta otra restricción efectiva. En una realización
acotada $H=H_{\rm clk}+D_0^*WD_0$ con $W>0$, la evolución
$P(t)=\exp(-itH/\hbar)$ determina $H=i\hbar\dot P(0)$.
Fijados $H_{\rm clk},D_0,\hbar$ y el reloj, la igualdad de las evoluciones
implica $D_0^*(W'-W)D_0=0$. Para una ruta con extremo inicial fijado,
$D_0$ es triangular invertible y se deduce $W'=W$. Para un grafo general
se determina la forma sobre $\operatorname{Ran}D_0$; extender después el
dominio de residuos modificaría el problema de eliminación.
Un solo retorno $P(\tau)$ admite alias espectrales, mientras la fase
levantada o los tiempos compatibles que tienden a cero los distinguen.

En la coordenatización de capacidad nonádica, sean
$\rho=\log_{10}9$, $\delta=1-\rho$,
$p_j=\lfloor j/\delta\rfloor$ y $v_j=p_j-j$. La capacidad de una ventana es
\[
 C(j,n)=p_{j+n}-p_j-n.
\]
Para $n=35$ sus valores son $729$ o $730$, según el origen; el retorno de
capacidad módulo nueve distingue el primero dentro de esa clase.
La ventana $j=20\to55$ satisface $p_j:437\to1201$ y
$v_j:417\to1146$. Su fase de capacidad vuelve a $3$ y el cociente aumenta
de $46$ a $127$. Las seis capacidades de $20$ y las veintinueve de $21$
suman $729$; la memoria acumulada sigue siendo parte del estado.

Para la primera llegada $T(k)=\lceil k/\rho\rceil-1$, el recuento es
$N(k,C)=T(k+C)-T(k)-C$. Escribir
$\epsilon(k)=\lceil k/\rho\rceil-k/\rho$ da exactamente
\[
 N(k,C)=(\rho^{-1}-1)C+\epsilon(k+C)-\epsilon(k).
\]
Las ventanas de capacidad $729$ con recuentos $34$ y $35$ conservan
términos de frontera diferentes. Incluso $k=21$ y $k=2289$ comparten
$(k\bmod108,T(k)\bmod108)=(21,22)$ y acumulan respectivamente $1$ y
$109$ vacancias. La coincidencia de calendarios finitos conserva menos
información que el estado completo. Estos cálculos interpretan el retorno;
el lector longitudinal de \eqref{md:rigidez:longitud-def} tiene su propio
constructor y no selecciona su escala mediante esas ventanas.

\subsection{Normalización geométrica y evaluación}
\label{md:rigidez:evaluacion}

El bloque central de APP
$B_c=\left(\begin{smallmatrix}7&2\\2&7\end{smallmatrix}\right)$
satisface $B_c^{\mathsf T}g_cB_c=45g_c$,
$g_c=\operatorname{diag}(1,-1)$. Su normalización positiva unimodular es
$B_c/\sqrt{45}$: un factor positivo adicional $a$ daría determinante $a^2$.
La incidencia areal--volumétrica
$T_{\rm av}=\left(\begin{smallmatrix}1&1\\1&2\end{smallmatrix}\right)$
conserva
$G_{\rm av}=\left(\begin{smallmatrix}2&1\\1&-2\end{smallmatrix}\right)$
por multiplicación directa.
En una realización de Minkowski la forma causal y la soldadura dimensional
tienen papeles distintos. La cocadena
$\beta_m=\ell_m\sigma_mW_\square e_\lambda$,
$\ell_m=\ell_0/9^m$, y la relación $\ell_0=ct_0$ conservan la escala
junto con el cono. Alterar la soldadura conservando sólo la firma cambia
esa realización.

La unidad expansiva de Pell con norma uno y traza cuatro satisface
$x+x^{-1}=4$, $x>1$, y por ello $x=2+\sqrt3$.
La órbita $r_n=r_0x^n$ conserva $r_nr_{-n}=r_0^2$.
Cuando $r_0=L_\alpha$, el producto conserva la misma sección areal.
La norma, la forma marcada, el retorno y la sección longitudinal restringen
objetos diferentes; sus funciones se mantienen al componerlos.

\Needspace{10\baselineskip}
En la sección SI declarada, $L_*=1\,\mathrm m$,
$S_*=10^{-34}\,\mathrm{J\,s}$ y $c=299792458\,\mathrm{m\,s^{-1}}$,
la sustitución estructural produce
\[
 L_\alpha=1.616254982625126330\ldots\,10^{-35}\,\mathrm m,\qquad
 \hbar_{\rm ret}=1.054571816915039061\ldots\,10^{-34}\,\mathrm{J\,s},
\]
\begin{equation}
 G_{\rm ret}=6.674299659414022706\ldots\,10^{-11}
 \,\mathrm{m^3\,kg^{-1}\,s^{-2}}.
 \label{md:rigidez:valor}
\end{equation}
El evaluador sustituye las coordenadas regionales y el registro $K$,
reconstruye $\alpha$, el retorno de $H_5$ y $L_\alpha$; la comparación
metrológica se realiza después. La referencia CODATA 2022 es
$6.67430(15)\,10^{-11}$ en esas unidades \cite{codata2022}.
La diferencia es $-3.40586\,10^{-18}$, equivalente a
$-0.00227$ incertidumbres estándar. Las cifras calculadas describen la
evaluación matemática en esta sección, no una incertidumbre experimental.
Las secciones de acción previa y retornada conservan sus etiquetas y
$\hbar_aG_a=c^3L_\alpha^2$.

El significado estructural de $G$ queda así determinado: es el coeficiente
que convierte, en la realización radial común, la lectura energética en
la longitud conjugada preservando simultáneamente acción y área.
La memoria, las marcas de hoja y las bases forman parte de la construcción
que fija ese coeficiente. Su conservación conjunta explica la unicidad
demostrada.
\endgroup
% Preparación autónoma de VII, 2026-09-30. Fragmento sin inputs adicionales.
% RESULTADOS_RECUPERADOS, pruebas reunidas sin modificación de los originales:
% /Users/ruben/Documents/excelencia academica/output/CUMPLIMIENTO_EDITORIAL_20260928/fuentes/VIII_ES/manuscrito/29_retorno_y_memoria_traslacional.tex
% /Users/ruben/Documents/excelencia academica/output/CUMPLIMIENTO_EDITORIAL_20260928/fuentes/VIII_ES/manuscrito/30_pantalla_y_respuesta.tex
% /Users/ruben/Documents/excelencia academica/output/CUMPLIMIENTO_EDITORIAL_20260928/fuentes/VIII_ES/manuscrito/30b_variacion_energia_memoria.tex
% /Users/ruben/Documents/excelencia academica/output/CUMPLIMIENTO_EDITORIAL_20260928/fuentes/VIII_ES/manuscrito/30c_composicion_corriente_conexion.tex
% /Users/ruben/Documents/excelencia academica/output/CUMPLIMIENTO_EDITORIAL_20260928/fuentes/VIII_ES/manuscrito/30d_realizacion_geometrica.tex
% /Users/ruben/Documents/excelencia academica/output/CUMPLIMIENTO_EDITORIAL_20260928/fuentes/VIII_ES/manuscrito/31_corriente_y_respuesta_cartan.tex
% /Users/ruben/Documents/excelencia academica/output/CUMPLIMIENTO_EDITORIAL_20260928/fuentes/VIII_ES/manuscrito/33_corriente_espinorial_y_densidad.tex
% Dependencias internas ya presentes: núcleo común, rutas, lectores y unidades.
% Las inversas de 31 están demostradas en 10, hmtcuatro:gea:torsion.
% Legendre, fuente completa y restricciones permanecen en 20.
% Se incluye de 33 sólo lo utilizado por 10+20, no sus estados cosmológicos.
\section{Antecedentes geométricos y materiales de la acción conjunta}
\label{hmtcuatro:ant:seccion}

Las semillas y las operaciones de APP, la orientación y el régimen
del TRIT y el transporte del TPK producen las rutas enriquecidas sobre
las que actúan los lectores del núcleo común. Se conservan conjuntamente
hoja, orientación, acarreo, frontera y memoria; el retorno de fase no
reinicia ese estado. Las construcciones del continuo pertenecen al mismo
objeto y no se convierten aquí en sectores independientes. La
realización que sigue transporta a la geometría y a la materia esa
estructura ya construida. Sus unidades y acoplamientos son salidas
HMT anteriores: ninguna comparación numérica interviene como generador.

Para que la acción y su reducción puedan leerse dentro de este
manuscrito, reunimos las operaciones de pantalla, la variación efectiva
del transporte y la corriente espinorial que aquéllas utilizan. La
eliminación de Cartan--Holst se demuestra en
\ref{hmtcuatro:gea:torsion}; la acción total y su reducción canónica,
en \ref{hmtcuatro:accion-retroaccion-restricciones}. Las fórmulas de
esta sección especifican los datos y los cálculos que entran en ambas.

\subsection{Incidencia, pantalla y soldadura}
\label{hmtcuatro:ant:pantalla}

La representación cúbica de la frontera de una ruta tiene las seis
caras \(F=\{(i,\epsilon):1\leq i\leq3,\ \epsilon=\pm1\}\).
Sobre \(H_F=\mathbb R^F\), la oposición es
\(Re_{i,\epsilon}=e_{i,-\epsilon}\). Con
\[
 u=\sum_{i,\epsilon}e_{i,\epsilon},\qquad
 P_u=\frac{uu^{\mathsf T}}6,\qquad
 P_-=\frac{I-R}2,\qquad
 P_0=\frac{I+R}2-P_u ,
\]
los tres proyectores son ortogonales, suman \(I\) y tienen rangos
\(1,3,2\). La incidencia entre caras y vértices
\(\nu\in\{-1,1\}^3\) es
\(M_{(i,\epsilon),\nu}=\mathbf1_{\{\nu_i=\epsilon\}}\).

\begin{proposition}[Pantalla de la incidencia orientada]
\label{hmtcuatro:ant:prop-pantalla}
Se cumple
\[
 MM^{\mathsf T}=12P_u+4P_-,\qquad
 \ker M^{\mathsf T}=P_0H_F .
\]
El cociente \(Q=H_F/P_0H_F\), identificado con
\((P_u+P_-)H_F\), es de dimensión cuatro. La forma
\(B=P_--P_u\) sobre \(Q\) tiene firma \((-+++)\).
\end{proposition}
\begin{proof}
Cada cara contiene cuatro vértices; dos caras opuestas no comparten
ninguno y dos caras de ejes distintos comparten dos. Por tanto,
\[
 MM^{\mathsf T}=4I+2(uu^{\mathsf T}-I-R)
              =12P_u+4P_- .
\]
La identidad
\(\|M^{\mathsf T}x\|^2=\langle x,MM^{\mathsf T}x\rangle\)
determina el núcleo. En la base ortonormal
\[
 t=u/\sqrt6,\qquad d_i=(e_{i,+}-e_{i,-})/\sqrt2
\]
de \(Q\), la matriz de \(B\) es \(\operatorname{diag}(-1,1,1,1)\).
\end{proof}

Sea \(W_\square=\sqrt6P_u+\sqrt2P_-\). Entonces
\[
 W_\square e_{i,\epsilon}=t+\epsilon d_i=:n_{i,\epsilon},
 \qquad B(n_{i,\epsilon},n_{i,\epsilon})=0,\qquad
 MM^{\mathsf T}|_Q=2W_\square^*W_\square|_Q .
\]
Estas igualdades se obtienen aplicando los proyectores a la base de
caras. Las rotaciones propias del cubo preservan los proyectores,
la orientación temporal \(t\) y la forma \(B\).

La ruta enriquecida proporciona a cada arista \(a\) su cara,
su signo \(\sigma_m(a)\) y el transporte
\(U_m(a):Q_{s(a)}\to Q_{t(a)}\). Su soldadura celular es
\begin{equation}
 \beta_m(a)=\ell_m\sigma_m(a)n_{\lambda(\underline a)}^{s(a)},
 \quad \ell_m=\ell_0\,9^{-m},\quad
 \beta_m(\bar a)=-U_m(a)\beta_m(a).
 \label{hmtcuatro:ant:soldadura}
\end{equation}
La inversión \(U_m(\bar a)=U_m(a)^{-1}\) implica que invertir dos
veces recupera la soldadura. Se conserva además el registro ordenado
del acarreo: no se sustituye por el vector de llegada.

\begin{proposition}[Naturalidad de la soldadura]
\label{hmtcuatro:ant:refinamiento-soldadura}
Sea \(a=a_1\cdots a_9\) un refinamiento compatible: las incidencias
hijas transportadas al origen del predecesor conservan cara y signo,
y \(\ell_{m+1}=\ell_m/9\). Si
\(\mathcal U_j:Q_{m,s(a)}\to Q_{m+1,s(a_j)}\) incluye la
identificación por refinamiento y el transporte hasta \(s(a_j)\),
entonces
\[
 \beta_m(a)=\sum_{j=1}^9\mathcal U_j^{-1}\beta_{m+1}(a_j).
\]
\end{proposition}
\begin{proof}
Cada sumando en la fibra de origen es
\((\ell_m/9)\sigma_m(a)n_{\lambda(\underline a)}\). Su suma da
\eqref{hmtcuatro:ant:soldadura}. Invertir la ruta invierte el orden
de los transportes y cada transporte; la fórmula de inversión de
la soldadura da el signo global. La concatenación conserva en el
mismo orden la memoria de las nueve incidencias. La igualdad se
itera por composición de refinamientos compatibles.
\end{proof}

\subsection{Representación del transporte y realización geométrica}
\label{hmtcuatro:ant:geometria}

El registro de una ruta \(\gamma\) cuenta las inversiones y los
cambios efectivos de hoja mediante el cociclo
\(c(\gamma)=(c_g(\gamma),c_s(\gamma))\in\mathbb Z^2\).
El recuento es acumulativo en las rutas hacia adelante; su
extensión al grupoide es orientada y asigna incrementos opuestos
a la ruta inversa.
Para rutas composables,
\(c(\gamma_2\gamma_1)=c(\gamma_2)+c(\gamma_1)\), y
\(c(\bar\gamma)=-c(\gamma)\). Sea \(R_4\) un elemento de orden
cuatro en \(\operatorname{Spin}^+(3,1)\) cuya representación
vectorial \(\Lambda(R_4)\) fija el vector temporal unitario \(u_0\).
La representación afín recibida es
\[
 \rho_{\rm C}^{\rm disc}(\gamma)
 =\bigl(R_4^{c_g(\gamma)},\ell_0c_s(\gamma)u_0\bigr).
\]
Con el producto
\((R,a)(S,b)=(RS,a+\Lambda(R)b)\), la fijación de \(u_0\) y la
aditividad de \(c\) prueban directamente que \(\rho_{\rm C}^{\rm disc}\)
conserva identidad, concatenación e inversión. Cuando la componente
rotacional retorna a la identidad, la traslación de memoria permanece
determinada por \(c_s(\gamma)\), que no se reinicia.

La normalización
\[
 N_{\ell_0}(R,a)=(\Lambda(R),a/\ell_0),\qquad
 \rho_{\rm C}=N_{\ell_0}\circ\rho_{\rm C}^{\rm disc}
\]
también es un homomorfismo, pues
\((a+\Lambda(R)b)/\ell_0=a/\ell_0+\Lambda(R)(b/\ell_0)\).
El levantamiento \(R_4^{c_g(\gamma)}\) se conserva cuando actúan
espinores; no se recupera sólo a partir de \(\Lambda(R_4^{c_g(\gamma)})\),
cuya representación tiene núcleo \(\{1,-1\}\).

En una realización de pantalla \(J_{{\rm scr},m}\) que transporta
la forma y la soldadura, \(e_m=J_{{\rm scr},m}\beta_m\).
La realización suave considerada consiste en una región orientada
\(\mathcal U\) de dimensión cuatro, una cotetrada no degenerada
\(e\), una conexión de Lorentz \(\omega\) y una realización de rutas
\(\mathfrak R_{\rm C}\) con la condición de compatibilidad
\begin{equation}
 \operatorname{Hol}_{\mathcal A_{\ell_0}}
          (\mathfrak R_{\rm C}(\gamma))
 =\rho_{\rm C}(\operatorname{Hol}_{\rm TPK}(\gamma)),\qquad
 \mathcal A_{\ell_0}=
 \begin{pmatrix}\omega&\ell_0^{-1}e\\0&0\end{pmatrix}.
\label{hmtcuatro:ant:compatibilidad-holonomia}
\end{equation}
El argumento de \(\rho_{\rm C}\) conserva el cociclo de la
holonomía enriquecida; no se sustituye por su sola fase reducida.
La escala \(\ell_0>0\) es constante en la coordenatización. Éstos
son los datos y la condición de la realización suave utilizada;
preservar la composición de caminos no identifica entre sí caminos
distintos de iguales extremos ni elimina su memoria. La métrica
es \(g=\eta_{IJ}e^I\otimes e^J\), \(\eta=\operatorname{diag}(-1,1,1,1)\).

\begin{proposition}[Curvatura, covariancia y Bianchi]
\label{hmtcuatro:ant:curvatura}
Con \(F_\omega=\mathrm d\omega+\omega\wedge\omega\) y
\(T=\mathrm de+\omega\wedge e\),
\[
 \mathcal F_{\ell_0}=
 \begin{pmatrix}F_\omega&\ell_0^{-1}T\\0&0\end{pmatrix},
 \qquad D_\omega T=F_\omega\wedge e,\quad D_\omega F_\omega=0 .
\]
Un cambio de marco \(g_L:\mathcal U\to SO^+(3,1)\) actúa por
\(e'=g_Le\),
\(\omega'=g_L\omega g_L^{-1}-\mathrm dg_L\,g_L^{-1}\);
entonces \(T'=g_LT\), \(F_{\omega'}=g_LF_\omega g_L^{-1}\).
\end{proposition}
\begin{proof}
La multiplicación de matrices de formas da el bloque diagonal
\(\mathrm d\omega+\omega\wedge\omega\) y el traslacional
\(\ell_0^{-1}(\mathrm de+\omega\wedge e)\). Al sustituir los
campos transformados, los términos con \(\mathrm dg_L\) se
cancelan. Expandir \(D_\omega T\) deja
\((\mathrm d\omega+\omega\wedge\omega)\wedge e\).
En \(D_\omega F_\omega\), las derivadas de \(\omega\) y los
productos cúbicos se cancelan por pares.
\end{proof}

La hoja trítica \(\tau\in\{+1,0,-1\}\) conserva su régimen en la
realización de Cartan con generadores \(J_{ab},P_a^{(\tau)}\),
acción vectorial de Lorentz y
\([P_a^{(\tau)},P_b^{(\tau)}]=-\tau J_{ab}\). Para
\(\mathcal A_\tau=\frac12\omega^{ab}J_{ab}+\ell^{-1}e^aP_a^{(\tau)}\),
\[
 \mathcal F_\tau=
 \tfrac12(F_\omega^{ab}-\tau\ell^{-2}e^a\wedge e^b)J_{ab}
 +\ell^{-1}T^aP_a^{(\tau)} .
\]
En efecto, el corchete de Lorentz da \(F_\omega\), el mixto
da \(D_\omega e\) y el traslacional da el término
\(-\tau e\wedge e/(2\ell^2)\). La condición
\eqref{hmtcuatro:ant:compatibilidad-holonomia} corresponde a la
hoja central; las hojas extremas utilizan su representación
\(\rho_{{\rm C},\tau}\) y su condición de holonomía propias.

\subsection{Diferencial del transporte y corriente material}
\label{hmtcuatro:ant:corriente-transporte}

En una realización celular finita, sea
\(U_a:\mathcal H_{s(a)}\to\mathcal H_{t(a)}\) un transporte
invertible entre fibras hermíticas. No se presupone que todo
transporte de Lorentz sea unitario para el producto energético.
Se toma una orientación por arista. Para
\[
 (Df)_a=f_{t(a)}-U_af_{s(a)},\quad r=Df,\quad
 q=\mathsf W r,\quad v_a=U_af_{s(a)},\quad
 E=\langle r,\mathsf W r\rangle ,
\]
el peso \(\mathsf W=\mathsf W^*>0\) puede acoplar aristas.
Este peso no es \(W_\square\). Las variaciones corresponden a
la representación de la conexión y al problema de borde elegido.

\begin{proposition}[Variación completa]
\label{hmtcuatro:ant:variacion-completa}
Para una curva diferenciable de datos, con
\(A_a=\dot U_aU_a^{-1}\), se cumple
\begin{equation}
 \dot E=2\operatorname{Re}\sum_a
 \langle q_a,\dot f_{t(a)}-U_a\dot f_{s(a)}-A_av_a\rangle
 +\langle r,\dot{\mathsf W}r\rangle .
 \label{hmtcuatro:ant:variacion}
\end{equation}
A campo y peso fijos, el gradiente completo del transporte es
\(G_a=-2q_av_a^*\):
\[
 \delta_U E=\sum_a\operatorname{Re}\operatorname{Tr}(G_a^*A_a).
\]
En la subclase unitaria \(A_a^*=-A_a\), basta su parte
antihermítica \(J_a=v_aq_a^*-q_av_a^*\).
\end{proposition}
\begin{proof}
La regla del producto da
\(\dot E=2\operatorname{Re}\langle\mathsf W r,\dot r\rangle+
\langle r,\dot{\mathsf W}r\rangle\).
Derivar cada residuo produce \eqref{hmtcuatro:ant:variacion}.
Como \(\operatorname{Tr}(v_aq_a^*A_a)=q_a^*A_av_a\), se obtiene
\(G_a\). La parte hermítica de \(G_a\) es ortogonal, para
\(\operatorname{Re}\operatorname{Tr}(X^*Y)\), a las variaciones
antihermíticas. Esta reducción sólo se aplica en esa subclase.
\end{proof}

Para una ruta con \(U_\gamma=U_N\cdots U_1\), sea
\(B_k=U_N\cdots U_{k+1}\). La regla del producto demuestra
\begin{equation}
 \dot U_\gamma U_\gamma^{-1}
 =\sum_k B_kA_kB_k^{-1}.
 \label{hmtcuatro:ant:variacion-ruta}
\end{equation}
Consecuentemente un gradiente \(G_\gamma\) se transporta hacia
la arista \(k\) como
\(G_k=B_k^*G_\gamma(B_k^{-1})^*\). La ciclicidad de la traza
prueba esta fórmula incluso si \(B_k\) no es unitario.

Si las coordenadas reales de variación de conexión son
\(\theta_b\), se escribe su diferencial efectivo como
\(A_a(\theta)=\sum_b B_{ab}\theta_b\). Para
\(S_{\rm mem}=\mathfrak a E\), a campo, peso y sección de acción
\(\mathfrak a\) fijos,
\begin{equation}
 \delta S_{\rm mem}=-\tfrac12\sum_b\theta_b s_b,\qquad
 s_b=4\mathfrak a\,\operatorname{Re}
          \sum_a\langle q_a,B_{ab}v_a\rangle .
 \label{hmtcuatro:ant:corriente-celular}
\end{equation}
Sustituir \(A_a(\theta)\) en la variación demuestra la igualdad
y fija el factor \(4\). Si el emparejamiento celular orientado es
una matriz invertible \(\mathsf M\), definida por
\(\delta S_{\rm mem}=-\frac12\theta^{\mathsf T}\mathsf M\sigma\),
la corriente es \(\sigma=\mathsf M^{-1}s\). Esta matriz conserva
los volúmenes y la orientación; no puede sustituirse por una
división escalar sin fijar una base en la que corresponda.
Si también varían \(f,\mathsf W,\mathfrak a\), se conservan todos
los términos de \eqref{hmtcuatro:ant:variacion} y \(E\delta\mathfrak a\).

Para un peso por arista, su inversión conserva la respuesta completa
al transformar conjuntamente
\[
 U_{\bar a}=U_a^{-1},\qquad
 r_{\bar a}=-U_a^{-1}r_a,\qquad
 \mathsf W_{\bar a}=U_a^*\mathsf W_aU_a .
\]
La sustitución da
\(r_{\bar a}^*\mathsf W_{\bar a}r_{\bar a}=r_a^*\mathsf W_ar_a\).
Su derivada conserva necesariamente
\(\dot{\mathsf W}_{\bar a}
=U_a^*(A_a^*\mathsf W_a+\dot{\mathsf W}_a+\mathsf W_aA_a)U_a\):
eliminar este término alteraría la corriente. Si el peso acopla varias
aristas, se aplica la misma congruencia a la matriz completa con el
transporte diagonal por bloques; se conservan también sus bloques
cruzados.

\begin{proposition}[Compatibilidad variacional por refinamiento]
\label{hmtcuatro:ant:refinamiento-corriente}
Para mapas fijos \(J_m,K_m\), supóngase que, a lo largo de la
variación considerada,
\[
 D_{m+1}J_m=K_mD_m,\qquad
 K_m^*\mathsf W_{m+1}K_m=\mathsf W_m .
\]
Entonces \(E_{m+1}(J_mf)=E_m(f)\) y sus diferenciales coinciden.
Si además las acciones coinciden y \(P_m\) transporta las
variaciones de conexión al refinamiento, sus corrientes cumplen
\[
 s_m=P_m^{\mathsf T}s_{m+1},\qquad
 \mathsf M_m\sigma_m=P_m^{\mathsf T}\mathsf M_{m+1}\sigma_{m+1}.
\]
\end{proposition}
\begin{proof}
Sustituir \(D_{m+1}J_m=K_mD_m\) en la energía y después la
congruencia del peso prueba la primera igualdad para toda la
curva. Derivarla y utilizar
\(-\frac12\theta^{\mathsf T}s_m
=-\frac12(P_m\theta)^{\mathsf T}s_{m+1}\)
prueba las restantes. Si \(J_m\) varía, su diferencial aporta
\(2\operatorname{Re}\langle D_{m+1}J_mf,
\mathsf W_{m+1}D_{m+1}\dot J_m f\rangle\); no se descarta.
\end{proof}

La corriente de memoria así obtenida procede de su acción.
La corriente espinorial siguiente procede de la acción afín de
los campos. Se componen cuando son términos de una misma acción;
no se identifican sólo por compartir el transporte. En particular,
eliminar una contorsión de la que dependen también el campo
preparado, la incidencia o su mínimo exige conservar esos
diferenciales. La fórmula lineal de Cartan utilizada a continuación
corresponde al problema material afín expresamente declarado.

\subsection{Representación espinorial y corriente afín}
\label{hmtcuatro:ant:espinores}

Se fija la estructura de espín de la realización geométrica. Las
matrices heredadas, después de complejificar el módulo real, son
\begin{equation}
 \begin{gathered}
 J=\begin{pmatrix}0&-1\\1&0\end{pmatrix},\quad
 R=\begin{pmatrix}0&1\\1&0\end{pmatrix},\quad
 Z=\begin{pmatrix}1&0\\0&-1\end{pmatrix},\\
 g^0=J\otimes I_2,\quad g^1=R\otimes I_2,\quad
 g^2=Z\otimes R,\quad g^3=Z\otimes Z .
 \end{gathered}
 \label{hmtcuatro:ant:clifford}
\end{equation}
Sus cuadrados son \(-I_4,I_4,I_4,I_4\) y las matrices
distintas anticonmutan. De ello se sigue
\(\{g^I,g^J\}=2\eta^{IJ}I_4\). Definimos
\[
 g_I=\eta_{IJ}g^J,\quad
 \mathbb B_{IJ}=\tfrac14[g_I,g_J],\quad
 A_D=-ig^0,\quad \bar\psi=\psi^\dagger A_D,\quad
 \Omega=\tfrac12\omega^{IJ}\mathbb B_{IJ}.
\]
La expansión del conmutador de tres matrices da
\([\mathbb B_{IJ},g^K]=\delta_J^Kg_I-\delta_I^Kg_J\).
Por tanto \(\Omega\) representa la misma conexión métrica, con
\(D\psi=\mathrm d\psi+\Omega\psi\) y
\(D\bar\psi=\mathrm d\bar\psi-\bar\psi\Omega\).
La matriz \(A_D\) es hermítica y \(A_Dg^I\) es antihermítica;
éstas fijan el adjunto y el factor de la acción en firma \((-+++)\).

Sea \(\eta_I=\iota_{E_I}\operatorname{vol}_e\), de modo que
\(e^J\wedge\eta_I=\delta_I^J\operatorname{vol}_e\). La acción afín es
\begin{equation}
 S_D=\int_M\left\{
 \frac{\hbar}{2}
 [\bar\psi g^ID\psi-(D\bar\psi)g^I\psi]\wedge\eta_I
 -mc\,\bar\psi\psi\operatorname{vol}_e\right\}.
 \label{hmtcuatro:ant:accion-dirac}
\end{equation}
La masa y las secciones positivas \(c,\hbar\) son las lecturas
recibidas, y \([\psi]=L^{-3/2}\). La notación alternativa
\(\Gamma^I=-ig^I\), con término principal
\(i\hbar\Gamma^ID_I\), cambia a la firma de Clifford opuesta;
no se mezclan ambas convenciones.

\begin{proposition}[Corriente espinorial normalizada]
\label{hmtcuatro:ant:corriente-spin}
Con \(B^{IJK}=\bar\psi g^{[I}g^Jg^{K]}\psi\), donde la
antisimetrización está normalizada, la convención
\(\delta_\omega S_D=-\frac12\int\delta\omega^{JK}\wedge\sigma_{JK}\)
da
\begin{equation}
 \sigma_{JK}
 =-\frac{\hbar}{2}\bar\psi\{g^I,\mathbb B_{JK}\}\psi\,\eta_I
 =-\frac{\hbar}{2}B^I{}_{JK}\eta_I .
 \label{hmtcuatro:ant:sigma-spin}
\end{equation}
La corriente es independiente de la contorsión al fijar \(e,\psi\).
\end{proposition}
\begin{proof}
Se tiene
\(\delta\Omega=\frac12\delta\omega^{JK}\mathbb B_{JK}\).
Los dos términos cinéticos de \eqref{hmtcuatro:ant:accion-dirac}
dan
\[
 \delta_\omega S_D=\frac{\hbar}{4}\int
 \delta\omega^{JK}\wedge\eta_I\,
          \bar\psi\{g^I,\mathbb B_{JK}\}\psi .
\]
Comparar con la convención variacional determina signo y factor.
Al desarrollar el anticonmutador mediante Clifford, las
contracciones de grado uno se cancelan y queda
\(\{g^I,\mathbb B_{JK}\}=g^{[I}g_Jg_{K]}\).
La acción es lineal en \(\Omega\), lo que prueba la independencia.
\end{proof}

Si se incluyen conexiones internas, actúan sobre el factor
material del módulo espinorial, mientras que
\(\mathbb B_{IJ}\) actúa sobre el factor de espín. Así,
\([\,\mathbb B_{IJ}\otimes I,I\otimes\rho(A)\,]=0\).
Variar \(\omega\) produce la misma fórmula, contraída sobre
las componentes internas y sumada sobre los multipletes presentes.
Esto preserva sus representaciones quirales, de color y de familia;
no añade multipletes a los ya construidos.

\subsection{Contracción torsional y cruces entre especies}
\label{hmtcuatro:ant:contraccion}

La inversión real de Cartan--Holst y de la aplicación
\(T^I\mapsto T^I\wedge e^J-e^I\wedge T^J\), demostrada en
\ref{hmtcuatro:gea:torsion}, fija sin ambigüedad la contorsión
\(\mathfrak k_*\) para la corriente afín
\eqref{hmtcuatro:ant:sigma-spin}. En esas fórmulas se conserva
\(\gamma\in\mathbb R\setminus\{0\}\), una cotetrada no degenerada
y las secciones \(G,c,\hbar,\gamma\) fijas durante la variación.
La eliminación utiliza una sola vez la interacción lineal y da
\(-\frac14\mathfrak k_*^{IJ}\wedge\sigma_{IJ}\).

\begin{theorem}[Coeficiente de la contracción espinorial]
\label{hmtcuatro:ant:coeficiente}
Sea \(\kappa=8\pi G/c^4\). La interacción efectiva es
\[
 \mathcal L_{\rm spin}^{\rm ef}
 =C_\gamma q(B)\operatorname{vol}_e,\qquad
 C_\gamma=\frac{3\kappa c\hbar^2}{16}
                    \frac{\gamma^2}{1+\gamma^2},
\]
donde
\[
 q(B)=\frac1{3!}B^{IJK}B_{IJK}
 =-(B^{012})^2-(B^{013})^2-(B^{023})^2+(B^{123})^2 .
\]
\end{theorem}
\begin{proof}
Se calcula primero
\[
 \mathcal Y^{IJ}
 =\kappa c\frac{\gamma^2}{1+\gamma^2}
       (-\star+\gamma^{-1}I)\sigma^{IJ},\quad
 \mathcal B^I=\sum_J\iota_{E_J}\mathcal Y^{IJ},\quad
 T^I=\mathcal B^I-\tfrac14e^I\wedge\sum_L\iota_{E_L}\mathcal B^L,
\]
y después
\(\mathfrak k_{IJK}=(T_{KIJ}-T_{IJK}-T_{JKI})/2\).
Estas son composiciones lineales de las inversas ya demostradas.
Para explicitar todos los coeficientes de la contracción, se
ordena el trivector como \((012,013,023,123)\) y se extrae el factor
\(\kappa c\hbar^2\gamma^2/(1+\gamma^2)\). La sustitución de
\(\sigma_{JK}/\hbar=-\frac12B^I{}_{JK}\eta_I\) en esas composiciones
y en \(-\frac14\mathfrak k^{IJ}\wedge\sigma_{IJ}\) da
\[
 \begin{array}{c|rrrr|c}
  \text{operador sobre }\sigma&012&013&023&123&
           \text{seis coeficientes de polarización}\\ \hline
  -\star&-3/16&-3/16&-3/16&3/16&0\\
  I&0&0&0&0&0
 \end{array}
\]
Cada entrada usa \(e^I\wedge\eta_J=\delta_J^I\operatorname{vol}_e\)
y las contracciones escritas arriba. Las cuatro entradas diagonales
y las seis polarizaciones enumeran la forma cuadrática completa.
La parte \(\gamma^{-1}I\) se anula y la restante es
\((3/16)q(B)\), lo que prueba la fórmula.
\end{proof}

Para varias especies, la variación se efectúa sobre su acción
común: \(\sigma=\sum_s\sigma_s\). Como la respuesta
\(\mathfrak k_*=\mathsf C_e(\sigma)\) es lineal,
\begin{equation}
 -\tfrac14\mathsf C_e\!\left(\sum_s\sigma_s\right)
       \wedge\sum_t\sigma_t
 =-\tfrac14\sum_{s,t}\mathsf C_e(\sigma_s)\wedge\sigma_t .
 \label{hmtcuatro:ant:cruces-especies}
\end{equation}
Esta expansión demuestra que los términos \(s\ne t\) permanecen.
Calcular respuestas separadas y conservar sólo \(s=t\) cambiaría
la acción total.

\subsection{Orden normal de la corriente total}
\label{hmtcuatro:ant:car}

En la realización fermiónica finita de una celda se utiliza
\(\mathcal F=\bigoplus_n\Lambda^n V\), con las relaciones
\(\{a_r,a_s^\dagger\}=\delta_{rs}\), \(\{a_r,a_s\}=0\) y el
vacío de ocupación como referencia del orden normal.
Escribimos \(d\Gamma(M)=\sum_{r,s}a_r^\dagger M_{rs}a_s\);
esta \(\Gamma\) denota segunda cuantización, no la holonomía nonádica.
Para el módulo espinorial de cuatro componentes, las matrices son
\[
 \begin{array}{c|cccc}
 abc&012&013&023&123\\ \hline
 M_{abc}=A_Dg^ag^bg^c&
 iJ\otimes R&iJ\otimes Z&iI_2\otimes J&iZ\otimes J
 \end{array}.
\]
Multiplicar \eqref{hmtcuatro:ant:clifford} da estas matrices;
son hermíticas, de traza cero y cuadrado \(I_4\).

\begin{proposition}[Contracción normal y conservación de los cruces]
\label{hmtcuatro:ant:orden-normal}
Para toda matriz \(M\),
\[
 :d\Gamma(M)^2:
 =d\Gamma(M)^2-d\Gamma(M^2).
\]
En el módulo de cuatro componentes,
\[
 \widehat{\mathscr Q}_{\rm sp}
 =\sum_{a=1}^4s_a d\Gamma(M_a)^2+2\widehat N,
 \qquad s=(-1,-1,-1,+1),\quad \widehat N=d\Gamma(I_4).
\]
El operador es autoadjunto y conserva la ocupación.
En una suma de multipletes se emplea la misma primera identidad
con las matrices totales \(\widetilde M_a\); se conservan las
contracciones entre especies.
\end{proposition}
\begin{proof}
En \(d\Gamma(M)^2\), la relación
\(a_sa_t^\dagger=\delta_{st}-a_t^\dagger a_s\) separa la
contracción \(d\Gamma(M^2)\) del término con dos creadores y
dos aniquiladores. Este último es, por definición, el producto
normal. Como \(M_a^2=I_4\) y \(\sum_a s_a=-2\), la sustracción
de contracciones vale \(+2\widehat N\). La hermiticidad de las
matrices hace autoadjuntos los sumandos y cada monomio conserva
la ocupación.

Si \(\widetilde M=\bigoplus_s M_s\), entonces
\(d\Gamma(\widetilde M)=\sum_s d\Gamma(M_s)\) y
\(d\Gamma(\widetilde M^2)=\sum_s d\Gamma(M_s^2)\).
Los bilineales de multipletes distintos conmutan; al elevar la
primera suma al cuadrado quedan
\(2d\Gamma(M_s)d\Gamma(M_t)\) para \(s<t\). La sustracción de
una ocupación no elimina estos cruces. Es la versión operatoria
de \eqref{hmtcuatro:ant:cruces-especies}.
\end{proof}

La invariancia bajo los cambios unitarios internos se obtiene
porque éstos actúan sobre los índices materiales y entrelazan las
matrices de corriente: la segunda cuantización transporta la
fórmula completa por conjugación. La densidad y su evaluación
sobre un estado utilizan, además, el volumen de la celda y el
estado material declarados. No se identifican la corriente media
y su segundo momento. Estos datos completan los antecedentes
utilizados por la acción conjunta, sin introducir una segunda
eliminación de torsión ni modificar sus condiciones de borde.

% Preparación separada del 30-09-2026; no modifica manuscritos publicados.
% Propietarios: VI_ES/sections/04_atlas_familias.tex:140--163;
% md/04d_ocupacion_estadistica.tex:305--319;
% md/06e_color_qutrit_gauge_finito.tex:1--194.
% Carta qutrit y forma compacta recibidas con sus hipótesis explícitas.
% Inserción: después de 00_antecedentes_geometricos_materiales y antes de 10.
\section{Residuo cromático y representación interna de la acción conjunta}
\label{hmtcuatro:color:seccion}

La cadena APP--TRIT--TPK conserva el estado enriquecido, sus rutas,
orientaciones, hojas y memoria dentro de la estructura discreta conjunta
del continuo. El residuo que se lee a continuación es una publicación
finita de ese estado; no lo sustituye ni reinicia su historia. Las
coordenadas de fase, incluida la unidad compleja y la coordenada regional
\(\pi\), son salidas HMT recibidas antes de esta realización algebraica.
La representación se construye sobre el soporte residual con la carta,
la orientación y el carácter de fase declarados. No se infiere un grupo
de calibre a partir del solo número de etiquetas.

\subsection{Soporte residual y orientación conservada}
\label{hmtcuatro:color:residuo}

En la copia posicional de APP, sean
\[
 I_9=\{1,\ldots,9\},\qquad
 \mathcal C=I_9\times I_9,\qquad
 \mathcal C_q=\{(r,c)\in\mathcal C:r\text{ es par}\}.
\]
El observable residual y la inversión celular son
\[
 \kappa_{\rm col}(r,c)=r-c\pmod3,\qquad
 \rho_{\rm c}(r,c)=(10-r,10-c).
\]
El símbolo \(\kappa_{\rm col}\) no designa el coeficiente gravitatorio
de la acción. La elevación de este observable a las secciones
persistentes se efectúa por la proyección al atlas y conserva los
restantes registros; no selecciona una celda por su masa evaluada.

\begin{proposition}[Fibras ternarias y balance cromático]
\label{hmtcuatro:color:balance}
El funcional \(\kappa_3:\mathbb F_3^2\to\mathbb F_3\),
\(\kappa_3(r,c)=r-c\), tiene tres fibras afines de cardinal tres.
Sobre \(\mathcal C_q\), los tres residuos tienen cardinales
\(12,12,12\). La inversión celular niega el residuo.
\end{proposition}
\begin{proof}
Para cada \(j\in\mathbb F_3\), la fibra es
\(\{(c+j,c):c\in\mathbb F_3\}\), una traslación del núcleo
\(\{(c,c)\}\). En \(\mathcal C_q\) existen cuatro valores pares
de \(r\); para cada uno, \(c=1,\ldots,9\) recorre tres veces
cada clase módulo tres. Cada residuo posee \(4\cdot3=12\)
representantes. Finalmente,
\[
 \kappa_{\rm col}(\rho_{\rm c}(r,c))
 =(10-r)-(10-c)=-(r-c)\pmod3.
\]
La clase cero queda fija y las clases uno y dos se intercambian.
\end{proof}

Este balance no identifica sabor, profundidad y color. Tampoco afirma
que las tres rectas que se construirán sean tres estados completos:
las rutas y la memoria que comparten una lectura residual permanecen
en la fibra enriquecida.

\subsection{Carta qutrit, par de Weyl y álgebra de color}
\label{hmtcuatro:color:weyl}

Se fija la carta
\[
 V_{\rm col}=\mathbb C[\mathbb F_3]\simeq\mathbb C^3,
 \qquad (e_0,e_1,e_2),\qquad
 \omega_3=e^{2\pi i/3},
\]
con base delta ortonormal, orientación \(j\mapsto j+1\) y carácter
\(j\mapsto\omega_3^j\). Los operadores son
\[
 Xe_j=e_{j+1\bmod3},\qquad Ze_j=\omega_3^j e_j,
 \qquad X^3=Z^3=I,\quad ZX=\omega_3XZ.
\]
La última igualdad se comprueba aplicando ambos miembros a cada
\(e_j\). La aplicación del atlas a esta carta es
\[
 (r,c)\longmapsto\mathbb C e_{\kappa_{\rm col}(r,c)}.
\]
Así, \(X\) permuta las rectas y \(Z\) registra sus fases. Cambiar
la carta afín y transportar conjuntamente base, orientación y carácter
conjuga los operadores por la correspondiente matriz unitaria de
cambio de base. No equivale a cambiar sólo una etiqueta y dejar fijos
los demás datos.

\begin{theorem}[Álgebra matricial y forma real compacta]
\label{hmtcuatro:color:algebra}
Los nueve operadores \(W_{ab}=X^aZ^b\),
\((a,b)\in\mathbb F_3^2\), forman una base ortogonal de
\(\operatorname{End}_{\mathbb C}(V_{\rm col})\). Los ocho índices
no nulos generan linealmente \(\mathfrak{sl}_3(\mathbb C)\), cuya
forma real de matrices antihermíticas es \(\mathfrak{su}(3)\).
\end{theorem}
\begin{proof}
Si \(a\ne0\), la matriz \(X^aZ^b\) no tiene entradas diagonales.
Si \(a=0\), su traza es \(1+\omega_3^b+\omega_3^{2b}\), que
vale cero para \(b\ne0\) y tres para \(b=0\). La relación de
Weyl da entonces
\[
 \operatorname{tr}(W_{ab}^{\dagger}W_{cd})
 =3\delta_{ac}\delta_{bd}.
\]
La independencia y la dimensión nueve prueban la primera afirmación;
los ocho restantes forman la base del subespacio de traza cero.
Además,
\[
 W_{ab}^{\dagger}=\omega_3^{ab}W_{-a,-b}.
\]
Los ocho índices no nulos forman cuatro pares adjuntos. Tomando
\(\mathcal R=\{(1,0),(0,1),(1,1),(1,2)\}\), las matrices
\[
 A_u=W_u-W_u^{\dagger},\qquad
 B_u=i(W_u+W_u^{\dagger}),\qquad u\in\mathcal R,
\]
son antihermíticas y de traza cero. Su independencia real se deduce
de la independencia compleja de los ocho \(W_u,W_{-u}\): en cada
par, los coeficientes de una combinación real nula fuerzan a cero
los dos coeficientes reales. Forman una base real de dimensión ocho.
Por tanto,
\[
 \{A\in\mathfrak{sl}_3(\mathbb C):-A^{\dagger}=A\}
 =\mathfrak{su}(3).
\]
\end{proof}

La forma hermítica y la orientación compleja fijadas determinan el
grupo de realización
\[
 SU(V_{\rm col})=\{U:U^{\dagger}U=I,\ \det U=1\}.
\]
Su álgebra tangente es exactamente la anterior: diferenciar ambas
ecuaciones en la identidad da \(A^{\dagger}=-A\) y
\(\operatorname{tr}A=0\); recíprocamente, \(e^{tA}\) satisface
ambas. Todo elemento del grupo se escribe como exponencial de una
matriz de esa álgebra. En efecto, diagonalícese unitariamente \(U\)
con autovalores \(e^{i\theta_j}\); como el determinante es uno,
\(\sum_j\theta_j=2\pi n\). Restar \(2\pi n\) a uno de los
ángulos no altera \(U\) y produce un logaritmo antihermítico de
traza cero. Esta realización continua no se identifica con el grupo
finito generado únicamente por \(X,Z\). Conserva la forma compacta
y no introduce un valor para el acoplamiento fuerte.

\subsection{Transporte cromático y factores internos independientes}
\label{hmtcuatro:color:transporte}

En una celularización finita orientada compatible con las rutas,
se asignan \(U_e\in SU(V_{\rm col})\) y
\(U_{\bar e}=U_e^{-1}\). Para \(g_v\in SU(V_{\rm col})\),
\[
 U_e^g=g_{t(e)}U_eg_{s(e)}^{-1},\qquad
 \psi_v^g=g_v\psi_v.
\]
Por sustitución,
\[
 U_e^g\psi_{s(e)}^g-\psi_{t(e)}^g
 =g_{t(e)}(U_e\psi_{s(e)}-\psi_{t(e)}).
\]
El signo opuesto, utilizado en \(D_T\), tiene la misma covariancia.
La transformación de un producto orientado alrededor de una cara
es conjugación por el calibre del vértice base: los factores de
los vértices intermedios se cancelan consecutivamente. En particular,
para \(\beta\ge0\),
\[
 S_{\rm col}^{(\mathcal K,\beta)}[U]
 =\frac\beta3\sum_f(3-\operatorname{Re}\operatorname{tr}U_f)
\]
es invariante y no negativa. La invariancia usa la traza cíclica;
la positividad resulta de \(\operatorname{Re}\operatorname{tr}U_f\le3\).
El parámetro y la celularización conservan su condición de datos de
esta realización. Esta prueba no los determina por comparación con
valores cromodinámicos.

La materia de la acción conjunta utiliza este factor, no otra
partición del atlas. Para los quarks izquierdos, su fibra interna
tiene factores \(V_{\rm col}\otimes\mathbb C^2\otimes F\),
donde \(F\) conserva las familias; para cada tipo derecho se usa
\(V_{\rm col}\otimes F\). El factor de color es trivial en los
leptones. El factor espinorial se conserva aparte. Para matrices
\(A\) de color y \(B\) débiles se cumple
\[
 [A\otimes I,I\otimes B]=0,
\]
pues los dos productos son \(A\otimes B\). Los caracteres de
hipercarga actúan como escalares sobre cada multiplete y también
conmutan. Esta identidad permite componer las conexiones internas
y la de espín en la misma derivada sin confundir color, quiralidad
y sabor.

Los mapas Yukawa de la acción posterior son la identidad en el
factor cromático y mapas en los factores Higgs--familia. Por ello
\(Y_q(H)(g_{\rm col}\otimes I)
=(g_{\rm col}\otimes I)Y_q(H)\). La compatibilidad débil y de
hipercarga se demuestra allí con sus columnas y sus pesos explícitos;
no se obtiene del censo \(12+12+12\).
Así se conserva la cadena residual--qutrit--álgebra--transporte
que requiere la acción común, sin afirmar en este antecedente un
límite cuántico de Yang--Mills ni reemplazar sus propietarios.

% Fragmento integrable en VII: Relaciones estructurales entre las constantes físicas.
% Insertar después de desarrollos/normalizacion_conjunta.tex.
% Fuentes y cobertura íntegra: editorial/INSERCION_VII.md.
% Sin preámbulo ni macros; usa amsmath y los entornos ya definidos por VII.
\section{Gravedad, energía y autoinercia del sistema conjunto}
\label{hmtcuatro:gea:seccion}

La normalización conjunta permite formular una pregunta operatoria:
cómo intervienen la geometría, la materia y la memoria en una misma
realización. Se reúnen tres enlaces. El carácter positivo conserva
simultáneamente sus lecturas energética, másica y radial; la variación
del transporte conjunto produce corrientes geométricas e internas de
una misma energía; y la respuesta autoinercial se expresa mediante el
radio gravitatorio de ese carácter. Se demuestran la conservación de
estos enlaces bajo reducción de memoria y la propagación de las
restricciones de la acción acoplada durante su evolución regular.

El antecedente es el núcleo común del artículo. APP produce las dos
hojas con sus residuos, cocientes e incidencias; TRIT conserva régimen
y orientación: $+1$ corresponde a retorno y memoria, $0$ al umbral
presente y $-1$ a apertura. TPK selecciona, transporta, pliega, retorna
y conserva el estado enriquecido y su historia. La prolongación
$w_6\to w_{12}\to w_{18}\to w_{24}\to w_{30}\to R_{36}\to G_9$
mantiene prefijo, acarreo, supervivencia, ruta y frontera. La dinámica,
su holonomía nonádica y la publicación reducida tienen el orden
$U_t\to\Gamma_9\to K_{\rm ph}$: retorna la fase, avanza la memoria
y no se reinicia el estado. Las cinco construcciones consustanciales
de la estructura discreta del continuo no se separan al realizar sus lectores.

Las constantes y escalas empleadas son salidas de esta genealogía.
La cuadrirrelación $(\pi,\varphi,e,\alpha)$ conserva su origen común
y su orden constructivo interno; acción, velocidad constitutiva y
Barbero--Immirzi se reciben con sus lectores. El lenguaje operatorio
y variacional compone y reconoce esas salidas, sin introducir valores
metrológicos para seleccionar estados o coeficientes.

\subsection{Energía, masa y radio del mismo carácter}
\label{hmtcuatro:gea:caracter}

Se conserva la realización marcada de la sección anterior:
\[
\Omega=\{(j,k,q,u,v):j,k\in\mathbb Z/12\mathbb Z,\ 
k-j\in\{0,3,6,9\},\ 1\leq q\leq8,\ 
u<v,\ u,v\in\{1,2,4,5,7,8\}\}.
\]
Su cardinal $N_\Omega=12\cdot4\cdot8\binom62=5760$ cuenta incidencias,
no estados de materia ni lapsos. Para el inversor transportado
$B:V_{\rm in}\to V_{\rm out}$, con $B^*B=BB^*=I/4$, el proyector
$P=uu^*$ del modo común tiene $|u_i|^2=1/N_\Omega$. El archivo retiene
\[
C_0=(I-P)B,\qquad M_0=PB,\qquad C_0+M_0=B,\qquad\mathsf U=2B.
\]
Las imágenes de $C_0,M_0$ son ortogonales y $\mathsf U$ es unitario.
La expectativa marcada $\Delta_\Omega(A)=\sum_iP_iAP_i$,
$P_i=e_ie_i^*$, es única entre las expectativas bimodulares hacia
el álgebra diagonal que fijan los $P_i$: sobre una unidad matricial,
la bimodularidad exige
$\mathcal E(E_{ij})=P_i\mathcal E(E_{ij})P_j$, que es cero si $i\ne j$
y vale $E_{ii}$ si $i=j$.

De $C_0C_0^*=(I-P)/4$ y $P_iPP_i=P_i/N_\Omega$ resulta
\begin{equation}
\Delta_\Omega(C_0C_0^*)=r_\Omega I,\qquad
\Delta_\Omega(M_0M_0^*)=\frac{I}{4N_\Omega},\qquad
r_\Omega=\frac{5759}{23040}.
\label{hmtcuatro:gea:retorno}
\end{equation}
Ambas lecturas suman $I/4$. La regla constitutiva longitudinal aplica
linealmente el primer coeficiente a una sección de longitud:
\begin{equation}
L=\alpha^{16}r_\Omega L_*,\qquad D=L\mathsf U.
\label{hmtcuatro:gea:longitud}
\end{equation}
La norma local $\alpha^{16}$ y la sección $L_*$ son antecedentes
explícitos; tomar una raíz de intensidad sería otro lector.
La regla métrica radial es $\|R_Xv\|^2=\|XD^*v\|^2$.
Para $X=X^*>0$ en la fibra finita, la polarización y la raíz positiva
única dan $R_X^2=DX^2D^*=L^2\mathsf UX^2\mathsf U^*$ y
$R_X=LY$, con $Y=\mathsf UX\mathsf U^*>0$.
El espectro no se selecciona a partir de una respuesta gravitatoria deseada.

La misma sección $\hbar>0$ y el reloj conjugado $\omega L=c$
determinan $E_L=\hbar c/L$ y los lectores
\begin{equation}
H_X=E_LY,\qquad M_X=H_X/c^2,\qquad
R_X=LY,\qquad\Lambda_X=LY^{-1}.
\label{hmtcuatro:gea:lectores}
\end{equation}
La notación $H_X$ queda reservada a este lector del carácter.
El Hamiltoniano conjunto $\mathbb H$ introducido después tiene otro
dominio y no se identifica con $H_X$ por compartir el nombre de energía.

\begin{proposition}[Coeficiente radial común]
\label{hmtcuatro:gea:coeficiente}
En la identificación explícita de $R_X$ como radio gravitatorio reducido,
$R_X=GM_X/c^2$, existe un único coeficiente compatible:
\begin{equation}
G=\frac{c^3L^2}{\hbar},\qquad
R_X=\frac{G}{c^4}H_X=\frac{G}{c^2}M_X,\qquad
R_X\Lambda_X=L^2I,\quad H_X\Lambda_X=\hbar cI.
\label{hmtcuatro:gea:radial}
\end{equation}
\end{proposition}
\begin{proof}
Multiplicar los lectores prueba los dos productos. La identificación
radial equivale a $LY=G\hbar Y/(c^3L)$.
Cancelar el carácter invertible determina $G$; la sustitución inversa
verifica la igualdad para cada carácter positivo. Si dos coeficientes
la satisfacen, su diferencia multiplicada por $Y$ es cero y ambos coinciden.
\end{proof}

Las reglas longitudinal y métrica fijan la realización; la convención
de radio reducido fija el significado físico de $G$. El radio
Schwarzschild de la misma masa es $2R_X$. La igualdad conserva
superposiciones y elementos no diagonales en cualquier base transportada,
no sólo valores medios. Para una familia diferenciable y una conexión
de marco común $\mathcal A_t$, con $G,c$ fijos,
\begin{equation}
\mathcal D_tR_X=\frac{G}{c^4}\mathcal D_tH_X,\qquad
\mathcal D_tM_X=c^{-2}\mathcal D_tH_X,\qquad
\mathcal D_tO=\dot O+[\mathcal A_t,O].
\label{hmtcuatro:gea:derivada}
\end{equation}
Se obtiene diferenciando la identidad completa, incluido el marco.
No implica inmovilidad del estado.

Para modos positivos, o para operadores en factores tensoriales distintos,
\begin{equation}
m(y)=\frac{\hbar}{cL}y,\qquad
\frac{Gm_1m_2}{\hbar c}=y_1y_2,\qquad
\frac{R(y)}{\bar\lambda(y)}=y^2,\quad \bar\lambda(y)=L/y.
\label{hmtcuatro:gea:adimensional}
\end{equation}
La sustitución cancela las secciones dimensionales; los operadores en
factores distintos conmutan. El tamaño adimensional depende del carácter,
sin otro parámetro libre por especie. No se deduce de ello una ley
espacial de fuerza ni una evaluación numérica sin evaluar el carácter.

El calendario conserva $t_P=L/c$, $t_0=\pi L/(54c)$ y
$108t_0=2\pi t_P$. Por tanto $L=(54/\pi)\ell_0$, con $\ell_0=ct_0$;
las dos longitudes no se identifican. Entre secciones
$\hbar_b=\rho\hbar_a$, $\rho>0$, con $L,c,y_i$ fijos, se tiene
$G_b=G_a/\rho$, $m_{i,b}=\rho m_{i,a}$ y la cantidad
$Gm_1m_2/(\hbar c)$ permanece invariante. Se comparan secciones
conjuntas, no una variación temporal con las demás coordenadas
dimensionales mantenidas artificialmente fijas.

\subsection{Memoria estática, dinámica y dominio espectral}
\label{hmtcuatro:gea:memoria}

En dimensión finita, o con bloques acotados y coercivos, escribamos
\[
Y=\begin{pmatrix}A&F\\F^*&C\end{pmatrix}>0,\quad
Y_{\rm ef}=A-FC^{-1}F^*,\quad \eta=y+C^{-1}F^*b.
\]
Completar el cuadrado demuestra
\begin{equation}
\left\langle\binom b y,Y\binom b y\right\rangle
=\langle b,Y_{\rm ef}b\rangle+\langle\eta,C\eta\rangle.
\label{hmtcuatro:gea:schur-prueba}
\end{equation}
De aquí $Y_{\rm ef}>0$, la extensión minimizante
$y=-C^{-1}F^*b$ y la reconstrucción $y=\eta-C^{-1}F^*b$.
El residuo $\eta$ se conserva. Como
$\operatorname{Schur}(sY)=sY_{\rm ef}$ para $s>0$, se obtienen
$H_{X,\rm ef}=E_LY_{\rm ef}$ y $R_{X,\rm ef}=LY_{\rm ef}$.
Resolver $Y(x,y)=(f,0)$ muestra que el bloque de frontera de $Y^{-1}$
es $Y_{\rm ef}^{-1}$, luego $\Lambda_b=LY_{\rm ef}^{-1}$.
La compresión directa de $Y$ daría $A$: no es la misma operación.

Sea $\chi=L/E_L=G/c^4>0$ y $R_X=\chi H_X$.
En una descomposición compatible con los dominios, definimos
\[
\mathcal K_{H_X}(z)=H_{bb}-zI-V(H_{ii}-zI)^{-1}V^*.
\]
\begin{proposition}[Memoria dinámica completa]
\label{hmtcuatro:gea:resolvente-prop}
En el dominio resolvente del bloque interior,
\begin{equation}
\mathcal K_{R_X}(\chi z)=\chi\mathcal K_{H_X}(z).
\label{hmtcuatro:gea:resolvente}
\end{equation}
Si ambos complementos son invertibles, sus inversas llevan el factor
$\chi^{-1}$.
\end{proposition}
\begin{proof}
Cada bloque radial es $\chi$ veces el energético y
$(\chi H_{ii}-\chi zI)^{-1}=\chi^{-1}(H_{ii}-zI)^{-1}$.
La sustitución deja un factor $\chi$ en cada término, conservando el
orden. Invertir prueba la última afirmación.
\end{proof}
El parámetro energético $z$ se transporta a la longitud $\chi z$;
la identidad conserva cada orden de memoria espectral.
Para $|z|<\|H_{ii}^{-1}\|^{-1}$, la serie del resolvente da
\[
\mathcal K_{H_X}(z)=H_{\rm est}-zZ+O(z^2),\quad
H_{\rm est}=H_{bb}-VH_{ii}^{-1}V^*,\quad
Z=I+VH_{ii}^{-2}V^*\geq I.
\]
La positividad sigue de
$Z-I=(H_{ii}^{-1}V^*)^*(H_{ii}^{-1}V^*)$.
La extensión estática $(b,-H_{ii}^{-1}V^*b)$ tiene norma cuadrada
$\langle b,Zb\rangle$: energía y norma temporal proceden de la misma extensión.

Si $T=Z^{1/2}$ es acotado e invertible, el transporte es dual:
\begin{equation}
H_{X,c}=T^{-*}H_{X,\rm ef}T^{-1},\quad
R_{X,c}=T^{-*}R_{X,\rm ef}T^{-1},\quad
\Lambda_c=T\Lambda_bT^*.
\label{hmtcuatro:gea:dual}
\end{equation}
Multiplicar prueba $R_{X,c}=\chi H_{X,c}$,
$R_{X,c}\Lambda_c=L^2I$ y $H_{X,c}\Lambda_c=\hbar cI$,
sin conmutatividad de $T$ con el carácter.
Si $\psi=T(t)b$, se conserva $\dot\psi=\dot Tb+T\dot b$.
El término temporal de conexión del generador no se incorpora
retrospectivamente al carácter radial sin declarar su mapa.

\begin{proposition}[Extensión espectral y refinamiento compatible]
\label{hmtcuatro:gea:limite}
Sea $X$ autoadjunto, positivo e inyectivo, y $\mathsf U$ unitario.
Los lectores de \eqref{hmtcuatro:gea:lectores}, definidos por cálculo
espectral de $Y=\mathsf UX\mathsf U^*$, satisfacen
$R_X=\chi H_X$ en $\operatorname{Dom}(Y)$.
Sus productos con $\Lambda_X$ valen en $\operatorname{Dom}(Y^{-1})$
y tienen las prolongaciones acotadas de \eqref{hmtcuatro:gea:radial}.
Si inscripciones isométricas satisfacen $X'J=JX$,
$\mathsf U'J=K\mathsf U$ y conservan $L,E_L$, entonces
$R'_XK=KR_X$ y $H'_XK=KH_X$ en los dominios transportados.
\end{proposition}
\begin{proof}
Sobre $P_{\varepsilon,M}=\mathbf1_{[\varepsilon,M]}(Y)$, lectores e
inversos son acotados y las identidades son multiplicaciones de
funciones de $Y$. Para $\psi\in\operatorname{Dom}(Y)$,
\[
\|Y(P_{\varepsilon,M}-I)\psi\|^2
=\int_{(0,\infty)\setminus[\varepsilon,M]}
\lambda^2\,d\langle\psi,E_Y(\lambda)\psi\rangle\longrightarrow0.
\]
La inyectividad elimina masa espectral en cero y también
$P_{\varepsilon,M}\psi\to\psi$. Es convergencia en norma de grafo.
Si $\psi\in\operatorname{Dom}(Y^{-1})$, entonces
$Y^{-1}\psi\in\operatorname{Dom}(Y)$ y $YY^{-1}\psi=\psi$,
lo que prueba los productos y sus prolongaciones.
Sustituir los entrelazamientos en $R'_X=L\mathsf U'X'\mathsf U'^*$
prueba $R'_XK=KR_X$; la cuenta energética es igual.
\end{proof}
El sector nulo permanece separado. Se precisa así la extensión no
acotada de la normalización anterior. Una proyección ortogonal
arbitraria no queda habilitada para las fórmulas de bloques:
en operadores no acotados se requieren dominios compatibles o formas
cerradas con bloques controlados.

\subsection{Retroacción en el mismo transporte}
\label{hmtcuatro:gea:retroaccion}

El registro geométrico $x$, los enlaces internos $u$ y los campos $f$
proceden del mismo estado enriquecido. En una celularización orientada,
\begin{equation}
E(f,x,u)=\langle r,W(x,u)r\rangle,\quad r=D_Tf,\quad
(D_Tf)_a=f_{t(a)}-T_af_{s(a)},\quad
T_a=U_a^{\rm sp}(x)\otimes R_{\rm int}(u_a).
\label{hmtcuatro:gea:energia-memoria}
\end{equation}
$W=W^*>0$ conserva los bloques entre incidencias. Espín y
representación interna tienen factores distintos; la realización
quiral conserva sus bloques y el mapa Higgs--Yukawa equivariante.

\begin{proposition}[Diferencial de una misma energía]
\label{hmtcuatro:gea:variacion}
Con campos fijos y $q=Wr$,
\begin{align}
\delta E&=-2\operatorname{Re}\sum_a
\langle q_a,\delta T_af_{s(a)}\rangle+\langle r,\delta Wr\rangle,
\label{hmtcuatro:gea:diferencial}\\
\delta T_a&=\delta U_a^{\rm sp}\otimes R_{\rm int}(u_a)
+U_a^{\rm sp}\otimes\delta R_{\rm int}(u_a).\nonumber
\end{align}
\end{proposition}
\begin{proof}
Derivar la forma cuadrática da dos términos conjugados por
$W=W^*$ y el término $\langle r,\delta Wr\rangle$.
Usar $\delta r_a=-\delta T_af_{s(a)}$ prueba la primera fórmula;
la regla del producto prueba la segunda.
Ambas variaciones actúan sobre los mismos $r,q$, sin corrientes
elegidas independientemente.
\end{proof}

\subsubsection{Un espacio cuántico y un operador conjunto}
\label{hmtcuatro:gea:cuantica}

Se fija una celularización finita, un conjunto finito de registros
geométricos $\mathcal X$ y una Fock fermiónica finita $\mathcal F$.
Los enlaces recorren el producto compacto $\mathcal G^E$ recibido;
los dobletes $\Phi_v\in\mathbb C^2$ no tienen corte de amplitud.
Las medidas invariantes $d\mu_x=Z_x^{-1}e^{-S_x}d\mu_0$ tienen
densidades positivas acotadas con inversa acotada y no dependen de
$\Phi$ en esta carta. El espacio es
\[
\mathscr H=\bigoplus_{x\in\mathcal X}
L^2(\mathcal G^E\times\mathbb C^{2|V|},
\mu_x\otimes d\Phi;\mathcal F).
\]
Las identificaciones $A_xf=\sqrt{Z_x}e^{S_x/2}f$ son unitarias
desde la medida común y transportan el reloj a
$H_{\rm clk}^{\mu}=A(H_{\rm clk}\otimes I)A^*$.
El reloj mezcla registros geométricos del mismo espacio.
El operador se realiza mediante la forma de
\begin{equation}
\mathbb H=H_{\rm clk}^{\mu}+\hbar cK_{\rm gauge}
+d\Gamma(D_T^*WD_T)+H_{\Phi,\rm cin}+V_\Phi
+\widehat H_Y+\widehat H_{\rm tor}.
\label{hmtcuatro:gea:total}
\end{equation}
El mapa Yukawa es lineal en $\Phi,\bar\Phi$ y satisface
$Y(g\Phi)\rho_R(g)=\rho_L(g)Y(\Phi)$.
La contorsión eliminada no se conserva simultáneamente como variable
independiente en los términos base de esta carta efectiva.

La cinética de momentos escalares tiene forma
$\int\langle\nabla_\Phi\Psi,A(x,u)\nabla_\Phi\Psi\rangle$,
con $A$ acotada, uniformemente positiva e independiente de $\Phi$.
Se exige $A(x,u^g)=O_gA(x,u)O_g^{\mathsf T}$ para la acción real
$O_g$ de los dobletes; positividad sola no prueba esa covarianza.
La energía espacial escalar es un potencial cuadrático positivo
acotado por $C_{\rm grad}r^2$, $r^2=\sum_v|\Phi_v|^2$,
no un operador acotado sobre todo $L^2$.
Los volúmenes satisfacen
$0<w_{\min}\leq w_v(x)\leq w_{\max}$ y $\lambda_\Phi>0$.
Reloj, bloque gauge, transportes, $W$ y coeficientes torsionales
son acotados en la carta. El potencial conserva su valor absoluto:
\begin{equation}
V_\Phi=\sum_vw_v\lambda_\Phi(|\Phi_v|^2-v_0^2/2)^2
-\sum_vw_v\lambda_\Phi v_0^4/4.
\label{hmtcuatro:gea:potencial}
\end{equation}
La última contribución depende del volumen y no se sustrae.

\begin{theorem}[Forma cerrada y evolución conjunta]
\label{hmtcuatro:gea:forma}
La forma de \eqref{hmtcuatro:gea:total} es cerrada y acotada
inferiormente en
$\mathcal Q=\{\Psi\in\mathscr H:\nabla_\Phi\Psi\in L^2,\
r^2\Psi\in L^2\}$.
Determina un único operador autoadjunto asociado $\mathbb H$
y la evolución unitaria $e^{-it\mathbb H/\hbar}$.
\end{theorem}
\begin{proof}
Cauchy--Schwarz da
$\sum_vw_v\lambda_\Phi|\Phi_v|^4\geq ar^4$,
$a=\lambda_\Phi w_{\min}/|V|>0$.
Fock finita y Yukawa lineal dan
$\|\widehat H_Y(\Phi)\|\leq C_Yr$; la torsión está acotada.
Para $\varepsilon>0$,
\[
C_Yr\leq\varepsilon r^4+
\frac{3C_Y^{4/3}}{4^{4/3}\varepsilon^{1/3}},\qquad
br^2\leq\varepsilon r^4+\frac{b^2}{4\varepsilon}.
\]
La primera cota maximiza $C_Yr-\varepsilon r^4$; la segunda
completa el cuadrado en $r^2$. Aplicarlas con $\varepsilon=a/4$
a Yukawa y al cuadrático negativo da
\[
q(\Psi)\geq c_1\|\nabla_\Phi\Psi\|^2+
(a/2)\|r^2\Psi\|^2-C\|\Psi\|^2,\qquad c_1>0.
\]
La norma
$\|\Psi\|^2+\|\nabla_\Phi\Psi\|^2+\|r^2\Psi\|^2$
es completa por cierre de la derivada débil y la multiplicación.
Estas cotas y las superiores hacen equivalente la norma de forma,
tras sumar una constante, a esa norma completa.
El teorema de representación de formas cerradas produce el operador
asociado y su cálculo espectral da la evolución.
La unicidad se refiere a esta forma y dominio, no a cada dominio
formal alternativo.
\end{proof}

Covarianza de $D_T,W$, invariancia de medidas, identidad Yukawa y
contracción torsional interna conservan $q$ y $\mathcal Q$.
La unicidad del operador asociado implica que la representación
compacta conmuta con $\mathbb H$; su promedio de Haar reduce
el operador y su evolución. No es un promedio correctivo posterior.

Se retira el corte computacional con bloques completos de Peter--Weyl
y capas isotrópicas completas de Hermite.
Regularización y truncamiento suave aproximan derivadas y peso
$r^2$; las densidades acotadas preservan normas equivalentes.
Su unión es densa en $\mathcal Q$. Para las proyecciones ortogonales
$P_n$ en las medidas $\mu_x$ y los operadores de Galerkin
$\mathbb H_n$,
\[
(\mathbb H_n+\lambda)^{-1}P_nf\longrightarrow
(\mathbb H+\lambda)^{-1}f\qquad(\lambda>C).
\]
La ecuación coerciva de forma y su versión de Galerkin dan
ortogonalidad del error; continuidad y coercividad lo acotan por
el mejor error de aproximación, que tiende a cero.
Se retira el corte escalar y funcional de enlace de esta carta,
no el espacial, geométrico ni fermiónico.
Autoadjunción tampoco implica clase traza de $e^{-t\mathbb H}$.

\subsubsection{Intercambio y conservación}
\label{hmtcuatro:gea:intercambio}

Para $V=\sum_xP_x\otimes B_x$ y un elemento $h_{xy}$ del reloj,
multiplicar bloques da
\begin{equation}
[H_{\rm clk}\otimes I,V]_{xy}=h_{xy}(B_y-B_x).
\label{hmtcuatro:gea:conmutador}
\end{equation}
Si $h_{xy}\ne0$ y $B_y\ne B_x$, hay intercambio
geométrico--material. Si $\mathbb H=A+B+V$, para operadores
acotados o un núcleo común que justifique las derivadas,
\[
\frac d{dt}\langle A\rangle=\frac i\hbar\langle[\mathbb H,A]\rangle,
\quad
\frac d{dt}\langle B\rangle=\frac i\hbar\langle[\mathbb H,B]\rangle,
\quad
\frac d{dt}\langle V\rangle=\frac i\hbar\langle[\mathbb H,V]\rangle.
\]
La suma es cero por $[\mathbb H,\mathbb H]=0$.
Se conserva la energía incluyendo interacción, no cada sector
por separado. La geometría está en el espacio cuántico, no añadida
como fondo externo. Las historias ampliadas retienen memoria y retorno;
su refinamiento no equivale por definición al refinamiento espacial.

\subsection{La misma acción y sus corrientes}
\label{hmtcuatro:gea:accion}

La realización geométrica emplea cotetrada no degenerada $e$,
conexión de Lorentz $\omega$ y conexiones internas de la misma
materia. Sean $\Sigma^{IJ}=e^I\wedge e^J$,
$P_\gamma=\star+\gamma^{-1}I$,
$\gamma\in\mathbb R\setminus\{0\}$ y $\star^2=-I$ en bivectores.
Los acoplamientos permanecen fijos durante la variación.
Con $\kappa=8\pi G/c^4$, la misma acción es
\begin{equation}
S_{\rm tot}=\frac{\hbar}{16\pi L^2}
\int\Sigma_{IJ}\wedge(P_\gamma F_\omega)^{IJ}
-\frac{\Lambda\hbar}{8\pi L^2}\int\operatorname{vol}_e+S_m,
\label{hmtcuatro:gea:accion-total}
\end{equation}
donde $S_m=S_{\rm YM}+S_\Phi+S_{\rm Weyl}+S_Y$ contrae sus
índices con la misma $e$, que se varía.
Los coeficientes siguen de
$(2\kappa c)^{-1}=\hbar/(16\pi L^2)$.

Con $\delta_\omega S_m=-\tfrac12\int
\delta\omega^{IJ}\wedge\sigma_{IJ}$,
$\delta F=D_\omega\delta\omega$ y la integración por partes,
\begin{equation}
D_\omega(P_\gamma\Sigma)=\kappa c\sigma.
\label{hmtcuatro:gea:cartan}
\end{equation}
Se usan variaciones interiores o la completación de borde compatible.
Corriente de espín y fuente métrica son derivadas de la misma
materia, no dos entradas independientes.

\subsubsection{Eliminación torsional y términos cruzados}
\label{hmtcuatro:gea:torsion}

La inversa real es
$P_\gamma^{-1}=\gamma^2(-\star+\gamma^{-1}I)/(1+\gamma^2)$:
el producto de los dos polinomios en $\star$ antes de normalizar
vale $(1+\gamma^{-2})I$.
Denótese por $\mathcal Y^{IJ}=\kappa c(P_\gamma^{-1}\sigma)^{IJ}$
la forma bivectorial de grado tres, distinta del carácter $Y$.
Para $T^I=D_\omega e^I$,
$D_\omega\Sigma^{IJ}=T^I\wedge e^J-e^I\wedge T^J$.
Con el marco dual $E_I$, su inversa es
\[
\mathcal B^I=\sum_J\iota_{E_J}\mathcal Y^{IJ},\qquad
\vartheta=\tfrac14\sum_I\iota_{E_I}\mathcal B^I,\qquad
T^I=\mathcal B^I-e^I\wedge\vartheta.
\]
En efecto, si $t=\sum_J\iota_{E_J}T^J$, las contracciones dan
$\mathcal B^I=T^I+e^I\wedge t$ y después $4t$.
El mapa tiene núcleo nulo entre espacios de dimensión veinticuatro,
luego es un isomorfismo. Para
$\omega=\mathring\omega(e)+\mathfrak k$,
\[
T_{IJK}=\mathfrak k_{IKJ}-\mathfrak k_{IJK},\qquad
\mathfrak k_{IJK}=\tfrac12(T_{KIJ}-T_{IJK}-T_{JKI}).
\]
La combinación cíclica y la antisimetría en los dos primeros índices
prueban estas fórmulas.
Para materia afín en $\mathfrak k$, con corriente independiente
de ella, hay una única $\mathfrak k_*[e,\sigma]$ lineal en
la corriente total $\sigma=\sum_s\sigma_s$.

La expansión de curvatura conserva el borde
\[
S_\partial[e,\mathfrak k]=\frac1{2\kappa c}
\int_{\partial M}\Sigma_{IJ}\wedge(P_\gamma\mathfrak k)^{IJ}.
\]
La densidad interior dependiente de contorsión es
$\mathcal Q_{e,\gamma}(\mathfrak k)
-\tfrac12\mathfrak k^{IJ}\wedge\sigma_{IJ}$,
donde $\mathcal Q$ es homogénea de grado dos.
Su estacionariedad evaluada en
$\delta\mathfrak k=\mathfrak k_*$ da
$2\mathcal Q(\mathfrak k_*)=
\tfrac12\mathfrak k_*^{IJ}\wedge\sigma_{IJ}$.
Por tanto,
\begin{equation}
\mathcal L_{\rm spin}^{\rm ef}
=-\tfrac14\mathfrak k_*^{IJ}\wedge\sigma_{IJ}
=-\mathcal Q_{e,\gamma}(\mathfrak k_*).
\label{hmtcuatro:gea:spin}
\end{equation}
La linealidad conserva los cruces entre especies.
Se elimina la contorsión una vez; no se añaden simultáneamente
su interacción lineal eliminada y el cuártico efectivo.

En la especialización espinorial, el coeficiente recibido
$C_\gamma=3\kappa c\hbar^2\gamma^2/[16(1+\gamma^2)]$ da
\begin{equation}
\frac{C_\gamma}{\hbar}
=\frac{3\pi}{2}L^2\frac{\gamma^2}{1+\gamma^2}.
\label{hmtcuatro:gea:coeficiente-torsion}
\end{equation}
Basta sustituir $\kappa=8\pi G/c^4$ y $\hbar G=c^3L^2$,
sin cambiar la rama de $\gamma$.
En la Fock de la carta, la corriente total conserva la contracción
\[
\widehat Q_{{\rm tot},v}=\sum_{a=1}^4s_a
\{d\Gamma(\widetilde M_{a,v})^2-d\Gamma(\widetilde M_{a,v}^2)\},
\quad s=(-1,-1,-1,+1),
\]
\[
\widehat H_{\rm tor}
=-\sum_v\frac{N_vcC_\gamma}{\mathcal V_v}\widehat Q_{{\rm tot},v}.
\]
Las matrices $\widetilde M_{a,v}$ representan las componentes
de corriente. $N_v$ es el lapse, distinto de $N_\Omega$ y de
una ocupación. La identidad de orden normal retira la contracción
de una partícula, no los cruces entre especies.
La signatura permanece en $s_a$.

\subsubsection{Fuente métrica y Legendre de la misma acción}
\label{hmtcuatro:gea:legendre}

La eliminación estacionaria preserva la variación de las demás
variables: la regla de la cadena contiene un término en
$\delta\mathfrak k_*$ multiplicado por la ecuación de conexión,
que es cero. La acción reducida es
\[
S_{\rm red}=\frac1{2\kappa c}\int(R[g]-2\Lambda)\operatorname{vol}_e
+S_m^{\rm LC}+S_{\rm spin}^{\rm ef}+S_\partial.
\]
Definimos la fuente de acción por
$\delta_g(S_m^{\rm LC}+S_{\rm spin}^{\rm ef})
=-\tfrac12\int\mathcal T^S_{\mu\nu}\delta g^{\mu\nu}
\operatorname{vol}_e$.
La variación gravitatoria interior produce
\begin{equation}
G_{\mu\nu}+\Lambda g_{\mu\nu}
=\kappa c\mathcal T^S_{\mu\nu}
=\kappa(T^{\rm LC}_{\mu\nu}+U^{\rm spin}_{\mu\nu}),
\quad T^{\rm LC}+U^{\rm spin}=c\mathcal T^S.
\label{hmtcuatro:gea:metrica}
\end{equation}
Se incluye la variación de volumen del término constante del
potencial. El factor $c$ corresponde a la cotetrada en la recta
de longitud; no se mezclan las fuentes de acción y físicas.

Con $a=(2\kappa c)^{-1}$ y
$\mathcal B_{IJ}=(P_\gamma\Sigma)_{IJ}$, el potencial geométrico
es $\theta_{\rm geom}=a\mathcal B_{IJ}\wedge\delta\omega^{IJ}$.
Descomponer $\omega=\omega_0dt+\underline\omega$ y
$F_{0a}=\dot\omega_a-D_a\omega_0$ identifica el término cinético;
integrar por partes conserva Gauss y borde.
Con $e_0=Nn+N^ae_a$ y la Legendre de la misma materia,
\begin{equation}
S_{\rm tot}=\int dt\,[\Theta_{\rm red}(\dot z)
-\mathcal C[N]-\mathcal V[N^a]-\mathcal G_L[\omega_0]
-\mathcal G_{\rm int}[A_0]]+S_\partial.
\label{hmtcuatro:gea:canonica}
\end{equation}
Lapse, shift y calibre mantienen sus direcciones degeneradas.
Las restricciones de segunda clase se reducen antes de usar
$\Theta_{\rm red}$ o se conservan expresamente.
Los espinores de primer orden no reciben una inversa ficticia
de velocidades.

Para contorsión y su momento, las restricciones son
$\pi_{\mathfrak k}=0$ y
$\partial\mathcal Q/\partial\mathfrak k-\sigma/2=0$.
Su matriz tiene bloques
$\left(\begin{smallmatrix}0&-M\\M^{\mathsf T}&B\end{smallmatrix}\right)$,
con $M$ invertible por la reconstrucción anterior; por tanto es
invertible. El corchete de Dirac efectúa esa eliminación y conserva
separadamente las restricciones fermiónicas graduadas.
En una dirección bosónica regular,
\[
H_\Phi=N\left(\frac{p^2}{2w}+wV+H_{\rm grad}\right),\quad
p=\frac wN D_t\phi,\quad
L_\Phi=\frac{w}{2N}|D_t\phi|^2-NwV-NH_{\rm grad}.
\]
El volumen $w=\sqrt{\det q}$ no se congela.
Para $K_{ij}=(\dot q_{ij}-\mathcal L_{N^a}q_{ij})/(2N)$,
\[
\Pi^{ij}=\frac{\sqrt q}{2\kappa c}(K^{ij}-q^{ij}K),\qquad
K_{ij}=\frac{2\kappa c}{\sqrt q}
(\Pi_{ij}-\tfrac12q_{ij}\Pi).
\]
La traza y la parte sin traza prueban la inversa.
Si el momento total contiene una contribución material de marco,
se usa $\Pi=p-p_m$, conservando $p_m$.
Así se recupera la misma acción y su fuente. La Legendre no
demuestra por sí sola que una cuantización arbitraria de
\eqref{hmtcuatro:gea:canonica} coincida con
\eqref{hmtcuatro:gea:total}.

\subsection{Propagación de restricciones con la fuente total}
\label{hmtcuatro:gea:restricciones}

Sea $E_{\mu\nu}=G_{\mu\nu}+\Lambda g_{\mu\nu}
-\kappa c\mathcal T^S_{\mu\nu}$.
Sobre las ecuaciones materiales, la invariancia por difeomorfismos
del funcional reducido implica
$\nabla^\mu\mathcal T^S_{\mu\nu}=0$:
una variación interior generada por $\xi$ deja la integral de
$\mathcal T^S_{\mu\nu}\nabla^\mu\xi^\nu$, mientras las ecuaciones
materiales anulan los demás términos. Integrar por partes y variar
$\xi$ prueba la divergencia nula, antes de imponer $E=0$.
Bianchi y los acoplamientos fijos dan $\nabla_\mu E^{\mu\nu}=0$.

\begin{proposition}[Conservación de los datos de restricción]
\label{hmtcuatro:gea:propagacion}
Una solución regular de las ecuaciones materiales y evolutivas
métricas conserva las restricciones satisfechas inicialmente
durante su intervalo de existencia regular.
\end{proposition}
\begin{proof}
En coordenadas gaussianas
$g=-d\tau^2+q_{ij}(\tau,x)dx^idx^j$ se imponen $E_{ij}=0$ y
se escribe
$C=E_{00}$, $C_i=E_{0i}$,
$K_{ij}=\dot q_{ij}/2$ y $\mathcal K=q^{ij}K_{ij}$.
Se tiene $E^{00}=C$, $E^{0i}=-C^i$, $E^{ij}=0$,
$\Gamma^i{}_{0j}=K^i{}_j$ y
$\Gamma^\mu{}_{\mu0}=\mathcal K$.
La componente temporal de la divergencia y la espacial dan
\[
\partial_\tau C-D_iC^i+\mathcal KC=0,\qquad
\partial_\tau C^i+\mathcal KC^i+2K^i{}_jC^j=0.
\]
Al bajar el índice, $\dot q_{ij}=2K_{ij}$ cancela el último
término espacial. Quedan exactamente
\begin{equation}
\partial_\tau C-D_iC^i+\mathcal KC=0,\qquad
\partial_\tau C_i+\mathcal KC_i=0.
\label{hmtcuatro:gea:propagacion-sistema}
\end{equation}
La segunda ecuación homogénea mantiene $C_i=0$ y la primera
queda $\dot C+\mathcal KC=0$, conservando $C=0$.
$q$ y $\mathcal K$ son los de la solución acoplada, no un fondo prescrito.
\end{proof}
La carta gaussiana prueba localmente una afirmación tensorial,
sin congelar $q^{-1}$ ni sustituir la fuente.
No se afirma existencia global sin singularidades para cada dato.
La propagación variacional tampoco es por sí sola una identidad
de conmutadores cuánticos para deformaciones normales arbitrarias.

\subsection{Frontera y operador autoinercial}
\label{hmtcuatro:gea:autoinercia}

Sea $b:\mathcal X_\infty^{\rm enr}\to B_{\rm ef}$,
$B_{\rm ef}=\operatorname{im}(b)$.
Una respuesta $I_v$ factoriza unívocamente como
$I_v=\overline I_v\circ b$ si y sólo si es constante en las
fibras de $b$. La necesidad es inmediata; para la suficiencia,
$\overline I_v(\beta)=I_v(x)$ con $b(x)=\beta$ es independiente
del representante. La imagen efectiva da unicidad.
La dependencia machiana fuerte añade estados con igual restricción
local y respuestas distintas.

La bisagra del mismo estado produce
$r_{\rm av}=\pi/\varphi^2$ y $r_{\rm av}^J=\varphi/\pi^2$.
En el sector de amplitud común,
$\mathcal L_{\rm ar}=\pi\mu_{\rm ar}\mathcal M(x)$ y
$\mathcal L_{\rm vol}=\mu_{\rm vol}\mathcal M(x)$,
$\mathcal M(x)>0$, con
$\mu_{\rm ar}/\mu_{\rm vol}=\varphi^{-2}$.
Dividir por $\mathcal L_{\rm vol}$ da lectores del mismo tipo
$A_\partial=r_{\rm av}$ y $V_\partial=1$.
El cociente cancela la amplitud, no la historia conservada en el archivo.

Para una orientación $a$,
$q_a(x)=h_a(x)/(108t_0(x))=\hbar_a(x)/t_P(x)$
desciende a $Q:B_{\rm ef}\to\mathcal L_E^+$
si y sólo si $b(x)=b(x')$ implica
$h_a(x)/t_0(x)=h_a(x')/t_0(x')$.
Es el criterio de factorización aplicado al reloj.
Lo satisfacen las cartas fijas y las variables cuyo cociente
desciende. En la carta radial $Q=\hbar_ac/L$.
Con energías de esa misma recta, $E_1E_2/Q^2$ es adimensional
e independiente de base.

Para proyecciones ortogonales complementarias,
\begin{align}
w_A&=\frac{A_\partial r_{\rm av}}
{A_\partial r_{\rm av}+V_\partial r_{\rm av}^J}
=\frac{\pi^4}{\pi^4+\varphi^5},\qquad w_V=1-w_A,
\label{hmtcuatro:gea:pesos}\\
W_{\rm av}&=w_AP_A+w_VP_V,\qquad
\overline I_v=\frac{E_1E_2}{Q^2}W_{\rm av}.
\label{hmtcuatro:gea:respuesta}
\end{align}
Multiplicar numerador y denominador de
$r_{\rm av}^2/(r_{\rm av}^2+r_{\rm av}^J)$ por
$\varphi^4\pi^2$ prueba la evaluación.
Se conserva la diferencia entre lector areal y ponderación:
la marca regional se genera una vez.
$0<w_A,w_V<1$ hace positivo e invertible a $W_{\rm av}$ y
\[
\langle z,\overline I_vz\rangle
=\frac{E_1E_2}{Q^2}
(w_A\|P_Az\|^2+w_V\|P_Vz\|^2)>0\quad(z\ne0).
\]
Intercambiar completamente
$(A_\partial,r_{\rm av},P_A)$ y
$(V_\partial,r_{\rm av}^J,P_V)$ permuta los sumandos.
Un transporte unitario conjuga la respuesta y conserva su espectro;
implementar el intercambio dentro de un espacio exige iguales
dimensiones sectoriales.
Una reescala común de $A_\partial,V_\partial$ no cambia los pesos
y no prueba dependencia global.
Con energías y pesos fijos, $Q\mapsto qQ$ cambia la respuesta
por $q^{-2}$; el testigo fuerte conserva además la misma
restricción local.

\begin{theorem}[Autoinercia y radio del mismo carácter]
\label{hmtcuatro:gea:autoinercial-radial}
Con el mismo reloj y orientación,
\begin{equation}
\overline I_v=\frac{G_aE_1E_2}{\hbar_ac^5}W_{\rm av}
=\frac{G_am_1m_2}{\hbar_ac}W_{\rm av}.
\label{hmtcuatro:gea:mach-gravedad}
\end{equation}
Para una sonda $y_p>0$, $E_p=Qy_p$ y
$\bar\lambda_p=L/y_p$, la extensión espectral satisface
\begin{equation}
\widehat I_{p,Y}=y_pY\otimes W_{\rm av}
=\frac{E_p}{Q^2}H_X\otimes W_{\rm av}
=\frac{R_X}{\bar\lambda_p}\otimes W_{\rm av}
=\frac{G_am_p}{\hbar_ac}M_X\otimes W_{\rm av}.
\label{hmtcuatro:gea:identidad-autoinercial}
\end{equation}
\end{theorem}
\begin{proof}
$Q^2=\hbar_a^2c^2/L^2$ y $\hbar_aG_a=c^3L^2$ dan
$Q^{-2}=G_a/(\hbar_ac^5)$; después se usa $E_i=m_ic^2$.
Para $Y=\sum_jy_jP_j$, aplicar el lector a cada energía $Qy_j$
produce $\sum_jy_py_jP_j\otimes W_{\rm av}$, independientemente
de base propia. El cálculo espectral extiende el lector lineal
al carácter autoadjunto positivo general.
La descomposición por $P_A,P_V$ da los operadores
$y_pw_AY$ y $y_pw_VY$, con sus multiplicidades:
son autoadjuntos positivos en el dominio de $Y\otimes I$.
Sustituir los lectores radiales prueba las otras expresiones.
\end{proof}

Si una realización adicional identifica efectivamente su operador
completo con $H_X=QY$, conserva esa identificación y sus dominios.
Para $H_X=\sum_rH_r$ en un núcleo común,
$\widehat I_{p,Y}=(E_p/Q^2)\sum_rH_r\otimes W_{\rm av}$.
No se requiere positividad de cada sumando; sí la del carácter
del operador identificado. No se suprimen términos ni se cambia
el origen de energía para imponerla.
El teorema no asigna automáticamente este lector a $\mathbb H$.

En los bloques de \eqref{hmtcuatro:gea:schur-prueba}, el inverso
interior de $y_pY\otimes W_{\rm av}$ es
$y_p^{-1}C^{-1}\otimes W_{\rm av}^{-1}$.
La multiplicación da
\begin{equation}
\operatorname{Schur}(\widehat I_{p,Y})
=y_pY_{\rm ef}\otimes W_{\rm av}
=\frac{R_{X,\rm ef}}{\bar\lambda_p}\otimes W_{\rm av}.
\label{hmtcuatro:gea:schur-autoinercia}
\end{equation}
No se borran los bloques cruzados.
Para $a_s=y_pw_s>0$, $s=A,V$, la misma cuenta da
$\mathcal K_{a_sY}(a_sz)=a_s\mathcal K_Y(z)$
en el dominio resolvente interior, conservando la memoria completa.
Si $Y'K=KY$ y $W'_{\rm av}J=JW_{\rm av}$, con
$y_p$ y sección intensiva fijos,
\begin{equation}
\widehat I'_{p,Y'}(K\otimes J)
=(K\otimes J)\widehat I_{p,Y}.
\label{hmtcuatro:gea:naturalidad}
\end{equation}
Se prueba por sustitución.
Los cortes espectrales de la proposición
\ref{hmtcuatro:gea:limite} convergen en norma de grafo;
los dos pesos fijos transmiten la convergencia a la respuesta.
La naturalidad es la de estos entrelazadores efectivos, no la
de una familia distinta de mapas espaciales identificada sólo por nombre.

\subsection{Autoinercia del estado y aceleración de su proyección}
\label{hmtcuatro:gea:proyeccion}

En una realización diferenciable, el estado conjunto
$\Gamma=(x,\mathcal O,m)$ tiene conexión $\nabla^{\rm tot}$
en $\mathcal N$. Una proyección suave
$\pi_S:\mathcal N\to\mathcal N_S$, con conexión $\nabla^S$,
conserva los transportes y lectores que realiza.
El marco observado es
$\mathfrak F_{\mathcal O}(\Gamma)
=\operatorname{Res}_{\mathcal O}(\Gamma,\nabla^{\rm tot})$.
Para una curva dos veces diferenciable y
$\gamma_S=\pi_S\circ\Gamma$, la autoinercia significa
$\nabla^{\rm tot}_{\dot\Gamma}\dot\Gamma=0$.
El defecto es
\[
\mathcal K_{\pi_S}(\dot\Gamma,\dot\Gamma)
=\nabla^S_{\dot\gamma_S}\dot\gamma_S
-d\pi_S(\nabla^{\rm tot}_{\dot\Gamma}\dot\Gamma).
\]
\begin{proposition}[Aceleración publicada]
\label{hmtcuatro:gea:aceleracion}
Para una trayectoria autoinercial,
\begin{equation}
a_S=\mathcal K_{\pi_S}(\dot\Gamma,\dot\Gamma),\qquad
(\mathcal K_{\pi_S})^a{}_{BC}
=\partial_B\partial_C\pi_S^a
+(\Gamma^S)^a{}_{bc}\partial_B\pi_S^b\partial_C\pi_S^c
-\partial_D\pi_S^a(\Gamma^{\rm tot})^D{}_{BC}.
\label{hmtcuatro:gea:defecto}
\end{equation}
\end{proposition}
\begin{proof}
La regla de la cadena da
$\ddot\gamma_S^a=\partial_B\pi_S^a\ddot\Gamma^B
+\partial_B\partial_C\pi_S^a\dot\Gamma^B\dot\Gamma^C$.
Sumar la conexión del subsistema y restar la imagen de la
aceleración total cancela los términos en $\ddot\Gamma$.
Los restantes son el tensor mostrado contraído con las velocidades.
La autoinercia anula la aceleración total y deja la igualdad.
\end{proof}
La respuesta geométrica factoriza por $b$ cuando conexión,
proyección y dato cinemático descienden a la frontera, o éste
se fija en la comparación. Si cambia la velocidad, se conserva
como argumento. Así, el conjunto autoinercial puede publicar
subsistemas acelerados sin un marco absoluto exterior.

Autoinercia no anula por sí sola curvatura ni holonomía.
Un dato transportado por un lazo retorna si y sólo si está en
el espacio fijo de su holonomía. El retorno del estado completo
requiere además una realización inyectiva en el sector comparado:
retorno visible y retorno de historia siguen siendo distintos.

\subsection{Resultado y controles de alcance}
\label{hmtcuatro:gea:conclusion}

El mismo carácter fija masa, energía y radio; su proporcionalidad
conserva memoria estática, espectral y normalización dual.
Los transportes geométrico e interno entran en el diferencial
de una sola energía. La acción conjunta produce las fuentes
geométrica y torsional y conserva las restricciones sobre su
solución regular. El operador autoinercial es el radio del
carácter dividido por la longitud conjugada de la sonda,
mientras la aceleración es el defecto de proyección del estado.

Los dos últimos objetos tienen tipos distintos:
$\widehat I_{p,Y}$ no se identifica con $\mathcal K_{\pi_S}$
ni con el tensor de Einstein. La incorporación probada tampoco
se sustituye por la exigencia diferente de anular una curvatura
cuántica total para deformaciones normales arbitrarias.
Ese enunciado posee sus generadores y dominios; no se afirma aquí
ni se declara ausente del corpus por no ser esta conclusión.

Los falsadores son concretos. Cambiar sólo $G$ o $\hbar$ con
$L,Y$ fijos rompe \eqref{hmtcuatro:gea:radial}; sustituir
$L$ por $\ell_0$ sin $54/\pi$ altera el reloj; borrar $F$
modifica la reducción; reescalar juntos los lectores areal y
volumétrico no prueba separación machiana; un mapa que no
entrelace $Y$ no satisface \eqref{hmtcuatro:gea:naturalidad}.
Los ensayos racionales son controles finitos reproducibles,
no valores objetivo del generador ni sustitutos de las pruebas
de dominio y límite.

% Fragmento integrable; no documento autónomo y no orden de compilación.
% Fecha: 2026-09-30. No modifica fuentes publicadas.
% Procedencia: reunión de RETROACCION_PCH_Y_RESTRICCIONES_VARIACIONALES.md,
% PRECUANTIZACION_PCH_Y_FIDELIDAD.md y CIERRE_CANONICO_TOTAL_Y_REPRESENTACION.md.
% Prerrequisitos de inserción: núcleo común APP--TRIT--TPK; VIII/30d,
% VIII/31 y VIII/33; lectores de acción, gamma y acoplamientos.
% Los nuevos labels usan exclusivamente el prefijo hmtcuatro:.
\section{Acción conjunta, retroacción y cierre canónico de las restricciones}
\label{hmtcuatro:accion-retroaccion-restricciones}

La secuencia causal de este desarrollo es
\[
 \mathrm{APP}\longrightarrow\mathrm{TRIT}\longrightarrow\mathrm{TPK}
 \longrightarrow\text{estado enriquecido}
 \longrightarrow\text{estructura discreta conjunta del continuo}.
\]
Se conservan las hojas, la orientación, los residuos, el acarreo, las rutas,
la frontera y la memoria de esa construcción. La escala de acción y los
acoplamientos son salidas HMT recibidas antes de emplear el lenguaje
variacional y simpléctico. No se seleccionan por valores objetivo ni por
la conclusión de curvatura que se va a demostrar. El retorno de fase
tampoco reinicia el estado enriquecido.

La realización geométrica previamente construida proporciona una
cotetrada no degenerada \(e\), su conexión métrica y el levantamiento
espinorial del transporte. Aquí esa geometría se varía junto con los
campos materiales. Reuniremos en una sola demostración la acción común,
su reducción torsional, la transformación de Legendre, la primera clase,
la propagación de las restricciones y la planitud característica de la
conexión canónica que resulta de la misma acción. Esta última realización
no se identifica por su nombre con otro Hamiltoniano interactuante.

\subsection{Estructura de partida y acción con fuente completa}
\label{hmtcuatro:datos-accion-total}

Sea \(M\) una región lisa orientada de dimensión cuatro con estructura
de espín, firma \((-+++)\) y una foliación espacial regular. La cotetrada
toma valores en la recta de longitud de la realización recibida.
Escribimos
\[
 \Sigma^{IJ}=e^I\wedge e^J,\qquad
 F_\omega=\mathrm d\omega+\omega\wedge\omega,\qquad
 \mathsf P_\gamma=\star+\gamma^{-1}I,\qquad \star^2=-I.
\]
La estrella de esta fórmula actúa sobre bivectores internos; no es la
dualidad exterior de formas. Durante la variación permanecen fijas las
secciones recibidas \(c>0,\hbar>0,G>0,\gamma\in\mathbb R\setminus\{0\}\)
y \(\Lambda\). Con \(\kappa=8\pi G/c^4\), la acción total es
\begin{equation}
 \begin{split}
 S_{\rm tot}={}&\frac1{2\kappa c}
   \int_M\Sigma_{IJ}\wedge(\mathsf P_\gamma F_\omega)^{IJ}
 -\frac{\Lambda}{\kappa c}\int_M\operatorname{vol}_e\\
 &+S_{\rm YM}[e,A]+S_H[e,A,H]
   +S_W[e,\omega,A,\Psi]+S_Y[e,H,\Psi]+S_{\partial}.
 \end{split}
 \label{hmtcuatro:accion-total}
\end{equation}

\begin{definition}[Unificación variacional y canónica]
\label{hmtcuatro:def-unificacion}
La unificación variacional y canónica HMT es la composición de las realizaciones geométrica, cromática y electrodébil del mismo estado enriquecido en la acción común \eqref{hmtcuatro:accion-total}, cuyas variaciones producen fuentes conjuntas y cuya reducción conserva una misma álgebra de restricciones y su propagación. La componente electromagnética pertenece a la realización electrodébil. El teorema~\ref{hmtcuatro:teorema-canonico-total} establece esta composición en el dominio regular y con las condiciones de borde declaradas.
\end{definition}
Los acoplamientos electrodébiles $g,g'$ y el autoacoplamiento $\lambda_H$ se reciben de la construcción de \ref{amp-higgs-seccion} y del teorema~\ref{amp-higgs-inversa-teorema}; la variación presente mantiene fijas esas secciones. Su comparación entre secciones es una operación distinta, conservada en el desarrollo del horizonte de compatibilidad.


El símbolo \(H\) sin subíndice designa aquí el campo Higgs, no un
Hamiltoniano. Todas las contracciones espaciotemporales de la materia
utilizan la misma \(e\). El problema de borde es el recibido de la
variación Cartan--Holst: variaciones interiores, datos de borde fijos
compatibles o una completación que anule el mismo término fronterizo.

Para fijar el contenido de \eqref{hmtcuatro:accion-total}, las conexiones
internas actúan sobre el fibrado quiral
\[
 (S_L\otimes E_L)\oplus(S_R\otimes E_R).
\]
No se sustituye ese fibrado por un acoplamiento vectorial. En la
normalización entera \(y=6Y\), los multipletes recibidos tienen pesos
\(1,4,-2,-3,-6\), respectivamente para \(q_L,u_R,d_R,\ell_L,e_R\);
el doblete \(H\) tiene peso \(3\). El color actúa en el factor
\(\mathbb C^3\) de los quarks y trivialmente en los leptones. El factor
de familias y sus mapas \(Y_u,Y_d,Y_e\) se conservan. No se añade en
este bloque un compañero derecho a un multiplete que no lo tenga.

Con \(\widetilde H=i\sigma_2\overline H\), el mapa material es
\begin{equation}
 \begin{split}
 \mathsf Y_q(H)(u,d)&=
       \widetilde H\otimes Y_u u+H\otimes Y_d d,\\
 \mathsf Y_\ell(H)e&=H\otimes Y_e e .
 \end{split}
 \label{hmtcuatro:yukawa}
\end{equation}
Se entiende la identidad sobre el factor de color en la primera línea.
El término de acción es
\(-(\bar\Psi_L\mathsf Y(H)\Psi_R+\mathrm{h.c.})\), con las secciones
dimensionales recibidas. La equivariancia no es sólo nominal: para
\(U\in SU(2)\), \(i\sigma_2\overline U=U i\sigma_2\), de modo que
\(\widetilde H\mapsto U\widetilde H\). Los pesos de las tres columnas
son \(-3+4=1\), \(3-2=1\) y \(3-6=-3\). El color se entrelaza por la
identidad y los mapas de familia conmutan con la acción interna.
Por tanto, para cada transformación interna \(u\),
\begin{equation}
 \mathsf Y(uH)\rho_R(u)=\rho_L(u)\mathsf Y(H).
 \label{hmtcuatro:equivariancia-yukawa}
\end{equation}
Si la realización material precedente contiene además otros
multipletes, sus mapas se incorporan con su identidad de equivariancia
propia; esta prueba no inventa estados nuevos.

En una carta de unidades \(\hbar=c=1\), la densidad material que resume
estas acciones tiene la forma
\begin{equation}
 \begin{split}
 \mathcal L_m={}&-\frac14 F_{B,\mu\nu}F_B^{\mu\nu}
 -\frac14\sum_a F_{W,\mu\nu}^aF_W^{a,\mu\nu}
 -\frac14\sum_a F_{C,\mu\nu}^aF_C^{a,\mu\nu}\\
 &-(\nabla_\mu H)^\dagger\nabla^\mu H-V(H)
 +\mathcal L_W(\Psi_L,\nabla_L)+\mathcal L_W(\Psi_R,\nabla_R)\\
 &-\bigl(\bar\Psi_L\mathsf Y(H)\Psi_R+\mathrm{h.c.}\bigr).
 \end{split}
 \label{hmtcuatro:densidad-material}
\end{equation}
Las curvaturas incluyen sus acoplamientos; éstos y las normalizaciones
de componentes son los de sus propietarios. La escritura en esa carta
no elimina la sección de acción ni recalibra los coeficientes.
En particular, no se introduce aquí un valor numérico nuevo para el
acoplamiento fuerte. La conexión de espín de \(\nabla_L,\nabla_R\)
es la de \(e,\omega\); las conexiones \(B,W,C\) actúan en las
representaciones recién descritas. La constante de \(V\) se conserva,
pues su variación de volumen pertenece a la fuente gravitatoria.

La parte espinorial conserva las convenciones explícitas del desarrollo
precedente: \(\{g^I,g^J\}=2\eta^{IJ}I\),
\(\mathbb B_{IJ}=\frac14[g_I,g_J]\), \(A_D=-ig^0\),
\(\bar\psi=\psi^\dagger A_D\) y
\(D\psi=\mathrm d\psi+\frac12\omega^{IJ}\mathbb B_{IJ}\psi+A\psi\).
Su cinética simétrica es
\[
 \frac{\hbar}{2}
 [\bar\psi g^I D\psi-(D\bar\psi)g^I\psi]\wedge\eta_I,
 \qquad \eta_I=\iota_{E_I}\operatorname{vol}_e .
\]
La suma se toma sobre los campos realmente presentes, con sus
proyecciones quirales, no sobre una duplicación vectorial del catálogo.
Los términos Yukawa son algebraicos y no dependen de \(\omega\).
La corriente total queda fijada por
\begin{equation}
 \delta_\omega S_m=-\frac12\int_M\delta\omega^{IJ}\wedge\sigma_{IJ},
 \qquad
 \sigma_{JK}=\sum_s-\frac{\hbar}{2}
   \bar\psi_s\{g^I,\mathbb B_{JK}\}\psi_s\,\eta_I.
 \label{hmtcuatro:corriente-total}
\end{equation}

\subsection{Una demostración de compatibilidad, propagación y cierre}
\label{hmtcuatro:demostracion-unica}

\begin{theorem}[Compatibilidad canónica de la acción conjunta]
\label{hmtcuatro:teorema-canonico-total}
Considérese \eqref{hmtcuatro:accion-total} en el dominio no degenerado
anterior, con materia afín en la contorsión, y en una carta regular de
la transformación de Legendre restringida. Reténganse las restricciones
fermiónicas de primer orden y elimínense las de segunda clase mediante
su corchete de Dirac. Para parámetros de soporte interior, o con el
borde compatible sin cargos impropios, se cumplen conjuntamente:
la ecuación métrica tiene la fuente material total de esa misma acción;
sus generadores de difeomorfías y calibre son de primera clase; las
restricciones se propagan sobre toda solución acoplada regular; y el
potencial canónico total induce una conexión precuántica plana en las
direcciones características de la superficie regular de restricciones.

La afirmación es sobre esta acción y esta realización canónica.
No identifica automáticamente su conexión con un Hamiltoniano
interactuante distinto ni prueba su límite espacial--quiral.
\end{theorem}

\begin{proof}
\par\medskip\noindent\textbf{1. Eliminación torsional sin duplicación.}\quad
Escribamos \(\omega=\mathring\omega(e)+\mathfrak k\).
La expansión \(F_\omega=F_{\mathring\omega}
+D_{\mathring\omega}\mathfrak k+\mathfrak k\wedge\mathfrak k\)
y \(D_{\mathring\omega}\Sigma=0\) separan la contribución de borde
\[
 S_{\partial,\mathfrak k}
 =\frac1{2\kappa c}\int_{\partial M}
                  \Sigma_{IJ}\wedge(\mathsf P_\gamma\mathfrak k)^{IJ}
\]
de la densidad algebraica
\begin{equation}
 \mathcal Q_{e,\gamma}(\mathfrak k)
       -\frac12\mathfrak k^{IJ}\wedge\sigma_{IJ},\qquad
 \mathcal Q_{e,\gamma}(\mathfrak k)=\frac1{2\kappa c}
       \Sigma_{IJ}\wedge
       \mathsf P_\gamma(\mathfrak k\wedge\mathfrak k)^{IJ}.
 \label{hmtcuatro:cuadratica-contorsion}
\end{equation}
La variación original da
\(D_\omega(\mathsf P_\gamma\Sigma)=\kappa c\,\sigma\).
La inversa de Holst y la inversa torsión--contorsión del desarrollo
precedente son invertibles para el coframe y \(\gamma\) declarados.
Así determinan una única \(\mathfrak k_*[e,\sigma]\), lineal en
\(\sigma\). La variación radial de
\eqref{hmtcuatro:cuadratica-contorsion} en su estacionario da
\(2\mathcal Q(\mathfrak k_*)=\frac12\mathfrak k_*^{IJ}\wedge\sigma_{IJ}\).
En consecuencia,
\begin{equation}
 S_{\rm red}=\frac1{2\kappa c}\int_M(R-2\Lambda)\operatorname{vol}_e
       +S_m^{\rm LC}
       -\frac14\int_M\mathfrak k_*^{IJ}\wedge\sigma_{IJ}
       +S_{\partial,\rm red}.
 \label{hmtcuatro:accion-reducida}
\end{equation}
La primera identidad de Bianchi anula el término Holst de
Levi--Civita en esta escritura. La regla de la cadena conserva las
ecuaciones de \(e,A,H,\Psi\), pues el coeficiente de
\(\delta\mathfrak k_*\) se anula por estacionariedad.
Como \(\sigma=\sum_s\sigma_s\), la contribución efectiva incluye
\(\mathfrak k_*[\sigma_s]\wedge\sigma_t\) para \(s\ne t\).
No se reemplaza por cuadrados de especies aisladas ni se conserva
simultáneamente la misma \(\mathfrak k\) como variable independiente.

\par\medskip\noindent\textbf{2. Fuente completa y retroacción.}\quad
Definamos el tensor de acción por la variación de todos los términos
materiales de \eqref{hmtcuatro:accion-reducida}:
\begin{equation}
 \delta_g S_{m,\rm red}=-\frac12\int_M
       \mathcal T^S_{\mu\nu}\,\delta g^{\mu\nu}\operatorname{vol}_e,
 \qquad
 E_{\mu\nu}:=G_{\mu\nu}+\Lambda g_{\mu\nu}
                   -\kappa c\,\mathcal T^S_{\mu\nu}.
 \label{hmtcuatro:fuente-completa}
\end{equation}
La ecuación métrica es \(E_{\mu\nu}=0\).
El factor \(\kappa c\) corresponde a la cotetrada de longitud;
si \(T^{\rm phys}=c\mathcal T^S\), se escribe \(\kappa T^{\rm phys}\).
Se varían también los campos materiales y las conexiones internas.
Por ello los campos determinan la fuente y la geometría que resuelve
la ecuación determina a su vez sus conexiones, volumen y evolución.
Ésta es retroacción, no evolución sobre un fondo fijado.
La variación incluye el volumen del potencial Higgs y la dependencia
métrica de la corriente torsional, además de sus términos cruzados.

\par\medskip\noindent\textbf{3. Legendre restringida y reducción canónica.}\quad
El potencial canónico total se obtiene de
\(\delta L=\mathcal E\delta z+\mathrm d\theta(\delta z)\),
integrando \(\theta\) sobre la hoja espacial y efectuando la reducción
de segunda clase. Su parte geométrica antes de esa reducción es
\[
 \theta_{\rm geom}(\delta)=
   \frac1{2\kappa c}(\mathsf P_\gamma\Sigma)_{IJ}
                       \wedge\delta\omega^{IJ}.
\]
La descomposición \(F_{0a}=\dot\omega_a-D_a\omega_0\), junto con la
misma descomposición de las acciones materiales, produce
\begin{equation}
 S_{\rm can}=\int d\tau\,
 \bigl[\Theta_{\rm red}(\dot z)-\mathcal C[N]-\mathcal D[\xi]
       -\mathcal G_{\rm L}[\lambda]-\mathcal G_{\rm int}[\chi]\bigr]
       +S_{\partial,\rm red}.
 \label{hmtcuatro:legendre-total}
\end{equation}
No se invierten artificialmente las velocidades ausentes de lapse,
shift o conexiones temporales.

La contorsión algebraica tiene restricciones
\(\pi_{\mathfrak k}=0\) y
\(\partial\mathcal Q/\partial\mathfrak k-\sigma/2=0\).
Su matriz de corchetes tiene el bloque
\[
 \begin{pmatrix}0&-M\\ M^T&B\end{pmatrix},
 \qquad M=\frac{\partial^2\mathcal Q}{\partial\mathfrak k^2}.
\]
La inversa torsional hace invertible \(M\), y por tanto esta matriz.
Son restricciones de segunda clase. Se usa
\begin{equation}
 \{F,G\}_D=\{F,G\}
       -\{F,\chi_a\}(C^{-1})^{ab}\{\chi_b,G\}.
 \label{hmtcuatro:corchete-dirac}
\end{equation}
En las variables que no dependen de la pareja contorsión--momento,
el bloque inferior derecho de \(C^{-1}\) es cero. No aparece por esa
eliminación un corchete suplementario entre ellas. Se retienen,
separadamente, las restricciones espinoriales de primer orden; en la
fase graduada se emplea el corchete graduado y los generadores pares.
\(\Theta_{\rm red}\) designa el resultado completo, no sólo su parte
geométrica.

La inversión de Legendre en direcciones regulares recupera la misma
acción. Por ejemplo, \(N[p^2/(2w)+wV+H_{\rm grad}]\), con
\(w=\sqrt{\det q}\), da \(p=(w/N)D_\tau\phi\), y por sustitución
\[
 L_H=\frac{w}{2N}|D_\tau\phi|^2-NwV-NH_{\rm grad}.
\]
\(D_\tau\) conserva shift y conexión temporal; \(w\) no se congela.
Para \(K_{ij}=(\dot q_{ij}-\mathcal L_\xi q_{ij})/(2N)\), la
contribución gravitatoria al momento y su inversa son
\begin{equation}
 \Pi^{ij}=\frac{\sqrt q}{2\kappa c}(K^{ij}-q^{ij}K),\qquad
 K_{ij}=\frac{2\kappa c}{\sqrt q}
                  (\Pi_{ij}-\tfrac12q_{ij}\Pi).
 \label{hmtcuatro:legendre-geometrica}
\end{equation}
De ellas resulta
\[
 \mathcal C_{\rm grav}
 =\frac{2\kappa c}{\sqrt q}
       (\Pi^{ij}\Pi_{ij}-\tfrac12\Pi^2)
  -\frac{\sqrt q}{2\kappa c}(R^{(3)}-2\Lambda).
\]
Si la carta de espinores aporta un desplazamiento material del momento
total \(p\), se usa \(\Pi=p-p_m\); no se descarta \(p_m\).
El carácter indefinido de la forma de DeWitt no impide esta inversa algebraica.
No se exige una inversa de velocidades a la cinética fermiónica de
primer orden.

\par\medskip\noindent\textbf{4. Identidad Noether--Legendre y primera clase.}\quad
La coincidencia de los generadores con las restricciones de la misma
acción puede verse antes de imponer sus ecuaciones.
Pónganse \(a=(2\kappa c)^{-1}\), \(B=\mathsf P_\gamma\Sigma\),
\(u^I=\iota_\zeta e^I\) y \(\lambda^{IJ}=\iota_\zeta\omega^{IJ}\).
De \(\mathcal L_\zeta\omega=\iota_\zeta F+D_\omega\lambda\) se obtiene
\begin{equation}
 \begin{split}
 j_\zeta^{\rm geom}
 ={}&\mathrm d(a B_{IJ}\lambda^{IJ})
   -a(D_\omega B_{IJ})\lambda^{IJ}\\
 &-2a u^Ie^J\wedge(\mathsf P_\gamma F)_{IJ}
       +2a\Lambda u^I\eta_I .
 \end{split}
 \label{hmtcuatro:noether-legendre}
\end{equation}
La parte material añade exactamente sus contribuciones a las
ecuaciones de \(e,\omega,A\), sus generadores verticales y su borde.
Descomponer la misma expresión en una hoja espacial da
\eqref{hmtcuatro:legendre-total}. La reducción estacionaria y la
reducción de Dirac transportan esas transformaciones, en vez de
sustituirlas por momentos de una acción externa.

Con la convención canónica usual para las restricciones, los corchetes
totales son, módulo los generadores Gauss conservados,
\begin{equation}
 \begin{aligned}
 \{\mathcal D[\xi],\mathcal D[\eta]\}_D
    &\simeq\mathcal D[[\xi,\eta]],\\
 \{\mathcal D[\xi],\mathcal C[N]\}_D
    &\simeq\mathcal C[\mathcal L_\xi N],\\
 \{\mathcal C[N],\mathcal C[M]\}_D
    &\simeq\mathcal D[
       q^{ij}(N\partial_jM-M\partial_jN)\partial_i].
 \end{aligned}
 \label{hmtcuatro:algebra-restricciones}
\end{equation}
En efecto, las transformaciones tangenciales tienen el corchete de
Lie, y transportan el lapse como escalar. Para las normales,
diferenciar su ortogonalidad con los vectores tangentes y
\(g(n,n)=-1\) determina la variación tangencial de \(n\);
con \(K_{ij}\) de \eqref{hmtcuatro:legendre-geometrica} se obtiene
\(\delta_N n=\operatorname{grad}_qN\). El conmutador de deformaciones,
convertido al corchete canónico con estas convenciones, da la última
línea de \eqref{hmtcuatro:algebra-restricciones}. El uso de transporte
covariantizado conserva además las rotaciones de marco y de calibre.
Las transformaciones son tangentes a la superficie de restricciones
de la acción invariante, de donde resulta la primera clase.
Las funciones \(q^{ij}\) dependen del estado y no se congelan.
Con soporte interior no se añade un cargo central de borde; un
generador impropio con cargo no nulo no se declara restricción cero.

\par\medskip\noindent\textbf{5. Propagación en la geometría que responde a la materia.}\quad
La invariancia por difeomorfías de \(S_{m,\rm red}\), evaluada sobre
sus ecuaciones materiales, da por integración por partes
\(\mathring\nabla^\mu\mathcal T^S_{\mu\nu}=0\). La identidad usa el
tensor completo de \eqref{hmtcuatro:fuente-completa}; no postula
conservación separada de cada contribución.
Bianchi implica entonces
\(\mathring\nabla_\mu E^{\mu\nu}=0\), aun antes de imponer
todas las ecuaciones métricas.

En una carta gaussiana local,
\(ds^2=-d\tau^2+q_{ij}dx^idx^j\), impongamos las ecuaciones
evolutivas \(E_{ij}=0\) y pongamos
\(C=E_{00}\), \(C_i=E_{0i}\),
\(K_{ij}=\dot q_{ij}/2\), \(K=q^{ij}K_{ij}\).
Aquí \(K\) es sólo la traza de la segunda forma fundamental.
Las dos componentes de la identidad divergente son
\begin{equation}
 \partial_\tau C-D_iC^i+KC=0,\qquad
 \partial_\tau C_i+KC_i=0.
 \label{hmtcuatro:propagacion}
\end{equation}
Para comprobar los signos se usan
\(E^{00}=C\), \(E^{0i}=-C^i\), \(E^{ij}=0\),
\(\Gamma^i{}_{0j}=K^i{}_j\) y
\(\Gamma^\mu{}_{\mu0}=K\). Al bajar el índice de la ecuación
espacial, \(\dot q_{ij}=2K_{ij}\) cancela los dos términos de conexión.
Si inicialmente \(C_i=0\), su ecuación homogénea mantiene \(C_i=0\);
la primera queda \(\dot C+KC=0\) y conserva \(C=0\).
\(q\) y \(K\) son los de la solución acoplada. La elección gaussiana
prueba la afirmación local; su forma tensorial permite cambiar lapse
y shift. Análogamente, la identidad gauge de Noether relaciona la
divergencia covariante de la ecuación de calibre con las ecuaciones
materiales y propaga Gauss cuando se cumplen las ecuaciones espaciales.
No se infiere existencia global ni ausencia de singularidades de
toda solución.

\par\medskip\noindent\textbf{6. La misma acción induce el cierre característico.}\quad
Sea \(\mathcal P_{\rm red}\) la carta canónica regular obtenida y
\(\mathcal Z=\{C_A=0\}\) su superficie regular de restricciones
totales, incluida Gauss. Usamos un símbolo distinto de \(\Sigma^{IJ}\)
para esa superficie. Se conserva el potencial \(\Theta=\Theta_{\rm red}\)
de \eqref{hmtcuatro:legendre-total}, incluidos los términos materiales
y mixtos. Fijamos
\begin{equation}
 \Omega=-\mathrm d_{\mathcal P}\Theta,\qquad
 \iota_{X_f}\Omega=\mathrm d_{\mathcal P}f,\qquad
 \{f,g\}_D=\Omega(X_f,X_g).
 \label{hmtcuatro:convenciones-simpecticas}
\end{equation}
En una pareja canónica, \(\Theta=p\,\mathrm dq\) da
\(X_f=f_p\partial_q-f_q\partial_p\), \(\{q,p\}_D=1\),
\(X_f(g)=-\{f,g\}_D\) y
\([X_f,X_g]=-X_{\{f,g\}_D}\). En la fase graduada esta escritura
se aplica a las funciones y restricciones pares pertinentes.

En la línea de fase trivializada sobre esta carta, construimos
\begin{equation}
 \nabla_X=X-\frac{i}{\hbar}\Theta(X)
       =\delta_X+\frac{i}{\hbar}H_X^{\rm can},
 \qquad H_X^{\rm can}=-\Theta(X).
 \label{hmtcuatro:conexion-canonica}
\end{equation}
No se escoge una conjugación arbitraria para hacerla plana:
su coeficiente proviene del potencial de la acción.
Para cualesquiera campos \(X,Y\) de la carta, el cálculo da
\begin{equation}
 \begin{aligned}
 \mathcal F^{\rm can}_{X,Y}
 &=XH_Y^{\rm can}-YH_X^{\rm can}-H_{[X,Y]}^{\rm can}
       +\frac{i}{\hbar}[H_X^{\rm can},H_Y^{\rm can}]\\
 &=-X\Theta(Y)+Y\Theta(X)+\Theta([X,Y])\\
 &=-\mathrm d_{\mathcal P}\Theta(X,Y)=\Omega(X,Y).
 \end{aligned}
 \label{hmtcuatro:curvatura-calculada}
\end{equation}
Los \(H_X^{\rm can}\) son multiplicadores de la línea, por lo que
su conmutador es cero; las derivadas de los coeficientes y de los
campos no se suprimen.

La primera clase ya probada significa
\(\{C_A,C_B\}_D=f_{AB}{}^D C_D\).
Para \(V\) tangente a \(\mathcal Z\),
\(\Omega(X_{C_A},V)=\mathrm d_{\mathcal P}C_A(V)=0\).
Los campos de las restricciones son, por tanto, característicos.
Bajo la regularidad declarada generan la distribución característica
de la superficie coisótropa. La tensorialidad de \(\Omega\) extiende
el cálculo a todas sus combinaciones suaves y prueba
\begin{equation}
 \left.\mathcal F^{\rm can}_{X,Y}\right|_{\mathcal Z}=0
 \quad\hbox{para }X,Y\hbox{ característicos}.
 \label{hmtcuatro:cierre-caracteristico}
\end{equation}
Para dos lapsos se incluye en \([X,Y]\) la deformación tangencial
y las acciones Gauss de \eqref{hmtcuatro:algebra-restricciones}.
Si se pasa al cociente gauge, la igualdad se expresa módulo esas
acciones verticales. No se afirma que \(\Omega\) sea cero en toda
la fase ni que toda holonomía global de una hoja sea trivial.

\par\medskip\noindent\textbf{7. Expresión precuántica del mismo cierre.}\quad
La realización diferencial asociada a
\eqref{hmtcuatro:conexion-canonica}, sobre secciones suaves donde
están definidos los productos, es
\begin{equation}
 Q_{\rm pre}(f)=-i\hbar\nabla_{X_f}+f
            =-i\hbar X_f+f-\Theta(X_f).
 \label{hmtcuatro:operador-precuantico}
\end{equation}
Con \(R^\nabla=(i/\hbar)\Omega\), la expansión del conmutador y
\(X_f(g)=-\{f,g\}_D\) dan
\begin{equation}
 [Q_{\rm pre}(f),Q_{\rm pre}(g)]
           =i\hbar Q_{\rm pre}(\{f,g\}_D).
 \label{hmtcuatro:representacion-precuantica}
\end{equation}
La regla de producto que conserva las funciones de estructura es
\begin{equation}
 Q_{\rm pre}(aC)
   =aQ_{\rm pre}(C)+C\bigl(Q_{\rm pre}(a)-a\bigr).
 \label{hmtcuatro:funciones-estructura}
\end{equation}
Se deduce de \(X_{aC}=aX_C+CX_a\) y de
\eqref{hmtcuatro:operador-precuantico}. No se elimina el segundo
término fuera de \(\mathcal Z\).
Sobre \(\mathcal Z\), \(Q_{\rm pre}(C_A)=-i\hbar\nabla_{X_{C_A}}\);
así \eqref{hmtcuatro:representacion-precuantica} y
\eqref{hmtcuatro:cierre-caracteristico} son las expresiones operatoria
y geométrica del mismo cierre canónico, no dos incorporaciones
independientes de gravedad.
\end{proof}

\subsection{Dominio del resultado y comparación de realizaciones}
\label{hmtcuatro:alcance-canonico}

La prueba reúne una acción con gravedad dinámica y fuente conjunta,
su Legendre y su conexión característica. El calificativo
\emph{arbitrario} para las deformaciones se refiere a los lapsos y
desplazamientos admisibles que generan direcciones características en
esa región regular; no a cualquier vector de fase, cualquier borde
o cualquier solución singular. La planitud local conserva la posible
monodromía global y no identifica dos historias TPK por tener iguales
extremos.

El coeficiente \(H_X^{\rm can}\) de una conexión de línea no es por
definición el operador de reloj y respuesta \(H_{\rm clk}+D^*WD\),
ni su extensión interactuante. \(Q_{\rm pre}(C)\) incluye además
una derivación en la fase; tampoco es un nombre alternativo de esos
operadores. Para identificar una realización concreta mediante un
mapa \(I\) se requiere la igualdad tipada
\[
 \nabla_N^{\rm op}I=I\nabla_N^{\rm can}
\]
sobre sus dominios efectivos. Si \(I\) depende de la geometría, su
derivada forma parte de esta igualdad.

La distinción tiene un control exacto: en la carta
\(\Theta=p\,\mathrm dq\),
\begin{equation}
 Q_{\rm pre}(p^2/2)=-i\hbar p\,\partial_q-p^2/2.
 \label{hmtcuatro:control-polarizacion}
\end{equation}
Al actuar sobre la función \(1\), produce \(-p^2/2\).
Por tanto no preserva por mera restricción la polarización de
funciones independientes de \(p\), y no es allí
\(-\hbar^2\partial_q^2/2\). El control no refuta el Hamiltoniano
recibido: impide sustituir una representación por otra sin su mapa.

No se ha supuesto una medida de Liouville infinito-dimensional ni
se deduce autoadjunción de las identidades diferenciales anteriores.
Si inclusiones de refinamiento entrelazan las conexiones y conservan
un núcleo común, sus curvaturas se entrelazan por composición; un
límite exige además los controles de dominio y convergencia propios.
No se presenta aquí esa condición como comprobada para el operador
interactuante y el límite espacial--quiral conjunto. Las construcciones
cuánticas anteriores se conservan; el resultado presente precisa
exactamente el cierre canónico que se acaba de demostrar.

% Exposición académica de las dos narrativas conservadas en el cuaderno.
% Procedencia y correspondencia de contenidos: gestion/mapa_narrativa.md.
% H_lazo y rho_lazo no designan la holonomía nonádica Gamma_9.
\section{Compatibilidad estructural de la acción, la masa y la información operacional}
\label{sec:narrativa-estructural}

La composición de las ecuaciones permite identificar una relación más precisa
que la equivalencia entre masa y energía o la proporcionalidad entre energía
y frecuencia: una misma estructura de ruta determina una frecuencia propia,
una relación entre longitudes y una ponderación térmica. El par angular
participa en estas tres lecturas porque contiene el coeficiente de acción y
determina los canales de transporte. La exposición que sigue conserva
conjuntamente la construcción de los objetos, sus operaciones de lectura y
las condiciones de sus realizaciones físicas.

Esta composición permite, además, caracterizar la huella que deja el orden
de una secuencia de operaciones. Dicha huella admite una expresión mediante
el par angular y una aplicación inversa que restituye el coeficiente de
acción. En el ámbito constitutivo, la respuesta del vacío y el momento
espectral de Catalán se reúnen mediante una misma colección funcional. El
contenido explicativo de estas relaciones reside en las restricciones que
imponen conjuntamente, y no en la mera acumulación de identidades.

La genealogía APP--TRIT--TPK proporciona los objetos que se componen. Las
lecturas posteriores conservan aspectos distintos de una misma estructura:
escala, orientación, fase, memoria, incidencia y respuesta. Su análisis
conjunto establece qué relaciones permanecen determinadas al cambiar de
representación, qué información elimina una proyección y qué observaciones
complementarias permiten restituirla. Las demostraciones de las composiciones
y de sus dominios se desarrollan en las secciones demostrativas.

\subsection{La acción en la composición angular y en la estructura electrónica}
\label{subsec:narrativa-accion-angular}

La sección de Planck interviene antes de la evaluación energética final.
Su coeficiente participa en la geometría angular utilizada posteriormente
por la ley de masas. El enlace, con las unidades conservadas, es
\begin{equation}
 \hbar_{\mathrm{ret}}
   =\eta_{\mathrm{ret}}\,10^{-34}\mathcal U_S,
 \qquad
 A_{\deg}=1000\alpha,
 \qquad
 C^*_{\deg}=2(\eta_{\mathrm{ret}}+\alpha).
 \label{eq:narrativa-par-angular}
\end{equation}
El coeficiente $\eta_{\mathrm{ret}}$ es adimensional y por ello puede
componerse con $\alpha$ en la expresión angular. La constante dimensional
$\hbar_{\mathrm{ret}}$ conserva su unidad de acción. Esta distinción permite
seguir una misma dependencia sin sumar magnitudes de dimensiones diferentes.

La acción desempeña así dos funciones relacionadas. Por una parte, fija una
escala dimensional; por otra, participa en la estructura que determina el
reparto del transporte entre las dos orientaciones. Después de la conversión
única de grados a radianes, el ángulo y el contraángulo producen los canales
$q_+$ y $q_-$. Estos canales intervienen en las incidencias hexádico--octádicas,
en el funcional de Barbero--Immirzi, en la respuesta constitutiva y en los
términos de la ley de masas.

La sustitución de una constante por su expresión estructural hace visibles
dependencias que permanecen implícitas cuando la ecuación se escribe mediante
constantes aisladas. La expresión ampliada transmite la operación que
determina el valor, las coordenadas de las que depende y las condiciones
bajo las cuales puede reutilizarse. Su función explicativa procede de esta
información transmitida a la etapa siguiente.

En el electrón, la expresión completa conserva el exponente
$\varphi/\pi^2$, las correcciones construidas con $\alpha$ y $\Delta_4$,
y el factor de retorno
\begin{equation}
 \left(\frac{H_5}{\eta_{\mathrm{ret}}}\right)^{80-54\alpha+6\Delta_4}.
 \label{eq:narrativa-factor-electronico}
\end{equation}
El mismo $\eta_{\mathrm{ret}}$ que interviene en el contraángulo participa
en la lectura electrónica completa. La geometría angular y la corrección
de acción del electrón comparten, por tanto, una dependencia efectiva.
Estudiar conjuntamente ambas ecuaciones permite determinar qué relaciones
quedan impuestas por la intervención de una misma coordenada estructural
en operaciones distintas.

\subsection{Frecuencia propia y normalización estructural de la masa}
\label{subsec:narrativa-masa-frecuencia}

La ley general de masas conserva su prefactor electrónico. Se designa por
$\Xi_\Gamma$ el factor estructural completo de una ruta $\Gamma$: reúne la
normalización electrónica, el carácter de la ruta con su firma y sus torres,
y la potencia correspondiente de la razón de acción
$H_5/\eta_{\mathrm{ret}}$. Esta notación abrevia una composición de
operaciones ya definidas; su contenido permanece ligado a esas operaciones.

El reloj común tiene frecuencia angular
$\omega_0=2\pi/(108t_0)$, con $t_0>0$. La composición de la ley de masas con la sección
de acción da
\begin{equation}
 \boxed{
 m_\Gamma=\frac{\hbar_{\mathrm{ret}}\omega_0}{c^2}\,\Xi_\Gamma,
 \qquad
 E_\Gamma=\hbar_{\mathrm{ret}}\omega_0\Xi_\Gamma,
 \qquad
 \Omega_\Gamma=\frac{E_\Gamma}{\hbar_{\mathrm{ret}}}
             =\omega_0\Xi_\Gamma.}
 \label{eq:narrativa-masa-energia-frecuencia}
\end{equation}
Aquí $E_\Gamma$ es la energía de reposo y $\Omega_\Gamma$ su frecuencia
energética asociada.

Esta lectura separa el ritmo común de la contribución propia de cada
estructura. Dos secciones pueden compartir el período elemental y presentar
masas y frecuencias propias diferentes: sus firmas, retornos y memorias
producen factores $\Xi_\Gamma$ diferentes. La diversidad de las lecturas
específicas queda referida a la estructura de la sección, mientras el reloj
común establece su normalización temporal.

En el cociente $E_\Gamma/\hbar_{\mathrm{ret}}$ se cancela el factor
absoluto de acción que aparece explícitamente en
\eqref{eq:narrativa-masa-energia-frecuencia}. Permanece, sin embargo, su
participación estructural en los canales y en la razón de retorno contenida
en $\Xi_\Gamma$. La cancelación de un factor en una lectura posterior
debe distinguirse de la eliminación de las relaciones mediante las cuales
ese factor se construyó y volvió a intervenir.

Los cocientes entre masas coinciden con los cocientes entre estas frecuencias
propias. Una modificación común de las unidades de acción, tiempo o velocidad
deja inalterados tales cocientes. La estructura relativa tiene, por ello,
una exigencia de contraste propia: sus relaciones entre secciones deben
mantenerse coherentes, con independencia de la normalización común elegida.

La frecuencia asociada a una fase levantada conserva, además, el número de
vueltas de la historia. Dos fases iguales después de su reducción módulo
$2\pi$ pueden corresponder a levantamientos diferentes. El retorno de fase
conserva una regularidad, mientras la historia distingue el recorrido que
produjo ese retorno. La identificación del estado completo con la fase
reducida eliminaría precisamente esta distinción.

\subsection{Dilatación recíproca de las lecturas espaciales de la masa}
\label{subsec:narrativa-reciprocidad-espacial}

En la realización gravitatoria circular, la acción y la gravitación
satisfacen
\begin{equation}
 G=\frac{c^3L_\alpha^2}{\hbar_{\mathrm{ret}}},
 \qquad R_{108}=L_\alpha=\frac{c}{\omega_0}.
 \label{eq:narrativa-gravedad-accion}
\end{equation}
La composición con la masa de
\eqref{eq:narrativa-masa-energia-frecuencia} determina
\begin{equation}
 \boxed{
 r_{g,\Gamma}=R_{108}\Xi_\Gamma,
 \qquad
 \bar\lambda_\Gamma=\frac{R_{108}}{\Xi_\Gamma},
 \qquad
 r_{g,\Gamma}\bar\lambda_\Gamma=R_{108}^{2}.}
 \label{eq:narrativa-longitudes-reciprocas}
\end{equation}
Se utiliza la convención de radio gravitatorio reducido del desarrollo;
la construcción de $L_\alpha$ y la normalización conjunta están
demostradas en la sección~\ref{md:rigidez}. El reloj y el radio son
dos expresiones compatibles de esa misma longitud.
$\bar\lambda_\Gamma$ designa la longitud de Compton reducida. Estas
identidades recíprocas se formulan para $\Xi_\Gamma>0$; la lectura
$R_{108}/\Xi_\Gamma$ requiere ese dominio másico positivo.

El mismo factor que aumenta la frecuencia propia dilata la lectura
gravitatoria y contrae la longitud asociada a la acción y la masa. Las dos
transformaciones conservan un área común. En esta realización, el factor
másico actúa como parámetro de dilatación recíproca: una longitud crece con
él y la otra disminuye inversamente.

Esta relación proporciona un contenido matemático preciso a la vinculación
entre masa y configuración espacial. La jerarquía de la composición es
esencial: primero se determina la ruta y su factor estructural; después se
obtienen sus lecturas temporales, energéticas y espaciales. La estructura
se transmite a las lecturas mediante operaciones determinadas, en lugar de
reconstruirse retrospectivamente a partir de una masa medida.

La condición circular pertenece al resultado
\eqref{eq:narrativa-longitudes-reciprocas}. Las dos longitudes son lecturas
definidas por esa realización; su interpretación como radios materiales de
una partícula requeriría la identificación física correspondiente. La
igualdad establece aquí una reciprocidad exacta entre las magnitudes
definidas, con la coordenatización de acción, el reloj y la velocidad conservados.

La incorporación del momento mantiene esta masa estructural en la ecuación
de propagación. En la realización libre correspondiente, la frecuencia
propia determina el término másico de la dispersión, mientras el momento
describe la variación espacial. Así quedan diferenciadas la estructura
persistente de la sección y su estado de movimiento.

\subsection{Orientación de los canales y respuesta constitutiva del vacío}
\label{subsec:narrativa-vacio-orientacion}

La composición de permitividad, permeabilidad y funcional de Barbero--Immirzi
permite distinguir el valor de una respuesta de la estructura que lo
determina. El lector constitutivo utiliza dos canales etiquetados, con
respuestas $r_+$ y $r_-$. La velocidad normalizada depende de su producto;
la impedancia depende de su cociente. La velocidad conserva una relación
común entre los canales, mientras la impedancia distingue el reparto de la
respuesta entre sus orientaciones.

Dentro de la colección normalizada, la velocidad y el funcional de
Barbero--Immirzi permiten recuperar conjuntamente las dos respuestas
etiquetadas. El funcional aporta información orientada que la velocidad
aislada no determina. La restitución se apoya en la composición completa
de sus lectores y conserva la distinción entre cada canal y su agregado.

Sean $x$ e $y$ las coordenadas en radianes del ángulo y del contraángulo.
Para $x>0$ y $|y|<x$, la respuesta satisface
\begin{equation}
 \widehat c(x,y)\geq\widehat c(x,0),
 \label{eq:narrativa-convexidad-velocidad}
\end{equation}
con igualdad únicamente cuando $y=0$. Más precisamente, manteniendo $x$
fijo, $\log\widehat c$ es una función par y estrictamente convexa de $y$.
La demostración de esta propiedad pertenece al desarrollo constitutivo
posterior.

La velocidad pierde el signo de la orientación, pero conserva un efecto
de la separación entre los canales. Cerca de la configuración simétrica,
ese efecto aparece a segundo orden; la impedancia distingue la orientación
ya a primer orden. Por ello, sustituir ambos canales por su media antes de
evaluar la respuesta cambia el resultado. Una lectura escalar puede ocultar
una orientación y conservar, al mismo tiempo, una consecuencia de ella.

Este comportamiento especifica qué información transmite cada operación.
La pérdida del signo, la conservación de una separación y la recuperación
del reparto constituyen propiedades diferentes. La estructura mantiene
consecuencias observables incluso cuando una de sus lecturas se expresa
mediante un único número.

La desigualdad \eqref{eq:narrativa-convexidad-velocidad} describe la colección
matemática de respuesta en su dominio. Una evolución física debe situarse
en los estados y trayectorias efectivamente producidos por TPK. La
parametrización analítica permite estudiar la sensibilidad del lector;
la identificación de una variación paramétrica con una trayectoria física
conserva su obligación de realización.

\subsection{Boltzmann, frecuencia estructural y ponderación térmica ternaria}
\label{subsec:narrativa-boltzmann}

La sección térmica relaciona la energía del ciclo con la constante de
Boltzmann y con su temperatura correspondiente:
\begin{equation}
 k_B\Theta_{108}\log3=E_{\mathrm{clk}}.
 \label{eq:narrativa-boltzmann-ciclo}
\end{equation}
Como la energía de una sección es
$E_\Gamma=E_{\mathrm{clk}}\Xi_\Gamma$, el peso térmico evaluado a esa
temperatura resulta
\begin{equation}
 \boxed{\exp\!\left(-\frac{E_\Gamma}{k_B\Theta_{108}}\right)
       =3^{-\Xi_\Gamma}.}
 \label{eq:narrativa-peso-termico}
\end{equation}
La base ternaria aparece por composición de la normalización informacional
con la acción y el reloj. La estructura de ruta que determina la masa
determina también el exponente del peso térmico.

El peso y la probabilidad desempeñan funciones distintas. Las probabilidades
se obtienen al normalizar los pesos y conservar las multiplicidades de los
estados. Asimismo, varias historias equivalentes requieren una identificación
de estados antes de recibir una multiplicidad física: el número de historias
y la degeneración del estado corresponden a objetos diferentes.

La composición permite reunir las lecturas en una identidad común:
\begin{equation}
 \boxed{
 \Xi_\Gamma
 =\frac{\Omega_\Gamma}{\omega_0}
 =\frac{r_{g,\Gamma}}{R_{108}}
 =\frac{R_{108}}{\bar\lambda_\Gamma}
 =\frac{E_\Gamma}{k_B\Theta_{108}\log3}.}
 \label{eq:narrativa-cuatro-lecturas}
\end{equation}
Se obtienen cuatro formas coordinadas de leer una misma componente
estructural: temporal, espacial, gravitatoria y térmica. Cada igualdad
conserva las condiciones de su realización, especialmente la condición
circular gravitatoria y la temperatura del ciclo.

La identidad \eqref{eq:narrativa-cuatro-lecturas} especifica la relación
que comparten las magnitudes y proporciona un criterio para contrastar su
compatibilidad. La masa relativa puede caracterizarse mediante la frecuencia
relativa; la misma coordenada determina una dilatación recíproca de
longitudes y una ponderación térmica. Las funciones de lectura son distintas,
pero su dependencia queda vinculada por la genealogía común.

La relación térmica identifica el estado y la normalización en los que
interviene $\Theta_{108}$. Esta identificación mantiene separadas la energía
de la sección y su eventual transferencia como calor; también mantiene
localizado el significado de la temperatura del ciclo, sin extenderla al
universo como temperatura única. La capacidad explicativa procede de las
condiciones explícitas de la composición.

\subsection{Elipse de acción y transporte de la forma geométrica}
\label{subsec:narrativa-elipse-transporte}

La elipse de acción tiene semiejes cuyo producto es $2\hbar$. El área
orientada acumulada durante una vuelta es $h=2\pi\hbar$. Una deformación
elíptica redistribuye los semiejes conservando su producto. En la
realización eléctrica y magnética se obtienen dos contribuciones
proporcionales a $\cos^2\theta$ y $\sin^2\theta$, cuya suma es
$\hbar\omega$. La estructura determina, por tanto, el reparto de la
energía durante el ciclo, además de su valor total.

El transporte de una elipse cuya forma cambia exige conservar la relación
entre la deformación y la acción. Sea $\zeta$ la coordenada de deformación,
determinada por la razón entre el contraángulo y el ángulo. Para una
deformación diferenciable, el transporte exacto introduce
\begin{equation}
 H_{\mathrm{tot}}
 =\frac{\omega}{2}
   \left(e^{-\zeta}Q^2+e^\zeta P^2\right)
  +\frac{\dot\zeta}{2}QP.
 \label{eq:narrativa-hamiltoniano-transporte}
\end{equation}
El primer término describe la evolución en una elipse de forma fija.
El segundo procede del transporte de una forma variable. La demostración
de la identidad de transporte establece que este término compensa
exactamente la variación geométrica y conserva la acción cuadrática
correspondiente. Su omisión o la inversión de su signo destruyen esa
compensación.

La interpretación exige distinguir dos situaciones. Cuando se cambia
únicamente de coordenadas, el término adicional es una conexión de
representación. Un intercambio físico requiere, en cambio, una variación
relativa entre sistema y referencia. La reexpresión de un Hamiltoniano
conserva el contenido de la representación; la atribución de una interacción
necesita determinar las operaciones relativas que actúan sobre el sistema.

El par angular rector conserva otra dependencia decisiva: fijar $\alpha$
y la sección de acción fija también $\zeta$. Una evolución de la forma
debe, por ello, extraerse de los estados y operaciones efectivos. La
deformación paramétrica por sí sola describe una colección de representaciones;
su identificación como evolución HMT transmite las relaciones que producen
el ángulo y el contraángulo.

El transporte conjunto de la forma y de la sección de acción conserva la
acción normalizada, mientras el área sin normalizar cambia de acuerdo con
el cociente entre las dos secciones. Esta construcción distingue tres
operaciones: redistribuir una acción conservada, transportar una sección de
acción y cambiar la representación o las unidades. La distinción permite
examinar interacciones conservando un balance que incluya el origen de la
variación, en lugar de atribuir generación de energía a una mera deformación
de coordenadas.

\subsection{Compatibilidad conjunta y contraste relacional}
\label{subsec:narrativa-compatibilidad-conjunta}

El eje de investigación es la determinación de las condiciones de
compatibilidad que una misma estructura impone a sus diferentes
manifestaciones físicas. Dos líneas de composición se complementan.

La primera reúne relaciones estáticas: los observables complementarios
restituyen los canales, y la estructura recuperada determina la acción,
la respuesta constitutiva y los factores de masa. El contraste pertinente
es conjunto. Una modificación que preserve una cifra debe conservar también
las demás relaciones que utiliza la misma construcción. La velocidad,
Barbero--Immirzi y la impedancia proporcionan un sistema de comprobación;
las razones de masas y de frecuencias proporcionan otro.

La segunda línea estudia las historias efectivas y compara sus resultados
después de transportar coherentemente el sistema y su referencia. Las
diferencias de fase relativa, respuesta constitutiva relativa o transferencia
entre modos tienen un contenido distinto de las diferencias debidas al
nombre de las coordenadas. La memoria, la orientación y los retornos
intervienen mediante las operaciones que una lectura escalar puede eliminar.

Una relación estructural alcanza capacidad explicativa y de contraste
cuando varias lecturas independientes deben concordar por proceder de una
misma estructura. Su valor reside en la rigidez de esa concordancia:
acción, ritmo, respuesta espacial y ponderación informacional dejan de
poder elegirse por separado dentro de la realización declarada. La
composición reúne esas restricciones y permite demostrar sus consecuencias
sin sustituir su genealogía por ajustes individuales de los resultados.

\subsection{Cierre de fase y holonomía relativa del estado}
\label{subsec:narrativa-holonomia-relativa}

El ciclo dodecafásico contiene el plano en el que se realiza la raíz
interna de $-1$. La inclusión de ese plano en las doce componentes es
explícita; el cuarto de giro utilizado a continuación procede de dicha
construcción. Sobre él actúa la forma de la elipse de acción determinada
por el par angular, cuya coordenada es
\begin{equation}
 \xi=\frac{C^*}{A},\qquad
 \zeta=\operatorname{artanh}\xi.
 \label{eq:narrativa-coordenada-forma}
\end{equation}
El cociente utiliza los dos ángulos expresados en la misma unidad y tiene
el mismo valor en grados y en radianes. La coordenada real de forma se
define para $|\xi|<1$. La forma conjugada corresponde al
cambio de signo de $\zeta$.

Se compone el cuarto de giro de una forma con el de su conjugada y con
sus inversos. Si $U_+$ y $U_-$ designan las dos operaciones, la
composición exacta es
\begin{equation}
 \boxed{
 H_{\mathrm{lazo}}=U_+U_-U_+^{-1}U_-^{-1}
 =\begin{pmatrix}\rho_{\mathrm{lazo}}&0\\[2pt]
                  0&\rho_{\mathrm{lazo}}^{-1}\end{pmatrix},
 \qquad
 \rho_{\mathrm{lazo}}=
 \left(\frac{1+\xi}{1-\xi}\right)^2.}
 \label{eq:narrativa-holonomia-lazo}
\end{equation}
La composición dilata una coordenada y contrae la conjugada por el factor
inverso: conserva el área y modifica el estado. Para $\xi\ne0$, su
resultado difiere de la identidad.

Las cuatro operaciones incluyen sus inversas y sus incrementos de fase
declarados suman cero. El orden de composición conserva, sin embargo, un
efecto. Una descripción que sólo retenga los incrementos que se añaden y
se compensan elimina ese efecto relativo. El cierre de la contabilidad
de fase y el retorno del estado completo son, por tanto, propiedades
distintas de la secuencia.

La igualdad admite un levantamiento al espacio original de doce
componentes. El operador actúa mediante la dilatación recíproca sobre el
plano considerado y como identidad sobre su complemento. El enlace con
el ciclo dodecafásico se realiza mediante una transformación explícita;
el plano conserva así su procedencia dentro del espacio de fases.

La notación $H_{\mathrm{lazo}}$ designa esta composición posterior y
mantiene diferenciada la holonomía nonádica $\Gamma_9$. La realización
ilustra de manera concreta una distinción estructural: el cierre de una
lectura del recorrido puede coexistir con una transformación persistente
del estado. La procedencia de los operadores y el alcance de su
levantamiento permanecen expuestos en el desarrollo demostrativo.

\subsection{Restitución del coeficiente de acción desde el transporte relativo}
\label{subsec:narrativa-restitucion-accion}

El par angular de \eqref{eq:narrativa-par-angular} puede expresarse con
$s_0=10^{-34}\mathcal U_S$ y
$\hbar_{\mathrm{ret}}=s_0\eta_{\mathrm{ret}}$. El coeficiente
$\eta_{\mathrm{ret}}$ conserva su carácter adimensional, mientras $s_0$
conserva la escala dimensional. Al incorporar el contraángulo completo
en \eqref{eq:narrativa-holonomia-lazo}, se obtiene
\begin{equation}
 \boxed{
 \rho_{\mathrm{lazo}}=
 \left(\frac{501\alpha+\eta_{\mathrm{ret}}}
              {499\alpha-\eta_{\mathrm{ret}}}\right)^2.}
 \label{eq:narrativa-lazo-alfa-accion}
\end{equation}
Esta expresión utiliza la rama positiva
$0<\eta_{\mathrm{ret}}<499\alpha$. Los números $501$ y $499$ proceden
de $500\pm1$, y $500$ procede de $1000/2$: son consecuencias de las
normalizaciones que ya intervienen en el par angular.

La relación entre $\alpha$ y el coeficiente de acción determina la
separación entre las respuestas conjugadas después del recorrido
compuesto. La acción participa tanto en la escala dimensional como en
la forma que gobierna el transporte. La fórmula hace explícita esta
doble intervención sin identificar ambas funciones.

La relación admite la inversión
\begin{equation}
 \boxed{
 \eta_{\mathrm{ret}}=
 \alpha\left[
 500\,\frac{\sqrt{\rho_{\mathrm{lazo}}}-1}
               {\sqrt{\rho_{\mathrm{lazo}}}+1}-1\right].}
 \label{eq:narrativa-inversa-accion}
\end{equation}
Con $\alpha$ y los ejes orientados conservados, la respuesta del lazo
permite restituir el coeficiente de acción. La lectura es estrictamente
creciente respecto de $\rho_{\mathrm{lazo}}$ en la rama declarada.

Se obtiene una caracterización operacional: el coeficiente de acción
puede restituirse mediante la dilatación relativa producida por las
formas conjugadas del par angular. Su definición generadora permanece
en la construcción anterior; la caracterización operacional establece
cómo recuperarlo en una representación distinta.

La escala $s_0$ sigue siendo necesaria para recuperar la constante
dimensional $\hbar_{\mathrm{ret}}$. La forma del transporte conserva
una relación adimensional; la escala de acción determina la
normalización física completa. Esta separación identifica con precisión
qué dato devuelve la inversión y qué estructura adicional transmite la
lectura dimensional.

\subsection{Iteración de la holonomía e información operacional de la historia}
\label{subsec:narrativa-informacion-operacional}

La repetición del lazo conserva el mismo programa de operaciones y
produce una ley explícita de acumulación:
\begin{equation}
 H_{\mathrm{lazo}}^n=
 \begin{pmatrix}\rho_{\mathrm{lazo}}^n&0\\
                 0&\rho_{\mathrm{lazo}}^{-n}\end{pmatrix}.
 \label{eq:narrativa-iteracion-lazo}
\end{equation}
Una preparación $(Q_0,P_0)$ se transforma en
\begin{equation}
 Q_n=\rho_{\mathrm{lazo}}^nQ_0,
 \qquad
 P_n=\rho_{\mathrm{lazo}}^{-n}P_0.
 \label{eq:narrativa-iteracion-estado}
\end{equation}
El producto $Q_nP_n$ permanece constante y el transporte conserva
el área simpléctica, mientras cambia la proporción entre las dos
coordenadas.

Una regla finita repetida puede dejar una huella acumulativa que
distinga sus iteraciones. Una preparación conocida y una referencia
conservada permiten recuperar $n$ mediante el incremento logarítmico
de una coordenada no nula, para $\rho_{\mathrm{lazo}}\ne1$.
Cuando se conserva
únicamente el estado final y se desconoce la preparación, la misma
restitución queda indeterminada.

La información accesible depende de la lectura retenida: balance,
amplitud, proporción orientada, preparación u operador completo. Cada
elección puede identificar historias que otra lectura distingue.
La comparación de
\begin{equation}
 U_+U_-U_+^{-1}U_-^{-1}
 \qquad\text{y}\qquad
 U_+U_+^{-1}U_-U_-^{-1}
 \label{eq:narrativa-orden-multiplicidades}
\end{equation}
proporciona un contraste directo. Ambos productos contienen las mismas
operaciones con las mismas multiplicidades. El segundo da la identidad;
el primero da $H_{\mathrm{lazo}}$. Una estadística de frecuencias de
aparición resulta insuficiente para predecir esta respuesta: debe
conservarse también el orden.

Este resultado precisa qué descripción estadística es insuficiente
para el objeto estudiado y qué observaciones restituyen la diferencia.
La teoría de información de Shannon admite distribuciones sobre
secuencias ordenadas y puede representar ese orden. La distinción
operacional se refiere al contenido retenido por el observable elegido.

La información operacional de una historia se caracteriza mediante
su capacidad de producir respuestas distinguibles para preparaciones
y detectores especificados. Esta definición permite precisar cuánto
conserva una lectura. Su alcance corresponde a la representación
utilizada: historias distintas pueden inducir un mismo operador
reducido y conservar diferencias en otras componentes de la genealogía
TPK.

La traza de $H_{\mathrm{lazo}}$ coincide con la de
$H_{\mathrm{lazo}}^{-1}$. Detecta la intensidad del efecto, pero elimina
su orientación. Una lectura que conserve los ejes y compare sus
respuestas distingue los dos sentidos. La elección de detector forma
parte, por tanto, de la restitución del recorrido.

La finitud del programa, el crecimiento de una memoria observable y
una tasa entrópica describen propiedades diferentes. Una tasa de
crecimiento combinatorio nula puede coexistir con estados distinguibles.
La entropía de la dinámica completa requiere especificar su espacio
de estados y, cuando corresponda, su medida. La composición demuestra
la acumulación y su restitución en la representación declarada;
la evaluación de una entropía termodinámica conserva sus variables
y condiciones físicas propias.

\subsection{Potencial espectral y reciprocidad diferencial de las respuestas}
\label{subsec:narrativa-potencial-espectral}

Los canales constitutivos conservan las coordenadas
$x=A_{\mathrm{rad}}$, $y=C^*_{\mathrm{rad}}$ y las incidencias $90/120$.
Su composición determina la identidad
\begin{equation}
 \boxed{
 \frac{\partial\log\widehat c}{\partial y}
 =\frac{\partial\log\widehat Z}{\partial x}.}
 \label{eq:narrativa-reciprocidad-diferencial}
\end{equation}
Las magnitudes $\widehat c$ y $\widehat Z$ designan la velocidad y la
impedancia normalizadas del lector.

La sensibilidad de la propagación a la separación orientada queda
ligada a la sensibilidad de la impedancia a la coordenada común.
Las dos respuestas comparten una función espectral y sus variaciones
deben mantener esa relación para seguir perteneciendo a la misma
construcción.

Existe una función $\mathscr F(x,y)$ cuyas derivadas son
\begin{equation}
 \partial_x\mathscr F=\log\widehat c,
 \qquad
 \partial_y\mathscr F=\log\widehat Z.
 \label{eq:narrativa-gradiente-potencial}
\end{equation}
Se trata de un potencial espectral definido desde los canales $90/120$,
con serie convergente y normalización explícita. Su Hessiano es
negativo definido en $x>|y|$, propiedad que permite demostrar
unicidad y sensibilidad de la restitución.

La igualdad entre las derivadas cruzadas establece una restricción
sobre la variación conjunta de las respuestas dentro de esta colección
de realización. Una modificación que destruya la identidad
\eqref{eq:narrativa-reciprocidad-diferencial} cambia el lector
constitutivo. El contraste afecta así a las relaciones de variación,
además de a los valores evaluados.

El potencial es de valor único: la integral de su diferencial sobre
un lazo cerrado en las dos coordenadas se anula. La composición
ordenada de transportes puede conservar, en cambio, una holonomía
no trivial. La proyección constitutiva y el transporte operatorio
retienen aspectos diferentes de la historia; su comparación exige
conservar sus dominios y operaciones de lectura.

\subsection{Momento de Catalán y respuesta contractiva de las incidencias}
\label{subsec:narrativa-catalan-potencial}

La conexión entre el momento orientado de Catalán y el potencial
constitutivo se expresa mediante el momento de profundidad dos
\begin{equation}
 D_2(z)=\sum_{n\geq1}\frac{z^n}{n^2}.
 \label{eq:narrativa-momento-dos}
\end{equation}
El momento de Catalán se obtiene como la parte imaginaria de
$D_2(is)$; en $s=1$ recupera la constante de Catalán. El potencial
constitutivo utiliza la misma función sobre las contracciones de
los canales, en $q_\pm^{90}$ y $q_\pm^{120}$, con las normalizaciones
correspondientes.

El cuarto de giro orientado y la respuesta contractiva $90/120$
se reúnen mediante un mismo momento espectral, evaluado sobre
soportes distintos. Catalán adquiere aquí significado por la
operación espectral que representa y por la relación de esa
operación con la respuesta constitutiva. La colección funcional
conserva información que excede la de su valor extremo.

La composición mantiene también las diferencias entre sus objetos.
El valor aislado de Catalán no determina todos los canales;
la relación utiliza el momento espectral completo, sus argumentos
y la orientación. La convergencia de la serie, las identidades
de sus derivadas y su compatibilidad con la amplificación forman
parte del desarrollo demostrativo de esta conexión.

\subsection{Restitución exacta y estabilidad ante precisión finita}
\label{subsec:narrativa-estabilidad-restitucion}

El par $(\widehat c,\widehat Z)$ recupera de manera única los dos
parámetros del lector en el dominio declarado. Las cotas de
sensibilidad relacionan, además, el error de las respuestas con
el error de los parámetros restituidos. La unicidad exacta y la
estabilidad de la inversión constituyen resultados distintos y
complementarios.

Esta distinción separa una información eliminada por la proyección
de una información conservada pero difícil de resolver con precisión
finita. En regiones de contracciones intensas, variaciones apreciables
de los parámetros pueden producir diferencias muy pequeñas en las
respuestas. La inversión sigue siendo única para datos exactos,
mientras una incertidumbre finita puede impedir distinguir tales
variaciones.

La restitución genealógica exige especificar tanto qué estructura
puede recuperarse como la estabilidad con la que puede recuperarse.
El margen de restitución designa la cota que cuantifica esa estabilidad
en un dominio. La elección de observable debe atender a la diferencia
que se pretende resolver: una lectura puede ser correcta y resultar
poco sensible a ella.

La traza del lazo ofrece un ejemplo de esta selección. Su lectura
elimina el sentido del recorrido; la respuesta sobre ejes orientados
lo conserva. La determinación del detector adecuado pertenece a
la demostración de recuperabilidad y al diseño del contraste.

\subsection{Síntesis: coordenadas de la sección y memoria de las operaciones}
\label{subsec:narrativa-sintesis}

La identidad \eqref{eq:narrativa-cuatro-lecturas} conserva como eje
la relación entre el factor estructural de una sección, su frecuencia
propia, las dos longitudes recíprocas y su lectura térmica. Permanecen
explícitas las condiciones: la misma coordenatización de acción, la realización
gravitatoria circular y la temperatura del ciclo.

La composición operatoria añade una distinción central. El factor
de la sección determina sus lecturas coordinadas; el transporte
determina el resultado de componer operaciones sobre ella y la
huella que conserva su orden. El factor másico y la dilatación
$\rho_{\mathrm{lazo}}$ son objetos diferentes. Su relación debe
pasar por los lectores correspondientes para distinguir una
transformación del estado de una identificación másica.

Planck interviene en la normalización de la acción y en la forma
angular del transporte. Boltzmann interviene en la relación entre
energía y ponderación térmica. La reciprocidad gravitatoria conserva
el producto de las longitudes en su realización declarada. Velocidad
e impedancia proporcionan lecturas complementarias de los canales.
Esta arquitectura admite conjuntamente una descripción de
compatibilidades y una descripción de la memoria operacional.

La capacidad explicativa de la estructura reside en determinar
varias respuestas, imponer relaciones entre ellas y especificar
qué historia puede restituirse desde sus observaciones. Las
coincidencias de valores quedan acompañadas por operaciones,
condiciones de compatibilidad y criterios de restitución.

La composición distingue el cambio pasivo de coordenadas, que
telescopa cuando retornan coordenatización y fase, del transporte relativo
cuyo operador final es $H_{\mathrm{lazo}}\ne I$. Un intercambio
físico requiere realizar las operaciones, sus inversas y la
referencia en un mismo sistema, incluyendo el controlador en
el balance. Este criterio conserva el contenido matemático
de una variación energética sin atribuirle producción de energía
por el mero cambio de lectura.

Las pruebas de composición, iteración, restitución del coeficiente
de acción, reciprocidad constitutiva y sensibilidad sostienen
estas caracterizaciones en sus dominios respectivos. Narración,
definiciones y demostraciones forman así una exposición continua:
la primera identifica el significado de las relaciones, las
segundas fijan los objetos y las terceras establecen las
operaciones que permiten componerlos y recuperarlos.

% Procedencia íntegra: CONTRAANGULO_ELECTRON_BARBERO_CATALAN.md, §§2--10.
% La gestión, el historial y los localizadores se conservan en gestion/.
\section{Composición y restitución de las lecturas angular, electrónica y constitutiva}
\label{md:compos:seccion}

\subsection{Coeficiente de acción y construcción del contraángulo}
\label{md:compos:contraangulo}

La abreviatura funcional $B(\alpha)$ reúne una evaluación en $\alpha$; su
argumento pertenece a la función y no representa el producto de una constante
independiente por $\alpha$. Para mantener explícitas las operaciones que
intervienen en esta abreviatura, se escribe
\begin{equation}
 B(\alpha)=\frac{9}{16}\alpha^2-\frac{5}{9}\alpha^3
 +\frac{7}{48}\alpha^4-\frac{1}{54}\alpha^5
 -\frac{90}{\pi}\mathcal C_\pi(\alpha),
 \label{md:compos:b-alpha}
\end{equation}
\begin{equation}
 \mathcal C_\pi(\alpha)=169\alpha^6+
 \frac{D_A\alpha^7}{1-D_A\alpha},
 \qquad D_A=\exp\!\left(-\frac{100\pi\alpha}{9}\right).
 \label{md:compos:correccion-regional}
\end{equation}
La identidad completa es
\begin{equation}
 C^*_{\deg}=\sqrt{\frac{1000\alpha}{\varphi}}+2B(\alpha).
 \label{md:compos:contraangulo-abreviado}
\end{equation}
En particular, $\alpha$ pertenece al radicando, y el término radical se
conserva junto a $2B(\alpha)$. Esta identidad reescribe el mapa construido
anteriormente y mantiene su objeto y su genealogía.

La definición original conserva separados el coeficiente adimensional y la
sección de acción:
\begin{equation}
 \eta_{\mathrm{ret}}=H_5-\frac{90}{\pi}\mathcal C_\pi,
 \qquad
 \hbar_{\mathrm{ret}}=\eta_{\mathrm{ret}}\,10^{-34}\mathcal U_S.
 \label{md:compos:seccion-accion}
\end{equation}
Por tanto,
\begin{equation}
 \boxed{C^*_{\deg}=2\left(
 \frac{\hbar_{\mathrm{ret}}}{10^{-34}\mathcal U_S}+\alpha\right)
 =2(\eta_{\mathrm{ret}}+\alpha).}
 \label{md:compos:contraangulo-accion}
\end{equation}
La expresión utiliza el coeficiente adimensional de la sección de Planck
reducida. La constante $h$ se obtiene después mediante $h=2\pi\hbar$.
La composición consiste en extraer el coeficiente de acción, componerlo con
$\alpha$ y aplicar la clausura bidireccional. Así se mantiene el tipo
dimensional de cada operación.

La construcción de masas expresa partícula y antipartícula como dos
secciones orientadas de una misma banda genealógica. La operación
dodecafásica es
\begin{equation}
 C_M(\tau)=-\operatorname{rev}(\tau),
 \qquad \tau_e=+++---+++---,
 \qquad C_M(\tau_e)=\tau_e,
 \label{md:compos:conjugacion-dodecafasica}
\end{equation}
prolongada al registro y a la representación de carga. La ruta central es
fija bajo $C_M$, mientras la representación de carga distingue electrón y
positrón. Esta composición da contenido preciso a la lectura bidireccional.
El factor dos angular conserva su función de clausura; la conjugación física
conserva además los operadores indicados. El retorno a $216$ restituye el
levantamiento orientado: su estructura no consiste en asignar $108$
actualizaciones a la partícula y otras $108$ a la antipartícula.

\subsubsection{Restitución radial del coeficiente de acción}
\label{md:compos:quinientos}

Escribimos $q_{\mathrm{rad}}=q_\varepsilon$ y $\epsilon=\varepsilon$ para
la razón radial y la orientación de \eqref{rec:eq:eta-orientada}.
El invariante radial orientado $\chi_{\mathrm{rad}}=\epsilon\log q_{\mathrm{rad}}$
permite escribir la restitución previamente construida como
\begin{equation}
 \frac{\eta_{\mathrm{ret}}}{\alpha}
 =-500\tanh\!\left(\frac{\chi_{\mathrm{rad}}}{2}\right)-1.
 \label{md:compos:restitucion-radial}
\end{equation}
Con $\xi=C^*_{\deg}/A_{\deg}$ y $A_{\deg}=1000\alpha$, se obtiene también
\begin{equation}
 \eta_{\mathrm{ret}}=\alpha(500\xi-1).
 \label{md:compos:restitucion-xi}
\end{equation}
Son recuperadores posteriores de la misma sección; su construcción permanece
en $H_5$ y en el retorno regional.

\subsection{Coordenatizaciones angulares y conservación de la acción}
\label{md:compos:cartas}

La coordenatización en grados tiene una vuelta de $360$ unidades y conserva las
particiones $12\cdot30=4\cdot90=3\cdot120=360$. El par angular construido es
\begin{equation}
 A_{\deg}=1000\alpha,
 \qquad C^*_{\deg}=2(\eta_{\mathrm{ret}}+\alpha).
 \label{md:compos:par-angular}
\end{equation}
La conversión se efectúa una vez:
\begin{equation}
 x=\frac{\pi}{180}A_{\deg},\qquad
 y=\frac{\pi}{180}C^*_{\deg},\qquad
 q_\pm=\exp(-x\mp y).
 \label{md:compos:canales}
\end{equation}
Los exponentes utilizan representantes reales levantados y conservan el
número de vueltas. La acción satisface
\begin{equation}
 \hbar\theta_{\mathrm{rad}}
 =\frac{h\theta_{\deg}}{360}.
 \label{md:compos:accion-cartas}
\end{equation}
En el funcional bilineal de Barbero, el transporte de coordenatización da
\begin{equation}
 \frac{\pi A_{\deg}C^*_{\deg}}{180}
 =\frac{180xy}{\pi}.
 \label{md:compos:bilineal-cartas}
\end{equation}
La transformación bilineal forma parte de la misma conversión que se aplica
a los argumentos exponenciales; omitirla cambiaría el funcional.

\subsection{Lectura electrónica completa y retorno de acción}
\label{md:compos:electron}

Se conservan las dos etapas de la misma coordenada electrónica:
\begin{equation}
 \varepsilon_e^{(0)}=\frac{\sqrt{3}}{4}
 \left[\exp\!\left(\frac{\varphi}{\pi^2}\right)-22\alpha^3\right]
 (1+15\Delta_4),\qquad
 \Delta_4=\frac{\pi-e\log\pi}{270}
 \left(1+\frac{\pi}{729}\right),
 \label{md:compos:electron-basal}
\end{equation}
\begin{equation}
 \boxed{\varepsilon_e^\beta=\varepsilon_e^{(0)}
 \left(\frac{H_5}{\eta_{\mathrm{ret}}}\right)^{80-54\alpha+6\Delta_4}.}
 \label{md:compos:electron-completo}
\end{equation}
El centro $(5,5)$, su inversión, las deformaciones $(8,2)/(2,8)$ y el
defecto unitario de cociente son antecedentes del valor. La norma
$\sqrt{3}/4$ procede del conmutador entre proyecciones aditivas y
multiplicativas; el carácter exponencial utiliza la orientación
$\varphi/\pi^2$; $22=27-5$ conserva un complemento regional y $15=3\cdot5$
las incidencias correspondientes. Los lectores de la ruta producen
$(80,54,6)$. El cociente $H_5/\eta_{\mathrm{ret}}$ incorpora el retorno de
acción a la misma sección central.

La función del contraángulo se explicita sustituyendo su inversa demostrada:
\begin{equation}
 \boxed{\varepsilon_e^\beta=\varepsilon_e^{(0)}
 \left(\frac{H_5}{C^*_{\deg}/2-\alpha}\right)^{80-54\alpha+6\Delta_4}.}
 \label{md:compos:electron-contraangulo}
\end{equation}
Esta expresión conserva la generación electrónica anterior. La misma
sección de acción que participa en el contraángulo determina el factor que
lleva la lectura basal a la completa. El cociente cancela la base dimensional
común; la realización absoluta de masa conserva su lector temporal y su coordenatización
dimensional. La sustitución es una composición de dependencias construidas,
sin introducir un ángulo externo como generador electrónico.

\subsection{Funcional de Barbero y transición hexada--octada}
\label{md:compos:barbero}

La incidencia mantiene el mismo par marcado al ampliar el soporte:
\begin{equation}
 \mathcal F_6=\mathsf H\times\binom{\mathsf H}{2}
 \hookrightarrow
 \mathcal F_8=\mathsf O\times\binom{\mathsf H}{2},
 \qquad |\mathcal F_6|=90,\qquad |\mathcal F_8|=120.
 \label{md:compos:incidencia}
\end{equation}
Los cardinales son $6\cdot15$ y $8\cdot15$; el referente compartido permite
comparar ambas incidencias. La diferencia normalizada es $-1/12$. La cadena
dodecafásica conserva doce contribuciones hexádicas y una contribución
octádica de frontera con signo contrario. El lector
\begin{equation}
 S_m(q)=\frac{q^m}{1-q^{3m}}
 =\sum_{j\geq0}q^{m+3mj}
 \label{md:compos:serie-retorno}
\end{equation}
conserva altura y retornos. Así se forma $12S_{90}-S_{120}$ y su diferencia
entre las dos orientaciones. El funcional completo es
\begin{equation}
 \gamma=\frac{\pi A_{\deg}C^*_{\deg}}{180}
 +12\bigl[S_{90}(q_-)-S_{90}(q_+)\bigr]
 -\bigl[S_{120}(q_-)-S_{120}(q_+)\bigr].
 \label{md:compos:funcional-barbero}
\end{equation}
Barbero es, en esta realización, la evaluación orientada conjunta del par
angular y de los retornos de la transición de incidencias. Su participación
en la escala de área se compone después con $\ell_P^2=G\hbar/c^3$ y con el
operador de ocupación. La entropía de entrelazamiento conserva además el
estado reducido y sus probabilidades; su reconstrucción requiere esos datos
y no queda determinada por el área escalar aislada.

\subsubsection{Eliminación de la coordenada común electrón--Barbero}
\label{md:compos:eliminacion}

Definamos $\kappa_e=80-54\alpha+6\Delta_4$. En la rama positiva, con
$\kappa_e\neq0$, la ecuación electrónica implica exactamente
\begin{equation}
 \eta_{\mathrm{ret}}=H_5
 \left(\frac{\varepsilon_e^{(0)}}{\varepsilon_e^\beta}\right)^{1/\kappa_e}.
 \label{md:compos:eliminacion-eta}
\end{equation}
Para expresar la composición manteniendo visibles sus operaciones, pongamos
$x=50\pi\alpha/9$, $y=\pi(\eta_{\mathrm{ret}}+\alpha)/90$ y
$\mathcal R(q)=12S_{90}(q)-S_{120}(q)$. Entonces
\begin{equation}
 \gamma=\frac{100\pi}{9}\alpha(\eta_{\mathrm{ret}}+\alpha)
 +\mathcal R(e^{-x+y})-\mathcal R(e^{-x-y}).
 \label{md:compos:barbero-eta}
\end{equation}
La sustitución de \eqref{md:compos:eliminacion-eta} en esta identidad produce
una relación explícita entre la lectura electrónica completa y Barbero,
conservando $\alpha$, $H_5$, $\Delta_4$ y las incidencias. Es una consecuencia
de compatibilidad de la construcción. La rama física angular utilizada
mantiene $0<y<x$; la operación no introduce una masa metrológica como dato
generador.

En la extensión con etiqueta de orientación $\epsilon=\pm1$,
$C^*_{\deg,\epsilon}=2\epsilon(\eta_{\mathrm{ret}}+\alpha)$, el intercambio
$\epsilon\mapsto-\epsilon$ deja fijo $\eta_{\mathrm{ret}}$, conserva la
expresión electrónica y permuta $q_+$ y $q_-$. El funcional de Barbero cambia
de signo. Esta covariancia distingue una lectura escalar de acción de una
evaluación orientada. La identificación con conjugación de carga conserva
además $C_M$ y la representación electromagnética de hojas; el signo angular
aislado no sustituye esos operadores.

\subsection{Momento orientado de Catalán y respuesta constitutiva}
\label{md:compos:catalan}

El ciclo dodecafásico realiza un plano orientado de cuarto de giro. Su
coeficiente resolvente es $k_4(s)=s/(1+s^2)$. El momento de profundidad dos
queda fijado por
\begin{equation}
 \mathcal C_2(s)=\sum_{k\geq0}
 \frac{(-1)^ks^{2k+1}}{(2k+1)^2},\qquad
 \mathcal D=s\frac{d}{ds},\qquad
 \mathcal D^2\mathcal C_2(s)=k_4(s),\qquad
 \mathcal C_2(1)=G_{\mathrm{Cat}}.
 \label{md:compos:momento-catalan}
\end{equation}
La identidad diferencial se utiliza para $0<s<1$; el valor extremo se
entiende con las condiciones de convergencia y restitución del momento ya
establecidas. La definición intrínseca de Catalán es el valor extremo de ese
momento orientado de profundidad. Su colección conserva más información que
el valor extremo. La respuesta constitutiva utiliza dicha colección:
\begin{equation}
 f(s)=\frac{1+s+s^2}{s(1+s)}\mathcal D^2\mathcal C_2(s)
 =\frac{1-s^3}{1-s^4},\qquad
 r_\pm=f(s_\pm),\qquad s_\pm=q_\pm^{30}.
 \label{md:compos:catalan-respuesta}
\end{equation}
Las potencias $3$ y $4$ corresponden a $90=3\cdot30$ y $120=4\cdot30$.
\begin{proposition}[Restitución funcional con condición de frontera]
\label{md:compos:inversa-catalan}
La colección constitutiva $f$ determina el momento orientado mediante
\[
 \mathcal C_2(s)=\int_0^s\log\frac{s}{t}\,
 \frac{1+t}{1+t+t^2}f(t)\,dt,\qquad 0<s<1.
\]
Es la única solución de
$\mathcal D^2\mathcal C_2(s)=s(1+s)f(s)/(1+s+s^2)$
que tiende a cero en el origen.
\end{proposition}
\begin{proof}
Para la colección construida,
$(1+t)f(t)/(1+t+t^2)=1/(1+t^2)$. La integral converge, y su
derivada logarítmica es
$\mathcal D\mathcal C_2(s)=\int_0^s(1+t^2)^{-1}\,dt$.
Aplicar de nuevo $\mathcal D$ da $s/(1+s^2)$. La integral es
$O(s)$ en el origen. Dos soluciones difieren en $a+b\log s$;
el límite nulo obliga $a=b=0$.
Para $s<1$, integrar la serie geométrica en compactos y usar
$\int_0^s t^{2k}\log(s/t)\,dt=s^{2k+1}/(2k+1)^2$
recupera la serie de \eqref{md:compos:momento-catalan}.
Su convergencia absoluta dominada permite $s\uparrow1$ y restituye Catalán.
\end{proof}
La condición de frontera fija las dos constantes de integración.
La restitución utiliza la función constitutiva con su dominio; el
valor extremo de Catalán y los dos valores de canal son lecturas de
esa función ya construida. En el extremo, la identidad doblemente
diferenciada conserva su interpretación por límite de Abel.

El operador constitutivo $V=r_+P_++r_-P_-$ determina
\begin{equation}
 \widehat\varepsilon=r_+^2,\qquad \widehat\mu=r_-^2,
 \qquad \widehat Z=\frac{r_-}{r_+},\qquad
 \widehat c=(r_+r_-)^{-1}.
 \label{md:compos:lectores-constitutivos}
\end{equation}
Con $\psi=f^{-1}$ y
\begin{equation}
 \mathcal H_B(s)=\frac{12s^3}{1-s^9}
 -\frac{s^4}{1-s^{12}}-\frac{(\log s)^2}{20\pi},
 \label{md:compos:h-barbero}
\end{equation}
la factorización construida toma la forma
\begin{equation}
 \gamma=\mathcal H_B(\psi(r_-))-\mathcal H_B(\psi(r_+)).
 \label{md:compos:factorizacion-barbero}
\end{equation}
La parte logarítmica restituye exactamente la contribución angular de
Barbero, conservada en grados o transportada íntegramente a radianes.

\subsection{Restitución mediante velocidad normalizada y Barbero}
\label{md:compos:difeomorfismo}

\paragraph{Proposición.}
En la colección constitutiva normalizada de dos canales etiquetados, el mapa
\begin{equation}
 (r_+,r_-)\longmapsto(\widehat c,\gamma),\qquad
 (3/4,1)^2\longrightarrow(1,16/9)\times\mathbb R
 \label{md:compos:mapa-global}
\end{equation}
es un difeomorfismo. Su inversa recupera los dos autovalores etiquetados.
Los proyectores de hoja pertenecen a la representación de partida: dos
escalares no recuperan una base espacial arbitraria.

\paragraph{Demostración.}
La expresión regular
$f(s)=(1+s+s^2)/(1+s+s^2+s^3)$ da
\begin{equation}
 f'(s)=-\frac{s^2(s^2+2s+3)}{(1+s+s^2+s^3)^2}<0.
 \label{md:compos:derivada-f}
\end{equation}
Así, $f$ es analítica y estrictamente decreciente, con $f(0^+)=1$ y
$f(1^-)=3/4$. Para $0<s<1$,
\begin{equation}
 \mathcal H_B'(s)=
 \frac{36s^2(1+2s^9)}{(1-s^9)^2}
 -\frac{4s^3(1+2s^{12})}{(1-s^{12})^2}
 -\frac{\log s}{10\pi s}>0.
 \label{md:compos:derivada-h}
\end{equation}
En efecto, el cociente del primer término positivo entre el segundo término
positivo antes de restarlo es
\begin{equation}
 \frac{9}{s}\frac{1+2s^9}{1+2s^{12}}
 \left(\frac{1-s^{12}}{1-s^9}\right)^2>9.
 \label{md:compos:cota-derivada-h}
\end{equation}
El último sumando de \eqref{md:compos:derivada-h} también es positivo.
Por tanto, $F_B=\mathcal H_B\circ\psi$ tiene derivada estrictamente negativa,
$F_B(3/4^+)=+\infty$ y $F_B(1^-)=-\infty$.

Fijado $c=\widehat c$, escribimos $z=\log\widehat Z$ y
$r_\pm=c^{-1/2}\exp(\mp z/2)$. El intervalo admisible es
\begin{equation}
 |z|<a(c):=\min\left\{\log c,\log\frac{16}{9c}\right\}.
 \label{md:compos:dominio-z}
\end{equation}
La función $\Gamma_c(z)=F_B(r_-)-F_B(r_+)$ satisface
\begin{equation}
 \Gamma_c'(z)=\frac12
 \bigl[r_-F_B'(r_-)+r_+F_B'(r_+)\bigr]<0.
 \label{md:compos:derivada-gamma}
\end{equation}
Cuando $z$ aumenta hasta $a(c)$, o $r_-$ alcanza $1$, o $r_+$ alcanza $3/4$,
o suceden ambas cosas; $\Gamma_c$ tiende a $-\infty$. En el extremo opuesto
tiende a $+\infty$. Existe, por tanto, una única solución para cada $\gamma$
real. La derivada no nula y el teorema de la función implícita proporcionan
una inversa suave global. Esto prueba el enunciado. En la orientación
original $y>0$, se tiene $z<0$ y $\gamma>0$.

La velocidad normalizada conserva el producto de las respuestas; Barbero
distingue de forma monótona la distribución orientada entre ellas. Juntas
restituyen la información espectral que cada lectura aislada no determina.

El enunciado concierne a la colección de lectores constitutivos. La imagen
efectivamente producida por el generador HMT completo es un subconjunto de
esa colección: restringir a dicha imagen conserva la inyectividad y el orden
causal. La normalización y las marcas de hoja son datos necesarios de la
restitución. La prueba no introduce valores CODATA para elegir canales ni
afirma que cada punto sintético de la colección sea un estado generado.

\subsection{Restitución sobre la imagen generada}
\label{md:compos:restitucion-completa}

Una vez recuperados $r_\pm$ desde $(\widehat c,\gamma)$, se obtiene
$s_\pm=\psi(r_\pm)$. Entonces
\begin{equation}
 A_{\deg}=-\frac{3}{\pi}\log(s_+s_-),\qquad
 C^*_{\deg}=\frac{3}{\pi}\log\frac{s_-}{s_+}.
 \label{md:compos:restitucion-angular}
\end{equation}
Sobre la imagen generada y en su orientación positiva,
\begin{equation}
 \alpha=\frac{A_{\deg}}{1000},\qquad
 \eta_{\mathrm{ret}}=\frac{C^*_{\deg}}{2}-\alpha>0.
 \label{md:compos:restitucion-alpha-eta}
\end{equation}
Con las coordenadas regionales $\pi,\varphi,e$ y las secciones dimensionales
ya construidas, esta restitución devuelve $H_5$, el factor
$R_{\mathrm{act}}=H_5/\eta_{\mathrm{ret}}$ y la lectura electrónica completa
de la sección~\ref{md:compos:electron}. También devuelve la sección de
Planck. Si se conservan además la base posicional de doce bloques, el origen
y el acarreo, la identidad
\begin{equation}
 \kappa_{\mathrm{ag}}=\pi+e-\varphi-4-\alpha
 \label{md:compos:restitucion-k}
\end{equation}
permite la restitución del registro $K$ descrita por su lector. Se trata de
un recorrido de reconstrucción; su sentido inverso no constituye una nueva
flecha generadora desde valores físicos externos.

La gravedad y la conversión térmica conservan las secciones del mismo
estado. El teorema~\ref{md:rigidez:teorema} y
\eqref{md:rigidez:reloj} determinan primero $L_\alpha$ y su
representación circular $R_{108}=L_\alpha$:
\begin{equation}
 G_a=\frac{c_{\mathrm{int}}^3R_{108}^2}{\hbar_a},\qquad
 \frac{G_a\hbar_a}{c_{\mathrm{int}}^3}=R_{108}^2,
 \qquad
 k_B\Theta_{\mathrm{clk},a}\log3=\frac{h_a}{108t_0}.
 \label{md:compos:gravedad-termica}
\end{equation}
El radio cronológico, la velocidad física y la base térmica conservan sus
constructores. En particular, el reloj de la representación circular queda
ligado a $L_\alpha$ y $c_{\mathrm{int}}$. La primera
composición explicita la reciprocidad acción--gravedad; la segunda reúne
acción, duración e información trítica, con individualización de temperatura
y $k_B$ en sus bases declaradas.

\subsection{Definiciones estructurales de las lecturas compuestas}
\label{md:compos:definiciones}

\begin{description}
 \item[$\alpha$.] Coordenada escalar de clausura dodecafásica compatible con
 las regiones y con el registro aritmético--geométrico $K$; transmite esa
 compatibilidad al sector angular y de acción.
 \item[Contraángulo.] Coordenada angular de la clausura bidireccional que
 combina el coeficiente de la sección de acción de retorno con $\alpha$,
 en la coordenatización de $360$ unidades.
 \item[Planck reducida.] Sección de acción cuyo coeficiente regional y escala
 determinantal se construyen antes del contraángulo; en la elipse, medida de
 área orientada por unidad de fase paramétrica.
 \item[Electrón.] Sección central persistente; su lectura escalar compone
 incompatibilidad aditiva--multiplicativa, transferencia regional, defecto
 de clausura y retorno de acción. La ecuación completa conserva la sección,
 además de su valor de masa.
 \item[Barbero--Immirzi.] Funcional orientado del transporte angular y de los
 retornos de la incidencia hexada--octada; en la colección constitutiva,
 coordenada que junto con la velocidad normalizada restituye el espectro
 etiquetado.
 \item[Catalán.] Valor extremo del momento orientado de profundidad dos del
 cuarto de giro dodecafásico. Su colección se enlaza diferencialmente con la
 respuesta del vacío.
 \item[$G$.] Coeficiente común de la realización radial que relaciona
 energía y longitud conjugada conservando simultáneamente
 $H\Lambda=\hbar cI$ y $R\Lambda=L_\alpha^2I$.
 La composición longitudinal y polar determina su normalización;
 su expresión dimensional es $c^3L_\alpha^2/\hbar$.
 \item[Boltzmann.] Transductor entre la recta térmica y la energía de
 decisión, con normalización informacional trítica y base térmica declaradas.
\end{description}
Estas formulaciones intrínsecas se basan en los objetos y ecuaciones
construidos. Su alcance es el de esta realización; las identificaciones
físicas conservan sus demostraciones correspondientes y no se deducen de
la sola denominación estructural.

\subsection{Transporte del bloque electrónico nilpotente}
\label{md:compos:nilpotentes}

\paragraph{Proposición.}
Sean
\begin{equation}
 S=\begin{pmatrix}1&0\\0&-1\end{pmatrix},\qquad
 J=\begin{pmatrix}0&1\\-1&0\end{pmatrix},\qquad
 n_\pm=\frac{S\pm J}{2},\qquad
 D=\operatorname{diag}(r,r^{-1}),\quad r>0.
 \label{md:compos:bloque-nilpotente}
\end{equation}
Defínanse $n_{\pm,D}=Dn_\pm D^{-1}$ y $g=D^{-T}D^{-1}$. Entonces
\begin{equation}
 n_{\pm,D}^2=0,\qquad
 \{n_{+,D},n_{-,D}\}=I,\qquad
 (n_{+,D})^{\dagger_g}=n_{-,D},
 \label{md:compos:identidades-nilpotentes}
\end{equation}
donde $A^{\dagger_g}=g^{-1}A^Tg$. La elipse transportada desde
$Y^TY=2\hbar$ mediante $X=DY$ satisface $X^TgX=2\hbar$.

\paragraph{Demostración.}
El cálculo matricial directo da
$S^2=I$, $J^2=-I$ y $SJ+JS=0$. Por ello,
\[
 n_\pm^2=\tfrac14(S^2+J^2\pm SJ\pm JS)=0,
 \qquad
 n_+n_-+n_-n_+=\tfrac12(S^2-J^2)=I.
\]
Además, $S^T=S$ y $J^T=-J$ implican $n_+^T=n_-$. La conjugación por $D$
preserva productos, cuadrados y la identidad. Para el adjunto métrico,
\[
 g^{-1}(Dn_+D^{-1})^Tg
 =DD^T D^{-T}n_-D^T D^{-T}D^{-1}
 =Dn_-D^{-1}=n_{-,D}.
\]
Finalmente, $D^TgD=I$, de donde
$X^TgX=Y^TD^TgDY=Y^TY=2\hbar$. Quedan probadas las afirmaciones.

La comprobación algebraica exacta tiene lugar en
$\mathbb Q[r,r^{-1}]$; la restricción $r>0$ fija la realización elíptica
orientada utilizada. Esta composición conserva las funciones particulares
de los distintos factores $1/2$ del desarrollo: una coincidencia numérica
entre ellos no constituye una identificación de sus operadores.

% Procedencia íntegra: HIBRIDACION_ACCION_MASA_PROPAGACION.md, §§2--7.
% La gestión, el historial y los localizadores se conservan en gestion/.
\section{Acción, masa y propagación: compatibilidad estructural y transporte}
\label{md:hib:seccion}

\subsection{Sección de acción compartida y ley de masas}
\label{md:hib:accion-masas}

Sean
\begin{equation}
 s_0=10^{-34}\mathcal U_S,\qquad
 \hbar_{\mathrm{ret}}=\eta_{\mathrm{ret}}s_0,\qquad
 R_{\mathrm{act}}=\frac{H_5}{\eta_{\mathrm{ret}}}.
 \label{md:hib:accion}
\end{equation}
Se mantienen las coordenadas construidas en grados:
\begin{equation}
 A_{\deg}=1000\alpha,\qquad
 C^*_{\deg}=2(\eta_{\mathrm{ret}}+\alpha).
 \label{md:hib:angulos}
\end{equation}
Sólo después se transportan a radianes mediante
$x=\pi A_{\deg}/180$ e $y=\pi C^*_{\deg}/180$. Los canales son
$q_\pm=\exp(-x\mp y)$, con $x>|y|$. La forma elíptica se expresa por
\begin{equation}
 \zeta=\operatorname{artanh}(y/x).
 \label{md:hib:forma}
\end{equation}
La composición utiliza primero el coeficiente adimensional
$\eta_{\mathrm{ret}}$ y conserva así los tipos de acción y de ángulo.
En esta notación, $\eta_{\mathrm{ret}}$ es el coeficiente de acción,
$\eta_h(\Gamma)$ pertenece a la torre de una ruta y $\zeta$ expresa la
deformación de la elipse; sus funciones se mantienen separadas.

Se trabaja con $t_0,c_{\mathrm{int}},\eta_{\mathrm{ret}}>0$, carácter
convergente y factor $\Xi_\Gamma>0$. El lector másico construido determina
\begin{equation}
 \omega_0=\frac{2\pi}{108t_0},\qquad
 E_{\mathrm{clk}}=\hbar_{\mathrm{ret}}\omega_0,\qquad
 m_{\mathrm{clk}}=\frac{E_{\mathrm{clk}}}{c_{\mathrm{int}}^2}.
 \label{md:hib:reloj}
\end{equation}
Para una ruta con realización escalar,
\begin{equation}
 m_\Gamma^\beta=m_{\mathrm{clk}}\Xi_\Gamma,\qquad
 \Xi_\Gamma=b_e\,
 \chi_{\mathrm{ref}}(\sigma_\Gamma-\sigma_e,\eta_\Gamma-\eta_e)
 R_{\mathrm{act}}^{\kappa_\Gamma}.
 \label{md:hib:factor-ruta}
\end{equation}
$\Xi_\Gamma$ abrevia el factor estructural de ruta. Se conserva expresamente
\begin{equation}
 b_e=\frac{\varepsilon_e^{(0)}\mathcal U_E}{E_{\mathrm{clk}}}.
 \label{md:hib:normalizacion-electron}
\end{equation}
El estado central conserva su normalización propia respecto del cuanto
cronológico:
\begin{equation}
 \varepsilon_e^{(0)}=\frac{\sqrt3}{4}
 \left[\exp\!\left(\frac{\varphi}{\pi^2}\right)-22\alpha^3\right]
 (1+15\Delta_4),\qquad
 \varepsilon_e^\beta=\varepsilon_e^{(0)}
 R_{\mathrm{act}}^{80-54\alpha+6\Delta_4}.
 \label{md:hib:electron}
\end{equation}
El carácter refinado incluye los términos
\begin{equation}
 n_Ax+\frac{s_C}{6}y+k_{\Delta_4}\Delta_4+b_\zeta\zeta_{270}
 +\sum_h N_h(q_-^h-q_+^h),\qquad
 N_h=\eta_h+b_{120}\mathbf 1_{\{h=120\}},
 \label{md:hib:caracter}
\end{equation}
con $h$ en el semigrupo generado por $90$ y $120$. El término de altura
$120$ se incluye una vez. Los coeficientes primitivos de retorno se
construyen mediante
\begin{equation}
 \begin{aligned}
 a_n(\Gamma)&=\operatorname{Tr}(R_\Gamma U_\Gamma^n),\\
 p_n(\Gamma)&=\frac1n\sum_{d\mid n}\mu_{\mathrm{Mob}}(d)
                         a_{n/d}(\Gamma),\\
 \eta_h(\Gamma)&=p_h(\Gamma)-p_h(\Gamma_e),
 \end{aligned}
 \label{md:hib:retornos-primitivos}
\end{equation}
donde $\mu_{\mathrm{Mob}}$ es la función de Möbius aritmética.
La acción participa en la escala absoluta, en la razón $R_{\mathrm{act}}$
y en $y$, que determina los canales. En consecuencia, la cancelación de
una escala dimensional conserva las restantes dependencias estructurales.
Al variar el reloj, $b_e=\varepsilon_e^{(0)}\mathcal U_E/E_{\mathrm{clk}}$
varía con él: conservar la lectura electrónica exige transportar el
carácter completo, en vez de fijar por separado ese cociente.
En la realización longitudinal integrada se utiliza
$t_0=\pi_{\rm HMT}L_\alpha/(54c_{\mathrm{int}})$ de
\eqref{md:rigidez:reloj}.

\subsubsection{Frecuencia y longitud de cada sección masiva}
\label{md:hib:frecuencia-longitud}

Definiendo $E_\Gamma=m_\Gamma c_{\mathrm{int}}^2$ como energía de reposo,
$\Omega_\Gamma=E_\Gamma/\hbar_{\mathrm{ret}}$,
$\nu_\Gamma=\Omega_\Gamma/(2\pi)$ y
$R_{108}=c_{\mathrm{int}}/\omega_0$, se obtiene
\begin{equation}
 \begin{aligned}
 E_\Gamma&=\hbar_{\mathrm{ret}}\omega_0\Xi_\Gamma,&
 \Omega_\Gamma&=\omega_0\Xi_\Gamma,\\
 \nu_\Gamma&=\frac{\Xi_\Gamma}{108t_0},&
 \bar\lambda_\Gamma&=
 \frac{\hbar_{\mathrm{ret}}}{m_\Gamma c_{\mathrm{int}}}
 =\frac{R_{108}}{\Xi_\Gamma}.
 \end{aligned}
 \label{md:hib:frecuencias}
\end{equation}
\begin{proof}
Se sustituye $m_\Gamma=m_{\mathrm{clk}}\Xi_\Gamma$ y se efectúan las
cancelaciones en una misma coordenatización dimensional. Para rutas positivas $X,Y$,
la propiedad de carácter y la normalización común dan
\begin{equation}
 \frac{\Omega_Y}{\Omega_X}=\frac{m_Y}{m_X}
 =\chi_{\mathrm{ref}}(\sigma_Y-\sigma_X,\eta_Y-\eta_X)
 R_{\mathrm{act}}^{\kappa_Y-\kappa_X}.
 \label{md:hib:cocientes}
\end{equation}
\end{proof}
La cancelación es dimensional; permanecen los canales y la memoria
relativa. El sector nulo queda fuera del dominio del inverso de masa.
En una fibra no escalar se conserva el operador o su medida espectral,
en lugar de escoger un representante puntual arbitrario.

\subsubsection{Frecuencia levantada y compatibilidad de normalizaciones}
\label{md:hib:tres-frecuencias}

La construcción cronológica también produce la frecuencia levantada de
holonomía
\begin{equation}
 \omega_j=\frac{\theta_j+2\pi w_j}{108t_0},\qquad
 \lambda_j=r_j\exp(i\theta_j),
 \label{md:hib:frecuencia-levantada}
\end{equation}
donde $w_j$ conserva el devanado de la historia. Su lectura energética es
\begin{equation}
 E_j=\hbar_{\mathrm{ret}}\omega_j
 \chi_{\mathrm{full}}R_{\mathrm{act}}^{\kappa_j}.
 \label{md:hib:energia-levantada}
\end{equation}
Esta frecuencia, $\omega_0$ y $\Omega_\Gamma$, corresponden a lectores
distintos. Si ambas expresiones energéticas representan el mismo modo y
el mismo exponente, su compatibilidad exacta exige
\begin{equation}
 \frac{\omega_j}{\omega_0}\chi_{\mathrm{full}}
 =b_e\chi_{\mathrm{ref,rel}}.
 \label{md:hib:compatibilidad-normalizaciones}
\end{equation}
Las dos normalizaciones se identifican mediante esta igualdad y no se
multiplican entre sí. La fase reducida $\theta_j$ por sí sola no recupera
$w_j$. Por ello, un retorno de fase conserva una distinción con la historia
cronológica completa. El devanado conserva su constructor; esta distinción
no autoriza elegirlo libremente ni producir energía sin una dinámica.

\subsection{Reciprocidad espacial y composición gravitatoria}
\label{md:hib:gravedad}

La sección gravitatoria circular expresa el resultado del
teorema~\ref{md:rigidez:teorema}:
$G_a=c_{\mathrm{int}}^3L_\alpha^2/\hbar_a$ y $R_{108}=L_\alpha$.
La longitud procede del retorno inscrito y precede a esta
identificación circular. Manteniendo
$\hbar_a=\hbar_{\mathrm{ret}}$ y la convención reducida
$r_g=Gm/c_{\mathrm{int}}^2$, la masa anterior determina
\begin{equation}
 \boxed{r_{g,\Gamma}=R_{108}\Xi_\Gamma,\qquad
 \bar\lambda_\Gamma=\frac{R_{108}}{\Xi_\Gamma},\qquad
 r_{g,\Gamma}\bar\lambda_\Gamma=R_{108}^2.}
 \label{md:hib:reciprocidad-radial}
\end{equation}
\begin{proof}
La sustitución de $G_a$ y de $m_\Gamma$ da
\[
 \frac{G_am_\Gamma}{c_{\mathrm{int}}^2}
 =\frac{c_{\mathrm{int}}^3R_{108}^2}{\hbar_{\mathrm{ret}}}
 \frac{\hbar_{\mathrm{ret}}\omega_0\Xi_\Gamma}
 {c_{\mathrm{int}}^4}
 =R_{108}\Xi_\Gamma.
\]
La segunda igualdad procede de \eqref{md:hib:frecuencias}; su producto
prueba la tercera.
\end{proof}
La misma coordenada de ruta aumenta una lectura radial y contrae la otra.
En logaritmos,
\begin{equation}
 \log\frac{r_{g,\Gamma}}{R_{108}}=\log\Xi_\Gamma,
 \qquad
 \log\frac{\bar\lambda_\Gamma}{R_{108}}=-\log\Xi_\Gamma.
 \label{md:hib:radios-logaritmicos}
\end{equation}
El centro geométrico de ambas lecturas queda fijado por el área $R_{108}^2$.

En la versión operatoria acotada, sea $X$ autoadjunto y acotado, con
$X\geq\delta I$ para algún $\delta>0$, en una fibra masiva. Los operadores
\begin{equation}
 M=m_{\mathrm{clk}}X,\qquad
 \Lambda=R_{108}X^{-1},\qquad R_g=R_{108}X
 \label{md:hib:operadores-radiales}
\end{equation}
satisfacen $R_g\Lambda=\Lambda R_g=R_{108}^2I$ en todo el espacio.
La afirmación utiliza el cálculo funcional del mismo operador. El núcleo
nulo queda fuera del inverso. El enunciado presente conserva el régimen
acotado uniformemente positivo; la extensión no acotada exige demostrar
por separado los dominios, los cortes compatibles y el paso al límite.
La sección~\ref{md:rigidez:variaciones} demuestra además la conservación
de ambas identidades bajo variaciones no conmutativas y reducción
compatible de memoria.

La interpretación propuesta caracteriza el factor másico relativo como
parámetro de dilatación recíproca de estas dos lecturas. La condición circular
pertenece al resultado. Las longitudes así relacionadas no se identifican
por esta sola igualdad con radios materiales medidos ni con un horizonte
electrónico. La identidad radial aislada se utiliza por su función en la
composición, sin atribuirle prioridad histórica independiente.

\subsection{Respuesta constitutiva, orientación y magnitud}
\label{md:hib:constitutivas}

Con $T=\exp(-xI-yR)$, $R=P_+-P_-$ y proyectores de hoja conservados,
el lector temporal es
\begin{equation}
 \begin{aligned}
 V&=(I-T^{90})(I-T^{120})^{-1}=r_+P_++r_-P_-,\\
 r_\pm&=f(s_\pm),\qquad
 s_\pm=\exp[-30(x\pm y)],\qquad
 f(s)=\frac{1-s^3}{1-s^4}.
 \end{aligned}
 \label{md:hib:operador-constitutivo}
\end{equation}
Se recuperan $\widehat\varepsilon=r_+^2$, $\widehat\mu=r_-^2$,
$\widehat Z=r_-/r_+$ y $\widehat c=(r_+r_-)^{-1}$.
El enlace constitutivo de velocidad establece
$c_{\mathrm{int}}=c_{\mathrm{vac}}=\widehat c\,u_C$ en la misma coordenatización.
La colección orientada de Catalán participa en $f$ mediante su momento
resolvente, conservando la estructura anterior al valor extremo.

La descomposición exacta
\begin{equation}
 V=\widehat c^{-1/2}
 \exp\!\left[-\frac12(\log\widehat Z)R\right]
 \label{md:hib:descomposicion-v}
\end{equation}
distingue dilatación común y reparto entre hojas. El intercambio de hojas
conserva $\widehat c$ e invierte $\widehat Z$; el funcional de Barbero cambia
de signo. La sección~\ref{md:compos:difeomorfismo} demuestra que
$(\widehat c,\gamma)$ recupera $r_+,r_-$ en la colección normalizada.
Su restricción a la imagen generada conserva esta recuperabilidad; el
enunciado de inversión no declara generada toda pareja sintética.

\subsubsection{Convexidad y evaluación de los canales separados}
\label{md:hib:convexidad}

\begin{proposition}
\label{md:hib:proposicion-convexidad}
Para $x>0$ fijo y $|y|<x$, la función
$y\mapsto\log\widehat c(x,y)$ es par y estrictamente convexa.
En consecuencia,
\begin{equation}
 \widehat c(x,y)\geq\widehat c(x,0)
 =\frac{1}{f(e^{-30x})^2},
 \label{md:hib:cota-velocidad}
\end{equation}
con igualdad exactamente para $y=0$.
\end{proposition}
\begin{proof}
Sea $g(u)=\log f(e^{-u})$, para $u>0$. Se tiene
\begin{equation}
 g'(u)=\frac{3}{e^{3u}-1}-\frac{4}{e^{4u}-1}>0,
 \label{md:hib:derivada-g}
\end{equation}
\begin{equation}
 g''(u)=-\frac{9}{4\sinh^2(3u/2)}
 +\frac{16}{4\sinh^2(2u)}<0.
 \label{md:hib:segunda-g}
\end{equation}
La primera desigualdad procede del decrecimiento de
$a/(\exp(au)-1)$ en $a>0$. Para la segunda, la función
$a/\sinh(au/2)$ decrece estrictamente: su derivada tiene el signo de
$\sinh v-v\cosh v$, negativo para $v>0$. Como
\[
 \log\widehat c(x,y)=-g(30(x+y))-g(30(x-y)),
\]
su segunda derivada respecto de $y$ es
\begin{equation}
 -900\bigl[g''(30(x+y))+g''(30(x-y))\bigr]>0.
 \label{md:hib:segunda-velocidad}
\end{equation}
La paridad fija el mínimo único en cero.
\end{proof}
También se obtiene
\begin{equation}
 \begin{aligned}
 \log\widehat c(x,y)
 &=-2g(30x)-900g''(30x)y^2+O(y^4),\\
 \log\widehat Z(x,y)
 &=-60g'(30x)y+O(y^3).
 \end{aligned}
 \label{md:hib:desarrollos-canales}
\end{equation}
La velocidad pierde el signo de orientación, pero conserva una respuesta
cuadrática a la separación de canales. La impedancia conserva información
orientada desde el primer orden. Promediar $x+y$ y $x-y$ antes de aplicar el
lector elimina aquella contribución. La afirmación es una consecuencia
demostrada de la respuesta constitutiva, distinta de una medición nueva.

Sobre la rama positiva de acción,
$\eta_{\mathrm{ret}}=(90/\pi)(y-x/500)$; por tanto,
\begin{equation}
 \hbar_{\mathrm{ret}}=s_0\frac{90}{\pi}(y-x/500),\qquad
 E_k=s_0\frac{90u_C|k|}{\pi}(y-x/500)\widehat c(x,y).
 \label{md:hib:energia-angular}
\end{equation}
Para $x$ fijo, $|k|>0$ y $x/500<y<x$, esta última función crece
estrictamente en $y$. Es una propiedad de la colección de realización
radiativa homogénea. Su aplicación a una trayectoria generada conserva el
constructor de $y$; la variación de un parámetro de esta colección no equivale
por sí sola a una variación experimental de una constante universal.
Bajo inversión de orientación se transporta también la etiqueta de acción,
manteniendo la rama positiva correspondiente.

\subsection{Acción, reloj e información térmica}
\label{md:hib:termica}

La sección térmica conserva
$k_B\Theta_{108}\log3=E_{\mathrm{clk}}$, con
$\Theta_{108}=\Theta_{\mathrm{clk,ret}}$ en la sección de retorno.
Definimos
$\beta_{108}=(k_B\Theta_{108})^{-1}$. Para
$E_\Gamma=E_{\mathrm{clk}}\Xi_\Gamma$ se deduce exactamente
\begin{equation}
 \beta_{108}E_\Gamma=(\log3)\Xi_\Gamma,
 \qquad
 \boxed{\exp(-\beta_{108}E_\Gamma)=3^{-\Xi_\Gamma}.}
 \label{md:hib:peso-termico}
\end{equation}
La expresión es un peso de Gibbs no normalizado. Para una colección finita
de estados, o una colección con suma de partición convergente, y
multiplicidades $g_\Gamma$, se obtiene
\begin{equation}
 Z_{\mathrm{part}}=\sum_\Gamma g_\Gamma3^{-\Xi_\Gamma},\qquad
 p_\Gamma=\frac{g_\Gamma3^{-\Xi_\Gamma}}{Z_{\mathrm{part}}},\qquad
 \frac{p_Y/g_Y}{p_X/g_X}=3^{-(\Xi_Y-\Xi_X)}.
 \label{md:hib:particion}
\end{equation}
Se asume una realización canónica a $\Theta_{108}$ y un origen energético
común. Cuando interviene un potencial químico, se incorpora a la energía
correspondiente. El índice $\Gamma$ enumera estados o grupos físicos
distintos; varias historias equivalentes no producen automáticamente una
multiplicidad. Las entropías adimensionales son
\begin{equation}
 \begin{aligned}
 H_{\mathrm{micro}}
 &=\log Z_{\mathrm{part}}+(\log3)\sum_\Gamma p_\Gamma\Xi_\Gamma,\\
 H_{\mathrm{grupos}}
 &=H_{\mathrm{micro}}-\sum_\Gamma p_\Gamma\log g_\Gamma.
 \end{aligned}
 \label{md:hib:entropias}
\end{equation}
Para afirmar la finitud de la entropía se requiere, además,
\begin{equation}
 \sum_\Gamma g_\Gamma\Xi_\Gamma3^{-\Xi_\Gamma}<\infty.
 \label{md:hib:condicion-entropia}
\end{equation}
La finitud de $Z_{\mathrm{part}}$ garantiza por sí sola la normalización.
En colecciones finitas se cumplen ambas condiciones.

La fórmula enlaza la estructura másica con frecuencia propia y ponderación
térmica en el mismo marco. Conserva la distinción entre energía y calor,
entre coordenada real $\Xi_\Gamma$ y conteo entero de trits, y entre una
realización a $\Theta_{108}$ y una afirmación sobre la temperatura del
universo. Para ocupaciones bosónicas o fermiónicas libres por modo, con
potencial químico nulo, la misma sustitución produce
\begin{equation}
 \overline n_\Gamma=\frac{1}{3^{\Xi_\Gamma}\mp1}.
 \label{md:hib:ocupaciones}
\end{equation}
El signo menos corresponde a Bose y el signo más a Fermi. El total del
grupo degenerado es $g_\Gamma\overline n_\Gamma$. En Bose se exige
$\Xi_\Gamma>0$ y el modo nulo se trata separadamente. Son realizaciones
estadísticas posteriores aplicadas a un espectro ya construido.

\subsubsection{Compatibilidad conjunta de cuatro lecturas}
\label{md:hib:cuatro-lecturas}

Componiendo las secciones anteriores, bajo sus mismas hipótesis,
\begin{equation}
 \boxed{\Xi_\Gamma=\frac{\Omega_\Gamma}{\omega_0}
 =\frac{r_{g,\Gamma}}{R_{108}}
 =\frac{R_{108}}{\bar\lambda_\Gamma}
 =\frac{E_\Gamma}{k_B\Theta_{108}\log3}.}
 \label{md:hib:compatibilidad-cuatro}
\end{equation}
En la colección circular, $\Xi=1$ identifica la sección de frecuencia
$\omega_0$, radios coincidentes $R_{108}$ y energía $E_{\mathrm{clk}}$.
Su peso térmico no normalizado es $1/3$. Es un punto de compatibilidad de
las lecturas. La existencia de una especie del atlas con $\Xi=1$, o de
una segunda especie generada para cada sección recíproca $1/\Xi$, requiere
su pertenencia producida al atlas; no se infiere de esta igualdad.

\subsection{Identidad energética de la elipse y transporte exacto}
\label{md:hib:elipse}

En la elipse de acción,
\begin{equation}
 a_S=\sqrt{2\hbar}\,e^{\zeta/2},\qquad
 b_S=\sqrt{2\hbar}\,e^{-\zeta/2},\qquad
 Q=a_S\cos\theta,\qquad P=b_S\sin\theta,
 \label{md:hib:elipse-coordenadas}
\end{equation}
con $\zeta=\operatorname{artanh}(C^*_{\deg}/A_{\deg})$, se tiene
$a_Sb_S=2\hbar$, área orientada
$\mathcal A(\theta)=\hbar\theta$ y $\mathcal A(2\pi)=h$.
La fase paramétrica $\theta$ y el ángulo polar de un punto de la elipse
conservan sus significados geométricos respectivos.

La realización LC declarada utiliza
\begin{equation}
 L_\zeta=\frac{Z_*}{\omega}e^{-\zeta},\qquad
 C_\zeta=\frac{1}{Z_*\omega}e^\zeta.
 \label{md:hib:lc}
\end{equation}
En coordenadas de acción $Q=\sqrt{Z_*}\,q$ y
$P=\Phi/\sqrt{Z_*}$,
\begin{equation}
 H_0=\frac{\omega}{2}
 \left(e^{-\zeta}Q^2+e^\zeta P^2\right),\qquad
 U_E=\hbar\omega\cos^2\theta,\qquad
 U_B=\hbar\omega\sin^2\theta.
 \label{md:hib:energia-elipse}
\end{equation}
Se obtiene $H_0=\hbar\omega$ por composición de las dos contribuciones.
El reloj y el transductor pertenecen a la realización declarada; la sección
$\hbar$ conserva su construcción anterior al circuito. La impedancia
$Z_{\mathrm{LC}}/Z_*=e^{-\zeta}$ y la impedancia constitutiva del vacío
$f(s_-)/f(s_+)$ son lectores distintos. Su antecedente angular común no
establece una igualdad entre ambos.

\subsubsection{Transporte discreto de la forma cuadrática}
\label{md:hib:transporte-discreto}

\begin{proposition}
\label{md:hib:proposicion-discreta}
Sean $S_\zeta=\operatorname{diag}(e^{\zeta/2},e^{-\zeta/2})$,
$g_\zeta=S_\zeta^{-T}S_\zeta^{-1}$ y
$\mathsf R(-\theta)$ una rotación. Para valores de forma $\zeta_n$ y
fases $\theta_n$, el transporte
\begin{equation}
 X_{n+1}=S_{\zeta_{n+1}}\mathsf R(-\theta_n)S_{\zeta_n}^{-1}X_n
 =T_nX_n
 \label{md:hib:transporte}
\end{equation}
conserva exactamente
$I_n=X_n^Tg_{\zeta_n}X_n/2$ y tiene determinante uno.
\end{proposition}
\begin{proof}
La sustitución y la identidad $\mathsf R^T\mathsf R=I$ dan
\begin{equation}
 T_n^Tg_{\zeta_{n+1}}T_n=g_{\zeta_n}.
 \label{md:hib:isometria-transportada}
\end{equation}
La forma cuadrática se conserva al evaluar ambos lados sobre $X_n$.
Los factores de $T_n$ tienen determinante uno, de donde
$\det T_n=1$.
\end{proof}
Es una composición del transporte radial previamente construido. Si
$\zeta_n,\theta_n$ se extraen de una ruta TPK, la identidad se aplica a
esa secuencia; la proposición algebraica conserva separada la selección
de la ruta.

\subsubsection{Transporte diferenciable y término de conexión}
\label{md:hib:conexion}

Si $\zeta(t)$ es de clase $C^1$ y $\omega(t)>0$ es continua, la relación
$X=S_\zeta Y$ y la evolución $\dot Y=-\omega J_0Y$, con
\begin{equation}
 J_0=\begin{pmatrix}0&-1\\1&0\end{pmatrix},
 \label{md:hib:j-canonico}
\end{equation}
dan
\begin{equation}
 \dot Q=\frac{\dot\zeta}{2}Q+\omega e^\zeta P,\qquad
 \dot P=-\omega e^{-\zeta}Q-\frac{\dot\zeta}{2}P.
 \label{md:hib:ecuaciones-transporte}
\end{equation}
Con las convenciones $\dot Q=\partial_PH$ y
$\dot P=-\partial_QH$, el Hamiltoniano es, salvo un sumando dependiente
sólo del tiempo,
\begin{equation}
 \boxed{H_{\mathrm{tot}}=
 \frac{\omega}{2}\left(e^{-\zeta}Q^2+e^\zeta P^2\right)
 +\frac{\dot\zeta}{2}QP.}
 \label{md:hib:hamiltoniano-conexion}
\end{equation}
El invariante
\begin{equation}
 I=\frac12\left(e^{-\zeta}Q^2+e^\zeta P^2\right)
 \label{md:hib:invariante}
\end{equation}
se conserva exactamente, sin requerir una aproximación adiabática.
\begin{proof}
Al derivar se obtiene
\[
 \dot I=\frac{\dot\zeta}{2}
 \left(e^\zeta P^2-e^{-\zeta}Q^2\right)
 +e^{-\zeta}Q\dot Q+e^\zeta P\dot P.
\]
La sustitución de \eqref{md:hib:ecuaciones-transporte} cancela los términos
proporcionales a $\dot\zeta$ y los dos términos $\omega QP$.
Así, $\dot I=0$. Si se omite el término mixto de
\eqref{md:hib:hamiltoniano-conexion}, queda
\begin{equation}
 \dot I=\frac{\dot\zeta}{2}
 \left(e^\zeta P^2-e^{-\zeta}Q^2\right).
 \label{md:hib:defecto-sin-conexion}
\end{equation}
El signo mixto contrario duplica ese defecto.
\end{proof}
Sobre la elipse, $QP=I\sin(2\theta)$ y $\dot\theta=-\omega$, por lo que
\begin{equation}
 H_{\mathrm{tot}}=I\omega+\frac{I\dot\zeta}{2}\sin(2\theta).
 \label{md:hib:hamiltoniano-elipse}
\end{equation}
La matriz de la forma cuadrática tiene determinante
$\omega^2-\dot\zeta^2/4$. Para $\omega>0$, es definida positiva
exactamente cuando $|\dot\zeta|<2\omega$, y semidefinida positiva en
igualdad. Este umbral concierne a la positividad: la existencia del
transporte y la estabilidad física conservan sus condiciones respectivas.
La conservación de $I$ no implica conservación de $H_{\mathrm{tot}}$,
que depende explícitamente del tiempo.

La versión cuántica exige la simetrización del término mixto,
$\dot\zeta(QP+PQ)/4$. Esta expresión indica la operación requerida sin
presuponer una representación cuántica adicional.

\subsubsection{Compatibilidad con la sección de acción generadora del contraángulo}
\label{md:hib:compatibilidad-accion}

En la construcción rectora,
\begin{equation}
 \zeta=\operatorname{artanh}
 \left[\frac{2(\eta_{\mathrm{ret}}+\alpha)}{1000\alpha}\right].
 \label{md:hib:forma-genealogica}
\end{equation}
Fijar $\alpha$ y $\eta_{\mathrm{ret}}$ fija también $\zeta$.
El resultado con $\zeta(t)$ variable es, por tanto, una extensión de
transporte entre formas. Para representar una trayectoria física del par
rector, la variación debe proceder de la genealogía correspondiente, y no
de una deformación elegida de manera independiente.

También puede componerse exactamente el transporte cuando cambia la
sección positiva de acción entre dos estados, conservando la misma coordenatización
dimensional. Sean
\begin{equation}
 \hbar_n>0,\qquad
 \rho_{\mathrm{act},n}=\frac{\hbar_{n+1}}{\hbar_n},\qquad
 \widetilde T_n=\sqrt{\rho_{\mathrm{act},n}}\,
 S_{\zeta_{n+1}}\mathsf R(-\theta_n)S_{\zeta_n}^{-1}.
 \label{md:hib:transporte-conforme}
\end{equation}
Entonces
\begin{equation}
 \det\widetilde T_n=\rho_{\mathrm{act},n},\qquad
 \widetilde T_n^Tg_{\zeta_{n+1}}\widetilde T_n
 =\rho_{\mathrm{act},n}g_{\zeta_n},\qquad
 \frac{I_{n+1}}{\hbar_{n+1}}=\frac{I_n}{\hbar_n}.
 \label{md:hib:accion-normalizada}
\end{equation}
Las igualdades siguen de \eqref{md:hib:isometria-transportada} al multiplicar
por el factor escalar. Es un transporte conformemente simpléctico: conserva
la acción normalizada y conserva el área de acción sin normalizar
exactamente cuando $\rho_{\mathrm{act},n}=1$.
Su límite diferenciable tiene traza $\dot\hbar/\hbar$. Un generador
hamiltoniano lineal en un plano con forma simpléctica fija tiene traza cero;
por ello, una variación material de $\hbar$ exige especificar la forma
transportada o los grados de libertad con los que se intercambia acción.
El término mixto de deformación, por sí solo, conserva traza cero.

El resultado es un criterio de compatibilidad de las realizaciones. Para
interpretar variaciones numéricas debidas a un cambio de unidades se
transportan primero las bases: el cambio de coordenatización mantiene la sección
física. La composición no afirma una variabilidad experimental de la
constante de Planck.

\subsubsection{Transformación de representación y deformación material}
\label{md:hib:interpretacion-transporte}

Si $S_\zeta$ es un cambio pasivo de coordenadas, el término mixto es una
conexión de representación. Se transportan también los observables,
\begin{equation}
 L_X(S_\zeta Y)=L_Y(Y).
 \label{md:hib:observables-transportados}
\end{equation}
Esta igualdad conserva la lectura física al cambiar la expresión del
Hamiltoniano, sin introducir una interacción adicional ni trabajo físico.

Una deformación material exige obtener $\zeta$ de la actualización TPK y
mostrar un cambio relativo entre sistema y referencia, ambos transportados
coherentemente. Puede contrastarse una fase relativa, una respuesta
$Z_{\mathrm{sistema}}/Z_{\mathrm{referencia}}$ o una transferencia entre
modos definidos por el detector. Deben distinguirse dos historias TPK no
equivalentes por cambio de representación y conservar sus condiciones de
lectura. Cuando se atribuya un intercambio energético, el balance comprende
campo, memoria y dispositivo. El teorema de transporte delimita esta
distinción y conserva separada la ley dinámica que la haga física.

\subsection{Momento, contrastes y alcance de la composición}
\label{md:hib:momento}

La realización de Clifford--Dirac recibe la masa de ruta. En una misma
sección de acción, su coeficiente se escribe
\begin{equation}
 \frac{m_\Gamma c}{\hbar}
 =\frac{\omega_0\Xi_\Gamma}{c}.
 \label{md:hib:coeficiente-dirac}
\end{equation}
En una realización libre y homogénea, con la relación relativista de
dispersión declarada, ésta se reexpresa como
\begin{equation}
 \Omega^2=c^2|k|^2+\omega_0^2\Xi_\Gamma^2.
 \label{md:hib:dispersion}
\end{equation}
La reescritura conserva las hipótesis de esa realización y sus antecedentes
de conexión espacio-temporal. Los momentos de retorno $q_-^h-q_+^h$
son evaluaciones de las torres y se distinguen de los momentos mecánicos.

El factor $\Xi_\Gamma$ reaparece como peso másico, frecuencia relativa,
dilatación radial y exponente térmico. La respuesta constitutiva determina
la conversión espacial común. Barbero conserva información orientada
complementaria a la velocidad. La acción fija la medida areal; el reloj
introduce la tasa temporal; la forma y su transporte especifican el reparto
y su evolución.

Los contrastes correspondientes comprenden cocientes de masas y frecuencias;
compatibilidad de velocidad, impedancia y Barbero; conservación de la forma
cuadrática bajo transporte; y pesos térmicos una vez fijadas temperatura y
multiplicidades. Estas relaciones preservan más datos que una coincidencia
decimal aislada. Sus pruebas mantienen el balance energético y las
condiciones de las entropías consideradas; no implican entropía
termodinámica nula.

% Ampliación para insertar después del bloque térmico de hibridacion.tex.
% Procedencia, dependencias y alcance de los controles: INFORME.md.
\section{Lectura marcada y conjugación termométrica de la acción}
\label{md:r27:termica}

La acción orientada, el calendario, los lectores de capacidad y el archivo de
memoria proceden de la misma dinámica enriquecida APP--TRIT--TPK. Las dos hojas
conservan residuo y cociente; el régimen trítico conserva orientación; el
transporte $U_t$ y su holonomía $\Gamma_9$ preceden a la lectura reducida
$K_{\rm ph}$. El retorno de fase conserva una historia que sigue creciendo.
La presente construcción compone esos resultados en un registro de referencia
y ensayo, dentro de la estructura discreta conjunta del continuo ya construida.
Su regla termométrica se especifica mediante dos operaciones sobre ese registro.

\subsection{Capacidad, referencia espectral y origen retenido}
\label{md:r27:capacidad}
La capacidad de un prefijo de índice $t\ge0$ y su resto son
\begin{equation}
 C_t=\lfloor\log_{10}(9^{t+1})\rfloor,\qquad
 9^{t+1}=10^{C_t}R_t,\qquad 1\le R_t<10.
 \label{md:r27:capacidad-resto}
\end{equation}
El conteo se evalúa por potencias enteras. Las desigualdades
$10^{20}\le9^{21}<10^{21}$ y $10^{20}\le9^{22}<10^{21}$,
junto con $10^{j-1}\le9^j<10^j$ para $1\le j\le21$, prueban que
la primera vacancia positiva tiene índice $21$ y capacidad $20$.
En la lectura trítica de capacidad, los residuos $0,1,2$ corresponden a
$+1,-1,0$, respectivamente. Las capacidades en $t=21,22,23,24$ son
$20,21,22,23$; el primer retorno posterior del régimen es, por tanto,
$t=24$. El prefijo contiene $25$ posiciones y expresa capacidad $23$.
Así se conservan las funciones diferenciadas de $(20,24,25,23)$.
El regreso del régimen conserva una fase nonádica distinta; la fase retorna
en $t=30$, con capacidad $29$.

El portador del lector es $\mathcal D=\mathcal D_0\oplus\mathcal D_1
\oplus\mathcal D_2$, con rangos $(1,380,649)$ y graduación $N=0,1,2$.
La construcción conserva dos graduaciones: $N$ distingue el prefijo y $L$
la ocupación del espacio exterior o simétrico. Para $m=20$, $n=24$,
sean $\mathcal F_2^J$ el fijo exterior de grado dos de dos hojas intercambiadas
y $\mathcal B_{\le2}^J$ el fijo simétrico de grado a lo sumo dos. El portador es
$\mathbb C\Omega_{\rm ext}\oplus\mathcal F_2^J\oplus\mathcal B_{\le2}^J$.
La involución $J$ intercambia las dos hojas mediante el intercambio graduado;
en el sector exterior cada uno de sus autoespacios tiene dimensión $m$.
Los caracteres de los subespacios fijos son
\begin{equation}
 \begin{aligned}
 \chi_F(z)&=\frac{(1+z)^{2m}+(1-z^2)^m}{2},\\
 \chi_B(z)&=\frac{P_n(z)^2+P_n(z^2)}{2},\\
 P_n(z)&=\prod_{j\ge1}(1-z^j)^{-n}.
 \end{aligned}
 \label{md:r27:caracteres-fock}
\end{equation}
En efecto, el proyector $(I+J)/2$ promedia el carácter total y la traza
del intercambio. En el espacio exterior, los autovalores $+1,-1$ de $J$
producen $(1+z)^m(1-z)^m$; en el producto simétrico, sólo los pares idénticos
contribuyen a la traza del intercambio, produciendo $P_n(z^2)$.
Para los coeficientes hasta grado dos basta el polinomio
$P_n(z)=1+nz+n(n+3)z^2/2+O(z^3)$; los productos formales son finitos en
cada grado. Extraer los coeficientes demuestra
\begin{equation}
 \begin{split}
 \dim\mathcal F_2^J&=m(m-1)=380,\\
 (\dim\mathcal B_0^J,\dim\mathcal B_1^J,\dim\mathcal B_2^J)
 &=(1,n,n(n+2))=(1,24,624).
 \end{split}
 \label{md:r27:grados-fock}
\end{equation}
Por tanto, la descomposición simultánea $(N,L;\text{multiplicidad})$ es
\[
 \begin{gathered}
 (0,0;1),(1,2;380),(2,0;1),\\
 (2,1;24),(2,2;624).
 \end{gathered}
\]
El rango $649=1+24+624=25^2+24$ conserva tres ocupaciones; el vacío externo
y el vacío del sumando simétrico permanecen ortogonales. El carácter conjunto
es $1+380xy^2+x^2(1+24y+624y^2)$.
La asignación de $N=0,1,2$ a los cortes $24,27,30$, de capacidades
$23,26,29$, especifica el lector utilizado en esta realización; mantiene $L$
como observable concomitante. En particular,
$\operatorname{Tr}q^L=2+24q+1004q^2$ y
$\operatorname{Tr}q^N=1+380q+649q^2$ conservan operaciones distintas.

El origen de la referencia espectral reside en
\eqref{md:antmem:cuantica:eq:5.1}--\eqref{md:antmem:cuantica:eq:5.10}.
El acoplamiento nonádico $8:1$ y el lazo ordenado
$H_1=\left(\begin{smallmatrix}2&1\\1&1\end{smallmatrix}\right)$,
con la preparación de dos modos puros idénticos y la coordenatización de covarianza
transportada conjuntamente, dan
$\nu^2=1+(8/81)[\operatorname{tr}(H_1H_1^{\mathsf T})-2]=121/81$.
La suma de cuadrados de las entradas de $H_1$ es siete, de modo que
$\nu=11/9$. El teorema~\ref{md:antmem:gauss:thm:espectro-gaussiano}
da entonces $r=(\nu-1)/(\nu+1)=1/10$ y autovalores $(1-r)r^j$.
Se utiliza esa preparación demostrada, incluida su reducción de memoria.
La división $j=3n+b$, $0\le b<3$, conserva cociente y residuo y satisface
\begin{equation}
 (1-r)r^{3n+b}=(1-r^3)(r^3)^n\frac{r^b}{1+r+r^2}.
 \label{md:r27:factorizacion}
\end{equation}
Por consiguiente, el cociente tiene razón $q=r^3=1/1000$ y el residuo conserva
la distribución $\operatorname{diag}(100,10,1)/111$. El aumento del cociente
tras completar los tres residuos conserva el acarreo. Si la energía elemental
es $\epsilon$, la energía del índice es $\epsilon(3n+b+1/2)$: el salto entre
bloques es $3\epsilon$ a la misma temperatura de la referencia.

Definimos el operador positivo de referencia y su intensidad:
\begin{equation}
 \begin{gathered}
 \eta=\bigoplus_{n=0}^{2}r^{23+3n}I_{d_n},\qquad
 (d_0,d_1,d_2)=(1,380,649),\\
 b_*:=\operatorname{Tr}\eta=\frac{1380649}{10^{29}}.
 \end{gathered}
 \label{md:r27:eta}
\end{equation}
En particular, $b_*=1{,}380649\times10^{-23}$ y $0<b_*<1$.
El origen $23$ está retenido en $\eta$; normalizar únicamente su estado
elimina ese factor de las probabilidades relativas.

La sección de acción de \eqref{md:hib:reloj} proporciona
$\epsilon_a=h_a/(108t_0)$ y $\beta_{\rm clk}=\ln3/\epsilon_a$.
Para la elevación energética aditiva de este lector de capacidad,
\begin{equation}
 \begin{split}
 H_{\rm tot}(t)&=2(t+1)\epsilon_a,\\
 H_{\rm mem}(t)&=\frac{\epsilon_a}{\ln3}\ln R_t,\\
 H_{\rm cap}(t)&=H_{\rm tot}(t)-H_{\rm mem}(t)
 =\frac{\epsilon_a\ln10}{\ln3}C_t.
 \end{split}
 \label{md:r27:energia-capacidad}
\end{equation}
La prueba consiste en tomar logaritmos en \eqref{md:r27:capacidad-resto}.
Por sustitución, $e^{-\beta_{\rm clk}H_{\rm cap}(t)}=10^{-C_t}$.
Sobre $\mathcal D$, esto da
$H_{\rm cap}=(\epsilon_a\ln10/\ln3)(23I+3N)$ y
$\operatorname{Tr}e^{-\beta_{\rm clk}H_{\rm cap}}=b_*$.
El resto permanece registrado como contribución energética; la compensación
es una lectura distinta de eliminar un factor por traza parcial.

\subsection{Identidad tensorial y especificación de los lectores}
\label{md:r27:lectores}
\begin{lemma}[Información incremental del registro producto]
Para toda matriz positiva $\eta$ y toda matriz de densidad $\rho$, con
$\mathscr H(X)=-\operatorname{Tr}(X\ln X)$ y $0\ln0=0$, se cumple
\begin{equation}
 \begin{gathered}
 \mathscr H(\eta\otimes\rho)-\mathscr H(\eta)
 =\operatorname{Tr}(\eta)s(\rho),\\
 s(\rho)=-\operatorname{Tr}(\rho\ln\rho).
 \end{gathered}
 \label{md:r27:tensor}
\end{equation}
\end{lemma}
\begin{proof}
Sobre los soportes, el logaritmo del producto tensorial es
$\ln\eta\otimes I+I\otimes\ln\rho$. La factorización de la traza y
$\operatorname{Tr}\rho=1$ dan
$\mathscr H(\eta\otimes\rho)=\mathscr H(\eta)+\operatorname{Tr}(\eta)s(\rho)$.
La continuidad de $x\ln x$ en cero extiende la igualdad a núcleos no triviales.
\end{proof}

Para \eqref{md:r27:eta}, formamos un registro normalizado y su marca:
\begin{equation}
 \widehat\eta=\eta\oplus(1-b_*)|f\rangle\langle f|,
 \qquad\Omega_\rho=\widehat\eta\otimes\rho,\qquad
 P=I_{\mathcal D}\oplus0,
 \label{md:r27:registro}
\end{equation}
donde $P$ actúa por la identidad sobre el factor de ensayo. La bandera $f$
completa la probabilidad de esta realización. Cualquier archivo microscópico
adicional se transporta como tal; la bandera conserva exclusivamente esta
distinción de selección. Para un Hamiltoniano fijo $H$, se especifican
\begin{equation}
 \begin{split}
 \mathcal S_P(\rho)&=-\operatorname{Tr}
 [P\Omega_\rho(I\otimes\ln\rho)]=b_*s(\rho),\\
 \mathcal E_P(\rho)&=
 \frac{\operatorname{Tr}[P\Omega_\rho(I\otimes H)]}
 {\operatorname{Tr}(P\Omega_\rho)}=\operatorname{Tr}(\rho H).
 \end{split}
 \label{md:r27:par-lector}
\end{equation}
La primera operación lee información antes del condicionamiento; la segunda
lee energía condicionada a la marca. Ambas operan sobre el mismo registro,
y la probabilidad de la marca sigue siendo $b_*$. La referencia permanece
fija al variar el ensayo.

\begin{theorem}[Conjugación termométrica del lector marcado]
Sea $H$ un Hamiltoniano autoadjunto fijo sobre un espacio finito y sea
$\rho_\beta=e^{-\beta H}/Z(\beta)$, con $\beta>0$ y
$\operatorname{Var}_{\rho_\beta}(H)>0$. Para el par
\eqref{md:r27:par-lector}, la temperatura conjugada satisface
\begin{equation}
 \Theta_P:=\left(\frac{d\mathcal S_P}{d\mathcal E_P}\right)^{-1},
 \qquad b_*\Theta_P=\beta^{-1}.
 \label{md:r27:transductor}
\end{equation}
\end{theorem}
\begin{proof}
Escribamos $E(\beta)=\operatorname{Tr}(\rho_\beta H)$.
Derivar la suma espectral finita da
$Z'/Z=-E$ y $E'=-\operatorname{Var}_{\rho_\beta}(H)$.
Como $s(\rho_\beta)=\beta E+\ln Z$, se tiene $s'=\beta E'$.
La referencia fija permite derivar $\mathcal S_P=b_*s$ sin término adicional,
y la varianza positiva permite dividir por $E'$. Por tanto,
$d\mathcal S_P/d\mathcal E_P=b_*\beta$, que prueba el enunciado.
\end{proof}

La normalización es parte de la definición operativa. Si información y energía
se leen ambas condicionadas, sus valores son $(s,E)$ y su razón diferencial es
$\beta$. Si ambas se leen antes de condicionar, sus valores son $(b_*s,b_*E)$
y la misma razón es $\beta$. El par \eqref{md:r27:par-lector} produce
$b_*\beta$ porque conserva la intensidad en la primera operación y utiliza
energía por ensayo seleccionado en la segunda. Mantener ese par fija su
coeficiente: reemplazarlo por $cb_*$, con $c\ne1$, cambia la ecuación.
Si la referencia variase, $d(b_*s)=b_*\,ds+s\,db_*$; ese protocolo tendría
otra respuesta y exige conservar el segundo término.

\subsection{Selección variacional y balance de información}
\begin{theorem}[Principio variacional en la realización marcada]
Para $\Theta>0$, el funcional
\begin{equation}
 \Phi_\Theta(\rho)=\mathcal E_P(\rho)-\Theta\mathcal S_P(\rho)
 =\operatorname{Tr}(\rho H)-b_*\Theta s(\rho)
 \label{md:r27:funcional}
\end{equation}
tiene un único mínimo sobre las matrices de densidad, dado por
\begin{equation}
 \rho_\Theta^*=
 \frac{e^{-H/(b_*\Theta)}}{\operatorname{Tr}e^{-H/(b_*\Theta)}}.
 \label{md:r27:gibbs}
\end{equation}
Además,
\begin{equation}
 \Phi_\Theta(\rho)-\Phi_\Theta(\rho_\Theta^*)
 =b_*\Theta D(\rho\Vert\rho_\Theta^*)\ge0.
 \label{md:r27:variacional}
\end{equation}
\end{theorem}
\begin{proof}
El estado \eqref{md:r27:gibbs} es fiel en dimensión finita.
Sustituir $\ln\rho_\Theta^*=-H/(b_*\Theta)-\ln Z_\Theta$ en
$D(\rho\Vert\rho_\Theta^*)=\operatorname{Tr}\rho(\ln\rho-\ln\rho_\Theta^*)$
da la igualdad. Para verificar la positividad, diagonalícense $\rho$ y el
estado fiel $\sigma=\rho_\Theta^*$, con autovalores $r_i,s_j$ y matriz
bistocástica de solapamientos $p_{ij}$. La concavidad de $\ln$ da
$D(\rho\Vert\sigma)\ge\sum_i r_i\ln(r_i/t_i)$,
donde $t_i=\sum_jp_{ij}s_j>0$ y $\sum_i t_i=1$.
La desigualdad $-\ln u\ge1-u$ demuestra que esta última suma es no negativa.
La igualdad exige $r_i=t_i$ para todos los índices y que cada vector propio
de $\rho$ esté contenido en un autoespacio de $\sigma$ de valor $t_i$;
por tanto, exige $\rho=\sigma$. Esto prueba el mínimo y su unicidad.
La convención sobre el núcleo incluye estados singulares; la fidelidad de
$\sigma$ impide que uno de ellos alcance la igualdad.
\end{proof}
En esta realización, el mismo coeficiente interviene en la conjugación
diferencial y en el equilibrio variacional. La generalización a sistemas
infinitos conserva por separado las condiciones de traza, diferenciabilidad
y finitud de energía y entropía; los teoremas anteriores tienen dominio finito.

\subsection{Composición con acción, área y reciprocidad gravitatoria}
\label{md:r27:composicion}
Se utilizan las secciones orientadas ya construidas, en una misma coordenatización:
\begin{equation}
 \begin{gathered}
 \epsilon_a=\frac{h_a}{108t_0}=\frac{\hbar_a c_{\rm int}}{L_\alpha},
 \qquad q_A=\gamma a_0L_\alpha^2,\\
 \mathcal T_{I/A}=\frac{\epsilon_a}{q_A},\qquad
 G_a=\frac{c_{\rm int}^3L_\alpha^2}{\hbar_a}.
 \end{gathered}
 \label{md:r27:secciones}
\end{equation}
La igualdad longitudinal conserva el reloj de
\eqref{md:rigidez:reloj}; $q_A>0$ fija la orientación areal utilizada.
El funcional $\gamma$ es el de \eqref{md:compos:funcional-barbero} y mantiene
las incidencias hexada--octada y los canales $90/120$. Su aparición en el área
conserva esa dependencia completa. La sección~\ref{md:rigidez:area}
construye $a_0$ como carácter trítico angular normalizado y demuestra
$\mathcal A=q_A(N_{90}+N_{120})$ sobre treinta y seis enlaces etiquetados.
Así, $q_A$ es el incremento de área por unidad de ocupación en esa sección,
y no un factor inferido de una entropía prefijada.

\begin{corollary}[Compatibilidad de los lectores térmico, areal y gravitatorio]
Con la referencia trítica $\beta_{\rm clk}=\ln3/\epsilon_a$ y el mismo par
\eqref{md:r27:par-lector}, se cumple
\begin{equation}
 \begin{gathered}
 b_*\Theta_{{\rm clk},P}
 =\frac{\epsilon_a}{\ln3}
 =\frac{\mathcal T_{I/A}\gamma a_0L_\alpha^2}{\ln3},\\
 G_a b_*\Theta_{{\rm clk},P}\ln3=c_{\rm int}^4L_\alpha.
 \end{gathered}
 \label{md:r27:area-gravedad}
\end{equation}
Para la sección de horizonte $\tau_R=\hbar_a c_{\rm int}/(2\pi R)$,
\begin{equation}
 b_*T_{R,P}=\tau_R,\qquad
 \frac{T_{R,P}}{\Theta_{{\rm clk},P}}
 =\frac{\ln3}{2\pi}\frac{L_\alpha}{R}.
 \label{md:r27:horizonte}
\end{equation}
\end{corollary}
\begin{proof}
Aplicar \eqref{md:r27:transductor} al reloj da la primera igualdad.
Sustituir $q_A$ en $\mathcal T_{I/A}q_A=\epsilon_a$ da la segunda.
Multiplicar por $G_a$ y usar \eqref{md:r27:secciones} da
$G_a\epsilon_a=c_{\rm int}^4L_\alpha$. Para el horizonte se aplica el mismo
teorema con $\beta_R=1/\tau_R$ y se divide por la igualdad del reloj.
\end{proof}
Las conversiones entre secciones $i,j$ conservan
$\Theta_{i,P}/\Theta_{j,P}=\tau_i/\tau_j$. Todo ciclo de conversiones
se cierra por cancelación telescópica. La temperatura del horizonte y la del
reloj coinciden sólo cuando coinciden sus secciones energéticas.

En una coordenatización interna, $b_*$ es el coeficiente del transductor definido por
\eqref{md:r27:par-lector}. La unidad térmica se especifica mediante la
conjugación entre esas lecturas y la sección energética elegida. La comparación
con la aplicación denominada Boltzmann en otro protocolo transporta también
sus lectores, sus rectas y sus bases; la cifra sola conserva el alcance escalar.
En particular, esta ampliación introduce una realización termométrica
explícita y sus pruebas. Su regla se incorpora como tal, sin atribuirla
retrospectivamente a todo operador térmico anterior ni identificar por nombre
su unidad interna con una unidad metrológica externa.

% Fragmento ES para inserción, no documento autónomo.
% Destino: VII/manuscrito/main.tex, después de relaciones_termica.tex.
% La sección recibe las salidas y pruebas de normalizacion_conjunta.tex,
% composicion.tex y relaciones_termica.tex; no altera sus generadores.
\section{Composición de horizonte y restitución térmico-geométrica}
\label{ins29:vii:horizonte}

La composición de las constantes permite distinguir qué conserva cada
lectura física del mismo estado generado. Recibimos las secciones de
acción y longitud, el par angular y la velocidad constitutiva ya
construidos desde APP--TRIT--TPK. Abreviamos \(\eta=\eta_{\rm ret}\)
y denotamos por \(s_0\) la sección de acción cuya expresión en la coordenatización
SI es \(10^{-34}\,\mathrm{J\,s}\). Así,
\[
 \hbar=\eta s_0,\qquad
 L_\alpha=r_N\alpha^{16}L_*,\quad r_N=\frac{5759}{23040},\qquad
 G=\frac{c^3L_\alpha^2}{\eta s_0}.
\]
El retorno inscrito determina \(r_N\); el normando local determina la
potencia \(16\); la sección \(L_*\) conserva la coordenatización longitudinal.
La función \(\eta=H_5-90\mathcal C_\pi/\pi\), con sus coeficientes y
dependencias regionales, es la obtenida en
\eqref{md:compos:seccion-accion}. Las sustituciones siguientes reciben esas
salidas y conservan las unidades y los dominios de sus lectores.

\subsection{Área fija, masa fija y conservación del producto térmico}

Para el contraste de horizonte de Einstein, la expresión de
Bekenstein--Hawking es
\(S_{\rm BH}=k_Bc^3\mathcal A/(4G\hbar)\).
El coeficiente \(1/4\) pertenece aquí a esta ley de contraste.
En un horizonte de Schwarzschild, estacionario, sin carga ni rotación,
\(R_S=2GM/c^2\), \(\mathcal A=4\pi R_S^2\) y
\(k_BT_H=\hbar c/(4\pi R_S)\).

\begin{proposition}[Dos restricciones de la composición de horizonte]
\label{ins29:vii:area_masa}
En la realización indicada, las salidas estructurales dan
\begin{align*}
 S_{\mathcal A}&=\frac{k_B\mathcal A}{4L_\alpha^2},&
 T_{\mathcal A}&=\frac{\eta s_0c}{2k_B\sqrt{\pi\mathcal A}},\\
 S_M&=\frac{4\pi k_Bc^2L_\alpha^2M^2}{\eta^2s_0^2},&
 T_M&=\frac{\eta^2s_0^2}{8\pi k_BL_\alpha^2M},
 \qquad S_MT_M=\frac{Mc^2}{2}.
\end{align*}
En la comparación entre secciones que conserva \(c,L_\alpha,k_B\),
la lectura a área fija mantiene \(S_{\mathcal A}\) y hace
\(T_{\mathcal A}\) proporcional a \(\eta\); la lectura a masa fija
da \(S_M\propto\eta^{-2}\) y \(T_M\propto\eta^2\).
\end{proposition}
\begin{proof}
La identidad \(G\hbar=c^3L_\alpha^2\) da inmediatamente la primera
entropía. Sustituir \(R_S=\sqrt{\mathcal A/(4\pi)}\) en la temperatura
da la primera temperatura. A masa fija,
\(R_S=2cL_\alpha^2M/(\eta s_0)\); su área y su inverso producen las
otras dos fórmulas. Multiplicarlas cancela los factores comunes y deja
\(Mc^2/2\).
\end{proof}

La cancelación de la acción en la entropía a área fija expresa una
compatibilidad entre sus lecturas gravitatoria y longitudinal. A masa
fija, la misma sección de acción interviene en la geometría del horizonte,
y su dependencia permanece en la entropía. El balance debe conservar la
restricción elegida durante toda la comparación.
El parámetro radial \(R\) de la realización anterior
\(\hbar c/(2\pi R)\) coincide en este contraste con \(2R_S\):
ambas expresiones representan entonces la misma temperatura. Esta
identificación conserva el factor dos entre radio geométrico y escala
de aceleración.

\subsection{Separación angular y restitución de la sección de acción}

La realización modular de frontera conserva dieciocho ocupaciones
tríticas en cada sector de incidencia. Escribimos \(X=N_{90}\),
\(Y=N_{120}\), \(N=X+Y\) y \(M=X-Y\). Su separación declarada es
\[
 \Delta_{\rm mod}=
 \frac{(A_{\deg}-C^*_{\deg})\pi/180}{9\cdot90}
 \left(1-\frac7{12}\frac\pi{729}\right),\qquad
 d=30\Delta_{\rm mod},\quad f_\pi=1-\frac{7\pi}{8748}.
\]
La coordenada \(d\) conserva tanto la conversión angular \(\pi/180\)
como la diferencia incidencial \(120-90=30\).
En esta realización, la distribución sobre configuraciones microscópicas
es proporcional a \(\exp[-d(3X+4Y)]\).
Recibimos \(A_{\deg}=1000\alpha\),
\(C^*_{\deg}=2(\eta+\alpha)\) y ponemos
\(\xi=C^*_{\deg}/A_{\deg}=\tanh s\), en la rama
\(1/500<\xi<1\) de acción positiva. La rapidez de forma \(s\) es
adimensional; las dos coordenadas del cociente angular están expresadas
en grados.

\begin{proposition}[Curva de compatibilidad y lector de restitución]
\label{ins29:vii:curva}
Con \(d_*=50\pi f_\pi\alpha/243\), se cumple
\[
 d=d_*(1-\xi)=d_*(1-\tanh s),\qquad
 0<d<\frac{499}{500}d_*.
\]
La distribución anterior satisface, para \(m=1,\ldots,36\),
\[
 \frac{\Pr(M=m)}{\Pr(M=-m)}=e^{dm}.
\]
Por consiguiente, cada razón de probabilidades determina el mismo \(d\)
y restituye las coordenadas
\[
 \xi=1-\frac d{d_*},\qquad
 \eta=499\alpha-\frac{2430d}{\pi f_\pi},\qquad
 C^*_{\deg}=1000\alpha\left(1-\frac d{d_*}\right).
\]
\end{proposition}
\begin{proof}
La sustitución del par angular da
\(d=\pi f_\pi(499\alpha-\eta)/2430\).
Como \(\eta=500\alpha\xi-\alpha\), resulta la primera igualdad y,
por la rama declarada, el intervalo. Para la segunda, cada configuración
de ocupaciones \((X,Y)\) tiene una compañera \((Y,X)\) con la misma
multiplicidad. El cociente de sus pesos es
\(e^{-d(3X+4Y)+d(3Y+4X)}=e^{d(X-Y)}\).
Sumar sobre \(X-Y=m\) preserva ese factor. Aplicar el logaritmo y
despejar \(\xi\) y \(\eta\) proporciona las tres restituciones.
\end{proof}

Esta composición establece una comprobación conjunta: las treinta y seis
razones, divididas logarítmicamente por su respectivo \(m\), deben dar
la misma coordenada. La restitución actúa sobre salidas ya generadas y
mantiene su orden causal. Su resultado enlaza el sesgo de las ocupaciones
con la forma elíptica y la sección de acción, sin fijar un radio ni una
unidad energética adicional.

Los canales del funcional de Barbero se restituyen a continuación:
\[
 q_-=e^{-27d/f_\pi},\qquad
 q_+=D_A/q_-,\qquad D_A=e^{-100\pi\alpha/9}.
\]
En efecto, \(27d/f_\pi=\pi(A_{\deg}-C^*_{\deg})/180\), mientras
\(q_+q_-=D_A\). Sustituir estos canales y \(C^*_{\deg}(d)\) en
\eqref{md:compos:funcional-barbero} determina \(\gamma(d)\).
El coeficiente del área conserva entonces su composición
\(\gamma(d)a_0(d)L_\alpha^2\). Al recibir un autovalor de área
\(\mathcal A_n=\gamma a_0L_\alpha^2n\), el contraste de horizonte
da \(S_{\rm BH}(\mathcal A_n)/k_B=\gamma a_0n/4\).
La entropía del estado reducido conserva, además, sus probabilidades y
la bipartición. Ambas evaluaciones quedan disponibles con su objeto
respectivo: geometría del horizonte y distribución de entrelazamiento.

% Ampliación para insertar después de hibridacion.tex y del módulo térmico.
% La ecuación electrónica completa y las torres permanecen en sus secciones.
\section{Historia electrónica y conservación de la masa bajo transporte}
\label{md:r27:electron}

La ecuación \eqref{md:compos:electron-completo} evalúa una sección central
persistente y conserva su retorno de acción. La trayectoria, su memoria y
su masa son lecturas diferenciadas del mismo estado APP--TRIT--TPK. El
transporte de una ruta puede modificar eventos y memoria conservando el valor
del lector másico. Para decidirlo se utiliza el carácter completo de
\eqref{md:hib:caracter}, junto con el exponente de retorno.

\subsection{Separación y restitución de las dos historias centrales}
Los registros expresados de ambas orientaciones son
\begin{equation}
 e_{++}=(2,2,5,-4,-3,-2)^2,\qquad
 e_{--}=(4,3,2,-2,-2,-5)^2,
 \label{md:r27:historias}
\end{equation}
donde el exponente $2$ indica concatenación de dos bloques. Comparten total
nulo y el mismo histograma absoluto. Para un registro $e\in\mathbb Z^{12}$,
sean
\begin{equation}
 \begin{gathered}
 p_i=e_i+e_{i+6},\quad \mathfrak a_i=e_i-e_{i+6},\\
 s_C=\sum_{i=0}^{5}p_i,\quad M_i=p_i-p_5\quad(0\le i<5).
 \end{gathered}
 \label{md:r27:registro-electron}
\end{equation}
\begin{proposition}[Restitución integral del registro]
El dato $(s_C,M_0,\ldots,M_4,\mathfrak a_0,\ldots,\mathfrak a_5)$ determina el registro mediante
\begin{equation}
 \begin{gathered}
 p_5=\frac{s_C-\sum_{i=0}^{4}M_i}{6},\quad p_i=p_5+M_i,\\
 e_i=\frac{p_i+\mathfrak a_i}{2},\quad e_{i+6}=\frac{p_i-\mathfrak a_i}{2}.
 \end{gathered}
 \label{md:r27:inversa-electron}
\end{equation}
Su imagen integral se caracteriza por la divisibilidad por seis del primer
numerador y la congruencia $p_i\equiv \mathfrak a_i\pmod2$ para cada $i$.
\end{proposition}
\begin{proof}
Sumar $p_i=p_5+M_i$ para los cinco primeros índices y añadir $p_5$ da
$s_C=6p_5+\sum M_i$. Sumar y restar las dos definiciones de $p_i,\mathfrak a_i$
da las fórmulas restantes. Las condiciones indicadas son necesarias para
obtener enteros y suficientes para reconstruirlos; la sustitución devuelve
todos los datos originales.
\end{proof}
Para \eqref{md:r27:historias}, $\mathfrak a=0$ y
$M_{++}=(8,8,14,-4,-2)$, mientras $M_{--}=(18,16,14,6,6)$.
La restitución distingue ambas historias aunque lectores escalares concretos
compartan su valor. Los símbolos $\mathfrak a_i$ de esta descomposición son diferencias
de eventos y se distinguen del carácter areal $a_0$ por su dominio.

\subsection{Forma espectral del lector de retorno}
Para $\tau\in\mathbb R^{12}$, con índices cíclicos, definimos
\begin{equation}
 C_s=\sum_{j=0}^{11}\tau_j\tau_{j+s},\quad
 T_k=\sum_{j=0}^{11}\tau_j e^{-2\pi i kj/12},\quad
 \theta_k=2\pi k/12.
 \label{md:r27:fourier}
\end{equation}
La representación por caracteres se aplica después de producir el registro.
El lector de retorno conserva
$\Omega=-2C_1-4C_2-6C_3$ y $P_6=C_6/2$.
\begin{proposition}[Descomposición armónica exacta]
\begin{equation}
 \begin{split}
 C_s&=\frac1{12}\sum_{k=0}^{11}|T_k|^2\cos(s\theta_k),\\
 \Omega&=\frac1{12}\sum_{k=0}^{11}|T_k|^2
 [-2\cos\theta_k-4\cos2\theta_k-6\cos3\theta_k],\\
 P_6&=\frac1{24}\sum_{k=0}^{11}|T_k|^2(-1)^k.
 \end{split}
 \label{md:r27:espectro}
\end{equation}
\end{proposition}
\begin{proof}
Desarrollar $|T_k|^2$ produce
$\sum_{j,l}\tau_j\tau_l e^{-i\theta_k(j-l)}$.
La suma $\sum_{k=0}^{11}e^{i\theta_km}$ vale doce si $12$ divide $m$ y
cero en otro caso. Aplicar esta identidad al desplazamiento $s$ recupera
$C_s$. Como $\tau$ es real, las partes imaginarias se cancelan en pares
$k,12-k$ y queda el coseno. Las otras dos fórmulas resultan por combinación
lineal y por $\cos(6\theta_k)=(-1)^k$.
\end{proof}
Para $\tau_e=(+++---)^2$, sólo quedan las potencias
$|T_2|^2=64$, $|T_6|^2=16$, $|T_{10}|^2=64$. Sus pesos en $\Omega$
son $7,4,7$, de modo que $\Omega=80$ y $P_6=6$.
Se recupera así la parte armónica de
$\kappa_e=80-54\alpha+6\Delta_4$ en la ecuación electrónica íntegra.
El registro alternante $(+-)^6$, con igual histograma, tiene
$|T_6|^2=144$ y $\Omega=48$: el orden es una dependencia efectiva del lector.
Las traslaciones cíclicas conservan $|T_k|^2$, aunque cambien la fase de
$T_k$; esta lectura espectral preserva correlaciones, no toda la historia.

\subsection{Criterio de masa transportada y sus lecturas asociadas}
Sea $\mathcal H_{\mathcal F}$ la fibra de un dominio finito de rutas y
$\{|\Gamma\rangle\}$ su base. Se fija uno de los lectores de retorno
$r=D$ o $r=\Omega$ y una coordenatización común con $R_{\rm act}>0$.
La ley de masas de \eqref{md:hib:factor-ruta} determina
\begin{equation}
 \begin{gathered}
 \frac{m_\Gamma^\beta}{m_e^\beta}
 =\chi_{\rm ref}(\sigma_\Gamma-\sigma_e,\eta_\Gamma-\eta_e)
 R_{\rm act}^{\kappa_\Gamma-\kappa_e},\\
 \Lambda_\Gamma=\ln(m_\Gamma^\beta/m_e^\beta).
 \end{gathered}
 \label{md:r27:masa-relativa}
\end{equation}
El logaritmo actúa sobre una razón adimensional. En una fibra no escalar
se mantiene el operador y su realización diagonal declarada; el dominio
finito puede ser una órbita o multisección, sin seleccionar una celda arbitraria.

\begin{theorem}[Conservación del lector másico en un dominio estable]
Sea $F$ un transporte TPK que actúa biyectivamente en el dominio fijado y
$V|\Gamma\rangle=|F\Gamma\rangle$. Para
$M|\Gamma\rangle=m_e^\beta e^{\Lambda_\Gamma}|\Gamma\rangle$, se cumple
\begin{equation}
 [M,V]|\Gamma\rangle=m_e^\beta
 (e^{\Lambda_{F\Gamma}}-e^{\Lambda_\Gamma})|F\Gamma\rangle.
 \label{md:r27:conmutador-masa}
\end{equation}
Por tanto, $[M,V]=0$ equivale a $\Lambda_{F\Gamma}=\Lambda_\Gamma$ para
todas las rutas. Con los parámetros de la coordenatización fijos, la condición es
\begin{equation}
 \begin{aligned}
 &\delta n_A x+\frac{\delta s_C}{6}y+
 \delta k_{\Delta_4}\Delta_4+\delta b_\zeta\zeta_{270}\\
 &\qquad+\sum_h\delta N_h(q_-^h-q_+^h)
 +\delta\kappa\ln R_{\rm act}=0,
 \end{aligned}
 \label{md:r27:criterio-masa}
\end{equation}
donde $\delta$ compara la ruta transportada con la inicial.
\end{theorem}
\begin{proof}
Aplicar $MV$ y $VM$ a cada vector de la base da
\eqref{md:r27:conmutador-masa}. La positividad de $m_e^\beta$ y la
inyectividad de la exponencial real dan la equivalencia. Tomar el logaritmo
de \eqref{md:r27:masa-relativa} y restar los caracteres de ambas rutas da
\eqref{md:r27:criterio-masa}. En torres infinitas se exige para cada ruta
$\sum_h|N_h(q_-^h-q_+^h)|<\infty$; esta convergencia absoluta legitima
la resta. La finitud del dominio de rutas conserva los operadores acotados
utilizados en esta prueba.
\end{proof}

El transporte se produce primero mediante TPK; la igualdad decide después
qué conserva. Los eventos individuales y la memoria pueden cambiar y
compensarse en \eqref{md:r27:criterio-masa}. Al mantener la sección de acción,
el reloj y la realización circular de \eqref{md:hib:reciprocidad-radial},
la conmutación con $M$ conserva también las lecturas de frecuencia propia,
radio gravitatorio, longitud de Compton reducida y peso térmico que son
funciones del mismo operador positivo. Para el inverso de masa se trabaja
fuera del sector nulo. En la realización marcada de
\eqref{md:r27:area-gravedad}, la cuarta lectura adopta la forma
\begin{equation}
 \Xi_\Gamma=\frac{\Omega_\Gamma}{\omega_0}
 =\frac{r_{g,\Gamma}}{L_\alpha}
 =\frac{L_\alpha}{\bar\lambda_\Gamma}
 =\frac{E_\Gamma}{b_*\Theta_{{\rm clk},P}\ln3}.
 \label{md:r27:lecturas-conservadas}
\end{equation}
La prueba es sustitución de las igualdades ya demostradas en una misma coordenatización.
Este criterio describe conservación durante un transporte declarado; los
selectores físicos de especie, de representación y de interacción conservan
sus condiciones propias. La trayectoria y su lectura armónica añaden
información a la masa escalar sin reemplazar esas condiciones.

\subsection{Transporte de la sección electrónica al sector electrodébil}
\label{amp-ew-seccion}

La sección central de \eqref{md:hib:electron} conserva conjuntamente el
prefactor, el corrector y la memoria de retorno. Sus dependencias proceden
de las hojas APP, de la orientación TRIT y de la ruta registrada por TPK;
el carácter angular y las torres actúan después sobre las multisecciones
de la misma estructura discreta del continuo. Escribimos
\begin{equation}
 r_e=R_{\rm act}^{\kappa_e},\qquad
 E_e^\beta=\mathcal U_E\varepsilon_e^{(0)}r_e,\qquad
 \kappa_e=80-54\alpha+6\Delta_4.
 \label{amp-ew-electron}
\end{equation}
Aquí $\mathcal U_E$ es la sección de la recta energética y $r_e$ abrevia
el retorno electrónico. La década de acción permanece en su recta propia;
la coordenada energética realizada se transporta una sola vez.

Las firmas angulares de \eqref{amp-higgs-firmas} producen, en la misma
coordenatización, $E_s^\angle=E_e^\beta\exp(n_sx+s_sy/6)$. En particular,
\begin{equation}
 E_W^\angle=\mathcal U_E\varepsilon_e^{(0)}r_e
                  e^{94x-y/6},\qquad D=e^{-2x}.
 \label{amp-ew-w}
\end{equation}
Se utiliza esta realización angular de W. Los refinamientos de otra torre
conservan sus factores propios cuando se compone con ellos.

\begin{proposition}[Escala electrónica y acoplamiento de Fermi de árbol]
En la realización electrodébil de árbol, con las coordenadas energéticas
expresadas en una misma unidad y $g^2=4\pi\alpha/(1-D)$, las relaciones
$v=2E_W^\angle/g$ y $G_F^{(0)}=1/(\sqrt2v^2)$ dan
\begin{equation}
 G_F^{(0)}=
 \frac{\pi\alpha\,e^{-188x+y/3}}
 {\sqrt2(1-D)\mathcal U_E^2(\varepsilon_e^{(0)})^2r_e^2}.
 \label{amp-ew-fermi}
\end{equation}
Con las firmas, los canales y la coordenatización energética fijos, el paso de la
etapa basal a la etapa retornada satisface
\begin{equation}
 \frac{(G_F^{(0)})_\beta}{(G_F^{(0)})_0}=r_e^{-2},\qquad
 \frac{v_\beta}{v_0}=r_e.
 \label{amp-ew-escala}
\end{equation}
Los cocientes $E_W^\angle/E_Z^\angle$, $E_H^\angle/E_W^\angle$ y
$\lambda_H=(g^2/8)(E_H^\angle/E_W^\angle)^2$ son independientes de
ese factor común $r_e$.
\end{proposition}
\begin{proof}
La sustitución de $v$ da $G_F^{(0)}=g^2/[4\sqrt2(E_W^\angle)^2]$.
Aplicar la expresión de $g^2$ y elevar \eqref{amp-ew-w} al cuadrado
produce \eqref{amp-ew-fermi}: los coeficientes $188$ y $1/3$ resultan
respectivamente de $2\cdot94$ y $2/6$. Cambiar $E_e^0$ por
$E_e^\beta=r_eE_e^0$ multiplica todas las energías angulares por $r_e$.
La definición de $v$ conserva la primera potencia; la de $G_F^{(0)}$
conserva la inversa de su cuadrado. En cada cociente de energías se
cancela el mismo factor positivo, y la fórmula de $\lambda_H$ conserva
esa cancelación. Los acoplamientos $g,g'$ dependen de $\alpha,x$ y
permanecen fijos en esta comparación.
\end{proof}

La magnitud $G_F^{(0)}$ tiene dimensión de energía inversa al cuadrado.
La realización de árbol declarada especifica el alcance físico de esta
composición; las correcciones radiativas tienen sus operadores y su orden
posteriores. Comparar las dos etapas del lector conserva sus coordenatizaciones y
sus firmas, en lugar de postular una variación temporal independiente de
las constantes correlacionadas.

\subsection{Conservación de los refinamientos de ruta}
\label{amp-ew-refinamiento}
La composición con la ley completa preserva el carácter y la torre que
corresponden a cada ruta. Para dos autosecciones positivas comparables,
en una misma coordenatización y con un mismo lector $r\in\{D,\Omega\}$, sea
\begin{equation}
 \rho_{W/e}^{\,r}
 =\frac{\chi_{\rm full}^{\,r}(\sigma_W,\eta_W)}
        {\chi_{\rm full}^{\,r}(\sigma_e,\eta_e)}
 R_{\rm act}^{\kappa_r(\Gamma_W)-\kappa_e}.
 \label{amp-ew-razon-completa}
\end{equation}
La energía correspondiente es $E_W=E_e^\beta\rho_{W/e}^{\,r}$.
En la realización de Fermi compatible con esa sección se obtiene
\begin{equation}
 G_F^{(0)}=
 \frac{\pi\alpha}
 {\sqrt2(1-D)(E_e^\beta)^2(\rho_{W/e}^{\,r})^2}.
 \label{amp-ew-fermi-completo}
\end{equation}
La prueba sustituye la razón completa en la identidad cuadrática
anterior. En la coordenatización angular explícita, esa razón adopta la forma de
\eqref{amp-ew-w}; en una realización refinada retiene también las
diferencias de torres y de lectores de \eqref{md:hib:cocientes}.
Si $E_s=E_e^\beta\chi_sf_s$, entonces
$E_s/E_t=(\chi_s/\chi_t)(f_s/f_t)$. Esta igualdad muestra exactamente
qué factor se cancela y qué información específica permanece.

\subsection{Congruencia másica e inversión espectral sin conmutación}
\label{amp-ew-congruencia}
Sea $\widehat M^0>0$ el operador basal sobre una fibra finita positiva,
con su carácter y sus torres incorporados, y sea $\widehat B^r$
el lector autoadjunto de retorno. El transporte de acción produce
$C=R_{\rm act}^{\widehat B^r/2}>0$ y la ley
$\widehat M^\beta=C\widehat M^0C$.

\begin{proposition}[Reciprocidad operatoria de la longitud de acción]
Con $\hbar_a,c>0$ en una misma coordenatización dimensional,
\begin{equation}
 \begin{aligned}
 (\widehat M^\beta)^{-1}
    &=C^{-1}(\widehat M^0)^{-1}C^{-1},\\
 \widehat L_C^\beta:=\frac{\hbar_a}{c}(\widehat M^\beta)^{-1}
    &=C^{-1}\widehat L_C^0C^{-1},\qquad
 \widehat L_C^0=\frac{\hbar_a}{c}(\widehat M^0)^{-1}.
 \end{aligned}
 \label{amp-ew-inversa-operador}
\end{equation}
Estas identidades admiten $[C,\widehat M^0]\ne0$.
\end{proposition}
\begin{proof}
La positividad de $C$ implica invertibilidad. Para $u\ne0$,
$\langle u,C\widehat M^0Cu\rangle
=\langle Cu,\widehat M^0Cu\rangle>0$; por tanto, el operador final
también es invertible. La multiplicación directa en ambos órdenes da
\[
 (C\widehat M^0C)\bigl(C^{-1}(\widehat M^0)^{-1}C^{-1}\bigr)=I,
 \qquad
 \bigl(C^{-1}(\widehat M^0)^{-1}C^{-1}\bigr)(C\widehat M^0C)=I.
\]
Multiplicar por $\hbar_a/c$ prueba la segunda identidad. En la
descomposición espectral del operador final,
$\widehat M^\beta=\sum_jm_jP_j$ con $m_j>0$, se obtiene
$\widehat L_C^\beta=\sum_j\hbar_a/(cm_j)P_j$. La función recíproca
es inyectiva sobre los autovalores positivos: conserva sus proyectores
y multiplicidades finales.
\end{proof}

La congruencia puede cambiar los proyectores del operador basal; la
inversión espectral conserva los del operador final. El sector nulo
queda fuera de esta inversión. Una preparación que ocupe varios
autoespacios conserva su medida espectral y sus multiplicidades.
En rango uno, la identidad recupera la longitud electrónica
$\bar\lambda_{C,e}=\hbar_ac/E_e^\beta$ y, después,
$a_{{\rm B},e}=\bar\lambda_{C,e}/\alpha$. El resultado conecta la
escala electrónica, la respuesta electrodébil y la longitud de acción
mediante operaciones explícitas sobre sus estructuras compartidas.

% Procedencia: COMPATIBILIDAD_HOLONOMIA_E_INFORMACION.md, secciones 2--9.
% Desarrollo científico conservado; gestión y correspondencia en archivo separado.
\section{Compatibilidad constitutiva, holonomía relativa e información operacional}
\label{md:comp:capitulo}

\subsection{Reciprocidad diferencial de velocidad e impedancia}
\label{md:comp:reciprocidad}

El antecedente constitutivo de la sección~\ref{iii:sec:constitutiva} construye, sobre las hojas
conservadas de su representación,
\begin{equation}
 \begin{gathered}
 T=e^{-xI-y\mathcal R},\qquad
 r_\pm=f\bigl(e^{-30(x\pm y)}\bigr),\\
 V=(I-T^{90})(I-T^{120})^{-1}=r_+P_++r_-P_-.
 \end{gathered}
 \label{md:comp:constitutiva}
\end{equation}
donde
\[
 f(s)=\frac{1-s^3}{1-s^4},\qquad x>|y|.
\]
Los proyectores de hoja permanecen fijados. La cámara física marcada
\(x>y>0\) y su conjugada están contenidas en esta colección analítica.
Definimos las coordenadas logarítmicas normalizadas
\begin{equation}
 c_\ell=\log\widehat c=-\log r_+-\log r_-,\qquad
 z_\ell=\log\widehat Z=\log r_--\log r_+.
 \label{md:comp:logaritmos}
\end{equation}
Estas letras no designan nuevas constantes ni el contraángulo \(C^*\).
Sean
\[
 u=30(x+y),\qquad v=30(x-y),\qquad
 g(w)=\log f(e^{-w}),\qquad a=g'(u),\quad b=g'(v).
\]

\begin{proposition}[Reciprocidad constitutiva diferencial]
\label{md:comp:prop-reciprocidad}
En el dominio \(x>|y|\), se tiene
\begin{equation}
 g'(w)=\frac{3}{e^{3w}-1}-\frac{4}{e^{4w}-1}>0,
 \qquad
 D(c_\ell,z_\ell)=-30
 \begin{pmatrix}a+b&a-b\\a-b&a+b\end{pmatrix}.
 \label{md:comp:jacobiano}
\end{equation}
Los autovalores del jacobiano son \(-60a\) y \(-60b\). Por tanto, es
simétrico y definido negativo, su determinante es \(3600ab>0\), y
\begin{equation}
 \boxed{\partial_y\log\widehat c=\partial_x\log\widehat Z.}
 \label{md:comp:reciprocidad-cruzada}
\end{equation}
\end{proposition}
\begin{proof}
Se derivan las identidades
\(c_\ell=-g(u)-g(v)\) y \(z_\ell=g(v)-g(u)\).
La positividad de \(g'\) procede del decrecimiento estricto, en \(k>0\),
de \(k/(e^{kw}-1)\). También resulta de la expresión racional positiva
de la sección~\ref{md:comp:restitucion}. La matriz obtenida tiene los
autovalores y el determinante indicados; la igualdad de sus entradas
no diagonales da \eqref{md:comp:reciprocidad-cruzada}.
\end{proof}

La igualdad relaciona la respuesta de la velocidad a la separación
orientada con la respuesta de la impedancia a la coordenada común.
Ambas sensibilidades proceden de una misma función espectral. Es una
reciprocidad de esta colección constitutiva; no se identifica sin prueba
con una ley termodinámica o de electromagnetismo macroscópico.

\subsection{Potencial espectral y conexión de profundidad dos con Catalán}
\label{md:comp:potencial}

\subsubsection{Construcción del potencial}
\label{md:comp:construccion-potencial}

Desde los canales ya construidos, definimos
\begin{equation}
 \mathscr F(x,y)=\sum_{\sigma=\pm1}\sum_{n\geq1}
 \left[
 \frac{e^{-120n(x+\sigma y)}}{120n^2}
 -\frac{e^{-90n(x+\sigma y)}}{90n^2}
 \right].
 \label{md:comp:serie-potencial}
\end{equation}

\begin{proposition}[Potencial de las dos respuestas logarítmicas]
\label{md:comp:prop-potencial}
La serie \eqref{md:comp:serie-potencial} define en \(x>|y|\) un potencial
que satisface
\begin{equation}
 \boxed{\partial_x\mathscr F=c_\ell,\qquad
        \partial_y\mathscr F=z_\ell.}
 \label{md:comp:gradiente-potencial}
\end{equation}
Su hessiano es el jacobiano de
\eqref{md:comp:jacobiano}; en consecuencia, \(\mathscr F\) es
estrictamente cóncava. La constante aditiva queda fijada por
\(\mathscr F\longrightarrow0\) cuando \(x-|y|\longrightarrow\infty\).
\end{proposition}
\begin{proof}
En todo compacto del dominio \(x>|y|\), las cantidades \(x\pm y\)
tienen un mínimo positivo. Una potencia de \(n\) multiplicada por una
exponencial decreciente es sumable; por ello, la serie y sus derivadas
de cualquier orden fijo convergen uniformemente en ese compacto.
Derivando término a término y utilizando
\(-\log(1-z)=\sum_{n\geq1}z^n/n\), se obtienen las dos identidades
de \eqref{md:comp:gradiente-potencial}. El hessiano coincide entonces
con \eqref{md:comp:jacobiano}, que es definido negativo.
La normalización en el infinito es la de la serie exhibida.
\end{proof}

\(\mathscr F\) es un potencial matemático de la respuesta. No se le
asignan unidades de energía ni se declara un potencial termodinámico.

\subsubsection{Realización por un mismo momento espectral}
\label{md:comp:momento-espectral}

Sea
\[
 D_2(z)=\sum_{n\geq1}\frac{z^n}{n^2}.
\]
Su reconocimiento como dilogaritmo es posterior a esta definición por
serie. Para \(|z|<1\), satisface
\((z\partial_z)^2D_2(z)=z/(1-z)\). En esta notación,
\begin{equation}
 \mathscr F=
 \sum_{\sigma\in\{+,-\}}
 \left[\frac{D_2(q_\sigma^{120})}{120}
       -\frac{D_2(q_\sigma^{90})}{90}\right].
 \label{md:comp:potencial-dos-momentos}
\end{equation}
El momento orientado de Catalán construido en la sección~\ref{iii:sec:catalan} satisface
\begin{equation}
 \mathcal C_2(s)=\operatorname{Im}D_2(is)
 =\sum_{k\geq0}\frac{(-1)^ks^{2k+1}}{(2k+1)^2},\qquad
 \mathcal C_2(1)=G_{\rm Cat}.
 \label{md:comp:catalan-momento}
\end{equation}

\begin{proof}
Las potencias pares de \(i\) son reales y las impares tienen parte
imaginaria \((-1)^k\). Esto da la serie de
\eqref{md:comp:catalan-momento}, que converge absolutamente incluso
en \(s=1\). Para \(0\leq s<1\), aplicar término a término dos veces
\(s\partial_s\) devuelve \(s/(1+s^2)\), el coeficiente orientado del
resolvente del cuarto de giro ya presente en el corpus.
La identidad diferenciada en \(s=1\), si se utiliza, se interpreta
por continuación analítica o límite de Abel: la serie dos veces
diferenciada no converge ordinariamente en ese extremo.
La identificación \(\mathcal C_2(1)=G_{\rm Cat}\) procede directamente
de la serie original absolutamente convergente.
\end{proof}

Por tanto, el momento de Catalán y el potencial constitutivo se reúnen
mediante un mismo momento espectral de profundidad dos, evaluado en
soportes diferentes: el cuarto de giro orientado y las contracciones
de los sectores \(90/120\). Esta composición no identifica Catalán
con \(\mathscr F\), ni afirma que su único valor extremo determine
todos los canales. La colección analítica y sus argumentos conservan
información que un valor aislado no contiene.

\subsubsection{Compatibilidad con amplificación}
\label{md:comp:amplificacion-potencial}

Para \(T_m=T\otimes I_{9^m}\) y
\(\tau_m=\operatorname{Tr}/(2\cdot9^m)\), se tiene exactamente
\begin{equation}
 \mathscr F_m=2\tau_m\left[
 \frac{D_2(T_m^{120})}{120}
 -\frac{D_2(T_m^{90})}{90}\right]=\mathscr F.
 \label{md:comp:potencial-amplificado}
\end{equation}
La multiplicidad \(9^m\) cancela la normalización. El mismo argumento
conserva las derivadas. Esto acredita la amplificación de la
representación, no cualquier refinamiento imaginable del TPK.
No se sustituye el operador temporal por una lista de autovalores
sin marcas de hoja.

Para cualquier lazo suave cerrado en esta colección analítica,
\begin{equation}
 \oint(c_\ell\,dx+z_\ell\,dy)=0.
 \label{md:comp:integral-cerrada}
\end{equation}
Esta igualdad se refiere a la proyección constitutiva de dos
coordenadas. No implica que la historia enriquecida o su holonomía
operatoria sean triviales.

\subsection{Restitución exacta y sensibilidad finita}
\label{md:comp:restitucion}

\begin{proposition}[Restitución analítica de las coordenadas angulares]
\label{md:comp:prop-restitucion}
Con \(L=\log(16/9)\), el mapa
\((x,y)\longmapsto(c_\ell,z_\ell)\) es un difeomorfismo real analítico
del dominio \(x>|y|\) sobre
\begin{equation}
 0<c_\ell+z_\ell<L,\qquad 0<c_\ell-z_\ell<L.
 \label{md:comp:imagen-logaritmica}
\end{equation}
\end{proposition}
\begin{proof}
La función \(g\) crece estrictamente de \(\log(3/4)\) a \(0\).
Las combinaciones
\[
 c_\ell+z_\ell=-2g(u),\qquad c_\ell-z_\ell=-2g(v)
\]
recuperan \(u,v\) mediante \(g^{-1}\). Entonces
\(x=(u+v)/60\) y \(y=(u-v)/60\).
La inyectividad se mantiene al restringir a la imagen efectivamente
producida por HMT; no se declara que HMT produzca todo el cono.
\end{proof}

Con \(s=e^{-w}\), la derivada tiene la forma racional
\begin{equation}
 a(w)=g'(w)=
 \frac{s^3(s^2+2s+3)}{(1+s+s^2)(1+s+s^2+s^3)},
 \qquad 0<a(w)<\frac12.
 \label{md:comp:derivada-racional}
\end{equation}
La función es estrictamente decreciente, con \(a(0+)=1/2\), y
\(a(w)\sim3e^{-3w}\) en el infinito. Para probar el decrecimiento,
se utiliza
\begin{equation}
 a'(w)=-\frac{9}{4\sinh^2(3w/2)}
        +\frac{16}{4\sinh^2(2w)}<0,
 \label{md:comp:decrecimiento}
\end{equation}
pues \(k/\sinh(kw/2)\) decrece estrictamente en \(k>0\).
Además,
\begin{equation}
 \frac12e^{-3w}<a(w)<3e^{-3w}.
 \label{md:comp:cotas-exponenciales}
\end{equation}
La cota inferior resulta al multiplicar los denominadores y utilizar
\begin{equation}
 \begin{split}
 &2(s^2+2s+3)-(1+s+s^2)(1+s+s^2+s^3)\\
 &\hspace{1cm}=(1-s)(5+7s+6s^2+3s^3+s^4)>0.
 \end{split}
 \label{md:comp:polinomio-cota}
\end{equation}
La superior resulta de
\[
 3(1+s+s^2)(1+s+s^2+s^3)-(s^2+2s+3)>0.
\]

Si \(m=30(x-|y|)\) y \(M=30(x+|y|)\), la norma del inverso del
jacobiano es \(1/[60a(M)]\). Para no confundir el mapa vectorial con
el potencial escalar, escribimos
\(\mathbf F=(c_\ell,z_\ell)\).

\begin{proposition}[Cotas finitas de restitución]
\label{md:comp:prop-sensibilidad}
Sobre una región convexa en la que ambos argumentos \(u,v\) pertenecen
a \([m_0,M_0]\subset(0,\infty)\), se verifican
\begin{equation}
 60a(M_0)\|p-q\|
 \leq\|\mathbf F(p)-\mathbf F(q)\|
 \leq60a(m_0)\|p-q\|.
 \label{md:comp:cotas-finitas}
\end{equation}
\end{proposition}
\begin{proof}
Se integra el hessiano de \(\mathscr F\), es decir, el jacobiano
de \(\mathbf F\), sobre el segmento. Para la cota inferior se utiliza
\[
 -\langle\mathbf F(p)-\mathbf F(q),p-q\rangle
 \geq60a(M_0)\|p-q\|^2
\]
y la desigualdad de Cauchy--Schwarz. La cota superior utiliza la
norma máxima del jacobiano. Son cotas de errores finitos en el
dominio especificado.
\end{proof}

En \(y=0\) y cuando \(x\longrightarrow\infty\), ambas direcciones
se contraen por igual: el número de condición matricial es uno,
pero la norma del inverso crece como \(e^{90x}/180\).
Invertibilidad exacta y recuperación robusta frente a error absoluto
son, por tanto, propiedades diferentes. Cuando \(|y|\longrightarrow x\)
con \(x>0\) fijo, no hay singularidad diferencial: una derivada tiende
a \(1/2\) y la otra a \(a(60x)>0\).

\subsection{El cuarto de giro procede del ciclo dodecafásico}
\label{md:comp:cuarto-giro}

Se recupera la construcción de \eqref{iii:eq:quarter-intertwiner}. En \(\mathbb R^{12}\),
sean
\[
 C_{12}e_j=e_{j+1\bmod12},\qquad
 \Pi_4=P_4^{\rm cyc}\frac{I-C_{12}^6}{2},
\]
y
\begin{equation}
 \begin{gathered}
 u_0=(1,0,-1,0,1,0,-1,0,1,0,-1,0)^{\mathsf T},\qquad
 v_0=C_{12}u_0,\\
 W=(u_0,v_0)/\sqrt6.
 \end{gathered}
 \label{md:comp:isometria-ciclica}
\end{equation}
El símbolo \(\Pi_4\) distingue el proyector de las coordenadas de
acción \(Q,P\). Se verifican
\begin{equation}
 \begin{gathered}
 W^{\mathsf T}W=I,\qquad WW^{\mathsf T}=\Pi_4,\qquad
 C_{12}W=WJ_0,\\
 J_0=\begin{pmatrix}0&-1\\1&0\end{pmatrix},\qquad J_0^2=-I.
 \end{gathered}
 \label{md:comp:entrelazador-ciclico}
\end{equation}
En particular, \(C_{12}^3W=-WJ_0\). El plano tiene así un entrelazador
explícito con el ciclo; no se ha introducido un giro externo para
ajustar una respuesta.

El antecedente radial de \eqref{rec:eq:entrelazamiento} realiza
\[
 S_\zeta=\operatorname{diag}(e^{\zeta/2},e^{-\zeta/2}),
 \qquad J_\zeta=S_\zeta J_0S_\zeta^{-1}.
\]
Sobre ese plano definimos
\begin{equation}
 U_\zeta(\theta)=S_\zeta e^{-\theta J_0}S_\zeta^{-1}
 =\begin{pmatrix}
 \cos\theta&e^\zeta\sin\theta\\
 -e^{-\zeta}\sin\theta&\cos\theta
 \end{pmatrix}.
 \label{md:comp:transporte-plano}
\end{equation}
El signo de \(\theta\) fija la orientación elegida. La razón generada
\(\xi=C^*/A\), con \(|\xi|<1\) y ambos ángulos expresados en la misma
unidad, satisface \(\zeta=\operatorname{artanh}\xi\).
En la orientación conjugada,
\((\varepsilon,\xi)\mapsto(-\varepsilon,-\xi)\) conserva la sección
positiva de acción mediante el recuperador orientado del corpus;
\(\zeta\) cambia de signo.

\subsection{Lazo relativo de dos transportes}
\label{md:comp:lazo}

Sean
\begin{equation}
 U=U_{\zeta_s}(\theta),\qquad
 V=U_{\zeta_r}(\vartheta),\qquad
 H_{\rm lazo}=UVU^{-1}V^{-1}.
 \label{md:comp:def-lazo}
\end{equation}
Los productos se leen con acción de derecha a izquierda.
La letra \(\vartheta\) designa aquí la segunda fase, para distinguirla
de la coordenada regional \(\varphi\). Pongamos
\[
 \delta=\zeta_s-\zeta_r,\qquad
 \chi=\sinh\delta\,\sin\theta\,\sin\vartheta.
\]

\begin{proposition}[Traza y espectro del lazo relativo]
\label{md:comp:prop-lazo}
La composición \eqref{md:comp:def-lazo} satisface
\begin{equation}
 \boxed{\begin{gathered}
 \det H_{\rm lazo}=1,\qquad
 \operatorname{tr}H_{\rm lazo}=2+4\chi^2,\\
 H_{\rm lazo}=I\ \Longleftrightarrow\ \chi=0.
 \end{gathered}}
 \label{md:comp:traza-lazo}
\end{equation}
Sus autovalores son
\begin{equation}
 \lambda_\pm=(\sqrt{1+\chi^2}\pm|\chi|)^2
 =\exp\bigl(\pm2\operatorname{arsinh}|\chi|\bigr).
 \label{md:comp:espectro-lazo}
\end{equation}
\end{proposition}
\begin{proof}
Conjugando simultáneamente por \(S_{\zeta_r}^{-1}\), el operador \(V\)
es \(R(-\vartheta)\), y \(U\) tiene términos
\(e^{\pm\delta}\sin\theta\). La multiplicación directa da, en esa
referencia,
\[
 UV-VU=-2\chi\operatorname{diag}(1,-1).
\]
Como
\[
 H_{\rm lazo}-I=(UV-VU)U^{-1}V^{-1},
\]
se obtiene la equivalencia con \(\chi=0\). La multiplicación, o la
identidad de trazas para matrices unimodulares, da
\(\operatorname{tr}H_{\rm lazo}=2+4\chi^2\).
El polinomio característico es
\(\lambda^2-(2+4\chi^2)\lambda+1\), y sus raíces son las de
\eqref{md:comp:espectro-lazo}.
\end{proof}

Esto es una holonomía relativa de la composición declarada.
No se identifica con \(\Gamma_9\) ni con una actualización completa
del TPK por semejanza de nombres.

\subsubsection{Especialización enteramente determinada por el par conjugado}
\label{md:comp:lazo-conjugado}

Sean \(U_+=U_\zeta(\pi/2)\) y
\(U_-=U_{-\zeta}(\pi/2)\). Ambos satisfacen
\(U_+^2=U_-^2=-I\), pero
\begin{equation}
 \begin{aligned}
 U_+U_-&=-\operatorname{diag}(e^{2\zeta},e^{-2\zeta}),\\
 H_{\rm lazo}=(U_+U_-)^2
 &=\operatorname{diag}(e^{4\zeta},e^{-4\zeta}).
 \end{aligned}
 \label{md:comp:lazo-diagonal}
\end{equation}
Al escribir \(\xi=\tanh\zeta=C^*/A\), con ambos ángulos expresados
en la misma unidad, se obtiene
\begin{equation}
 \begin{aligned}
 H_{\rm lazo}
 &=\operatorname{diag}\left(
 \left[\frac{1+\xi}{1-\xi}\right]^2,
 \left[\frac{1-\xi}{1+\xi}\right]^2\right),\\
 \operatorname{tr}H_{\rm lazo}-2
 &=\frac{16\xi^2}{(1-\xi^2)^2}.
 \end{aligned}
 \label{md:comp:lazo-razon-angular}
\end{equation}
La fórmula de \(H_{\rm lazo}\) compone el cuarto de giro y las formas
conjugadas del mismo par. Su expresión es exacta en la razón angular
generada \(\xi\).

En la rama positiva de acción del corpus,
\(A_{\deg}=1000\alpha\) y
\(C^*_{\deg}=2(\eta_{\rm ret}+\alpha)\), con
\(\hbar_{\rm ret}=s_0\eta_{\rm ret}\). Por tanto,
\begin{equation}
 \begin{aligned}
 \xi&=\frac{\eta_{\rm ret}+\alpha}{500\alpha},\\
 \rho_{\rm lazo}=e^{4\zeta}
 &=\left(\frac{501\alpha+\eta_{\rm ret}}
                 {499\alpha-\eta_{\rm ret}}\right)^2,
 \qquad 0<\eta_{\rm ret}<499\alpha.
 \end{aligned}
 \label{md:comp:multiplicador-lazo}
\end{equation}
Los coeficientes \(501\) y \(499\) son \(500\pm1\);
\(500\) procede del factor \(1000/2\) del par angular y no es un
nuevo parámetro. Con los ejes orientados conservados, la restitución es
\begin{equation}
 \boxed{
 \eta_{\rm ret}=\alpha\left[
 500\frac{\sqrt{\rho_{\rm lazo}}-1}
          {\sqrt{\rho_{\rm lazo}}+1}-1\right].}
 \label{md:comp:restitucion-accion}
\end{equation}
La derivada
\[
 \frac{d\eta_{\rm ret}}{d\rho_{\rm lazo}}
 =\frac{500\alpha}
 {\sqrt{\rho_{\rm lazo}}(\sqrt{\rho_{\rm lazo}}+1)^2}
\]
es positiva. La cámara \(\eta_{\rm ret}>0\) corresponde a
\(\rho_{\rm lazo}>(501/499)^2\); el borde
\(\eta_{\rm ret}\longrightarrow499\alpha\) corresponde a
\(\rho_{\rm lazo}\longrightarrow\infty\).
Es una lectura inversa del coeficiente adimensional de acción desde
el transporte relativo y la coordenada \(\alpha\) ya construida.
Conserva la escala \(s_0\) como antecedente: una transformación
unimodular de forma no determina por sí sola esa escala dimensional.
La orientación conjugada invierte \(\rho_{\rm lazo}\) y transporta
la etiqueta de acción; no convierte \(\eta_{\rm ret}\) en acción negativa.

\subsubsection{Iteración del lazo y acumulación orientada}
\label{md:comp:iteracion}

Sea
\[
 \rho_{\rm lazo}=e^{4\zeta}
 =\left(\frac{1+\xi}{1-\xi}\right)^2.
\]
Para todo entero \(n\),
\begin{equation}
 \begin{gathered}
 H_{\rm lazo}^n
 =\operatorname{diag}(\rho_{\rm lazo}^{n},\rho_{\rm lazo}^{-n}),
 \\
 Q_n=\rho_{\rm lazo}^{n}Q_0,\quad
 P_n=\rho_{\rm lazo}^{-n}P_0.
 \end{gathered}
 \label{md:comp:iteracion-diagonal}
\end{equation}
La prueba es la multiplicación de matrices diagonales, extendida a
\(n<0\) mediante el inverso. Se conservan el producto
\(Q_nP_n=Q_0P_0\) y el área simpléctica, mientras el estado cambia
para \(\zeta\ne0\) y \(X_0\ne0\). En el programa de cuatro
operaciones, los incrementos de fase declarados suman cero en cada
repetición; esa contabilidad de fases no determina el transporte
total del estado.

Para \(\zeta\ne0\) y \(Q_0\ne0\), se recupera exactamente
\[
 n=\frac{\log|Q_n/Q_0|}{4\zeta}.
\]
Si \(Q_0=0\), la condición \(P_0\ne0\) permite usar
\[
 n=-\frac{\log|P_n/P_0|}{4\zeta}.
\]
La lectura requiere conocer la preparación y conservar una referencia.
La extensión multiplicativa
\(H_{\rm lazo}^{n+m}=H_{\rm lazo}^{n}H_{\rm lazo}^{m}\)
corresponde a un incremento logarítmico aditivo \(4(n+m)\zeta\).
Invertir el lazo cambia el signo de ese incremento.

Así, un programa fijo y repetido admite una respuesta que registra
su iteración. Esto es una realización matemática de acumulación
operacional en la composición estudiada. No se identifica \(n\)
con todo el acarreo o la memoria de \(\Gamma_9\), ni se deduce una
entropía nula para la dinámica completa a partir de la periodicidad
de su programa de fases. Si la amplitud inicial es desconocida y
sólo se dispone del estado final, \(n\) no queda determinado por
esa lectura aislada.

\subsubsection{Levantamiento explícito a las doce fases}
\label{md:comp:levantamiento}

Sean
\[
 \widetilde S_\zeta=I-\Pi_4+WS_\zeta W^{\mathsf T},\qquad
 L_\pm=\widetilde S_{\pm\zeta}C_{12}^3
       \widetilde S_{\pm\zeta}^{-1}.
\]
En \(\operatorname{Ran}\Pi_4\) actúan \(U_+,U_-\), respectivamente;
en el complemento, ambos actúan como \(C_{12}^3\). Por tanto,
\begin{equation}
 L_+L_-L_+^{-1}L_-^{-1}
 =I-\Pi_4+WH_{\rm lazo}W^{\mathsf T}.
 \label{md:comp:levantamiento-lazo}
\end{equation}
La acción complementaria se cancela. La traza del operador de doce
fases es \(10+\operatorname{tr}H_{\rm lazo}\).
El entrelazamiento y la composición son exactos y están comprobados
racionalmente usando los vectores enteros \(u_0,v_0\).
Este levantamiento reúne el antecedente dodecafásico y el transporte
elíptico; no afirma que su secuencia sea por sí misma la actualización
autónoma original del TPK.

\subsubsection{Refinamiento y límite}
\label{md:comp:limite}

La amplificación
\(H_{{\rm lazo},m}=H_{\rm lazo}\otimes I_{9^m}\) conserva
\[
 \frac{1}{9^m}\operatorname{Tr}H_{{\rm lazo},m}
 =\operatorname{Tr}H_{\rm lazo},
\]
así como la norma y el espectro, salvo las multiplicidades.
Con las inclusiones \(v\mapsto v\otimes e_0\), se cumple
\[
 H_{{\rm lazo},m+1}\iota_m=\iota_m H_{{\rm lazo},m}.
\]
El operador uniforme define en el límite inductivo hilbertiano una
extensión acotada, con la misma norma. Es el límite de esta
amplificación; no se sustituye por él otro sistema de refinamientos
del corpus.

\subsection{Qué distingue un cambio de coordenatización de un efecto relativo}
\label{md:comp:carta-efecto}

El transporte de coordenadas
\[
 T_n=S_{\zeta_{n+1}}R(-\theta_n)S_{\zeta_n}^{-1}
\]
telescopa:
\begin{equation}
 T_{N-1}\cdots T_0
 =S_{\zeta_N}R\left(-\sum_{n=0}^{N-1}\theta_n\right)
 S_{\zeta_0}^{-1}.
 \label{md:comp:telescopaje}
\end{equation}
Si retornan coordenatización y fase, resulta \(I\). El conmutador de la
sección~\ref{md:comp:lazo} compone evoluciones de formas distintas
en una referencia común; no es ese cambio pasivo.

Bajo un cambio común de representación \(P\), los operadores
\(U,V,H_{\rm lazo}\) se conjugan por \(P\).
La traza y el espectro de \(H_{\rm lazo}\) permanecen invariantes.
Si además \(X\mapsto PX\) y
\(g\mapsto P^{-\mathsf T}gP^{-1}\), se conserva
\(X^{\mathsf T}gX\) y toda lectura energética correctamente
transportada. Así, un \(H_{\rm lazo}\ne I\) no desaparece por
renombrar coordenadas.

Esto constituye un protocolo matemático relacional cerrado.
Su implementación experimental requiere que las evoluciones e
inversas sean operaciones físicas disponibles y que sistema,
referencia y controlador tengan una realización común.
No se afirma aquí una fuerza nueva ni una fuente de energía.

\subsubsection{Conservación de área y conservación de una misma energía}
\label{md:comp:area-energia}

\begin{proposition}[Ausencia de una métrica positiva común]
\label{md:comp:prop-metrica}
Cada cuarto de giro conserva su propia métrica:
\(g_\zeta=\operatorname{diag}(e^{-\zeta},e^\zeta)\) para \(U_+\),
y \(g_{-\zeta}\) para \(U_-\). Cuando \(\zeta\ne0\), no existe
una métrica positiva común.
\end{proposition}
\begin{proof}
Las condiciones
\[
 G=\begin{pmatrix}a&b\\b&d\end{pmatrix}>0,
 \qquad U_+^{\mathsf T}GU_+=G
\]
obligan a \(b=0\) y \(d=ae^{2\zeta}\).
La condición \(U_-^{\mathsf T}GU_-=G\) exige
\(d=ae^{-2\zeta}\). Ambas son compatibles sólo con \(\zeta=0\).
\end{proof}

Por tanto, conservar el área simpléctica no equivale a conservar
una misma forma energética durante toda composición de estos
transportes. En la referencia
\(g_r=S_r^{-\mathsf T}S_r^{-1}\), con \(\omega_r>0\),
\(X\ne0\) y \(E_r(X)=\omega_r X^{\mathsf T}g_rX/2\), se tiene
\begin{equation}
 \begin{gathered}
 \frac{E_r(H_{\rm lazo}X)}{E_r(X)}
 =\frac{x^{\mathsf T}\overline{H}_{\rm lazo}^{\mathsf T}
                   \overline{H}_{\rm lazo}x}{x^{\mathsf T}x},\\
 \overline{H}_{\rm lazo}=S_r^{-1}H_{\rm lazo}S_r,\qquad
 x=S_r^{-1}X.
 \end{gathered}
 \label{md:comp:cociente-energia}
\end{equation}
Los extremos son los cuadrados de los valores singulares.
En el caso conjugado diagonal son \(e^{\pm8|\zeta|}\), mientras
las dilataciones del estado son \(e^{\pm4\zeta}\).
Esta diferencia de exponentes es necesaria. Cambiar físicamente
de operación exige contabilizar el controlador; la ganancia de
una lectura no acredita creación de energía.

\subsubsection{El orden requiere una lectura orientada}
\label{md:comp:orden-orientado}

Intercambiar \(U_+\) y \(U_-\) invierte \(H_{\rm lazo}\).
La traza y el conjunto de ganancias extremas son iguales para
\(H_{\rm lazo}\) y \(H_{\rm lazo}^{-1}\). Esos escalares detectan
el efecto del lazo, pero no su orientación. Dos entradas
independientes con sus respuestas vectoriales completas reconstruyen
\(H_{\rm lazo}\); una medida escalar no basta en general.

En el caso conjugado, con ejes de referencia marcados, sean
\[
 G_Q=\frac{E_r(H_{\rm lazo}e_Q)}{E_r(e_Q)},\qquad
 G_P=\frac{E_r(H_{\rm lazo}e_P)}{E_r(e_P)}.
\]
La lectura
\begin{equation}
 \mathcal O=\frac14\log(G_Q/G_P)=4\zeta
 \label{md:comp:lector-orientado}
\end{equation}
recupera la orientación y cambia de signo al invertir el lazo.
Los ejes y la métrica deben transportarse juntos bajo un cambio
de representación. Así se distingue la orientación sin atribuirla
a una traza que la elimina.

\subsection{Información de historia frente a estadística de ocupación}
\label{md:comp:informacion}

Los productos
\[
 U_+U_-U_+^{-1}U_-^{-1}
 \qquad\text{y}\qquad
 U_+U_+^{-1}U_-U_-^{-1}
\]
contienen las mismas cuatro operaciones con las mismas multiplicidades.
El segundo vale \(I\); el primero es \(H_{\rm lazo}\ne I\) para
\(\zeta\ne0\). Por tanto, cualquier estadística que dependa
únicamente de las multiplicidades asigna el mismo resultado a dos
protocolos con acciones diferentes sobre el estado.

La conclusión exacta es que esa estadística no constituye una
descripción suficiente para predecir la respuesta de la composición.
No se refuta Shannon: una distribución sobre secuencias ordenadas
puede representar el orden. Se distingue la información que se ha
conservado en el observable elegido.

\begin{definition}[Información operacional de una historia]
\label{md:comp:def-informacion}
En esta representación, proponemos llamar información operacional de
una historia a la clase de su transporte ordenado distinguible por las
lecturas admitidas. Sea \(\mathcal D\) una colección de detectores
definidos sobre operadores. Puede incluir
\[
 D_{X,L}(H_{\rm lazo})=L(H_{\rm lazo}X),
 \qquad D_{\rm tr}(H_{\rm lazo})=\operatorname{tr}H_{\rm lazo},
\]
con preparación \(X\) y lector lineal \(L\).
Dos historias son equivalentes si
\[
 D(H_{{\rm lazo},1})=D(H_{{\rm lazo},2})
 \qquad\text{para todo }D\in\mathcal D.
\]
\end{definition}

Ampliar la colección puede separar clases antes indistinguibles.
Con respuestas vectoriales completas en una base se reconstruye todo
operador lineal; con sólo traza puede persistir
\(H_{\rm lazo}\sim H_{\rm lazo}^{-1}\).
Así se distingue el detector de una respuesta del estado del detector
espectral del operador.

Esta definición no identifica el transporte con toda la genealogía
TPK: historias distintas pueden producir el mismo operador en una
representación. La información latente en otras hojas, incidencias y
prolongaciones se conserva en sus propios lectores.

La dimensión estadística y la operacional son compatibles.
La tasa de crecimiento combinatorio, o entropía topológica, cuenta
el crecimiento asintótico de configuraciones admisibles; una tasa de
entropía de Shannon requiere además una distribución.
Un mapa de observación determina qué transformaciones pueden
reconstruirse. Ninguna de las pruebas anteriores declara entropía
termodinámica universal nula. El aporte es precisar qué descripción
es suficiente para una tarea de predicción o restitución.

\subsection{Consecuencias y nuevas definiciones propuestas}
\label{md:comp:consecuencias}

Los resultados anteriores reúnen cinco nociones, con sus dominios
y lectores conservados:
\begin{enumerate}
 \item \textbf{Potencial espectral constitutivo.}
 \(\mathscr F\), definido por los momentos \(90/120\), tiene como
 sus dos derivadas el logaritmo de la velocidad normalizada y el
 logaritmo de la impedancia normalizada. Es una definición matemática
 posterior a los canales, sin interpretación energética añadida.

 \item \textbf{Reciprocidad constitutiva diferencial.}
 La igualdad de sensibilidades cruzadas demostrada en la
 sección~\ref{md:comp:reciprocidad} proporciona una restricción
 conjunta verificable, además de los valores.

 \item \textbf{Defecto de cierre relativo.}
 \(\operatorname{tr}H_{\rm lazo}-2\) detecta el lazo no trivial.
 Se acompaña de una lectura orientada, como \(\mathcal O\), cuando
 se requiere distinguir el orden. La traza sola no define toda la memoria.

 \item \textbf{Restitución operacional.}
 Es la recuperación de la acción de la historia desde una colección
 especificada de observables. Se distingue de la restitución de dos
 canales y de la reconstrucción integral del estado.

 \item \textbf{Margen de restitución.}
 La cota \(60a(M_0)\) de la sección~\ref{md:comp:restitucion}
 cuantifica estabilidad absoluta en un dominio, mientras la
 inyectividad sólo afirma unicidad exacta.
\end{enumerate}

Las coordenadas de masa, frecuencia, radios y temperatura del
desarrollo precedente permanecen intactas. Este desarrollo añade
dos niveles: compatibilidad diferencial entre lectores escalares
y memoria de orden en composiciones operatorias.
No se identifica el factor de dilatación de \(H_{\rm lazo}\) con
una masa nueva sin pasar por el lector de masas correspondiente.

\clearpage
\begingroup
\fontsize{10.2}{12.2}\selectfont
\setlength{\parskip}{1pt plus .5pt minus .5pt}
\titlespacing*{\paragraph}{0pt}{6pt}{.75em}
\section{Conclusiones}
\label{md:conclusiones}
El resultado transversal de este trabajo es la determinación conjunta de
estructuras que sus evaluaciones escalares presentan por separado. Se han
construido mapas explícitos de compatibilidad y restitución, identificado
los invariantes de sus transformaciones y demostrado qué observaciones
recuperan la información que conservan. La composición reúne las constantes
mediante las operaciones que comparten.
La acción participa en una escala dimensional, en el contraángulo y en el
cociente de retorno electrónico. El par angular determina canales
orientados; las incidencias hexada--octada producen lectores constitutivos
y el funcional de Barbero--Immirzi. La composición conserva esta genealogía
y permite formular propiedades conjuntas que exceden la evaluación
individual de cada lector. El punto de partida permanece en la composición
APP--TRIT--TPK: las relaciones inversas actúan sobre sus salidas construidas
y preservan las orientaciones, los dominios y los datos dimensionales que
utilizan.

\paragraph{Memoria cronológica y respuesta operatoria.}

El análisis cronológico precisa qué información conserva una prolongación.
En el factor de capacidad y con la medida de fase declarada, la entropía
de los bloques crece sin cota mientras su tasa por posición tiende a cero.
El balance, que admite dos valores, pierde una cantidad de información
cronológica también creciente. La actividad cuadrática del lector de
corriente permanece positiva aunque su media sea nula. Estas propiedades
se refieren a observaciones diferentes del mismo proceso y pueden
coexistir sin identificar la tasa simbólica con una temperatura.

La restitución operatoria añade una distinción dinámica. En la
representación nonádica de lectura y memoria, la sensibilidad al orden
se distribuye en las contribuciones \(1/81\) y \(8/81\) del
conmutador interno. Reintroducir la memoria reúne ambas y produce
\(1/9\). La historia retenida participa, por tanto, en la composición
de la respuesta: conocer el balance presente y conocer la capacidad
de responder a una continuación son informaciones diferentes.

\paragraph{Determinación espectral y estabilidad de la reconstrucción.}
La restitución espectral establece un resultado de determinación conjunta. En la colección
constitutiva normalizada, la velocidad fija el producto de las dos
respuestas y el funcional de Barbero distingue de forma monótona su reparto
orientado. Ambos datos recuperan los autovalores etiquetados. Sobre la
imagen generada, la inversión se compone con la restitución angular y del
coeficiente de acción. Las marcas de hoja y las bases dimensionales se
conservan: son datos de la representación que la recuperación utiliza.

La caracterización diferencial precisa esa reconstrucción. Los logaritmos de velocidad e
impedancia son las derivadas de un mismo potencial espectral, definido
mediante los canales de incidencia. Su Hessiano negativo definido produce
una reciprocidad exacta de sensibilidades y permite acotar la inversión
en dominios controlados. La unicidad con datos exactos y la sensibilidad
ante un error finito quedan así caracterizadas por resultados distintos.
Esta precisión resulta especialmente relevante cuando una contracción
intensa conserva la inyectividad y reduce, al mismo tiempo, la separación
observable entre estados.

\paragraph{Composición constitutiva y reconstrucción electrodébil.}
La inversa cúbica transporta la multiplicación de contracciones a una
ley asociativa sobre las respuestas. Su coordenada logarítmica es
aditiva y conserva el sentido de la composición en cada hoja. La
reconstrucción electrodébil precisa otra complementariedad: el cociente
de acoplamientos de calibre recupera la contracción común, mientras
Higgs determina la coordenada orientada en el dominio especificado.
Los proyectores marcados convierten esa recuperación escalar en
restitución del operador constitutivo. Se mantiene así una conexión
demostrada entre lectores distintos de los mismos canales.

\paragraph{La constante de Catalán y el potencial constitutivo.}
El momento de profundidad dos establece una conexión entre dos evaluaciones
de una misma estructura funcional. Sobre el
cuarto de giro orientado, su parte imaginaria produce la colección cuyo
valor extremo es Catalán; sobre los canales contractivos, el mismo
momento interviene en el potencial constitutivo. Los argumentos, la
orientación y la normalización de cada evaluación permanecen especificados.
La estructura funcional transmite más información que cualquiera de sus
valores extremos considerado aisladamente.
La fórmula integral de la proposición~\ref{md:compos:inversa-catalan}
recupera además el momento desde la respuesta constitutiva. La condición
de límite nulo en el origen fija unívocamente las constantes de integración.
Producción y restitución quedan así articuladas sobre una misma colección
funcional, con su dominio y su valor extremo.

\newpage
\paragraph{Reconstrucción de la acción y memoria del orden.}
La composición relativa demuestra que el orden deja
una huella operatoria recuperable. Los cuartos de giro de formas
conjugadas y sus inversos producen un lazo unimodular cuyo factor es
\[
 \rho_{\rm lazo}=
 \left(\frac{501\alpha+\eta_{\rm ret}}
 {499\alpha-\eta_{\rm ret}}\right)^2,
 \qquad 0<\eta_{\rm ret}<499\alpha.
\]
La rama y las referencias declaradas permiten recuperar
\(\eta_{\rm ret}\). Su escala dimensional sigue siendo el antecedente
necesario para reconstruir la sección de Planck. El levantamiento al
ciclo de doce fases establece dónde actúa esa transformación, y la
iteración expresa su acumulación mediante potencias recíprocas. La
holonomía nonádica conserva su propio dominio y su ley de actualización;
la composición estudiada constituye un transporte posterior bien definido.

La información operacional adquiere así una definición precisa: dos
historias son distinguibles respecto de las preparaciones y los
detectores admitidos cuando sus transportes producen respuestas
diferentes. Las multiplicidades de operaciones no determinan en general
esa respuesta; el orden puede ser indispensable. La traza y las respuestas
vectoriales tampoco conservan los mismos datos. Una colección suficiente
de lecturas reconstruye el operador, mientras otras identifican
transformaciones de orientación opuesta. Esta definición hace explícita
la relación entre observación, memoria y restitución sin confundir una
representación con la genealogía completa del estado.

\paragraph{Reciprocidad másica y articulación de las magnitudes físicas.}
La normalización conjunta tiene una prueba operatoria previa:
el retorno inscrito produce $L_\alpha$, y el transporte polar del mismo
carácter determina los lectores $R$ y $\Lambda$.
Las identidades $H\Lambda=\hbar cI$ y $R\Lambda=L_\alpha^2I$ fijan
$G=c^3L_\alpha^2/\hbar$ como coeficiente único de la realización radial.
La prueba conserva variaciones no conmutativas y la eliminación compatible
de memoria, incluido el transporte dual de la longitud de acción.
Su evaluación retornada en la sección SI es
$6.674299659414022706\ldots\,10^{-11}\,\mathrm{m^3\,kg^{-1}\,s^{-2}}$;
el contraste CODATA se efectúa después de esa evaluación.

La composición reúne cuatro lecturas de una misma coordenada
estructural positiva de masa. Bajo las condiciones circular y térmica
declaradas, la frecuencia propia, el radio gravitatorio reducido, la
longitud de Compton reducida y la energía normalizada por la decisión
trítica satisfacen
\[
 \Xi_\Gamma=
 \frac{\Omega_\Gamma}{\omega_0}
 =\frac{r_{g,\Gamma}}{R_{108}}
 =\frac{R_{108}}{\bar\lambda_\Gamma}
 =\frac{E_\Gamma}{k_B\Theta_{108}\log3}.
\]
El prefactor electrónico, el carácter y las torres pertenecen al
constructor de \(\Xi_\Gamma\). La identidad conserva sus funciones y
sus normalizaciones. El factor másico se caracteriza, en esa realización,
como parámetro de dilatación recíproca de las dos longitudes; su producto
mantiene el área \(R_{108}^2\). La ponderación térmica del ciclo se
expresa mediante \(3^{-\Xi_\Gamma}\), con las condiciones de
normalización y convergencia correspondientes.
Aquí $R_{108}=L_\alpha$; la forma circular conserva la longitud ya
construida. Al variar el reloj se transporta también el cociente electrónico
$b_e$, preservando la misma lectura de masa. Las ventanas cronológicas,
las fases reducidas y los costes terminales describen aspectos del estado;
su conservación aislada permite cambios que la composición completa distingue.

\paragraph{Escala electrónica, Fermi y longitud de acción.}
La sustitución de la sección electrónica completa conserva su prefactor
y su retorno dentro de la realización electrodébil. La escala común
aparece con potencia inversa doble en Fermi de árbol y con potencia
directa en el valor de expectación de Higgs; se cancela en los
cocientes angulares y en el autoacoplamiento. Los factores específicos
de cada especie permanecen en sus razones. La congruencia positiva
transporta, además, la longitud de Compton por la congruencia inversa,
sin exigir conmutación entre masa basal y retorno. Estas pruebas
distinguen la escala absoluta, la información angular y la estructura
espectral conservada por cada operación.

\paragraph{Restitución electrónica y conservación bajo transporte.}
La lectura electrónica conserva una dependencia verificable del orden.
La suma y las diferencias entre posiciones opuestas reconstruyen el
registro, incluidas sus condiciones de integridad; las autocorrelaciones
y el espectro separan registros con iguales multiplicidades. El
coeficiente $80$ corresponde a la evaluación del registro central
especificado, y cambia al alterar su orden. Para una colección finita
de multisecciones, el operador de masa es positivo y su conmutador
con un transporte determina exactamente cuándo éste conserva la masa.
Ese criterio permite trasladar la misma conservación a las lecturas
temporal, espacial, gravitatoria y térmica, con sus referencias fijas.

\newpage
\paragraph{Traza marcada y respuesta termométrica.}
La composición de capacidades y ocupaciones produce
\[
 b_*=10^{-23}+380\,10^{-26}+649\,10^{-29}
     =1{,}380649\times10^{-23}.
\]
La graduación de los $649$ estados conserva los sectores $1$, $24$
y $624$ y sus respectivos índices de memoria. El espectro de razón
$1/10$ procede del transporte de memoria ya demostrado. Sobre esa
referencia se construye una realización termométrica explícita:
la información marcada es $b_*s(\rho)$ y la energía es la lectura
condicionada al sector marcado. Para variaciones de Gibbs con
Hamiltoniano y referencia fijos, la temperatura satisface
$b_*\Theta=\beta^{-1}$. El coeficiente queda así caracterizado por
una operación térmica concreta. Condicionar ambas magnitudes de la
misma manera produce otra normalización; variar la referencia añade
el término diferencial correspondiente. La prueba identifica esas
dependencias y conserva la coordenatización dimensional necesaria para la
lectura física de Boltzmann. Al componer con el reloj trítico, el
área de entrelazamiento y la reciprocidad gravitatoria, se obtiene
la relación conjunta entre energía de decisión, acción y área.

\paragraph{Contenido físico de la transformación.}
El análisis distingue de forma constructiva el transporte
pasivo y la respuesta relativa. El primero telescopa al retornar coordenatización
y fase. La segunda conserva un operador final distinto de la identidad
cuando el conmutador es no trivial. La conservación simpléctica de área
coexiste con la ausencia de una métrica energética positiva común para
las dos formas conjugadas. Interpretar una ganancia energética exige
incluir la realización del controlador y la referencia. El resultado
matemático determina qué debe contrastarse y qué datos deben conservarse
para atribuir una respuesta al sistema físico.

\paragraph{Definición estructural y unidad del resultado.}
Los desarrollos permiten precisar las definiciones utilizadas. El
coeficiente de acción caracteriza una sección regional que participa tanto
en la construcción angular como en su restitución por transporte. El
funcional de Barbero--Immirzi distingue el reparto orientado de los canales
constitutivos; unido a la velocidad normalizada, determina su espectro.
La constante de Catalán aparece como una evaluación extrema del momento
orientado cuya colección interviene también en el potencial de respuesta.
La coordenada de masa parametriza, en la realización especificada, una
dilatación recíproca que conserva el producto de las longitudes. La
constante de Boltzmann articula las rectas de temperatura y energía de
decisión mediante el transductor positivo y la normalización trítica
declarados. Cada definición añade al valor una operación, un dominio y una
función dentro de la composición.

\paragraph{Separación crítica y restitución de lectores.}
El escalado paramétrico queda referido a la sucesión que selecciona la propia construcción: las primeras raíces de retorno del lector dodecafásico uniforme. APP--TRIT--TPK, el continuo conjunto, la orientación y la memoria permanecen como su antecedente. La identificación de centros y el control del resto de \ref{hmtfg:escalado:15}--\ref{hmtfg:escalado:17} culminan en \eqref{hmtfg:articulacion:limite}: la constante de Feigenbaum es el invariante asintótico de separación de esos parámetros. El mismo desarrollo permite restituir el perfil desde la medida de doce nodos mediante Jacobi y recuperar, por inversión constitutiva, los canales orientados y el momento cuyo extremo es Catalán. El balance reúne así generación, selección y restitución, conservando el dominio y la información que distinguen cada lector.

% Adición a las conclusiones de VII; las conclusiones anteriores se conservan.
\paragraph{Integración de las cuatro interacciones y cierre canónico total.}
La gravitación y las interacciones fuerte, débil y electromagnética
actúan sobre una misma realización material de la estructura
APP--TRIT--TPK. El transporte conserva las hojas, la orientación,
la memoria y la incidencia del estado enriquecido; sus realizaciones
de espín y calibre se componen sobre las mismas fibras materiales.
La componente electromagnética pertenece a la realización
electrodébil: no se introduce como una interacción adicional
desconectada de ella. La acción total utiliza los coeficientes y
lectores ya construidos, incluidos la sección de Planck, la longitud
radial y el funcional de Barbero--Immirzi.

El contenido de esta integración no se limita a reunir términos en
una expresión. La variación de la misma acción produce la corriente
total de espín y la fuente métrica conjunta. La eliminación de la
contorsión conserva los términos cruzados entre especies; la
transformación de Legendre conserva los momentos materiales y
produce las restricciones totales. Las pruebas de las
secciones~\ref{hmtcuatro:gea:seccion}
y~\ref{hmtcuatro:accion-retroaccion-restricciones}
establecen la retroacción, la primera clase y la propagación de
esas restricciones en la realización regular especificada.

En particular, sea $\Theta_{\rm tot}$ el potencial canónico total
obtenido en esa transformación, $\mathcal Z=\{C_A=0\}$ la superficie
regular de restricciones y $X_N=N^AX_{C_A}$,
$X_M=M^BX_{C_B}$ dos deformaciones características admisibles.
La conexión inducida por la acción tiene
$H_N^{\rm tot,can}=-\Theta_{\rm tot}(X_N)$. Su curvatura, calculada
sin omitir las derivadas de los coeficientes de deformación, es
\begin{equation}
\begin{aligned}
\mathcal F^{\rm tot,can}_{N,M}
&=\delta_NH_M^{\rm tot,can}-\delta_MH_N^{\rm tot,can}
 -H_{[X_N,X_M]}^{\rm tot,can}
 +\frac{i}{\hbar}[H_N^{\rm tot,can},H_M^{\rm tot,can}]\\
&=-\mathrm d\Theta_{\rm tot}(X_N,X_M)
 =N^AM^B f_{AB}{}^{D}C_D.
\end{aligned}
\label{hmtcuatro:conclusion-curvatura-total}
\end{equation}
Por tanto,
\[
 \left.\mathcal F^{\rm tot,can}_{N,M}\right|_{\mathcal Z}=0.
\]
La fórmula incluye la deformación tangencial y las acciones de
calibre del corchete total; si se trabaja en el cociente, expresa
la misma igualdad módulo esas direcciones verticales. La
representación precuántica conserva este cierre mediante la
identidad de conmutadores demostrada en
\eqref{hmtcuatro:representacion-precuantica}.

Así, la incorporación gravitatoria es parte de la misma acción,
de sus fuentes y de su álgebra de restricciones, no una ecuación
yuxtapuesta después de construir las otras interacciones. La
anulación probada es la de la curvatura característica canónica;
no exige que sean nulas la curvatura del espacio-tiempo ni las
intensidades de los campos de calibre. Su dominio conserva el
coframe no degenerado, la carta regular y las condiciones de
borde utilizadas en la prueba. Tampoco suprime la memoria ni la
monodromía global de las historias que dieron lugar a la
realización.


La composición de horizonte distingue la dependencia de la acción a área
fija y a masa fija. La sección~\ref{ins29:vii:horizonte} demuestra además
una curva de compatibilidad entre el sesgo térmico, la forma elíptica y los
canales de Barbero. Las razones de probabilidades restituyen la misma
sección de acción y proporcionan una comprobación conjunta de esas lecturas.


Producción, compatibilidad y restitución constituyen, por tanto, tres
niveles complementarios de una misma descripción. La genealogía fija
los objetos, las relaciones conjuntas restringen sus lecturas y los
detectores determinan la estructura recuperable. Las nuevas
caracterizaciones se apoyan en las operaciones expuestas y conservan
sus condiciones, lo que permite incorporarlas a desarrollos sectoriales
sin disolver la unidad del argumento.

\endgroup
\clearpage\appendix
% Correspondencia editorial: gestion/mapa_transportes_auxiliares.md.
% Realizaciones auxiliares: no sustituyen el eje de compatibilidad y restitución.
\section{Transporte compensado, memoria geométrica y contraste energético}
\label{md:aux:transporte}

\subsection{Origen operatorio y dominio de la realización}
\label{md:aux:origen}

La construcción conserva APP--TRIT--TPK y el estado enriquecido. APP
reúne las hojas aritméticas con sus residuos y cocientes; TRIT conserva
el régimen y la orientación; TPK compone selección, transporte y
actualización de memoria. Las representaciones matriciales utilizadas
a continuación son realizaciones posteriores. La holonomía nonádica
$\Gamma_9$ y su memoria mantienen sus operadores propios.

Los antecedentes empleados son las representaciones cuadráticas
elíptica, hiperbólica y parabólica del TRIT, el diferencial de
transportes invertibles, el funcional energético
$E=\langle Df,WDf\rangle$ con $W>0$, el retorno con memoria
traslacional, la elipse de acción con producto de semiejes $2\hbar$
y su realización LC, la polarización de vacancias y el acoplamiento
reversible de dos modos. El retorno traslacional y el conmutador
que sigue son objetos distintos. El lazo elíptico--parabólico y
su balance de área orientada se utilizarán después como colección
de contraste.

La composición compensada y el protocolo de contraste reúnen estas
realizaciones. La teoría algebraica de trazas de conmutadores en
$\operatorname{SL}(2)$ constituye un antecedente matemático
establecido.\footnote{Goldman, \emph{Trace Coordinates on
Fricke spaces}, \href{https://arxiv.org/abs/0901.1404}{arXiv:0901.1404}.}
El objeto de este apéndice es su composición concreta con los
lectores declarados. Los parámetros $a,b$ son parámetros de las
realizaciones, no constantes físicas ajustadas.

El alcance termina en esta realización local. La identificación de
rutas TPK que admitan un dispositivo material con esos parámetros
requiere componer sus lectores efectivos. Una matriz invertible
arbitraria no acredita por sí sola una ruta TPK, y la identificación
del defecto de conmutación con torsión de Cartan requiere su
conexión y su soldadura.

\subsection{Compensación de operaciones y conservación de área}
\label{md:aux:conmutador}

Las matrices de régimen son
\begin{equation}
 J_{\mathrm e}=\begin{pmatrix}0&1\\-1&0\end{pmatrix},\qquad
 J_{\mathrm h}=\begin{pmatrix}0&1\\1&0\end{pmatrix}.
 \label{md:aux:regimenes}
\end{equation}
Satisfacen $J_{\mathrm e}^2=-I$, $J_{\mathrm h}^2=I$ y
$[J_{\mathrm e},J_{\mathrm h}]=2\operatorname{diag}(1,-1)$.
Se definen
\begin{equation}
 R(a)=e^{aJ_{\mathrm e}},\qquad B(b)=e^{bJ_{\mathrm h}},\qquad
 C_{\mathrm{aux}}(a,b)=R(a)B(b)R(a)^{-1}B(b)^{-1}.
 \label{md:aux:def-compensado}
\end{equation}
La acción sobre columnas se lee de derecha a izquierda. El protocolo
contiene cada transformación y su inversa, pero la sucesión no es
el retroceso completo del recorrido. Un retroceso completo tendría,
por ejemplo, la forma $RBB^{-1}R^{-1}=I$. El símbolo
$C_{\mathrm{aux}}$ distingue este conmutador elíptico--hiperbólico
del contraángulo $C^*$ y del lazo de formas conjugadas
$H_{\mathrm{lazo}}$ estudiado en el cuerpo principal.

\begin{proposition}
\label{md:aux:prop-traza}
Para $a,b\in\mathbb R$,
\begin{equation}
 \det C_{\mathrm{aux}}=1,\qquad
 \operatorname{tr}C_{\mathrm{aux}}=2+4\sin^2a\sinh^2b.
 \label{md:aux:traza-compensado}
\end{equation}
\end{proposition}
\begin{proof}
Póngase $c_a=\cos a$, $s_a=\sin a$, $u_b=\cosh b$ y $v_b=\sinh b$.
Entonces
\[
 R=\begin{pmatrix}c_a&s_a\\-s_a&c_a\end{pmatrix},\qquad
 B=\begin{pmatrix}u_b&v_b\\v_b&u_b\end{pmatrix},
\]
con $c_a^2+s_a^2=1$ y $u_b^2-v_b^2=1$. Todas las matrices tienen
determinante uno. La multiplicación da
\begin{equation}
 \begin{aligned}
 (C_{\mathrm{aux}})_{11}&=u_b^2+2c_as_au_bv_b-(c_a^2-s_a^2)v_b^2,\\
 (C_{\mathrm{aux}})_{12}&=-2s_a^2u_bv_b-2c_as_av_b^2,\\
 (C_{\mathrm{aux}})_{21}&=-2s_a^2u_bv_b+2c_as_av_b^2,\\
 (C_{\mathrm{aux}})_{22}&=u_b^2-2c_as_au_bv_b-(c_a^2-s_a^2)v_b^2.
 \end{aligned}
 \label{md:aux:entradas-compensado}
\end{equation}
La suma diagonal es
$2u_b^2-2(c_a^2-s_a^2)v_b^2=2+4s_a^2v_b^2$, como se afirma.
\end{proof}

Cuando $s_av_b\ne0$, la traza supera dos y el conmutador es
hiperbólico, con autovalores positivos recíprocos
\begin{equation}
 e^{\lambda},\ e^{-\lambda},\qquad
 \lambda=2\operatorname{arsinh}|\sin a\sinh b|>0.
 \label{md:aux:espectro-compensado}
\end{equation}
Se obtiene de $\cosh\lambda=\operatorname{tr}(C_{\mathrm{aux}})/2$
y $\cosh(2\operatorname{arsinh}x)=1+2x^2$.
Si $s_av_b=0$, entonces $b=0$ o $R=\pm I$, y
$C_{\mathrm{aux}}=I$.
La expansión local es
\begin{equation}
 C_{\mathrm{aux}}=I+2ab\operatorname{diag}(1,-1)
 +O(|a|^2|b|+|a||b|^2).
 \label{md:aux:expansion-compensado}
\end{equation}
El residuo comienza por el producto de las dos operaciones. Aplicar
cualquiera de ellas seguida de su propia inversa produce la identidad.

Una matriz real bidimensional de determinante uno conserva la forma
de área. Por ello, $C_{\mathrm{aux}}$ transforma la elipse de
acción en otra elipse de área $h$. Si $\Sigma>0$ es una matriz
de covarianza,
\begin{equation}
 \Sigma'=C_{\mathrm{aux}}\Sigma C_{\mathrm{aux}}^{\mathsf T},
 \qquad \det\Sigma'=\det\Sigma.
 \label{md:aux:covarianza-area}
\end{equation}
Pueden cambiar la distribución de semiejes y su orientación.
Los autovalores $e^{\pm\lambda}$ describen direcciones propias e
iteraciones; los valores singulares determinan otra lectura y, en
general, no coinciden con ellos. La variación de semiejes euclídeos
en una sola aplicación requiere la métrica correspondiente.

La conservación del área de acción es compatible con una
acumulación de deformación orientada. El observable geométrico
$\lambda$ depende de $a,b$ y no constituye una constante universal
adicional. La conservación del determinante o del área, por sí
sola, no establece conservación de energía física.

\subsection{Lector positivo del residuo e inversión del recorrido}
\label{md:aux:residuo}

En una fibra de retorno, con campo $f$ y peso hermítico $W>0$
declarados, el lector energético toma la forma
\begin{equation}
 E_C(f)=\langle(I-C_{\mathrm{aux}})f,
                  W(I-C_{\mathrm{aux}})f\rangle.
 \label{md:aux:energia-residuo}
\end{equation}
Para $s_av_b\ne0$, el conmutador carece de autovalor uno y
$E_C(f)>0$ para todo $f\ne0$. El retroceso completo produce residuo
cero. La normalización dimensional pertenece al funcional de
partida: sus coordenadas numéricas requieren su realización antes
de interpretarse como julios, calor, masa o corriente de conducción.

La inversión completa del recorrido invierte $C_{\mathrm{aux}}$.
La traza y $\lambda$ identifican el operador y su inverso. La
restitución de dirección requiere una lectura adicional: el vector
residual, las componentes de transporte o una sonda con resolución
de fase. Esta comparación limita de manera concreta la suficiencia
de una lectura exclusivamente escalar.

Un ejemplo racional, independiente de datos físicos, es
\begin{equation}
 R=\begin{pmatrix}3/5&4/5\\-4/5&3/5\end{pmatrix},\qquad
 B=\begin{pmatrix}5/4&3/4\\3/4&5/4\end{pmatrix}.
 \label{md:aux:ejemplo-racional}
\end{equation}
En él se obtiene $\operatorname{tr}C_{\mathrm{aux}}=86/25$,
$\det C_{\mathrm{aux}}=1$ y autovalores
$(43\pm6\sqrt{34})/25$. Los controles racionales del suplemento
incluyen este ejemplo y $169$ pares, con casos nulos, retroceso
completo y conservación del determinante de covarianza. Son
comprobaciones finitas de las identidades demostradas.

\subsection{Acoplamiento reversible y lectura reducida de entrelazamiento}
\label{md:aux:acoplamiento}

Para $H\in\operatorname{Sp}(2,\mathbb R)$, la realización de dos
modos utiliza el acoplamiento
\begin{equation}
 U_H=\frac19\begin{pmatrix}
 8I+H&\sqrt8(H-I)\\
 \sqrt8(H-I)&I+8H
 \end{pmatrix}.
 \label{md:aux:acoplamiento-uh}
\end{equation}
Con dos gaussianas puras idénticas preparadas desde la elipse de
acción y con transporte conjunto de la coordenatización inicial, el lector
reducido satisface
\begin{equation}
 \nu(H)^2=1+\frac8{81}
    \bigl\{\operatorname{tr}(HH^{\mathsf T})-2\bigr\},
 \qquad\operatorname{Tr}\rho_{\mathrm{vis}}^2=\nu(H)^{-1}.
 \label{md:aux:mezcla-generica}
\end{equation}
La preparación y los lectores se mantienen fijados. La
transformación conjunta es reversible y conserva la pureza global.

\begin{corollary}
\label{md:aux:cor-mezcla}
Para $C_{\mathrm{aux}}(a,b)$,
\begin{equation}
 \operatorname{tr}(C_{\mathrm{aux}}C_{\mathrm{aux}}^{\mathsf T})-2
 =16\sin^2a\sinh^2b\cosh^2b,
 \label{md:aux:traza-cuadratica}
\end{equation}
y
\begin{equation}
 \boxed{\nu_C^2=1+\frac{32}{81}\sin^2a\sinh^2(2b).}
 \label{md:aux:mezcla-compensado}
\end{equation}
\end{corollary}
\begin{proof}
Las cuatro entradas de \eqref{md:aux:entradas-compensado} se
escriben como $d+k$, $l-m$, $l+m$, $d-k$, donde
\[
 d=1+2s_a^2v_b^2,\quad k=2c_as_au_bv_b,\quad
 l=-2s_a^2u_bv_b,\quad m=2c_as_av_b^2.
\]
La suma de sus cuadrados es
$2(d^2+k^2+l^2+m^2)=2+16s_a^2u_b^2v_b^2$.
La sustitución en \eqref{md:aux:mezcla-generica} y
$\sinh(2b)=2\sinh b\cosh b$ prueban el resultado.
\end{proof}
El ejemplo \eqref{md:aux:ejemplo-racional} produce
$\nu_C^2=17/9$ y pureza $3/\sqrt{17}$.

La cancelación de $\hbar$ y de la deformación inicial utiliza
la preparación declarada y el transporte conjunto de coordenatización.
Un protocolo con $s_av_b\ne0$ produce mezcla reducida; la
inversión completa del acoplamiento puede deshacerla. Mezcla
reducida, energía almacenada y calor conservan sus funciones
como lecturas distintas.

\subsubsection{Contraste firmado en una colección elíptico--parabólica}
\label{md:aux:familia-hn}

La colección construida de lazos es
\begin{equation}
 H_n=\begin{pmatrix}1+n&n^2\\n&n^2-n+1\end{pmatrix},
 \qquad n\in\mathbb Z.
 \label{md:aux:hn}
\end{equation}
Su balance alternante de memoria depende de $n^2$, mientras el
cociente de determinantes de covarianza es
\begin{equation}
 \mathcal R_n=\frac{\det V_{\mathrm{vis},n}}{\det V_0}
 =1+\frac8{81}n^2(2n^2-2n+5).
 \label{md:aux:covarianza-hn}
\end{equation}
Al restar las dos orientaciones se cancelan los términos pares:
\begin{equation}
 \boxed{\mathcal R_{-n}-\mathcal R_n=\frac{32n^3}{81}.}
 \label{md:aux:diferencia-direccional}
\end{equation}
Para $n=1$ y $n=-1$, las respuestas son $121/81$ y $153/81$,
con purezas $9/11$ y $3/\sqrt{17}$, respectivamente.
Esta colección satisface $H_{-n}\ne H_n^{-1}$ en general;
tampoco es la iteración $H_1^n$. El cambio del parámetro
orientado debe distinguirse del retroceso completo de las
operaciones. Su lectura alternante de área conserva una
identificación distinta de la del funcional de Barbero.

El criterio de mezcla emplea el conmutador $[D,J]$, con
$D=\sqrt8(H-I)/9$. Si $D$ conmuta con la estructura compleja
compatible con el estado inicial, puede transportar memoria
sin mezclar ese producto gaussiano. Si no conmuta, el lector
determina la mezcla. La ausencia de mezcla para ese estado
es una propiedad distinta de la ausencia de disipación material.

La resta \eqref{md:aux:diferencia-direccional} es un corolario
de la colección y de su lector, no una constante universal.
La fórmula de $C_{\mathrm{aux}}$ compone el mismo lector con
otra colección admitida por el teorema. Un contraste material
debe determinar transportes y acoplamientos antes de medir;
la programación arbitraria de esas matrices acredita su
emulación y conserva un alcance diferente.

\subsection{Segundo momento de la corriente de vacancias}
\label{md:aux:vacancias}

El lector de capacidad fija
\[
 \delta=\log_{10}(10/9),\qquad z_n\in\{0,1\},\qquad
 P_N=\sum_{n<N}(z_n-\delta),\qquad |P_b-P_a|<1.
\]
En intervalos iguales $t_0$, la corriente centrada es
$I_n=A_I(z_n-\delta)$, con $A_I=u_Q/t_0$.
El símbolo $A_I$ distingue esta amplitud del ángulo principal.
Para $L=b-a$,
\begin{equation}
 \frac1{A_I^2}\sum_{n=a}^{b-1}I_n^2
 =L\delta(1-\delta)+(1-2\delta)(P_b-P_a).
 \label{md:aux:segundo-momento}
\end{equation}
La prueba consiste en desarrollar $(z_n-\delta)^2$, utilizar
$z_n^2=z_n$ y sumar. El desequilibrio firmado permanece
acotado, mientras
$I_{\mathrm{rms}}/|A_I|\longrightarrow\sqrt{\delta(1-\delta)}$.
La amplitud $A_I$ es el salto entre los dos niveles de corriente.
El contraste conserva la realización temporal y debe contabilizar
cualquier acoplamiento que filtre la sucesión.

Un balance firmado nulo puede coexistir con una lectura
cuadrática positiva. Si se añade una resistencia ideal $R_E>0$
y se mantiene esa corriente por intervalos, el modelo posterior
de contraste da
\begin{equation}
 Q_{R_E}=R_Et_0\sum_n I_n^2.
 \label{md:aux:contraste-resistivo}
\end{equation}
La resistencia pertenece a esta realización posterior. Su
función consiste en contrastar una posible reducción de calor
con un balance explícito.

\subsection{Condiciones de un contraste material}
\label{md:aux:contraste-material}

El protocolo separa dos cuestiones. Para la memoria de orden
se comparan recorridos compensados, intercambios de orden y
retrocesos completos mediante una respuesta remanente vectorial
o resuelta en fase. Una intensidad o una carga neta aisladas
pueden perder la diferencia buscada. Para la recuperación de
energía se miden entrada, energía útil o devuelta, variación
de almacenamiento y calor, incluidos preparación y restauración
del dispositivo. El retorno de fase conserva separada la
variación de energía almacenada.

En un contraste escalar de un puerto, la relación
$v_2(t)=v_1(T-t)$ conserva la distribución de amplitudes y el
espectro de potencia. En régimen periódico, un mismo sistema
lineal estacionario predice igual potencia media. Una ley
instantánea $i=g(v)$ también conserva la integral de $vi$.
Una diferencia reproducible examinaría dependencia histórica,
pero su atribución requiere comparar histéresis, calentamiento
y modelos dinámicos alternativos. En varios puertos deben
conservarse además los espectros cruzados.

Una reducción de disipación se contrasta a igual transferencia
útil, fidelidad y ritmo. La atenuación debe discriminar energía,
amplitud de señal, población de portadores y carga conservada
antes de atribuirle un origen gravitatorio. La distinción entre
almacenamiento, recuperación y pérdidas cuenta con precedentes
experimentales en recuperación electrocalórica.\footnote{Defay
et al., \emph{Enhanced electrocaloric efficiency via energy
recovery} (2018),
\href{https://doi.org/10.1038/s41467-018-04027-9}{doi:10.1038/s41467-018-04027-9}.
Este antecedente no constituye evidencia experimental del
transporte HMT aquí compuesto.}

El resultado del apéndice es una identidad general dentro
de la colección matricial especificada, un lector positivo del
residuo, un control de inversión y una ley exacta del segundo
momento de la corriente. Su alcance distingue esos resultados
de una nueva constante universal, un prototipo sin calentamiento
o una demostración de decaimiento electrónico gravitatorio.
La vinculación prospectiva con un dispositivo debe fijarse
antes de obtener los datos de contraste.

\section{Acción, memoria reducida y trabajo recuperable}
\label{md:aux:trabajo}

\subsection{Composición de la acción con una realización energética}
\label{md:aux:accion}

La composición recibe las salidas de acción y los transportes
de memoria de APP--TRIT--TPK, conservando el estado enriquecido
como antecedente. Su dominio es la realización gaussiana y
energética declarada. La pasividad y la ergotropía son conceptos
establecidos; se emplean para estudiar el acoplamiento específico
y los lectores ya construidos.

La ecuación de acción mantiene explícitos sus coeficientes:
\begin{equation}
 H_5=\frac12\sqrt{\frac{1000\alpha}{\varphi}}-\alpha
 +\frac9{16}\alpha^2-\frac59\alpha^3
 +\frac7{48}\alpha^4-\frac1{54}\alpha^5,
 \label{md:aux:h5}
\end{equation}
\begin{equation}
 D_A=e^{-100\pi\alpha/9},\qquad
 \mathcal C_\pi=169\alpha^6+\frac{D_A\alpha^7}{1-D_A\alpha},
 \label{md:aux:correccion-accion}
\end{equation}
\begin{equation}
 \begin{aligned}
 \hbar_{\mathrm{pre}}&=s_0H_5,\\
 \hbar_{\mathrm{ret}}&=s_0\left(H_5-\frac{90}{\pi}\mathcal C_\pi\right),
 \qquad h_s=2\pi\hbar_s,\qquad s_0=10^{-34}\mathcal U_S.
 \end{aligned}
 \label{md:aux:secciones-accion}
\end{equation}
La década y la coordenatización dimensional conservan antecedentes
distintos. La fórmula mantiene las coordenadas $\pi,\varphi$,
la orientación y la escala del lector, además de $\alpha$.
La igualdad dimensional conserva su unidad de acción.

Fijada una sección positiva, se escribe $\hbar=\hbar_s$.
La elipse tiene semiejes $\sqrt{2\hbar}e^{\zeta/2}$ y
$\sqrt{2\hbar}e^{-\zeta/2}$, donde
$\zeta=\operatorname{artanh}(C^*/A)$. Su área orientada cumple
\begin{equation}
 \mathcal A(\theta)=\hbar\theta,\qquad
 \mathcal A(2\pi)=h.
 \label{md:aux:area-fase}
\end{equation}
En la realización LC, con reloj $\omega$ e impedancia de
referencia declarados,
\begin{equation}
 \mathcal H_\zeta(Q,P)
 =\frac\omega2\left(e^{-\zeta}Q^2+e^\zeta P^2\right).
 \label{md:aux:hamiltoniano-lc}
\end{equation}
Sobre la órbita se obtiene $\mathcal H_\zeta=\hbar\omega=hf$,
con $f=\omega/(2\pi)$. La convención hamiltoniana da
$\dot\theta=-\omega$ y, por ello,
$d\mathcal A/dt=-\hbar\omega$. La energía positiva coincide
con la magnitud de esa velocidad areal en la realización y
sección fija consideradas. El área acumulada y la acción
lagrangiana de una trayectoria arbitraria son objetos distintos.

La ecuación interna determina la escala de acción por unidad
de fase; el reloj fija su ritmo. La deformación de la elipse
y el orden del transporte conservan información adicional a
$hf$. La preparación gaussiana empleada después es distinta
de un punto de la órbita clásica: su energía inicial es
$hf/2$, en lugar de $hf$.

\subsection{Preparación y lectores del resultado energético}
\label{md:aux:hipotesis-gaussianas}

Se preparan dos modos gaussianos puros, centrados e idénticos,
con covarianza
\begin{equation}
 V_0=\frac\hbar2MM^{\mathsf T},\qquad
 M=\operatorname{diag}(e^{\zeta/2},e^{-\zeta/2}).
 \label{md:aux:preparacion}
\end{equation}
Actúa $U_H$ de \eqref{md:aux:acoplamiento-uh}, transportado a
la coordenatización inicial mediante $M\oplus M$. El estado conjunto
permanece puro. Los lectores energéticos finales tienen una
misma frecuencia $\omega$, conservan la métrica inicial
$G_\zeta=M^{-\mathsf T}M^{-1}$ y no incluyen energía de
interacción residual entre los modos. Las operaciones de
extracción comienzan y terminan con esos mismos lectores.
La accesibilidad física de las operaciones admitidas es una
hipótesis de la cota de trabajo, distinta de la construcción
de un dispositivo que las realice.

Se definen
\begin{equation}
 \begin{gathered}
 \tau=\operatorname{tr}(HH^{\mathsf T})-2\geq0,\\
 W_v=\frac{8I+HH^{\mathsf T}}9,\qquad
 W_m=\frac{I+8HH^{\mathsf T}}9.
 \end{gathered}
 \label{md:aux:covarianzas-reducidas}
\end{equation}
Las covarianzas son $V_j=(\hbar/2)MW_jM^{\mathsf T}$.
La pureza visible se denota por $\mu=\operatorname{Tr}\rho_v^2$
y su autovalor simpléctico normalizado por $\nu=\mu^{-1}$.
En este apéndice, $\nu$ tiene esta función; la frecuencia
ordinaria se escribe $f$ y la angular $\omega$.

\subsection{Identidad entre energía visible y pureza}
\label{md:aux:energia-pureza}

\begin{proposition}
\label{md:aux:prop-energia-pureza}
Con la preparación y los lectores anteriores, la energía
visible por encima del estado inicial satisface
\begin{equation}
 \frac{\Delta E_v}{hf}=\frac\tau{36},\qquad
 \nu^2=1+\frac{8\tau}{81},\qquad
 \boxed{\mu^{-2}-1=\frac{32}{9}\frac{\Delta E_v}{hf}.}
 \label{md:aux:identidad-energia-pureza}
\end{equation}
\end{proposition}
\begin{proof}
El estado es centrado y el lector cuadrático permanece fijo.
Por tanto,
\[
 E_v=\frac\omega2\operatorname{tr}(G_\zeta V_v)
 =\frac{\hbar\omega}{4}\operatorname{tr}W_v.
\]
Como $\operatorname{tr}W_v=2+\tau/9$, al restar
$\hbar\omega/2$ se obtiene la primera igualdad.
La covarianza reducida satisface
$\det W_v=1+8\tau/81$ y $\mu=(\det W_v)^{-1/2}$.
La eliminación de $\tau$ prueba la identidad encuadrada.
La demostración no requiere simetría de $H$ ni coincidencia
entre sus autovalores y sus valores singulares.
\end{proof}

La ley elimina los parámetros particulares del recorrido.
El coeficiente $32/9$ procede del reparto del acoplamiento
y la escala $h$ se conserva hasta la normalización.
Para comprobar la especificidad de ese reparto, considérese
$W_p=pI+(1-p)HH^{\mathsf T}$, con $0<p<1$. Se obtiene
\begin{equation}
 \mu^{-2}-1=4p\frac{\Delta E}{hf}.
 \label{md:aux:reparto-general}
\end{equation}
La elección declarada es $p=8/9$. Cuando $\Delta E>0$,
un contraste puede estimar
\begin{equation}
 p_{\mathrm{obs}}=
 \frac{\mu^{-2}-1}{4\Delta E/(hf)}
 \label{md:aux:reparto-observado}
\end{equation}
y compararlo con $8/9$ sin ajustar $p$ a esos mismos datos.
Programar un mezclador con ese reparto verifica una
emulación; la atribución del reparto a un mecanismo natural
conserva su prueba independiente.

La identidad utiliza gaussianas centradas y la preparación
especificada. Un desplazamiento coherente añade energía sin
cambiar la pureza. El ruido térmico inicial y los cambios de
métrica alteran la relación. Estas variaciones delimitan el
enunciado y proporcionan falsadores de una extensión a estados
arbitrarios.

La combinación convexa $W_v$ es una combinación de matrices
de covarianza producida al reducir un estado gaussiano conjunto.
Una mezcla probabilística de dos estados puros constituye
otro objeto, generalmente no gaussiano; su pureza no viene
dada por el determinante de esa covarianza.

\subsection{Trabajo recuperable bajo operaciones locales}
\label{md:aux:ergotropia}

La ergotropía es el máximo descenso de energía mediante
operaciones unitarias cíclicas con lector energético fijo.
La definición no incluye automáticamente el coste neto de
construcción, preparación o control del aparato.

El espectro reducido tiene poblaciones
\begin{equation}
 \lambda_j=(1-r)r^j,\qquad r=\frac{\nu-1}{\nu+1}.
 \label{md:aux:espectro-gaussiano}
\end{equation}
Están ordenadas decrecientemente con la energía del oscilador.
El estado pasivo tiene así energía $E_{v,\mathrm{pas}}=hf\nu/2$.
Lo alcanza una operación gaussiana:
\begin{equation}
 S=\sqrt\nu\,W_v^{-1/2},\qquad
 \det S=1,\qquad SW_vS^{\mathsf T}=\nu I.
 \label{md:aux:pasivacion}
\end{equation}
El orden de las poblaciones demuestra también la minimalidad
entre todas las unitarias, además de las gaussianas.

Por consiguiente,
\begin{equation}
 \boxed{
 \frac{\mathcal W_v}{hf}
 =\frac{(\nu-1)(9\nu-7)}{32},\qquad
 \frac{E_{v,\mathrm{pas}}-hf/2}{hf}=\frac{\nu-1}{2}.}
 \label{md:aux:trabajo-visible}
\end{equation}
Las dos cantidades suman
$\Delta E_v/(hf)=9(\nu^2-1)/32$. El término pasivo es
energía inaccesible bajo las operaciones locales declaradas;
su presencia no equivale a un flujo de calor realizado.

\subsubsection{Acceso conjunto a las correlaciones}
\label{md:aux:acceso-conjunto}

El segundo modo cumple $\det W_m=\det W_v=\nu^2$ y
$\Delta E_m=8\Delta E_v$. Por ello,
\begin{equation}
 \Delta E_{\mathrm{tot}}=9\Delta E_v=hf\frac\tau4.
 \label{md:aux:energia-conjunta}
\end{equation}
La inversión conjunta del acoplamiento restituye el producto
inicial. El estado total es puro y ese producto es el estado
fundamental del lector conjunto; el acceso conjunto permite
recuperar idealmente $\Delta E_{\mathrm{tot}}$.
Las operaciones locales independientes dejan dos estados
pasivos, cada uno de energía $hf\nu/2$. Se deduce
\begin{equation}
 \boxed{\mathcal W_{\mathrm{conjunta}}-\mathcal W_{\mathrm{local}}
 =hf(\nu-1)=hf(\mu^{-1}-1).}
 \label{md:aux:ventaja-conjunta}
\end{equation}

La identidad cuantifica el trabajo adicional accesible
mediante las correlaciones. La preparación inicial requirió
energía y el balance de un ciclo completo incluye ese aporte.
La relación compara lecturas energéticas y operaciones sobre
el estado, sin identificar información y energía como una
misma magnitud física.

En este enunciado, ``local'' significa unitarias de producto
$U_v\otimes U_m$; quedan fuera de esa clase las mediciones,
la comunicación clásica y la realimentación. El incremento
procede de la restricción operativa sobre el estado
correlacionado y de sus espectros marginales, no de una
energía de interacción residual. El reparto $1:8$ concierne
a las energías por encima del vacío, no a las energías
totales ni a las ergotropías.

\subsection{Casos exactos y orientación del transporte}
\label{md:aux:casos-energia}

Para $H_n$ de \eqref{md:aux:hn},
\begin{equation}
 \tau_n=n^2(2n^2-2n+5).
 \label{md:aux:tau-hn}
\end{equation}
En $n=1$ se obtiene
\begin{equation}
 \begin{gathered}
 \nu=11/9,\qquad\mu=9/11,\qquad
 \Delta E_v=\frac5{36}hf,\\
 \mathcal W_v=\frac1{36}hf,\qquad
 E_{v,\mathrm{pas}}-hf/2=\frac19hf.
 \end{gathered}
 \label{md:aux:caso-positivo}
\end{equation}
El conjunto tiene $\Delta E_{\mathrm{tot}}=5hf/4$, trabajo
local $37hf/36$ y ventaja conjunta $2hf/9$.
El caso $n=-1$ comparte los autovalores de $H_1$,
pero tiene $\tau=9$ y
\begin{equation}
 \nu=\frac{\sqrt{17}}3,\qquad
 \Delta E_v=\frac{hf}{4},\qquad
 \mathcal W_v=\frac{9-2\sqrt{17}}{12}hf.
 \label{md:aux:caso-negativo}
\end{equation}
La energía visible añadida difiere en razón $9/5$ a
frecuencia y sección de acción iguales. El cambio $n\mapsto-n$
no es el retroceso completo: $H_{-n}\ne H_n^{-1}$ en general.
La inversión matricial verdadera conserva $\tau$ en
dimensión dos. La asimetría paramétrica de esta colección
conserva así su distinción respecto de una reversión completa.

Para $C_{\mathrm{aux}}$,
$\tau_C=4\sin^2a\sinh^2(2b)$. La composición energética da
\begin{equation}
 \boxed{\Delta E_v(C_{\mathrm{aux}})
 =\frac{hf}{9}\sin^2a\sinh^2(2b),}
 \label{md:aux:energia-compensado}
\end{equation}
\begin{equation}
 \mathcal W_{\mathrm{conjunta}}-\mathcal W_{\mathrm{local}}
 =hf\left(\sqrt{1+\frac{32}{81}\sin^2a\sinh^2(2b)}-1\right).
 \label{md:aux:trabajo-compensado}
\end{equation}
La conexión con Planck permanece explícita: en la sección
previa, $hf=2\pi s_0H_5(\alpha,\varphi)f$; en la de retorno,
$H_5$ se sustituye por $\eta_{\mathrm{ret}}$. La composición
mantiene fijada la sección utilizada; no supone una
variación experimental de $\alpha$ ni una alternancia libre
entre secciones dimensionales.

\subsection{Boltzmann y temperatura equivalente del estado pasivo}
\label{md:aux:temperatura}

Para $\nu>1$, el estado pasivo tiene temperatura equivalente
\begin{equation}
 k_BT_{\mathrm{pas}}
 =\frac{hf}{\log[(\nu+1)/(\nu-1)]}.
 \label{md:aux:temperatura-pasiva}
\end{equation}
Cuando se utiliza el reloj del transductor,
$hf=k_B\Theta_{\mathrm{clk}}\log3$, resulta
\begin{equation}
 \frac{T_{\mathrm{pas}}}{\Theta_{\mathrm{clk}}}
 =\frac{\log3}{\log[(\nu+1)/(\nu-1)]}.
 \label{md:aux:razon-temperaturas}
\end{equation}
Para $n=1$ se recupera $\log3/\log10$. La entropía es
la del espectro gaussiano, con ocupación $(\nu-1)/2$.
La composición energética conserva esta lectura térmica
previamente construida. La temperatura equivalente del
estado pasivo no implica que el estado reducido previo
estuviese en equilibrio ni que se haya transferido calor.

\subsection{Especificidad y condiciones de contraste}
\label{md:aux:contraste-trabajo}

La relación general entre coherencia, correlaciones y
extracción de trabajo pertenece a la literatura existente.
La composición examinada especifica una ley conjunta:
coeficiente $32/9$, reparto energético $1:8$, colección
orientada $H_n$, escala de acción y diferencia entre
acceso local y conjunto.

Se distinguen dos niveles experimentales. La emulación
realiza $U_H$ y comprueba sus identidades; verifica esa
realización del modelo. La predicción física exige
determinar prospectivamente, desde la construcción, el
dispositivo, los grados de libertad, la frecuencia, la
preparación y el acoplamiento que realizan sus objetos,
sin ajustar $8/9$ a los datos del contraste. Se miden
por separado energía, covarianza, pureza, costes de
control y energía de restauración.

Un contraste independiente obtiene temperatura y pureza
por lectores independientes de la ley que pretende
comprobar. El trabajo recuperado se compara con un
balance que incluya preparación, controles y restauración
de la memoria.

Las comparaciones externas incluyen estudios de
ergotropía gaussiana y pasivación,\footnote{Kua, Serafini
y Genoni, \emph{Daemonic ergotropy of Gaussian quantum
states and the role of measurement-induced purification
via general-dyne detection},
\href{https://arxiv.org/abs/2506.22288}{arXiv:2506.22288}.
Su papel es la comparación con fórmulas generales de
energía, pureza y pasividad; no valida el acoplamiento HMT
ni sus constantes.}
la extracción de ergotropía coherente con costes
energéticos explícitos,\footnote{\emph{Experimental
Extraction of Coherent Ergotropy and Its Energetic Cost
in a Superconducting Qubit},
\href{https://arxiv.org/abs/2506.16881}{arXiv:2506.16881}.}
y el control coherente entre ciclos de máquinas
cuánticas.\footnote{Hou et al., \emph{Combining energy
efficiency and quantum advantage in cyclic machines},
\href{https://doi.org/10.1038/s41467-025-60179-5}{doi:10.1038/s41467-025-60179-5}.
El experimento citado no mide el transporte nonádico
compuesto en este apéndice.}

El resultado es una composición algebraica con
estructura de partida explícita. Sus pruebas fijan
identidades y restricciones operativas. Una ventaja
experimental observada, una nueva constante universal,
una violación de la termodinámica o un mecanismo material
de transporte sin calor requerirían sus respectivas
demostraciones y contrastes; no son conclusiones de las
identidades auxiliares aquí establecidas.

\section{Enumeración exhaustiva y determinación finita de los operadores de transporte}
\label{iii:nucleo-finito}

Este apéndice desarrolla dos operaciones finitas con datos de partida
diferentes. La primera enumera las condiciones iniciales de los dos cursores
y construye sus emisiones mediante la recurrencia aritmética declarada.
La segunda determina los operadores de transporte a partir del registro
regional de transiciones. La separación de sus dominios permite precisar
qué produce cada demostración: la enumeración construye el catálogo de
emisiones; la reconstrucción matricial determina los operadores compatibles
con el registro. El teorema de reconstrucción no se emplea como demostración
de la producción anterior de ese registro.

Las primitivas de la primera operación son el soporte $I_9^2$, las cuatro
orientaciones cardinales, el selector trítico, la lectura anterior al
desplazamiento, las inversiones de orientación y la regla de emisión.
Sus definiciones se reúnen a continuación. Las de la segunda son las
cuatro matrices de transiciones de~\eqref{iii:nf-registro}.
Su reconstrucción única se demuestra en el
teorema~\ref{iii:nf-unicidad}; el calendario resultante se expresa
en~\eqref{iii:nf-calendario}. Estas determinaciones conservan
los datos de partida declarados para cada operación.

\subsection{Dominio completo de condiciones iniciales}

Sea $I_9=\{0,\ldots,8\}$ y sea
\[
 \mathcal D=\{N,E,S,O\},\qquad
 v_N=(-1,0),\quad v_E=(0,1),\quad
 v_S=(1,0),\quad v_O=(0,-1).
\]
Una condición inicial de cursor es $s=(i,j,d)\in
\mathcal S=I_9^2\times\mathcal D$. Los dos cursores son independientes
en el dominio inicial
\begin{equation}
 \Omega=\mathcal S_+\times\mathcal S_\times,
 \qquad |\mathcal S|=9^2\cdot4=324,
 \qquad |\Omega|=324^2=104\,976.
 \label{iii:nf-dominio}
\end{equation}
El índice $i$ corresponde a la etiqueta positiva $i+1$ de APP. Para
enteros positivos, la reducción positiva es $\rho_9(n)=1+[(n-1)]_9$,
donde $[a]_b$ denota el representante en $\{0,\ldots,b-1\}$.
Las evaluaciones residuales de las dos hojas son
\[
 \sigma(i,j)=\rho_9(i+j+2),\qquad
 \varpi(i,j)=\rho_9((i+1)(j+1)).
\]
El lector multiplicativo finito utiliza
\begin{equation}
 \begin{array}{c|rrrrrrrrr}
 a&1&2&3&4&5&6&7&8&9\\ \hline
 \ell_9(a)&0&1&0&2&5&0&4&3&0
 \end{array}.
 \label{iii:nf-log}
\end{equation}
Sobre las unidades, $\ell_9$ es el exponente discreto en base $2$.
Sobre $3,6,9$, el valor $0$ constituye la extensión declarada del
observable de emisión; la pertenencia a ese sector permanece en el
registro de trayectoria.

Para $1\le t\le54$, defínase
\[
 r_t=1+[(t-1)]_9,
 \qquad
 \tau_t=
 \begin{cases}
 +1,&r_t\in\{1,4,7\},\\
 -1,&r_t\in\{2,5,8\},\\
 0,&r_t\in\{3,6,9\}.
 \end{cases}
\]
Cada cursor conserva su posición $z_t$ inmediatamente antes del paso
$t$ y su vector de orientación $v_t$. El cursor aditivo está activo
cuando $\tau_t=+1$; el multiplicativo, cuando $\tau_t=-1$. Para cada
uno de ellos, con indicador de actividad $a_t$, la actualización es
\begin{equation}
 z_{t+1}=[z_t+a_tv_t]_9,
 \qquad
 v_{t+1}=
 \begin{cases}
 -v_t,&t\in\{27,54\},\\
 v_t,&t\notin\{27,54\}.
 \end{cases}
 \label{iii:nf-actualizacion}
\end{equation}
La lectura del paso $t$ se realiza en $z_t$, antes del desplazamiento.
El evento neutral conserva ambas posiciones. Esta convención fija
inequívocamente los extremos de las ventanas y el efecto de las
inversiones cardinales.

En las ventanas $W_m=\{9m-8,\ldots,9m\}$, $1\le m\le6$, sean
\begin{align}
 S_m(s_+)&=\sum_{\substack{t\in W_m\\\tau_t=+1}}
                 \sigma(z_t^+),&
 E_m(s_\times)&=\sum_{\substack{t\in W_m\\\tau_t=-1}}
                 \ell_9(\varpi(z_t^\times)),
 \label{iii:nf-firmas}\\
 s_m&=[S_m]_{10},&e_m&=[E_m]_{10}.\notag
\end{align}
En cada ventana hay tres eventos de cada tipo. Las firmas $s$ y $e$
son las seis componentes de estos residuos, no los enteros $S_m,E_m$
anteriores a su reducción.

\subsection{Censo de firmas y de emisiones}

La emisión de una pareja de condiciones iniciales queda determinada por
\begin{equation}
 U_m=100s_m+10[s_m+e_m]_{10}+[3s_m+5e_m+1]_{10},
 \qquad 1\le m\le6.
 \label{iii:nf-emision}
\end{equation}
El término final $1$ es $[7\cdot3]_{10}$, correspondiente a los tres
eventos neutrales. Los coeficientes de esta regla forman parte de la
recurrencia declarada. Las constantes reales que se estudian posteriormente
no intervienen en la enumeración de su dominio ni en su evaluación.

\begin{proposition}[Factorización inyectiva del catálogo finito]
\label{iii:nf-inyectividad}
La emisión distingue cada pareja de firmas $(s,e)$. Hay $18$ firmas
aditivas y $26$ multiplicativas; sus pares producen $468$ emisiones.
Las multiplicidades de estas emisiones son $144$, $432$ y $1944$,
con $432$, $18$ y $18$ emisiones respectivamente.
\end{proposition}

\begin{proof}
La cifra de centenas de $U_m$ recupera $s_m$. Si $b_m$ es su cifra de
decenas, entonces $e_m=[b_m-s_m]_{10}$. La cifra de unidades verifica
la tercera congruencia de \eqref{iii:nf-emision}. Por tanto, el
acoplamiento es inyectivo.

El censo se obtiene recorriendo, en el orden cartesiano completo,
$i=0,\ldots,8$, $j=0,\ldots,8$ y $d=N,E,S,O$, y aplicando los $54$
pasos de \eqref{iii:nf-actualizacion}. Las $324$ condiciones aditivas
se distribuyen en las siguientes $18$ firmas, cada una con $18$
preimágenes. En las tablas, una sucesión de seis cifras representa
las seis componentes ordenadas de una firma.
\begin{center}
\begin{tabular}{ccc}
$122564$&$122898$&$212645$\\
$212988$&$221456$&$221889$\\
$456221$&$465898$&$546988$\\
$564122$&$645212$&$654889$\\
$889221$&$889654$&$898122$\\
$898465$&$988212$&$988546$
\end{tabular}
\end{center}
Para el cursor multiplicativo, las firmas y sus multiplicidades son
\begin{center}
\begin{tabular}{rr@{\qquad}rr}
Firma&Multiplicidad&Firma&Multiplicidad\\ \hline
$000000$&$108$&$555555$&$24$\\
$177393$&$8$&$555717$&$8$\\
$177555$&$8$&$555771$&$8$\\
$177771$&$8$&$555933$&$8$\\
$339555$&$8$&$717555$&$8$\\
$339771$&$8$&$717717$&$8$\\
$339933$&$8$&$717933$&$8$\\
$393177$&$8$&$771177$&$8$\\
$393393$&$8$&$771339$&$8$\\
$393555$&$8$&$771555$&$8$\\
$555177$&$8$&$933339$&$8$\\
$555339$&$8$&$933555$&$8$\\
$555393$&$8$&$933717$&$8$
\end{tabular}
\end{center}
Estas tablas resultan de las sumas finitas \eqref{iii:nf-firmas}, con
dominio, orden de lectura e inversiones ya especificados. Sus sumas
de multiplicidades son $18\cdot18=324$ y
$24\cdot8+24+108=324$. El dominio cartesiano
\eqref{iii:nf-dominio} permite realizar cualquier pareja de firmas.
La inyectividad proporciona $18\cdot26=468$ emisiones y las
multiplicidades
\[
 (18\cdot24,\ 18\cdot8)=(432,144),\qquad
 (18,\ 18\cdot24)=(18,432),\qquad
 (18,\ 18\cdot108)=(18,1944),
\]
donde cada par indica número de emisiones y multiplicidad de su fibra.
Finalmente,
\[
 432\cdot144+18\cdot432+18\cdot1944=104\,976.
\]
La enumeración agota así el dominio completo de condiciones iniciales.
\end{proof}

La reducción ternaria orientada es
\[
 \Psi(U)=([-U_1]_3,\ldots,[-U_6]_3).
\]
Aplicada a las $18\cdot26$ emisiones de las tablas anteriores,
produce el siguiente censo de fibras. Aquí $r$ cuenta emisiones
distintas, no condiciones iniciales.
\begin{equation}
 \begin{array}{c|rrrrrr}
 r&1&2&3&4&5&6\\ \hline
 \#\{w:\#\Psi^{-1}(w)=r\}&132&66&18&3&6&18
 \end{array}.
 \label{iii:nf-censo-ternario}
\end{equation}
La suma de la segunda fila es $243$, mientras que su suma ponderada
por $r$ es $468$. Se distingue, por ello, la imagen de $243$ palabras
del ambiente $\mathbb F_3^6$, que contiene $729$ palabras. El certificado
adjunto conserva todas las firmas con sus preimágenes, las $468$
emisiones y su reducción ternaria; las fórmulas anteriores proporcionan
un procedimiento completo para reconstruirlas.

\subsection{Registro de transiciones y reconstrucción de los levantamientos}

En este apartado se fija el registro regional del propietario citado.
Sean $b_0^\chi,\ldots,b_4^\chi\in\mathbb F_3^6$, con
$\chi\in\{\mathsf C,\mathsf P,\mathsf A\}$, sus tres sucesiones.
El registro organiza las transiciones $b_0\to b_1\to b_2$ en
$X_0,Y_0$ y las transiciones $b_2\to b_3\to b_4$ en $X_1,Y_1$,
con vectores fila y el mismo orden regional:
\begingroup\small
\begin{equation}
\begin{aligned}
X_0&=\begin{pmatrix}
0&1&0&2&1&1\\0&1&2&2&2&2\\2&0&1&1&0&1\\
1&2&1&2&2&1\\1&2&1&2&0&0\\1&1&2&2&0&2
\end{pmatrix},&
Y_0&=\begin{pmatrix}
0&1&2&2&2&2\\0&1&0&2&1&1\\1&2&1&2&2&1\\
1&0&2&0&1&1\\1&1&2&2&0&2\\1&2&1&0&2&0
\end{pmatrix},\\[1ex]
X_1&=\begin{pmatrix}
0&1&0&2&1&1\\0&0&2&1&1&1\\1&0&2&0&1&1\\
0&1&2&2&2&2\\1&2&1&0&2&0\\0&1&0&2&1&0
\end{pmatrix},&
Y_1&=\begin{pmatrix}
0&0&2&1&1&1\\1&1&0&2&2&1\\0&1&2&2&2&2\\
1&0&2&0&1&1\\0&1&0&2&1&0\\0&1&0&2&0&0
\end{pmatrix}.
\end{aligned}
\label{iii:nf-registro}
\end{equation}
\endgroup
La producción precedente se expresa en su propietario mediante la
composición de selección, transporte y actualización:
\[
 x_{j+1}^{\chi}
 =\mathcal U_{9j+9}\circ\cdots\circ\mathcal U_{9j+1}(x_j^{\chi}),
 \qquad b_j^\chi=\Pi_6(x_j^\chi),
 \qquad\mathcal U_t=\mathsf{Upd}_t\circ\mathsf{Tra}_t\circ\mathsf{Sel}_t.
\]
El resultado siguiente parte del registro \eqref{iii:nf-registro}.
Su demostración permite reconstruir los operadores y verificar la
compatibilidad de las transiciones; no atribuye a la inversión matricial
la producción coordenada de los estados $x_j^\chi$.

\begin{proposition}[Unicidad sobre el registro regional]
\label{iii:nf-unicidad}
Los sistemas $X_aL_a=Y_a$, $a=0,1$, tienen una única solución sobre
$\mathbb F_3$. Las dos soluciones son
\begingroup\small
\begin{equation}
L_0=\begin{pmatrix}
2&2&2&1&2&1\\2&1&2&2&1&1\\1&0&0&1&0&2\\
0&0&1&1&1&0\\2&2&2&0&2&1\\2&1&2&1&0&0
\end{pmatrix},\qquad
L_1=\begin{pmatrix}
2&1&2&0&2&2\\1&2&1&1&2&0\\1&2&1&1&1&2\\
0&1&0&0&2&1\\2&0&2&1&0&1\\0&2&2&2&1&1
\end{pmatrix}.
\label{iii:nf-lifts}
\end{equation}
\endgroup
\end{proposition}

\begin{proof}
La eliminación en $\mathbb F_3$ da $\det X_0=1$ y
$\det X_1=2=-1$. Cada aplicación $L\mapsto X_aL$ es, por tanto,
un automorfismo del espacio de matrices de dimensión $36$; en una
vectorización por columnas está representada por $I_6\otimes X_a$.
Se obtiene $L_a=X_a^{-1}Y_a$. La multiplicación de las matrices de
\eqref{iii:nf-lifts} por las respectivas matrices $X_a$ reproduce las
dos matrices $Y_a$ de \eqref{iii:nf-registro}; existencia y unicidad
quedan establecidas sobre ese registro.
\end{proof}

Las tres sucesiones resultantes, con separación explícita en bloques,
son
\begin{align*}
 B^{\mathsf C}&=010211\mid012222\mid010211\mid002111\mid110221,\\
 B^{\mathsf P}&=201101\mid121221\mid102011\mid012222\mid102011,\\
 B^{\mathsf A}&=121200\mid112202\mid121020\mid010210\mid010200.
\end{align*}
Cada bloque conserva sus seis posiciones. La concatenación no se
identifica con un entero en base diez.

\subsection{Exhaustividad del calendario y estabilidad de la reconstrucción}

Un calendario binario es $c=(c_0,c_1,c_2,c_3)\in\{0,1\}^4$.
Defínase el número de transiciones incompatibles con el registro por
\begin{equation}
 d(c)=\sum_{\chi\in\{\mathsf C,\mathsf P,\mathsf A\}}
       \sum_{j=0}^{3}
       \mathbf1\{b_j^\chi L_{c_j}\ne b_{j+1}^\chi\}.
 \label{iii:nf-calendario}
\end{equation}
Las multiplicaciones de las sucesiones anteriores por
\eqref{iii:nf-lifts} dan
\[
 d(c)=3\,d_{\rm H}(c,(0,0,1,1)),
\]
donde $d_{\rm H}$ es la distancia de Hamming: en cada una de las
cuatro posiciones, el operador alternativo incumple las tres
transiciones regionales. El censo completo es
\begin{center}
\begin{tabular}{cr@{\qquad}cr@{\qquad}cr@{\qquad}cr}
$c$&$d(c)$&$c$&$d(c)$&$c$&$d(c)$&$c$&$d(c)$\\ \hline
$0000$&$6$&$0100$&$9$&$1000$&$9$&$1100$&$12$\\
$0001$&$3$&$0101$&$6$&$1001$&$6$&$1101$&$9$\\
$0010$&$3$&$0110$&$6$&$1010$&$6$&$1110$&$9$\\
$0011$&$0$&$0111$&$3$&$1011$&$3$&$1111$&$6$
\end{tabular}
\end{center}
En consecuencia, $0011$ es el único calendario que reproduce las
tres sucesiones completas.

La sensibilidad a las entradas de las matrices también admite una
prueba exhaustiva sin aproximaciones. Una modificación unitaria
sustituye $L_a$ por $L_a+\lambda E_{ij}$, con
$a\in\{0,1\}$, $1\le i,j\le6$ y
$\lambda\in\mathbb F_3\setminus\{0\}$. Hay
$2\cdot6\cdot6\cdot2=144$ modificaciones de este tipo. Como
$X_a$ es invertible,
\[
 X_a(L_a+\lambda E_{ij})-Y_a
 =\lambda X_aE_{ij}\ne0.
\]
Toda modificación destruye al menos una de las transiciones del
registro correspondiente. El certificado ejecutable enumera además
las $144$ modificaciones y conserva, para cada una, el número de
transiciones incumplidas.

La propiedad demostrada es la rigidez de los operadores respecto del
registro regional fijado. Para reproducir una construcción anterior
del propio registro se necesita la acción coordenada de sus operadores
de selección, transporte y actualización. Esa dependencia se conserva
separada del teorema de rigidez, del catálogo finito de emisiones y
de las realizaciones analíticas posteriores.

\subsection{Certificado de reproducción y residencia de los datos}

El programa \texttt{verificar\_nucleo\_finito.py} recorre las $324$
condiciones de cada cursor y forma el producto de las clases resultantes.
El catálogo documental se carga después de esta construcción para
comparar las emisiones y sus multiplicidades. En la segunda parte,
las cuatro matrices de transiciones se extraen de la copia literal
del propietario; las soluciones se calculan por eliminación exacta
módulo $3$ antes de compararlas con las matrices impresas. El recibo
\texttt{NUCLEO\_FINITO.json} conserva los dominios, las preimágenes,
las emisiones, los $16$ calendarios y las $144$ modificaciones.

La comprobación utiliza enteros y residuos exactos. No recibe valores
metrológicos ni expansiones reales de las constantes. Su alcance
es el de las operaciones finitas y los datos de partida especificados
en este apéndice. Las demostraciones de compatibilidad a profundidad
ilimitada y las identificaciones físicas conservan sus respectivos
dominios y obligaciones expositivas.

\section{Fuentes, conservación y reproducción}
El suplemento digital conserva el documento acumulativo íntegro, sus
treinta y nueve fuentes, los programas de comprobación y los recibos
correspondientes. La narrativa aprobada y su continuación se mantienen
literalmente en ese archivo. La presente edición transforma su
enunciación conversacional en exposición académica y registra la
correspondencia entre secciones, sin sustituir los originales.

Los antecedentes proceden de las fuentes de la serie del 21 de
septiembre de 2026. El núcleo formal se incorpora desde su clausura
documental en el artículo constitutivo; los desarrollos de acción,
incidencia, masa y temperatura se conservan con sus demostraciones.
El manifiesto indica cada archivo, intervalo extraído y huella. Estos
datos documentan procedencia y conservación, mientras las pruebas del
cuerpo fijan el alcance matemático de los resultados.

Los programas de composición de constantes, hibridación y
compatibilidad contrastan identidades exactas y ejemplos finitos.
Los diagnósticos numéricos se distinguen de los controles racionales.
Las pruebas de monotonía, convergencia, inversión y paso al límite
permanecen desarrolladas en el texto: los ejemplos ejecutables
complementan esas demostraciones.

La compilación utiliza LuaLaTeX. Las fuentes incluidas y sus figuras
permiten reproducir el PDF sin depender de los directorios originales
del autor. El archivo de instrucciones del paquete especifica el orden
de compilación y el modo de ejecución de las comprobaciones.

\clearpage\begin{thebibliography}{99}
\bibitem{mustovicary} B. Musto y J. Vicary, ``Quantum Latin Squares and Unitary Error Bases'', \emph{Quantum Information and Computation} \textbf{16}, 1318--1332 (2016). \href{https://doi.org/10.26421/QIC16.15-16-4}{doi:10.26421/QIC16.15-16-4}.

\bibitem{delascuevas2020} G. De las Cuevas, T. Drescher y T. Netzer, ``Quantum Magic Squares: Dilations and Their Limitations'', \emph{Journal of Mathematical Physics} \textbf{61}, 111704 (2020). \href{https://doi.org/10.1063/5.0022344}{doi:10.1063/5.0022344}.

\bibitem{delascuevas2023} G. De las Cuevas, T. Netzer e I. Valentiner-Branth, ``Magic Squares: Latin, Semiclassical, and Quantum'', \emph{Journal of Mathematical Physics} \textbf{64}, 022201 (2023). \href{https://doi.org/10.1063/5.0127393}{doi:10.1063/5.0127393}.

\bibitem{codata2022} P. J. Mohr, D. B. Newell, B. N. Taylor y E. Tiesinga, ``CODATA Recommended Values of the Fundamental Physical Constants: 2022'', \emph{Reviews of Modern Physics} \textbf{97}, 025002 (2025). \href{https://doi.org/10.1103/RevModPhys.97.025002}{doi:10.1103/RevModPhys.97.025002}. \href{https://physics.nist.gov/cuu/Constants/Table/allascii.txt}{Tabla de constantes}.

\bibitem{flm} I. B. Frenkel, J. Lepowsky y A. Meurman, \emph{Vertex Operator Algebras and the Monster}, Pure and Applied Mathematics, vol. 134, Academic Press, 1988.

\bibitem{borcherds} R. E. Borcherds, ``Monstrous Moonshine and Monstrous Lie Superalgebras'', \emph{Inventiones Mathematicae} \textbf{109}, 405--444 (1992). \href{https://doi.org/10.1007/BF01232032}{doi:10.1007/BF01232032}.

\bibitem{conwaysloane} J. H. Conway y N. J. A. Sloane, \emph{Sphere Packings, Lattices and Groups}, 3.\textsuperscript{a} ed., Grundlehren der mathematischen Wissenschaften 290, Springer, 1999. \href{https://doi.org/10.1007/978-1-4757-6568-7}{doi:10.1007/978-1-4757-6568-7}.

\bibitem{bipm2018} Conférence générale des poids et mesures, \emph{Resolution 1 of the 26th CGPM: On the revision of the International System of Units}, 2018. \href{https://www.bipm.org/en/committees/cg/cgpm/26-2018/resolution-1}{Resolución oficial}. \href{https://doi.org/10.59161/CGPM2018RES1E}{doi:10.59161/CGPM2018RES1E}.

\bibitem{Bose1924} Satyendra~Nath Bose.
\newblock Plancks gesetz und lichtquantenhypothese.
\newblock {\em Zeitschrift für Physik}, 26:178--181, 1924.

\bibitem{Dirac1926} Paul A.~M. Dirac.
\newblock On the theory of quantum mechanics.
\newblock {\em Proceedings of the Royal Society A}, 112:661--677, 1926.

\bibitem{Einstein1924} Albert Einstein.
\newblock Quantentheorie des einatomigen idealen gases.
\newblock {\em Sitzungsberichte der Preussischen Akademie der Wissenschaften},
  pages 261--267, 1924.

\bibitem{Fermi1926} Enrico Fermi.
\newblock Sulla quantizzazione del gas perfetto monoatomico.
\newblock {\em Rendiconti Lincei}, 3:145--149, 1926.

\bibitem{Pauli1925} Wolfgang Pauli.
\newblock Über den zusammenhang des abschlusses der elektronengruppen im atom
  mit der komplexstruktur der spektren.
\newblock {\em Zeitschrift für Physik}, 31:765--783, 1925.

\bibitem{Pathria2011} R.~K. Pathria and Paul~D. Beale.
\newblock {\em Statistical Mechanics}.
\newblock Elsevier, 3 edition, 2011.

\bibitem{bennett2003} C. H. Bennett, ``Notes on Landauer's principle,
reversible computation, and Maxwell's Demon'', \emph{Studies in History
and Philosophy of Modern Physics} \textbf{34}, 501--510 (2003).
\href{https://doi.org/10.1016/S1355-2198(03)00039-X}{doi:10.1016/S1355-2198(03)00039-X}.
\end{thebibliography}

\end{document}
